% !TeX program = xelatex
% ================================================================
%  《收获与播种》(Récoltes et Semailles) 中文译本
%  卷五 · 第四部 —— 葬礼(III): 或四种运算
%  Quatrième partie: L'Enterrement (III) ou les Quatre Opérations
%  原书 p.633 部扉起, 收 18.3-18.5.4 注(c1)“记忆的失灵——或新编历史”
%  第四段, 并续至全书末行(原书 p.929)
%  --------------------------------------------------------------
%  注: 本卷正文起始于 part4-ch18.tex(原 part2-ch18.tex 第 2190 行起,
%      即第四部部扉所在处), 承接第 18 章 18.3。
%  五卷本之一。共用导言区 common.tex, 封面 cover-art.tex。
%  构建: xelatex ×4 (目录页码需 4 遍才收敛)
% ================================================================
\PassOptionsToPackage{AutoFallBack}{xeCJK}
\documentclass[
    PaperSize=a4paper,
    HeadingFont=hei,
    Biblatex=false,
    BookTitleCN={收获与播种},
    Translator={数译},
    LogoFile=logo.pdf
]{mathtranslation}

% ================================================================
%  《收获与播种》中文译本 —— 各卷共用导言区
%  由 vol-prelude / vol-part1 / vol-part2 / vol-part3 / vol-part4
%  以及全书合订 main-all 共同 \input。基类 mathtranslation.cls v3.2 不改。
%  --------------------------------------------------------------
%  样式均在此改写, 以贴近原书:
%    · 章标题仿原书 "5. Travail et découverte"(块状左对齐黑体)
%    · 节/小节: 黑体左对齐, 编号后带句点
%    · 脚注: 人工编号照录原书——前奏与第一部分页内重新编号,
%      第二至四部各自连续编号并带(*)角标; 第一部分另有 \fnoteP
%      处理 "(n)" 式连续注号
%    · 边码: 原打字稿页码 \origpage{N}(前奏作 A1、P1、L1、罗马数字等)
%    · ⋄ 记号: 段首补写标记 \pmark 与上标形式 \dmark
% ================================================================

% 题名含空格, 不能走类选项(会被 kvoptions 吃掉), 在此赋值
\renewcommand{\BookTitleEN}{Récoltes et Semailles}
\renewcommand{\OriginalAuthor}{Alexandre Grothendieck}
\renewcommand{\TranslationDate}{\today}

% --------------------------------------------------------------
%  罕用字后备(如"祐"): Fandol 不载时退到 SimSun
% --------------------------------------------------------------
\setCJKfallbackfamilyfont{\CJKrmdefault}{{SimSun}}

% --------------------------------------------------------------
%  hyperref 元信息(各卷在 % ================================================================
%  《收获与播种》中文译本 —— 各卷共用导言区
%  由 vol-prelude / vol-part1 / vol-part2 / vol-part3 / vol-part4
%  以及全书合订 main-all 共同 \input。基类 mathtranslation.cls v3.2 不改。
%  --------------------------------------------------------------
%  样式均在此改写, 以贴近原书:
%    · 章标题仿原书 "5. Travail et découverte"(块状左对齐黑体)
%    · 节/小节: 黑体左对齐, 编号后带句点
%    · 脚注: 人工编号照录原书——前奏与第一部分页内重新编号,
%      第二至四部各自连续编号并带(*)角标; 第一部分另有 \fnoteP
%      处理 "(n)" 式连续注号
%    · 边码: 原打字稿页码 \origpage{N}(前奏作 A1、P1、L1、罗马数字等)
%    · ⋄ 记号: 段首补写标记 \pmark 与上标形式 \dmark
% ================================================================

% 题名含空格, 不能走类选项(会被 kvoptions 吃掉), 在此赋值
\renewcommand{\BookTitleEN}{Récoltes et Semailles}
\renewcommand{\OriginalAuthor}{Alexandre Grothendieck}
\renewcommand{\TranslationDate}{\today}

% --------------------------------------------------------------
%  罕用字后备(如"祐"): Fandol 不载时退到 SimSun
% --------------------------------------------------------------
\setCJKfallbackfamilyfont{\CJKrmdefault}{{SimSun}}

% --------------------------------------------------------------
%  hyperref 元信息(各卷在 % ================================================================
%  《收获与播种》中文译本 —— 各卷共用导言区
%  由 vol-prelude / vol-part1 / vol-part2 / vol-part3 / vol-part4
%  以及全书合订 main-all 共同 \input。基类 mathtranslation.cls v3.2 不改。
%  --------------------------------------------------------------
%  样式均在此改写, 以贴近原书:
%    · 章标题仿原书 "5. Travail et découverte"(块状左对齐黑体)
%    · 节/小节: 黑体左对齐, 编号后带句点
%    · 脚注: 人工编号照录原书——前奏与第一部分页内重新编号,
%      第二至四部各自连续编号并带(*)角标; 第一部分另有 \fnoteP
%      处理 "(n)" 式连续注号
%    · 边码: 原打字稿页码 \origpage{N}(前奏作 A1、P1、L1、罗马数字等)
%    · ⋄ 记号: 段首补写标记 \pmark 与上标形式 \dmark
% ================================================================

% 题名含空格, 不能走类选项(会被 kvoptions 吃掉), 在此赋值
\renewcommand{\BookTitleEN}{Récoltes et Semailles}
\renewcommand{\OriginalAuthor}{Alexandre Grothendieck}
\renewcommand{\TranslationDate}{\today}

% --------------------------------------------------------------
%  罕用字后备(如"祐"): Fandol 不载时退到 SimSun
% --------------------------------------------------------------
\setCJKfallbackfamilyfont{\CJKrmdefault}{{SimSun}}

% --------------------------------------------------------------
%  hyperref 元信息(各卷在 \input{common} 之后可用 \hypersetup 覆写)
% --------------------------------------------------------------
\hypersetup{
  pdftitle={收获与播种},
  pdfauthor={Alexandre Grothendieck / 数译 译},
  pdfsubject={Récoltes et Semailles — Réflexions et témoignage sur un passé de mathématicien}
}

% --------------------------------------------------------------
%  标题样式: 仿原书 —— 黑体左对齐, 编号后带句点
% --------------------------------------------------------------
\titleformat{\section}[block]
  {\normalfont\Large\headfont}{\thesection.}{0.5em}{}
\titlespacing*{\section}{0pt}{3.2ex plus 1.2ex minus .2ex}{1.6ex plus .2ex}
\titleformat{\subsection}[block]
  {\normalfont\large\headfont}{\thesubsection.}{0.5em}{}
\titlespacing*{\subsection}{0pt}{2.6ex plus 1ex minus .2ex}{1.2ex plus .2ex}
\titleformat{\subsubsection}[block]
  {\normalfont\normalsize\headfont}{\thesubsubsection.}{0.5em}{}
\titlespacing*{\subsubsection}{0pt}{2.2ex plus .8ex minus .2ex}{1ex plus .2ex}
\titleformat{name=\subsubsection,numberless}[block]
  {\normalfont\normalsize\headfont}{}{0em}{}
\setcounter{secnumdepth}{3}
\setcounter{tocdepth}{3}
% 为含长外文书名与行内公式的段落提供必要的断行余量
\setlength{\emergencystretch}{3em}

% 章标题样式: 原书章标题与节同级排印("13. A) HERITAGE..."),
% 故章标题取块状左对齐黑体, 字号介于节与部之间
\titleformat{\chapter}[block]
  {\normalfont\LARGE\headfont}{\thechapter.}{0.5em}{}
\titlespacing*{\chapter}{0pt}{3.6ex plus 1.4ex minus .2ex}{2ex plus .2ex}
% 目录中章号照原书用阿拉伯数字("13. A) ..." 而非 "第十三章")
\renewcommand{\CTEXthechapter}{\thechapter\hspace{.3em}}

% --------------------------------------------------------------
%  目录观感: 条目与引导点之间、引导点与页码之间各留 0.5em
%  内核 \@dottedtocline 让 \leaders 从标题末字起排, 第一个点几乎
%  贴着末字, 读起来像多标了一个句点; 页码前也没有余量。两处各补
%  0.5em。
% --------------------------------------------------------------
\makeatletter
\patchcmd{\@dottedtocline}{\nobreak\leaders}{\nobreak\hspace{0.5em}\leaders}%
  {\typeout{MT: TOC 引导点前已留 0.5em}}%
  {\typeout{MT-WARN: TOC 引导点前间距补丁未生效}}
\patchcmd{\@dottedtocline}{\hfill\nobreak\hb@xt@\@pnumwidth}%
  {\hfill\nobreak\hspace{0.5em}\hb@xt@\@pnumwidth}%
  {\typeout{MT: TOC 页码前已留 0.5em}}%
  {\typeout{MT-WARN: TOC 页码前间距补丁未生效}}
\makeatother

% --------------------------------------------------------------
%  正文半角标点后的间距
%  本书体例用半角 , ; : ? ! 与 ASCII 括号, 由 xeCJK 按 HalfRight
%  (右半留白) 处理。放宽到 0.45em, 与全角标点右半留白大致齐平。
% --------------------------------------------------------------
\xeCJKsetup{CJKecglue={\hskip 0.45em plus 0.06em minus 0.09em}}

% --------------------------------------------------------------
%  页眉页脚: 偶数页眉右对齐斜体节题; 页码页脚外侧
% --------------------------------------------------------------
\fancyhf{}
\fancyhead[RE]{\small\itshape\rightmark}
\fancyfoot[LE,RO]{\small\thepage}
\renewcommand{\headrulewidth}{0pt}
\renewcommand{\sectionmark}[1]{\markright{\thesection.\ #1}}
\fancypagestyle{plain}{%
  \fancyhf{}%
  \fancyfoot[LE,RO]{\small \thepage}%
  \renewcommand{\headrulewidth}{0pt}%
}

% --------------------------------------------------------------
%  原书版式记号宏
% --------------------------------------------------------------
% 边码: 原打字稿页码(G. 手稿系统里的自引页码), 排在页边
\usepackage{marginnote}
\newcommand{\origpage}[1]{\marginnote{\normalfont\scriptsize p.~#1}}

% ⋄ 记号: 段首/行内正常大小; 上标形式为作者后补文字的位置标记
\DeclareRobustCommand{\pmark}{\ensuremath{\diamond}\hspace{0.15em}}
\DeclareRobustCommand{\dmark}{\textsuperscript{\ensuremath{\diamond}}}

% 脚注人工编号: \fnote{795}{(*)}{文本} —— 编号与(*)角标均照录原书
\newcommand{\fnote}[3]{%
  \begingroup
  \renewcommand{\thefootnote}{#1#2}%
  \footnote{#3}%
  \endgroup
}

% 第一部分的脚注: 页内序号作上标(#1), 注文行首印 "序号(连续注号)",
% 正文在序号之后另行括注同一连续注号 —— \fnoteP{1}{(3)}{注文}
% 原书有三条注无连续注号而改标日期(ch07 注3 = 30 Septembre,
% ch09 注11/12 = 8 août), 第二参数留空, 日期写进注文之首。
\newcommand{\fnoteP}[3]{%
  \begingroup
  \renewcommand{\thefootnote}{#1}\footnotemark
  \renewcommand{\thefootnote}{#1#2}\footnotetext{#3}%
  \endgroup
}

% 正文呼号与注文行首标号不同(原书照录)
\newcommand{\fnotevariant}[4]{%
  \begingroup
  \renewcommand{\thefootnote}{#1#2}\footnotemark
  \renewcommand{\thefootnote}{#1#3}\footnotetext{#4}%
  \endgroup
}

% 大注块标题: \notehead{171³}{斜体评论} —— 正文随后另起
\newcommand{\notehead}[2]{%
  \par\bigskip
  \noindent{\headfont 注~#1}\ifblank{#2}{}{\hspace{0.5em}{\itshape #2}}%
  \par\nobreak\smallskip\nobreak
}

% ∗ ∗ ∗ 分隔
\newcommand{\asterisms}{%
  \par\vspace{1.2ex}%
  \noindent\hspace*{2em}$\ast$\hspace{1.2em}$\ast$\hspace{1.2em}$\ast$%
  \par\vspace{1.2ex}%
}

% 无编号子目的目录条目
\newcommand{\tocsubsub}[1]{%
  \phantomsection
  \addcontentsline{toc}{subsubsection}{#1}%
}

% 部扉: 仿原书部题页
\newcommand{\partdivider}[3]{%
  \clearpage
  \thispagestyle{empty}%
  \phantomsection
  \addcontentsline{toc}{part}{#1\quad #3}%
  \vspace*{6.5cm}
  \begin{center}
    {\fontsize{16pt}{26pt}\selectfont\headfont #1\par}
    \vspace{0.9cm}
    {\fontsize{22pt}{34pt}\selectfont\headfont #2\par}
  \end{center}
  \clearpage
}

% 献辞页: 仿原书部题后的题献, 独占一页且不入目录
\newcommand{\dedication}[1]{%
  \thispagestyle{empty}%
  \vspace*{4cm}%
  \begin{center}#1\end{center}%
  \clearpage
}

% --------------------------------------------------------------
%  各卷"译者说明"的共用体例段(各卷的范围段自写, 后接本段)
% --------------------------------------------------------------
\newcommand{\CommonNoteBody}{%
为便于与原书及学界通行的打字稿页码对勘, 本册保留:

\begin{itemize}
  \item 原书全部章节编号(章~1--18, 节如~5.2、18.2, 子目如~18.1.1)与原书小节标题;
  \item 原书脚注的编号与角标——
        \begin{itemize}
          \item 前奏(第~1--4 章)与第一部的注在各页内重新编号, 正文中的上标即该
                页序号; 第一部紧随其后的括注(如“(3)”)是该部分的连续注号, 与第~12
                章“第一部分的注”里的“注~3”一一对应, 注文行首照原书印作“1(3)”式。
                其中三条注原书未标连续注号而改标日期(ch7 注~3 = 30 Septembre、
                ch9 注~11/12 = 8 août), 本册照录其日期;
          \item 第二部第~13--16~章各自重新编号, 分别为~1--41、1--86、1--132、1--29;
                第三部第~17~章自~1~起至~12, 第~18~章自~1~重新起算并延续至~794;
                末段自~\textbf{795(*)}~续至~1053(*)。各段与原书一一对应;
        \end{itemize}
  \item 页边码“p.~1”“p.~A3”“p.~1066”式小注——这是作者引用自己打字稿原稿时
        所用的页码, 与本排印本页码不同。前奏四章的边码各有前缀: 第~1 章作
        \texttt{A1…A6}, 第~2 章作 \texttt{P1…P65}, 第~3 章作 \texttt{L1…L56},
        第~4 章作罗马数字 \texttt{I…XXII}, 自第一部起作纯数字 \texttt{1…171};
  \item 作者用以标记段落的~$\diamond$~记号, 以及标示后来补写文字的上标
        ~\textsuperscript{$\diamond$}~记号;
  \item 原书大小注(“注~44'”“注~46$^1$”“注~171³”等)的编号与层级。
\end{itemize}

第~12 章“第一部分的注”把第~5 至第~11 章各页脚注汇为一编, 并给每条注配上节题
(如“12.19. 耶稣与十二使徒”), 因此部分注文与正文脚注重复——这是原书的体例, 非
重出之误。这一部分属随笔式的自省文字, 同一句法语在正文脚注与章末汇编注中
各有措辞, 本册有意保留各自的译笔, 不强求统一。

除上述标记外, 译文力求贴近作者驳杂、跳跃而辛辣的文风; 数学名词的译法见各卷
术语表(如~\emph{topos}~不译, \emph{faisceaux pervers}~作“反常层”)。文中的日期
括注(如“(5~月~30~日)”“(1984年3月补写)”)为作者历次补写时的纪日, 一律照录。文中
另有一些对注号、页码的引用(如“n°~1695”“打字稿第~227~页”)沿袭作者打字稿的编号
系统, 与本排印本不尽一致, 亦无法在本书中落实, 悉照原文录出。

全套五卷(另有全书合订一册)的分工是:

\begin{itemize}
  \item \textbf{卷一·前奏}——主题的呈现: 或四个乐章的前奏(原书 p.23--122);
  \item \textbf{卷二·第一部}——愚妄与更新(原书 p.123--254);
  \item \textbf{卷三·第二部}——葬礼(I): 或中国皇帝的新衣(原书 p.255--412);
  \item \textbf{卷四·第三部}——葬礼(II): 或阴阳之钥(原书 p.413--632);
  \item \textbf{卷五·第四部}——葬礼(III): 或四种运算(原书 p.633--929)。
\end{itemize}
}
 之后可用 \hypersetup 覆写)
% --------------------------------------------------------------
\hypersetup{
  pdftitle={收获与播种},
  pdfauthor={Alexandre Grothendieck / 数译 译},
  pdfsubject={Récoltes et Semailles — Réflexions et témoignage sur un passé de mathématicien}
}

% --------------------------------------------------------------
%  标题样式: 仿原书 —— 黑体左对齐, 编号后带句点
% --------------------------------------------------------------
\titleformat{\section}[block]
  {\normalfont\Large\headfont}{\thesection.}{0.5em}{}
\titlespacing*{\section}{0pt}{3.2ex plus 1.2ex minus .2ex}{1.6ex plus .2ex}
\titleformat{\subsection}[block]
  {\normalfont\large\headfont}{\thesubsection.}{0.5em}{}
\titlespacing*{\subsection}{0pt}{2.6ex plus 1ex minus .2ex}{1.2ex plus .2ex}
\titleformat{\subsubsection}[block]
  {\normalfont\normalsize\headfont}{\thesubsubsection.}{0.5em}{}
\titlespacing*{\subsubsection}{0pt}{2.2ex plus .8ex minus .2ex}{1ex plus .2ex}
\titleformat{name=\subsubsection,numberless}[block]
  {\normalfont\normalsize\headfont}{}{0em}{}
\setcounter{secnumdepth}{3}
\setcounter{tocdepth}{3}
% 为含长外文书名与行内公式的段落提供必要的断行余量
\setlength{\emergencystretch}{3em}

% 章标题样式: 原书章标题与节同级排印("13. A) HERITAGE..."),
% 故章标题取块状左对齐黑体, 字号介于节与部之间
\titleformat{\chapter}[block]
  {\normalfont\LARGE\headfont}{\thechapter.}{0.5em}{}
\titlespacing*{\chapter}{0pt}{3.6ex plus 1.4ex minus .2ex}{2ex plus .2ex}
% 目录中章号照原书用阿拉伯数字("13. A) ..." 而非 "第十三章")
\renewcommand{\CTEXthechapter}{\thechapter\hspace{.3em}}

% --------------------------------------------------------------
%  目录观感: 条目与引导点之间、引导点与页码之间各留 0.5em
%  内核 \@dottedtocline 让 \leaders 从标题末字起排, 第一个点几乎
%  贴着末字, 读起来像多标了一个句点; 页码前也没有余量。两处各补
%  0.5em。
% --------------------------------------------------------------
\makeatletter
\patchcmd{\@dottedtocline}{\nobreak\leaders}{\nobreak\hspace{0.5em}\leaders}%
  {\typeout{MT: TOC 引导点前已留 0.5em}}%
  {\typeout{MT-WARN: TOC 引导点前间距补丁未生效}}
\patchcmd{\@dottedtocline}{\hfill\nobreak\hb@xt@\@pnumwidth}%
  {\hfill\nobreak\hspace{0.5em}\hb@xt@\@pnumwidth}%
  {\typeout{MT: TOC 页码前已留 0.5em}}%
  {\typeout{MT-WARN: TOC 页码前间距补丁未生效}}
\makeatother

% --------------------------------------------------------------
%  正文半角标点后的间距
%  本书体例用半角 , ; : ? ! 与 ASCII 括号, 由 xeCJK 按 HalfRight
%  (右半留白) 处理。放宽到 0.45em, 与全角标点右半留白大致齐平。
% --------------------------------------------------------------
\xeCJKsetup{CJKecglue={\hskip 0.45em plus 0.06em minus 0.09em}}

% --------------------------------------------------------------
%  页眉页脚: 偶数页眉右对齐斜体节题; 页码页脚外侧
% --------------------------------------------------------------
\fancyhf{}
\fancyhead[RE]{\small\itshape\rightmark}
\fancyfoot[LE,RO]{\small\thepage}
\renewcommand{\headrulewidth}{0pt}
\renewcommand{\sectionmark}[1]{\markright{\thesection.\ #1}}
\fancypagestyle{plain}{%
  \fancyhf{}%
  \fancyfoot[LE,RO]{\small \thepage}%
  \renewcommand{\headrulewidth}{0pt}%
}

% --------------------------------------------------------------
%  原书版式记号宏
% --------------------------------------------------------------
% 边码: 原打字稿页码(G. 手稿系统里的自引页码), 排在页边
\usepackage{marginnote}
\newcommand{\origpage}[1]{\marginnote{\normalfont\scriptsize p.~#1}}

% ⋄ 记号: 段首/行内正常大小; 上标形式为作者后补文字的位置标记
\DeclareRobustCommand{\pmark}{\ensuremath{\diamond}\hspace{0.15em}}
\DeclareRobustCommand{\dmark}{\textsuperscript{\ensuremath{\diamond}}}

% 脚注人工编号: \fnote{795}{(*)}{文本} —— 编号与(*)角标均照录原书
\newcommand{\fnote}[3]{%
  \begingroup
  \renewcommand{\thefootnote}{#1#2}%
  \footnote{#3}%
  \endgroup
}

% 第一部分的脚注: 页内序号作上标(#1), 注文行首印 "序号(连续注号)",
% 正文在序号之后另行括注同一连续注号 —— \fnoteP{1}{(3)}{注文}
% 原书有三条注无连续注号而改标日期(ch07 注3 = 30 Septembre,
% ch09 注11/12 = 8 août), 第二参数留空, 日期写进注文之首。
\newcommand{\fnoteP}[3]{%
  \begingroup
  \renewcommand{\thefootnote}{#1}\footnotemark
  \renewcommand{\thefootnote}{#1#2}\footnotetext{#3}%
  \endgroup
}

% 正文呼号与注文行首标号不同(原书照录)
\newcommand{\fnotevariant}[4]{%
  \begingroup
  \renewcommand{\thefootnote}{#1#2}\footnotemark
  \renewcommand{\thefootnote}{#1#3}\footnotetext{#4}%
  \endgroup
}

% 大注块标题: \notehead{171³}{斜体评论} —— 正文随后另起
\newcommand{\notehead}[2]{%
  \par\bigskip
  \noindent{\headfont 注~#1}\ifblank{#2}{}{\hspace{0.5em}{\itshape #2}}%
  \par\nobreak\smallskip\nobreak
}

% ∗ ∗ ∗ 分隔
\newcommand{\asterisms}{%
  \par\vspace{1.2ex}%
  \noindent\hspace*{2em}$\ast$\hspace{1.2em}$\ast$\hspace{1.2em}$\ast$%
  \par\vspace{1.2ex}%
}

% 无编号子目的目录条目
\newcommand{\tocsubsub}[1]{%
  \phantomsection
  \addcontentsline{toc}{subsubsection}{#1}%
}

% 部扉: 仿原书部题页
\newcommand{\partdivider}[3]{%
  \clearpage
  \thispagestyle{empty}%
  \phantomsection
  \addcontentsline{toc}{part}{#1\quad #3}%
  \vspace*{6.5cm}
  \begin{center}
    {\fontsize{16pt}{26pt}\selectfont\headfont #1\par}
    \vspace{0.9cm}
    {\fontsize{22pt}{34pt}\selectfont\headfont #2\par}
  \end{center}
  \clearpage
}

% 献辞页: 仿原书部题后的题献, 独占一页且不入目录
\newcommand{\dedication}[1]{%
  \thispagestyle{empty}%
  \vspace*{4cm}%
  \begin{center}#1\end{center}%
  \clearpage
}

% --------------------------------------------------------------
%  各卷"译者说明"的共用体例段(各卷的范围段自写, 后接本段)
% --------------------------------------------------------------
\newcommand{\CommonNoteBody}{%
为便于与原书及学界通行的打字稿页码对勘, 本册保留:

\begin{itemize}
  \item 原书全部章节编号(章~1--18, 节如~5.2、18.2, 子目如~18.1.1)与原书小节标题;
  \item 原书脚注的编号与角标——
        \begin{itemize}
          \item 前奏(第~1--4 章)与第一部的注在各页内重新编号, 正文中的上标即该
                页序号; 第一部紧随其后的括注(如“(3)”)是该部分的连续注号, 与第~12
                章“第一部分的注”里的“注~3”一一对应, 注文行首照原书印作“1(3)”式。
                其中三条注原书未标连续注号而改标日期(ch7 注~3 = 30 Septembre、
                ch9 注~11/12 = 8 août), 本册照录其日期;
          \item 第二部第~13--16~章各自重新编号, 分别为~1--41、1--86、1--132、1--29;
                第三部第~17~章自~1~起至~12, 第~18~章自~1~重新起算并延续至~794;
                末段自~\textbf{795(*)}~续至~1053(*)。各段与原书一一对应;
        \end{itemize}
  \item 页边码“p.~1”“p.~A3”“p.~1066”式小注——这是作者引用自己打字稿原稿时
        所用的页码, 与本排印本页码不同。前奏四章的边码各有前缀: 第~1 章作
        \texttt{A1…A6}, 第~2 章作 \texttt{P1…P65}, 第~3 章作 \texttt{L1…L56},
        第~4 章作罗马数字 \texttt{I…XXII}, 自第一部起作纯数字 \texttt{1…171};
  \item 作者用以标记段落的~$\diamond$~记号, 以及标示后来补写文字的上标
        ~\textsuperscript{$\diamond$}~记号;
  \item 原书大小注(“注~44'”“注~46$^1$”“注~171³”等)的编号与层级。
\end{itemize}

第~12 章“第一部分的注”把第~5 至第~11 章各页脚注汇为一编, 并给每条注配上节题
(如“12.19. 耶稣与十二使徒”), 因此部分注文与正文脚注重复——这是原书的体例, 非
重出之误。这一部分属随笔式的自省文字, 同一句法语在正文脚注与章末汇编注中
各有措辞, 本册有意保留各自的译笔, 不强求统一。

除上述标记外, 译文力求贴近作者驳杂、跳跃而辛辣的文风; 数学名词的译法见各卷
术语表(如~\emph{topos}~不译, \emph{faisceaux pervers}~作“反常层”)。文中的日期
括注(如“(5~月~30~日)”“(1984年3月补写)”)为作者历次补写时的纪日, 一律照录。文中
另有一些对注号、页码的引用(如“n°~1695”“打字稿第~227~页”)沿袭作者打字稿的编号
系统, 与本排印本不尽一致, 亦无法在本书中落实, 悉照原文录出。

全套五卷(另有全书合订一册)的分工是:

\begin{itemize}
  \item \textbf{卷一·前奏}——主题的呈现: 或四个乐章的前奏(原书 p.23--122);
  \item \textbf{卷二·第一部}——愚妄与更新(原书 p.123--254);
  \item \textbf{卷三·第二部}——葬礼(I): 或中国皇帝的新衣(原书 p.255--412);
  \item \textbf{卷四·第三部}——葬礼(II): 或阴阳之钥(原书 p.413--632);
  \item \textbf{卷五·第四部}——葬礼(III): 或四种运算(原书 p.633--929)。
\end{itemize}
}
 之后可用 \hypersetup 覆写)
% --------------------------------------------------------------
\hypersetup{
  pdftitle={收获与播种},
  pdfauthor={Alexandre Grothendieck / 数译 译},
  pdfsubject={Récoltes et Semailles — Réflexions et témoignage sur un passé de mathématicien}
}

% --------------------------------------------------------------
%  标题样式: 仿原书 —— 黑体左对齐, 编号后带句点
% --------------------------------------------------------------
\titleformat{\section}[block]
  {\normalfont\Large\headfont}{\thesection.}{0.5em}{}
\titlespacing*{\section}{0pt}{3.2ex plus 1.2ex minus .2ex}{1.6ex plus .2ex}
\titleformat{\subsection}[block]
  {\normalfont\large\headfont}{\thesubsection.}{0.5em}{}
\titlespacing*{\subsection}{0pt}{2.6ex plus 1ex minus .2ex}{1.2ex plus .2ex}
\titleformat{\subsubsection}[block]
  {\normalfont\normalsize\headfont}{\thesubsubsection.}{0.5em}{}
\titlespacing*{\subsubsection}{0pt}{2.2ex plus .8ex minus .2ex}{1ex plus .2ex}
\titleformat{name=\subsubsection,numberless}[block]
  {\normalfont\normalsize\headfont}{}{0em}{}
\setcounter{secnumdepth}{3}
\setcounter{tocdepth}{3}
% 为含长外文书名与行内公式的段落提供必要的断行余量
\setlength{\emergencystretch}{3em}

% 章标题样式: 原书章标题与节同级排印("13. A) HERITAGE..."),
% 故章标题取块状左对齐黑体, 字号介于节与部之间
\titleformat{\chapter}[block]
  {\normalfont\LARGE\headfont}{\thechapter.}{0.5em}{}
\titlespacing*{\chapter}{0pt}{3.6ex plus 1.4ex minus .2ex}{2ex plus .2ex}
% 目录中章号照原书用阿拉伯数字("13. A) ..." 而非 "第十三章")
\renewcommand{\CTEXthechapter}{\thechapter\hspace{.3em}}

% --------------------------------------------------------------
%  目录观感: 条目与引导点之间、引导点与页码之间各留 0.5em
%  内核 \@dottedtocline 让 \leaders 从标题末字起排, 第一个点几乎
%  贴着末字, 读起来像多标了一个句点; 页码前也没有余量。两处各补
%  0.5em。
% --------------------------------------------------------------
\makeatletter
\patchcmd{\@dottedtocline}{\nobreak\leaders}{\nobreak\hspace{0.5em}\leaders}%
  {\typeout{MT: TOC 引导点前已留 0.5em}}%
  {\typeout{MT-WARN: TOC 引导点前间距补丁未生效}}
\patchcmd{\@dottedtocline}{\hfill\nobreak\hb@xt@\@pnumwidth}%
  {\hfill\nobreak\hspace{0.5em}\hb@xt@\@pnumwidth}%
  {\typeout{MT: TOC 页码前已留 0.5em}}%
  {\typeout{MT-WARN: TOC 页码前间距补丁未生效}}
\makeatother

% --------------------------------------------------------------
%  正文半角标点后的间距
%  本书体例用半角 , ; : ? ! 与 ASCII 括号, 由 xeCJK 按 HalfRight
%  (右半留白) 处理。放宽到 0.45em, 与全角标点右半留白大致齐平。
% --------------------------------------------------------------
\xeCJKsetup{CJKecglue={\hskip 0.45em plus 0.06em minus 0.09em}}

% --------------------------------------------------------------
%  页眉页脚: 偶数页眉右对齐斜体节题; 页码页脚外侧
% --------------------------------------------------------------
\fancyhf{}
\fancyhead[RE]{\small\itshape\rightmark}
\fancyfoot[LE,RO]{\small\thepage}
\renewcommand{\headrulewidth}{0pt}
\renewcommand{\sectionmark}[1]{\markright{\thesection.\ #1}}
\fancypagestyle{plain}{%
  \fancyhf{}%
  \fancyfoot[LE,RO]{\small \thepage}%
  \renewcommand{\headrulewidth}{0pt}%
}

% --------------------------------------------------------------
%  原书版式记号宏
% --------------------------------------------------------------
% 边码: 原打字稿页码(G. 手稿系统里的自引页码), 排在页边
\usepackage{marginnote}
\newcommand{\origpage}[1]{\marginnote{\normalfont\scriptsize p.~#1}}

% ⋄ 记号: 段首/行内正常大小; 上标形式为作者后补文字的位置标记
\DeclareRobustCommand{\pmark}{\ensuremath{\diamond}\hspace{0.15em}}
\DeclareRobustCommand{\dmark}{\textsuperscript{\ensuremath{\diamond}}}

% 脚注人工编号: \fnote{795}{(*)}{文本} —— 编号与(*)角标均照录原书
\newcommand{\fnote}[3]{%
  \begingroup
  \renewcommand{\thefootnote}{#1#2}%
  \footnote{#3}%
  \endgroup
}

% 第一部分的脚注: 页内序号作上标(#1), 注文行首印 "序号(连续注号)",
% 正文在序号之后另行括注同一连续注号 —— \fnoteP{1}{(3)}{注文}
% 原书有三条注无连续注号而改标日期(ch07 注3 = 30 Septembre,
% ch09 注11/12 = 8 août), 第二参数留空, 日期写进注文之首。
\newcommand{\fnoteP}[3]{%
  \begingroup
  \renewcommand{\thefootnote}{#1}\footnotemark
  \renewcommand{\thefootnote}{#1#2}\footnotetext{#3}%
  \endgroup
}

% 正文呼号与注文行首标号不同(原书照录)
\newcommand{\fnotevariant}[4]{%
  \begingroup
  \renewcommand{\thefootnote}{#1#2}\footnotemark
  \renewcommand{\thefootnote}{#1#3}\footnotetext{#4}%
  \endgroup
}

% 大注块标题: \notehead{171³}{斜体评论} —— 正文随后另起
\newcommand{\notehead}[2]{%
  \par\bigskip
  \noindent{\headfont 注~#1}\ifblank{#2}{}{\hspace{0.5em}{\itshape #2}}%
  \par\nobreak\smallskip\nobreak
}

% ∗ ∗ ∗ 分隔
\newcommand{\asterisms}{%
  \par\vspace{1.2ex}%
  \noindent\hspace*{2em}$\ast$\hspace{1.2em}$\ast$\hspace{1.2em}$\ast$%
  \par\vspace{1.2ex}%
}

% 无编号子目的目录条目
\newcommand{\tocsubsub}[1]{%
  \phantomsection
  \addcontentsline{toc}{subsubsection}{#1}%
}

% 部扉: 仿原书部题页
\newcommand{\partdivider}[3]{%
  \clearpage
  \thispagestyle{empty}%
  \phantomsection
  \addcontentsline{toc}{part}{#1\quad #3}%
  \vspace*{6.5cm}
  \begin{center}
    {\fontsize{16pt}{26pt}\selectfont\headfont #1\par}
    \vspace{0.9cm}
    {\fontsize{22pt}{34pt}\selectfont\headfont #2\par}
  \end{center}
  \clearpage
}

% 献辞页: 仿原书部题后的题献, 独占一页且不入目录
\newcommand{\dedication}[1]{%
  \thispagestyle{empty}%
  \vspace*{4cm}%
  \begin{center}#1\end{center}%
  \clearpage
}

% --------------------------------------------------------------
%  各卷"译者说明"的共用体例段(各卷的范围段自写, 后接本段)
% --------------------------------------------------------------
\newcommand{\CommonNoteBody}{%
为便于与原书及学界通行的打字稿页码对勘, 本册保留:

\begin{itemize}
  \item 原书全部章节编号(章~1--18, 节如~5.2、18.2, 子目如~18.1.1)与原书小节标题;
  \item 原书脚注的编号与角标——
        \begin{itemize}
          \item 前奏(第~1--4 章)与第一部的注在各页内重新编号, 正文中的上标即该
                页序号; 第一部紧随其后的括注(如“(3)”)是该部分的连续注号, 与第~12
                章“第一部分的注”里的“注~3”一一对应, 注文行首照原书印作“1(3)”式。
                其中三条注原书未标连续注号而改标日期(ch7 注~3 = 30 Septembre、
                ch9 注~11/12 = 8 août), 本册照录其日期;
          \item 第二部第~13--16~章各自重新编号, 分别为~1--41、1--86、1--132、1--29;
                第三部第~17~章自~1~起至~12, 第~18~章自~1~重新起算并延续至~794;
                末段自~\textbf{795(*)}~续至~1053(*)。各段与原书一一对应;
        \end{itemize}
  \item 页边码“p.~1”“p.~A3”“p.~1066”式小注——这是作者引用自己打字稿原稿时
        所用的页码, 与本排印本页码不同。前奏四章的边码各有前缀: 第~1 章作
        \texttt{A1…A6}, 第~2 章作 \texttt{P1…P65}, 第~3 章作 \texttt{L1…L56},
        第~4 章作罗马数字 \texttt{I…XXII}, 自第一部起作纯数字 \texttt{1…171};
  \item 作者用以标记段落的~$\diamond$~记号, 以及标示后来补写文字的上标
        ~\textsuperscript{$\diamond$}~记号;
  \item 原书大小注(“注~44'”“注~46$^1$”“注~171³”等)的编号与层级。
\end{itemize}

第~12 章“第一部分的注”把第~5 至第~11 章各页脚注汇为一编, 并给每条注配上节题
(如“12.19. 耶稣与十二使徒”), 因此部分注文与正文脚注重复——这是原书的体例, 非
重出之误。这一部分属随笔式的自省文字, 同一句法语在正文脚注与章末汇编注中
各有措辞, 本册有意保留各自的译笔, 不强求统一。

除上述标记外, 译文力求贴近作者驳杂、跳跃而辛辣的文风; 数学名词的译法见各卷
术语表(如~\emph{topos}~不译, \emph{faisceaux pervers}~作“反常层”)。文中的日期
括注(如“(5~月~30~日)”“(1984年3月补写)”)为作者历次补写时的纪日, 一律照录。文中
另有一些对注号、页码的引用(如“n°~1695”“打字稿第~227~页”)沿袭作者打字稿的编号
系统, 与本排印本不尽一致, 亦无法在本书中落实, 悉照原文录出。

全套五卷(另有全书合订一册)的分工是:

\begin{itemize}
  \item \textbf{卷一·前奏}——主题的呈现: 或四个乐章的前奏(原书 p.23--122);
  \item \textbf{卷二·第一部}——愚妄与更新(原书 p.123--254);
  \item \textbf{卷三·第二部}——葬礼(I): 或中国皇帝的新衣(原书 p.255--412);
  \item \textbf{卷四·第三部}——葬礼(II): 或阴阳之钥(原书 p.413--632);
  \item \textbf{卷五·第四部}——葬礼(III): 或四种运算(原书 p.633--929)。
\end{itemize}
}


% --------------------------------------------------------------
%  封面参数
% --------------------------------------------------------------
% ================================================================
%  《收获与播种》统一封面 —— TikZ 原生矢量绘制(不依赖任何外部图片)
%  母题: 暮色中的麦田 —— 犁沟向地平线的灭点收束, 前景麦穗剪影,
%        一位播种者把一把种子撒向天空, 种子沿弧线落入田里。
%  用法(在 \begin{document} 之前, 依次):
%      % ================================================================
%  《收获与播种》中文译本 —— 各卷共用导言区
%  由 vol-prelude / vol-part1 / vol-part2 / vol-part3 / vol-part4
%  以及全书合订 main-all 共同 \input。基类 mathtranslation.cls v3.2 不改。
%  --------------------------------------------------------------
%  样式均在此改写, 以贴近原书:
%    · 章标题仿原书 "5. Travail et découverte"(块状左对齐黑体)
%    · 节/小节: 黑体左对齐, 编号后带句点
%    · 脚注: 人工编号照录原书——前奏与第一部分页内重新编号,
%      第二至四部各自连续编号并带(*)角标; 第一部分另有 \fnoteP
%      处理 "(n)" 式连续注号
%    · 边码: 原打字稿页码 \origpage{N}(前奏作 A1、P1、L1、罗马数字等)
%    · ⋄ 记号: 段首补写标记 \pmark 与上标形式 \dmark
% ================================================================

% 题名含空格, 不能走类选项(会被 kvoptions 吃掉), 在此赋值
\renewcommand{\BookTitleEN}{Récoltes et Semailles}
\renewcommand{\OriginalAuthor}{Alexandre Grothendieck}
\renewcommand{\TranslationDate}{\today}

% --------------------------------------------------------------
%  罕用字后备(如"祐"): Fandol 不载时退到 SimSun
% --------------------------------------------------------------
\setCJKfallbackfamilyfont{\CJKrmdefault}{{SimSun}}

% --------------------------------------------------------------
%  hyperref 元信息(各卷在 % ================================================================
%  《收获与播种》中文译本 —— 各卷共用导言区
%  由 vol-prelude / vol-part1 / vol-part2 / vol-part3 / vol-part4
%  以及全书合订 main-all 共同 \input。基类 mathtranslation.cls v3.2 不改。
%  --------------------------------------------------------------
%  样式均在此改写, 以贴近原书:
%    · 章标题仿原书 "5. Travail et découverte"(块状左对齐黑体)
%    · 节/小节: 黑体左对齐, 编号后带句点
%    · 脚注: 人工编号照录原书——前奏与第一部分页内重新编号,
%      第二至四部各自连续编号并带(*)角标; 第一部分另有 \fnoteP
%      处理 "(n)" 式连续注号
%    · 边码: 原打字稿页码 \origpage{N}(前奏作 A1、P1、L1、罗马数字等)
%    · ⋄ 记号: 段首补写标记 \pmark 与上标形式 \dmark
% ================================================================

% 题名含空格, 不能走类选项(会被 kvoptions 吃掉), 在此赋值
\renewcommand{\BookTitleEN}{Récoltes et Semailles}
\renewcommand{\OriginalAuthor}{Alexandre Grothendieck}
\renewcommand{\TranslationDate}{\today}

% --------------------------------------------------------------
%  罕用字后备(如"祐"): Fandol 不载时退到 SimSun
% --------------------------------------------------------------
\setCJKfallbackfamilyfont{\CJKrmdefault}{{SimSun}}

% --------------------------------------------------------------
%  hyperref 元信息(各卷在 \input{common} 之后可用 \hypersetup 覆写)
% --------------------------------------------------------------
\hypersetup{
  pdftitle={收获与播种},
  pdfauthor={Alexandre Grothendieck / 数译 译},
  pdfsubject={Récoltes et Semailles — Réflexions et témoignage sur un passé de mathématicien}
}

% --------------------------------------------------------------
%  标题样式: 仿原书 —— 黑体左对齐, 编号后带句点
% --------------------------------------------------------------
\titleformat{\section}[block]
  {\normalfont\Large\headfont}{\thesection.}{0.5em}{}
\titlespacing*{\section}{0pt}{3.2ex plus 1.2ex minus .2ex}{1.6ex plus .2ex}
\titleformat{\subsection}[block]
  {\normalfont\large\headfont}{\thesubsection.}{0.5em}{}
\titlespacing*{\subsection}{0pt}{2.6ex plus 1ex minus .2ex}{1.2ex plus .2ex}
\titleformat{\subsubsection}[block]
  {\normalfont\normalsize\headfont}{\thesubsubsection.}{0.5em}{}
\titlespacing*{\subsubsection}{0pt}{2.2ex plus .8ex minus .2ex}{1ex plus .2ex}
\titleformat{name=\subsubsection,numberless}[block]
  {\normalfont\normalsize\headfont}{}{0em}{}
\setcounter{secnumdepth}{3}
\setcounter{tocdepth}{3}
% 为含长外文书名与行内公式的段落提供必要的断行余量
\setlength{\emergencystretch}{3em}

% 章标题样式: 原书章标题与节同级排印("13. A) HERITAGE..."),
% 故章标题取块状左对齐黑体, 字号介于节与部之间
\titleformat{\chapter}[block]
  {\normalfont\LARGE\headfont}{\thechapter.}{0.5em}{}
\titlespacing*{\chapter}{0pt}{3.6ex plus 1.4ex minus .2ex}{2ex plus .2ex}
% 目录中章号照原书用阿拉伯数字("13. A) ..." 而非 "第十三章")
\renewcommand{\CTEXthechapter}{\thechapter\hspace{.3em}}

% --------------------------------------------------------------
%  目录观感: 条目与引导点之间、引导点与页码之间各留 0.5em
%  内核 \@dottedtocline 让 \leaders 从标题末字起排, 第一个点几乎
%  贴着末字, 读起来像多标了一个句点; 页码前也没有余量。两处各补
%  0.5em。
% --------------------------------------------------------------
\makeatletter
\patchcmd{\@dottedtocline}{\nobreak\leaders}{\nobreak\hspace{0.5em}\leaders}%
  {\typeout{MT: TOC 引导点前已留 0.5em}}%
  {\typeout{MT-WARN: TOC 引导点前间距补丁未生效}}
\patchcmd{\@dottedtocline}{\hfill\nobreak\hb@xt@\@pnumwidth}%
  {\hfill\nobreak\hspace{0.5em}\hb@xt@\@pnumwidth}%
  {\typeout{MT: TOC 页码前已留 0.5em}}%
  {\typeout{MT-WARN: TOC 页码前间距补丁未生效}}
\makeatother

% --------------------------------------------------------------
%  正文半角标点后的间距
%  本书体例用半角 , ; : ? ! 与 ASCII 括号, 由 xeCJK 按 HalfRight
%  (右半留白) 处理。放宽到 0.45em, 与全角标点右半留白大致齐平。
% --------------------------------------------------------------
\xeCJKsetup{CJKecglue={\hskip 0.45em plus 0.06em minus 0.09em}}

% --------------------------------------------------------------
%  页眉页脚: 偶数页眉右对齐斜体节题; 页码页脚外侧
% --------------------------------------------------------------
\fancyhf{}
\fancyhead[RE]{\small\itshape\rightmark}
\fancyfoot[LE,RO]{\small\thepage}
\renewcommand{\headrulewidth}{0pt}
\renewcommand{\sectionmark}[1]{\markright{\thesection.\ #1}}
\fancypagestyle{plain}{%
  \fancyhf{}%
  \fancyfoot[LE,RO]{\small \thepage}%
  \renewcommand{\headrulewidth}{0pt}%
}

% --------------------------------------------------------------
%  原书版式记号宏
% --------------------------------------------------------------
% 边码: 原打字稿页码(G. 手稿系统里的自引页码), 排在页边
\usepackage{marginnote}
\newcommand{\origpage}[1]{\marginnote{\normalfont\scriptsize p.~#1}}

% ⋄ 记号: 段首/行内正常大小; 上标形式为作者后补文字的位置标记
\DeclareRobustCommand{\pmark}{\ensuremath{\diamond}\hspace{0.15em}}
\DeclareRobustCommand{\dmark}{\textsuperscript{\ensuremath{\diamond}}}

% 脚注人工编号: \fnote{795}{(*)}{文本} —— 编号与(*)角标均照录原书
\newcommand{\fnote}[3]{%
  \begingroup
  \renewcommand{\thefootnote}{#1#2}%
  \footnote{#3}%
  \endgroup
}

% 第一部分的脚注: 页内序号作上标(#1), 注文行首印 "序号(连续注号)",
% 正文在序号之后另行括注同一连续注号 —— \fnoteP{1}{(3)}{注文}
% 原书有三条注无连续注号而改标日期(ch07 注3 = 30 Septembre,
% ch09 注11/12 = 8 août), 第二参数留空, 日期写进注文之首。
\newcommand{\fnoteP}[3]{%
  \begingroup
  \renewcommand{\thefootnote}{#1}\footnotemark
  \renewcommand{\thefootnote}{#1#2}\footnotetext{#3}%
  \endgroup
}

% 正文呼号与注文行首标号不同(原书照录)
\newcommand{\fnotevariant}[4]{%
  \begingroup
  \renewcommand{\thefootnote}{#1#2}\footnotemark
  \renewcommand{\thefootnote}{#1#3}\footnotetext{#4}%
  \endgroup
}

% 大注块标题: \notehead{171³}{斜体评论} —— 正文随后另起
\newcommand{\notehead}[2]{%
  \par\bigskip
  \noindent{\headfont 注~#1}\ifblank{#2}{}{\hspace{0.5em}{\itshape #2}}%
  \par\nobreak\smallskip\nobreak
}

% ∗ ∗ ∗ 分隔
\newcommand{\asterisms}{%
  \par\vspace{1.2ex}%
  \noindent\hspace*{2em}$\ast$\hspace{1.2em}$\ast$\hspace{1.2em}$\ast$%
  \par\vspace{1.2ex}%
}

% 无编号子目的目录条目
\newcommand{\tocsubsub}[1]{%
  \phantomsection
  \addcontentsline{toc}{subsubsection}{#1}%
}

% 部扉: 仿原书部题页
\newcommand{\partdivider}[3]{%
  \clearpage
  \thispagestyle{empty}%
  \phantomsection
  \addcontentsline{toc}{part}{#1\quad #3}%
  \vspace*{6.5cm}
  \begin{center}
    {\fontsize{16pt}{26pt}\selectfont\headfont #1\par}
    \vspace{0.9cm}
    {\fontsize{22pt}{34pt}\selectfont\headfont #2\par}
  \end{center}
  \clearpage
}

% 献辞页: 仿原书部题后的题献, 独占一页且不入目录
\newcommand{\dedication}[1]{%
  \thispagestyle{empty}%
  \vspace*{4cm}%
  \begin{center}#1\end{center}%
  \clearpage
}

% --------------------------------------------------------------
%  各卷"译者说明"的共用体例段(各卷的范围段自写, 后接本段)
% --------------------------------------------------------------
\newcommand{\CommonNoteBody}{%
为便于与原书及学界通行的打字稿页码对勘, 本册保留:

\begin{itemize}
  \item 原书全部章节编号(章~1--18, 节如~5.2、18.2, 子目如~18.1.1)与原书小节标题;
  \item 原书脚注的编号与角标——
        \begin{itemize}
          \item 前奏(第~1--4 章)与第一部的注在各页内重新编号, 正文中的上标即该
                页序号; 第一部紧随其后的括注(如“(3)”)是该部分的连续注号, 与第~12
                章“第一部分的注”里的“注~3”一一对应, 注文行首照原书印作“1(3)”式。
                其中三条注原书未标连续注号而改标日期(ch7 注~3 = 30 Septembre、
                ch9 注~11/12 = 8 août), 本册照录其日期;
          \item 第二部第~13--16~章各自重新编号, 分别为~1--41、1--86、1--132、1--29;
                第三部第~17~章自~1~起至~12, 第~18~章自~1~重新起算并延续至~794;
                末段自~\textbf{795(*)}~续至~1053(*)。各段与原书一一对应;
        \end{itemize}
  \item 页边码“p.~1”“p.~A3”“p.~1066”式小注——这是作者引用自己打字稿原稿时
        所用的页码, 与本排印本页码不同。前奏四章的边码各有前缀: 第~1 章作
        \texttt{A1…A6}, 第~2 章作 \texttt{P1…P65}, 第~3 章作 \texttt{L1…L56},
        第~4 章作罗马数字 \texttt{I…XXII}, 自第一部起作纯数字 \texttt{1…171};
  \item 作者用以标记段落的~$\diamond$~记号, 以及标示后来补写文字的上标
        ~\textsuperscript{$\diamond$}~记号;
  \item 原书大小注(“注~44'”“注~46$^1$”“注~171³”等)的编号与层级。
\end{itemize}

第~12 章“第一部分的注”把第~5 至第~11 章各页脚注汇为一编, 并给每条注配上节题
(如“12.19. 耶稣与十二使徒”), 因此部分注文与正文脚注重复——这是原书的体例, 非
重出之误。这一部分属随笔式的自省文字, 同一句法语在正文脚注与章末汇编注中
各有措辞, 本册有意保留各自的译笔, 不强求统一。

除上述标记外, 译文力求贴近作者驳杂、跳跃而辛辣的文风; 数学名词的译法见各卷
术语表(如~\emph{topos}~不译, \emph{faisceaux pervers}~作“反常层”)。文中的日期
括注(如“(5~月~30~日)”“(1984年3月补写)”)为作者历次补写时的纪日, 一律照录。文中
另有一些对注号、页码的引用(如“n°~1695”“打字稿第~227~页”)沿袭作者打字稿的编号
系统, 与本排印本不尽一致, 亦无法在本书中落实, 悉照原文录出。

全套五卷(另有全书合订一册)的分工是:

\begin{itemize}
  \item \textbf{卷一·前奏}——主题的呈现: 或四个乐章的前奏(原书 p.23--122);
  \item \textbf{卷二·第一部}——愚妄与更新(原书 p.123--254);
  \item \textbf{卷三·第二部}——葬礼(I): 或中国皇帝的新衣(原书 p.255--412);
  \item \textbf{卷四·第三部}——葬礼(II): 或阴阳之钥(原书 p.413--632);
  \item \textbf{卷五·第四部}——葬礼(III): 或四种运算(原书 p.633--929)。
\end{itemize}
}
 之后可用 \hypersetup 覆写)
% --------------------------------------------------------------
\hypersetup{
  pdftitle={收获与播种},
  pdfauthor={Alexandre Grothendieck / 数译 译},
  pdfsubject={Récoltes et Semailles — Réflexions et témoignage sur un passé de mathématicien}
}

% --------------------------------------------------------------
%  标题样式: 仿原书 —— 黑体左对齐, 编号后带句点
% --------------------------------------------------------------
\titleformat{\section}[block]
  {\normalfont\Large\headfont}{\thesection.}{0.5em}{}
\titlespacing*{\section}{0pt}{3.2ex plus 1.2ex minus .2ex}{1.6ex plus .2ex}
\titleformat{\subsection}[block]
  {\normalfont\large\headfont}{\thesubsection.}{0.5em}{}
\titlespacing*{\subsection}{0pt}{2.6ex plus 1ex minus .2ex}{1.2ex plus .2ex}
\titleformat{\subsubsection}[block]
  {\normalfont\normalsize\headfont}{\thesubsubsection.}{0.5em}{}
\titlespacing*{\subsubsection}{0pt}{2.2ex plus .8ex minus .2ex}{1ex plus .2ex}
\titleformat{name=\subsubsection,numberless}[block]
  {\normalfont\normalsize\headfont}{}{0em}{}
\setcounter{secnumdepth}{3}
\setcounter{tocdepth}{3}
% 为含长外文书名与行内公式的段落提供必要的断行余量
\setlength{\emergencystretch}{3em}

% 章标题样式: 原书章标题与节同级排印("13. A) HERITAGE..."),
% 故章标题取块状左对齐黑体, 字号介于节与部之间
\titleformat{\chapter}[block]
  {\normalfont\LARGE\headfont}{\thechapter.}{0.5em}{}
\titlespacing*{\chapter}{0pt}{3.6ex plus 1.4ex minus .2ex}{2ex plus .2ex}
% 目录中章号照原书用阿拉伯数字("13. A) ..." 而非 "第十三章")
\renewcommand{\CTEXthechapter}{\thechapter\hspace{.3em}}

% --------------------------------------------------------------
%  目录观感: 条目与引导点之间、引导点与页码之间各留 0.5em
%  内核 \@dottedtocline 让 \leaders 从标题末字起排, 第一个点几乎
%  贴着末字, 读起来像多标了一个句点; 页码前也没有余量。两处各补
%  0.5em。
% --------------------------------------------------------------
\makeatletter
\patchcmd{\@dottedtocline}{\nobreak\leaders}{\nobreak\hspace{0.5em}\leaders}%
  {\typeout{MT: TOC 引导点前已留 0.5em}}%
  {\typeout{MT-WARN: TOC 引导点前间距补丁未生效}}
\patchcmd{\@dottedtocline}{\hfill\nobreak\hb@xt@\@pnumwidth}%
  {\hfill\nobreak\hspace{0.5em}\hb@xt@\@pnumwidth}%
  {\typeout{MT: TOC 页码前已留 0.5em}}%
  {\typeout{MT-WARN: TOC 页码前间距补丁未生效}}
\makeatother

% --------------------------------------------------------------
%  正文半角标点后的间距
%  本书体例用半角 , ; : ? ! 与 ASCII 括号, 由 xeCJK 按 HalfRight
%  (右半留白) 处理。放宽到 0.45em, 与全角标点右半留白大致齐平。
% --------------------------------------------------------------
\xeCJKsetup{CJKecglue={\hskip 0.45em plus 0.06em minus 0.09em}}

% --------------------------------------------------------------
%  页眉页脚: 偶数页眉右对齐斜体节题; 页码页脚外侧
% --------------------------------------------------------------
\fancyhf{}
\fancyhead[RE]{\small\itshape\rightmark}
\fancyfoot[LE,RO]{\small\thepage}
\renewcommand{\headrulewidth}{0pt}
\renewcommand{\sectionmark}[1]{\markright{\thesection.\ #1}}
\fancypagestyle{plain}{%
  \fancyhf{}%
  \fancyfoot[LE,RO]{\small \thepage}%
  \renewcommand{\headrulewidth}{0pt}%
}

% --------------------------------------------------------------
%  原书版式记号宏
% --------------------------------------------------------------
% 边码: 原打字稿页码(G. 手稿系统里的自引页码), 排在页边
\usepackage{marginnote}
\newcommand{\origpage}[1]{\marginnote{\normalfont\scriptsize p.~#1}}

% ⋄ 记号: 段首/行内正常大小; 上标形式为作者后补文字的位置标记
\DeclareRobustCommand{\pmark}{\ensuremath{\diamond}\hspace{0.15em}}
\DeclareRobustCommand{\dmark}{\textsuperscript{\ensuremath{\diamond}}}

% 脚注人工编号: \fnote{795}{(*)}{文本} —— 编号与(*)角标均照录原书
\newcommand{\fnote}[3]{%
  \begingroup
  \renewcommand{\thefootnote}{#1#2}%
  \footnote{#3}%
  \endgroup
}

% 第一部分的脚注: 页内序号作上标(#1), 注文行首印 "序号(连续注号)",
% 正文在序号之后另行括注同一连续注号 —— \fnoteP{1}{(3)}{注文}
% 原书有三条注无连续注号而改标日期(ch07 注3 = 30 Septembre,
% ch09 注11/12 = 8 août), 第二参数留空, 日期写进注文之首。
\newcommand{\fnoteP}[3]{%
  \begingroup
  \renewcommand{\thefootnote}{#1}\footnotemark
  \renewcommand{\thefootnote}{#1#2}\footnotetext{#3}%
  \endgroup
}

% 正文呼号与注文行首标号不同(原书照录)
\newcommand{\fnotevariant}[4]{%
  \begingroup
  \renewcommand{\thefootnote}{#1#2}\footnotemark
  \renewcommand{\thefootnote}{#1#3}\footnotetext{#4}%
  \endgroup
}

% 大注块标题: \notehead{171³}{斜体评论} —— 正文随后另起
\newcommand{\notehead}[2]{%
  \par\bigskip
  \noindent{\headfont 注~#1}\ifblank{#2}{}{\hspace{0.5em}{\itshape #2}}%
  \par\nobreak\smallskip\nobreak
}

% ∗ ∗ ∗ 分隔
\newcommand{\asterisms}{%
  \par\vspace{1.2ex}%
  \noindent\hspace*{2em}$\ast$\hspace{1.2em}$\ast$\hspace{1.2em}$\ast$%
  \par\vspace{1.2ex}%
}

% 无编号子目的目录条目
\newcommand{\tocsubsub}[1]{%
  \phantomsection
  \addcontentsline{toc}{subsubsection}{#1}%
}

% 部扉: 仿原书部题页
\newcommand{\partdivider}[3]{%
  \clearpage
  \thispagestyle{empty}%
  \phantomsection
  \addcontentsline{toc}{part}{#1\quad #3}%
  \vspace*{6.5cm}
  \begin{center}
    {\fontsize{16pt}{26pt}\selectfont\headfont #1\par}
    \vspace{0.9cm}
    {\fontsize{22pt}{34pt}\selectfont\headfont #2\par}
  \end{center}
  \clearpage
}

% 献辞页: 仿原书部题后的题献, 独占一页且不入目录
\newcommand{\dedication}[1]{%
  \thispagestyle{empty}%
  \vspace*{4cm}%
  \begin{center}#1\end{center}%
  \clearpage
}

% --------------------------------------------------------------
%  各卷"译者说明"的共用体例段(各卷的范围段自写, 后接本段)
% --------------------------------------------------------------
\newcommand{\CommonNoteBody}{%
为便于与原书及学界通行的打字稿页码对勘, 本册保留:

\begin{itemize}
  \item 原书全部章节编号(章~1--18, 节如~5.2、18.2, 子目如~18.1.1)与原书小节标题;
  \item 原书脚注的编号与角标——
        \begin{itemize}
          \item 前奏(第~1--4 章)与第一部的注在各页内重新编号, 正文中的上标即该
                页序号; 第一部紧随其后的括注(如“(3)”)是该部分的连续注号, 与第~12
                章“第一部分的注”里的“注~3”一一对应, 注文行首照原书印作“1(3)”式。
                其中三条注原书未标连续注号而改标日期(ch7 注~3 = 30 Septembre、
                ch9 注~11/12 = 8 août), 本册照录其日期;
          \item 第二部第~13--16~章各自重新编号, 分别为~1--41、1--86、1--132、1--29;
                第三部第~17~章自~1~起至~12, 第~18~章自~1~重新起算并延续至~794;
                末段自~\textbf{795(*)}~续至~1053(*)。各段与原书一一对应;
        \end{itemize}
  \item 页边码“p.~1”“p.~A3”“p.~1066”式小注——这是作者引用自己打字稿原稿时
        所用的页码, 与本排印本页码不同。前奏四章的边码各有前缀: 第~1 章作
        \texttt{A1…A6}, 第~2 章作 \texttt{P1…P65}, 第~3 章作 \texttt{L1…L56},
        第~4 章作罗马数字 \texttt{I…XXII}, 自第一部起作纯数字 \texttt{1…171};
  \item 作者用以标记段落的~$\diamond$~记号, 以及标示后来补写文字的上标
        ~\textsuperscript{$\diamond$}~记号;
  \item 原书大小注(“注~44'”“注~46$^1$”“注~171³”等)的编号与层级。
\end{itemize}

第~12 章“第一部分的注”把第~5 至第~11 章各页脚注汇为一编, 并给每条注配上节题
(如“12.19. 耶稣与十二使徒”), 因此部分注文与正文脚注重复——这是原书的体例, 非
重出之误。这一部分属随笔式的自省文字, 同一句法语在正文脚注与章末汇编注中
各有措辞, 本册有意保留各自的译笔, 不强求统一。

除上述标记外, 译文力求贴近作者驳杂、跳跃而辛辣的文风; 数学名词的译法见各卷
术语表(如~\emph{topos}~不译, \emph{faisceaux pervers}~作“反常层”)。文中的日期
括注(如“(5~月~30~日)”“(1984年3月补写)”)为作者历次补写时的纪日, 一律照录。文中
另有一些对注号、页码的引用(如“n°~1695”“打字稿第~227~页”)沿袭作者打字稿的编号
系统, 与本排印本不尽一致, 亦无法在本书中落实, 悉照原文录出。

全套五卷(另有全书合订一册)的分工是:

\begin{itemize}
  \item \textbf{卷一·前奏}——主题的呈现: 或四个乐章的前奏(原书 p.23--122);
  \item \textbf{卷二·第一部}——愚妄与更新(原书 p.123--254);
  \item \textbf{卷三·第二部}——葬礼(I): 或中国皇帝的新衣(原书 p.255--412);
  \item \textbf{卷四·第三部}——葬礼(II): 或阴阳之钥(原书 p.413--632);
  \item \textbf{卷五·第四部}——葬礼(III): 或四种运算(原书 p.633--929)。
\end{itemize}
}
                       % 共用导言区
%      % ================================================================
%  《收获与播种》统一封面 —— TikZ 原生矢量绘制(不依赖任何外部图片)
%  母题: 暮色中的麦田 —— 犁沟向地平线的灭点收束, 前景麦穗剪影,
%        一位播种者把一把种子撒向天空, 种子沿弧线落入田里。
%  用法(在 \begin{document} 之前, 依次):
%      % ================================================================
%  《收获与播种》中文译本 —— 各卷共用导言区
%  由 vol-prelude / vol-part1 / vol-part2 / vol-part3 / vol-part4
%  以及全书合订 main-all 共同 \input。基类 mathtranslation.cls v3.2 不改。
%  --------------------------------------------------------------
%  样式均在此改写, 以贴近原书:
%    · 章标题仿原书 "5. Travail et découverte"(块状左对齐黑体)
%    · 节/小节: 黑体左对齐, 编号后带句点
%    · 脚注: 人工编号照录原书——前奏与第一部分页内重新编号,
%      第二至四部各自连续编号并带(*)角标; 第一部分另有 \fnoteP
%      处理 "(n)" 式连续注号
%    · 边码: 原打字稿页码 \origpage{N}(前奏作 A1、P1、L1、罗马数字等)
%    · ⋄ 记号: 段首补写标记 \pmark 与上标形式 \dmark
% ================================================================

% 题名含空格, 不能走类选项(会被 kvoptions 吃掉), 在此赋值
\renewcommand{\BookTitleEN}{Récoltes et Semailles}
\renewcommand{\OriginalAuthor}{Alexandre Grothendieck}
\renewcommand{\TranslationDate}{\today}

% --------------------------------------------------------------
%  罕用字后备(如"祐"): Fandol 不载时退到 SimSun
% --------------------------------------------------------------
\setCJKfallbackfamilyfont{\CJKrmdefault}{{SimSun}}

% --------------------------------------------------------------
%  hyperref 元信息(各卷在 \input{common} 之后可用 \hypersetup 覆写)
% --------------------------------------------------------------
\hypersetup{
  pdftitle={收获与播种},
  pdfauthor={Alexandre Grothendieck / 数译 译},
  pdfsubject={Récoltes et Semailles — Réflexions et témoignage sur un passé de mathématicien}
}

% --------------------------------------------------------------
%  标题样式: 仿原书 —— 黑体左对齐, 编号后带句点
% --------------------------------------------------------------
\titleformat{\section}[block]
  {\normalfont\Large\headfont}{\thesection.}{0.5em}{}
\titlespacing*{\section}{0pt}{3.2ex plus 1.2ex minus .2ex}{1.6ex plus .2ex}
\titleformat{\subsection}[block]
  {\normalfont\large\headfont}{\thesubsection.}{0.5em}{}
\titlespacing*{\subsection}{0pt}{2.6ex plus 1ex minus .2ex}{1.2ex plus .2ex}
\titleformat{\subsubsection}[block]
  {\normalfont\normalsize\headfont}{\thesubsubsection.}{0.5em}{}
\titlespacing*{\subsubsection}{0pt}{2.2ex plus .8ex minus .2ex}{1ex plus .2ex}
\titleformat{name=\subsubsection,numberless}[block]
  {\normalfont\normalsize\headfont}{}{0em}{}
\setcounter{secnumdepth}{3}
\setcounter{tocdepth}{3}
% 为含长外文书名与行内公式的段落提供必要的断行余量
\setlength{\emergencystretch}{3em}

% 章标题样式: 原书章标题与节同级排印("13. A) HERITAGE..."),
% 故章标题取块状左对齐黑体, 字号介于节与部之间
\titleformat{\chapter}[block]
  {\normalfont\LARGE\headfont}{\thechapter.}{0.5em}{}
\titlespacing*{\chapter}{0pt}{3.6ex plus 1.4ex minus .2ex}{2ex plus .2ex}
% 目录中章号照原书用阿拉伯数字("13. A) ..." 而非 "第十三章")
\renewcommand{\CTEXthechapter}{\thechapter\hspace{.3em}}

% --------------------------------------------------------------
%  目录观感: 条目与引导点之间、引导点与页码之间各留 0.5em
%  内核 \@dottedtocline 让 \leaders 从标题末字起排, 第一个点几乎
%  贴着末字, 读起来像多标了一个句点; 页码前也没有余量。两处各补
%  0.5em。
% --------------------------------------------------------------
\makeatletter
\patchcmd{\@dottedtocline}{\nobreak\leaders}{\nobreak\hspace{0.5em}\leaders}%
  {\typeout{MT: TOC 引导点前已留 0.5em}}%
  {\typeout{MT-WARN: TOC 引导点前间距补丁未生效}}
\patchcmd{\@dottedtocline}{\hfill\nobreak\hb@xt@\@pnumwidth}%
  {\hfill\nobreak\hspace{0.5em}\hb@xt@\@pnumwidth}%
  {\typeout{MT: TOC 页码前已留 0.5em}}%
  {\typeout{MT-WARN: TOC 页码前间距补丁未生效}}
\makeatother

% --------------------------------------------------------------
%  正文半角标点后的间距
%  本书体例用半角 , ; : ? ! 与 ASCII 括号, 由 xeCJK 按 HalfRight
%  (右半留白) 处理。放宽到 0.45em, 与全角标点右半留白大致齐平。
% --------------------------------------------------------------
\xeCJKsetup{CJKecglue={\hskip 0.45em plus 0.06em minus 0.09em}}

% --------------------------------------------------------------
%  页眉页脚: 偶数页眉右对齐斜体节题; 页码页脚外侧
% --------------------------------------------------------------
\fancyhf{}
\fancyhead[RE]{\small\itshape\rightmark}
\fancyfoot[LE,RO]{\small\thepage}
\renewcommand{\headrulewidth}{0pt}
\renewcommand{\sectionmark}[1]{\markright{\thesection.\ #1}}
\fancypagestyle{plain}{%
  \fancyhf{}%
  \fancyfoot[LE,RO]{\small \thepage}%
  \renewcommand{\headrulewidth}{0pt}%
}

% --------------------------------------------------------------
%  原书版式记号宏
% --------------------------------------------------------------
% 边码: 原打字稿页码(G. 手稿系统里的自引页码), 排在页边
\usepackage{marginnote}
\newcommand{\origpage}[1]{\marginnote{\normalfont\scriptsize p.~#1}}

% ⋄ 记号: 段首/行内正常大小; 上标形式为作者后补文字的位置标记
\DeclareRobustCommand{\pmark}{\ensuremath{\diamond}\hspace{0.15em}}
\DeclareRobustCommand{\dmark}{\textsuperscript{\ensuremath{\diamond}}}

% 脚注人工编号: \fnote{795}{(*)}{文本} —— 编号与(*)角标均照录原书
\newcommand{\fnote}[3]{%
  \begingroup
  \renewcommand{\thefootnote}{#1#2}%
  \footnote{#3}%
  \endgroup
}

% 第一部分的脚注: 页内序号作上标(#1), 注文行首印 "序号(连续注号)",
% 正文在序号之后另行括注同一连续注号 —— \fnoteP{1}{(3)}{注文}
% 原书有三条注无连续注号而改标日期(ch07 注3 = 30 Septembre,
% ch09 注11/12 = 8 août), 第二参数留空, 日期写进注文之首。
\newcommand{\fnoteP}[3]{%
  \begingroup
  \renewcommand{\thefootnote}{#1}\footnotemark
  \renewcommand{\thefootnote}{#1#2}\footnotetext{#3}%
  \endgroup
}

% 正文呼号与注文行首标号不同(原书照录)
\newcommand{\fnotevariant}[4]{%
  \begingroup
  \renewcommand{\thefootnote}{#1#2}\footnotemark
  \renewcommand{\thefootnote}{#1#3}\footnotetext{#4}%
  \endgroup
}

% 大注块标题: \notehead{171³}{斜体评论} —— 正文随后另起
\newcommand{\notehead}[2]{%
  \par\bigskip
  \noindent{\headfont 注~#1}\ifblank{#2}{}{\hspace{0.5em}{\itshape #2}}%
  \par\nobreak\smallskip\nobreak
}

% ∗ ∗ ∗ 分隔
\newcommand{\asterisms}{%
  \par\vspace{1.2ex}%
  \noindent\hspace*{2em}$\ast$\hspace{1.2em}$\ast$\hspace{1.2em}$\ast$%
  \par\vspace{1.2ex}%
}

% 无编号子目的目录条目
\newcommand{\tocsubsub}[1]{%
  \phantomsection
  \addcontentsline{toc}{subsubsection}{#1}%
}

% 部扉: 仿原书部题页
\newcommand{\partdivider}[3]{%
  \clearpage
  \thispagestyle{empty}%
  \phantomsection
  \addcontentsline{toc}{part}{#1\quad #3}%
  \vspace*{6.5cm}
  \begin{center}
    {\fontsize{16pt}{26pt}\selectfont\headfont #1\par}
    \vspace{0.9cm}
    {\fontsize{22pt}{34pt}\selectfont\headfont #2\par}
  \end{center}
  \clearpage
}

% 献辞页: 仿原书部题后的题献, 独占一页且不入目录
\newcommand{\dedication}[1]{%
  \thispagestyle{empty}%
  \vspace*{4cm}%
  \begin{center}#1\end{center}%
  \clearpage
}

% --------------------------------------------------------------
%  各卷"译者说明"的共用体例段(各卷的范围段自写, 后接本段)
% --------------------------------------------------------------
\newcommand{\CommonNoteBody}{%
为便于与原书及学界通行的打字稿页码对勘, 本册保留:

\begin{itemize}
  \item 原书全部章节编号(章~1--18, 节如~5.2、18.2, 子目如~18.1.1)与原书小节标题;
  \item 原书脚注的编号与角标——
        \begin{itemize}
          \item 前奏(第~1--4 章)与第一部的注在各页内重新编号, 正文中的上标即该
                页序号; 第一部紧随其后的括注(如“(3)”)是该部分的连续注号, 与第~12
                章“第一部分的注”里的“注~3”一一对应, 注文行首照原书印作“1(3)”式。
                其中三条注原书未标连续注号而改标日期(ch7 注~3 = 30 Septembre、
                ch9 注~11/12 = 8 août), 本册照录其日期;
          \item 第二部第~13--16~章各自重新编号, 分别为~1--41、1--86、1--132、1--29;
                第三部第~17~章自~1~起至~12, 第~18~章自~1~重新起算并延续至~794;
                末段自~\textbf{795(*)}~续至~1053(*)。各段与原书一一对应;
        \end{itemize}
  \item 页边码“p.~1”“p.~A3”“p.~1066”式小注——这是作者引用自己打字稿原稿时
        所用的页码, 与本排印本页码不同。前奏四章的边码各有前缀: 第~1 章作
        \texttt{A1…A6}, 第~2 章作 \texttt{P1…P65}, 第~3 章作 \texttt{L1…L56},
        第~4 章作罗马数字 \texttt{I…XXII}, 自第一部起作纯数字 \texttt{1…171};
  \item 作者用以标记段落的~$\diamond$~记号, 以及标示后来补写文字的上标
        ~\textsuperscript{$\diamond$}~记号;
  \item 原书大小注(“注~44'”“注~46$^1$”“注~171³”等)的编号与层级。
\end{itemize}

第~12 章“第一部分的注”把第~5 至第~11 章各页脚注汇为一编, 并给每条注配上节题
(如“12.19. 耶稣与十二使徒”), 因此部分注文与正文脚注重复——这是原书的体例, 非
重出之误。这一部分属随笔式的自省文字, 同一句法语在正文脚注与章末汇编注中
各有措辞, 本册有意保留各自的译笔, 不强求统一。

除上述标记外, 译文力求贴近作者驳杂、跳跃而辛辣的文风; 数学名词的译法见各卷
术语表(如~\emph{topos}~不译, \emph{faisceaux pervers}~作“反常层”)。文中的日期
括注(如“(5~月~30~日)”“(1984年3月补写)”)为作者历次补写时的纪日, 一律照录。文中
另有一些对注号、页码的引用(如“n°~1695”“打字稿第~227~页”)沿袭作者打字稿的编号
系统, 与本排印本不尽一致, 亦无法在本书中落实, 悉照原文录出。

全套五卷(另有全书合订一册)的分工是:

\begin{itemize}
  \item \textbf{卷一·前奏}——主题的呈现: 或四个乐章的前奏(原书 p.23--122);
  \item \textbf{卷二·第一部}——愚妄与更新(原书 p.123--254);
  \item \textbf{卷三·第二部}——葬礼(I): 或中国皇帝的新衣(原书 p.255--412);
  \item \textbf{卷四·第三部}——葬礼(II): 或阴阳之钥(原书 p.413--632);
  \item \textbf{卷五·第四部}——葬礼(III): 或四种运算(原书 p.633--929)。
\end{itemize}
}
                       % 共用导言区
%      % ================================================================
%  《收获与播种》统一封面 —— TikZ 原生矢量绘制(不依赖任何外部图片)
%  母题: 暮色中的麦田 —— 犁沟向地平线的灭点收束, 前景麦穗剪影,
%        一位播种者把一把种子撒向天空, 种子沿弧线落入田里。
%  用法(在 \begin{document} 之前, 依次):
%      \input{common}                       % 共用导言区
%      \input{cover-art}                    % 本文件
%      \CoverSetup{prelude}{卷一}{前奏}%
%                 {主题的呈现——或四个乐章的前奏}{p.23–122}
%  正文中再 \makecover 即可。配色关键字:
%      prelude / part1 / part2 / part3 / part4 / all
%  五卷共用同一幅画, 只换配色与卷次文字, 以保证"一套书"的观感。
% ================================================================

\RequirePackage{tikz}

% --------------------------------------------------------------
%  配色: 六套(前奏 / 四部 / 全书), 各为同一暮色母题的变调
% --------------------------------------------------------------
\newcommand{\CoverPalette}[1]{%
  \ifstrequal{#1}{prelude}{%
    \definecolor{cvSkyTop}{RGB}{10,26,48}%
    \definecolor{cvSkyMid}{RGB}{46,84,116}%
    \definecolor{cvSkyBot}{RGB}{232,164,92}%
    \definecolor{cvSun}{RGB}{255,224,168}%
    \definecolor{cvFieldTop}{RGB}{150,110,52}%
    \definecolor{cvFieldBot}{RGB}{58,36,18}%
    \definecolor{cvFurrow}{RGB}{40,22,10}%
    \definecolor{cvSil}{RGB}{30,18,10}%
    \definecolor{cvCream}{RGB}{247,240,228}%
  }{%
  \ifstrequal{#1}{part1}{%
    \definecolor{cvSkyTop}{RGB}{22,70,100}%
    \definecolor{cvSkyMid}{RGB}{92,150,168}%
    \definecolor{cvSkyBot}{RGB}{240,200,104}%
    \definecolor{cvSun}{RGB}{255,244,200}%
    \definecolor{cvFieldTop}{RGB}{176,140,58}%
    \definecolor{cvFieldBot}{RGB}{84,54,22}%
    \definecolor{cvFurrow}{RGB}{45,28,12}%
    \definecolor{cvSil}{RGB}{34,22,10}%
    \definecolor{cvCream}{RGB}{250,245,232}%
  }{%
  \ifstrequal{#1}{part2}{%
    \definecolor{cvSkyTop}{RGB}{46,20,30}%
    \definecolor{cvSkyMid}{RGB}{128,52,54}%
    \definecolor{cvSkyBot}{RGB}{226,120,64}%
    \definecolor{cvSun}{RGB}{255,206,150}%
    \definecolor{cvFieldTop}{RGB}{132,84,44}%
    \definecolor{cvFieldBot}{RGB}{48,26,18}%
    \definecolor{cvFurrow}{RGB}{36,16,14}%
    \definecolor{cvSil}{RGB}{26,12,12}%
    \definecolor{cvCream}{RGB}{248,236,224}%
  }{%
  \ifstrequal{#1}{part3}{%
    \definecolor{cvSkyTop}{RGB}{30,20,54}%
    \definecolor{cvSkyMid}{RGB}{92,62,124}%
    \definecolor{cvSkyBot}{RGB}{216,132,140}%
    \definecolor{cvSun}{RGB}{255,214,196}%
    \definecolor{cvFieldTop}{RGB}{128,90,74}%
    \definecolor{cvFieldBot}{RGB}{44,26,30}%
    \definecolor{cvFurrow}{RGB}{34,20,30}%
    \definecolor{cvSil}{RGB}{22,14,22}%
    \definecolor{cvCream}{RGB}{246,238,238}%
  }{%
  \ifstrequal{#1}{part4}{%
    \definecolor{cvSkyTop}{RGB}{12,28,28}%
    \definecolor{cvSkyMid}{RGB}{44,86,72}%
    \definecolor{cvSkyBot}{RGB}{186,168,86}%
    \definecolor{cvSun}{RGB}{240,232,168}%
    \definecolor{cvFieldTop}{RGB}{110,104,50}%
    \definecolor{cvFieldBot}{RGB}{30,34,20}%
    \definecolor{cvFurrow}{RGB}{26,32,20}%
    \definecolor{cvSil}{RGB}{14,20,14}%
    \definecolor{cvCream}{RGB}{240,238,226}%
  }{%
    % 默认: 全书(麦浪金褐)
    \definecolor{cvSkyTop}{RGB}{28,40,32}%
    \definecolor{cvSkyMid}{RGB}{128,110,64}%
    \definecolor{cvSkyBot}{RGB}{232,178,88}%
    \definecolor{cvSun}{RGB}{255,236,180}%
    \definecolor{cvFieldTop}{RGB}{168,124,52}%
    \definecolor{cvFieldBot}{RGB}{62,40,20}%
    \definecolor{cvFurrow}{RGB}{42,26,12}%
    \definecolor{cvSil}{RGB}{28,18,10}%
    \definecolor{cvCream}{RGB}{250,242,226}%
  }}}}}%
}

% --------------------------------------------------------------
%  一键设定:  \CoverSetup{配色}{卷次}{卷名}{副题}{原书页码范围}
% --------------------------------------------------------------
\newcommand{\CoverSetup}[5]{%
  \CoverPalette{#1}%
  \newcommand{\VolNum}{#2}%
  \newcommand{\VolName}{#3}%
  \newcommand{\VolSub}{#4}%
  \newcommand{\VolPages}{#5}%
}

% --------------------------------------------------------------
%  一支麦穗:  \cvWheatEar{x}{底 y}{高}{倾角}
%  茎略弯, 麦粒成对左右斜出, 顶端一芒三针
% --------------------------------------------------------------
\newcommand{\cvWheatEar}[4]{%
  \begin{scope}[shift={(#1,#2)}, rotate=#4]
    \draw[cvSil, line width=0.85mm, line cap=round]
      (0,0) .. controls (0.5,{0.35*#3}) and (0.8,{0.72*#3}) .. ({0.55*#3},#3);
    \foreach \cvK in {0,...,6}{%
      \pgfmathsetmacro{\cvT}{0.50+0.50*\cvK/6}
      \pgfmathsetmacro{\cvGY}{\cvT*#3}
      \pgfmathsetmacro{\cvGX}{0.55*\cvT*\cvT*#3}
      \pgfmathsetmacro{\cvGS}{1.0-0.34*\cvT}
      \begin{scope}[shift={(\cvGX,\cvGY)}, rotate=28]
        \fill[cvSil] (0,{1.9*\cvGS}) ellipse [x radius={0.92*\cvGS}, y radius={2.3*\cvGS}];
      \end{scope}
      \begin{scope}[shift={(\cvGX,\cvGY)}, rotate=-28]
        \fill[cvSil] (0,{-1.9*\cvGS}) ellipse [x radius={0.92*\cvGS}, y radius={2.3*\cvGS}];
      \end{scope}
    }
    \begin{scope}[shift={({0.55*#3},#3)}]
      \fill[cvSil] (0,1.5) ellipse [x radius=0.78, y radius=1.95];
    \end{scope}
    \draw[cvSil, line width=0.26mm, line cap=round]
      ({0.55*#3},{#3+2.3}) -- ({0.55*#3-1.7},{#3+9.5});
    \draw[cvSil, line width=0.26mm, line cap=round]
      ({0.55*#3},{#3+2.3}) -- ({0.55*#3+2.1},{#3+10.5});
    \draw[cvSil, line width=0.26mm, line cap=round]
      ({0.55*#3},{#3+2.3}) -- ({0.55*#3},{#3+11.5});
  \end{scope}
}

% --------------------------------------------------------------
%  封面
% --------------------------------------------------------------
\makeatletter
\renewcommand{\makecover}{%
  \begin{titlepage}
    \thispagestyle{empty}
    \begin{tikzpicture}[remember picture, overlay, x=1mm, y=1mm,
                        shift={(current page.south west)}]
      % ===== 天空 =====
      \shade[top color=cvSkyTop, middle color=cvSkyMid, bottom color=cvSkyBot]
        (0,105) rectangle (210,297);
      % ===== 田野 =====
      \shade[top color=cvFieldTop, bottom color=cvFieldBot]
        (0,0) rectangle (210,105);
      % ===== 落日与光晕 =====
      \begin{scope}
        \clip (0,105) rectangle (210,297);
        \fill[cvSun, opacity=0.14] (152,132) circle (46);
        \fill[cvSun, opacity=0.24] (152,132) circle (30);
      \end{scope}
      \fill[cvSun, opacity=0.95] (152,132) circle (16.5);
      % ===== 犁沟: 向地平线灭点收束 =====
      \begin{scope}
        \clip (0,0) rectangle (210,105);
        \foreach \cvI in {1,...,20}{%
          \pgfmathsetmacro{\cvDX}{3.1*pow(\cvI,1.55)}
          \draw[cvFurrow, opacity=0.30, line width=0.32mm] ({105+\cvDX},0) -- (105,105);
          \draw[cvFurrow, opacity=0.30, line width=0.32mm] ({105-\cvDX},0) -- (105,105);
        }
        \foreach \cvY in {8,18,31,47,68,92}{%
          \draw[cvFurrow, opacity=0.15, line width=0.3mm]
            (0,\cvY) .. controls (60,{\cvY+3.5}) and (150,{\cvY-3}) .. (210,{\cvY+1.5});
        }
      \end{scope}
      % 地平线
      \draw[cvSil, opacity=0.45, line width=0.45mm] (0,105) -- (210,105);
      % ===== 播种者剪影(立于地平线) =====
      \begin{scope}[cvSil]
        \draw[line width=1.9mm, line cap=round, line join=round]
          (55.5,116) -- (50,105)  (55.5,116) -- (60.5,105.5);
        \draw[line width=2.3mm, line cap=round, line join=round]
          (55.5,116) -- (55,123.5);
        \draw[line width=1.6mm, line cap=round, line join=round]
          (55.2,123.2) -- (61.5,126.6) -- (68,129.4);
        \draw[line width=1.5mm, line cap=round, line join=round]
          (55.2,123.2) -- (51.4,119.8) -- (48.6,116.2);
        \fill (55.6,127.3) circle (2.35);
      \end{scope}
      % ===== 种子: 沿贝塞尔弧撒出 =====
      \draw[cvCream, opacity=0.13, line width=0.5mm]
        (68,129.4) .. controls (100,168) and (160,162) .. (196,104);
      % 掌心那一把(尚未撒出)
      \foreach \cvI in {1,...,5}{%
        \pgfmathsetmacro{\cvJX}{1.1*sin(\cvI*53)}
        \pgfmathsetmacro{\cvJY}{1.1*cos(\cvI*83)}
        \fill[cvCream, opacity=0.9] ({68+\cvJX},{129.4+\cvJY}) circle (1.0);
      }
      \foreach \cvI in {1,...,26}{%
        \pgfmathsetmacro{\cvT}{0.035+0.92*\cvI/26}
        \pgfmathsetmacro{\cvU}{1-\cvT}
        \pgfmathsetmacro{\cvSX}{\cvU*\cvU*\cvU*68 + 3*\cvU*\cvU*\cvT*100
                                + 3*\cvU*\cvT*\cvT*160 + \cvT*\cvT*\cvT*196}
        \pgfmathsetmacro{\cvSY}{\cvU*\cvU*\cvU*129.4 + 3*\cvU*\cvU*\cvT*168
                                + 3*\cvU*\cvT*\cvT*162 + \cvT*\cvT*\cvT*104}
        \pgfmathsetmacro{\cvR}{1.55-0.85*\cvT}
        \pgfmathsetmacro{\cvJX}{1.7*sin(\cvI*97)}
        \pgfmathsetmacro{\cvJY}{1.7*cos(\cvI*71)}
        \fill[cvCream, opacity={0.88-0.28*\cvT}] ({\cvSX+\cvJX},{\cvSY+\cvJY}) circle (\cvR);
      }
      % ===== 前景麦穗(两排) =====
      \foreach \cvI in {1,...,11}{%
        \pgfmathsetmacro{\cvWX}{4+19*(\cvI-1)+1.6*sin(\cvI*2.1)}
        \pgfmathsetmacro{\cvWH}{23+7*sin(\cvI*1.7)}
        \pgfmathsetmacro{\cvWL}{3.2*sin(\cvI*2.9)}
        \begin{scope}[opacity=0.45]
          \cvWheatEar{\cvWX}{30}{\cvWH}{\cvWL}
        \end{scope}
      }
      \foreach \cvI in {1,...,10}{%
        \pgfmathsetmacro{\cvWX}{9+20.4*(\cvI-1)+1.5*sin(\cvI*1.3)}
        \pgfmathsetmacro{\cvWH}{41+10*sin(\cvI*2.3)}
        \pgfmathsetmacro{\cvWL}{4.5*sin(\cvI*3.7)}
        \cvWheatEar{\cvWX}{30}{\cvWH}{\cvWL}
      }
      % ===== 压暗一层, 统一色调并保证文字可读 =====
      \fill[black, opacity=0.10] (0,105) rectangle (210,297);
      % ===== 底部信息带 =====
      \fill[cvSil, opacity=0.72] (0,0) rectangle (210,30);
      \draw[cvCream, opacity=0.40, line width=0.4mm] (0,30) -- (210,30);
      % ===== 文字 =====
      \node[anchor=north west, inner sep=0, text=cvCream] at (22,268)
        {\fontsize{38pt}{45pt}\selectfont\headfont \ziju{0.08}\BookTitleCN};
      \node[anchor=north west, inner sep=0, text=cvCream] at (22.6,242)
        {\fontsize{15pt}{20pt}\selectfont\itshape \BookTitleEN};
      \draw[cvCream, opacity=0.55, line width=0.5mm] (22,233) -- (120,233);
      \node[anchor=north west, inner sep=0, text=cvCream] at (22,224)
        {\fontsize{11pt}{15pt}\selectfont 一位数学家对往昔的反思与见证};
      \node[anchor=north west, inner sep=0, text=cvCream, opacity=0.78] at (22.4,208)
        {\fontsize{10pt}{14pt}\selectfont \ziju{0.14}\VolNum};
      \node[anchor=north west, inner sep=0, text=cvCream] at (22,199)
        {\fontsize{30pt}{37pt}\selectfont\headfont \ziju{0.1}\VolName};
      \node[anchor=north west, inner sep=0, text=cvCream, opacity=0.92] at (22,180)
        {\fontsize{13pt}{18pt}\selectfont \VolSub};
      \node[anchor=north west, inner sep=0, text=cvCream] at (22,19)
        {\fontsize{11.5pt}{15pt}\selectfont \OriginalAuthor\quad 著};
      \node[anchor=north west, inner sep=0, text=cvCream, opacity=0.85] at (22,7)
        {\fontsize{9.5pt}{13pt}\selectfont \Translator\ 译\quad ·\quad
         \ifdefempty{\TranslationDate}{\today}{\TranslationDate}};
      \node[anchor=north east, inner sep=0, text=cvCream, opacity=0.85] at (188,13)
        {\fontsize{10pt}{14pt}\selectfont 原书 \VolPages};
    \end{tikzpicture}
  \end{titlepage}%
}
\makeatother
                    % 本文件
%      \CoverSetup{prelude}{卷一}{前奏}%
%                 {主题的呈现——或四个乐章的前奏}{p.23–122}
%  正文中再 \makecover 即可。配色关键字:
%      prelude / part1 / part2 / part3 / part4 / all
%  五卷共用同一幅画, 只换配色与卷次文字, 以保证"一套书"的观感。
% ================================================================

\RequirePackage{tikz}

% --------------------------------------------------------------
%  配色: 六套(前奏 / 四部 / 全书), 各为同一暮色母题的变调
% --------------------------------------------------------------
\newcommand{\CoverPalette}[1]{%
  \ifstrequal{#1}{prelude}{%
    \definecolor{cvSkyTop}{RGB}{10,26,48}%
    \definecolor{cvSkyMid}{RGB}{46,84,116}%
    \definecolor{cvSkyBot}{RGB}{232,164,92}%
    \definecolor{cvSun}{RGB}{255,224,168}%
    \definecolor{cvFieldTop}{RGB}{150,110,52}%
    \definecolor{cvFieldBot}{RGB}{58,36,18}%
    \definecolor{cvFurrow}{RGB}{40,22,10}%
    \definecolor{cvSil}{RGB}{30,18,10}%
    \definecolor{cvCream}{RGB}{247,240,228}%
  }{%
  \ifstrequal{#1}{part1}{%
    \definecolor{cvSkyTop}{RGB}{22,70,100}%
    \definecolor{cvSkyMid}{RGB}{92,150,168}%
    \definecolor{cvSkyBot}{RGB}{240,200,104}%
    \definecolor{cvSun}{RGB}{255,244,200}%
    \definecolor{cvFieldTop}{RGB}{176,140,58}%
    \definecolor{cvFieldBot}{RGB}{84,54,22}%
    \definecolor{cvFurrow}{RGB}{45,28,12}%
    \definecolor{cvSil}{RGB}{34,22,10}%
    \definecolor{cvCream}{RGB}{250,245,232}%
  }{%
  \ifstrequal{#1}{part2}{%
    \definecolor{cvSkyTop}{RGB}{46,20,30}%
    \definecolor{cvSkyMid}{RGB}{128,52,54}%
    \definecolor{cvSkyBot}{RGB}{226,120,64}%
    \definecolor{cvSun}{RGB}{255,206,150}%
    \definecolor{cvFieldTop}{RGB}{132,84,44}%
    \definecolor{cvFieldBot}{RGB}{48,26,18}%
    \definecolor{cvFurrow}{RGB}{36,16,14}%
    \definecolor{cvSil}{RGB}{26,12,12}%
    \definecolor{cvCream}{RGB}{248,236,224}%
  }{%
  \ifstrequal{#1}{part3}{%
    \definecolor{cvSkyTop}{RGB}{30,20,54}%
    \definecolor{cvSkyMid}{RGB}{92,62,124}%
    \definecolor{cvSkyBot}{RGB}{216,132,140}%
    \definecolor{cvSun}{RGB}{255,214,196}%
    \definecolor{cvFieldTop}{RGB}{128,90,74}%
    \definecolor{cvFieldBot}{RGB}{44,26,30}%
    \definecolor{cvFurrow}{RGB}{34,20,30}%
    \definecolor{cvSil}{RGB}{22,14,22}%
    \definecolor{cvCream}{RGB}{246,238,238}%
  }{%
  \ifstrequal{#1}{part4}{%
    \definecolor{cvSkyTop}{RGB}{12,28,28}%
    \definecolor{cvSkyMid}{RGB}{44,86,72}%
    \definecolor{cvSkyBot}{RGB}{186,168,86}%
    \definecolor{cvSun}{RGB}{240,232,168}%
    \definecolor{cvFieldTop}{RGB}{110,104,50}%
    \definecolor{cvFieldBot}{RGB}{30,34,20}%
    \definecolor{cvFurrow}{RGB}{26,32,20}%
    \definecolor{cvSil}{RGB}{14,20,14}%
    \definecolor{cvCream}{RGB}{240,238,226}%
  }{%
    % 默认: 全书(麦浪金褐)
    \definecolor{cvSkyTop}{RGB}{28,40,32}%
    \definecolor{cvSkyMid}{RGB}{128,110,64}%
    \definecolor{cvSkyBot}{RGB}{232,178,88}%
    \definecolor{cvSun}{RGB}{255,236,180}%
    \definecolor{cvFieldTop}{RGB}{168,124,52}%
    \definecolor{cvFieldBot}{RGB}{62,40,20}%
    \definecolor{cvFurrow}{RGB}{42,26,12}%
    \definecolor{cvSil}{RGB}{28,18,10}%
    \definecolor{cvCream}{RGB}{250,242,226}%
  }}}}}%
}

% --------------------------------------------------------------
%  一键设定:  \CoverSetup{配色}{卷次}{卷名}{副题}{原书页码范围}
% --------------------------------------------------------------
\newcommand{\CoverSetup}[5]{%
  \CoverPalette{#1}%
  \newcommand{\VolNum}{#2}%
  \newcommand{\VolName}{#3}%
  \newcommand{\VolSub}{#4}%
  \newcommand{\VolPages}{#5}%
}

% --------------------------------------------------------------
%  一支麦穗:  \cvWheatEar{x}{底 y}{高}{倾角}
%  茎略弯, 麦粒成对左右斜出, 顶端一芒三针
% --------------------------------------------------------------
\newcommand{\cvWheatEar}[4]{%
  \begin{scope}[shift={(#1,#2)}, rotate=#4]
    \draw[cvSil, line width=0.85mm, line cap=round]
      (0,0) .. controls (0.5,{0.35*#3}) and (0.8,{0.72*#3}) .. ({0.55*#3},#3);
    \foreach \cvK in {0,...,6}{%
      \pgfmathsetmacro{\cvT}{0.50+0.50*\cvK/6}
      \pgfmathsetmacro{\cvGY}{\cvT*#3}
      \pgfmathsetmacro{\cvGX}{0.55*\cvT*\cvT*#3}
      \pgfmathsetmacro{\cvGS}{1.0-0.34*\cvT}
      \begin{scope}[shift={(\cvGX,\cvGY)}, rotate=28]
        \fill[cvSil] (0,{1.9*\cvGS}) ellipse [x radius={0.92*\cvGS}, y radius={2.3*\cvGS}];
      \end{scope}
      \begin{scope}[shift={(\cvGX,\cvGY)}, rotate=-28]
        \fill[cvSil] (0,{-1.9*\cvGS}) ellipse [x radius={0.92*\cvGS}, y radius={2.3*\cvGS}];
      \end{scope}
    }
    \begin{scope}[shift={({0.55*#3},#3)}]
      \fill[cvSil] (0,1.5) ellipse [x radius=0.78, y radius=1.95];
    \end{scope}
    \draw[cvSil, line width=0.26mm, line cap=round]
      ({0.55*#3},{#3+2.3}) -- ({0.55*#3-1.7},{#3+9.5});
    \draw[cvSil, line width=0.26mm, line cap=round]
      ({0.55*#3},{#3+2.3}) -- ({0.55*#3+2.1},{#3+10.5});
    \draw[cvSil, line width=0.26mm, line cap=round]
      ({0.55*#3},{#3+2.3}) -- ({0.55*#3},{#3+11.5});
  \end{scope}
}

% --------------------------------------------------------------
%  封面
% --------------------------------------------------------------
\makeatletter
\renewcommand{\makecover}{%
  \begin{titlepage}
    \thispagestyle{empty}
    \begin{tikzpicture}[remember picture, overlay, x=1mm, y=1mm,
                        shift={(current page.south west)}]
      % ===== 天空 =====
      \shade[top color=cvSkyTop, middle color=cvSkyMid, bottom color=cvSkyBot]
        (0,105) rectangle (210,297);
      % ===== 田野 =====
      \shade[top color=cvFieldTop, bottom color=cvFieldBot]
        (0,0) rectangle (210,105);
      % ===== 落日与光晕 =====
      \begin{scope}
        \clip (0,105) rectangle (210,297);
        \fill[cvSun, opacity=0.14] (152,132) circle (46);
        \fill[cvSun, opacity=0.24] (152,132) circle (30);
      \end{scope}
      \fill[cvSun, opacity=0.95] (152,132) circle (16.5);
      % ===== 犁沟: 向地平线灭点收束 =====
      \begin{scope}
        \clip (0,0) rectangle (210,105);
        \foreach \cvI in {1,...,20}{%
          \pgfmathsetmacro{\cvDX}{3.1*pow(\cvI,1.55)}
          \draw[cvFurrow, opacity=0.30, line width=0.32mm] ({105+\cvDX},0) -- (105,105);
          \draw[cvFurrow, opacity=0.30, line width=0.32mm] ({105-\cvDX},0) -- (105,105);
        }
        \foreach \cvY in {8,18,31,47,68,92}{%
          \draw[cvFurrow, opacity=0.15, line width=0.3mm]
            (0,\cvY) .. controls (60,{\cvY+3.5}) and (150,{\cvY-3}) .. (210,{\cvY+1.5});
        }
      \end{scope}
      % 地平线
      \draw[cvSil, opacity=0.45, line width=0.45mm] (0,105) -- (210,105);
      % ===== 播种者剪影(立于地平线) =====
      \begin{scope}[cvSil]
        \draw[line width=1.9mm, line cap=round, line join=round]
          (55.5,116) -- (50,105)  (55.5,116) -- (60.5,105.5);
        \draw[line width=2.3mm, line cap=round, line join=round]
          (55.5,116) -- (55,123.5);
        \draw[line width=1.6mm, line cap=round, line join=round]
          (55.2,123.2) -- (61.5,126.6) -- (68,129.4);
        \draw[line width=1.5mm, line cap=round, line join=round]
          (55.2,123.2) -- (51.4,119.8) -- (48.6,116.2);
        \fill (55.6,127.3) circle (2.35);
      \end{scope}
      % ===== 种子: 沿贝塞尔弧撒出 =====
      \draw[cvCream, opacity=0.13, line width=0.5mm]
        (68,129.4) .. controls (100,168) and (160,162) .. (196,104);
      % 掌心那一把(尚未撒出)
      \foreach \cvI in {1,...,5}{%
        \pgfmathsetmacro{\cvJX}{1.1*sin(\cvI*53)}
        \pgfmathsetmacro{\cvJY}{1.1*cos(\cvI*83)}
        \fill[cvCream, opacity=0.9] ({68+\cvJX},{129.4+\cvJY}) circle (1.0);
      }
      \foreach \cvI in {1,...,26}{%
        \pgfmathsetmacro{\cvT}{0.035+0.92*\cvI/26}
        \pgfmathsetmacro{\cvU}{1-\cvT}
        \pgfmathsetmacro{\cvSX}{\cvU*\cvU*\cvU*68 + 3*\cvU*\cvU*\cvT*100
                                + 3*\cvU*\cvT*\cvT*160 + \cvT*\cvT*\cvT*196}
        \pgfmathsetmacro{\cvSY}{\cvU*\cvU*\cvU*129.4 + 3*\cvU*\cvU*\cvT*168
                                + 3*\cvU*\cvT*\cvT*162 + \cvT*\cvT*\cvT*104}
        \pgfmathsetmacro{\cvR}{1.55-0.85*\cvT}
        \pgfmathsetmacro{\cvJX}{1.7*sin(\cvI*97)}
        \pgfmathsetmacro{\cvJY}{1.7*cos(\cvI*71)}
        \fill[cvCream, opacity={0.88-0.28*\cvT}] ({\cvSX+\cvJX},{\cvSY+\cvJY}) circle (\cvR);
      }
      % ===== 前景麦穗(两排) =====
      \foreach \cvI in {1,...,11}{%
        \pgfmathsetmacro{\cvWX}{4+19*(\cvI-1)+1.6*sin(\cvI*2.1)}
        \pgfmathsetmacro{\cvWH}{23+7*sin(\cvI*1.7)}
        \pgfmathsetmacro{\cvWL}{3.2*sin(\cvI*2.9)}
        \begin{scope}[opacity=0.45]
          \cvWheatEar{\cvWX}{30}{\cvWH}{\cvWL}
        \end{scope}
      }
      \foreach \cvI in {1,...,10}{%
        \pgfmathsetmacro{\cvWX}{9+20.4*(\cvI-1)+1.5*sin(\cvI*1.3)}
        \pgfmathsetmacro{\cvWH}{41+10*sin(\cvI*2.3)}
        \pgfmathsetmacro{\cvWL}{4.5*sin(\cvI*3.7)}
        \cvWheatEar{\cvWX}{30}{\cvWH}{\cvWL}
      }
      % ===== 压暗一层, 统一色调并保证文字可读 =====
      \fill[black, opacity=0.10] (0,105) rectangle (210,297);
      % ===== 底部信息带 =====
      \fill[cvSil, opacity=0.72] (0,0) rectangle (210,30);
      \draw[cvCream, opacity=0.40, line width=0.4mm] (0,30) -- (210,30);
      % ===== 文字 =====
      \node[anchor=north west, inner sep=0, text=cvCream] at (22,268)
        {\fontsize{38pt}{45pt}\selectfont\headfont \ziju{0.08}\BookTitleCN};
      \node[anchor=north west, inner sep=0, text=cvCream] at (22.6,242)
        {\fontsize{15pt}{20pt}\selectfont\itshape \BookTitleEN};
      \draw[cvCream, opacity=0.55, line width=0.5mm] (22,233) -- (120,233);
      \node[anchor=north west, inner sep=0, text=cvCream] at (22,224)
        {\fontsize{11pt}{15pt}\selectfont 一位数学家对往昔的反思与见证};
      \node[anchor=north west, inner sep=0, text=cvCream, opacity=0.78] at (22.4,208)
        {\fontsize{10pt}{14pt}\selectfont \ziju{0.14}\VolNum};
      \node[anchor=north west, inner sep=0, text=cvCream] at (22,199)
        {\fontsize{30pt}{37pt}\selectfont\headfont \ziju{0.1}\VolName};
      \node[anchor=north west, inner sep=0, text=cvCream, opacity=0.92] at (22,180)
        {\fontsize{13pt}{18pt}\selectfont \VolSub};
      \node[anchor=north west, inner sep=0, text=cvCream] at (22,19)
        {\fontsize{11.5pt}{15pt}\selectfont \OriginalAuthor\quad 著};
      \node[anchor=north west, inner sep=0, text=cvCream, opacity=0.85] at (22,7)
        {\fontsize{9.5pt}{13pt}\selectfont \Translator\ 译\quad ·\quad
         \ifdefempty{\TranslationDate}{\today}{\TranslationDate}};
      \node[anchor=north east, inner sep=0, text=cvCream, opacity=0.85] at (188,13)
        {\fontsize{10pt}{14pt}\selectfont 原书 \VolPages};
    \end{tikzpicture}
  \end{titlepage}%
}
\makeatother
                    % 本文件
%      \CoverSetup{prelude}{卷一}{前奏}%
%                 {主题的呈现——或四个乐章的前奏}{p.23–122}
%  正文中再 \makecover 即可。配色关键字:
%      prelude / part1 / part2 / part3 / part4 / all
%  五卷共用同一幅画, 只换配色与卷次文字, 以保证"一套书"的观感。
% ================================================================

\RequirePackage{tikz}

% --------------------------------------------------------------
%  配色: 六套(前奏 / 四部 / 全书), 各为同一暮色母题的变调
% --------------------------------------------------------------
\newcommand{\CoverPalette}[1]{%
  \ifstrequal{#1}{prelude}{%
    \definecolor{cvSkyTop}{RGB}{10,26,48}%
    \definecolor{cvSkyMid}{RGB}{46,84,116}%
    \definecolor{cvSkyBot}{RGB}{232,164,92}%
    \definecolor{cvSun}{RGB}{255,224,168}%
    \definecolor{cvFieldTop}{RGB}{150,110,52}%
    \definecolor{cvFieldBot}{RGB}{58,36,18}%
    \definecolor{cvFurrow}{RGB}{40,22,10}%
    \definecolor{cvSil}{RGB}{30,18,10}%
    \definecolor{cvCream}{RGB}{247,240,228}%
  }{%
  \ifstrequal{#1}{part1}{%
    \definecolor{cvSkyTop}{RGB}{22,70,100}%
    \definecolor{cvSkyMid}{RGB}{92,150,168}%
    \definecolor{cvSkyBot}{RGB}{240,200,104}%
    \definecolor{cvSun}{RGB}{255,244,200}%
    \definecolor{cvFieldTop}{RGB}{176,140,58}%
    \definecolor{cvFieldBot}{RGB}{84,54,22}%
    \definecolor{cvFurrow}{RGB}{45,28,12}%
    \definecolor{cvSil}{RGB}{34,22,10}%
    \definecolor{cvCream}{RGB}{250,245,232}%
  }{%
  \ifstrequal{#1}{part2}{%
    \definecolor{cvSkyTop}{RGB}{46,20,30}%
    \definecolor{cvSkyMid}{RGB}{128,52,54}%
    \definecolor{cvSkyBot}{RGB}{226,120,64}%
    \definecolor{cvSun}{RGB}{255,206,150}%
    \definecolor{cvFieldTop}{RGB}{132,84,44}%
    \definecolor{cvFieldBot}{RGB}{48,26,18}%
    \definecolor{cvFurrow}{RGB}{36,16,14}%
    \definecolor{cvSil}{RGB}{26,12,12}%
    \definecolor{cvCream}{RGB}{248,236,224}%
  }{%
  \ifstrequal{#1}{part3}{%
    \definecolor{cvSkyTop}{RGB}{30,20,54}%
    \definecolor{cvSkyMid}{RGB}{92,62,124}%
    \definecolor{cvSkyBot}{RGB}{216,132,140}%
    \definecolor{cvSun}{RGB}{255,214,196}%
    \definecolor{cvFieldTop}{RGB}{128,90,74}%
    \definecolor{cvFieldBot}{RGB}{44,26,30}%
    \definecolor{cvFurrow}{RGB}{34,20,30}%
    \definecolor{cvSil}{RGB}{22,14,22}%
    \definecolor{cvCream}{RGB}{246,238,238}%
  }{%
  \ifstrequal{#1}{part4}{%
    \definecolor{cvSkyTop}{RGB}{12,28,28}%
    \definecolor{cvSkyMid}{RGB}{44,86,72}%
    \definecolor{cvSkyBot}{RGB}{186,168,86}%
    \definecolor{cvSun}{RGB}{240,232,168}%
    \definecolor{cvFieldTop}{RGB}{110,104,50}%
    \definecolor{cvFieldBot}{RGB}{30,34,20}%
    \definecolor{cvFurrow}{RGB}{26,32,20}%
    \definecolor{cvSil}{RGB}{14,20,14}%
    \definecolor{cvCream}{RGB}{240,238,226}%
  }{%
    % 默认: 全书(麦浪金褐)
    \definecolor{cvSkyTop}{RGB}{28,40,32}%
    \definecolor{cvSkyMid}{RGB}{128,110,64}%
    \definecolor{cvSkyBot}{RGB}{232,178,88}%
    \definecolor{cvSun}{RGB}{255,236,180}%
    \definecolor{cvFieldTop}{RGB}{168,124,52}%
    \definecolor{cvFieldBot}{RGB}{62,40,20}%
    \definecolor{cvFurrow}{RGB}{42,26,12}%
    \definecolor{cvSil}{RGB}{28,18,10}%
    \definecolor{cvCream}{RGB}{250,242,226}%
  }}}}}%
}

% --------------------------------------------------------------
%  一键设定:  \CoverSetup{配色}{卷次}{卷名}{副题}{原书页码范围}
% --------------------------------------------------------------
\newcommand{\CoverSetup}[5]{%
  \CoverPalette{#1}%
  \newcommand{\VolNum}{#2}%
  \newcommand{\VolName}{#3}%
  \newcommand{\VolSub}{#4}%
  \newcommand{\VolPages}{#5}%
}

% --------------------------------------------------------------
%  一支麦穗:  \cvWheatEar{x}{底 y}{高}{倾角}
%  茎略弯, 麦粒成对左右斜出, 顶端一芒三针
% --------------------------------------------------------------
\newcommand{\cvWheatEar}[4]{%
  \begin{scope}[shift={(#1,#2)}, rotate=#4]
    \draw[cvSil, line width=0.85mm, line cap=round]
      (0,0) .. controls (0.5,{0.35*#3}) and (0.8,{0.72*#3}) .. ({0.55*#3},#3);
    \foreach \cvK in {0,...,6}{%
      \pgfmathsetmacro{\cvT}{0.50+0.50*\cvK/6}
      \pgfmathsetmacro{\cvGY}{\cvT*#3}
      \pgfmathsetmacro{\cvGX}{0.55*\cvT*\cvT*#3}
      \pgfmathsetmacro{\cvGS}{1.0-0.34*\cvT}
      \begin{scope}[shift={(\cvGX,\cvGY)}, rotate=28]
        \fill[cvSil] (0,{1.9*\cvGS}) ellipse [x radius={0.92*\cvGS}, y radius={2.3*\cvGS}];
      \end{scope}
      \begin{scope}[shift={(\cvGX,\cvGY)}, rotate=-28]
        \fill[cvSil] (0,{-1.9*\cvGS}) ellipse [x radius={0.92*\cvGS}, y radius={2.3*\cvGS}];
      \end{scope}
    }
    \begin{scope}[shift={({0.55*#3},#3)}]
      \fill[cvSil] (0,1.5) ellipse [x radius=0.78, y radius=1.95];
    \end{scope}
    \draw[cvSil, line width=0.26mm, line cap=round]
      ({0.55*#3},{#3+2.3}) -- ({0.55*#3-1.7},{#3+9.5});
    \draw[cvSil, line width=0.26mm, line cap=round]
      ({0.55*#3},{#3+2.3}) -- ({0.55*#3+2.1},{#3+10.5});
    \draw[cvSil, line width=0.26mm, line cap=round]
      ({0.55*#3},{#3+2.3}) -- ({0.55*#3},{#3+11.5});
  \end{scope}
}

% --------------------------------------------------------------
%  封面
% --------------------------------------------------------------
\makeatletter
\renewcommand{\makecover}{%
  \begin{titlepage}
    \thispagestyle{empty}
    \begin{tikzpicture}[remember picture, overlay, x=1mm, y=1mm,
                        shift={(current page.south west)}]
      % ===== 天空 =====
      \shade[top color=cvSkyTop, middle color=cvSkyMid, bottom color=cvSkyBot]
        (0,105) rectangle (210,297);
      % ===== 田野 =====
      \shade[top color=cvFieldTop, bottom color=cvFieldBot]
        (0,0) rectangle (210,105);
      % ===== 落日与光晕 =====
      \begin{scope}
        \clip (0,105) rectangle (210,297);
        \fill[cvSun, opacity=0.14] (152,132) circle (46);
        \fill[cvSun, opacity=0.24] (152,132) circle (30);
      \end{scope}
      \fill[cvSun, opacity=0.95] (152,132) circle (16.5);
      % ===== 犁沟: 向地平线灭点收束 =====
      \begin{scope}
        \clip (0,0) rectangle (210,105);
        \foreach \cvI in {1,...,20}{%
          \pgfmathsetmacro{\cvDX}{3.1*pow(\cvI,1.55)}
          \draw[cvFurrow, opacity=0.30, line width=0.32mm] ({105+\cvDX},0) -- (105,105);
          \draw[cvFurrow, opacity=0.30, line width=0.32mm] ({105-\cvDX},0) -- (105,105);
        }
        \foreach \cvY in {8,18,31,47,68,92}{%
          \draw[cvFurrow, opacity=0.15, line width=0.3mm]
            (0,\cvY) .. controls (60,{\cvY+3.5}) and (150,{\cvY-3}) .. (210,{\cvY+1.5});
        }
      \end{scope}
      % 地平线
      \draw[cvSil, opacity=0.45, line width=0.45mm] (0,105) -- (210,105);
      % ===== 播种者剪影(立于地平线) =====
      \begin{scope}[cvSil]
        \draw[line width=1.9mm, line cap=round, line join=round]
          (55.5,116) -- (50,105)  (55.5,116) -- (60.5,105.5);
        \draw[line width=2.3mm, line cap=round, line join=round]
          (55.5,116) -- (55,123.5);
        \draw[line width=1.6mm, line cap=round, line join=round]
          (55.2,123.2) -- (61.5,126.6) -- (68,129.4);
        \draw[line width=1.5mm, line cap=round, line join=round]
          (55.2,123.2) -- (51.4,119.8) -- (48.6,116.2);
        \fill (55.6,127.3) circle (2.35);
      \end{scope}
      % ===== 种子: 沿贝塞尔弧撒出 =====
      \draw[cvCream, opacity=0.13, line width=0.5mm]
        (68,129.4) .. controls (100,168) and (160,162) .. (196,104);
      % 掌心那一把(尚未撒出)
      \foreach \cvI in {1,...,5}{%
        \pgfmathsetmacro{\cvJX}{1.1*sin(\cvI*53)}
        \pgfmathsetmacro{\cvJY}{1.1*cos(\cvI*83)}
        \fill[cvCream, opacity=0.9] ({68+\cvJX},{129.4+\cvJY}) circle (1.0);
      }
      \foreach \cvI in {1,...,26}{%
        \pgfmathsetmacro{\cvT}{0.035+0.92*\cvI/26}
        \pgfmathsetmacro{\cvU}{1-\cvT}
        \pgfmathsetmacro{\cvSX}{\cvU*\cvU*\cvU*68 + 3*\cvU*\cvU*\cvT*100
                                + 3*\cvU*\cvT*\cvT*160 + \cvT*\cvT*\cvT*196}
        \pgfmathsetmacro{\cvSY}{\cvU*\cvU*\cvU*129.4 + 3*\cvU*\cvU*\cvT*168
                                + 3*\cvU*\cvT*\cvT*162 + \cvT*\cvT*\cvT*104}
        \pgfmathsetmacro{\cvR}{1.55-0.85*\cvT}
        \pgfmathsetmacro{\cvJX}{1.7*sin(\cvI*97)}
        \pgfmathsetmacro{\cvJY}{1.7*cos(\cvI*71)}
        \fill[cvCream, opacity={0.88-0.28*\cvT}] ({\cvSX+\cvJX},{\cvSY+\cvJY}) circle (\cvR);
      }
      % ===== 前景麦穗(两排) =====
      \foreach \cvI in {1,...,11}{%
        \pgfmathsetmacro{\cvWX}{4+19*(\cvI-1)+1.6*sin(\cvI*2.1)}
        \pgfmathsetmacro{\cvWH}{23+7*sin(\cvI*1.7)}
        \pgfmathsetmacro{\cvWL}{3.2*sin(\cvI*2.9)}
        \begin{scope}[opacity=0.45]
          \cvWheatEar{\cvWX}{30}{\cvWH}{\cvWL}
        \end{scope}
      }
      \foreach \cvI in {1,...,10}{%
        \pgfmathsetmacro{\cvWX}{9+20.4*(\cvI-1)+1.5*sin(\cvI*1.3)}
        \pgfmathsetmacro{\cvWH}{41+10*sin(\cvI*2.3)}
        \pgfmathsetmacro{\cvWL}{4.5*sin(\cvI*3.7)}
        \cvWheatEar{\cvWX}{30}{\cvWH}{\cvWL}
      }
      % ===== 压暗一层, 统一色调并保证文字可读 =====
      \fill[black, opacity=0.10] (0,105) rectangle (210,297);
      % ===== 底部信息带 =====
      \fill[cvSil, opacity=0.72] (0,0) rectangle (210,30);
      \draw[cvCream, opacity=0.40, line width=0.4mm] (0,30) -- (210,30);
      % ===== 文字 =====
      \node[anchor=north west, inner sep=0, text=cvCream] at (22,268)
        {\fontsize{38pt}{45pt}\selectfont\headfont \ziju{0.08}\BookTitleCN};
      \node[anchor=north west, inner sep=0, text=cvCream] at (22.6,242)
        {\fontsize{15pt}{20pt}\selectfont\itshape \BookTitleEN};
      \draw[cvCream, opacity=0.55, line width=0.5mm] (22,233) -- (120,233);
      \node[anchor=north west, inner sep=0, text=cvCream] at (22,224)
        {\fontsize{11pt}{15pt}\selectfont 一位数学家对往昔的反思与见证};
      \node[anchor=north west, inner sep=0, text=cvCream, opacity=0.78] at (22.4,208)
        {\fontsize{10pt}{14pt}\selectfont \ziju{0.14}\VolNum};
      \node[anchor=north west, inner sep=0, text=cvCream] at (22,199)
        {\fontsize{30pt}{37pt}\selectfont\headfont \ziju{0.1}\VolName};
      \node[anchor=north west, inner sep=0, text=cvCream, opacity=0.92] at (22,180)
        {\fontsize{13pt}{18pt}\selectfont \VolSub};
      \node[anchor=north west, inner sep=0, text=cvCream] at (22,19)
        {\fontsize{11.5pt}{15pt}\selectfont \OriginalAuthor\quad 著};
      \node[anchor=north west, inner sep=0, text=cvCream, opacity=0.85] at (22,7)
        {\fontsize{9.5pt}{13pt}\selectfont \Translator\ 译\quad ·\quad
         \ifdefempty{\TranslationDate}{\today}{\TranslationDate}};
      \node[anchor=north east, inner sep=0, text=cvCream, opacity=0.85] at (188,13)
        {\fontsize{10pt}{14pt}\selectfont 原书 \VolPages};
    \end{tikzpicture}
  \end{titlepage}%
}
\makeatother

\CoverSetup{part4}{卷五}{第四部}{葬礼(III)——或四种运算}{p.633–929}

\hypersetup{
  pdftitle={收获与播种·卷五·第四部(葬礼 III)},
  pdfauthor={Alexandre Grothendieck / 数译 译}
}

% --------------------------------------------------------------
\begin{document}

\frontmatter
\makecover

% ---------------------------------------------------------------
%  译者说明
% ---------------------------------------------------------------
\chapter*{译者说明}
\markboth{译者说明}{译者说明}
\addcontentsline{toc}{chapter}{译者说明}

本册是《收获与播种》中文译本的第五卷, 收原书的第四部“葬礼(III):或四种运算”
(Quatrième partie: L'Enterrement~(III) ou les Quatre Opérations), 自原书第
~633~页部扉起, 收第~18~章的小节~18.3“最后的义务(或那次来访)”、18.4“死之舞”、
18.5“四种运算(在遗骸上)”及其下各子目与大批按语, 并续至全书末行(原书第~929
页)。这是全书篇幅最大的一部, 也是“调查”部分的后半与本卷的收束。

本册脚注承前连续编号, 自~\textbf{795(*)}~起, 经~(*)~与~(**)~角标, 止于
~1053(*)。

\CommonNoteBody

\makecontents

\mainmatter

% ================================================================
%  第四部 (18.3 起, 原书 p.633-929)
%  —— 部扉在 part4-ch18 片段之首, 承第 18 章连续排印
% ================================================================
\markboth{收获与播种·第四部}{}

\setcounter{chapter}{18}
\setcounter{section}{2}

% 本卷不含 \chapter 命令(第 18 章起始在卷四), 需补一条章级目录条目, 使
% 18.3 以下各节在目录中仍归于第 18 章之下。第四部部扉在片段首行, 故临时
% 包一层 \partdivider, 让章条目排在部条目之后。
\let\mtOrigPartdivider\partdivider
\renewcommand{\partdivider}[3]{%
  \mtOrigPartdivider{#1}{#2}{#3}%
  \phantomsection
  \addcontentsline{toc}{chapter}{18\quad 葬礼仪式}%
}

\partdivider{第四部}{葬礼(III)\\[0.35em]{\normalsize 或四种运算}}{葬礼(III)或四种运算}

\section{3.~最后的义务(或那次来访)}

\subsection{(1) 义务已尽——或真话时刻}

\notehead{163}{}

\dmark{}(2月16日)恰好整整一个月前,我开始了一场即兴的反思\origpage{775}——它由读C.G.~Jung的自传引发。我本打算花上几天,把读后最初那几份强烈印象写下来——而到今天,这些印象我还没有巡完一圈!它们在阅读的进程中丰富、转变——凭这份阅读所触发的工夫、凭我写下读书注的工夫。我恰好来得及巡完Jung青年时代头四章所引发的印象——那几章出自Jung本人手笔。我正准备拿这些印象,与后几章所引发的、乍看并不总是相合的其他印象相对勘。可就在今天要着手时,我意识到:这条岔路(已近百页……)实在不该待在另一条“岔路”里——那条我唤作“阴阳之钥”的岔路,本身就已够长。(一个月前,我还以为这条岔路已近尾声\fnote{313}{(*)}{(3月26日)写这一行时,我还抱着一个印象:我正在起头的这条注,会属于“阴阳之钥”。直到随后几天,我才明白:反思的另一个阶段早已开始。“阴阳之钥”于是以前一条注“无穷的链——或传递(3)”(n\textsuperscript{o}~162$''$)告终。}。)诚然,我关于Jung的读书注,确确实实嵌在阴阳的辩证法里;它们也的确不经意间引我把许多先前仅仅掠过的东西讲清楚了——无论关于我的生活,还是关于生活本身。可我觉得,这仍不足以让人在一道括弧之内,再开一道尺寸如此骇人的括弧——而后一道括弧,本身又位于我自己的下葬那场漫长反思的末章、“葬礼仪式”之内。终于该拾起这场反思,把它进行到底了!

到头来,这些读书注我既不打算收进“阴阳之钥”,也不打算收进葬礼——它与葬礼只有一线相当微弱的牵连。可以把这些注看作一个例证:印证我曾在“表面与深处”与“书写的颂歌”(n\textsuperscript{o}s~101、102)等注里,以一般措辞试图表达的东西。我还在犹豫:是把它们作为第四部分收进《收获与播种》,还是把它们单独成文,放进《数学反思》第2卷\fnote{314}{(**)}{(3月26日)最终,这些读书注将构成《收获与播种》的(不是第四部,而是)第五部、也是最后一部;它多半会收入《数学反思》第3卷(而非第2卷),与另一些性质更数学的文本放在一起。关于葬礼的整套注——它们构成《收获与播种》写作中的“第三口气”,始于去年9月22日;我本想把这一套做成《收获与播种》的第三部分——将被分成两个不同的部分,分别名为“阴阳之钥”与“四种运算”,构成《收获与播种》的第三部与第四部。}\dmark{}诚然,关于Jung生平的这场反思——它实际展开的这副样子——确确实实是我追寻已逾一年的那场漫长反思\origpage{776}的一个不可分割的部分;于我,它就叫《收获与播种》——而我直接卷入其中,一如我在这些注的任何其他地方。因此,仅仅因为这场反思不曾打招呼、就在一场葬礼的半中间冒了出来,又“溢出”葬礼的中心主题稍多了一点,就把它的这一部分从《收获与播种》里分出去,那是做作。

眼下,我要借Jung自传反思里的这道断口,言归正传,把这场葬礼仪式——倘若有此可能——终于进行到底!

如今该到时候了:就我朋友皮埃尔去年10月来我这里的来访,作一份小小的汇报。10月21日那条注(“行为”,n\textsuperscript{o}~113)里,我提过他的到来——其时他已于前一晚携女儿娜塔莉(两岁)抵达。在客人们离开之后(10月25日那条注“失乐园”,n\textsuperscript{o}~116)我写道:“过几天,该到时候盘点一下这次来访给了我什么——一次我已不再指望的来访……”这些“几天”,变成了将近四个月——可我总算到了这里!

我本想就这次相遇作一份“现场”的讲述——它于我,是在发现葬礼、发现其实在与意义的这场冒险中的一个重要插曲。可这一回,我感到自己被一份缄默之虑拽住了:不忍把我的朋友此行留给我的全部印象——繁多而鲜活的——原样端出。诚然,把其中一份印象收进反思(12月26日那条注“否认(2)——或变形”,n\textsuperscript{o}~153),我并无这样的踌躇。但提及某个时刻对某位朋友有过的一份印象,与对那个确切的“时刻”作一份现场的描绘——弥散的印象骤然变得显豁、无可否认的那个时刻——这是全然不同的两回事。后者有点像趁一位朋友不觉得自己被观看的时刻给他拍一张照片,进而不曾确得他的同意便\dmark{}让它传看\origpage{777}。所以,我只限于给出这次来访留给我的几份印象,并且(一如《收获与播种》别处\fnote{315}{(*)}{不过有一个例外——即“玩笑——或‘权复形’”那条注(n\textsuperscript{o}~83)里,我在一次电话交谈中给J.-L.~Verdier拍的那张“照片”。我还记得,为对那桩小场景作“现场”的描绘,我不得不压下心里的一份犹豫——我有点儿觉得,自己给我的昔日学生设了个圈套,这全然不是“我的风格”。当然,他扯满风帆一头扎进了这个圈套——那还是最大、最显眼的圈套之一——我又满心欢喜、对自己满意得很。活该!})绝不拍那些不谨慎的照片!

首先我得把这次来访\textbf{定位}一下。我原本的打算,是去皮埃尔家中看他\fnote{316}{(**)}{这一意向,我在“我的朋友们”那条注(n\textsuperscript{o}~79)的开头、及随后第一条脚注里表述过。},让他读《收获与播种》,包括葬礼。5月初,我曾写信给他,说我希望近期见他,给他读一篇文字——它首先是写给“数学世界里我昔日的朋友与昔日的学生们”的——在其中我“把自己整个放了进去”——“我不记得曾像照料这一篇那样照料过一篇文字”。我当时以为打字当月即可完工,便提议6月上半月去看他。最终,由于打字的延误,还不算给葬礼(如它当时所计划的,即大体上如今葬礼第一部分的内容)作最后润饰的工夫,我的来访一推再推;而7、8月间,皮埃尔不在法国。况且,对于那部我一意要亲手交给他、要他在任何他人之先阅读的作品,消息公布时他不曾显出分毫好奇。最终,6月间,我把《收获与播种》的第一部“愚妄与更新”寄给了他,想着:在把葬礼劈头砸给他之前,让他先有个了解——说不定我对我自己的反思会在他那里触响什么机关,触发点什么——谁知道呢!那时我已病倒十来天,反正近期去巴黎已无从谈起。

然而我急于让他读葬礼——在其中,皮埃尔是以触目惊心的方式被卷入的;我本想让他赶在度假出发之前,来我这里一读。怀着这样的心情,我在6月底把完整的《引言》连同葬礼的目录寄给了他——我以为这会给他一记冲击,他也会\dmark{}有心赶在出发之前来看我\origpage{778},细细了解我对那场出名的葬礼要说的话、以及在这场葬礼里分派给他的角色。可事实呢,直到8月底,我再没有得到过他的一点音讯——以致我疑心他是否收到了我的寄件。好一场悬念!他回来后的第二封信(署8月25日)里,终于对引言与目录说了几句——措辞在我看来闪躲之极。“我过去觉得,你不大知道你的‘孤儿们’曾被怎样的爱环绕着……”,他这样写道,并随信附上一份评注书目以为佐证——那分明是一份显见的善意,要驱散一桩他显然感受为令人痛心的误会。随后的下一封信(9月12日)里,他宣布10月7日将迁居普林斯顿,并说会设法在此之前抽空来我这里一趟。接着又音讯全无,我以为他已动身去了普林斯顿——并没有:往IHES打电话,我得知他的行期推迟了。而一周之后,正当我已不大指望短时间内见到他——他本人已活生生站在面前,还带着小娜塔莉!

(2月17日)相遇是在一种气氛里进行的——照一切表象,它是再平和、再友好不过的。一个凑巧在场的浮泛观察者,会指天誓日:皮埃尔正在啃一篇数学手稿,时不时把他的观察与批评——一个很“在行”的数学家的建设性意见——提交给我。对皮埃尔本人来说,事情本该不言自明:他是专程赶来的(出于对毕竟是他的“老师”的我的敬重),牺牲了一位大忙人宝贵的两天,为的是竭己所能,驱散一桩不幸溜进我心里去的、恼人的误会——天知道是怎样一场倒霉的机缘巧合促成了它。他的诚意与我的诚意,当然都在一切怀疑之上;连提也不必提起,事情就那么明摆着。而他的角色,是就一切物质性的细节之点给我以启发——那些在我的注里显得不甚分明之点,或我或许出错之点。他的阅读推进到哪里,他随手记下自己的观察,列成一张单子,在动身那天交给了我——我有这份好脑筋,当即用几个关键词把它好好记了下来。况且,他确确实实在两天里读完了葬礼I的大部,至少读完了(借目录、借正文内部的交叉指涉而定位的)一切直接关涉他本人的注。好一份\dmark{}业绩——要是想一想\origpage{779}:写这些注,我可是花了整整两个月的全时……

这两天里,小娜塔莉是最乖的孩子里最乖的。我几乎可以说不曾听见她的嗓音——无论是说话、喊叫还是啼哭。她在我家里似乎并不难受,却不怎么露面。至于她爸爸,他是十足模范的爸爸——要喂饭、要遛弯、要哄一位毫不苛求、绝不添乱的小姑娘睡觉,他总在需要的时刻随叫随到。他带着她来,他对我说,是因为筹备迁居普林斯顿的种种大事之后,妈妈忙着收拾家务,腾不出手再照管娜塔莉。但在这条实际而无可奈何的理由之外,我相信自己感到了另一条理由,它留在没说出口的地方:小姑娘的在场,给我朋友所畏惧的那场相遇的气氛,添了一笔柔和——纵然他或许连在心里也不愿承认这份畏惧。而这份在场,同时又是一个鲜活的、夺目的标志:表明他是在怎样的默契心境里,在迁居美国的手忙脚乱之中赶来的——一份昭然的诚意、一份同样显见的好意的心境。

至于我,我毫无搅扰我朋友的心思,要逼他触碰什么——他若感到愿意深入哪一个问题,我便随时奉陪。可偏偏,他最要紧的执着,是对我注里检视的众多情境中的任何一个,都决计不\textbf{进入实质}——在那些情境里,他作为数学家的操守(或干脆说他的操守)被明明白白打上了问号。一个听到我们谈话的旁观者——那谈话有时甚至岔进数学的讨论(我们之间,这已有三年多不曾发生了\fnote{317}{(*)}{关于Deligne与我之间一切数学交流的中断,见“两个转折”那条注(n\textsuperscript{o}~66)。}!)——决计猜不到:在我朋友评注的文本里,竟可能有什么东西,稍稍带上个人色彩地把他卷了进去。至于我,我感到我的朋友死死抓住这副虚构——那费劲维持着的“最佳世界里的昭然善意”的虚构。他小心翼翼避开一切可能让它爆裂的东西——凡是会显出:他想在我们之间、排除万难设立的那份默契“共识”根本不是实在,而恰恰是一场虚构、一根可供攀抓的“稻草”……

\dmark{}这两天里,我分明感到这情势有多假:那份平和温良、乐呵呵的门面\origpage{780}底下,压着多少焦虑。正应了那句话:悬梁者家里的绳子——人人心里想着,无人提起!可我到底还是就此说了一句——我想那是动身那天,午饭后。大意是说:我们这场相遇,竟走了茶馆客套的路数,我很有些透不过气;毕竟,在他正在读的这些注里,单在早他近四个月就该收到的引言里,我已经用相当分明、相当有力的措辞,谈到好几桩\textbf{出自他手}的\textbf{行为}。对这一点,他当真无话可答吗?他回答我,眼神模糊,挂着一丝惨淡的、近乎可怜的浅笑,说他在尽力“自保”——(就我记得的)却不说他究竟想从什么里头“自保”。无疑,我的追查,该被他感受为对一个生命的粗暴侵入——那生命至今于他想必是最安顺、最没有问题的,其中的一切想必甚至显得出奇地\textbf{驯服};驯服到这般地步,或许,以致他终于忘了:事情本可以是另一个样子。\textbf{承担}他自己置身其中的情势——就是说,仅仅与之直面,如其所是地检视它——那将是他对自己、对世界的视景的一场幅度巨大的颠覆,是自我那副僵硬结构的一场崩塌;大多数人宁愿死上千百回、(若做得到)把世界投入火与血,也不肯冒这样一个跃入未知的风险。正是为了躲开这一切,我的朋友执意(而且多半,今天依然执意)“自保”。

我本不该为此惊讶——这类剧本,我见过它重演成百上千次:那是面对事物实在的大恐惧,而尤在其上的,是面对内在更新之风险的大恐惧。我本不该惊讶,诚然;可每一次,我仍旧一次次惊讶——当我看见最刺目的显明遭到拒斥,当人们受着千般熬煎、也施人千般熬煎,只为躲开那个我凭可靠认识分明知道的东西:万般福祉中最大的福祉……

不管怎样,在我这番想“脱轨”的不成功尝试之后,谈话转了个弯便止住了。这两分钟,我想,是那两天里仅有的一次\textsuperscript{317(*)}——我们的谈话取了一份\textbf{个人的}转向\dmark{},或说出了什么,越过了那份“共识”的虚构——尽管相反的证据明摆着,那虚构仍被维持着\origpage{781}!我怕,一如常有的情形,这一回我身上那份深情却又直截了当的“圆融”——它本可以帮助我的朋友,给一种气氛祛魅——却不曾到场:那气氛,尽管表象相反,早已紧绷到极点,而且紧绷已有数月。我把自己限在家务、园子与写作的诸务之中,任我的朋友自去读书;共同进餐时也一样:在我之内,对我的年轻朋友,悬着一份无言的\textbf{期待}——一份\textbf{回答}的期待:对我借他手中那部文本所说话语的回答。这份期待,他不可能感不到——而且他心里清楚:他那几条可怜的物质性说明,可“回答”不了它!无疑,倘若我以这样那样的方式抢先把话挑明——哪怕从一场他活该受领的、精心的臭骂开头——反倒会让他松一口气:好歹在空无之处,立起一份\textbf{接触}。

诚然,在已往的十五年里,每当我试着同他提起一件切身的、于我紧要的事,我撞见的总是全然的沉默,或(若当面)那套规矩的惊诧语调——最纯粹的“天鹅绒的脚掌”风格。这个游戏,我自然不再想玩;自1981年的“转折”起,我早已离场,并无回头的意思\fnote{318}{(*)}{见已引过的“两个转折”那条注,n\textsuperscript{o}~66.}。但同样真的还有:这一回,在我们之间,分明有一个独一无二的“时刻”,它或许值得对一条规矩(或一种已成第二天性的习惯……)作一回通融——那条规矩说:不要去顶撞他人对谈论某事的推拒。有时候,“稍稍强按着头”或许是有益的(在一定的分寸之内)——就像带一个怕牙医(那怕是无来由的)的孩子去看牙医……

我说这些,可不是要怜惜我那位可怜的朋友皮埃尔——怪他没有从我这里得到他本可以指望的全部善意鼓励,诸如此类!毕竟,我有我的限度,一如人人都有,这很正常;况且,替那些把自己放进了某种情势(哪怕是身不由己)的人缓冲震荡——那情势迟早会以这样那样的方式落回他们头上——这未必是我的角色,更不是我的义务。

\dmark{}何况,10月22日晚,把皮埃尔与娜塔莉送到奥朗日车站之后,我毫无“白遇一场”、“错失良机”之感。我还不曾天真到指望有山呼海啸的收获——两个人就一桩把他们彼此深深牵连的问题谈到实质,这是何其稀罕!没有对话,这是明摆着的——可我仍感到自己学到了许多。已经有那些“物质性细节”了,其中不止一条相当有滋味,它们就某些发生过的“行动”的“脚本”问题、及其情境,给最后的~i~点上了最后的点。本注的续篇里我还会回到这里\fnote{319}{(*)}{见“把一个个~i~点清”那条注(n\textsuperscript{o}~164),它紧随本注之后。}。更要紧的是:在这两天里,我以新的眼睛观察了我的朋友——借我在葬礼反思里对他新学到的东西。可以说,我与他“重新认识”了一遍——在他与我的关系里、与物的关系里、与女儿的关系里……那一章仍是保留地——正是在这里,我今天这几条注开头所提到的那份自然的分寸,于我成为必须。

而在领会葬礼的视角里,还有另一条理由——比前两条更精微——使这场相遇的举行成为要紧。我相信,从我决意去巴黎会见我的朋友那一刻起,我就感到了这份要紧;可当时我多半说不出为什么,只知道:当有牵涉双方的大事之时,与当事人当面交谈,总是要紧的。可恰恰,我们不曾谈这些事——而我却有一种印象:关于葬礼的\textbf{实在},我尚该学的东西,我已经学到了。

我也可以这样说。在这场相遇之前,构成葬礼的那一整套情境、事实与举止,显得何其不可信、何其滑稽、何其谵妄:尽管几周几个月间积攒下来的、一切可触而无可否认的“证据”,尽管我已献给它的近三百页注——在我心底的某个地方,\textbf{我依旧无法相信}\fnote{320}{(**)}{面对我们健全官能的见证——当它们过于猛烈地撼动通行共识、或我们珍爱的看法时——这份\textbf{不相信},在“中国皇帝的新衣”那条注(n\textsuperscript{o}~77$'$)里已被唤起过。显然,写那条注,于我是抵达(至少部分)这份面对显明的不相信的手段:把手指点在这份积习的反应上。可这样做时,我把自己与这份不相信拉开距离——把它呈现为凡俗众生(成年人)的不相信,而把自己认同于“相信自己眼睛见证的小小孩”(“纵然他所见的相当闻所未闻、从未见过、被所有人无视与否认”)。这多半正是我写那条注时无意识的用意:与一份对我自己官能的不相信态度拉开距离,与一份“随大流”的群居本能拉开距离。这样的态度与本能,确确实实存在于我身上,一如在每个人身上;可(一如所有人)它们多半停留于无意识。于是,这像是把那使我自我疏离的东西从我心里驱除的一场尝试——而这尝试的主要结果,我相信,是把我想与之拉开距离的东西\textbf{压进了}无意识更深处。那份阴险的怀疑——它像一道暗藏的裂隙,在我对事物的认识里运作——并没有因此消除;那份倒霉的不相信,也谈不上“超越”(“至少部分”,原文如此)!\par 我再次发觉:在反思的那个时刻,反思还停留在低于我所称“禅观”的地方——禅观,是这样一场反思:其中那些幽暗而精微的内在运动(如这份隐秘的不相信,以及我写那条注时真正的动机——“驱除”这份碍事的不相信)始终处在一个警醒关注的对象之下。}!

而这样的事临到我,也\dmark{}并非头一遭——远非如此\origpage{783}:一份执拗的怀疑盘踞一时,那是抵抗旧物入档的残余——一副旧的事物的视景,它往往更舒适些,或更合于通行共识,比起接替它的那一副。有时,这份怀疑也不只是惰性的表达——对事物视景中创造性改变的惰性——它同样映出一个健康的成分、旧视景里\textbf{实在}的一面:那一面,或许曾被连同其余,过急地抛下了船!不管怎样,一如怀疑现身的每一次,该做的好事是意识到它(这并不总那么容易,鉴于那根深蒂固的、要“封住”不受欢迎的怀疑的反射),而做了这个,再细心检视它。我不记得有哪一次:我留心检视过一份怀疑,却不曾从中学到些有趣的东西(甚或,对我重要的东西),而且足以让一切怀疑烟消云散\fnote{321}{(*)}{更确切的说法是:怀疑转变成了\textbf{认识},认识取代了它的位置。这与人们驱赶(或“超越”!)一份怀疑时所发生的,毫无共同之处——后者使怀疑从目光里消失,而它其实已退隐(或遭流放……)于更深的、看不见的层。它离被解决(并转变为认识)比以往更远;它一如既往地运作着,像一道暗藏的\textbf{裂隙},一份不安,一桩仍被回避的工夫的信号。可与上一条脚注的评注对读。}。任何怀疑,都是一桩有待去做的工夫的无可置疑的信号。

\dmark{}就本个案而言——即我那份不曾言出的、全然无理性的怀疑:怀疑一场所谓的“葬礼”\origpage{784}的实在本身——我该承认:在这场与朋友的相遇之前,我甚至不曾抵达一切工夫的头一个先决条件:我不曾真正意识到它。它停留在一个单纯弥散的不安的状态,而且不自道其名——错在我没有审问它!不安及其意义,我是事后才发觉的——恰在它凭这场与朋友的相遇消散之时。此外我相信:无论他采取怎样的态度,这效果都会产生——无论是一种殷勤合作的姿态,急于向我提供一切缺失的“物质性细节”(事情正是如此),还是相反地——不妨说——一种激烈的、或许狂怒的否认,否认那些最昭然的事实。无论哪种情形,葬礼的\textbf{心灵}实在都不可能不向我显形——这一次,凭直接的知觉(而非凭文献的“归纳”、凭我所知其他事实的印证等等):眼看着我的对坐者干脆\textbf{无视}“最佳可能世界”版本里那些荒诞绝伦的荒诞——而恰恰是那些荒诞的巨大,当初使我在心底深处怀疑过那场葬礼的实在!

只举一个例子:竟需要Deligne本人亲口告诉我,他确确实实是从Zoghman~Mebkhout本人嘴里得知“好上帝定理”的——却不愿在自己与Beilinson、Bernstein合写的文章里指涉他\fnote{322}{(*)}{见“当值无名之辈与好上帝定理”那条注(n\textsuperscript{o}~48$'$)与“不义——或归来的意义”那条注(n\textsuperscript{o}~75),以及紧随后者的那些注——它们与它一同构成第VII行列“研讨会——或Mebkhout的层与反常性”。},出于对Kashiwara的(!)顾虑——(作为门外汉)他拿不准两人在这条定理里各占多少\fnote{323}{(**)}{当然,在Beilinson、Bernstein与Deligne发展那所谓“反常层”(为免称它们“Mebkhout的层”)的形式主义的文章里,指涉Kashiwara的,与指涉Zoghman~Mebkhout的一样,全无——而它整个是从那位始终不被提名的Mebkhout的哲学出发的。况且,Kashiwara在好上帝定理(别名:Mebkhout)里的角色,Deligne比我清楚:Kashiwara的可构性定理,使Mebkhout得以定义一个函子——它从一个由“连续”系数(微分算子的复形)构成的三角范畴,导向另一个由“离散”(可构)系数构成的范畴——这件事,世上在他之前无人想过去做,更无人猜到会有一个范畴等价。这恰恰是我花了十年左右(1956--66)所发展的对偶形式主义里那“缺失的环节”;而我的那些上同调家弟子们——Deligne领头——在我1970年离开之后,争先恐后把它埋掉了。}——竟需要我亲耳听\dmark{}Deligne说出这样的话,才亲眼看见这一奇怪的组合:细节上的诚意,与实质上、要害处异乎寻常而昭彰的自欺\origpage{785}。我没有觉得有必要让我的朋友留意他那别致的做法(它在“牵线人”那条注(n\textsuperscript{o}~75$'$)里已被揭示——他明明读过它!):为了那个“本该在其中找到位置”的结果,他怎样安排,好显出:那条定理光彩的作者,不是别人,正是他(或至少,那篇显赫文章的三位作者之一)!对这桩怪事——那个我唤作“反常层会议”的会议,主要是在Mebkhout前些年所发展的工作与哲学的余波里成形的(对此Deligne并不摆出争辩的样子\fnote{324}{(*)}{Deligne只限于摆出对我的视景稍加细化的样子,说依他看,MacPherson的想法在1981年6月吕米尼会议(即所谓“反常层会议”)上的影响,比Mebkhout的还要大。我身在局外,无从就事论事地讨论这个;而这显见得是个细节之点,几乎减轻不了所发生之事的巨大。况且Deligne并不否认:没有Mebkhout前些年的开拓性工作、没有他在全然孤独之中发展的那套哲学,无论那场会议,还是它所标志的代数簇上同调理论的那场大规模更新,都不会发生。

我得以理解到:MacPherson那个关于簇的“相交上同调”的想法——它独立于Mebkhout的想法而发展——一直有些束之高阁,直到Mebkhout的“哲学”以一道新的、意想不到的光把它照亮(这是Deligne发现的)。那是Mebkhout层理论(被错称作“反常层”,顶替了某个会议……)的全速起步。这个起步,是那场会议的主要\textbf{事件},而且(似乎)是我们理解代数簇上同调的历史上的一个转折。这份新理解的拱顶石,在我看正是那道“好上帝定理”——它自六十年代初便“悬在空中”,而我、以及后来的Deligne,都不曾把它提炼出来。}),而他的名字却严丝合缝地不见于登在《Astérisque》上的会议文集\fnote{325}{(**)}{“严丝合缝地缺席”这个说法,照字面为真——至少对文集的第1卷(由引言与Beilinson、Bernstein、Deligne的文章构成),它是会议的主干部分。第2卷里两篇文章的书目(一篇出自Brylinski,一篇出自马尔格朗日)里,有两处对Mebkhout的“暂停!”式提及——没有一处涉及“好上帝定理”的发明权。}——他也没有任何解释可提。他看那情形,像看一场不巧的\textbf{巧合}:他与任何人都不负责任。总之,我所称的葬礼,在我朋友皮埃尔那里,化约为二十或三十来桩这样的“巧合”。

\dmark{}我在这里重新遇到一个我在他身上熟知的把戏——而且不止在他身上:一个装傻的把戏,带着世上最无辜的神情,笃定自己永远不会被抓个正着\origpage{786}。而确实,有些时日了,我不再浪费时间,去说服什么人(比如)相信:某些所谓的“巧合”并非单纯的巧合。有时候指一指显眼的东西是有用的;可这做过之后,再去说服什么人相信:它们确确实实是\textbf{东西},而非胡思乱想——“你们到那儿找什么呀!”——那就是浪费时间。想说服自欺,是浪费时间:它是自觉的还是不自觉的,一样;它摆出愚蠢的脸孔还是精明的脸孔——还是一样。

然而在我们的相遇里变了的,给我的朋友添上一份他竭力控制、竭力遮掩的焦虑的,是:这一回,这把戏不再局限于两人之间、无人得见的那点无后果的小运动——而且对手还是一个\textbf{亡故者}!这一回,牌摊在桌面上,\textbf{这是一场公开的戏}。赌局已开,赌的是那位大名鼎鼎的教团将要认领、将要容忍什么。诚然,十或十五年来,它已经容忍了许多、认领了许多;它或许还会继续,谁知道呢?一如我的朋友皮埃尔,它或许并不在乎再多二三十桩“巧合”……

(2月18日)10月22日晚,当我终于把皮埃尔与娜塔莉送到奥朗日车站,我觉得自己蠢透了。皮埃尔一副样子:照着他给自己定的日程,一板一眼、巨细无遗地尽完了他的义务;而我感到一阵郁郁的挫败——这场终于举行的相遇(谈论它已有数月),竟什么也不曾说出、什么也不曾触及。

天黑了,小家伙(坐在后座)该睡了——紧赶着开车到车站,得四十来分钟。有很长一段时间,谁也没说话。是我先打破沉默的——那份在我心里寻找出口的不满推着我;那不满,无疑是对我自己,胜过对任何他人。可这不妨碍:我此行就是存心要撩一撩我的朋友。我对他说:我还没跟自己想明白,要不要对斯普林格出版社提起一桩诉讼,逼它把那部海盗卷册SGA~4$\frac{1}{2}$\dmark{}从流通中撤下来\origpage{787}——它登在《Lecture Notes》丛书里\fnote{326}{(*)}{关于这部卷册,尤见“同伙”“白板”“绿灯”“颠倒”四条注(n\textsuperscript{o}s~63$'''$、67、68、68$'$)。}。说实话,连这念头当初是何时掠过我的,我也说不大清了;我这会儿把它随口抛出来,也算顺带探一探我的朋友(德语所谓“ihm auf den Zahn fühlen”——探他的口风)。他没怎么反应,说真的;与其说我在与他对话,不如说我在独白——重拾一根我早在4月或5月就丢下的线头。我随即发觉,顺着它想下去:一桩单纯的法律角力,说到底成不了什么气候——把SGA~4$\frac{1}{2}$从它眼下的题名与装帧底下撤出流通,这事情的意义,除非创议出自别人而非我:或出自斯普林格,或者——谁知道呢——更好,出自Deligne本人。我只得接着说:由他Deligne来做一个这样的公开姿态——说到底,是对某些针对我的行径的一种悔罪——这在我看算不得奢侈。它会给一片亟需净化的气氛消消毒!

我的朋友以音节应答我的独白,此处彼处点缀一二。他让人听出:把手里那一大堆SGA~4$\frac{1}{2}$的存货贱价处理掉,斯普林格未必乐意——我回敬说:换一副封面就够了,它有先例,毫无问题\fnote{327}{(**)}{那是我与斯普林格出版社的第一次不幸遭遇:它出版了Hartshorne的讲义(据一门我在其中发展了局部上同调形式主义的课),标明的作者却是Hartshorne。那是《Lecture Notes》的n\textsuperscript{o}~41卷“局部上同调”,不得不更换封面。当时斯普林格还有这份礼貌:为这场错认致歉,并勤勉地补救错误。此后的家风,已大不相同……},那花不了它几个钱。而就算它贱卖存货——《Lecture Notes》一千多种里的一种,亏点损益,算得了什么!更别提Deligne:假如他真在乎,那几百万旧法郎——填补亏空所必需的——他总拿得出来……

有些话我没有说出口,可它藏在话里(而且肯定被听出来了):这里头攸关的东西,其价值或许超过我们中哪一个的一两个月的薪水。可我到底还是说了:在这类事情里,头一位要紧的,不是看\textbf{怎样}去做一件事(或者反过来,把去做的\textbf{障碍}一条条数出来),而是首先把自己\dmark{}\textbf{想做}什么弄明白\origpage{788}。这一步做了,其余便成了后勤问题,自然“跟上”(当它肯“跟上”的时候)。怪我这位寡言的对话者不肯把真实的感受说破,我把这当作不言自明:他分明看得出,把一桩亟需“清扫”的情势“清扫”一下会是好事——只是他对能为此做什么拿不定主意,无疑是为了“保全面子”,诸如此类。其实我压根“没搔着痒处”!我最终察觉这一点,是在站台上等火车的时候。那是当Deligne回到这个话头,带一点讪讪的神气对我说:最终他宁愿由\textbf{我}去联系斯普林格,谈SGA~4$\frac{1}{2}$的事。显见得,他无意掺和其中,甚至此刻都不肯就这部著作的命运表个态——而它明明是他著述的(自然,还有我的“协作”\fnote{328}{(*)}{关于这一点,见已引过的“颠倒”那条注,n\textsuperscript{o}~68$'$。}。)。

直到那时我才明白:我们在路上时,我的那番思量,的确是一场独白——而对我的朋友皮埃尔,那桩“SGA~4$\frac{1}{2}$—SGA~5”的“行动”里或许有什么不太“合规”的东西,仍旧不明不白。而我偏偏在这万般话题里的这一个上岔了出去,为自己的不满寻找出口,这决非偶然。正是这一行动——它连着一场对一件美好作品的按部就班的屠杀,在那作品里,我放进了自己最好的一部分\fnote{329}{(**)}{见“屠杀”那条注(n\textsuperscript{o}~87)与紧随其后的两条注。}——给我的触动最深:凭其中一口气——暴力的气(在屠杀里)与安然放肆的气(对着被屠杀的东西)。而我再次被触动,凭那份做作(我在朋友身上知之甚稔):仿佛这一切与他毫不相干,他与我对这个那个所怀的“想法”毫不相干。

火车很快就要到了;而这是头一回:借一份终于浮出水面的\textbf{情绪},我要用几句话,进入一桩于我紧要之事的\textbf{实质}。要把此刻的感受当面说出,用不了多久。那是真的情感——一个人的体面之感受了伤,出手的是他所亲近、而那亲近者戏耍了他——这不再是边上那点半科学化的文学,可以拿支铅笔一丝不苟地随读随注。

\dmark{}他一时愣住了,却仍竭力撑住他那副不动声色的架子\origpage{789}。我只得对他说了类似这样的话:“那么,你觉得那是个好题名了——SGA~4$\frac{1}{2}$,暗示这是排在SGA~5\textbf{之前}的东西——而在那里,早在十一年前,你学到了那些直到今天你天天在用的数学!”他回答我,语气像在背书:他把它题作SGA~4$\frac{1}{2}$,只是为了标明一种\textbf{逻辑}依赖关系,而非时间在先。

就这样,我有幸亲耳听到、并出自当事人自己之口:这出“闹剧”——它大到何等地步,以至于当初我白纸黑字读到它时——先是在他的笔下(在“SGA~4$\frac{1}{2}$”里),随后是在Illusie笔下(在那部名叫SGA~5的卷册里,它“合乎逻辑地”跟在我那位前任的卷册之后……)——竟难以相信自己的眼睛所见证的东西!

我只得对他说:他与我一样清楚,SGA~5整个“站得住”,不需要任何先决、任何猜想,它在逻辑上、在任何意义上都不依赖那些后来的添加。我对他说话时,直直盯着他的眼睛,他回答时也如此。他以同样平板的嗓音重复他的功课:SGA~5在逻辑上确乎依赖SGA~4$\frac{1}{2}$——可我从他闪躲的眼睛里看见:事情究竟如何,他与我一样清楚。他的眼睛,身不由己地,比他的嘴诚实。

于是,在我们之间,“真话时刻”到底还是来了——可没有任何仪器,相机或录音机,能把它侦测出来。当时,唯有他与我知道正在发生什么。

火车在随后几分钟里进了站,我记得。反正,对那一天来说,彼此再没有什么可说的了。

\subsection{(2) 把一个个~i~点清}

\notehead{164}{}

(2月20–21日)为给Deligne上一次(去年10月)来我家的访问作一份回顾的收尾,我想在这里把他乐意就若干点提供的说明清点一遍——那些点,在我关于葬礼I的反思注里或模糊,或有误。这也将给我一个机会,补充另外若干说明——由Deligne所提供的那些引发出来的\dmark{}\origpage{790}。

\textbf{I~Motif(《Lecture Notes~900》卷)。}

1. Deligne向我说明:LN~900卷\fnote{330}{(*)}{关于这部“值得纪念的卷册”的详情,见“一个梦的回忆——或Motif的诞生”与“下葬——或新父亲”两条注(n\textsuperscript{o}s~51、52)。}的主旨,是在一个数域~$K\subset\mathbb{C}$上——$\mathbb{Q}$的有限扩张——发展一套交换的、Motif的类域论。换言之,即确定“$\bar{K}$在$K$上的Motif Galois群,使之交换”。就此我记起:率先提出这个问题的正是我(事出有因!),在六十年代末。对一个选定的Motif概念,借“Motif范畴上的Betti自由函子”、借$K$在复数域$\mathbb{C}$中的给定嵌入,这个问题有一个确切的意义。其实,我当时给自己提出的问题还更一般些:确定“$\bar{K}/K$的\textbf{亚交换}Motif Galois群”——由完整的Motif Galois群导出:使之交换的,不是整个pro-代数群,而只是它的单位元分支。该当得到~Gal($\bar{K}$/$K$)~的一个全然典范的扩张:以诸乘法群~$L^*$——$L$取遍~$C/K$的有限子扩张——所对应的那些环面取投影极限而成的pro-环面为核。我记得Serre对这个问题极感兴趣,可无论他还是我(连同Deligne——我当然也把他拉了进来)都没能即席给出一个说得过去的“候选者”。这个问题随后坠入全然的遗忘,一如它所由出的Motif的yoga。这份沉默直到1979年才被Langlands的文章打破(Deligne在他1984年5月28日的信里,以一份Motif评注书目向我提示了它\fnote{331}{(**)}{即R.P.~Langlands的文章“Automorphic Representations, Shimura Varieties and Motives. Ein Märchen”(Corvallis会议),载Proc.~Symp.~Pure~Math.~33~(1979), AMS, vol.~II, pp.~205--246.}),在那篇文章里,我的Motif Galois群的想法头一次在文献里得到显式的表述。由于我无缘收到这篇文章的单行本,我不知道文中是否提到我这号人。Motif在文献里的下一次露面,似乎就是LN~900——在那里,凡指向我这个人、说我与该卷的主题及其主要问题有任何干系的提及,一概不见\fnote{332}{(***)}{(4月8日)我近来得知:Motif被用于Deligne1979年的一篇文章里(它登在与前一条脚注所引Langlands文章相同的那一卷里)。
(5月12日:这条“结尾”成了子注“先期起棺”(n\textsuperscript{o}~168~(iv)))。}。

2. Deligne向我说明:与我此前所认出的\dmark{}(凭某种“文风家风”……)相反\origpage{791},LN~900里Deligne--Milne的那篇文章——它从头重述了N.R.~Saavedra所发展的塔纳卡范畴的Galois理论——几乎全部出自Milne之手\fnote{333}{(*)}{关于Deligne--Milne的这篇文章,见“下葬——或新父亲”那条注(n\textsuperscript{o}~52),亦见后一条注“白板”(n\textsuperscript{o}~67)里的评注。}。Deligne还向我解释了Saavedra著作里的那个错误:它使人(若想拥有纤维函子的Galois--Poincaré理论的形式主义)不得不强化Saavedra对所谓“塔纳卡”范畴的定义。Deligne--Milne文章里的工作,确实仅限于做这个调整——错误一经发现,调整便是明摆着的。这倒顺带提出了一个极有趣的问题:对那些“真正的”塔纳卡范畴的$\otimes$-范畴,给出一个可操作的内部刻画(不妨更醒目地称作$\otimes$-\textbf{Galois--Poincaré范畴},因为正是为它们,才能发展一套Galois--Poincaré广群的理论\fnote{334}{(**)}{“(Galois--Poincaré)广群”这一称呼,好处在于暗示它与拓扑空间或topos的基本广群概念的近亲关系。不过按技术讲,‘(Galois--Poincaré)gerbe’的称呼会更贴切。所涉的是“纤维函子”的gerbe:它不只定义在所考虑的$\otimes$-范畴的基域~$k$上,而且定义在~$k$上概形的fpqc~site的任意对象上(对这site中形如~Spec($k'$)~的对象——$k'$为~$k$的扩张,乃至~$k$的\textbf{有限}扩张——予以特别的关注)。})。这个问题在该文中未曾触及,而且至今未得到令人满意的解决。显见得,那里要紧的不是提出或解决有趣的数学问题,而恰恰是给Saavedra的文章备下一份替身文献。(关于这一点,见“白板”那条注(n\textsuperscript{o}~67)的结尾。)\fnote{335}{(***)}{(5月12日)近来读到了所引的Saavedra那部书,如今看来:这部书,连同这个概念的名字(“塔纳卡范畴”)——那概念系我约1964年引入,并给了这部书它的书名——是一场\textbf{骗局}。我将在“第六颗钉子(钉进棺材)”那组注(n\textsuperscript{o}s~176$_{1}$至~176$_{7}$)里,对它作一番细密的拆解。}

3. 在葬礼I里,我不止一次强调:Deligne六十年代末所发展的Hodge-Deligne理论,只是迈向“$\mathbb{C}$上有限型概形上的‘Hodge-Deligne系数’”理论、及为此种系数的“六种运算的形式主义”的头一步。我曾(并且依然)深信:倘若不是Deligne存心作对——冲着我引入的某些强力观念(如六种运算的形式主义)——Hodge-Deligne理论今天早该“瓜熟蒂落”。Deligne强调:单是在~$\mathbb{C}$上有限型概形上定义一个Hodge-Deligne系数的范畴,就撞上了严峻的困难\dmark{}\origpage{792},而他自认无力克服。(他愈是如此,就愈该从理论一开始,就把这个问题明白地表述出来,连同与它休戚与共的、为此种系数的六种运算的形式主义——而Deligne始终避免这样做。)照他的说法,Mebkhout的视角与Mebkhout的层\fnote{336}{(*)}{这些层,即Deligne当年以“反常层”之名引入的那些。(关于这一点,见“不义——或归来的意义”与“反常性”两条注,n\textsuperscript{o}s~75、76.)他并不别扭,在我们的交谈里乐意把它们称作“Mebkhout的层”……}该当提供一条通往正确定义的进路。(而倘若没有这份存心作对,Deligne当然不会等Mebkhout去发展那套哲学——它逆着他的前辈们而行——并拿它去做那件显见得根本性的工作:十五年来它一直撂在那里,直到今天,除了我在《收获与播种》里亲手指出,文献中连提都无人提起!)

4. 我记错了:我曾以为,一个Motif的“按权过滤”是我引入的——它(对每个~$\ell$)反映为该Motif的$\ell$-adic实现上相应的过滤(以Frobenius特征值的绝对值为项定义的过滤)。其实,Deligne提醒我:我只以“虚拟”的权概念工作过(那等于以虚拟的Motif工作——某个合适的“Grothendieck群”的元素……)。这个重要的事实——我工作时所用的虚拟概念,该当对应一份典范的过滤,按“权递增”——是Deligne发现的\fnote{337}{(**)}{使Deligne相信一个Motif必有这样一份过滤(且必唯一)的启发式理由是:存在阿贝尔簇被环面作的非平凡扩张(其Motif的~$H^1$由此给出一个权~2~的Motif被权~1~的Motif作的非平凡扩张),但反方向不存在。这听来单薄——可我自己也差不多当场就被说服了:它太诱人,不可能是假的!更严肃的一条理由,在来自有限型域~$K$上Motif的$\ell$-adic表示的层面上,是去证明:一个权~i~的Galois模被另一个权~j~的Galois模作的扩张,当~i~<~j~时必为平凡。这个论断——它将证明与一个Motif相关联的$\ell$-adicGalois模(已是相当贴近Motif本身的对象……)存在一份典范的“按权递增”过滤——究竟当年是Deligne还是我证出来的,我已记不得了。}。这一发现——它一如“猜想性的Motif理论”,同样只是“猜想”——随即供出了那把钥匙:由此可在复数域上,给\textbf{Hodge-Deligne结构}(又称“混合Hodge结构”)作出成形的定义,把它当作Motif之上、及其Hodge实现之上那些“早已知道”的结构的“Hodge式”转写。

\dmark{}\origpage{793}按技术讲,我的想法在Hodge-Deligne结构的定义里的影响是双重的。一方面,经由一个Motif的权的概念——Deligne把它恰当地精确化为一副“\textbf{按权过滤}”的结构。另一方面,自五十年代起,我就一直强调:一张光滑代数簇~$X$的代数\textbf{De Rham上同调}——$X$不必真——之重要:它是一个比朴素的Hodge上同调(诸~$H^q(X,\Omega^p)$的直和)更丰富的不变量;后者借那条著名的谱序列与前者相连,谱序列关联于De Rham上同调的一份典范过滤(\textbf{De Rham过滤})。我是第一个定义代数De Rham上同调的人(在当年,谁也不曾想过去看一个微分算子复形——如De Rham复形——的整体超上同调),也是第一个坚持它的分次\textbf{过滤}结构的人——以与Hodge上同调的双分次结构相抗;自Hodge以来,双分次结构一直站在前台。在$X$真的情形(即拥有Hodge理论的情形,该理论蕴涵上述谱序列在特征零下退化),且在基域$\mathbb{C}$上,人们取这份过滤与复共轭过滤二者的“交”,便从De Rham上同调的过滤结构复原出双分次结构(借De Rham上同调的“实结构”——它同构于Betti上同调~$H^*(X,\mathbb{C})$)。随后我证明了(其时除我之外,还无人相信非真情形的De Rham上同调):对复数域上一张光滑概形~$X$,De Rham上同调(它有一个“纯粹代数”的意义)典范同构于复Betti上同调(由超越途径定义)。

话虽如此:一旦设定——$\mathbb{C}$上存在一个Motif概念(不必半单),以及一个$\mathbb{C}$-概形~$X$(当然不必真)的Motif上同调,以及一个$\mathbb{C}$上Motif的“Hodge实现”概念(合用的、有待找到的)——它(依我的想法)该当给光滑~$X$的Motif上同调配上一副“广义Hodge结构”(有待定义),以De Rham上同调~$H_{DR}(X)$为其底集合;那么,人们在这份De Rham上同调上读出的最初两副结构——即De Rham过滤(我早在五十年代引入)与按权过滤(Deligne引入,出自我关于虚拟权的想法,并细化了Serre的想法——那些想法本身出自Weil猜想)——恰恰就落在Deligne所引入的“混合Hodge结构”这个概念上。

\dmark{}\origpage{794}当然,这层想法的亲缘(164$_{1}$),Deligne知道得一清二楚。按行业的伦理(那份伦理,我不曾成功地传给他),他本该在他引入混合Hodge结构的那篇工作里把它明白指出\fnote{338}{(*)}{即文章“Théorie de Hodge II”(Pub.~Math.~IHES~40~(1971), pp.~5--58)。相反,在1970年尼斯大会上“Hodge~I”的摘要公告里,Serre与我倒被并在同一行里提及——一如我在“牺牲者”那条注(n\textsuperscript{o}~78$'$,第308页)里指出的。关于这方面的评注,见该注的子注~n\textsuperscript{o}~$78'^{1}$、$78'^{2}$。}。他却宁可在这篇工作里对它讳莫如深——那篇工作,也正是\textbf{他的博士论文};一如他判定,在这一特定场合,连那个曾是他师父的人的名字,也一并讳莫如深为宜。

5. 在那份关于Motif的评注书目(附于他去年8月25日的信)里,Deligne指明:“人们[!]之所以犹豫要不要在它上面构造——即在关于Motif的那几份‘经典文本’\fnote{339}{(**)}{即Kleiman、Manin、Demazure关于Motif的那几份零星的(‘经典’)文本,出版至1970年为止。它们几乎不出Motif的初始观念,对于我所发展的那份‘yoga’的精微——我曾试着把它传给愿意听的人——给不出丝毫印象。尤其,其中对Motif Galois群只字未提——而它恰恰是发展这份yoga的一条初始的根本动机。(见‘一个梦的回忆——或Motif的诞生’那条注,n\textsuperscript{o}~51.)}——之上构造——原因之一,是其中对代数闭链存在性猜想的使用:那些猜想并无真实的凭据;而Motif,它们,于我是无可怀疑的。”

对这份说明,我将回答:这些“经典文本”丝毫代表不了六十年代末的“技艺之现状”,差得还远;而Deligne学得这份“现状”,靠的也不是这些文本!他心里十分清楚:我的“标准猜想”只是诸多可能进法之中的一种——为在域上给一个(半单的)Motif概念作一份临时的“成形构造”——它丝毫不限制他得自我的那些想法的分量与内在动力。(关于这一点,见“一个梦的回忆——或Motif的诞生”那条注(n\textsuperscript{o}~51)的子注~n\textsuperscript{o}~51$^1$。)一手图二利:我离开之后,他一面竭力把标准猜想败坏成“无从下手”而乏味的东西,一面又竭力败坏某一种Motif的进法——仿佛那进法是我的,并且仿佛它代表一条死胡同,(照他的说法)与那些无望的猜想拆不开地绑在一起——以至于,在LN~900那部总算做了真正该做的工作的卷册里,把我的名字\dmark{}\origpage{795}知趣地讳莫如深,反倒算是对我仁慈……\fnote{340}{(*)}{还在他住在我家的头一天,Deligne就抢在任何我可能就此提出的问题之前,带着他最漂亮的微笑对我说:“你当真以为,并不是人人都已经知道,Motif是你引入的吗!”。而令人惊讶的恰恰是:不管我的朋友为让人忘掉它做了多少,我得以确认:这至今仍是人所共知的。可由于我的想法没有成文的文献,Deligne便大有余地去制造一种印象:仿佛我的贡献照例不过提出过一个模糊的一般观念(何况按它本来的样子根本无法使用,因为它依赖那些“今天与当年同样无从下手”的猜想……)——模糊到这般地步,以致它实在配不上一位做真功夫的严肃数学家,劳神为它哪怕作一次纯粹装点门面的引用……}。

6. 在同一份“评注书目”里,我读到:

\begin{quote}
从这个“经典”的观点看\fnote{341}{(**)}{见倒数第二条脚注。},文献里有一处令人惋惜的空缺:你对$\mathbb{F}_p$上Motif的塔纳卡$\otimes$-范畴的猜想式刻画——唯一等价,差非唯一的同构——连同那几个纤维函子(结晶的与$\ell$-adic的),参见Tate,~Classes d'isogénie des variétés abéliennes sur un corps fini,~Sém.~Bourbaki~352~(1968)。
\end{quote}

这是鳄鱼的眼泪:那处“令人惋惜的空缺”,除了我……之外,该由谁负责?正是我的朋友Pierre~Deligne本人——因为除我之外,世上知道那份“猜想式刻画”的数学家,恐怕唯有他一个……把它收进同一部LN~900,凑足分量,原只在他一念之间!况且,就我现在记得的,这份刻画毫无“猜想”可言——除了一点:得假定人们拥有一个所谓“$\mathbb{F}_p$上的Motif”的范畴,满足几条合理的条件——那是任何当得起这个名字的范畴都理应满足的。如果我没记错,所引Tate-本田的那条指涉蕴涵着:该范畴由Tate的Motif(及其逆)、与$\mathbb{F}_p$上定义的阿贝尔簇,乘法地生成。好东西是有一些的(我还略去许多),它们曾被我亲手交到那位辉煌的前学生手里,却被精心埋到了今天……

\textbf{II~平展上同调(“SGA~4$\frac{1}{2}$”、SGA~5、SGA~7、离散Riemann-Roch)。}

1. Deligne就葬礼I向我作的最初几条评论之一\dmark{},关乎一条猜想性定理的坎坷命运\origpage{796}——那条定理,是我在SGA~5里提炼出来的,名为“离散Riemann-Roch定理”。在“屠杀”那条注(n\textsuperscript{o}~87)的子注~n\textsuperscript{o}~87$^1$里,我对它作过相当细密的表述。Deligne向我指明:当他把我的猜想表述转达给MacPherson时,他自认扮演的是“邮差”、中介的角色。他没有给我的表述添任何新成分——把我的表述译成同调语言、从而使它对奇异空间也有意义的主意,出自MacPherson,而非Deligne。他对我说,当他收到MacPherson那篇文章的单行本——它以超越的论证,在解析-复的情形、在同调的语境里证明了我的猜想——他吃了一惊:猜想竟顶着“Deligne--Grothendieck猜想”的名字。他曾想过写信给MacPherson,纠正这个误会,可(连他自己也说不出为什么)最终没有写……

2. 与我所设想、所暗示的相反:在SGA~5口头讨论班进行之际,Deligne并未承担撰写该班一篇或几篇报告的承诺——例如关于结合于一个代数闭链的上同调类的报告(那一篇,他在讨论班十一年之后终于写成,收进由他编成的那部题为“SGA~4$\frac{1}{2}$”的卷册,不经任何手续\fnote{342}{(*)}{这一拆散SGA~5讨论班(还有许多其他)以肥了那部题为“SGA~4$\frac{1}{2}$”的卷册的行动,履行着两种职能,都朝着“颠倒”角色而去:把我弄成Deligne的“协作者”,并为“SGA~4$\frac{1}{2}$”对SGA~5的优先权主张(那主张,已由那个骗人的名字SGA~4$\frac{1}{2}$所暗示,又在引言里“字里行间”显豁出来——无论Deligne为SGA~4$\frac{1}{2}$写的引言,还是Illusie为SGA~5写的引言)提供支撑(在SGA~5里,经由那份被盗版的报告,指涉SGA~4$\frac{1}{2}$之处比比皆是)。关于这一点,亦见“颠倒”那条注(n\textsuperscript{o}~68$'$)里的评注——在那里,我终于识破那部海盗卷册所得的怪名字的意味,以及我的那篇代数闭链的报告何以在那卷册里现身。})。就此我提了一个问题:他是否不认为——能“现场”地在SGA~5里学到那些基本技术,是他此后全部著作的倚仗——这份特权,不给他强加一份\textbf{义务}或责任:尽力使这些技术经由SGA~5的迅速刊行,交到数学公众的手里。Deligne回答我:他\textbf{不这么认为}。关于Motif的哲学——它曾是他代数簇上同调(其著作的中心主题……)的主要灵感之源——同样的问题,我克制着没有问他。

\dmark{}\origpage{797}3. 正是Deligne主动向Verdier征询同意,要把Verdier关于导出范畴的工作里那份出名的“第0状态”收进“SGA~4$\frac{1}{2}$”。Verdier起初谢绝,认为那毫无意义(确切的措辞我记不得了)。最终是Illusie说服Verdier点了头。Verdier的第一反应,在我看来再自然不过,合乎单纯的数学常识。况且,导出范畴,Verdier多年前已决定埋葬——以一份大规模的逐件加工的形式,它有朝一日该构成他的博士论文——于是,发表一份大纲草图(它久已被文献大大覆盖)便显得滑稽突兀。我自信懂得Deligne与Illusie何以对发表这份“第0状态”如此执着——其中不提我的名字。至于Verdier何以收回他那最初合乎常识的反应,我自信有所感觉,并在“赊账的学位论文与全险”那条注(n\textsuperscript{o}~81)里谈到。

\par\noindent 4. 在“白板”那条注(n\textsuperscript{o}~67)里,我早已指出“这个讨论班”一语的歧义——它出现在SGA~4$\frac{1}{2}$的引言(第2页)里:“为着对L函数的应用,这个讨论班另含一个证明,它补足Frobenius态射的特殊情形。”这个歧义的说法,照上下文与文气,十有八九会被读作意指“SGA~4$\frac{1}{2}$”,从而暗示:母讨论班SGA~5并不含有L函数有理性的一份“完全”的证明。Deligne向我指明:在他心里,“这个讨论班”指的正是“SGA~5”。说实话,这份说明在我这里什么也没说明。我很清楚:SGA~5里有一个“完全”的证明——Deligne与我同样清楚——不过那是一条迹公式的证明,而且(与他所暗示的相反)它远远溢出“Frobenius态射的特殊情形”。但在Deligne的笔下,含混与歧义比比皆是,乃至昭然的假话,而它们全都指向同一个方向——这决非偶然:制造一种印象,关于我的著作、或Mebkhout及其他与我这个人相连者的著作,足以败坏它们的印象,同时抬举他自己的声望,或凭空捏造出这样的声望\fnote{343}{(*)}{尤其是暗示:Motif的枢纽想法、平展上同调、“好上帝定理”及与之相随的Mebkhout哲学,都是他的发明。(3月26日)就本个案与“这个讨论班”,亦见子注“双重含义——或行骗的艺术”(n\textsuperscript{o}~169$_7$)。}。

\dmark{}\origpage{798}5. 我借这个机会,就SGA~7~II(那部署为P.~Deligne与N.~Katz指导的讨论班)再添几条评注——关于它,我在那条(当时无名的\fnote{344}{(*)}{(3月26日)其间我已补上这个缺:把这条注收进目录,题为“一场屠杀的序曲”。})注~n\textsuperscript{o}~56里,已作过相当细密的表述。稍细一点的审视让我看到:在这一场里,N.~Katz没有客气,悄悄给那辆由Deligne利落驾着的灵车推轮子,而且方式很多。

Katz接受与Deligne一同署名为卷册与讨论班的合著者——这与四年前口头讨论班上实际发生的事,全然对不上。SGA~7讨论班的总体设计(它延续1967--69两个年头)出自于我;讨论班的招牌,是与Deligne共同指导。Katz在其中列名为一名合作者-报告人,另有若干位。但既然Katz接受了署名合著者(其中五篇报告出自他的手笔,但主要结果没有一条归于他),那么把卷册的整体操持、以及卷册里对我这个人的藏匿,也视他与Deligne同等的共同负责,便顺理成章。

我头一个想到的,是卷册引言(Deligne署名)里的藏匿:文中呈现的任何主题或结果,竟无一处让人想到与我有什么干系——而讨论班摆在突出位置的两条“关键结果”之一(即Lefschetz束的理论),早在我筹办SGA~7之前就由我发展,而且正是我设想办一个单值群主题讨论班的动机之一。在Katz呈现这一理论的报告(第XVIII讲)里——题名“Lefschetz束的上同调研究,作者N.~Katz”——我的名字不按惯例出现在题名里(“据A.~Grothendieck”),而是出现在N.~Katz名字之后一条简短的脚注里:“据Grothendieck的(简略)笔记”。那“简略”的考语,像是后加的,好冲淡这桩不巧的“Grothendieck笔记”在此扮演过角色的事实。可那些笔记纵然“简略”,它们代表的却是好几天工夫的结晶——面对一桩先验绝不显然的任务:把一些经由超越途径陈述并证明的结果,转写进一个技术上全然不同的语境。一如平展对偶、或一如Nielsen--Wecken的理论\dmark{}\origpage{799}\fnote{345}{(*)}{Deligne比他的朋友N.~Katz更少拘束,倒是压根不觉得值得提一句:在他所谓的“Nielsen--Wecken方法”里,我也有几分干系——关于这一点,见“白板”那条注(n\textsuperscript{o}~67)的子注~n\textsuperscript{o}~67$^1$。},经典的论证原样用不上,一切须从头重做:拿经典\textbf{结果}当引路的线头,而把它们传统的“证明”(如果那配得上这个称呼的话)整个忘掉。正常的倒是:即便有我那些细密的笔记相助,Katz也得费一番力气才能入水,一如在他之前我也曾费过——可这丝毫不意味着(至少,不按通行认可的游戏规则)他就是平展上同调里Lefschetz束理论的作者!

顺着这股劲头,在同一报告的引言(第225页)里,Katz摆出把Raynaud太太呈现为那定理作者的架势——特征p的代数曲线上、“与p互素”的驯顺基本群的构造定理。如果我记忆不差,正是这条定理(我1958年证明的,其时我尚未结识我未来的学生)——连同“蛮横的Lefschetz定理”——构成这份理论的深刻技术成分;而在不可约性定理的证明里,能以全副力量使用它,我曾深感快意。

在Katz第XXI讲的引言(第364--365页)里,描述完该报告的主要定理——关乎射影空间里的完全相交——之后,有这样一段:

\begin{quote}
存在一些归功于A.~Grothendieck的启发式论证,它们倚仗结晶上同调的yoga,以本质相同的方法,使一般的表述对一切光滑射影的X显得可信。
\end{quote}

这条评注暗示:仿佛我是从正文的方法(它出自一位未具名的作者——那多半只能是卷册的两位作者之一)获得启发,在它上面绣出那些“启发式论证”,以推广所证的结果。我记得的恰恰相反:是我的那些“启发式论证”(我远在讨论班之前,循着对Griffiths定理与Lefschetz束的反思\fnote{346}{(**)}{其实正是这些反思——一如我关于抽象代数几何里消没闭链理论的反思(我的另一场“超越理论的纯代数转写”!)——是SGA~7讨论班的起源。},独自发展出来)恰恰“走得通”\dmark{}\origpage{800}——何况不带任何猜想性成分——在X为完全相交的情形。况且,在前一报告(同样出自Katz)里——它专讲Griffiths定理——引言里说:“这里给出的证明(出自Grothendieck)是Griffiths原有证明——或多或少超越性的——在纯代数项下的转写。”这条评注给人的印象是:特征任意的Griffiths定理,证明不止一种,任凭挑选,而我有幸被选中了我的那种。其实,据我所知,此外并无别种证明。况且,凭我当时被迫下的工夫,我怀疑这证明是Griffiths证明的单纯“转写”——一如平展上同调里任何一条大关键定理的证明,都不是某个已有证明的“转写”;或者(干脆说到底)掌握概形的平展上同调,从来不是把寻常熟悉的普通上同调理论“译成纯代数项”的问题。

我清点了N.~Katz报告文本里指向我这个人的全部三条指涉(而在Deligne的八篇报告里,总共只有一条!)。在我看来,三条映出的是同一份蓄意。收尾之前我指出:卷册最后一篇报告——N.~Katz作,专讲特征p上一条L函数的“模p同余公式”——其文本里没有我的名字\fnote{347}{(*)}{这并不全对——它倒是出现了(于是这是指向我这个人第四条指涉):在第410页,与Deligne并提,感谢我们俩向作者解释过他所呈现的同余公式形状的几种等价改写。滑稽的是:他为这些辉煌变体开列的三条编号文献,在报告里一条也不存在——以致这份致谢,俨然一桩善意的恶作剧!(这已不是我在葬礼里遇到的第一桩……)}——连L函数那个寻常的上同调表达式也没提我。其实,正是结晶上同调项下的类似表达式(它还停留于猜想),多年前引我猜出了这条同余公式。我把这个猜想告诉过Deligne;他找到了一个惊人简单的证明,凭他那条对称Künneth公式(陈述于SGA~4~XVII~5.4.21)。我猜想:对这类事情全然在行的Katz,同样清楚这条猜想的来历,只是不觉得值得提。(他在文本里给出的,是与Deligne不同的证明,而且拙劣得多。)

滑稽的是,在SGA~7~II这最后一篇报告的引言末尾,人们读到:Deligne的证明“当收进SGA~5的重编本”(而SGA~5,它还无缘见识自己的头一“版”呢\dmark{}\origpage{801}!)。这可以让人猜想:早于“SGA~4$\frac{1}{2}$—SGA~5”那场行动五年,Deligne或许仍有意(那本属正常)在SGA~5将来的刊行版本里,收入他自1966年以来对平展上同调理论——发展于SGA~4、SGA~5——所作的种种补充\fnote{348}{(*)}{我猜想:正是SGA~7里那些藏匿未遭任何反应(来自任何一位在行的人),才鼓励了Deligne在节节攀爬中迈出下一步:“SGA~4$\frac{1}{2}$—SGA~5”那场行动的大骗局。}。

\textbf{III~Mebkhout的哲学(1981年6月吕米尼会议;Beilinson、Bernstein、Deligne论“反常层”的文章)。}

这里我把前一条注里就此所记的,再录一遍以备记忆。

1. Deligne对我说,他是在一次Bourbaki讨论班上与Mebkhout交谈时得知“好上帝定理”的\fnote{349}{(**)}{见“当值无名之辈与好上帝定理”那条注,n\textsuperscript{o}~48$'$。}——反正在1980年夏之前。这与我从Mebkhout那里听来的相印证:那条定理,是Deligne于1980年10月转告Bernstein与Beilinson的,随即被他们用进他们对Kazhdan-Lusztig猜想的证明里\fnote{350}{(***)}{见5月28日的脚注——附于“不义——或归来的意义”那条注(n\textsuperscript{o}~75);亦见“一种不义与无力的感觉”那条注(n\textsuperscript{o}~44$''$)。}。Deligne补充说:他在与Bernstein、Beilinson合写的文章里没有引Mebkhout,因为拿不准这条定理里Kashiwara该占几分\fnote{351}{(****)}{关于这方面的评注,见前一条注“义务已尽——或真话时刻”,尤其p.~784,及论“Kashiwara”的那条脚注。}。

2. Deligne并不否认:没有Mebkhout前些年的工作,1981年6月的吕米尼会议(他自己是会上最大的明星)不会举行。他只坚持补一句:MacPherson想法的角色,在他看“还要更根本”。至于Mebkhout的名字不见于会议文集,是否有几分怪异、几分反常——他从不曾暗示。

\textbf{IV~上同调的对偶形式主义、导出范畴(“好的文献”、“导出范畴的第0状态”)。}

\dmark{}\origpage{802}1. Deligne向我指明:Verdier那篇文章\fnote{352}{(*)}{即“好的文献们”那条注——这名字显然非它莫属!(n\textsuperscript{o}~82)所引的那篇文章。
(5月12日)关于Deligne这个叫人难以相信的版本,见“荣耀满天飞——或暧昧”那条注(n\textsuperscript{o}~170(ii)),第930、931页。}——它把(别的且不论、并且不提我名)结合于一个闭链的同调类与上同调类的形式主义(我1965/66年在SGA~5里发展的)认作己出——他是在1977年SGA~4$\frac{1}{2}$刊行之后才知道的,至少比那篇文章发表晚一年。这么说来,我曾经的印象——Verdier1976年那手漂亮的操作,是Deligne及其同伙次年那场大得多的行动的一只“试探气球”——倒像被它推翻了。

Deligne对我说:浏览Verdier那篇文章时,对他一目了然:那不过是把我SGA~5里的某些想法铺陈一遍。Verdier终于肯负责提供一份文献,他甚至为此满心欢喜。(“SGA~5的刊行本也许才是更相称的文献”——这个念头,想必从不曾掠过他的心头……)循着这层意思发问,Deligne回答我:他不曾留意Verdier的文章里没有我的名字——又补一句:他承认,他连想都不曾想到要问这个问题。我得到的印象是:他在无声地示意:这类事情在他操心的事里排最末,不值得停留……

2. 在Beilinson、Bernstein、Deligne那篇文章(葬礼I里多次引过;它由Deligne执笔整理,由他在吕米尼会议上宣读\fnote{353}{(**)}{关于这场“值得纪念的会议”与这篇文章,见“不义——或归来的意义”那条注,n\textsuperscript{o}~75.})里,平展上同调的对偶(我1963年发展的)被称作“Verdier对偶”\fnote{354}{(***)}{这场操作分几步完成。经我提议,Verdier在1963年之后,于寻常拓扑空间的语境里,发展了一套“六运算”的对偶理论,循着我在代数凝聚与平展语境里主持发展的那份总设计。这份对偶,不出所料,被我的上同调家弟子们洗礼作“Verdier对偶”或“Poincaré--Verdier对偶”,只字不提我这微不足道的人。另一方面,在1976年的那份“好文献”里,Verdier又在解析的语境里,不加指名地重述了我在五十年代于凝聚框架里发展的那份形式主义的一部分(对它不须改动分毫)。于是这份对偶,在解析框架里,又取名为“Verdier对偶”,或有时“Serre-Verdier对偶”——依旧不提我这个人——连Mebkhout也跟着大势走!可(妙的是一记回马枪)显然之极:代数的凝聚对偶,不过是那份超越的解析理论的“纯代数转写”,一如平展对偶是那份超越的拓扑理论的如此“转写”。于是势所必至:连它们也一并洗礼作“Verdier对偶”(Serre与Poincaré,情有可原,因为离得远)。据Deligne对我说,这正是“人人”争先恐后做了的。幕落……}。

\dmark{}\origpage{803}这个奇怪的称呼,我问过Deligne。他回答我(这一回,带着一丝窘迫):因为“人人”都这么叫。这个“人人”都是谁、以及这算什么理由——我没有细问;毕竟,这门理论该归功于谁,他Deligne一清二楚。

这让我想起一件久已触动我的事。至少在与我交谈、或给我写信时,Deligne从不使用“导出范畴”这个说法而不加上“Verdier的”。这每次都给我一份不快,可我(在发现葬礼之前)从不曾停下来探究其意,更谈不上把点重新点清。倘若当初肯费心对“SGA~4$\frac{1}{2}$”投去稍带好奇的一瞥,对Verdier那部在其中被起出的“论文”的“第0状态”——我无疑会停下来。(后一事的细节,见上II~3.)

\textbf{V~葬礼悼词}

1. IHES的那本周年小册子——我的那篇葬礼悼词\fnote{355}{(*)}{见“葬礼悼词(1)——或恭维”与“葬礼悼词(2)——或力量与光环”两条注,n\textsuperscript{o}s~104、105.}就载于其中——并非出自其创建人暨首任院长Léon~Motchane之手(我原以为是)。册子作者是谁——Deligne把身份告知了我——在此倒无关紧要。他向我确认:涉及我的那段,确确实实出自他的手笔;而那一段,连同涉及他Deligne本人的那一段(出自册子作者之手),在付印之前,都确确实实得到过他的“绿灯”。他写我的文字,起初更长,被册子作者删节了(经他同意)。涉及他自己的文字,Deligne也审阅并修改过。所以,这些文字确确实实代表着Deligne的立场——关于他的著作,与我的著作。

2. 我问过Deligne:我是否猜错了——我曾认定,在他任何一篇发表的东西里,他从不曾暗示:他或曾从我口中\dmark{}\origpage{804}学到过什么。他向我确认无误——只有一处保留。那涉及他写给(布鲁塞尔)国家科学研究基金会的一篇小传,为着“五年奖”的颁发。那项奖(我记得是1974年)授给他,奖励他对Weil猜想的证明。诚然(他补充说),这篇小传不算数学出版物,流布也极其有限。至于我,我根本不知道它存在。经我请求,随后几天里他给我送来一份复印件;这篇小传,我打算在下一注里再谈。

Deligne向我确认的这份对我的系统性否认,似乎不构成他的任何问题。他不觉得其中有什么怪异、值得停留的东西。鉴于他这份心态,我没有起意在这个方向问他任何问题——我不认为多问能榨出更多。为这份回顾收尾,我只添一句:凡属严格意义的“物质性事实”,我对Deligne的诚意并无怀疑——那诚意在我看是显见的\fnote{356}{(**)}{(5月12日)凭回顾,对这份印象毕竟浮现了某些保留,如前文一条脚注((*)第802页)所提到的那些。还显出来:有两条粗疏的事实性错误在我的注里,Deligne不曾向我指出——它们本很难逃过他的眼睛。(我曾漏看:他早在1970年Hodge~I里就披露了“权的yoga”的一部分,而且他1979年起就谈到Motif.)}。唯一的例外,是他的那个断言:SGA~5(1965/66年的)讨论班在逻辑上依赖SGA~4$\frac{1}{2}$\fnote{357}{(*)}{诚然,这一断言之来,并非Deligne自发地主动前来,给我送来“物质性说明”,以启发我、以显豁他那全然的诚意;而是在我刚刚当面道出我对“SGA~4$\frac{1}{2}$—SGA~5”那场令人难以相信的行动的感受之际,在“保全面子”这份未及预料的需要的压力之下。关于这一点,见前一条注“义务已尽——或真话时刻”的末一部分(2月18日)。}——那些结果自1973年起,在Deligne讲述其Weil猜想证明的各报告的边缘上发展而成。诚然,靠“俘获”母讨论班SGA~5的某些报告(尤其那篇关于结合于一个闭链的上同调类的),凭Illusie(他负责SGA~5的编辑)与许多其他人的默许,他赢得了这份辉煌的战果:SGA~5通篇塞满对SGA~4$\frac{1}{2}$的指涉,足以给(一位不甚留神、或不甚在行的)读者一种印象:SGA~5确确实实依赖SGA~4$\frac{1}{2}$,而后者在一切方面都摆出一份“更早”文本的架势。这手腾挪,在我们这门科学的编年史里想必独一无二;依我看,它足以把本世纪的七十年代,从数学所经历的一切其他时代里标出来。

\subsubsection*{注~164$_{1}$}
\notehead{164$_{1}$}{}

\dmark{}\origpage{805}关于那份出自Weil猜想的“权的哲学”,这份“亲缘”在我看来可总结如下。

a) 一如“我的孤儿”那条注的子注~n\textsuperscript{o}~46$^9$所说:作为Weil猜想背后那份“哲学”的一部分,Serre曾传给我某种“虚拟权的yoga”——在域上有限型概形的$\ell$-adic上同调的层面上。他不曾试图给出它一个精确的显式表述;对不同~$\ell$~所发生之事之间的关系,则全然神秘。

b) 自六十年代初起,引导我发展一份“Motif的yoga”的两条主要动机之一,恰恰是:把对不同~$\ell$~的那些“虚拟权结构”彼此连起来。(关于这一点,见“一个梦的回忆——或Motif的诞生”那条注(n\textsuperscript{o}~46),尤其第208页。)由此便清楚:这份结构该当在一个Motif一切可能的“实现”上重现——不只在$\ell$-adic实现上——尤其(在基域$\mathbb{C}$上)在De Rham-Hodge实现上。

c) 这份虚拟权的哲学——其终极的源头是Motif——经我告知Deligne之后,他为这份yoga添上一份重要的精确化:揣想一个Motif上的虚拟权结构,联系于一份(必然典范的)\textbf{按权递增}的\textbf{过滤}。于是这份过滤该当在Motif的一切实现上重现——无论$\ell$-adic实现,还是(在域$\mathbb{C}$上)De Rham-Hodge实现。Deligne的这份“揣想”,成了他“混合”Hodge结构理论(我称之为“Hodge-Deligne结构”)的起点,也是他为它们\textbf{成形}的定义的两条本质技术成分之一(另一条是De Rham过滤——我早在五十年代引入)。他为描述$\mathbb{C}$上任意有限型分离概形的“Hodge上同调”所作的尝试,其成功,可以被视作我们眼下就Motif上按权过滤存在性这份“揣想”之有效所拥有的主要的(甚或唯一的)“凭据”。

当然,阐明“$\mathbb{C}$上有限型概形上的Hodge系数”概念,本是我围绕Motif那份大工作纲领的一部分——Deligne对它是第一手的、日复一日地知情的:其方式,须使X上的一个Motif对应一份“Hodge实现”,并且对X上光滑而纯粹的Motif(例如那些取一张X上真而光滑的概形的“第i维X上Motif上同调”而得的Motif\dmark{}\origpage{806}),人们重新找回“Hodge结构族”的概念(六十年代尤其由Griffiths研究过)。而且,X变动时,这些“Hodge系数”的范畴须满足一份六运算的形式主义——在Motif的层面上映照同一份形式主义——Deligne的贡献,代表朝这份纲领之完成迈出的头一步:即(本质上)对X缩为一点的情形给出Hdg(X)范畴的描述\fnote{358}{(*)}{为求完善,须给Deligne的定义补上引入一个合适的\textbf{三角}范畴~Hdg$^*$(它莫非也是Hdg的导出范畴?)。他漏掉了这一步,在我看来是(此后还有别的)最早的信号之一:对导出范畴的yoga与六种运算的冷淡——它一直肆虐到1981年“反常层会议的转折”。},以及“实现”函子的描述——即本质上,在$\mathbb{C}$的有限型分离概形上构造一套上同调理论,取值于Hodge-Deligne结构这个范畴。

\section{4.~死之舞}

\subsection{(1) 为一副模糊的骸骨作安魂曲}

\notehead{165}{}

(2月22日)自他去年10月的来访算起——甚至更早,自他8月底的那几封信算起\fnote{359}{(**)}{见“义务已尽——或真话时刻”那条注(n\textsuperscript{o}~163)——我在那里给这次来访“定位”;亦见8月底的那两封信(它们是在近两个月的沉默之后到的——那沉默,继我寄出引言与葬礼目录之后)。}——我的朋友皮埃尔,成了昔日学生与好小子中的头一份:满载着一份显然动人的好意,要驱散那桩不巧溜进我们之间的误会,要让我感到他的善意与诚意。双方讲定:在《收获与播种》由我的大学(USTL)作预定预印之前,他对他读我注子的内容、乃至其存在,守口如瓶。他是否把话守全了,我不知道;反正好几路回音传到我这里\fnote{360}{(***)}{比如:我收到Illusie一份未署日期的预印本(我猜是最新的),出自某个未具名讨论班的一篇报告(文中特意说明:它不对应任何一次口头报告)。标题里——难以置信,却是真的——我的名字赫然在列:“据A.~Grothendieck的Barsotti-Tate群形变”,作者Luc~Illusie!引言里还有大段大段的“Grothendieck”——我疑心自己在做梦。分明是有什么事发生了……随信附有一封短信,他就若干“Grothendieck风格”的同伦代数之点向我讨教,并自问:“搞K理论的人们”(即Quillen等)为什么用层工作,而不用我二十多年前引入的装备库里的那些(伪凝聚或完美)复形。人们也的确要问个为什么……回信里,我不得不示意:这样的问题,轮不到他、也轮不到我任何一位昔日学生来问我。此后我再没有他的音讯。},我颇有印象:他准是两头都递过话——暗示:或许该是时候,给那位老师傅递几分体贴了(那位老师傅:人们偶尔在小圈子里谈起,却小心翼翼地避免在公开场合提名……)。

我另外有一种印象:骨子里,我的朋友并不相信(或至少不愿相信),我会\dmark{}\origpage{807}确确实实把葬礼付印,与《收获与播种》的第一部同时。第一部,倒是一副十足的“心软爸爸”的形象:唯恐点谁的名会让谁难受,对自己想到的种种自家的缺失,又随时准备当众认账。读那部“愚妄与更新”——我朋友度假动身之前、我给他寄出葬礼引言之前,我得过它的一点回音——不曾让他有分毫不安;恰恰相反——它多半还助长了一份自我满意的做派,那份做派,在他身上我早已熟稔——那丝居高临下、或至少是对那位显然已故的师父的护短的做派。葬礼,可就全然不是同一回事了:牌,骤然间干脆摊上了桌面!我猜想,读那篇引言,该给了他一记冲击——可惜当时我不在场,否则或许会有些什么发生。反正,他给自己留了回神的时间,才来见我——一阵风似的,在他迁居美国前五分钟。而他来得满载好意,相遇进行得何其合家欢、何其“心软”,以致——不妨说“以归谬法”——这反倒排除了:那位心软爸爸当真会把某一份不像他的文本当回事(那份文本,不说也罢,最好忘掉……),甚或把它散发给一群与他朋友皮埃尔同样体面、样样“正派”的人们,以及那位如他所一直认识的前亡故者……\fnote{361}{(*)}{不过,把我在葬礼上的全部各注公之于众——与《收获与播种》第一部一样——我的这份意图,从不曾有过分毫的迟疑;对此,我当然不曾留下任何暧昧。}

一如他所允诺的,回到比雷斯之后没几天,我的朋友就把我上面说的那篇小传\dmark{}\origpage{808}寄给了我——他1974年(或1975年)为(比利时)国家科学研究基金会而写的那篇\fnote{362}{(*)}{这篇小传头一次被提到,是在“神经中的神经——或侏儒与巨人”那条注(n\textsuperscript{o}~148)的最后一条脚注里。亦见前一条注n\textsuperscript{o}~164的结尾(V~2部分)。}。文字相当短,小小的两页;我那时饶有兴味地读过,刚才又重读了一遍(这已是第三遍,我想)。可乍一看,我并没有这份印象:这文本带来了什么新东西,值得我在葬礼里为它停留。诚然,藏匿的技巧——在我朋友身上我本已识得够多——在这里,在一篇百余行的紧凑文字里,得到了格外触目的演示。我的名字在其中出现四次(一如Serre,而Weil出现三次)——却没有任何东西让人想到:他与我打过交道的路数,或许不止是我那个讨论班(主题未说明)的一名无名听者(1965--66)。在提及我的四段文字里,有三段把我与另一位数学家并提(两次Serre,一次兰金),为的是避免给人印象:我曾在他身边扮演过稍稍特别的角色。这又是一门在别处久经考验的技巧\fnote{363}{(**)}{我这里想到那条一行字的简略指涉——在(尼斯大会上)Deligne宣布其Hodge理论结果的公告里——它把Serre(不提名)与“Grothendieck的猜想性Motif理论”并在一口气里。详情与评注,见“牺牲者”那条注(n\textsuperscript{o}~78$'$)的子注~n\textsuperscript{o}~$78'^{1}$。}。既然引起来不长,我就不揣冒昧,把我的微不足道的人现身的三个段落全文引出,以照亮那些不象我这样手头有这篇小传的读者。

第三个段落,接着(方才已作过的)对1965--66这一年的回顾——“在高等师范学校理想的氛围里,作为一名外籍寄宿生度过”——说下去\fnote{364}{(***)}{出于一个我无从把握的缘故,Henri~Cartan在这里没有具名。或许因为:凭我对他怀有的一份蓄意(见“格格不入者”那条注,n\textsuperscript{o}~67$'$)给他撑腰,Deligne执意避免任何“他或曾是谁的学生”的表象。“高师生”的身份,会即刻唤起“Cartan的学生”的联想;点名Cartan,只会坐实这份联想。}:

\begin{quote}
“在巴黎,我听了Grothendieck的讨论班与J.P.~Serre的课。每周三个小时的课\dmark{}\origpage{809};然而,尽管工作愉快而拼命,一周的其余时间也才勉强够我消化它们(165$_{1}$)。从Grothendieck那里,我学到了代数几何的现代技术;从Serre那里,学到了数论那令人着迷的美(165$_{2}$)。Serre的课专讲椭圆曲线理论,其中交织着……”,\end{quote}

随后继续铺陈Serre课程的魅力与花样。不在行的读者会想:作者所说那份“愉快而拼命的工作”,对象是这些课——每周三个钟头(言下之意:消化一个Grothendieck讨论班的“最大的自然一般性”,根本用不着下工夫……165$_{1}$)。

第五段,谈他对Weil猜想的证明,人们读到:

\begin{quote}
“我最可观的成功,是证明了‘Weil猜想’(……)。我无疑是由于既熟悉Grothendieck的著作,又——在一个全然不同的领域里——熟悉兰金关于模形式的著作,才做到这一点的。”
\end{quote}

人们将叹赏那个踌躇的“无疑”(放在那里,真是大师手笔!),叹赏那句“在一个全然不同的领域里”(暗示:我的著作与模形式毫无干系\fnote{365}{(*)}{诚然,“模形式”代表着我数学教养里一处令人惋惜的空洞(空洞还多着呢),一如解析数论——我对它至今不曾“挂上钩”。可我的见识毕竟足以知道:没有来自代数几何的想法——它给这门理论提供其“几何”的内容——对模形式的理解几乎无从谈起,而模形式理论最深刻的问题,又与\textbf{Motif}的(长期隐而不显的)在场密不可分。如下文将见,Motif同样不声不响地,出现在这篇小传的下一段里(别名:葬礼悼词(3)!)。}),尤其叹赏那个“既熟悉……又”——我就有这份荣幸,由它引荐登场——好把我那份宏大的奠基工作\fnote{366}{(**)}{关于概形的概念、与平展上同调形式主义的发展——Deligne小心避免提它们,除非在前一条引文里,借那个客气的、无我之称的婉辞“代数几何的现代技术”。},与一条从兰金那里借来的“一时一地”的技术想法,放在同一个天平上。

最后,在接下来那段谈DeligneHodge理论工作的文字里,有这样的句子:

\begin{quote}
“受到算术的启发\dmark{}\origpage{810}——尤其受到Grothendieck对Weil猜想之深刻意义的理解的启发——我把他的理论(以非平凡的方式)推广到任意簇的情形,并且(与Sullivan合作)推广到“形”的其他不变量,而不止上同调。这份理论的根已经颇老,有皮卡关于“两个独立变量的代数函数”的论著(约1890年);但如今人们对它的认识,大概不过剩一副模糊的骸骨。”\end{quote}

须要我费心把这段抄下来,我才恍然:“Grothendieck对Weil猜想之深刻意义的理解”云云,是我那位辉煌的前学生不点Motif之名的大师级“暂停!”手法——而你又抓不到他把它们讳莫如深的把柄!毫无疑问,“他的[亦即,\textbf{我的}]理论”——我此刻才反问起它来(整个这一段,前几遍阅读都逃过了我的注意)——只能是那份出名的Motif理论:四年以来,提它的名字都谈不上(而往后八年,更无人会提!)。那措辞模糊到这般地步——说穿了,除了一小撮在行的人(他们多半不曾象如今的我,有机缘读到这篇预备的葬礼悼词),谁也读不懂——以致连点破一句都嫌多余:他“推广”的那门“理论”,其实整个是猜想的!照上下文看,所云的“推广”,也只能是指Hodge-Deligne理论。这是我的朋友给自己的一点象征性的快意:他在这里断言(断言的地点,连同措辞那闪躲的模糊,保证他永远不会被驳)——Hodge-Deligne理论(它至今仍在童年)“推广”了我让他看的那幅宏大的Motif图景。可在这幅图景里,一门抵达全熟的“Hodge理论”,不过是全图诸多“层”中的一层\fnote{367}{(*)}{(2月27日)关于这方面的详情,尤见“墓上的旋律——或自负”那条注(n\textsuperscript{o}~167)。}。至于“形”的其他不变量,我早在六十年代就“了然于胸”(作为我的“Motif的yoga”的一部分):“任意”的代数簇(Deligne特意强调“任意”)有一个“Motif同伦型”,其高阶$\pi_i$($i\ge 2$)推广Motif的“几何”基本群,并且(对一个给定在数域$K$上的纤维函子)显形为$K$上的仿射pro-代数群。

\dmark{}\origpage{811}至于把皮卡抬出来作“这份理论的根”——在我看来,这整段是装点门面的赝品,插进来为着一石二鸟:既“显得有学问”,又顺势引出紧随其后的收尾一段\fnote{368}{(*)}{这收尾的一段,将是紧随本注的那条注(n\textsuperscript{o}~165)的对象。}。“模糊的骸骨”这个字眼,在我看来,同样是我朋友付给自己的另一份“象征性快意”的表达:他暗地里从那幅宏大的视景获得灵感,却把它始终埋着\fnote{369}{(**)}{Motif的视景,以两种方式保持“被埋”。一方面,对着\textbf{外界},对着数学公众:对Motif概念绝口不提(除了1970年Hodge~I里那半行的“暂停!”——参看注~$78'^{1}$)——直到1982年,这概念才在Deligne的(心照不宣的)父权之下“大张旗鼓地”被起出(见注~n\textsuperscript{o}~51及以下诸注)。另一方面,甚至在他自己的私人使用里,我看到:这份视景也被Deligne剥去了它真正的\textbf{元气}——那份使它\textbf{不止}是一种通行配方的汇编(为在代数簇的上同调里认路)的东西:一份足够宏大、足够深刻的\textbf{力之梦},足以充当灵感、地平线上的准线,予或许数代的算术几何学家。
Deligne(照旧不声不响)以“模糊的骸骨”这个字眼指称这份视景,令人惊心地显出他身上的\textbf{掘墓人心态}——在他与这个梦、与梦所由出的那位匠人的关系里,他始终守着这份心态。这已不是那种尚能感到一口气(如他当年曾感到)的心态,也不是体现一个梦的心态。体现一个梦,不是把它\textbf{拿来用}于自己的目的(同时还否认它……),而唯有\textbf{甘为其仆}。},却在心里把它当作——照总账——不过一副“模糊的骸骨”(照样是那副“暂停!”的做派,脸上却不露)。到头来,这些随手挥洒的藏匿,竟比我预想的更有滋味——我原本只打算顺带点一笔,算是了却良心。如今最触动我的,已不是(一如头几遍匆匆浮泛的阅读)那“暂停!”风格的完美——那风格早领教得倒胃。而是:这篇文字——它写于那篇葬礼悼词九年之前\fnote{370}{(***)}{见“葬礼悼词(1)——或恭维”与“葬礼悼词(2)——或力量与光环”两条注,n\textsuperscript{o}s~104、105.}——以触目的方式预示了后者;而且(依我看)在两个方面。一方面,凭那份照例的模糊:凡我这微不足道的人现身处,必得围着一团雾(与此相对,此处,Serre课程的回顾却伴着技术细节的铺陈)。另一方面,同向而行:对平展上同调或$\ell$-adic上同调——那份我从虚无中发展的、新的而根本的工具,没有它,Weil猜想再过一百年恐怕也证不出来——竟全然\dmark{}\origpage{812}无一词提及。其实,一如那篇葬礼悼词:“上同调”这个词,从未在与我的名字的关系里出口;同样无人暗示:Deligne对Weil猜想的证明,不过是一段长途的\textbf{最后一步}——那段路最长、也最新颖的部分,是另一个人在他之前走完的,那时我那位辉煌的学生尚未在数学的舞台上现身\fnote{371}{(*)}{这位“另一个人”的贡献,被Deligne用无我的字眼藏了起来,如“代数几何的现代技术”[或在他处,“有力的工具”]。}。

\subsubsection*{注~165$_{1}$}
\notehead{165$_{1}$}{}

一如我在下面几行里指出的:那份措辞不可抗拒地暗示,“每周三个小时的课”指的是刚谈到过的“J.P.~Serre的课”——而且两句之后还会再谈到它。其实,Serre每年只开一门课(在法兰西公学院),每周一个钟头。倘若试图消除歧义,把文本读成历年连年的Serre“课程”(与上下文所暗示的相反),又会撞上另一处不合:Serre每年换主题,绝不局限于椭圆曲线——可两句之后偏偏这么说。

我的朋友在这里挪用Serre这个人,想就那在他数学生涯的关键岁月里属于我的角色,混淆视听;有意思的是:据我所知,文献里唯一一处言明Deligne曾是我学生的——出自Serre的笔下,他借此(而不点破地)弥补了那位辉煌前学生亲手犯下的种种昭彰的省略。那是Serre1977年5月就Pierre~Deligne的工作所作的报告,为着颁发1978年菲尔兹奖章的国际委员会。该报告在1978年赫尔辛基大会颁发菲尔兹奖章之后公开。报告开头如此:

\begin{quote}
“Deligne最初的工作,直接受Grothendieck的启发——他是Grothendieck的学生——关乎代数几何的若干技术之点。我只限于提及:……”
\end{quote}

\dmark{}\origpage{813}再往下,Serre也提到我的想法与结果在Weil猜想之证明里的影响,以及(经由Motif)在Deligne模形式工作里的影响;但在Deligne--Mumford关于$(g,\nu)$型代数曲线之模多重簇的工作里,以及在Hodge-Deligne上同调的观念里——它与Motif的yoga及Weil猜想的关联,似乎逃过了他的眼睛——则没有提。(诚然,Deligne为此尽了遮掩之能事。)

菲尔兹奖章颁发之际关于Deligne的演说,照公认的惯例,原该是另一个场合:公开重提这份与我的关联——它此前一直被当事人捂着。出于一个我无从把握的缘故,受命介绍Deligne工作的数学家,不是J.P.~Serre,而是N.~Katz——SGA~7~II与Deligne“合著”的那位(关于这一点,见注~n\textsuperscript{o}~164(II~5))。不用说,N.~Katz对这层关联不着一字——尽管它于他一清二楚、而且是第一手的。(倒是,他顺手不露声色地弥补了那位桂冠得主关于我的若干处有些难堪的省略……)

\subsubsection*{注~165$_{2}$}
\notehead{165$_{2}$}{}

这里考语的拣选(给我的是“现代技术”,给Serre的是“令人着迷的美”)决非出于偶然。我在其中分明察觉我朋友的用意:把那份着迷——自我们相遇起(或许更早)便把他系于我这个人、系于我的著作的那份\textbf{着迷},他眼看它在他眼前日复一日地成形、铺展——(象征性地)打发掉。

在别的场合我也留意到我朋友的一份蓄意:把我的著作(尤其EGA即《代数几何原理》、SGA即“布瓦玛丽”代数几何讨论班)看成并呈现为某种“汇编”——搜罗一些或多或少技术性的结果,“人人”素来皆知,而我可嘉的努力在于把它们白纸黑字写下来,好终于补齐缺失的文献,从此不再谈起。可骨子里,他清楚该把账算在哪里:EGA与SGA的每一卷,呈现的都是我引入的想法——多少年间,我是它们唯一的持有者与辩护人;呈现的是无人梦到过的技术(除我之外),而我不得不以不倦的耐心把它们发展、检验、完善,直到它们磨合纯熟,随时可以进入“人所共知”的领地。这一点,没有人比他知道得更清楚;可与此同时,他彰示了十多年的那份蓄意\dmark{}\origpage{814},终于变成了一份“第二天性”,而他自己,成了它头一号(若非唯一的)受骗者。

就在几周前,我还为之一怔:我的朋友——自他10月造访寒舍以来,对我处处体贴——让人给我送来一份他与Heinze博士(在斯普林格负责“Ergebnisse der Mathematik”丛书)之间通信往来的抄件,谈EGA的重版计划(其中许多卷已售罄或濒临绝版)。在其回信中,Deligne毫无保留地推荐全部原样再版,“ne varietur”——几乎不作改动,理由是除了一处例外(EGA第三部第二部分,在那里用导出范畴叙述会更好(sic!)),这部论著“has aged very well”,它的巨大功绩正在于提供不可或缺的参考文献:“Thanks to it [EGA], in algebraic geometry (as opposed to analytic geometry, for instance) one can march securely on the ground without having to worry if this or that is indeed in the littérature”。(他接着提出若干建议,关于可能给某些卷册添附的附录,以及有能力提供这些附录的数学家……)

斯普林格出版社对我这个人的态度极具代表性:这场通信(关于我身为作者的著作的再版)是\textbf{同Deligne}进行的,而斯普林格甚至不曾先通报我这一计划。过了一个多月(在1月24日的信中),Heinze博士才顺带向我提起这事,算是了却良心——说Deligne教授先生“曾非常体贴地向我提供他1984年12月19日来信的抄件”(真是体贴……),并且“我们[斯普林格]愿听取您对此事[再版计划]的意见”(真是太荣幸了……)。我回复说:鉴于斯普林格出版社在重印方面的惯常做法(想想它在Lecture Notes里出版SGA~7与SGA~5,连知会我一声都不曾,更遑论征得我的同意),我认为向斯普林格出版社通报我那显然无足轻重的“意见”,纯属多余。事情便停在这里……

\subsection{(2) 信仰表白——或伪中之真}

\notehead{166}{}

(2月23日)归根结底,昨天在谈论我朋友皮埃尔的小传时,我还不曾抵达我的真正题旨。同“模糊的骸骨”(别名:Motif理论)的相遇是一段意外的插曲,当时我正准备顺势引出小传的最后一段——它紧接在方才所引的那最后一段之后\dmark{}\origpage{815}。现在终于迎来了“传记小记”的结语,也就是我打一开始就想谈到的地方:

\begin{quote}
“在结束时,我想强调与往昔数学家(从1800年至今)的著作保持接触对我何其宝贵,无论直接接触还是通过比我更博学的人(如A.~Weil与J.-P.~Serre)转介。我们‘是骑在巨人肩膀上的矮人’,最美丽的现代数学理论,都是由解决他们遗赠给我们的若干问题的希望所驱动的。

Pierre~Deligne”
\end{quote}

一如常见的情形,我对这几行——某种向往昔宣誓的信仰表白——的第一反应停留在表面,停留在字面含义——但我必定隐隐感到,在这之外另有玄机。这句引文(无疑出自某位著名数学家,而我理应“同所有人一样”知晓它)并不合我的意。我在其中感到一种刻意为之的谦逊、甚至谦卑,它十足是一副姿态,根本不符合事情的单纯实在。而且这一蓄意几乎近于荒谬:如果每一代人在格局上都比前一代“更小”,人类早就耗尽元气而灭绝了,沦为一群可悲的侏儒!我深知人身上的创造力在今天并不比一百年前、或一百个世纪前更少(无疑也不更多)。我也深知,单就数学而言,我所熟识的人们的这样那样的著作(包括我自己在内),完全配得上列在往昔最伟大的数学家名下。我也深知,\textbf{我}在做数学时的动机\fnote{372}{(*)}{顺带说一句,包括Pierre~Deligne本人!},以及我昔日在数学界的大多数朋友的动机,绝不在于“解决前人遗赠的若干问题的希望”!若非如此,我们的科学将无力自我更新——它将不再具有创造性。

在这篇借来的信仰表白中,更让我震惊、或者更确切地说是更让我\textbf{痛心}的,是我深知做出这一表白的人,比我所认识的任何其他人都更蒙受了令我赞叹的“禀赋”,而且我曾在他对待数学的进路中确认了一种“新鲜之感”,凭此他理应做出大事业,如极少有数学家曾有特权做出的那样\dmark{}\origpage{816}。我心中有一份痛惜,也有某种怨恨,因为在那个自诩于同往昔之人的交往中获得谦卑的人的姿态背后,我感到了一场\textbf{退却}。退却了那份在他体内的创造力——他似乎很久以前就已将它遗忘,而正是那份力量使他成为远\textbf{非其他}、远胜于骑在巨人肩膀上的矮人这一意象所暗示的存在\fnote{373}{(*)}{(2月25日)这种“退却”的印象,强烈地同葬礼悼词所引发的印象联系在一起。参见“葬礼悼词(2)——或力量与光环”那条注(n\textsuperscript{o}~105)末尾对此的唤起。}。

这是自初读小传以来,我第一次试图界定那次阅读最初在我心中唤起的情感。在随后的日子里,甚至并非出于我的刻意图谋,这层思考一直在发酵。尤其是最后这段话一直萦绕在我脑海中,像一件确乎不同寻常、难以释怀的事情。在这篇短小的传记文字结尾那表面荒谬的表白背后,我必定隐隐预感到一种\textbf{意义},它最初无疑停留在无意识层面,然后逐渐浮向表层,其间就我记忆所及并无严格意义上的反思。毕竟我很清楚,我的朋友皮埃尔并不比我更习惯于钻研往昔的著作。就算他读得更多,那也不是陈年古卷,而是那些在消息灵通的圈子里流传的最新抽印本与预印本,他总是最先拿到它们。我也同样清楚,我的朋友从我离开数学舞台以来(乃至在此之前)滋养其工作的灵感,绝非源自皮卡或上个世纪甚至本世纪其他受人尊崇的先驱。如果说他确实乐于“骑在某人的肩膀上”,那绝非在一份公开的、纯粹修辞性的信仰表白之中,而是秘密地、\textbf{切实地}骑在上面——毕竟在反思某场葬礼之际,我处于再合适不过的位置,去知晓\textbf{谁}是那个承受代价、吞下苦果的人!人们口头上用“往昔的伟人”替换了那位从不被提及的人\fnote{374}{(**)}{或者,如果实在无法回避,就假意以迂回的方式、“暂停!”的风格来提及他。}——而在上一段里,人们刚刚心照不宣地把Motif(别名“如今不过一副模糊的骸骨”)的父权归于那些伟人\dmark{}\origpage{817}——从而使这副替代面孔背后的\textbf{真正}身份变得更加昭彰。

我曾多次观察到,人身上有一种力量,似乎具有普遍的本性,驱使人在排除万难之中,往往以迂回和象征的方式,去表达那些无法公开显露的愿望与企图(无论清醒的还是潜意识的),从而给予它们一种出路和满足——这种满足(在“理性”的措辞与通行的共识看来)或许显得微不足道,却并不因此而不具实质意义。在某种意义上,这是一股推动我们(仿佛身不由己地)向愿意倾听的人宣告我们存在的真相的力量(而我们每个人心中都有一个听觉敏锐的孩子……),而且\textbf{哪怕}这样“宣告”出来的东西无论在他人面前还是在我们自己面前都是最严重的禁忌。表达这一力量的首选领地是梦,这也是为什么梦是通往我们自身认识的一把无比强有力的钥匙。然而,正是由于梦的亲密、个人性质——它只向我们自己显露我们自己——这一表达手段对我们来说绝不充分,因为它无法在\textbf{他人面前}、甚至在象征意义上在全世界面前证实我们存在的真相。正是因此,在每一个看似挑战理性的荒谬无稽背后,都隐藏着一种“意义”——或者更确切地说,荒谬无稽正是潜意识以无误的本能所选中的\textbf{特选表达手段},用以\textbf{宣告这一意义}:它既被隐藏,又被堂而皇之地昭示在所有人面前\fnote{375}{(*)}{关于借表面的荒谬无稽而大声宣扬的\textbf{意义}的另一个格外张扬的例子,见“玩笑——或‘权复形’”那条注(n\textsuperscript{o}~83)。亦见“表面与深处”那条注(n\textsuperscript{o}~101)末尾(第440页)以及紧随其后的“书写的颂歌”(n\textsuperscript{o}~102)中的评注。}!

这正是我在读到这个“荒谬之谈”之后的日子里朦胧感到的东西:一个“矮人”(天生本该是巨人)骑在一个“巨人”(其手段比那个自诩“矮人”的人远为质朴,而后者骑在他身上却又否认他……)的肩膀上。阻碍我清晰意识到这个荒谬所揭示的意义的原因之一\dmark{}\origpage{818}\fnote{376}{(**)}{另一个原因,而且在我看来是主要障碍的,是面对“相信自己眼睛的见证,哪怕所见之物闻所未闻且被所有人否认”时的一定程度的\textbf{怯懦}。最近在“义务已尽——或真话时刻”(n\textsuperscript{o}~163)那条注里我又遇到了它。尤见第782页脚注(**),我在那里探究了这种“不相信”。},无疑是我不愿在那张生硬的“巨人”画像中认出自己;或者更确切地说,是不愿在那张确实曾属于我的品牌形象中认出自己——而通过这个荒唐的荒谬无稽的意外转介,它突然在拷问着我!直到数周之后,在12月18日那条注“神经中的神经——或侏儒与巨人”(n\textsuperscript{o}~148)里,我才终于回到矮人与巨人的奇特意象上,这一回是在掌握实据的条件下,处在关于葬礼的反思语境完全准备好接纳它的时刻。这个意象(就在当天)随即显现为一个关键的“力量意象”,对于理解我朋友对我的关系具有决定性意义,更深层、更关键的是,它是理解我朋友同他自身关系(注定永远只能是局部的理解)的开端——也就是理解\textbf{他自身人格中分裂}所采取的特定形态。既然这场葬礼首先由我的前学生兼前继承人发起并实施\fnote{377}{(*)}{诚然,在这一“实施”中,他是与“教团”密切勾结而行动的,在某种程度上充当了实现集体意志的工具。(见“掘墓人——或全体教众”,n\textsuperscript{o}~97.)但这一力量意象完全可能同样存在于该教团的“集体无意识”层面,并在其众多成员的个体无意识中找到共鸣,特别是在曾是我学生的一些人身上(而不仅是Deligne一人)。(5月12日)自写下这些文字以来,这一直觉大有进展,如今已显而易见。关于这一点,见“信使(2)”那条注(n\textsuperscript{o}~181)。},那么正是这同一股力量,在我看来成了贯穿这场漫长调查中始终顽固运作的枢纽力量,成了它的真正\textbf{神经}。在它出现在笔记中之后的两周里,它一直处于反思的核心,贯穿12月18日(已引的“神经中的神经——或侏儒与巨人”)到1月3日“敌对的兄弟——或传递(2)”(n\textsuperscript{o}~156)之间接续的九条注。

这个看似平淡无奇的意象在我的反思中所扮演的枢纽力量意象角色的“有效性”,也就是这股力量在我朋友自身心灵中的\textbf{实在性}问题——作为深层冲突的表达与不可遏制之补偿行为的驱动力\dmark{}\origpage{819}\fnote{378}{(*)}{用“不可遏制”一词,我绝非暗示这股力量的存在具有不可抗拒的宿命性,从而免除我朋友的责任。这样一股力量在我们体内的运作,完全取决于我们乐于并执意回避对它的认识,以便坐享通过这种刻意的“无知”所“买来”的微小特权。代价是高昂的,诚然,但不去认识这一点正是“交易”本身的一部分。}——这一问题在我看来无法通过“证明”来解决,即无法通过一种旨在赢得任何善意且充分知情之对话者认同的“客观”进路来解决。对我而言,这种实在性无可怀疑,我的深切信念并非某种“论证性”步骤的结果。诚然,在刚才提到的两周反思中,它得到了深化(我在此不打算对其作“总结”或“盘点”)。但它从第一天起就已存在——从我在阅读之后第一次劳神将阅读带给我的灵感落于纸上那一刻起,仿佛听凭一个无声之音的口授\fnote{379}{(**)}{我感到这种由“无声之音”口授的意象在《收获与播种》中不止一次出现,每次都自然而然。这绝非某种文学套话的重复,而是(在我看来)发现过程的一个普遍方面,因情境不同而或显或隐。}。我随后“想起”了深处本已“知晓”的事情。我必定通过某种感知能力“知道”它——这种能力固然非同寻常,但比我们在对事物的\textbf{清醒}认识层面通常运用的能力精微得多。那些压制在我们内心“某处”被感知到、却与我们例行接受(或者说,\textbf{视而不见})周遭现实的逻辑不合之物的机制——无需多言,它们在我身上与在任何人身上一样运作。在这方面如果说我与他人有什么不同,那就是我终于意识到了它们在我体内的无声运作,尤其是自打我开始“禅观”以来:在不可遏制的好奇心驱动下,我有时会下功夫去\textbf{驻足}在那些我感受到的事情上,其效果便是将深层模糊隐匿之物\textbf{召唤上升}到意识表层,并使之成形。

而且初始的感知在给它赋予形状的\textbf{工作}过程中会发生转变,在光天化日之下成形\dmark{}\origpage{820}。这项工作同时是一场\textbf{沉淀},通过它,感知的转译(转译为明白晓畅的言语)逐渐摆脱最初玷污它的主观先验。在本例中,那些扭曲的先验之一(在方才所引的最后一条注中被察觉)便是我体内根深蒂固的倾向,促使我“以阳看待自己”,即便在显然是我存在的阴面——“我身上的女性”——提供理解钥匙的情形中亦然(或者至少提供了不可或缺的钥匙之一,或提供了一份细致理解所需的“视角”之一)。我在别处谈到过那些固然完全“主观”却无可置疑的迹象,它们向我昭示这样一项工作的\textbf{进展},并在我误入歧途或陷入暂时停滞时发出警告,一旦察觉,停滞便随即消解\fnote{380}{(*)}{关于这一点,见“孩子与海——或信与疑”那条注,n\textsuperscript{o}~103.}。

\subsection{(3) 墓上的旋律——或自负}

\notehead{167}{}

(2月25日)昨天的大部分时间都在给一位年轻同事Norman~Walter写长信,他似乎有志于投身Motif理论,不为决然谈不上鼓励人的大环境所吓倒。这一回是密密麻麻的八页纸(打字机打出),谈Motif范畴与最重要“系数范畴”的“六种运算”。这让我再次惊愕地意识到:自该问题被提出以来的二十年间(诚然不是在文献中……),除了相对于基概形~$X$互素的“$\ell$-adic系数”这唯一的例外之外,至今没有一种合用的概形上同调的系数“好范畴”被成形构造出来(sic!);而且那项工作,在三角范畴框架下(六运算形式主义不可或缺的框架)由Jouanolou完成,却从未发表。我自己甚至从未手持过这位学生博士论文的一份复本\fnote{381}{(**)}{Jouanolou的博士论文工作缺乏真正的信念(这使他区别于我离开前的所有其他学生),一拖再拖,答辩直到1970年之后才举行。无论是答辩委员会的组建还是答辩本身,我都不记得曾得到通知,更谈不上被联系过。答辩后Jouanolou也未觉得有必要给我寄一份论文复本。去年我写信向他索取,他(不带解释地)告知我深表遗憾已经一份不剩了……(5月12日)我的记忆在此出现了偏差——事实上Jouanolou的论文答辩早自1969年就举行了!关于此事的详情,见最后的注(在写下这些文字时尚未写就)“第六颗钉子(钉进棺材)”。}。这些都是普遍背离我六十年代所开创的基础纲领的鲜明迹象\dmark{}\origpage{821},我当年绝未料到这项工作不会借着已取得的势头继续推进,而是在我离开数学舞台后即刻被硬生生斩断(或被“电锯截断”……)。

当素数~$\ell$~在概形~$X$~上是\textbf{幂零的}时,比如所谓“$\ell$-adic系数”范畴~$\mathbb{Z}_\ell^*(X)$\fnote{382}{(*)}{在选定理论的基环说明之后的星号$*$意味着我们不仅处理朴素的可构层(此处为合适意义下的$\ell$-adic层),而且处理由层构成的“可构”复形,它们是合适三角范畴的对象(哪怕可构层范畴$\mathbb{Z}_\ell(X)$已知,三角范畴的描述也非易事)。当我们想处理Motif时(通常是指“iso-motif”,即差同源意义下的Motif,构成具有$\mathbb{Q}$-系数的范畴),用于“实现”这类(iso)motif的自然系数范畴自身必须具有$\mathbb{Q}$-系数,人们将采用$\mathbb{Q}_\ell(X)$、$\mathbb{Q}_\ell^*(X)$。当人们想同时处理所有的$\ell$时,最自然的是处理“adelic”层(或复形),其基环是adele环$\mathbb{A}$,它是由所有系数范畴$\mathbb{Z}_\ell^*(X)$的积张量$\mathbb{Q}$得到的。人们会注意到,当素数$\ell$不与概形$X$互素时,在$X$上“$\ell$-adic系数”的定义中,$\mathcal{O}(X)$的幂零元不能被忽略——它们干预在$X$于$\ell$处的奇异性中。对于$X$上的adelic系数更当如此,这使它们接近于下一段中将要谈到的De Rham-Mebkhout系数(目前尚属假想)。必须理解,这两种主要的系数类型,adelic系数与De Rham-Mebkhout系数(后者被赋予下文所提及的所有丰富结构),是同一个\textbf{Motif}的两种不同描述(弱化描述)或“实现”,由二者之一极其紧密地圈定。关于这种“忠实性”,我早在六十年代就提出了若干猜想,类似于所谓的“Tate猜想”(Deligne将它们连同其余一切一并埋葬了……)。我打算在致力于数学思想的《数学反思》卷中再回到这里。人们感受到这两种系数类型之间的强烈亲缘关系(adelic与De Rham-Mebkhout,后者“差同源”)。后者优于前者之处在于,使其在某些方面显得“更精细”的是,它们的自然基域是$\mathbb{Q}$,而对于前者则是大得多的adele环。},应当正是“结晶系数”范畴,带有Frobenius态射~$F$~与相应的\textbf{过滤}。这一三角范畴的成形构造,更不用说其上的六运算,仍在等待有人着手进行。至于通过我六十年代末隐约预见的一个“神秘函子”,将上述“普通的”$\ell$-adic情形与前述“结晶”情形进行“粘合”,以获得对~$\ell$~无限制的系数范畴~$\mathbb{Z}_\ell^*(X)$~的定义,即便在最简单的非平凡情形~$X=\mathrm{Spec}(\mathbb{Z}_\ell)$~下仍未完成\dmark{}\origpage{822}。至于一般概形上的De Rham-Hodge系数~$DRHdg^*(X)$\fnote{383}{(*)}{(5月12日)正如后文将看到的,这一临时起意的名称与记号被证明是不合宜的,应由记号$DRM^*(X)$或$\mathrm{Meb}^*(X)$代替,对偶于$DRD^*(X)$或$\mathrm{Del}^*(X)$,分别代表De Rham-Mebkhout与De Rham-Deligne系数。后者已被其生父遗弃,并在1985年由我在充分知情的情况下重新拾起,作为Grothendieck装备库中基础成分之一(通过Mebkhout)……},我并无清晰的描述思路,Deligne也未能给出真正令人满意的界定。这里的创新思想归功于Zoghman~Mebkhout——大家都知道他是在何等恶劣的逆境下工作的,以及在他的思想之深远影响(极为有限地)被承认之后,他本人遭遇了怎样的对待。无论如何,人们终于有了一条可靠的线索,借助于对绝对基~$\mathrm{Spec}(\mathbb{Z})$~的“晶体”施加有限性、完整性与正则性条件,并或许辅以“De Rham过滤”与“按权过滤”,来着手构建~$DRHdg^*(X)$~范畴——并有望构建出不仅局限于特征零、而且在给定的正特征下能够或多或少还原出“通常的”结晶系数(sic!)的东西。非同寻常的是,我似乎是世界上唯一感受到这项任务的人——Zoghman~Mebkhout本人在尝尽苦头之后,似乎也再无心思为\textbf{他的}哲学的奠基问题多花哪怕一天时间!我何苦多言,既然我看到Deligne在Hodge理论上以身作则,硬生生截断了“在我那个时代”曾激励他并催生出大有前途(却未兑现……)之进路的那股冲劲。我很清楚,复代数簇~$X$~之上的Hodge系数形式主义(甚至尚未成形)理应或多或少包含在刚才按我六十年代的语言习惯称之为“De Rham系数”或“De Rham-Hodge系数”的形式主义中——以唤起De Rham\textbf{过滤}对象与相应的分次对象(所谓的“Hodge对象”)之间的联系。鉴于Mebkhout哲学在理解这些系数范畴(诚然目前仍属假设)中的决定性作用,无疑最好将它们称为“\textbf{De Rham-Mebkhout系数}”(记号~$DRM^*(X)$),或勉强称为“De Rham-Hodge-Mebkhout系数”($DRHM^*(X)$)。当~$X$~在复数域~$\mathbb{C}$~上为有限型时,人们应当能够从De Rham-Mebkhout系数出发,或多或少“重言式地”复原出假设的Hodge系数范畴\dmark{}\origpage{823}$HDG^*(X)$(我绝不称之为Hodge-Deligne范畴,因为Deligne似乎竭尽全力隐藏这个问题,绝非揭示它!),以及其上的六种运算——只需在De Rham-Mebkhout系数上简单添加一层被称为“Betti结构”的额外结构(具有超越性质)。因此在我看来,目前摆在代数簇上同调“自然系数范畴”描述面前的主要问题\fnote{384}{(*)}{在某种意义上,这些问题是为Motif的yoga赋予其应有的精度与一般性(正如我自六十年代起所设想的那样)的预备性问题(或被默认解决的问题)。}如下:

1. 对于给定的素数~$\ell$~与\textbf{任意}概形~$X$(不必然“与~$\ell$~互素”),给出~$\ell$-adic系数范畴~$\mathbb{Z}_\ell^*(X)$~的描述,以及这些系数的六种运算形式主义。(这一问题或多或少等价于“神秘函子”问题。)

2. 对于任意概形~$X$,给出“De Rham-Mebkhout系数”范畴~$DRM^*(X)$~的描述,或对相对概形~$X/S$~给出类似范畴~$DRM^*(X/S)$~的描述(通过定义约定~$DRM^*(X)=DRM^*(X/\mathrm{Spec}(\mathbb{Z}))$),以及这些系数的六种运算形式主义。

对于问题2),根据人们希望在这些系数中引入的结构丰富程度,可能有几种变体。无论如何,“好上帝定理”(别名Mebkhout定理)先验地向我们表明(至少对于复数域上的有限型概形~$X$),必定存在Mebkhout式的结晶系数的六种变换形式主义,而无需在其上添加De Rham或/与按权的过滤。对于一般概形~$X$上的一个Motif(或一个“绝对系数”)所关联的De Rham-Mebkhout结晶复形~$K$,必然存在的第三类重要的额外结构,是对每个素数~$p$~给出一个“Frobenius态射”$K(p)^{(p)}\to K(p)$,其中~$K(p)$~表示对由模~$p$~约化得到的子概形~$X(p)$~的限制,上标$(p)$表示~$K(p)$~的“Frobenius化”,即它在Frobenius映射下的拉回。因此,根据人们打算在一个结晶复形中引入的额外结构(刚才提到的三种之中),先验地可以预见关于“De Rham-Mebkhout系数”概念的\textbf{八种}变体。唯有实地工作\dmark{}\origpage{824}才能向我们指明其中哪些变体确实能产生六种运算的形式主义。同样真实的是,出于Motif的需要,当我们旨在寻找能够尽可能贴近Motif并尽可能忠实丰富地描述其结构的简单“代数”对象时,先验地“最丰富的”系数显得“最好”。正是出于其巨大的丰富性,Hodge系数具有主要的魅力——甚至达到了这样的程度:如果Hodge猜想为真,人们可以希望在~$\mathbb{C}$~上完全重建Motif范畴,甚至对于~$\mathbb{C}$~上的任意有限型概形~$X$~亦然。

这提醒我注意到,某些结构可能是“多余的”,并从其他结构中推导出来(但确实是以一种极其隐蔽的方式,以至于很难用平实的语言明确表达)\fnote{385}{(*)}{作为朝同一方向的评注,我在此指出必须注意施加于某一类型“上同调系数”所关联的全部结构之上的或隐或显的相容性条件。我在此特别想到那些在“可Motif化”系数(即源自Motif的系数)上自动满足的相容性(具有或多或少的代数性质)。如果人们希望拥有一套“六种运算”的形式主义(哪怕独立于尽可能贴近Motif的愿望),很可能必须将它们强加于所考虑的系数之上。我特别想到Mebkhout理论中的完整性与正则性条件,以及(如果将De Rham过滤作为额外结构引入)联系De Rham过滤与Gauss-Manin联络的相容性条件。这些例子在我看来相当清楚地表明,描述上同调系数的“好范畴”这一根本任务,伴随着彻底探索和充分利用迄今为止在“代数簇上同调”上构想的所有结构以及联系这些结构的深层关系的约束。这从一开始就是Motif的yoga的主要宗旨:在一个差异性背后显现一份\textbf{统一},并同时提供一条可靠的线索,让人在不可穷尽的丰富性中认清方向。}。例如,在概形~$X$~关于另一个概形~$S$~的(相对)De Rham上同调上,我(在六十年代末)\fnote{386}{(**)}{(5月2日)实际上是在1966年。}揭示了一种无曲率的(绝对)典范联络的存在,我称之为\textbf{Gauss-Manin联络}。由此可知,Deligne赋予~$\mathbb{C}$~上光滑概形~$X$(乃至赋予~$\mathbb{C}$~上任意有限型概形~$X$)的Hodge-Deligne结构,典范地配有这样一个相对于素子域~$\mathbb{Q}$~的联络。只要Motif上同调本身已经能从其“Hodge实现”中复原出来\dmark{}\origpage{825},这意味着在任何可以称之为“Motif的”或“代数的”(即源自一个Motif)Hodge结构上,都将存在这样一个典范的Gauss-Manin联络。同样不难描述关联于一个Hodge-Deligne结构的其他更为精妙的典范结构,其存在性“源于Motif”:某些Galois群在~$\mathrm{Bet}(K)\otimes_\mathbb{Z}\mathbb{Z}_\ell$~上的作用(其中~$\mathrm{Bet}(K)$~是Hodge-Deligne结构~$K$~底层的“格”),以及在“模~$p$~约化”上的“Frobenius结构”(对几乎所有~$p$)。正是这种表面上毫无关联的结构的丰富多样性——其隐藏的纽带正是所有这些结构共有的\textbf{“motif”}——这种丰富性于我曾代表(且依然代表)代数簇上同调这一主题的特殊魅力,以及“Motif”的魅力,它们如同赋予这一具有无数变奏的主题以生命和意义的共同精致旋律\fnote{387}{(*)}{(3月26日)在对我昨天笔记开头(n\textsuperscript{o}~176)所涉及的各种类型上同调系数(用来“圈定Motif”)的(紧密相关的)问题以及一个上同调类(来自一个代数闭链)必须满足的“代数条件”进行了简短反思之后,我决定从《数学反思》第3卷(包含《收获与播种》之后的文字)开始,着手写一篇关于Motif、“系数”与标准猜想的实质性文本。我现在相信我已经掌握了对一个概形上的“该”Motif范畴进行成形描述的原则,至少在关键情形(通过极限过渡应当能够归结到的情形)下:在绝对基~$\mathbb{Z}$~上的有限型概形的情形。作为相对于我六十年代思想的唯一新成分,我引入了“Mebkhout范畴的理念”(由“好上帝定理”所表达)。此外,我假设“神秘函子”的问题已经解决(目前尚未解决!),它在我计划给出的完全描述中起着决定性作用。}。

如果说除我之外还有谁听到了并感受到了这首旋律,并在它喷涌展开之际长久地沉浸其中,那正是Pierre~Deligne。如果说我曾向谁托付过一样活着的东西——一样精致而强健、我倾注了自己最好心血、多年来以我的力量与爱所滋养的东西——那就是他。那是一样生来便要在光天化日之下绽放、生长与繁衍的东西——一颗种子,随时准备传递其内在的生命。昨天和今天的短暂接触,是一场与我久已失去联络之物的\textbf{重逢}——并非与词语或概念、亦非与僵死之物的重逢,而是与一样充盈着浓烈\textbf{生命}的东西的重逢\dmark{}\origpage{826}。它让我再次掂量出:我所留下的这件“东西”足够宏大、足够深邃,足以启发一位全副身心投入其中的数学家的整整一生,以及他身后的其他数学家——因为他的生命无疑不足以穷尽这项任务\fnote{388}{(*)}{(3月26日)我现在看来可能高估了这项任务的规模(但绝非其深度!)。关于这一点见同日的上一条脚注。}。

这是一种奇怪而令人欣慰的巧合:这次相遇发生在我刚刚经历了另一场同样不期而至的“相遇”之时——读到了我朋友的文字,在其中他避而不提其名,恰恰谈论着我托付于他的所有事物中最切中我心怀的那件东西:“如今人们对它的认识,大概不过剩一副模糊的骸骨……”

这几句话在过去的这三天里一直萦绕在我心头。我分明认出了那种自负——属于那个“任何东西都不够美丽以致他不屑为此欣喜”之人的自负。而且,并未刻意去想,关于“\textbf{坟墓}”的记忆又浮上心头\fnote{389}{(**)}{关于这种强烈且长久未曾言表的印象——它在我从巴黎“归来”并将《追寻栈》托付给Deligne之后一直纠缠着我——见“坟墓”那条注(n\textsuperscript{o}~71)。}。同一种印象在我心中复活了,带着同一个无声而执拗的意象。那件曾为我所珍爱的活着的事物,我当年曾以为托付给了充满爱意的双手——而在我失去它的这十五年间,它却在远离风、雨和阳光的恩泽的一座坟墓里衰败腐烂。今天我发现它已毫无血色,“一副骸骨……”,成了一个乐于\textbf{利用}它却绝不愿\textbf{付出}自己的人那居高临下的鄙夷对象。

\section{5.~四种运算(在遗骸上)}

\subsection{(0) 侦探——或玫瑰人生}

\notehead{167$'$}{}

\dmark{}\origpage{827}(4月22日)原定接在这里的注,很早以前就预定名为“四种运算”(这个名字将在下一条注的开头得到详尽的解释\fnote{390}{(*)}{(5月12日)在把原来那条旧注“沉默”(n\textsuperscript{o}~168)一分为四之后,“下一条注”指的是注~n\textsuperscript{o}~167$''$,即“四种运算——或一场调查的‘梳理’”。})。我原本打算把这一“梳理”工作——对一场当时在我看来已经结束的调查的梳理——写成一条注,顶多两条。从那时起已经过去了近两个月,鉴于意外波澜的涌现,直到此刻我甚至还没能完全巡览这个主题。在整整一年之后,发现葬礼的那一幕充满惊奇的剧本仿佛又以更高的调门重演了一遍!

最终,在目录中,出名的“四种运算”所指称的不再是一两注,而是一整套卷帙浩繁——恐怕略显芜杂——的\textbf{三十}条注与子注\fnote{391}{(**)}{(5月12日)自写下这些断言性的文字以来,这一数目又有所增加……}。它们分成八个部分(1)到(8),冠以(我希望是)富有提示性的名称——从(1)“藏宝”到(8)“第六颗钉子(钉进棺材)”。在此过程中,我不得不对四条注作了彻底的重构\fnote{392}{(***)}{(5月12日)这些注的篇幅已膨胀得难以容纳,最终每条又分成数条,成为注n\textsuperscript{o}s~168(i)--(iii)、169(i)--(v)、170(i)--(iii)、171(i)--(iv)。},它们曾构成“四种运算”在2月26日至3月1日期间写下的“初稿”。在3月22日(一个月前)那条注“门槛”(n\textsuperscript{o}~172)的开头,我就违背在《收获与播种》其余地方一贯秉持的精神作了解释。

所涉及的四条注是:\textbf{“沉默”,“操演”,“分赃”,“封神”}(n\textsuperscript{o}s~168, 169, 170, 171)\fnote{393}{(***)}{(5月12日)这些注的篇幅已膨胀得难以容纳,最终每条又分成数条,成为注n\textsuperscript{o}s~168(i)--(iii)、169(i)--(v)、170(i)--(iii)、171(i)--(iv)。},相继致力于勾勒一幅\dmark{}\origpage{828}全景,涵盖围绕我的著作、继而是Zoghman~Mebkhout的著作展开的四场“大操演”——隐匿与侵吞的操演。我建议读者最初只读这四条注,略过脚注(这里的脚注比《收获与播种》任何其他部分都更繁密),并略过正文中引用的子注(它们同样异乎寻常地繁多且厚重)。他随后可以沿着这一势头继续阅读接下来的四条主要注:\textbf{“门槛”“家庭相册”“节节加码(2)”“殡仪馆——为科学服务”}(n\textsuperscript{o}s~172--175),它们已经全无技术性可言。

好奇的读者若想更细致地了解这“四种运算”曲折迂回的迷宫,第二遍阅读时可以把脚注与子注也包括进去;甚至(若他尚未读过葬礼第一部,或觉得需要唤醒阅读的记忆),可以随时(一如我自己常做的)回头查阅文中大量指涉的葬礼I(或“中国皇帝的新衣”)的段落。

构成“四种运算”(或描述并评议它们)的三十条注,每一条的核心内容都没有技术性质。依我看,任何有兴趣、有头脑的读者,即便完全不是代数簇上同调的专家,甚至不是数学家、谈不上多少“科学素养”,也能读懂。对于仍然犹豫、怕一头栽进“诈骗艺术”种种奥秘的人,我尤其推荐以下子注:它们的内容在我看来最为丰厚,其意义显然超出“拆解”那些时而荒诞绝伦、却总是精心布置(专供巴不得受骗的人……)的“勾当”所能带来的兴趣。这些子注是“驱逐”(n\textsuperscript{o}~169$_1$),接着是“正经数学……”“……与无稽之谈”“遗产——或骗术与创造”(合起来是“公式”这一总题下五条子注中的头三条);最后是“封神”(n\textsuperscript{o}~171)那条注里关于Zoghman~Mebkhout奇异经历的四条子注:“眼光的绽放——或闯入者”“黑手党”“根”“掠夺的空白授权”(n\textsuperscript{o}s~171$_1$至171$_4$)。因此,我特别提请读者留意的是这八条子注(总数二十一条中的八条\fnote{394}{(*)}{(5月12日)其间已增至二十七条,还不算“第六颗钉子(钉进棺材)”(其中又有七条愉快可口的注)。})。

至于另外十三条子注,那些对其“文献价值”不感兴趣的读者,不妨在闲暇之余,权当阅读一部扣人心弦的侦探冒险传奇:在其中,一位半路出家的业余侦探(即区区在下)追踪并搜集种种“蛛丝马迹”——有的微妙难寻、若隐若现,有的则庞大显眼到无人再看得见的地步;而这些线索最终拼合为一幅色彩浓烈且不容置辩的(世态)风俗画卷:其中一位“面带微笑、和蔼可亲的当代凡尔杜先生(别名朗德吕)”\dmark{}\origpage{829},在四邻所有善良市民温柔(乃至赞叹)的注视下,对他那些纯洁无辜的受害者施以大卸八块与焚骨扬灰的极刑。他们久已习惯了这种稍显特殊的味道,显然不再让任何人感到不适。甚至不止一人效仿起他那位讨人喜欢又机敏过人的邻居,烟囱们争先恐后地嗡嗡作响、油烟四溢。

彻底开了眼的“侦探”,唯有蹑手蹑脚地退场:显见得,此地达成了一致的赞同,而一切都好极了,在这最佳的世界里……

\subsubsection*{四种运算——或一场调查的“梳理”}
\tocsubsub{四种运算——或一场调查的“梳理”}

\notehead{167$''$}{}

(2月26日)\fnote{395}{(*)}{这条注原本应当叫做“四种运算”,紧接在“墓上的旋律——或自负”(注n\textsuperscript{o}~167)之后,比前面那条(具有导论性质的)注“侦探——或玫瑰人生”(n\textsuperscript{o}~167$'$)早了近两个月。我建议先读后一条。}看来我对葬礼已经作了或多或少的巡视。当然是不完全的巡视,而且是暂时的——但就目前而言,我认为自己不会走得更远了。我感到如今需要的是抽身抽离的纵深,现在是收尾的时刻了。剩下的是盘点我在写作《收获与播种》这一即兴禅观过程中所学到的东西。

关于葬礼的反思构成了我绝大部分的工作。这一反思在两个截然不同的层面上有序推进。首先,在极其必要的“敬重之举”——即双重注“我的孤儿”与“对一份遗产的拒绝——或一个矛盾的价格”(n\textsuperscript{o}s~46, 47)之后,是对“全部辉煌中的”葬礼的渐进发现。七八年来我确实嗅到了这股风气——面对一部著作和某种风格的“轻微嘲讽之风”,以及面对那些表面上仍受其启发并在某种方式上\dmark{}\origpage{830}“顶着我的名字”的人的同样谨慎而无可挑剔的“闭门羹”。这是通过一种风尚和一个“无懈可击的共识”所体现的葬礼面向,在“掘墓人——或全体教众”以及先前的诸注(n\textsuperscript{o}s~93--97)中得到了审视,构成了第X行列即“灵车”。这一面向在过去岁月中由于未能花功夫进行周详的反思而一直处于弥散状态,在工作过程中得到了相当大的澄清,但对我而言并未增添真正崭新的事实。

相反,崭新的事实——我去年4月19日首次面对的事实,或者如果你愿意称之为“奇特事件”——是围绕我的著作、以及在离开数学舞台后唯一承担起“Grothendieck继承者”这一吃力不讨好且凶险角色的数学家Zoghman~Mebkhout的著作所展开的某场大规模行动。

去年4月19日所做的发现(1982年的Lecture Notes 900卷,在十二年的死寂沉寂之后\fnote{396}{(*)}{(4月19日)关于这“十二年”的更正,见子注“先期起棺”,n\textsuperscript{o}~168$^1$。}Motif在其中被起出,且完全未提及我这个人),成为了狭义上的调查的起点:一场关于我的著作所遭受命运的调查,首先是那些作为其首要和主要托管人的人——我的学生们——所给予的命运。这项调查揭露了大量事实,一个比一个出乎意料,随着时日推移,它们汇聚成一幅关于葬礼究竟为何物以及其主要行动者是谁的某种外部图景。这幅图景无疑并不完整,但它在极其精确和无可争辩的细节上足够丰富,足以满足我在这方面的全部好奇心。这是我刚才提到的反思的两个“层面”中的第一个。它基本上对应于关于葬礼反思的“第一口气”,从4月19日持续到6月10日左右,并以“疾病插曲”告终。

这也大体上对应于我笔记中的“葬礼I”(或“中国皇帝的新衣”)部分。此外还必须加上5月12日的“葬礼悼词(1)——或恭维”(n\textsuperscript{o}~104),它(无疑带有一点随意性)被划归到后续也是最后的行列“葬礼仪式”中,属于“葬礼II”的一部分。我还将紧接引述之注的下一条注——即“葬礼悼词(2)——或力量与光环”(n\textsuperscript{o}~105)\fnote{397}{(*)}{这条注早在5月12日写完前一条注“葬礼悼词(1)——或恭维”的次日就已预定。我那时意识到我刚刚略加审视的文本是一座真正的富矿,远未挖掘殆尽……(关于葬礼悼词的一些细节,见“封神”那条注的开头,n\textsuperscript{o}~171)。}\dmark{}\origpage{831}——归入这构成反思“第一层面”的“调查”之中,并且它还在下一条注“肌肉与肚肠(阳葬阴(1))”(n\textsuperscript{o}~106)的评注中延续。最后这两条注写于9月底至10月初。同样,在“葬礼悼词”谱系中,即在Deligne或多或少提及我的(极其罕见的)书面文献谱系中,人们可以将Deligne的小传最近所引发的两条注归入这项调查,即“为一副模糊的骸骨作安魂曲”与“信仰表白——或伪中之真”(n\textsuperscript{o}s~165, 166)。最后,还应加上“把一个个~i~点清”(n\textsuperscript{o}~164),它给出了一系列说明(主要是事实细节),大部分由Deligne本人在去年10月访问我处时提供\fnote{398}{(**)}{关于这次来访,见“义务已尽——或真话时刻”那条注(n\textsuperscript{o}~163)。}。

在结束了三个多月的一切智力活动的疾病插曲之后,反思的“第二口气”(或我刚才谈到的“第二层面”)是由理解4月和5月的调查所揭露的这一整套事实(其中某些确实非常重大甚至令人难以置信)的\textbf{意义}的努力所驱动的。这一反思的核心部分是“阴阳之钥”,在很大程度上独立于严格意义上的葬礼主题,尽管后者定期重现,每次都重新启动一场关于我个人、我的生活以及普遍存在的禅观。

显而易见,反思的这两个层面——“调查”与“禅观”——绝非相互独立或截然分开,而是彼此交融渗透。具体而言,这体现在贯穿葬礼第一部分的努力之中——去\textbf{理解}我日复一日所发现之事的意义;同时也体现在第二部分中依然出现的事实细节,它们补充了在初步“调查”中已获得的内容。

我目前的用意,是对在整个调查过程中逐日浮现的\textbf{事实}\dmark{}\origpage{832}作一个“盘点”或大纲式摘要——这些事实我至今从未劳神以稍具连贯性的方式加以整饬。因此这将是对我目前所知的针对我的著作\fnote{399}{(*)}{根据我所知的事实,这完全局限于我1955年至1970年间致力于发展关于概形上同调与(余)同调代数思想的著作部分。}以及Mebkhout著作的“大规模行动”的一场\textbf{梳理}。根据承受代价的是后者还是前者,以及根据被作为目标的是我著作的哪一部分,我实际上在其中区分出\textbf{四种}主要运算(总而言之即“四种运算”),我首先想对它们进行回顾。事实证明,它们在反思过程中引起我注意的顺序,也与它们在1970年我“离开”之后(甚至在此之前)启动的时间顺序相重合(最后两项仅有一点微小的倒置)。

\subsection{(1) 藏宝}

\subsubsection*{(a) 沉默(“Motifs”)}
\tocsubsub{(a) 沉默(“Motifs”)}

a1. \textbf{“Motifs”背景}

\notehead{168(i)}{}

\textbf{I~“Motifs”行动}

在从Serre的某些思想中获得启发、并出于为复代数簇的经典Betti上同调的各种已知(或预感到的)纯代数“化身”寻找一个共同“原则”(或“motif”)的愿望之下,我在六十年代初引入了“motif”的概念。在整个六十年代,尤其自1963年起\fnote{400}{(**)}{1963年是平展上同调(在1963/64年的SGA 4讨论班中发展)全速“启动”的一年,它终于为Motif反思的磨坊带来了充沛的水流,而此前这些反思多少带有一点思辨的意味。正是在接下来的这一年,我发展了“Motif Galois群”的形式主义,其详尽的概念基础(按照我提交给他的理论纲领)在N.~Saavedra的博士论文中得到发展,该论文直到1972年才发表(Springer Verlag, Lecture Notes n\textsuperscript{o}~265)。},在我整理基础理论的任务之余,我围绕这一主题发展出了一套既丰富又精确的“yoga”(或“哲学”)。这一宏大的理论当时仍属猜想性质,并且多半在未来的几代人中仍将保持为猜想性质\fnote{401}{(***)}{(8月8日)现在看来,这套理论并不像我曾以为的那样遥悬于“地平线上”——只要人们终于肯投身于它!关于这一点,见3月27日“守财奴与行将就木者”(n\textsuperscript{o}~177)中的评注。},然而它在当时(乃至直到今天)却提供了一个极可靠的向导,让人在涉及代数簇上同调的一切情境中认清方向,无论用于预见“人们有权期待什么”,还是用于启发该引入的“正确概念”,有时甚至是用于提供通向证明的进路。我在《收获与播种》引言\dmark{}\origpage{833}(“沉默的终结”,第XVIII页)中对此说道:

\begin{quote}
“在我有幸发现并公之于世的所有数学事物中,Motif的这一实在于我依然显得最为迷人,最富于奥秘——正处于‘几何’与‘算术’深层同一性的核心。而这一长久被忽视的实在所引向我的‘Motif的yoga’,也许是我在数学家生涯的这第一时期\fnote{402}{(*)}{如果我在此对‘我数学家生涯的这第一时期’作此限定,那是想到了‘阿那贝尔代数几何的yoga’,它在我看来具有相当的深度与广度。在《纲领概要》中对此有少许涉及,该文将作为《数学反思》的一部分附于《收获与播种》之后。}所提炼出的最强有力的发现工具。”
\end{quote}

除了为域上半单Motif范畴给出一种可能的显式构造(在众多其他构造之中)的临时草图之外,我在个人笔记中围绕这一主题所发展的思想一直停留在口头交流的阶段。我当时太专注于编写基础理论著作的大量其他任务\fnote{403}{(**)}{这首先是指EGA(代数几何原理,与Jean~Dieudonné合作)和SGA(布瓦玛丽代数几何讨论班),后者或由我单独撰写,或(特别是与学生们)合作撰写,依据我独创的指导思想与主控框架。在1959年至1969年间,这些无一例外成为标准参考文献的著作平均“产量”约为每年一千页。这一基础工作在我离开数学舞台之日便戛然而止。关于这一点见“阴,作为仆人,与新主人”(n\textsuperscript{o}~135)。},以至于抽不出所需的几个月闲暇来充实手稿笔记,以便为在我心中发展起来的内在视野编写一份足够“深入”以致在我看来可以发表的整体“主控大纲”。从1965年起直到1970年我离开数学舞台,我进行Motif(及其他)沉思的特选对话者,也是唯一完全领会了Motif的yoga并感受到其全部深远影响的人,就是Pierre~Deligne。

关于“Motif的yoga”的详细说明(比引文所属的引言部分更为详尽),读者可在“我的孤儿”(n\textsuperscript{o}~46)末尾找到,尤其是(关于该yoga的起源)在“一个梦的回忆——或Motif的诞生”(n\textsuperscript{o}~51)中。关于“Motif的yoga”融入六种运算形式主义(自我的离开以来直到今天依然被我的上同调弟子们\dmark{}\origpage{834}所忽视,尽管它是同调代数中的基础结构……),见“墓上的旋律——或自负”那条注(n\textsuperscript{o}~167)。关于围绕权的yoga(构成Motif的yoga的关键成分之一)与Hodge-Deligne理论(直接源于后一yoga)的思想亲缘(在文献中被彻底抹杀),见“把一个个~i~点清”那条注n\textsuperscript{o}~164(第II~4部分),以及紧随其后的子注(n\textsuperscript{o}~164$_{1}$)。

\subsubsection*{a2.~葬礼……}
\tocsubsub{a2.~葬礼……}

\notehead{168(ii)}{}

“Motifs”行动首先并在我离开数学舞台后即刻表现为:对Motif的yoga与“motif”一词本身的系统性\textbf{抹杀};继而在十二年的沉默之后\fnote{404}{(*)}{(4月8日)关于这“十二年”的更正,见本注“沉默”的子注“先期起棺”(n\textsuperscript{o}~168$^4$)。},随着该yoga的一个狭隘版本的起棺(于1982年),表现为对我这微不足道的亡故之人的\textbf{抹杀}——抹杀我与该yoga有任何瓜葛的事实。

对该yoga以“权的yoga”形式进行的首次公然抹杀,早在1968年就已发生,也就是早在我离开之前,出现在Deligne发表在《数学出版物》上关于谱序列退化的文章中。这首先在“罐头里的权与十二年秘密”那条注中谈及(写于发现起棺的“难忘卷册”之前),并在紧随其后的“驱逐”那条注的开头作了详尽阐述(注n\textsuperscript{o}s~49, 63)。

这场试探性的抹杀,在未遭任何反应之下\fnote{405}{(**)}{这种反应首先本该来自我。尽管那一篇文章呈现中的不诚实在我看来昭然若揭(参见引注,n\textsuperscript{o}~63),但我未能正视它,面对拿到文章读到有关段落时的“轻微不适”。关于我在反思过程中察觉到的自身某种顺从或暧昧所起的作用,见“暧昧”那条注,n\textsuperscript{o}~63$'$。至少在意识层面上,Deligne或我的任何其他学生可能存在不诚实念头从未掠过我的心头;或者说,我把它压制了下去,即便在不诚实显而易见并通过那从未被指认的“不适”向我发出警示之时。},在Deligne陈述其1968/69年发展的Hodge理论优美推广的Hodge~I、II、III文章中继续并加剧。虽然这一理论直接源于Motif的yoga(正如前文所回顾的),但在Hodge~II或Hodge~III中对此毫无提及——考虑到Hodge~II构成Deligne的博士论文,而他曾在我身边度过了其成长关键岁月\dmark{}\origpage{835},这一事实尤为触目\fnote{406}{(*)}{在Deligne与我之间存在某种默契,以抹杀其作为学生的身份——他是如此耀眼,以至于我无法自称曾是他的“导师”。我在“格格不入者”(n\textsuperscript{o}~67$'$)中揭示并检视了这一默契。}。

至于在1970年尼斯国际数学家大会上的简短“预告”Hodge~I,Deligne仅限于以半行晦涩的“暂停!”式提及指称“Grothendieck的猜想性Motif理论”(并同对Serre的一条假文献并列在一口气中,显然旨在掩人耳目\fnote{407}{(**)}{所涉及的是Serre关于Weil猜想的Kähler类似物的文章,它曾启发我关于“标准猜想”的工作。这是一篇出色的文章,在此绝非要加以贬低。但Deligne本人也很难解释这篇文章在何种程度上曾是其Hodge理论工作的“源泉”——无疑也从未有人想到去问他。近距离目睹了Deligne思想的萌发,我深知其源泉何在(参见前引注n\textsuperscript{o}~164$_{1}$),而且绝非Demazure关于Motif定义的初级ABC报告!Deligne引用这篇报告作为“Grothendieck的猜想性Motif理论”的凭据,以便给任何不深究的读者留下印象(而充分知情的人本就不多……),即所谓的“猜想性理论”仅限于Demazure报告中所写的内容,从而利用了该yoga在文献中缺乏更详尽发表痕迹的状况。})。这一做法在1974年温哥华国际数学家大会上关于“权的yoga”的报告中得以延续,会上Serre与我的名字都不再被提及。在这次交流中,如同在1970年尼斯大会的Hodge~I中一样,Deligne绝口不提他得自我的该yoga在Motif语境下(该语境被严格保密)的重要组成部分:权概念在“六种运算”下的行为,首先是通过$Rf_!$与$Rf_*$。这是众多例子中的一个,体现了一种如今司空见惯的做法,Deligne在我看来是其最早的推行者之一:即将某一数学领域中提出的“大问题”的知情权垄断给少数“圈内人”(甚至仅限其本人),以确保绝对霸权,而不是将其置于科学界共同支配之下以启发每个人\fnote{408}{(***)}{关于这种新心态(在写下这些文字时我从未遇到过痕迹),见“阴,作为仆人,与新主人”那条注,n\textsuperscript{o}~135,以及紧随其后的“操演”那条注末尾(4月28日)(x)。正是这种心态促使我给围绕我著作的前两场“大运算”的整个集合(n\textsuperscript{o}s~168至169$_8$)冠以“\textbf{藏宝}”之名。(x)这一末尾成为了“藏宝”那条注(n\textsuperscript{o}~169(v))。}。据我所知,\dmark{}\origpage{836}这一问题直到Deligne1980年在其“Weil~II”文章中(在$Rf_!$情形下)解决之前一直无人触碰,当然文中对我只字未提(正是向他传达了Motif语境下的相关猜想,而他所处理的$\ell$-adic语境不过是其反映,如同De Rham-Hodge系数语境一样……)。

在我了解(非常零碎地)Deligne工作或能够对其做出评价的范围内,我认为可以断言:他得自我的Motif的yoga构成了其全部工作的主要灵感源泉。他通过保持对Motif概念的绝对沉默直到1982年\fnote{409}{(*)}{(4月8日)关于更正,见前引子注“先期起棺”(n\textsuperscript{o}~168$^4$)。},将这一源泉隐匿了起来。唯一的例外(若无疏漏\fnote{410}{(**)}{如上一条脚注中所解释的,这一“暂停!”式提及的目的绝非为了承认或告知,而是为了(双重地)误导。至于Hodge-Deligne理论的思想谱系(在前引两条注中有所描述),我完全有理由相信除他之外无人知晓。})是那条“见证半行”,对于除他和我(以及或许Serre)之外的任何人来说都是无法理解的——正如稍早时他在关于谱序列退化的文章中晦涩地提及促使我猜想其退化结果“特例”的“权之考虑”(参见前注“驱逐”,n\textsuperscript{o}~63)。

\subsubsection*{a3.~……与起棺}
\tocsubsub{a3.~……与起棺}

\notehead{168(iii)}{}

随着Lecture Notes 900这一“难忘卷册”的出版,布景骤然改变\fnote{411}{(***)}{Springer Verlag, Lecture Notes in Mathematics, n\textsuperscript{o}~900, Hodge cycles, Motives, and Shimura Varieties, P. Deligne, J.S. Milne, A. Ogus, K.Y. Shih.}.Motif在其中大张旗鼓地被起出,初始yoga的一部分被公之于世。在这一卷册中,我的名字“顺带地”且仿佛出于极大偶然出现了两三次,绝不可能让读者怀疑我在其中发展的思想中起过任何作用。这些思想是以这样一种方式呈现的,即在读者心中无可怀疑:卷册的主要作者、辉煌的Pierre~Deligne刚刚发现了它们,并将它们新鲜端出。诚然,正如在尼斯或温哥华他未曾宣称是自己发现了首次在文献中阐明的权的yoga一样,这里也绝无任何一处\dmark{}\origpage{837}明确说是他发现了卷册中(表面上)首次阐述的所有这些优美思想——该卷册围绕一条确实出自他手笔的优美定理展开。这是他深谙的“暂停!”艺术,对此我首先在“暂停!”以及紧随其后的“中国皇帝的新衣”中作了评述(n\textsuperscript{o}s~77, 77$'$);亦见发现“难忘卷册”之前写下的更早的注:“下葬——或新父亲”、“新伦理——或抢夺的集市”以及“侵吞与轻蔑”(n\textsuperscript{o}s~52, 59, 59$'$)。

事实上,LN~900卷册中关于Motif的所有主要思想不仅早在六十年代就已为我所知(Deligne自1965年起完全有机会从我口中习得),而且该书的核心问题也早在六十年代末就由我提出(并且当然传达给了Deligne)。关于这方面的详情,见“把一个个~i~点清”那条注(n\textsuperscript{o}~164)(在该注的第一部分)。

正如我在《收获与播种》引言(“秘密的终结”,第XVII页)中所强调的,Deligne绝非我向其详述Motif的yoga的唯一之人,哪怕他是唯一深入将其内化之人。如果说在长达十年左右的时间里存在对这一yoga本身的彻底抹杀\fnote{412}{(*)}{根据Deligne好心提供给我的一份“Motif评注书目”,在我离开之后文献中关于Motif还出现过两次零星的工作,均直接源于六十年代的工作(在N.~Saavedra与我共同准备的论文中,以及在S.~Kleiman的一份报告中)。下一处指涉由R.P.~Langlands在1979年做出。之后便是1982年的LN~900。若无疏漏,在1982年之前Deligne发表的任何文本中都未出现“motif”一词,亦未在任何发表文本中(最多除注n\textsuperscript{o}s~165, 166中检视的传记小记之外)提及他曾从我口中习得任何东西……(4月8日)关于“若无疏漏”,见子注“先期起棺”(n\textsuperscript{o}~168$^4$)中的更正。},以及后来对我发现、发展和阐明它的角色的抹杀,这一抹杀若无我视为朋友的大批数学家的同谋是绝不可能完成的,特别是我的每一位(交换)“上同调学生”的同谋\fnote{413}{(**)}{我想可以说,我1970年之前的所有学生,唯一的例外是信女士(她自六十年代起就在越南工作),都知道(哪怕未必掌握)我关于Motif的想法,我曾在IHES就此做过一系列详尽的报告(1967年)。他们之中那些致力于代数簇上同调研究的人,在我看来对在主要“当事人”Deligne倡议下进行的下葬负有连带责任。这里主要涉及J.-L.~Verdier、L.~Illusie、P.~Berthelot,他们在所涉及的围绕我著作的其他三场运算中,表现得远比单纯的同谋更为主动。}。这一抹杀虽然服务于唯一一个人的可疑“利益”,却是由许多人共同的作为与不作为完成的。

\dmark{}\origpage{838}除Deligne和我的其他上同调学生之外,对这一“难忘卷册”LN~900负有最严重责任的,正是其\textbf{合著者},即\textbf{J.S.~Milne}、\textbf{A.~Ogus}与\textbf{K.Y.~Shih}。这些是我未曾谋面的数学家,我无意断定他们居心不良。对我而言,这丝毫不减免他们作为这部不同寻常卷册共同署名者的完全责任。

\subsubsection*{a4.~先期起棺}
\tocsubsub{a4.~先期起棺}

\notehead{168(iv)}{}

(4月8日)最近有人向我提示了Deligne1979年发表的文章“L函数的值与积分周期”(载于Proceedings of Symposia in Pure Mathematics, Vol.~33 (1979), part~2, pp.~313--346),与前文提及的R.P.~Langlands文章“自守表示、Shimura簇与Motif. 一个Corvallis童话”(pp.~205--246)刊于同一卷中。后一篇文章(而非Deligne的文章)出现在Deligne去年8月寄给我的Motif评注书目中,我曾一直以为正是Langlands这篇文章构成了在我离开后、1982年起棺之前文献中唯一谈及Motif之处(除了前文脚注中提及的Saavedra与Kleiman的报告)。

事实上,在Deligne引用的文章中,赫然有一章“第0章”,题为“Motif”,其开篇为:“我们回顾\textbf{归功于Grothendieck的Motif形式主义的一部分}”(着重号为我所加)。其阐述方式清晰表明,我为一个域上的(半单)Motif范畴给出的构造通则是多价的——文中接着指出“Grothendieck关于Motif的\textbf{定义之一}是通过……获得的”。在这方面该文是诚实的。确实,这里阐述的Motif“yoga”部分是最基础的部分,几乎在文献中早已存在(在Manin、Demazure、Kleiman、Saavedra的报告中),因此我的首创权\dmark{}\origpage{839}格外无可争议。(相反,在权的yoga以及后来的Motif Galois群中对我首创权的抹杀——以及对Serre首创权的抹杀——似乎未遭遇任何阻碍……)

正如我已经强调过的(在注“节节加码(2)”,n\textsuperscript{o}~174中),在1977年“下葬行动”达到暂时高潮(伴随“SGA~4$\frac{1}{2}$—SGA~5行动”)之后,出现了一段相对平静的时期,直到1981年反常层会议的“封神”,该会议标志着在对遗骸的瓜分中放弃了任何克制意图。(见注“封神”,n\textsuperscript{o}~171.)Deligne的这篇文章显然处于这一平静期的标志之下。我推测Langlands对Motif yoga的兴趣迫使Deligne最终泄露了Motif这一(早已公开的)秘密,而在当时的心理背景下,将死者的名字纯粹抹杀的时机尚未成熟。在随后的短短三年里,在这一对Motif的“怯懦起棺”与1982年LN~900卷册“大张旗鼓的起棺”之间,确实发生了一场令人惊叹的“节节加码”。

(4月22日)前一页所评述的(微小)发现在随后的日子里继续并大大扩充。我确实阅读了引述的R.P.~Langlands文章,并且最重要的是在次日读到了钉入我棺木的“第六颗钉子”\fnote{414}{(*)}{这是按发现顺序排列的第六颗“钉子”,但按制造年代算则是六颗中的第一颗,是由我朋友皮埃尔熟练地“打下”的,使用的是知名殡仪馆Springer Verlag GmbH(殡仪部“Lecture Notes in Mathematics”)提供的专利设备。},以我的(前学生)Neantro~Saavedra~Rivano所著、题为《塔纳卡范畴》的卷册形式出现。因此关于“Motifs行动”还有一段实质性的“后话”,我在题为“第六颗钉子(钉进棺材)”的一系列子注(n\textsuperscript{o}s~176$_{1}$至176$_{7}$)中予以展开。将这一后话移至调查“四种运算”的末尾在我看来更为妥当,因为在整场调查中浮现的新事实,尤其是在注“封神”(n\textsuperscript{o}~171)及其四条子注\fnote{415}{(**)}{(5月11日)自写下这些文字以来,引述的注已拆分为四条注,并增补了另外八条子注(n\textsuperscript{o}s~171(v)至(xii))。}中浮现的事实,对于准确安放这一“后话”并赋予其全部意义是必不可少的。

\subsubsection*{(b) 操演(“平展上同调”)}
\tocsubsub{(b) 操演(“平展上同调”)}

b1.~围绕“Weil猜想”的背景

\notehead{169(i)}{}

\dmark{}(2月27日)我现在来谈第二场“大运算”:\origpage{840}

\textbf{II~“平展上同调”行动。}

如同对待Motif一样,先用几句话说明背景将是有益的。

关于在任意域上的代数簇存在一种“上同调”理论——该理论将为这种簇(至少当它是射影且光滑时)赋予“上同调空间”,其系数域为特征零(例如$p$-adic域),其性质仿照复基域情形下(经由超越途径定义的)经典“Betti”上同调的已知性质——这一思想“在字里行间”存在于Weil著名的猜想陈述中(1949年)。无论如何,大约在1955年,Serre就是用上同调术语向我解释Weil猜想的——而且唯有在这些术语下,它们才真正具有吸引力。

当时无人知晓如何定义这样一种上同调,我甚至不确定除Serre和我之外是否还有其他人(如果这能算得上的话,连Weil本人也未必)具有其必然存在的深切信念。人们当时唯一能直接几何把握的只有$H^1$,通过阿贝尔簇理论及其有限阶点(由Weil发展),以及通过关联于非奇异射影代数簇的Albanese簇或Picard簇。对$H^1$的这一构造提示“自然”系数应当是$\ell$-adic域$\mathbb{Q}_\ell$,其中素数~$\ell$~\textbf{不同于}基域的特征。对于~$\ell$~等于特征的情形(当特征非零时),Serre在代数\textbf{曲线}情形下的初步结果提示人们应当可以取$k$(设为完美域)的Witt向量环的分式域作为系数。因此人们有望对\textbf{每一个}素数~$\ell$~拥有一套$\ell$-adic理论(对于$\ell=p$需打点折扣)——并且在合理意义下,它们应当“都给出相同的结果”。最后,当$k$为特征零时,鉴于人们拥有Hodge上同调空间(自Serre引入代数簇的“凝聚”上同调理论以来对任意$k$都有意义)与De Rham上同调空间(我受微分De Rham上同调启发引入),这些在当时提供了具有全部所期\dmark{}性质的上同调理论\origpage{841}\fnote{416}{(*)}{我早在五十年代就发展了结合于代数闭链的上同调类(De Rham或Hodge)形式主义。},并且它们应当同$\ell$-adic理论一样给出“相同的结果”。

这些问题处于我1955年至1970年(我离开数学舞台之年)间发表与未发表数学思考的核心。如果撇开我在凝聚上同调方面的工作不谈(“六运算”形式主义、Riemann-Roch-Grothendieck公式),大体可以说,我上同调工作的核心就在于给出这些问题的答案或大纲。至少在以Weil猜想为主要灵感源泉的视角下,我对上同调主题的思考具体化为四条大\textbf{潮流}或“\textbf{线索}”,紧密交织构成一张宏大织锦:

\textbf{线索~1.}——我(在合作者的协助下\fnote{417}{(**)}{发展平展上同调形式主义的主要合作者是M.~Artin与J.-L.~Verdier。$\ell$-adic系数则是在我的前学生P.~Jouanolou的论文中发展的(可惜他从未发表,我从未见过,如今已无从寻觅)。我打算在随附于《数学反思》的“历史”评注中给出关于平展上同调发展的全面描述。})发展了概形的\textbf{平展上同调}与\textbf{$\ell$-adic上同调}形式主义,对于与残余特征互素的~$\ell$,它具备拓扑空间熟悉的“离散”上同调的全部性质(乃至更多……)。除了三项技术性质的未决\fnote{418}{(***)}{这三项“未决问题”如下:
\textbf{a.}~对于特征与~$\ell$~互素的正则概形~$X$~的正则闭子概形~$Y$,“上同调纯性猜想”(平展版本)。当~$X$~与~$Y$~在正则基概形~$S$~上均光滑时(对于应用足够的情形),以及在特征零情形下(由Artin充分利用奇点消解证明),相关猜想已获证。
\textbf{b.}~更为严重的是对于诺特概形(需要时为优良概形)之态射~$f:X\to Y$,当~$f$~\textbf{非假定}为真时,关于更高直接像$R^i f_*$在$\ell$-adic“可构”框架下的\textbf{有限性定理}之有效性问题。我在基概形为代数闭域上的簇且在奇点消解与“上同调纯性”(参见a)假设下证明了该结果。然而我指出,在有限系数框架下(相对于$\ell$-adic系数而言),六运算的对偶形式主义(包括Poincaré对偶)在1963年由我在无需有限性条件下建立。
\textbf{c.}~优良正则概形上“双对偶定理”的有效性。类似情况。
Deligne(约在1973年?)对维数$\le 1$的正则概形~$S$~上的有限型概形态射给出的优雅证明大大改善了局面。在同一情形下,通过类似论证,Deligne也证明了双对偶定理。}问题之外,可以说自1963年起人们“原则上”就已拥有了\dmark{}\origpage{842}(伴随1963/64年SGA~4之后的SGA~5讨论班的展开)该上同调在平展上同调框架下的\textbf{完全构造}——以“六运算”对偶形式主义的形式。平展上同调的定义原则可追溯至1958年,我在1963年2月至3月证明了完整形式主义(包括“弱Lefschetz”型定理和平展语境下的上同调深度概念)所必需且充分的“关键结果”。

\textbf{线索~2.}——借助\textbf{Motif}的yoga,我发现了使人们能够将一个簇的各种$\ell$-adic(及其他)上同调视为该簇所共有的“Motif上同调”的不同“实现”的\textbf{哲学}。这一哲学在六十年代初萌生,伴随着直接受Weil猜想启发的“权的yoga”(以及Serre受其启发的关于簇的“虚Betti数”想法\fnote{419}{(*)}{关于这一点,见“我的孤儿”(n\textsuperscript{o}~46)的子注~n\textsuperscript{o}~46$^9$。})。它在1964年平展上同调启动的势头中丰富起来,引入了“Motif Galois群”的关键概念。

\textbf{线索~3.}——受Monsky-Washnitzer思想的启发(他们为特征~$p>0$~下的\textbf{光滑仿射}代数簇构建了常系数“$p$-adic”上同调理论),我在1968年给出了“$p$-adic上同调”的通用定义,我亦称之为\textbf{结晶上同调}\fnote{420}{(**)}{这一术语如今(且久已)由习惯确立,这一理论亦然。引导我走向这一理论的两项新思想(相对于Monsky与Washnitzer的思想而言)是:\textbf{晶体}(模等)的概念,与在“加厚”(初始无穷小加厚)之上的“生长”思想相联系;以及在所考虑的理想中引入\textbf{除幂}结构,以确保除幂“形式Poincaré引理”成立。奇怪的是,晶体的关键直觉被我的学生们弃之不顾,正如贯穿我全部上同调思考的线索——六运算形式主义一样。这在我看来是导致结晶上同调与Hodge-Deligne理论令人遗憾地陷入停滞的主要原因。此外在我看来显而易见的是,在这两种理论中,Zoghman~Mebkhout所发展的理论将扮演不可或缺的角色。我在这一方向上的建议(1978年致Berthelot)显然被置若罔闻……}。这一\dmark{}理论旨在涵盖非常数亦非局部常数的“系数”,并产生如同$\ell$-adic理论一样的“六运算”形式主义。从一开始就确立的是,对于\textbf{光滑}簇,这一上同调具有人们所期望的与De Rham上同调的联系,并且推广了Monsky-Washnitzer理论\fnote{421}{(*)}{P.~Berthelot以我的思想为起点的博士论文为光滑射影簇给出了其证明与对偶形式主义,至少足以给出有限域上此类簇的L函数的\textbf{结晶}上同调解释。然而距离掌握一套用于“结晶系数”的“六运算”形式主义,至今依然遥远。对于代数簇上的“Hodge系数”亦然!关于这一系数问题,见“墓上的旋律——或自负”(n\textsuperscript{o}~167)。该问题在六十年代对我而言十分明确,此后却被我的上同调弟子们埋葬至今。(4月23日)关于这一点见注“工地巡礼——或工具与视野”,n\textsuperscript{o}~170(i)。}\origpage{843}。

\dmark{}\textbf{线索~4.}——将\origpage{844}平展上同调及其近邻变体(涉及Zariski、fpqc、fppf等拓扑)、结晶上同调、超越语境下的“Betti”上同调以及更一般的任意拓扑空间之层上同调通过共同的“拓扑”直觉联系起来的统一几何概念,是“\textbf{site}”的概念,以及在其背后更内在、更隐秘的“\textbf{topos}”概念。自1964年起在随后的岁月里,这一概念逐渐走上前台。我在“我的孤儿”(n\textsuperscript{o}~46,第636--638页)中阐述了这一处于我工作核心、如今却被逐出几何学的概念的深远意义,在此仅引述如下段落:

\begin{quote}
“这一对概念[概形与topos]潜藏着代数几何与算术、以及拓扑学的大规模更新,通过将这些隔绝太久的‘世界’综合进一个共同的几何直觉之中。”\fnote{422}{(*)}{我在别处提议称这种尚在萌芽中的‘新科学’为\textbf{算术几何},它诞生于六十年代初Weil猜想的余波中,其中‘Motif的yoga’如同灵魂。我在此重申:‘在构建新代数几何之外,通过掌握平展拓扑,为一门仍在演进中的科学构建整体主控大纲,在我眼中是我对同时代数学的主要贡献。’}
\end{quote}

topos的语言与平展上同调形式主义在紧密相连的两个讨论班SGA~4(1963/64年)与SGA~5(1965/66年)中展开\fnote{423}{(**)}{SGA~4全面重修的第二版(三卷)于1972--73年发表在Lecture Notes n\textsuperscript{o}s~269, 270, 305中。关于SGA~5的波折,其遭受严重肢解的版本在口头讨论班结束后十一年的1977年发表于同套丛书第569卷。}。第一部\dmark{}与他人合作完成\fnote{424}{(*)}{site与topos语言的发展基于我1958年的最初想法,在M.~Artin、J.~Giraud、J.-L.~Verdier协助下完成。},除了topos语言外,发展了平展上同调的初步知识,包括对偶的关键起步陈述(六运算风格)。第二部中我几乎独力作战\fnote{425}{(**)}{唯一的例外是J.-P.~Serre作了若干关于有限群在代数簇上的作用以及与Artin导子相关联的Serre-Swan模的报告,这是我发展通用不动点公式所需的。对于所有其他报告,我是唯一的报告人。},极详尽地展开了\origpage{845}对偶形式主义,包括通向L函数上同调理论的不动点公式(构成Weil猜想的重要部分)。我在“遗骸……”那条注(n\textsuperscript{o}~88)中评价道:

\begin{quote}
“接续进行的两个讨论班SGA~4与SGA~5(于我犹如同一个讨论班)从虚无中发展了topos语言所代表的综合与发现的强有力工具,以及平展上同调这一极其精纯完备的工具——当时对其本质形式性质的理解甚至超过了普通空间的上同调理论。这构成我在数学中所做出的最深刻、最创新的贡献。这两部讨论班不可分割地联系在一起。它们在统一性中同时代表了\textbf{视野}与\textbf{工具}——topos,与上同调的完整形式主义。虽然视野在今天仍遭拒斥,但该工具二十多年来深刻更新了代数几何中最迷人的‘算术’面向。”
\end{quote}

\asterisms

\subsubsection*{b2.~四种操演}
\tocsubsub{b2.~四种操演}

\notehead{169(ii)}{}

\dmark{}\textbf{“平展上同调”行动}在于一方面将SGA~4与SGA~5中完成的奠基工作\textbf{贬损}\origpage{846}(贬为“无稽之谈”、瞎拼凑等),并顺势抹杀我在发现和发展这一工具中的角色;另一方面,\textbf{据为己有这一工具},即据为己有我在平展上同调中所发展的思想、技术与结果的发明权。这里的“受益者”又是Deligne\fnote{426}{(*)}{然而Verdier也获得了丰厚的回报,正如后文将看到的,他通过发表其“难忘文章”给肢解SGA~5发出了“开门红”,随后在1981年的反常层会议上大获其利。},正是其非凡的声望(既得益于其卓越的才能,也得益于其作为我默认“继承人”的地位)使得这样一场大规模的贬损与侵吞行动得以推行,且表面上不着一丝痕迹……

正是在1965/66年的SGA~5口头讨论班上,年轻的新来者Deligne第一次学到了概形语言、Grothendieck风格的同调代数以及平展上同调的新技术\fnote{427}{(**)}{这正是我在注“格格不入者”(n\textsuperscript{o}~67$'$)中回忆起的。Deligne正是在SGA~5中顺便学到了Poincaré对偶的‘六运算’形式主义。}——这些技术成为了他全部后续工作的基石。

在我这位前学生操纵的行动中(我称之为“SGA~4$\frac{1}{2}$—SGA~5行动”),我辨识出四种不可分割的“\textbf{操演}”:

\textbf{操演~1. 贬低}母讨论班\dmark{}SGA~4—SGA~5,斥之为“废话矿渣”及类似言辞:这在Deligne署名、名为“SGA~4$\frac{1}{2}$”(副题:平展上同调,载于Lecture Notes in Mathematics第569卷)的卷册前言中迂回且看似不经意地完成。关于Deligne学艺并获取其后续工作基础工具的双讨论班SGA~4—SGA~5所遭受的清算细节,见“白板”(n\textsuperscript{o}~67)。\origpage{847}

\textbf{操演~2. 破坏}我对SGA~5口头报告的整体撰写工作\fnote{428}{(*)}{正如我三条脚注前所指明的,对于SGA~5的全部报告,我都有详尽的手写笔记。只要分配下去,记录整理本该在两三年内完成。可现实是,这项工作被完全冷冻了十一年之久。}。通常这一整理本应在合理期限内(至多一两年)由我那些上同调学生完成,他们有幸学习了大部分技艺,并同讨论班其他听众一起成为了多年间唯一知晓这些思想技术的人。这也是他们熟悉这些实质内容与思想技术的最快方式,而在口头报告时这些内容往往让他们觉得有点“高深莫测”(当然总是敏捷过人的Deligne除外)。然而这一整理——或者更确切地说是\textbf{拒绝整理}——硬生生\textbf{拖延了十一年}——直到Deligne“亮绿灯”让Illusie最终处理这本不幸的SGA~5的撰写与出版的精确时刻(出于巧合),也就是确信它将在1977年发表在他自己手笔的一卷书\textbf{之后}的时刻!那卷书(写于1973年及随后几年)起初是为了普及平展上同调中证明Weil猜想(最后一环)所必需的“成分”,却在此机缘下被冠以奇特的名称“SGA~4$\frac{1}{2}$”。(这名字直到今天似乎除了我之外\dmark{}从未让任何人感到错愕、惊讶、乃至震惊……(169$_1$)\fnote{429}{(*)}{关于这一点,以及关于SGA这一缩写的最初真实含义,见“白板”(n\textsuperscript{o}~67)。})。\origpage{848}详情见注“绿灯”与“颠倒”(n\textsuperscript{o}s~68, 68$'$),在那里题为“SGA~4$\frac{1}{2}$”的卷册的含义开始向我显露,亦见注“沉默”与“团结”(n\textsuperscript{o}s~84, 85)。

\textbf{操演~3. 肢解原始讨论班}SGA~5,其出版版本(由我的前学生Luc~Illusie“操持”)已不过是一具被严重残害的“遗骸”。我在同名注“屠杀”(n\textsuperscript{o}~87)——关于葬礼反思中最长且最富启示性的注之一——中详尽检视了这场肆无忌惮的肢解,或者更确切地说,对托付给学生们的光辉讨论班的\textbf{屠杀}。

\textbf{操演~4. 击碎我关于平展上同调工作的统一性},这一工作由SGA~4与SGA~5两扇不可分割的门扉所体现,通过将其“切成两段”,“通过在两扇门扉之间强行插入一篇外来的、蔑视性的文本”\fnote{430}{(**)}{这句引文凭记忆引自注“遗骸……”(n\textsuperscript{o}~88)。},它顶着非同寻常的名字“SGA~4$\frac{1}{2}$”\fnote{431}{(***)}{副题:平展上同调——Pierre~Deligne著……这副题也同样言简意赅!}。\origpage{849}这个天才的名字言简意赅——关键在于能想得出来!单凭这一个名字,该卷册便自诩为平展上同调的\textbf{核心与根本文本},旨在\textbf{取代}“SGA~4与SGA~5冗长繁琐的报告”,后者“可被视为一系列离题赘述”,虽然“某些非常有趣”,但核心文本“应当允许使用者将其遗忘”。

此外,我这位天才的前学生兼朋友无需在此陷入冗长无用的辩解:单单这个简峻的名称“SGA~4$\frac{1}{2}$”\dmark{}便宣告并确立了这一文本相对于题为SGA~5的“离题赘述”而言具有无可辩驳的\textbf{在先性}(后者作为其必然结果,确实是在它之后出版的……),并同时确立了SGA~5相对于这一“在先”文本的(所谓的)\textbf{逻辑依赖性}。

关于SGA~5相对于这一伪作具有所谓的逻辑依赖性的离奇骗局,在该书的前言中被确凿宣称\fnote{432}{(*)}{我提醒读者,Deligne在上个月访问我处时,曾当面向我确认过这一点。},作者毫不脸红地宣布(且鉴于当今风气,在我之前显然无人觉得有何不妥……):“它的存在[‘SGA~4$\frac{1}{2}$’的存在]将使得不久后得以出版SGA~5\textbf{原貌}”(着重号为我所加)——请读作:处于被洗劫和大量盗窃后的\textbf{遗骸}状态……虽然我此前一周多已知晓我朋友的“Motifs”行动,但我仍花了整整两天时间(从4月26日的“白板”注到28日的“颠倒”注,n\textsuperscript{o}s~67, 68$'$)才领悟我这位高足这一显然荒唐的断言对我所代表的“谜团”的意义——并由此理解了看似平淡无奇的符号“SGA~4$\frac{1}{2}$”的含义,此前两天我甚至未曾在此停留。

相同的“逻辑依赖”骗局在Illusie为SGA~5所写的前言中被明确提示(169$_2$)\fnote{433}{(**)}{详情见本注的子注“好撒玛利亚人”(n\textsuperscript{o}~169$_2$)。}。对于不知情的读者而言,由于我的报告的迟到整理者们\fnote{434}{(***)}{(4月9日)经细致核实,所涉及的‘迟到整理者’主要是Illusie一人……}乐于在整理稿中塞满对“SGA~4$\frac{1}{2}$”的无数指涉,这一骗局显得更加可信\origpage{850}。这些指涉许多并非虚构,而是关联于原始讨论班的两篇报告(一篇由Illusie整理,另一篇——格外关键的——由Deligne整理\fnote{435}{(*)}{所涉及的是报告‘结合于闭链的上同调类,A.~Grothendieck著,P.~Deligne整理’。}),它们被不由分说地直接并入了题为“SGA~4$\frac{1}{2}$”的卷册中——既未通知我,亦未征求我的意见,而是仿佛在逝去的师父缺席下这本就理所当然属于他们一样……

\dmark{}这一强盗行径进一步使得我的前学生Deligne实现了\textbf{角色的绝妙倒错},得以在书本封面上(同样绝不与我商量……)将我呈现为他的\textbf{协作者}(为了平展上同调的发展!)\origpage{851}\fnote{436}{(*)}{这一幕安排(将我呈现为我学生Deligne的‘协作者’)格外滑稽,考虑到在整本书中Deligne独创的材料至多占三十到四十页……}——一个略显“糊涂”的协作者\fnote{437}{(**)}{在他寄给我的一份‘SGA~4$\frac{1}{2}$’提要中,我的名字甚至未被提及。},这倒是真的,但无论如何是个“协作者”……

至于那本被称为“SGA~4$\frac{1}{2}$”的盗版文本,除了两篇从原始SGA~5语境中剥离出来的报告,以及若干对算术应用格外重要的SGA~4—SGA~5结果的“摘要”,外加一章关于三角和应用的原创章节,以及最后Verdier“论文”的“第0状态”之外,它由对SGA~4—SGA~5中所发展的上同调形式主义的若干补充(固然十分有用\fnote{438}{(***)}{所涉及的是有限性结果(已在三条脚注前提及)以及关于具有有限系数的局部对偶定理。})组成。这本可以写成一篇约三十页(若算上‘三角和’一章则约五十页)内容略显杂糅的出色文章。如果我这位辉煌的前学生尚存一丝体面,将这几项补充各自放置在其所受启发并加以补充的两三篇SGA~5报告中\dmark{},本是不言而喻的\origpage{852}。然而相反,它们成了纯粹取消SGA~5第II讲的借口(得到了负责整理并以此‘弥补’的Illusie的祝福,将该讲转化为‘SGA~4$\frac{1}{2}$’中的一个附录),并顺理成章地将平展上同调的双对偶定理(由我在1963年根据五十年代发现的‘凝聚’类似物建立)重新命名为“Deligne定理”\fnote{439}{(*)}{\textbf{双对偶定理},或‘局部对偶定理’(两个名称均出于我手),在SGA~4$\frac{1}{2}$中被冠以Deligne之名,四年后Deligne又慷慨地将其转赠给他的朋友Verdier……}。

\subsubsection*{b3.~节节加码的各幕}
\tocsubsub{b3.~节节加码的各幕}

\notehead{169(iii)}{}

\dmark{}\origpage{853}(3月12日)“平展上同调”行动在从1966年到1977年的整整十一年间持续推进,介于SGA~5讨论班的结束与电锯卷册“SGA~4$\frac{1}{2}$”及随后的SGA~5屠杀版(所谓的“Illusie版”)\fnote{440}{(*)}{(3月12日)现在看来,将‘平展上同调’行动视为在1977年便告终止是不准确的。它在1981年反常层会议上达到了封神。}接踵出版之间。它的完成首先归功于我的五位“上同调”学生的共同参与,无论通过作为还是不作为:\textbf{P.~Deligne、L.~Illusie、J.-L.~Verdier、J.P.~Jouanolou、P.~Berthelot}\fnote{441}{(**)}{这种团结一致对这五位前学生中的每一个人来说,首先表现为在1970年我离开之后,在长达七年间对完成他们受托整理的SGA~5报告的彻底拖延。}。除Deligne之外,Illusie的责任在我看来最为沉重,因为正是他承担了屠杀版编辑的责任,从而充当了Deligne的顺从工具\dmark{}\origpage{854}\fnote{442}{(*)}{Illusie同样充当了\textbf{Verdier的同谋},为后者在导出范畴上的欺骗行为提供掩护。}。

Deligne侵吞平展上同调“真正”发明权的意图无可怀疑。整场“平展上同调”行动的精神为此提供了明证,这在我们科学的编年史中无疑是独一无二的。它同样在1975年Deligne的传记小记中被低调表达(其中对我交到他手中的上同调工具——它在其证明Weil猜想最后一环中起到了决定性作用\fnote{443}{(**)}{(3月12日)在这篇文本中,如同在Deligne的任何其他出版物中一样,绝无任何字句暗示他曾从我这里学到任何上同调工具……}——只字未提),并在八年后,在我称之为“葬礼悼词”(三篇)的这组简短而雄辩的文本(1983年)中被公然表达出来\fnote{444}{(***)}{在我关于葬礼的反思中,与葬礼悼词的相遇构成了最决定性的转折点之一。}。它们在“葬礼悼词(1)——或恭维”与“葬礼悼词(2)——或力量与光环”(n\textsuperscript{o}s~104, 105)这两条注中得到了其应得的审视(并在后来的注“阴的葬礼(阳葬阴(4))”,n\textsuperscript{o}~124中以更深刻的光照加以重温)。至于Deligne的自传性(且决非葬礼性的)“赞词”,我在“为一副模糊的骸骨作安魂曲”和“信仰表白——或伪中之真”(n\textsuperscript{o}s~165, 166)两条注中对其作了剖析\fnote{445}{(****)}{关于这一自传小记的详情,亦见前引注n\textsuperscript{o}~164末尾(第V~2部分)。}。

该行动在1977年达到顶点\fnote{446}{(*)}{(3月12日)这是一场纯粹暂时的‘登顶’!见本注‘操演’中今天所写的各条脚注(第853页)。},随着“SGA~4$\frac{1}{2}$(sic)—SGA~5”的相继出版\dmark{}\origpage{855}。这是一场长达十一年的针对我的著作与我个人的下葬中漫长\textbf{节节加码}的(暂时)终点,其每一步骤都因在学术界找到的默许鼓励而变得更加肆无忌惮。

——在前几个阶段里,面对其可疑本性,人们不是漠不关心就是无动于衷(甚或趋之若鹜……)。其中若干阶段我在上文已有提及,即方才回顾过的“Motifs”行动。我还注意到另外三段插曲,它们与“平展上同调”行动有更直接的关联,现在轮到我把它们检视一遍。

\textbf{插曲1.} 它关乎某一条“离散Riemann-Roch”型猜想的遭遇\dmark{}\origpage{856}。我于1966年在SGA~5口头讨论班上引入它,在最后一次报告里,我提炼并评述了若干未决问题与前所未有的猜想。这份报告在Illusie版里尸骨无存——文中对这条猜想绝口不提(并非无缘无故……),对当时提出的众多其他问题亦然。然而,讨论班七年后,这条猜想以解析语境的面目在MacPherson笔下重新现身,绝口不提任何SGA~5讨论班(或概形语境),并顶着“Deligne--Grothendieck猜想”这一不寻常的名字。所涉的是MacPherson那篇广为人知的文章\fnote{447}{(**)}{即R.~MacPherson《奇异代数簇的Chern类》,载Annals of Math.~(2)~100~(1974), pp.~423--432。他在该文中于解析语境证明了这条猜想。},他在其中于解析语境证明了这条猜想。在他去年10月来访时,Deligne向我指明:他在1972年仅限于把我的猜想\textbf{原封不动地}传达给MacPherson(那是他随SGA~5其他听众一起在口头讨论班上学到的)。他说他对MacPherson所起的名字感到惊讶,却并未因此劳神给他写信请他更正。关于这一点,见“把一个个~i~点清”那条注(n\textsuperscript{o}~164,第II~1部分);关于这条猜想本身更详细的说明,见“屠杀”那条注(n\textsuperscript{o}~87)的长子注~n\textsuperscript{o}~87$^1$\fnote{448}{(*)}{这条猜想将首次以其原初与完备的形式,出现在SGA~5的刊行本中(第VIII讲)——十一年后!而在题名之下,我的名字……}。

\textbf{插曲2.} 它关乎SGA~7讨论班的坎坷命运。该班致力于\textbf{平展上同调中的单值群}问题,1967至1969年间,在Deligne与我共同倡议和主持下进行。讨论班的起始想法与总体设计出自我手;Deligne在其中亦有数项贡献,最重要的一项是他在平展语境中对Picard-Lefschetz公式的证明。一如SGA~5,报告的整理工作也拖了数年——那恰是其不幸前任那场(非)整理之剧本的(开头部分的)重演!出版终究还是于1972与1973年在Lecture~Notes~n\textsuperscript{o}s~288, 340中实现了,由Deligne操持——其时我离开数学舞台已三年。经他倡议,讨论班被\textbf{硬切成两部分}:第一部署为我的名义指导,第二部署为他和N.~Katz(后者在讨论班第二年不过是一名普通的报告人)的共同指导\fnote{449}{(**)}{关于我在这场没有任何数学理由支持的\textbf{切断}中所看到的意义,见下文紧随其后的注(n\textsuperscript{o}~56),以及本注“操演”的子注“驱逐(2)”(n\textsuperscript{o}~169$_1$)。}。

在以我名义刊行的第一卷SGA~7~I中,我在讨论班开篇几场报告中系统阐述的消没闭链理论,被“砍”成Deligne一份二十页的摘要(其余报告由我本人及讨论班其他参加者在合理期限内完成整理)。至于以Deligne-Katz联合署名刊行的第二卷——我在其中对主要主题与结果之发展所起的作用,不减于第一卷——这一份贡献,却被系统性地抹杀了\dmark{}\origpage{857}。关于这方面详情,见“一场屠杀的序曲”那条注(我在其中试图厘清SGA~7这场微型行动的意义),尤其是“把一个个~i~点清”(第II~5部分),n\textsuperscript{o}s~56, 164。

这里我只限于提一提最大的一笔抹杀。它关乎我在平展上同调语境下,对“Lefschetz束”上同调理论及“不可约性定理”所作的转写。这批经典结果(凡是真被证明了的……)经由超越途径获得,这一转写(一如常有的情形)绝非机械照搬。我记得为此花了好几天乃至整整一星期。据我所知,至今没有别的已知证明可以处理其中的主要事实——除了我当年借助于谱序列、以及我于1958年确定的代数曲线“驯顺”基本群的“著名”结构所提炼出的那一个\fnote{450}{(*)}{在接下来将被引用的Katz报告的引言里,他似乎还给人以Raynaud夫人是那条定理作者的印象!}。这一理论被复录于SGA~7~II里Katz的一份报告(第XVIII讲)中——依据的是我提供给他的笔记。在卷册引言里,Lefschetz束理论(连同Deligne证明的Picard-Lefschetz公式)被呈为讨论班的“关键结果”之一,但对我在该卷所展开的任何主题中扮演的角色,却无一字提及。据我所知,文献中唯一一处哪怕暗示了我这层角色的,是Lefschetz束报告题名(以及作者名)之后一条简短而暧昧的脚注\fnote{451}{(**)}{这条脚注含混其辞,小心避免确认发明权——而当时的圈内人都心知肚明。}:“据Grothendieck的(简略)笔记”。

在Deligne同年(1973)发表于《数学出版物》的文章“Weil~I”中(169$_4$)\fnote{452}{(***)}{见子注“‘那条’猜想”(n\textsuperscript{o}~169$_4$),它源自此处的一条脚注。},Lefschetz束理论作为其Weil猜想证明中的一件重要技术成分出场。在这篇文章里,Deligne还没有胆量抹杀我在$\ell$-adic迹公式中的角色(那是他证明的另一件关键成分,其在知情圈内的发明权毕竟还太昭然\fnote{453}{(*)}{然而次年,在他那篇自传小记里(已在注n\textsuperscript{o}s~165, 166中检视),对上同调工具的任何提及都消失得无影无踪……});可是,当他郑重其事地表述他即将使用的Lefschetz理论的种种结果时,对我这个人却无一字提及。他仅限于指向SGA~7的相应报告——而一位不幸的读者,多半永远无从在那里觅得他朋友Katz那处难以捉摸的脚注……

\textbf{插曲3.} 我所知的“节节加码”最后一段插曲,发生在1976年——比“SGA~4$\frac{1}{2}$”行动的“高潮”早一年。它关乎Verdier那篇文章的命运\dmark{}\origpage{858}。

在这篇文章里——我曾把它称作“好的文献”——Verdier在解析空间上、对解析可构离散系数,复述了不差分毫地\textbf{我在1963年、尤其1965/66年在SGA~5的概形平展语境下}所发展出来的对偶形式主义。他处处只字不提它的来源——只除了以他特有的一类非常特别的玩笑方式……\fnote{454}{(**)}{将概形(或“空间”……)的\textbf{同调}定义为其在某“对偶化复形”中取值的超上同调——这一想法由我于1963年引入。}——并且对作者或许有所耳闻的任何SGA~5讨论班,也绝口不提。详情见“好的文献们”与“玩笑——或‘权复形’”两条注(永远是同样的“权”,错不了……)n\textsuperscript{o}s~82, 83。

正是从这份“难忘文章”开始,复解析空间上、对解析可构离散系数的对偶形式主义——它逐字逐句复录了我早在1963年、尤其1965/66年在SGA~5的概形平展语境下所发展的那套——便神不知鬼不觉地变成了“Verdier对偶”;人们尚须等待五年之后(在1981年6月吕米尼会议的亢奋中),同一套腾挪才也降临到平展对偶头上。但我在此提前说了(一如我方才对“难忘文章”那段插曲本身已提前的那样)那\textbf{第三}场大行动——这一回,Verdier是其主要的(若非唯一的)“受益者”——它将在下文谈到\fnote{455}{(*)}{见注“分赃”,n\textsuperscript{o}s~170(i)-(iii)。}。

\subsubsection*{b4.~放肆}
\tocsubsub{b4.~放肆}

\notehead{169(iv)}{}

\dmark{}\origpage{859}Verdier的这篇文章,于我如同投来一道意外的光照,照亮了SGA~5在我昔日某些学生手中所遭的命运。它让我看清:这些人从对我为他们最先发展于SGA~5的思想与技术的独占性知悉中,究竟能榨出何种“利益”。它也无可置疑地向我显明:我全体上同调学生与这类行动之间的同谋与团结。我把这篇文章称作“好的文献”时,不曾想到这名字取得如此贴切——它确实已经(如我从多方得到证实的那样)变成了一份标准的参考文献\fnote{456}{(*)}{至于\textbf{上同调}变体(在Verdier的文章中仅被一笔带过——他深知Deligne要把这一份留给自己……),见下一节。},他们当中谁也不可能不曾注意到它\dmark{}\origpage{860}。这正是我最终在注“沉默”与“团结”(n\textsuperscript{o}s~84, 85)中领悟到的东西。我知道我不必惊讶:在那部曾是SGA~5讨论班的Illusie版里,从头到尾对我的同调(及结合于闭链的同调类)形式主义毫无一字提及——确实也无须提,因为(十年之后)他的好伙计Verdier已经替代大家心满意足地提供了那缺失的文献……

Verdier提供的“好文献”,正如Deligne用以对Motif进行局部起棺的“难忘卷册”一样,在我看来纯粹就是彻头彻尾的剽窃。被称为“SGA~4$\frac{1}{2}$”的文本尚且不完全是这般模样\fnote{457}{(**)}{(3月21日)在合题为“公式”的那一系列子注(n\textsuperscript{o}s~169$_5$至169$_8$)里展开的反思向我表明,尽管某些“形式”仍被保留,但这种印象是错误的……}。在那里,某些形式仍被守着,遵循着那套必不可少的“暂停!”风格——它擅长于不断暗示虚假,却从不(或几乎从不……(169$_3$)\fnote{458}{(***)}{关于这一点见子注“特洛伊木马”(n\textsuperscript{o}~169$_3$),它源自此处本来用以阐明这“或几乎……”的一条脚注。})贸然把它明白\dmark{}\origpage{861}说出来。我与“SGA~4$\frac{1}{2}$”以及该风格在其中所取的特定形式(即轻蔑贬低的风格\fnote{459}{(*)}{这种“贬低”假装对一位“糊涂”(“虽说严密”……)且尽情杂乱的前辈所积累的“废话矿渣”一扫而空……})的首次交锋,是在“白板”那条注(n\textsuperscript{o}~67)里。

但这门行动尤其给我的触动,胜过一桩平庸的剽窃所能为的,在于它在放肆上所具有的某种维度。在我眼里,其余三场行动没有哪一场抵达了这一极端的维度\fnote{460}{(**)}{(3月11日)这种评价当然完全是主观的。顺带说,写下这行字时,我还有过一丝迟疑,想着那桩难以想象的反常层会议“行动”(即“行动IV”,下文还将谈到)。这场难忘的会议确确实实构成了借一个鲁莽的继承者(Zoghman~Mebkhout)的中介对我这个人下葬的一场真正的集体封神。正是在这一机缘下我意识到,这场封神同时是“平展上同调”行动的一个自然延续与最终结局,而“SGA~4$\frac{1}{2}$—SGA~5”那一幕,实际上不过是一场纯粹暂时的“登顶”。在后一场里,我的前学生Deligne仍免不了在此处彼处提及我这微不足道的人和我的著作,哪怕心不甘情不愿,并以此用傲慢的考语同其划清界限。而在1981年6月的吕米尼会议上相反,平展上同调虽处于公众注意的中心,我的名字(正如当值无名之辈Zoghman~Mebkhout的名字一样)却在任何时刻都不曾被吐露过……}。而且它对我的触动或许比其余三场中的任何一场都更强烈,因为它进一步像一桩暴力行为一样刺痛了我,像“为了好玩”而对一部我已完成且在其中倾注了整个自我的美好工作的一场屠杀——首要地,为了那些恰恰在事后乐于对它加以洗劫的人们而完成,好拿它充作他们自负的饲料,并且(顶着大雅之堂上名流君子的高尚门面)在其中张扬一种隐秘的狂妄以及那副自得轻蔑的神气\fnote{461}{(***)}{这种自负与轻蔑极其清晰地流露在这一卷名为“SGA~4$\frac{1}{2}$”(在我们这门科学的历史上在该类型中想必是独一无二的)的书本字里行间。它们也在这部卷册发表的当年(但调门更为隐蔽),出现在Pierre~Deligne对我的个人关系里。(见“两个转折”那条注,n\textsuperscript{o}~66.)我在那些曾是我学生的人之中这一个与那一个的怠慢轻浮中找到了它们——对我写给他谈及那些我挂心或使我伤心之事的信件置之不理。我轻盈而轻慢地,在对一项满怀爱意完成的工作的“Illusie版”(或屠杀版)的引言字里行间重新找到了它们,并在去年,在另一位学生的父式优越神气里(在“玩笑——或‘权复形’”那条注,n\textsuperscript{o}~83里谈到)重新找到了它们。}。

\subsubsection*{b5.~藏宝}
\tocsubsub{b5.~藏宝}

\notehead{169(v)}{}

\dmark{}\origpage{862}(2月28日)我方才检视的两场“行动”,连同第四场(所谓“反常层会议行动”)——下文还将谈及——皆有着共同的特征:在众多人的参与或纵容之下进行,而“受益者”(看来)仅此一人。这是这三场行动之间一个触目的共同点,它印证了在“掘墓人——或全体教众”(n\textsuperscript{o}~97)那条注里所继续的反思。

但我在这围绕Motif与平展上同调展开的前两场行动里,看到了一个更阴险的共同点,关乎激励着它们的某种精神。这里涉及的,是一种对待占有高水平、受限流通的科学信息的内心姿态;或者说到底,将信息局限在一个由利益同盟结合起来的少数人小圈子(乃至单一一个人)之内,只要把排他性的独占“收益”留给自己显得有利,便动用阻断流通的权力。

就这样,在我1970年“离开”之后,Deligne成了唯一一个(除我之外)深切同化了“Motif的yoga”并感受其全部深远影响的人——去作人们已知的那种使用。我的五位上同调学生(包括Deligne在内),或许还有两三位其他当年有毅力真正消化其精髓的SGA~5前听众,成了唯一独家支配我在这一讨论班中所发展的想法与技术的人。

在这两种情形里,无论在1965至1969年间与Deligne数不清的面对面长谈中,还是在1965/66年面对SGA~5听众的小圈子,若说确确实实是“首先为了他们、优于所有其他人”而在他们眼前长久地阐发并展开某种内在视景,我把这些于我珍贵的事物交到他们手中,却绝非作为某种“利益集团”的代表。于我,这是不言而喻的:我面对他们,就像面对同我一样的人——除了证明自己与为数学事物的共同认识作出贡献的自然愿望之外,还被一种服务精神所激励,面对着一个在空间与时间中皆无边界的“数学共同体”\fnote{462}{(*)}{关于这样一种“服务精神”,尤见下文同样引用的那条注“阴,作为仆人,与新主人”(n\textsuperscript{o}~135)。}。\dmark{}\origpage{863}而我交到他们手中的东西,我深知那绝非“稀罕物件”、绝非陈列品,而是活生生、灼热的事物,生来便要生长、便要播撒——而这也正是那些我所面对的人自始便感知到的\fnote{463}{(*)}{(4月10日)这并不妨碍他们之中的某些人在事后竭尽所能去贬损他们长久藏匿的东西——在他们起初(除Deligne外)费了九牛二虎之力才把握其意义与影响、并将其同化之后。我在这种贬损的调门里(它叠加在下文所谈的“藏宝”姿态之上),看到了一种双重补偿。一方面是在他们内心深处排遣一种不安(由于挪用了一样并不属于他们、而属于所有人的东西所引起的不安)——在他们自己的眼里假装贬低被挪用之物的价值。另一方面,则是面对被感受为一种超出他们之创造力化身的“父亲”的补偿(而他们无法承担那份同样安息在他们体内、一如安息在他们暗中引以为怨怼的那人身上的力量……)。我的“亡故”状态,以及直接继承人树立的榜样,创造了一种有利于“发泄”隐秘敌对的良机,“父亲”从此被感受为处于虚弱、低下的境地。}。

如果我面对他们,绝非像面对某种我托付了股票的股东(以不知什么共同“利益”的名义),而正是像面对以共同冒险与我相系的人们——因而那些人,会以充当我在他们耳边传递的“信息”的传递站为己任(哪怕随他们的心意在其中加入自己的东西,并在周遭引发回响……),正如我自己曾为了他们的利益而充当那份传递一样\fnote{464}{(**)}{因此,正是对这个“无边界的数学共同体”,我既直接面对它,又通过他们并穿越他们面对它。我在别处解释过(见第847页脚注(*)),为何在讨论班结束后的至少一年里,我不曾亲自负责将它誊清以供所有人支配。}。

相隔近二十年之后,我才意识到:在我们之间有着根本的误解——我们不曾“调谐在同一波长上”。我曾作为活生生的事物托付在自以为充满爱意之手里的东西,被像某种急于掩埋深藏的私藏宝藏一样囤积了起来。对私藏宝藏的占有代表着某种权力(鉴于其代价,这权力自然是可笑的……)——哪怕只是一种扣留的权力,一种阻拦(哪怕只在一时)一样生来便要绽放与播撒的活物绽放与播撒的权力。

我曾试图抓住在这场“误会”\fnote{465}{(***)}{在写下这些行与“误会”这个词时,一个联想不由自主地浮现在我心头:那是Zoghman~Mebkhout的信(引于注“教学的挫败——或创造与自负”,n\textsuperscript{o}~44$'$),其中谈到在我的学生与我之间存在“某种不理解”(然而把Deligne排除在外……)。我当时并不能肯定理解了他想说的是哪种“不理解”。这是否正是我在此谈到的这个“误会”——而他之所以将Deligne排除在外,是出于他蓄意(这位朋友的这种打算曾不止一次令我吃惊)只“从玫瑰色”看他?}中相互碰撞的两种本质迥异的姿态,在“阴,作为仆人,\dmark{}\origpage{864}与新主人”和“阴,作为仆人(2)——或慷慨”(n\textsuperscript{o}s~135, 136)这两条注中。我并不想在此摆出一副“服务姿态”模范化身的姿态,去同“等级姿态”相对立——后者把“知识”变成某种精英的独特标志,并(在风尚败坏的更深阶段)变成对他人施行任意权力的手段。正如在《愚妄与更新》(《收获与播种》第一部分)中的反思所表明的,现实要复杂得多。我曾在我本人身上、以及在我作为数学家的过去的某些行为中,察觉到今天所见到的普遍堕落的萌芽。而同样千真万确的是,我心中这股“服务冲动”曾是一大强大的行动引擎,尤其在推进我的数学文字工作、更具体地在不知疲倦地追求EGA与SGA两套奠基丛书的过程中\fnote{466}{(*)}{这项“不知疲倦的追求”往往直接违背了我心中同样强烈的另一种冲动:抛下所有牵制我的“任务”,奔向前方那不断召唤我(并且至今仍在召唤我……)的未知。}。

看来我未能把这股冲动的丝毫、或反映这股冲动的姿态,传递给我的学生。这项所开创的事业,在其体现为服务于一个共同体的姿态与性情的限度内,在我离开后便戛然而止\fnote{467}{(**)}{(4月10日)在重新誊写这几行时,我被这一局面的奇特讽刺深深震撼,其意义(正如整个下葬的意义一样)在此时此刻尚未被完全领会。正是那位全身心投入到服务于某个“数学共同体”的任务中的人,发现自己竟被从其自己的著作中驱逐了出去,而且得到了上述“共同体”默许且毫无保留的赞同——被恰恰把拒绝服务奉为等级信条与第二本能的人们所驱逐。
然而,一旦忆及我心中这股“服务冲动”所面向的“共同体”,绝非那个毫无保留地参与了我的下葬的社会学实体(带着其头面人物的“等级”等),而是前文所谈及的那个“在空间与时间中皆无边界的数学共同体”,这一表面的悖论在我看来便在很大程度上消解了。(关于这两个“共同体”之间的区分与混淆的评述,见后文注“尊重”(n\textsuperscript{o}~179)的第一条脚注。)}——仿佛突然被电\dmark{}\origpage{865}锯(或链锯……\fnote{468}{(*)}{关于“链锯”效应——它硬生生斩断了我几乎所有前学生的一项正在展开的著作的生机勃勃的势头——见“共同继承人……”与“……及链锯”两条注(n\textsuperscript{o}s~91, 92)。})截断了一样。

从我所离开的世界在此处彼处仍传达给我的回声里,我看到:这种自发的姿态(当年我与在我起步时接纳我的仁善前辈所共有的姿态),已经(正如那份仁善本身一样)在我曾所属的这个世界里成了一个陌生的异乡人。

\subsubsection*{b6.~驱逐}
\tocsubsub{b6.~驱逐}

\notehead{169$_1$}{}

(3月9日)\fnote{469}{(**)}{本子注源自主注“操演”(n\textsuperscript{o}~169)的一条脚注(见第848页脚注(*))。}“SGA”这一缩写,是“布瓦玛丽代数几何讨论班”(Séminaire de Géométrie Algébrique du Bois Marie)的代称。它指称(或至少,在六十年代曾指称)这样一些讨论班:1950至1969年间(自1962年起,与学生及其他合作者一道),我在其中展开了我那新的代数几何基础纲领,与EGA(《代数几何原理》)丛书那些(不那么“前沿”、文体也更规范的)文本并行\fnote{470}{(***)}{系与J.~Dieudonné合作撰写。}。这些讨论班自1962年起在“布瓦玛丽”举行——那是Bures-sur-Yvette的一处地名,IHES自1962年起即设于彼。说实话,头两期讨论班(1950至1962年间)是在巴黎一处因陋就简的场所(蒂耶尔研究所)进行的,听众从未超过十来个人,而我在他们面前严格地“单枪匹马”。“SGA”这个缩写就诞生于那些年头,当时还谈不上什么“布瓦玛丽”。我只是在后来,才给“代数几何讨论班”这个初始的名称添上这漂亮的别号,好让它少一点干巴巴的味道。

不言而喻,这些讨论班的后续,从SGA~1直到SGA~7,是按年代先后编号的。同样不言而喻:每一期讨论班的总体构想都出自于我。它受我那全局性与长远意图的启发:为代数几何奠定广阔的基础,并且——越来越同时地——为一种更广阔的“几何学”奠定基础;这种几何学,我从1963年前后起就强烈地感受到它,而它一直没有名字(今天我会以“算术几何”称之,即代数几何、拓扑\dmark{}与算术的综合\fnote{471}{(*)}{关于这一点,见第844页脚注(*)。})。这些讨论班中的最后一个\origpage{866}是SGA~7,它(与其前的各期不同)连续办了两个年度,1967--69,由我与Deligne共同主持。

那部名称玩着障眼法的卷册“SGA~4$\frac{1}{2}$”(如上文第847与851页所解释)由1973年之后的文本编成,因而也晚于SGA讨论班的最后一期——倘若不算那些从SGA~5里劫掠来的文本,以及Verdier一篇“博士论文”的那份著名的“第0状态”(第三场行动处将谈到它)。撇开日期问题不谈,构成这一卷册的文本之五花八门的性质,与我当年推进SGA丛书的精神毫无相干:SGA丛书的每一卷,都针对我的纲领中一个此前在任何地方都不曾展开过的部分,做出一项大规模的奠基工作——因此,它排绝一切“文摘”式的卷册,排绝对那些早已尽人皆知、打磨停当的结果的汇编,也排绝那些即使是新的、却属零星偶得的结果。退一步说,即便把Deligne那一卷命名为SGA~8(假定我会首肯)也不妥当:这样一个名称,会暗示(一种全无根据的想法)我在SGA~1至SGA~7各期讨论班中所推进的事业仍在延续。至于Deligne选定的“SGA~4$\frac{1}{2}$”这个标号,它就不只是“不妥当”了:它本身就是一场骗局与冒名顶替。在我看来,这一点对自1977年以来有机会读过此卷的众多数学家——他们同时又懂得SGA这个标号的含义,懂得它与我这个人、与我的著作、并由此与某种精神不可分离——本应一望即知。尽管如此,这一冒名顶替,在一个标准参考文献的名称里,已被“数学共同体”容忍了八年,表面上竟“不起一丝波澜”。我在其中——连同作为其自然延伸的1981年反常层会议——看到七、八十年代数学界的奇耻大辱;在我看来,这一耻辱在我们的科学史上没有先例。

作为这场把我从SGA丛书中驱逐出去的行动的先声,有一个插曲意在制造这样的印象:在SGA基础文本的发展中,我这个人只扮演一个偶然的、杂乱而次要的角色。这就是所谓的“SGA~7微型行动”。在注《操演》(n\textsuperscript{o}~169)的“第3幕”(一场节节加码)里谈到过它,而尤其(从我此处关心的视角看)是在\origpage{867}注《一场屠杀的序曲》(n\textsuperscript{o}~56)里。它说的是:把原初讨论班的一部分,以Deligne与Katz之名、将我排除在外,单独发表为SGA~7~II一卷(同时把我在其主体主题与若干关键结果的发展中所起的作用掉包隐去)。我在所引的那条注(n\textsuperscript{o}~56)中就此写道:“这一‘SGA~7’行动,绝非对SGA各期中所推进事业的延续;我把它感受为一种粗暴的‘锯击’(或电锯……),以一卷同我这个人公然划清界限的卷册,终结了SGA的序列——尽管它系于我的著作,并且同其他各卷一样带着它的印记。”SGA~7~I与SGA~7~II这两卷,尚未摆出那副对它们所从出的著作近乎掩饰的轻蔑与鄙夷的神气。倘若这一加码的步伐四年后终究迈了出去,那是因为先前的步子(其中包括这桩看似无奇的SGA~7微型行动)都“过去了”,从未(至少据我所知)在数学界激起半点反应。

我想以一段(想必是暂时的)富有教益的尾声,来结束这场由Deligne操办、得到“全体教众”默许的、把我从SGA中驱逐出去的行动。那是Byrne夫人的回答——她执掌Springer出版社的“Lecture Notes”丛书——口气十分“冷淡”。我曾写信给她,要求澄清一卷名为SGA~5的书:它于1977年以我的名义在“Lecture Notes”丛书出版,由Illusie编辑,而Springer出版社非但没有征求我的同意,甚至没有通知我它经手了这次出版。我从她的回信(一个月后才收到)中获悉:纠缠这类手续纯属多余,因为我自认为是该SGA~5卷的作者,这实属谬误——我不过是这期讨论班的“主持人”,仅此而已,它由Illusie编辑!(于是人们不免回过头来纳闷:那位已故的主持人,当年在那个讨论班里究竟干了些什么……)我抱着看看的念头,亲自写信给K.F.~Springer先生,述说我自1972年以来(那一年,SGA~7~I以同样的方式冠以我的名字出版——说真的,我对它的“作者”身份,并不比我对SGA~5的更多……)与Springer出版社打交道所经历的种种怪事。我至今仍在等他的回音……\fnote{472}{(*)}{(4月9日)关于这个故事的后续,见注“殡仪馆——为科学服务”(n\textsuperscript{o}~175)。}

\origpage{868}(3月16日)本子注得到了它当之无愧的名称:“驱逐(2)”。这个(2)提醒人们:先前已有一条同名注“驱逐”(n\textsuperscript{o}~63),我最近还有机会援引过它(连同“Motifs”那一行动)。那条注里(极其含蓄地……)提到的“驱逐”,是1970年我离开IHES那一幕中的驱逐;而那次离开,显然正合我那位年轻而杰出的朋友的心意——他彼时刚刚在取得的位子上安顿下来\fnote{473}{(*)}{关于我1970年离开IHES的插曲,见“未竟的收成”一节(n\textsuperscript{o}~28),以及注“有益的挣脱”“驱逐”“兄弟与配偶——或双重签名”(n\textsuperscript{o}s~42、63、134),最后见后引注的子注(n\textsuperscript{o}~134$_1$)。}。这两次“驱逐”——一次出于IHES,一次出于SGA丛书——之间的血脉渊源,在我看来一目了然。我在其中看到一种触目的递进,性质上仍是一次“节节加码”:第一次,只是把我这个人从一个机构中驱逐出去——我对那个机构固然依恋极深(我真的曾打算在那里终老,千真万确!),但我很快就脱身而出,不留一丝遗憾;第二次,则是把我从SGA中驱逐出去——而SGA(象征性地,甚至远不止象征性地)代表着我的数学家著作,一部直到今天我仍系念的著作。诚然,我被“驱逐”出IHES,十五年来早已尘埃落定;而我却始终怀疑,对一部我曾倾注十五年充实岁月的著作,事情是否同样如此。

我想到,当年我曾亲手为把我从SGA中驱逐出去提供过方便:我顺着自己自发的冲动,把那些曾在某些时刻全身心投入某期讨论班之发展的学生与合作者,呈现为与我同为该讨论班的“主持者”。这在我那个年代并非常例,在今天更绝无可能。我不知这样做对不对。一方面,这并不完全符合实情:我在其中所起的作用,同我的合作者们(纵然才华横溢、纵然如我一般全力以赴)所起的作用,绝无对称可言。因此,这种呈现方式,恰恰合乎我在注“上升”尤其是注“暧昧”(n\textsuperscript{o}s~63$'$、63$''$)中所考察的那种“暧昧”(或者说,对杰出的青年数学家的“纵容”)的方向。倘若由我设立的这种暧昧,竟鼓励了那些曾热烈参与SGA的人们中的一些人,在某个时刻来“驱逐”我(以或多或少局部的、或或多或少彻底的方式),那我怨他们便是我的不是!我不过是收获我自己播下的种子。

\origpage{869}但这丝毫不妨碍我在另一方面,对已然发生的事情作一份公开的陈述。另一方面也确实如此:我同某些合作者建立的那种关系,在他们看来可以是一种器重与信任的标记(它也确实是),并由此鼓励他们同我一样全身心投入这项任务。但我如今告诉自己,这样的器重与信任,完全可以同样清晰、同样鼓舞人心地表达出来,而不必因此染上暧昧。这就好比我在“购买”一份与任务相称的投入,所付出的代价是一份“好处”——而这份“好处”(回过头看)在我看来形迹可疑。因为,装出自己本不是的样子,是一种虚假的好处。而十分显然:在制造这样一种(即便不是彻头彻尾、至少也)边缘走样的假象这件事上,首当其冲负起责任的,是那个以长者身份出现的我,而不是任何别人。

渐渐地,这场反思在一道我落笔时全然不曾料到的“个案情形”的意外光照下,越来越趋近注“暧昧”的反思。我意识到:正如我同那位“年轻天才”(绝非无名之辈)Deligne的关系是虚假的——出于一种虚假的谦逊,我不肯承担那在我与他之间确实属于我的长者与“老师”的角色——同样,我同另一些才华横溢、在那项我当时视为“共同”的任务中全力以赴的年轻人之间的关系\fnote{474}{(*)}{直到注“绿灯”(n\textsuperscript{o}~68)的末尾,我才开始认识到:这绝非误会,而是一场彻底的幻灭。},也同样是虚假的。《下葬》中展开的反思,已经足够清楚地表明:倘若有所谓“共同任务”,那也只是一年顶两年的光阴——是让那个年轻人(姑且这么说)做出一篇博士论文的时间(这已经不算小了)。恰恰在我离开的1970年,响起了信号:这一整片明明灼手、前一天还实实在在“在我手中燃烧”的庞大“任务”,当即、几乎彻底地被弃置一旁\fnote{475}{(**)}{这一在离去的翌日便对一项纲领与灼手任务的即刻抛弃,在注“本能与时尚——或强者的法则”(n\textsuperscript{o}~48)中、尤其在与“共同继承人……”“……与电锯”两条注(n\textsuperscript{o}s~91、92)中曾被谈及;在那两条注里,我(依据传回我耳中的回声)试图逐一点检那些曾由我“1970年以前”的各色学生接手的主题,后来的下落。}。撇开Deligne关于Weil猜想的工作不谈,这同时也是每一个曾最深吸引我的重大主题长期停滞的开端\dmark{}\origpage{870}——这场停滞(除由不予点名的Mebkhout的哲学所触发的那场部分“重启”之外)至今仍在继续\fnote{476}{(*)}{(3月17日)这种“停滞”的印象,将在后面一条注里获得或许更具体的意义:我打算对那些曾列在我日程上、并在我离开之际被我的学生们以完美的齐一性弃置不顾的最“灼手”的主题,作一份简短的带评点的清单。(4月9日)关于这一点,见注“工地的巡礼——或工具与视野”(n\textsuperscript{o}~178)。}。

\subsubsection*{b7.~好撒玛利亚人}
\tocsubsub{b7.~好撒玛利亚人}

\notehead{169$_2$}{}

(3月13日)\fnote{477}{(**)}{本子注源自主注“操演”(n\textsuperscript{o}~169)的一条脚注(见第849页脚注(**))。至于把那套“暂停!”手法——专司把鱼鳔当灯笼卖给一位急不可耐、只求相信的“用户”——拆解得更周详些,见子注“特洛伊木马”与“公式”(n\textsuperscript{o}s~169$_3$与169$_5$--169$_8$)。}在SGA的这一篇序言里,Illusie热忱地感谢Deligne,其中尤其感谢后者“说服了他撰写……Lefschetz--Verdier公式的证明,从而为此讨论班的出版扫除了一道障碍”(着重号是我加的)。说白了:那障碍,便是Illusie十一年来不肯动笔撰写他承诺要撰写的文本的这一“缺乏信念”——而这缺乏,如上文所说,恰恰在好撒玛利亚人Deligne向好撒玛利亚人Illusie亮出“绿灯”、说他“可以放手去做”的那一刻,戛然而止……

这正是“假中之真”。至于这段文字显然力图暗示、却不必明说的那个“假”(循着那套久经考验、早已自成一派的手法……),则是:SGA~5讨论班须依赖那条公式(该公式在讨论班当时,仅凭奇点消解的假设方能成立;这一假设,后来已由Deligne在那部名为“SGA~4$\frac{1}{2}$”的“先行”卷册中提出的有限性结果,在最常见的情形下予以扫除)\fnote{478}{(***)}{关于这一点,见第841页脚注(***)与第850页脚注(*)。}。实际上,正如这两位朋友同我一样清楚无疑:Lefschetz--Verdier公式在SGA~5中所起的作用(一如它在我那$\ell$-进L函数上同调公式的证明中所起的作用)纯然是启发性的:它提供了\origpage{871}动机,促使人们去寻找并证明“显式”的不动点公式(即其中“局部项”可以被显式算出的那种)。于是,Illusie便同他的朋友一唱一和,制造出这样的印象:SGA~5确实(且是在一种他同他的朋友都不肯明确言明的意义上)从属于那部由此只能以“SGA~4$\frac{1}{2}$”为名的文本。

欲知详情,见注《屠杀》及其子注n\textsuperscript{o}~87$^2$。在这条注及其全部子注中,我终于得以弄清(迟悟总比不悟好):Illusie撰写的整篇序言,以及一般地说,Illusie版(或屠杀版)的整体呈现,乃是一份背信弃义的范本——它以一副漫不经心的姿态、伴着那份为其人格平添魅力的天真神气侍奉出来。

Illusie竭力营造的这份动人印象——倒霉的SGA~5讨论班之所以终得出版(晚了十一年,而且是付梓于我所知道的那个状态),全亏了好撒玛利亚人Deligne(自然还有第二位好撒玛利亚人Illusie)——看来“安然通过”,未遇任何阻碍。我在Serre1977年为菲尔兹奖国际评选委员会撰写的关于Deligne工作的报告中,重新找到了这一版本。我对Serre的十足诚意毫不怀疑;他本来也只是远远地追随那场口头讨论班的曲折周折——更何况此间早已逝水年华……他想必同所有人一样、不加质疑地,把Illusie序言里说出或暗示的东西当了真——那序言他总该有一天翻过,为的是看看(而他什么也没看出来!)……

有趣的是,据我所知,同一份Serre报告,也是全部文献中唯一一处(而且就在报告的第一句)提到Deligne曾是我学生的地方。而Deligne的任何一篇出版物,则决不会让哪怕一个读者猜想:这位作者或许从我这里学到过什么。

\subsubsection*{b8.~特洛伊木马}
\tocsubsub{b8.~特洛伊木马}

\notehead{169$_3$}{}

(3月10日)\fnote{479}{(*)}{本子注(属注《操演》)源自该注的一条脚注;见第860页脚注(***)。}在注“白板”的子注(n\textsuperscript{o}~67$_1$)里,我指出过两个例子:Deligne在其中越出了他素来的谨慎,当真“直言不讳地道出了假话”。为那位好奇、消息灵通、手头却没有那条注与其子注的读者着想,我在此提示:撇开他对SGA~4与SGA~5的那些“殷勤话”,\dmark{}以及对我这卑微\origpage{872}的人在各处或明或暗的大“遗忘”(这些在注《操演》及其脚注里已有这里那里的指涉),我所识破的那些明目张胆的骗局,集中在第2页的第3、4两段里(在“SGA~4、SGA~4$\frac{1}{2}$、SGA~5的阿里阿德涅之线”一节里——请欣赏这列壮美的队伍吧……)。这十七行文字,堪称“浑水摸鱼”之术的范本,大值得逐细剖析\fnote{480}{(*)}{关于所引两段中第二段的详细评注,见子注“双重含义——或行骗的艺术”(n\textsuperscript{o}~169$_7$)。}。

我在此只须指出:在所引第一段里,人们读到,为了“在étale上同调中建立一个与凝聚对偶相类似的对偶形式体系”,“Grothendieck曾使用奇点消解与纯性猜想”\fnote{481}{(**)}{原文在“纯性猜想”之后接着写道:“在SGA~4第XVI讲中于一个相对框架下[??]确立,并且——模去消解——在等特征情形下于SGA~4第XIX讲中确立。”这句“在一个相对框架下”(对任何不预先知情的读者都无从索解),正是掩盖下述事实的一种手法:该定理对一切特征的光滑代数簇早已成立。\par(3月17日)我直到此刻才注意到所引那段话结尾的妙处,它在头几遍阅读时竟“蒙混过关”:“若干发展见于SGA~5~I。在SGA~5~III中,人们展示这一形式体系[??]如何蕴含那条极为一般的Lefschetz--Verdier迹公式。”(着重号是我加的。)人们当会欣赏这“若干发展”之不着一字,而作者(他在别的场合是懂得精确的)随即接上“这一形式体系”(=若干发展?),说它“蕴含那条极为一般的迹公式”;好让紧接着的下一段一开篇便凸显:那条公式“在SGA~5的原始版本中”,“仅以猜想的方式得到确立”。我刚刚在SGA~5里核实了这些“若干发展”究竟是什么——第I讲的标题便告诉了我:“对偶化复形”,因而也包括双对偶定理。为什么说“若干发展”,而不说“对偶化复形的理论”或“双对偶定理”?其实并不更长,而且总不那么浑浊!这叫我想起:在那份著名的“有限性”报告里——亦即在“特洛伊木马”里——那位杰出的作者恰恰证明了一条“双对偶定理”,对我这卑微的人一字不提;而这条定理旋即(在由Illusie撰写的SGA~5第I讲引言里)被干脆命名为“Deligne定理”。的确,环环相扣……附记:关于这条双对偶定理(它竟被如此故作漫不经心地对待)的评注,见第852页那条长长的脚注(*)。}。紧跟着便补充说:在本卷中(多亏上苍与那位杰出的作者),这些“要点已由另一方法确立”(着重号是我加的),而该方法有效于“定义在0维或1维正则概形上的有限型概形”——也就是说,实际上有效于使用者会遇到的几乎所有情形。

如此,Deligne便竭力营造(甚至明白断言)这样一种印象:我所发展的étale对偶的整\origpage{873}个形式体系仍属猜想(至少在非零特征下如此),而这些“要点”终究只有靠他Deligne本人、靠本卷(也就是说,靠他的有限性结果——前面几条脚注已提及、他也随即指向的那些结果)方得确立。这可真是,瞧瞧!,正足以替那著名的“逻辑依赖”的虚构张目——SGA~5依赖于那部名为“SGA~4$\frac{1}{2}$”的文本(这层依赖由这名称本身、由“SGA~4—SGA~4$\frac{1}{2}$—SGA~5”这列壮美的队伍所设立),并由此替其引言中那句令人难以置信的断语(已在别处引用并评注过)开脱:“它的存在[‘SGA~4$\frac{1}{2}$’的存在]将使SGA~5不日得以按原貌出版。”

这便是Deligne的版本——它被悄悄塞进那部名为“SGA~4$\frac{1}{2}$”的电锯文本,并且毫无含糊地滑入我方才引用的“阿里阿德涅之线”那段文字之中。而实情是:早在1963年3月,我就已在étale框架中确立了六种运算的完整形式体系(远远越出通常的Poincaré对偶),唯一的限制性假设(显然不可少)是:须以“素于”所论概形之剩余特征的挠系数为系数\fnote{482}{(*)}{由此,“六个函子”以及关于它们的基本公式——其中最关键的是有限型态射的“对偶公式”(人们可把它看作迄今为止可以想见的、经典Poincaré对偶定理的最一般形态)——皆由我确立,自始至终不曾对系数强加任何有限性假设。况且,这一点Deligne比谁都清楚:恰恰是他,以我1963年的笔记为底本,把SGA~4中展开这一对偶形式体系(围绕上述对偶公式)的那份报告写成了详细的文本!}。唯独étale上同调的双对偶定理,我的证明才用到Deligne所指的那些假设。而这后一条定理,在我发现它之前,在上同调(各种“空间”的上同调)里属于前所未知的类型;况且,它在SGA~5讨论班中也只扮演过一个插曲性的角色,那是为了证明Lefschetz--Verdier公式\fnote{483}{(*)}{(3月17日)尽管如此,在所引第二段里,Deligne随即接着强调:该公式“仅以猜想的方式得到确立”,且“其中的局部项并未算出”(这一断语在数学上毫无意义,却有助于为SGA~5营造一片“废话矿渣”的印象,好让大家慈善地把它忘掉……)。我自己承认:一年前头几遍读这些段落时,我如遭雷击——这些奇怪地“不得要领”的评说,针对的又是一部人们本来就劝人遗忘的文本,我当时全然摸不着头脑。正是凭着事后的回顾,凭借一次细心的“按件”复核,那份掉包的意图才终于显形:它由一套精密而万全的“浑水摸鱼”式隐匿手法侍奉着,藏在我最初只当是表皮式恶意、沿着一支讨好的笔端信马由缰地挥洒的东西背后。欲知这套手法更周详的考察,见注《操演》的诸子注“公式”(n\textsuperscript{o}s~169$_5$--169$_9$)。},而该公式本身也只起纯启\origpage{874}发作用\fnote{484}{(**)}{正如我在下文子注“正经数学……”(n\textsuperscript{o}~169$_5$)中所强调的,这条公式曾在心理上十分重要:它为发展“显式”的不动点公式提供了动机。}。在Deligne那部伪作之中,上述双对偶定理的作用干脆等于零(若不算上它在其中于一些“帮衬性”的假设下得到证明,并且——在Illusie那副谄媚的笔端下、在其朋友的怂恿之中——摇身变成了“Deligne定理”……)。

此处无意贬低Deligne那些有限性结果的价值:它们确实填补了SGA~5的一项空白(如同其他许多空白之一),而这是事物之常。任何一门朝气蓬勃的数学理论都不曾“完成”!但人们不能不看到:Deligne利用了这份贡献——它再有用,也仍是有限度的(他还做过更深刻、更困难的事,而且并不费力气……)——把它吹胀到失度,使之成为一场规模空前的诈骗行动的“特洛伊木马”:这就是“étale上同调”行动。

这同一匹“特洛伊木马”,还重新出现在Deligne为名为“SGA~4$\frac{1}{2}$”的那部卷册写给Zentralblatt的“书评”里(前引之文;见第851页脚注(**))。在该书评的最后一段里,我读到:“人们证明,对于定义在一个1维正则概形S上的有限型概形,通常的那些上同调运算[更不必说那‘六种运算’,那是万万点不得名的!]把任何可构层变成可构层。”(着重号是我加的。)这话的表述方式,足以让人以为:在这位作者所呈献的这部杰作之前,人们手头\origpage{875}并没有任何一条针对étale上同调中那些著名的“通常运算”的有限性定理\fnote{485}{(*)}{这正合乎同一篇书评前文所说的“SGA~5的混乱状态”——而本卷本当“补其偏弊”。}。可我却有幸证明了第一条这样的有限性定理,而且是一切之中最关键的一条:关于函子$Rf_!$(带真支撑的上同调)的定理——而且这恰恰是在我发现这样一个函子于étale上同调中的定义之后的那几天里(如果我的记忆无误;当$f$为真时,它与“平凡的”$Rf_*$重合)。那是1963年2月,在我有幸结识我这位未来的学生之前,而在那时,除了我(勉强再算上Artin),还没有人十分有把握:étale上同调这东西究竟“存在”与否。它正是在那些日子里,开始真正存在的。

还剩下$Rf_*$的类似问题,它证明更为顽抗,而且事实上直到今天,仍未以它(想必)应有的全部一般性获得解决。其实,就在同一年(甚至同一个月),我已完成那些必要的“拆卸”(这在今天,随便谁来都能一挥而就……),表明从$Rf_*$的有限性出发,可以证明$Lf^!$与$R\mathcal{H}om$的有限性\fnote{486}{(**)}{至于六种运算中余下的两个,即$Lf^*$与$L\otimes$,它们把可构系数变成可构系数,这是平凡的。}。诚然,这从此成了étale上同调的“基本常识”,并且想必属于我那杰出的“前辈”“SGA~4$\frac{1}{2}$”注定要使人“遗忘”的那些“技术性离题”之列……

\subsubsection*{b9.~“那条”猜想}
\tocsubsub{b9.~“那条”猜想}

\notehead{169$_4$}{}


(3月12日)\fnote{487}{(***)}{本子注源自注“操演”(n\textsuperscript{o}~169)的一条脚注;见第857页脚注(***)。}自Deligne发表其著名文章“La Conjecture de Weil~I”(在其中他攻克了Weil猜想中我曾悬而未决的最后一环)以来,我不止一次注意到一件奇特的事,但直到最近几天才真正去探究:Deligne总是谈论\textbf{那条}Weil猜想,而此前人们习惯的说法始终是Weil\textbf{诸}猜想。

确实,在Weil1949年的经典论文《有限域上方程的解数》(载Bull. Amer. Math. Soc. 55 (1949), pp.~497--\origpage{876}508)中,这些猜想正是以一连串一个比一个更令人叹服的断言的形式呈现的;而我在五十年代中期从Serre口中习得它们时,也是作为这样一整套猜想群。诚然,在这个乍看五花八门的猜想集合里,有一种显然的灵感之\textbf{统一}:它首先来自与上同调形式主义相关联的直觉(经由Lefschetz公式),同时也(我至少如此揣想)来自Hodge理论。

通过创造并发展这样一套\textbf{上同调工具}——适用于基域任意的代数簇——我得以证明这些猜想中的大部分。我在Artin、Verdier及其他人的协助下完成了这项工作,把生命中紧凑的三年献给了一场实地精雕细琢的工夫——它凝结为两千页“不可卒读”的“废话矿渣”与“技术性离题”,正是这些,使得一位Deligne能在区区二十页密排文字里“砍下”最后一步……此外,受Serre发现的一个出色的Weil猜想“Kähler类似物”启发,我提炼出了(以我所谓关于代数闭链的“\textbf{标准猜想}”的形式)把Hodge理论\textbf{移植}到任意基域上去的原则——或更确切地说,把Hodge理论中从“代数”观点对复代数簇上代数闭链理论真正要紧的那部分加以移植。只要把这些猜想按其初始形式(或许过于乐观)稍加(且显然地)重新表述,它们至少在特征零中成立,并且在特征~$p>0$~中也“一定为真”(只要Weil诸猜想为真……)。

这决非巧合:同一位执意把Weil猜想“变为单数”的Deligne,也正是那个竭力抹杀其证明中他的老师所扮演角色的人;也正是他,还努力(鉴于普遍的冷漠,他成功了)把“标准猜想”败坏成一条死胡同、一条无法企及的路——说穿了,是一个\textbf{障碍}——好在通往Weil猜想的证明之路上,仰赖上帝(和他那微不足道的本人)把它绕了过去\fnote{488}{(*)}{(3月16日)关于这一双重“掉包-拆解”的若干细节,见《葬礼悼词》(注n\textsuperscript{o}~104、105),以及注171(x)开头关于这篇悼词的寥寥数语。关于掉包之术的更周详考察,见“公式”诸子注(n\textsuperscript{o}~169$_5$--169$_9$)之全体。(x)(5月11日)旧注“封神”的这一开头部分已与该注分离,成为一篇独立的注“珍宝”(n\textsuperscript{o}~170(iii))。}……

\subsubsection*{b10.~公式}
\tocsubsub{b10.~公式}

\subsubsection*{(a) 正经数学……}
\tocsubsub{(a) 正经数学……}

\notehead{169$_5$}{}

\dmark{}(3月17日)\origpage{877}对于定义在有限域~$k$~上的代数簇~$X$,出名的“Weil诸猜想”关乎与~$X$~相关联的“L函数”(所谓“Artin-Weil函数”)。它被定义为某一有理系数的形式级数,对该级数的了解等值于对~$X$~在~$k$~上及其全部有限扩张上的有理点个数的了解。诸猜想中的第一条断言是:该形式级数(常数项为1)是一个\textbf{有理函数}在$\mathbb{Q}$~上的级数展开。所有其余断言都关乎该有理函数的特殊形式与性质——在~$X$~为连通射影且非奇异的特殊情形下。这些猜想的核心是一条假定典范的公式,它把该有理函数呈为
\begin{quote}
$$L(t)=\frac{P_0(t)\,P_2(t)\cdots P_{2n}(t)}{P_1(t)\cdots P_{2n-1}(t)}$$
\end{quote}
其中诸~$P_i$($0\le i\le 2n$,$n=\dim X$)是常数项为1的整系数多项式。~$P_i$~的次数~$b_i$~应当充当~$X$~(更精确地说,对于域~$k$~的代数闭包~$\bar{k}$~上与之对应的簇~$\bar{X}$)的第~$i$~个“Betti数”。于是,当~$X$~由定义在特征零域~$K$~上的非奇异射影簇~$X_K$~经“特征~$p>0$~的约化”而来时,$b_i$~应等于该\textbf{复}代数簇的第~$i$~个Betti数(以超越途径定义)——后者经由~$K$~到~$\mathbb{C}$~的任一嵌入从~$X_K$~得来\fnote{489}{(*)}{Weil作这些猜想时,人们甚至还不知道如此定义的$b_i$与$K$到$\mathbb{C}$的选取无关。几年之后,这一点将由Serre的凝聚层上同调理论得出——该理论给了Hodge理论中更精细的不变量$h^{i,j}$一个“纯代数的”意义。}。该有理函数须满足一份\textbf{函数方程},等价于说:$P_{2n-i}$~的根恰是~$q^{n}/\xi_\alpha$,其中~$q=p^f$~是基域~$k$~的基数,$\xi_\alpha$~遍历~$P_i$~的根。(从义理上说,这该当“源于”簇~$\bar{X}$~那未具名、未定义的“上同调”之“Poincaré对偶”的存在。)

我相信Weil想必也猜到了:对~$i\le n$,$P_{2n-i}$~的零点恰是~$q^{n-i}\xi_\alpha$,其中~$\xi_\alpha$~遍历~$P_i$~的零点(或者——鉴于对偶条件——等价地说,$P_i$~的零点成对配组,每对的乘积等于~$q^i$)。这里的启发式“理由”是复非奇异射影簇\dmark{}上同调的另一条重要性质,\origpage{878}这一次由“Lefschetz定理”(所谓“蛮牛”版,即“困难”版)来表达。最后,Weil猜想中最末一条——Riemann猜想的“几何”类似物——说:$P_i$~的零点的倒数的绝对值都等于~$q^{-i/2}$(这条断言导致对~$X$~点数的高精度估计\fnote{490}{(*)}{由Weil猜想中的这最后一条,同时得出:$L$函数的写法$(L)$是\textbf{唯一的}。})。

对一般簇~$X$~的L函数的有理性,曾由Dwork于1950年以非上同调的“$p$-adic”方法确立。该方法的不便在于:L函数没有上同调解释,从而对其他猜想在非奇异射影~$X$~情形的攻克无从下手。在后者情形,拥有一个上同调形式主义(系数域~$R$~特征零)、含Poincaré对偶、含结合于闭链的上同调类形式主义(把相交变为杯积),就能本质上以“形式”的方式转录经典的“Lefschetz不动点公式”。把这一公式应用于~$\bar{X}$~的Frobenius自同态及其迭代,便将得到Weil所要求的形如(L)的表达式,其中诸~$P_i$~是~$R$~系数的多项式。这在Weil1949年陈述这些猜想时就该是清楚的,至少对五十年代的Serre和我而言是如此——这正是发展这套形式主义的初始动机。早在1963年3月,取~$R=\mathbb{Q}_\ell$($\ell\ne p$),这事就已做成。只剩两粒盐:

a) 先验地不清晰(尽管人们相信理当如此):多项式~$P_i(t)$——先验以$\ell$-adic整数环~$\mathbb{Z}_\ell$~为系数——实际上是\textbf{普通整数}系数的,而且不依赖于所取素数~$\ell$($\ell\ne p=\operatorname{car}k$).

b) 从非奇异射影~$X$~的L函数有理性,只有拥有奇点消解,才能推导出一般~$X$~的情形。

a) 提出的问题对Motif的yoga之绽放与发展、对与该yoga紧密相连的\textbf{标准猜想}的后来表述,扮演了关键角色。\origpage{879}它们也激励人们寻找一套\dmark{}\textbf{$p$-adic上同调理论}(后来由“结晶”理论实现),作为证明诸~$P_i$~系数之整性的可能进路——一旦(比如经由标准猜想的肯定解决)人们知道它们是有理的且不依赖于~$\ell$(\textbf{包括}$\ell=p$~的情形)。

不管怎样,到1963年人们已经拥有了L函数的表达式(L)(尽管它先验地依赖于~$\ell$~的选取)、函数方程、以及Betti数在特化下的良好行为。剩下的就是解决问题a)、证明关于~$P_i$~零点绝对值的断言、以及最后(为了凑足分量)关于~$P_i$~零点的“Lefschetz式”关系。这些正是十年后在Deligne的文章“La conjecture de Weil~I”(Pub.~Math.~de l'IHES~n\textsuperscript{o}~43~(1973), pp.~273--308)中所做的。

作为Deligne证明的成分,根本不需要比“普通”不动点公式更精致的东西——后者早在1963年初就已可用(不带任何“猜想性”成分)。如果我没记错,Deligne文章里仅有的另一项上同调成分,是我约1967或1968年发展的Lefschetz束的上同调理论(平展版),加上Picard-Lefschetz公式(由Deligne在平展框架下证明);两者都发表于SGA~7~II卷——对这部\textbf{前面已经谈到过}的卷册,我的名字不出所料地几乎消失了……

被称为“Lefschetz--Verdier”的“更精致”不动点公式,倒是扮演过一个重要的\textbf{心理}角色:它鼓励我去提炼L函数的上同调解释(L)——对任意簇(不必然射影非奇异)成立。Verdier的这条公式提醒我:不动点公式应当无须对~$X$~附加非奇异性条件(一如在普通Lefschetz公式情形中早已众所周知);而尤其,它把我的注意力引向\origpage{880}这样的事实:存在关于以\textbf{任意}(“可构”)\textbf{层}为系数的上同调的不动点公式——它把(以这种层为系数的上同调空间中的)交错和迹,解释为对应于自同态~$f:X\to X$~之不动点的“局部项”之和(当后者孤立时)。在这份启发式动机中,这条Lefschetz--Verdier公式在特征~$p>0$~下“仍属猜想”(因为还没有奇点消解,\dmark{}从而没有“双对偶定理”)这一事实,\textbf{完全无关紧要}\fnote{491}{(*)}{(3月20日)这一点是如此显著,以至于去年,我早已全然并且长久地忘记了这一事实,而在读到(出自Deligne的笔下)Lefschetz--Verdier公式“仅在SGA~5的原始版本中以猜想方式确立”时,惊得瞠目结舌。我在随后两天(3月18与19日)的反思中回到这一点。(在子注n\textsuperscript{o}~169$_6$与169$_7$中。)}。

一如常有的情形,关键的一步是找到“\textbf{正确的表述}”(在此处,是L函数的“上同调公式”的正确表述)。Verdier的公式启发我引入一个任意的(可构)$\ell$-adic层,取代此前一直隐而不显的常值层$\mathbb{Q}_\ell$。于是,仿照Weil对“普通”L函数的定义,须定义一个“以层~$\mathcal{F}$~为系数的”L函数。一旦想到这一点,定义便自动涌现:那就是我1964年12月Bourbaki报告(《Lefschetz公式与L函数的有理性》,Sém.~Bourbaki~279)中所给出的——在此无须重复。此外,Lefschetz--Verdier公式的合理“局部项”(用给定的系数层与Frobenius对应来表述)也同样自动涌现。最后(人要么大胆要么不大胆!),为什么索性在这里书写公式时,连“正统”Lefschetz--Verdier公式的\textbf{真性}假设也一并放弃,改用\textbf{真支撑上同调}?!

因此,关键的一步,这一回又是提炼出“\textbf{正确的}陈述”(此处即“\textbf{正确的}公式”),\textbf{足够一般}——从而足够\textbf{灵活},能经得起递归与“拆卸”,畅通无阻地通过。我(以及至今任何人都不)无法\textbf{直接}证明“普通”L函数的那条公式——对任意~$X$(甚至光滑但不真,或反之)——以常值$\ell$-adic层的真支撑上同调为项,而不经由层论推广。(一如在特征~$p>0$~下,我无法不经由把Riemann-Roch-Hirzebruch公式\textbf{首先}推广为对光滑代数簇\textbf{真态射}的层论公式,就\textbf{直接}证明它的\textbf{普通}版本——而且据我所知,至今也无人能够……)

\dmark{}\origpage{881}在那份Bourbaki报告里,我仅限于给出以普通$\ell$-adic层为“系数”的L函数公式的一般表述,并展示如何通过非常简单的拆卸,把问题归结为~$X$~为光滑射影曲线的情形。我心里清楚:一旦到达那一步,\textbf{就赢了}——因为对维数一的把握已经足够充分,使该公式的证明成为例行公事\fnote{492}{(*)}{我在此谈“例行工作”时,丝毫没有贬义。数学工作的十分之九,如果说不是更多的话,属于这一类——在我身上如此,在任何一位数学家的身上也如此,只要他有时机经历一些时刻,而那些时刻,恰恰是别的东西,是创造的时刻。在Verdier之后,我自己也曾花时间摇动那些可用技术的曲柄——精细而上了油的——去找到并证明一条令我满意(至少暂时如此)的一维不动点公式。那也像Verdier所做过的,是“例行工作”。}。我当时并未着手去提炼一维情形的好不动点公式并证明它;我觉得那更该轮到Verdier来干。他在次年给出了所谓“伍兹霍尔”不动点公式,足以覆盖Frobenius态射及L函数的应用。我读了他的表述,并不真正满意:因为他对其上同调对应所施加的条件(出于其证明的需要——该证明我未曾过目)在我看来有些人工——我倒希望有一条公式适用于代数曲线的\textbf{任何}自同态。SGA~5讨论班成了发展一条合我口味的这种公式的第一个好时机。(若无疏漏,那正是Illusie版第XII讲中奇迹般幸存下来的那条公式——尽管那场不幸的讨论班风波不断……)Weil诸猜想曾是初始动机与宝贵线索,引我“投身”于平展上同调(及其他)完整形式主义的发展。但我分明感到:这个上同调主题——它已在我工作的中心待了八九个年头,并将持续到1970年我离开——其意蕴比把我引到它跟前来的Weil诸猜想还要辽远。对我而言,Frobenius自同态不是上同调形式主义的“阿尔法与欧米伽”,而不过是众多自同态中的一员……

\origpage{882}依我看,Deligne“SGA~4$\frac{1}{2}$—SGA~5”行动的初始动机\dmark{}曾是\textbf{唯独}侵吞迹公式这一件,从而(作为“推论”)侵吞L函数。是路走到一半,这一图谋才扩大为对“平展上同调”整体的侵吞。我并且相信:这两个“份额”都太大;而且时至今日,尽管有“SGA~4$\frac{1}{2}$”、反常层会议以及诸如此类,“人们”(甚至那些消息并不怎么灵通的人)也都“知道”:创建$\ell$-adic上同调工具的并不是他,而且他也没有单凭一己之力证明“\textbf{那条}”Weil猜想。

尽管如此,为了给“平展上同调”这场行动收尾,我还想在这里稍稍追随我的朋友兼昔日学生Deligne的几番辗转——看他如何呈现那部名为“SGA~4$\frac{1}{2}$”卷册的“中心主题”\fnote{493}{(*)}{况且,在“SGA~4$\frac{1}{2}$”中没有任何一处说“报告”构成了“中心主题”,也没有一处说其主要宗旨是为“那条”Weil猜想提供平展上同调的主要成分。在为该卷撰写双重导言之际,一份规模及于整个平展与$\ell$-adic上同调的侵吞意图,想必已然在场。},即“那条”迹公式,由之导出L函数的上同调公式。它就是“迹公式报告”(在其书中被引作[报告],前引书,第76--109页)的主题。

Deligne在该卷中足足有\textbf{四个}地方,对迹公式作了多少带点“历史”味道的评注。该卷的读者,倘若事先并不知情——无论他读还是不读那四个段落(我们即将逐一检视)——从中得到的印象都将是:某位Grothendieck(一个有点马虎、且晚于“SGA~4$\frac{1}{2}$”一卷的讨论班的作者或主持者——人们尤其劝人千万别贸然去读它)关于L函数似乎也曾有过一些想法,当然难免有点纠缠不清,直到那部辉煌卷册的作者终于出面,给出可理解的陈述与站得住脚的证明。

在整部书中,对这位先生唯一准确的引证,是指向某一场Bourbaki报告(1964年),在一句“评注3.7。”(前引书,第88页)的转角处;而它出现在那里,是作为最后一棒,收在一串三条愈来愈技术的评注末尾\fnote{494}{(**)}{写下这些文字时,我正被一种强烈的“同一性”感觉所震撼:我在此探测的文风,与四年后为“以轻蔑”侵吞“好上帝定理”(别名Mebkhout)而展开的那一种。这些翻云覆雨,我是在“变戏法者”那条注里发现的(这名字当得起大写……),n\textsuperscript{o}~75。在那里,“敏感点”藏在一条评注4.1.9里(而不是3.7),更是乱得可以。进步是挡不住的……(3月22日)我不曾注意到,事实上在“SGA~4$\frac{1}{2}$”里还有对1974年那同一场Bourbaki报告的第二处指涉——一处以精湛的艺术奉上的指涉,就在“阿里阿德涅之线”里,一如将在子注“双重含义——或行骗的艺术”(n\textsuperscript{o}~169)中所见。}

其文如下:

\begin{quote}
“若承认$\mathbb{Q}_\ell$-层的形式主义……\textbf{很容易把证明归结}到$X_0$为光滑曲线、$F_0$光滑的情形。这在[2]§5中解释得很清楚(对3.1而言;3.2同样处理)。”
\end{quote}

\origpage{883}\dmark{}(着重号是我加的。)总之,这位除了以谀媚的代号[2]\fnote{495}{(*)}{轮到谁就是谁——1970年(在尼斯国际大会上),是Serre(在Deligne的报告“霍奇理论I”里)得到了代号[3]的待遇,而不被提名——就在那句神谕般的话里,人们头一回也是最后一回向所呈现的理论之若干“源头”示警……}之外再不获点名的先生,(诚然没有做出,但)解释了那点平凡的活儿——平凡到几乎不配在这条收尾评注里提一句,而且还好心让人以为:平凡归平凡,这里至少解释得很清楚。(我们从那位辉煌作者另一些评注中已经知道:清晰可不是这位含混先生的强项……)换言之:这一章“迹公式报告”的宗旨在于\textbf{做真正的工作},把那些平凡的边角留给干这个的人……

既然说到这里,不妨马上指出:正是在这一页上,有着我方才提及的四个段落之一——它们含有关于“那条”迹公式的历史性评注。那是第3.8节(紧随3.7的评注之后,合乎规矩)。在那里人们解释:证明3.2(即迹公式在该卷中唯一被显式讨论的情形——Frobenius对应的特殊情形)有“两种方法”。不用说,那位先生的名字在两者中都没有出现。人们区分方法A,称为“Lefschetz--Verdier”方法,和方法B,称为“Nielsen--Wecken”方法(这个名字对我而言也意味着什么……)。让我们看它怎么说:

“方法B:Nielsen--Wecken。一种受Nielsen--Wecken工作启发的方法,允许把3-2[即Frobenius的迹公式]归结为Weil证明过的一个特例;这将在以下诸段中解释。”

\origpage{884}事实上,第5节(pp.~100--106)的标题,恰如其分地,叫“\dmark{}\textbf{Nielsen--Wecken方法}”。前文说过该方法\textbf{受启发于}Nielsen--Wecken的工作——所以那位卷册作者把这方法冠名为“Nielsen--Wecken”的,想必纯粹出于谦虚。这就更清楚了,因为他们不是当下的人。如果读者碰巧翻阅了某份报告XII——文中从未引用它,而且是在一个人们建议他遗忘的讨论班中——的文献目录,他会知道:那是两位在四十年代初发表著作的人。如果他甚至读了他们的精美著作(那位辉煌作者,我敢打赌,从未亲手翻过),他会知道:他们的方法是三角剖分技术。那显然不是正文里用的方法。在没有相反声明的情形下,那卷册的谦逊作者,便同时也是该方法的作者。该方法没有标注日期——无疑又是出于谦虚,免得人们知道:果真是他第一个辛辛苦苦证明了这条出名的迹公式。

再看方法A,所谓“Lefschetz--Verdier”方法,人们怎么说它。这可不怎么鼓舞人心:

“若~$X_0$~是真的……Lefschetz--Verdier一般迹公式允许把3.2的右端表示为局部项之和,每个~$X^{p^n}$~的点一项。在SGA~5的\textbf{原始版本}中,这一公式仅在奇点消解的模意义下被证明[人们早就猜到会碰上麻烦!]。读者将在\textbf{最终版本}中找到无条件的证明[又谦虚得忘了提醒:正是多亏他,这一步才得救——反正人们是决计不会去翻那该死的SGA~5的]。对于曲线的情形(可以归结到该情形,见3.7),成分[???——我们放弃了……]反正都已可用。”

可如果它们已经可用(一位比其他人更清醒的读者,如果有的话,也许会问),为什么还要这样长篇大论地谈什么“仅在……模意义下被证明”的Lefschetz--Verdier公式等等等等?刚才不是说了:\textbf{真正}的工作在维数\textbf{一}中进行吗?答案:这就是所谓“乌贼战术”——喷墨汁,好在浑水里摸鱼!读者到了这一步,已经完全确信这方法肯定不是对的那一个。他以一双倦眼扫过下一段,那段给他补上了最后一击:

“为从Lefschetz--Verdier公式推导3.2,须能计算其局部项[天哪,这是哪门子苦役……!]。对一条曲线与Frobenius自同态[啊!他们泄气了!],这已由Artin与Verdier完成[还要两个人一起干!](见J.-L.\origpage{885}Verdier,\dmark{}平展上同调中的Lefschetz不动点定理,载于Local~Fields会议文集,Driebergen,Springer~Verlag~1967)以及SGA~5的\textbf{最终版本}[人们不免好奇:那原始版本到底长什么样,可怜的我们!]。”(此处及上文着重号皆为我所加,纯属恶意!)

显然是发善心,那位辉煌作者才免去指向那部注定被遗忘的讨论班的相应报告,或者哪怕暗示一声“那条”公式确确实实就在那里!一位不弃不馁而好奇的读者,若碰巧去那里面翻找,会发现一份报告XII,顶着不寻常的标题“代数几何中的Lefschetz与\textbf{Nielsen--Wecken}公式,A.~Grothendieck著[还是那位先生,我发誓!],I.~Bucur整理[不认识]”。那位先生和他的伙伴,八成是把他们辉煌前辈的报告抄了一遍,再兴致勃勃地堆上多余的细节……

在那份出名的“报告”里,没有任何东西能让读者想到:除了Lefschetz--Verdier公式(或者毋宁说该叫Lefschetz--Verdier-Deligne公式,反正也不怎么鼓舞人心——那位作者自己泄气的评注表明了这一点)之外,还存在一份显式的迹公式,适用于\textbf{任何}自同态,而非仅限于Frobenius。无论在援引Artin-Verdier的那一段,还是在(下面要引的)指向SGA~5的另一段(特别是不提名那位先生)里,都被暗示:这工作\textbf{仅}在Frobenius自同态的情形下被做过。人们跟Verdier是哥们儿(而且也向他证明了这一点),但迹公式嘛,这是明摆着的:给Verdier一个“暂停!”式引用可以(与Artin并提\fnote{496}{(*)}{我早已领教过Deligne这套久经考验的技艺:把水搅浑、把鱼淹死,好掉包某个某人(此处是Verdier,毕竟是哥们儿,而且人家将在别处获得实质性补偿)——办法是把他与另一个人一口气并列提名,这样就没法指责他不慷慨了!这便是所谓“以同化作稀释”的掉包法。这方法中的艺术,在于找到那位与待掉包者“配对”的先生。就我而言,我的朋友回回都找Serre……},淹没在一大段技术性的、索然无味的文字中间,读过即忘)——但毫无疑问而且没有弄错:迹公式,那是\textbf{他,Deligne}!

诚然,Deligne不止一根弦,把这些带着“历史”架势(sic!)的评注散布在\dmark{}\origpage{886}四个不同的地方,实有深意——好在这一处补上在另一处可能被人指摘遗漏(或夸大)的东西。那里他还可以退守到同一章的引言——一切都已经安排好了!那是一段七行的引言,值得一字不漏地引用\fnote{497}{(*)}{(3月20日)我在昨天的反思中回到这段导言。(参见子注“双重含义——或行骗的艺术”,n\textsuperscript{o}~169$_7$。)}。

\subsubsection*{(b) ……与无稽之谈}
\tocsubsub{(b) ……与无稽之谈}

\notehead{169$_6$}{}

(3月18日)昨天不得不在兴头上停笔,因为时候已实在太晚,而且显然当夜写不完“公式”了!在重新接上围绕那份公式的几番辗转之前,我想先借这个机会,就“Lefschetz--Verdier”这条优美的公式,稍许“把脚伸进菜碟里”。这条公式以一种完美的例证方式,阐明了一样在我看来至关紧要的东西——在《收获与播种》全程乃至《引言》中\fnote{498}{(**)}{见《引言》4,“一场追寻显然之物的旅程”。},我不止一次回过头来坚持谈过、但措辞或许始终有些过于“一般”的东西。

这条公式是一份\textbf{深刻}陈述的触目例证,而其证明是“平凡的”(169$_6$bis)。当Verdier告诉我,他提炼并证明了一条Lefschetz公式——针对任意(“真的”,当然)代数簇上的“上同调对应”(当时连定义都还没有!)与任意可构“系数”——我起初难以置信。或许“带有或多或少一般系数的Lefschetz公式”这一念头曾掠过我的脑海——至少我很久以前就为“局部常数”系数(即局部系统)写过这样一条。但\textbf{我不相信}对一般系数也行——那看上去太美好了,好得不像真的!Verdier说服我没花多久。当场写下公式并给我证明,多半花了十五分钟——还得算我慢,尤其涉及要确信某桩如此出乎意料的事情的时候!这正是人们可以称之为一“\textbf{平凡}证明”的东西——就“众所周知”的事物而言。而顺着当今刮的风(那阵风,J.H.C.~Whitehead早已嗅到了最初几缕\dmark{}\fnote{499}{(***)}{关于这一点,见“年轻人的势利——或纯洁的捍卫者”一注(n\textsuperscript{o}~27)。}),只需一步(大多数人轻快地迈出了这一步),就能把定理本身也归为“平凡”——十条或百条公式中的一条,从上同调形式主义里“自己掉下来”的\dmark{}\origpage{887}——此处的形式主义,是那个我在此前一年(1963)刚在平展框架中发展的\textbf{完整}体系:六种运算,\textbf{以及}双对偶定理。

我说Verdier(循Lefschetz开辟的路)发现的定理“深刻”,绝非因为(尽管这理由也成立)其证明所从出的那套形式主义本身“深刻”。那股时风,早就(而且恰恰有Verdier本人的无条件支持!)把这形式主义归档为“Grothendieck式的大厚片”,挥手即扫,同时又暗中使用那些“厚片”于每一步(而不点名)。至于该定理“仍属猜想”(某某先生以怜悯的口吻强调),或在一切特征中都已完全确立(它如今确实如此——多亏恰恰是顶着这位某某先生名字的“双对偶定理”)这个问题,在我看来,当我说它是一条深刻的定理时,同样无关宏旨——它实质性地丰富了我们对于各种“上同调主题”的理解(离散或连续系数,各种“簇”或“空间”……)。同样的,对普通Lefschetz公式——比如紧(或无论如何)微分流形及其有孤立不动点的自同态——也可这样说:“形式”证明,从上同调对偶形式主义出发,一页足矣,几行亦可。然而在两个情形里,都有过\textbf{创造}——某种崭新的、实质性的东西,它此前逃过了所有人的眼睛,它“还不存在”,而骤然间现身了……

在当前的情形里,“创造”确切地定位于何处?我相信,不止一位数学家——尤其不止一位曾是我的学生的人,他们当年本知道什么是创造,却早已把它忘到九霄云外——应当在这个案例上,或在与他们更近的某个类似案例上,好好沉思一番\fnote{500}{(*)}{当年“在圈内”的人并不多(而且如今也不多,鉴于事态的走向……)——但除了Verdier和我,总该还有三四个。Deligne嘛,他还没在这一带露面……}。我知道,如果当年我向我本人、或向某位当时在“上同调形式主义”行内的学生或同事,提议显式地陈述一条一般Lefschetz公式——\origpage{888}针对任意系数与任意“上同调对应”(定义由他们临时去下!)——在一张紧(抱歉,是真的)任意\dmark{}簇上,每一个人都会万无一失地到达:花上几个小时或几天,或至多几星期\fnote{501}{(*)}{当然,我在此假定所论的那个人确已对提出的问题“咬住了钩”,亦即我对此曾有过的“感觉”(若非如此,我也不会把它提出来!)确已“传了过去”,而那位学生或同事确确实实“发动”了。一件事要“传得过去”,可绝不是显而易见的——远非如此!}。一旦问题被提出(哪怕还相当模糊,而主要术语尚待定义……)并被\textbf{看见},去“解决”它(在此即:找到\textbf{正确的}表述,由已有的上同调形式主义所暗示)便成了一桩“\textbf{例行}”的事(Weil在同一意义上称之为“习题”)。这种“例行工作”确实需要嗅觉、最低限度的智力与想象力,但(一如我不止一次写过的)此时“是事物本身在向我们口授”如何着手——只要我们肯倾听它们。(而如果我们不会倾听数学事物,我们倒不如改行……)我所说的那道\textbf{火花}——它迸发出崭新之物——并不在这一类工作中\fnote{502}{(**)}{而火花就更不会迸发于那样一件打杂的工作里——也许是十年之后做的——它只不过证明:使某条证明“行得通”的诸项假设,确确实实在人们以为的地方得到了验证……}。

创造的时刻——触发发现进程的那道火花——在此处,是\textbf{问题被看见}并加以\textbf{“承担”}的那一刻:当意图从\textbf{真正去看}、从\textbf{追到底}、从弄个水落石出以“看清”这条人人自以为“理解”的Lefschetz公式之\textbf{确切}的“真实”有效域中诞生。迸出火花的,不是某种“精湛技艺”或“强大脑力”(通常意义上驾驭复杂技术或记忆嵌套局面的那种“力量”)。那是一份\textbf{天真}:人人都自以为理解Lefschetz公式,而我,可怜的人,尚无已理解之感,我想把这事弄个水落石出!在这一类情形中,一旦上路,\textbf{便已赢了}:事物向我们吹拂该做的事,我们去做。“追到底”在一种情形里,可以意味着证明“那个”正确的定理(在此处以已有的形式主义为项——该形式主义本身是“已确立的”还是“仍属猜想的”,无关宏旨)。在另一种情形,它可以\dmark{}\origpage{889}意味着:提炼出“\textbf{正确的}猜想”\fnote{503}{(*)}{两种情形——其一是“火花”(被“追到底”)把我们引向一条定理,另一是我们被引向一条猜想——并非天性不同。“追到底”意味着:使一份尚且弥散的直觉完全具体化,从它的所有侧面、用我们掌握的一切手段去探测它。一条定理在本性上并不比一条猜想更“完成”。有种种显然暂时的定理(甚至跛脚而歪斜的),也有种种猜想(例如Weil猜想全体)给人一种全然完成、完美无缺的印象。这并不妨碍这些Weil猜想曾是通向其他发展(起初和它们一样是猜想的)的起点,而那些发展更宏阔,并把它们涵盖在内。在这个意义上,可以说数学中没有任何东西,只要它还生气勃勃,是“完成了的”或“定局的”。};而这猜想常常自身是临时的,或许日后会被发现为假或不充分,须加以调整或拓宽——这同样无关宏旨。这条猜想是通往对事物(在此即数学事物)更深切的集体认识之路上的一个台阶,一个无法绕过的台阶\fnote{504}{(**)}{关于发现的动力学,以及“错误”在其中的关键角色,见(在《收获与播种》第一部分的)“错误与发现”一节(n\textsuperscript{o}~2)。}

深刻与丰饶是紧密相连的品质——后者在我看来是前者有形的标记。Verdier所发现的公式之丰饶的第一个标志,当年(如果不是在几天或几周之内的话——我说不准了)就到来了:这条公式曾是\textbf{主要}动机,引我写出以任意$\ell$-adic层“为系数”的L函数上同调公式。在\textbf{技术}上,我根本没有用到Lefschetz--Verdier公式——这一点无关紧要。确定无疑的是:没有这条公式作为引路线头,或者不如说:没有它向我吹拂“应当有什么东西”——一桩以层“为系数”的L函数——没有这声执拗的声音,我根本不会想到去提炼\textbf{正确的}概念及与之相伴的公式;诚然我日后无疑也会到达,但那须得\textbf{先}靠自己的力量发现另一条更具一般性、恰好“挡在路上”、\textbf{必须}被发现的公式。

在心理层面上,两种情境极为相似。正如Verdier须首先提炼“上同调对应”的概念,以精确化“Lefschetz公式问题”(超出“普通”公式之外),同样,我也须提炼“以层为系数的L函数”的概念,以\dmark{}精确化“L函数公式问题”(言下之意:\origpage{890}超出与真的、光滑的~$X$~相关联的“普通”L函数……)。“创造的时刻”——火花闪过的一刻——是我\textbf{看见这问题}的那一刻:定义这种广义L函数——而且我\textbf{承担了}它,把那个问题进行到底。一旦问题被清楚地看见,而且假定我能使“上同调形式主义”行内的随便哪个周围的人“接住”它,那么显然他必将情不自禁地解决它——以唯一自然而合理的方式,花上几天工夫(一如对我而言必曾如此),定义、陈述、证明,一气呵成\fnote{505}{(*)}{我在此把证明的最后一步单独挑出来——我当年把它搁置了(因为它想必不构成真正的问题),而它险些变得更长。}。诚然,那些归结到维数一的“拆卸”是“容易的”,甚至(如果你非要这么说)“平凡的”。发现并不在这类拆卸之中——那种随便谁来都做得跟我一样好(或不屑去做)的拆卸。\textbf{发现}在于一个\textbf{概念}:一个此前无人想到过的概念,尽管它是\textbf{显然的}:以层“为系数”的L函数。在这个概念及与之不可分的公式里,蕴含着(在素域$\mathbb{F}_p$~上有限型概形的语境中,或更一般地,在绝对基环$\mathbb{Z}$~上)把上同调中的“六种运算”——从函子$Rf_!$开始(即“\textbf{几何}”性质的操作)——解释为“L函数场”上的操作,即以“\textbf{算术}”为项,加以诠释的可能。这是1949年Weil猜想所开启的方向上的新一步:经由上同调主题,朝几何与算术之联姻迈进。

这两项发现在这份自诩为平展上同调与$\ell$-adic上同调\textbf{标准}参考书的文本——出自我曾有过的学生中最有天赋、最有名望者的手笔——里,变成了什么呢?

Lefschetz--Verdier公式——它启发了我,而我从未“使用”过它——变成了\textbf{稻草人}:被适时地高高举起,好让读者(他只求相信!)领会:某个讨论班悬在一根何其纤细而乏味(还是“猜想性的”,更不要说“局部项未被计算”)\origpage{891}的线上,而那个讨论班(“按照本卷的精神”)人们善意地绝口不以指涉(除非是为了贬损它……);同时不忘在此处彼处轻描淡写地提醒:如果那条不巧的(而且说到底毫无用处的)公式毕竟已不再“属于猜想”,那得归功于那部辉煌卷册的谦逊作者。

至于以层为系数的L函数的概念——它是构成全书核心的那份“报告”的中心概念——它在“报告”第1.6节(loc.~cit.~p.~80)悄无声息地登场,没有半句评注指示其动机或来历。定义嘛终究是定义,无须辩护。对这概念的来历——不得不说它有点装神弄鬼(尤其当你囫囵吞下的时候……)——心存疑问的读者,可以在Artin-Weil(但他们那个时代尚无$\ell$-adic层,它们显然是由作者在同一卷中\dmark{}引入的……)与(更可能的)这位辉煌作者之间做选择:后者正在利落地把他引向那条著名的“迹公式”。

后者在第3节(loc.~cit.~p.~86)被引入,该节以这样的字句开篇:“Grothendieck对L函数的\textbf{上同调解释}是如下定理:……”(接下来是公式3.1——注意:着重号是我加的)。除了那一章引言(下文还会回到它)之外,这是全章唯一一处\fnote{506}{(*)}{(4月9日)有一处例外(起初逃过了我的注意),带着一处“暂停!”式指涉(在第90页),指向“Grothendieck对导出范畴理论的一处本质运用”(用以在“非正统”情形中定义迹)。}提到某个名字的地方。因此还是那位先生,两页之后以缩写[2]被指涉(作为一个“清楚地解释了”若干“容易的约化”的人),他也给出了那条装神弄鬼的“解释”3.1——冷不防地抛在那里。这反正也没什么了不起,读者马上(而且毫不意外地)会看到:证明占不上半页(同在p.~86),而且本来就是“经典的”:它是那条给“报告”命名的出名的“\textbf{迹公式}”——即显然构成“真正定理”(3.2)的东西——的简单推论。后者的发明权,没有任何名字被提及——依旧是这份谦虚的癖好,偏偏在那些最辉煌的人身上!两页之后(如昨天所见\dmark{}\origpage{892}),Lefschetz、Verdier、Artin、Nielsen、Wecken的名字纷纷出笼——真是一场谦虚的豪宴——就是为了不说出是他!

我在这里想强调的、在我看来远远超出本案与这股骗术气息的事情,是这个。无论对那称之为(名副其实)“Lefschetz--Verdier”的公式,还是对L函数(“以层为系数”)的“上同调解释”,恰恰是\textbf{这一点}——即它们的发现是\textbf{创造之举}——在今天成为普遍低估的对象(即便不是轻慢的嘲弄),人们习惯以带贬义色彩的形容词来表达,如“\textbf{平凡}”“\textbf{幼稚}”“\textbf{显然}”“\textbf{容易}”“\textbf{猜想性}”,甚至“软数学”“梦想”“大话”“无稽之谈”,以及诸如此类的好听话,听凭各人即兴挥洒。而工作的另一部分——我始终知道(尤其,我觉得,从未\textbf{忘记})它是“额外的”、由事物的力量尾随而来的,如同“后勤”必然追随(只要人们肯躬身去做)——即那\textbf{技术}部分,它常被冠以“\textbf{困难}”之名,要“凭腕力”来做,我也曾称之为“例行工作”(不带任何贬义)——恰恰是工作的这一部分,被当今的共识所推崇,被捧到了排斥其余一切的地步。

对我而言,“困难”的概念全然相对:一件事在我尚未理解它时显得“困难”。我的工作便不是“凭腕力克服”困难,而是深入我的不理解之中,达到某种理解,把曾显得“困难”的东西变得“容易”\fnote{507}{(*)}{读者将注意到,这是对一项困难的“阴”的、“女性”进路——“上涨的海”那条路——的一番描述。我在此并不是说这是唯一可能的创造进路——还有那条所谓“锤子与凿子”的进路,“男性”的进路,唯一享尊荣(且不说今日是唯一被容许……)的一条。关于这两种可能的进路,见“上涨的海……”一注(n\textsuperscript{o}~122);关于时下对两种进路各自通行的态度,见“肌肉与肚肠(阳葬阴(1))”与“天赐的机缘——或封神”两注(n\textsuperscript{o}s~106、151),以及承后者而写的“否认(1)——或提醒”(n\textsuperscript{o}~152)。}。例如,我为“L函数公式”以及在其他场合所做的那些拆卸——它们今天被当作“平凡的”\dmark{}——对我而言并不比处理那些被称作“困难的”“不可约”情形更“容易”\origpage{893}。它们是工作的不同阶段,仅此而已\fnote{508}{(*)}{我所想到的情形,除“带系数的”函数L的一般公式之外,是我做过“拆卸”以回到维数(或相对维数)\textbf{一}之情形的地方,主要就是平展上同调中的两条基变换定理(对一真态射,以及经由一光滑态射)——它们构成那两条使该上同调“可活”(Deligne语)的关键陈述;还有$Rf_!$的“比较定理”,在平展上同调与超越上同调之间(对复域上有限型的概形)。(还有仿射态射的(所谓“弱”)Lefschetz定理。)从心理上说,一旦得以回到这类所谓“不可约”的情形,我就感到(多少)“赢了”,期待的定理必将“出来”,而经验在每一个这样的场合都确证了这感觉没有骗我。然而从技术上说,正是这些拆卸代表着那所谓“容易”的一步。凑巧的是一种当时令我甚为触动的“天意”:处理两条基变换定理各自的两个“不可约”情形所需的配料,都已由我发展出来(当时毫无自觉),为了第一条在SGA~1里,为了第二条在SGA~2里,早三年和早两年……}。某一阶段在另一阶段之后,或比它更长,并不意味着它更“困难”。在两个阶段里,都须\textbf{一份想法}:此处的“拆卸”的想法(这在以前从未有人想到要在这种情形下做——而且理由充分,因为那涉及针对任何非Frobenius对应的不动点公式!);彼处则是一个无疑更精微的想法,受一条比Lefschetz原始公式更精致的(出自Nielsen--Wecken\fnote{509}{(**)}{(4月10日)况且,Deligne正是和我、以及SGA~5的其他听众一样,从我口中得知这条“Nielsen--Wecken”公式及其在平展上同调中的转写的——这使他无须再去翻那三位作者的(德文)漂亮文章(发表于1941至1943年间),而后来的用处,其别致一如人们所知(见子注“真正的数学……”,n\textsuperscript{o}~169$_5$)。}的)不动点公式启发,通过对系数层的精心剖分来实施,以合适的导出范畴为项加以表达\fnote{510}{(***)}{在这项证明中,导出范畴的语言是不可或缺的。在我离开之后,直到名为SGA~4$\frac{1}{2}$的那一卷出版之年左右,我的上同调学生们对导出范畴建立了一种默契而有效的抵制——而它们曾是发展对偶形式主义(“六种运算”与双对偶)之概念关键工具的,在“凝聚”系数、继而“离散”系数的语境中。尽管它在Lefschetz--Verdier公式的证明中、同样也在平展语境中诸“经典”对偶公式的证明里扮演着关键角色,这一形式主义本身,作为一种数学结构、一个融贯的概念整体,却遭到了同样的抵制,而且抵制延续至今(且不说“六种运算”这个名字本身——它始终是个忌讳)。有可能是迹公式证明的种种需要,促使Deligne于1977年朝解除对导出范畴的抵制迈出了第一步:在海盗卷中出土了Verdier“论文”的一副骷髅般的“第0状态”文本(文中对我的名字只字未提)。关于这一点,见专论“行动III”的注“分赃”(n\textsuperscript{o}~170);至于“论文”这桩滑稽事件的详情,见“同伙”与“赊账的学位论文与全险”两注(n\textsuperscript{o}~63$'''$、81)。}\dmark{}。\origpage{894}第二个阶段确实更长:当须以它应得的全部一般性将其完善时(鉴于曲线上的自同态不止Frobenius一种),出现了整整一“地毯”的“斯坦伯格式”非交换迹,它们最终粘合在一起,须我细心发展\fnote{511}{(*)}{(4月23日)一般性被Illusie在其SGA~5屠杀版导言(第二段)中按“理”判为“多余”,同时毕恭毕敬地回声他那显赫的朋友Deligne——后者(不作进一步说明)指涉他会“剪除”的“无用细节”。这番贬损同时也一劳永逸地使他免于让读者疑心:在维数一存在一条比他为Frobenius所呈示的更一般的显式迹公式,其中他步步重演我的证明,却给人印象以为那是他的独创。见下一条子注“双重含义——或行骗的艺术”。}。那很漫长,也很“容易”——而它也是一桩\textbf{必须}做的事,这一点我十分清楚。但即便找到这类使工作“变容易”(或干脆使之成为可能……)的想法,在我看也属于“例行工作”的一部分。它给这工作平添一份魅力——那通常使它超越单纯摇动手柄的东西。

工作的\textbf{创造}性部分,却是那个\textbf{幼稚}的想法:那个人人早该看见的想法——若不是几年,就是几个世纪、几千年以来——而偏偏无人看见,虽然它一直在众目睽睽之下,人们每回都不得不绕一个大弯以免撞上它!

当你遭遇到这样一个想法——无论是你自己“撞上了它”(这可是实至名归的说法……),还是别人向你解释了它(如Verdier有一天向我解释了一个)——你觉得自己蠢透了:无法原谅自己此前竟没有看见它,当它恰恰是最自然、最显然、最“傻”的那一个——恕我直说……人们\textbf{早该}撞上它,这是肯定的;可偏偏没有……

看起来,当今(尤其当人们占着强势的时候……)在这种情形下,人们以圆滑来补偿\dmark{}:当\textbf{另一个人}(一位显赫的无名氏,或某个\origpage{895}早已入土的“亡者”……)不巧带着这样一个想法凑上前来(或某天曾凑上来……)。“可是我可怜的朋友,你跟我讲的这些\textbf{太平凡了}!为了向那不幸的人证明这有多平凡(顺便把他放回原位……),人们只消用说完一句话的工夫就把它重新做出来——你这就见识见识什么叫做数学!我们袖子里可比这些初出茅庐的(或被冷落一旁的……)多的是东西!只需稍加拉拽,呼——再一拉——阿布拉卡达布拉,好嘞!而且\textbf{瞧},一条像模像样的陈述就从帽子里出来了,甚至一整套理论,一点不掺水,这才是干活呐,没错!年轻人,回去穿好衣服,等你学会来这一套再回来吧!”

我在这里不假思索地,把我的“身后学生”Zoghman~Mebkhout的遭遇做了个速写:他当年不过是里尔(或天知道哪里)的一名微贱助教,落在我的“隐形学生”Pierre~Deligne——一所精选机构的头号台柱(其余且不论……)——的手中;那段遭遇发生在恩典之年1981年,而且延续至今……这便是“第四行动”,所谓“当值无名之辈”行动(或“反常层会议”,为免提名)——四场行动中最令人难以置信的一场。(关于它,见“封神”那条注,n\textsuperscript{o}~171.)

但同时,在写下前一段时,我隐约觉得在或多或少重写某段我已在另一场合写过的东西……

没过多久我就想起来了——那是在《收获与播种》第一部分,写于如今整整一年前,在“竞技数学”一节(名字不言自明)里,n\textsuperscript{o}~40(p.~105)。我在那里唤起的插曲与反常层会议插曲的区别在于:这一回,“当值无名之辈”的角色,由“那个踩着我花坛的小白脸”扮演,而那位高傲而“竞技”的大东家,不是我的某个可鄙前学生,而是我自己。诚然,我不认为自己走到了(象征性地,在此场合)把别人的想法据为己有的地步。但我无法凭良心起誓,而且须由当事人(二十年之后——但迟做总比不做好)告诉我\textbf{他}如何回忆那段插曲,它在我的记忆中有些模糊。他遭遇的不幸是:重新做了一些我早已烂熟的东西(其中包括经由“拆卸”从约化情形构造非约化概形的Picard概形……),而那“弄得不愉快”——这就是\dmark{}留给我的印象;但我不会起誓说他的进路(在一个不如我一般的框架里,这是肯定的)当真完全被我的进路所覆盖\fnote{512}{(*)}{我从不曾有机会把所说的皮卡概形“相对”构造(通过在幂零理想上“拆卸”)整理成稿并发表——那原定为《代数几何原理》后续一章的内容(而该章从未问世)。况且,当我谈对他人一个想法(无论大小)的“侵吞”时,未必是指通常意义上的抄袭,即把这个想法(哪怕是经过修改和完善的形式)奉上而不指明其来历——在我看来,这事正变得越来越平常。但侵吞也可以是那种漫不经心的轻蔑之侵吞,它的气息使发现的喜悦枯萎,仿佛只为挫败它而取乐,配上一副厌倦的“哦!只不过是这个……”的腔调。那腔调不言而喻地让人懂:我们刚刚被告知的东西,人家早就知道,不如说,从来就知道;而如果还不曾费心把它讲明白,那实在是因为它值不得费这个神……对于这类腔调(或其先祖),见(在《收获与播种》第一部分的)“挫败人的权力”一节(n\textsuperscript{o}~31)(在已引的“竞技数学”一注(n\textsuperscript{o}~40)中重提);以及(在70与80年代更严酷的氛围中)《葬礼I》中的“侵吞与轻蔑”(注n\textsuperscript{o}~59$'$)。}。

不管怎样,我必须在此再一次确认一份\textbf{亲缘}:它存在于六十年代至少在某些时刻属于我的那种态度,与我在我某些昔日学生身上所见到的那种态度之间。他们向我送回的那个人(从前的我)的一副确实被扭曲了的形象——一幅\origpage{896}多年来我一直想拒斥的形象。但如果《收获与播种》——它首先是对我数学过去的反思——曾有过\textbf{一个意义},那便是让我同样理解:即便我某些昔日学生乐于否认我,我也不该否认他们中的任何一个。经由他们回送到我这里的东西,属于我曾参与播种的那些收获,一如他们本身也参与了播种。而这份我以一支毫不留情的笔、在即将满三周之际做出的确认,不是对任何人的控诉,而恰恰是一份\textbf{确认}——它把我自己卷入得与任何人一样深。

\subsubsection*{(c) 遗产——或骗术与创造}
\tocsubsub{(c) 遗产——或骗术与创造}

\notehead{169$_6$bis}{}

(4月10日)\fnote{513}{(**)}{本子注源出前一条子注“……与无稽之谈”(n\textsuperscript{o}~169$_6$)的一条脚注;见第886页转引。}如众所周知,“平凡”一词在数学中的意义极为相对。在此,所谓“平凡”,我指的是:以人们理当“已知”的东西为项——即(在此情形)六种运算的形式主义,以及双对偶定理(后者在离散平展语境的特征~$p>0$~中,在Deligne找到证明之前仍属猜想……)。以这套形式主义为项,证明的原理可以在几分钟内(与陈述一起)以全然令人信服的方式\dmark{}\origpage{897}解释清楚。这当然不免于写一份正式的证明,那意味着验证若干烦冗的相容性。

在这种情形下的惯例,是定理的作者(尤其当它重要时)肯费心写下一份证明。在Verdier的情形,我毫不怀疑:这是他的全部成果中最深刻、就其意义而言最重要的一条——在他有幸(而且在此情形下理所当然)冠以自己名字的成果之中(借用Weil的惯用语)。然而他对待这条定理,一如对待导出范畴理论:既然荣誉横竖都能到手,他便不认为值得去做这项工作、把一份完整的证明交到所有人手中。

这是一个富有雄辩性的心态之标志——我在此彼处谈到过它,最近还在“操演”(n\textsuperscript{o}~169)那条注(2月28日)的末尾。我得以确认它已蔚然成风。诚然,那被称为(带上文保留)“Lefschetz--Verdier”的公式,确确实实是Verdier的一项创造之举——在他尚与我共事、尚对工作满怀热情的岁月。我看到两件事之间的直接关联:他从未尊重到去证明“他的”定理,以及\textbf{他的数学家生涯中再未有过类似的创造之举}。创造的时刻,唯有在我们“配得上它们”的时刻——亦即:处在接纳它们的状态——才降临我们……

这条被一位逃之夭夭的父亲弃置不顾的优美公式,此后还经历了奇特的坎坷。它先是成为我在SGA~5中最早报告之一(第III讲,1965年)的主题。Illusie承担了它的整理工作,却在十二年间不认为值得费这份工夫。而后,在他与Deligne的完美同谋中(我猜,还至少得到了Verdier的默认——Deligne将向他支付实质性的补偿),它变成了“特洛伊木马”(或如我下文所写的“稻草人”)的头颅,被熟练地操控,以使那场令人难以置信的冒名顶替——名为“SGA~4$\frac{1}{2}$”——显得可信。后者被凭空搭建,旨在埋葬他们三人共同的老师——也就是说,归根结底,该公式的“\textbf{祖父}”(那条公式,若非有我这微不足道的人与随我一同被埋葬的六种运算,恐怕再过一百年也写不出来……)。作为风俗画,这才是风俗画!

如果我那些亲爱的上同调前学生们,不是把自己挥霍在这类骗术里扮演矮子(他们不是),骑在一个巨人的\dmark{}\origpage{898}肩膀上(我也不是),而是——在这十五年里——让存在于他们身上的创造力自由发挥,一如让存在于我身上的创造力自由发挥,那么晶体系数、De Rham-Mebkhout 与 Hodge-Deligne诸理论,以及以“神秘函子”为钥匙的那一套,肯定早已抵达六种运算形式主义的“完全成熟阶段”了。而且(我有一两个星期以来就怀疑这一点……),那位曾是他们的师父的人的宏大梦想——那个“Motif”,本是用来成为旋律的,却在同样这几只手里变成了采邑、宝藏与“模糊的骨架”——它也早已化身为一首宏大的交响曲(绝非“猜想性的”,而同样是“完全成熟的”),并从今天起便成为\textbf{全人类的共同财富}。

\subsubsection*{(d) 双重含义——或行骗的艺术}

\tocsubsub{(d) 双重含义——或行骗的艺术}

\notehead{169$_7$}{}

(3月19日)但我必须回到我朋友Pierre~Deligne在介绍著名的“迹公式”时的“翻云覆雨”。咄咄怪事的是,他从不曾指明:就Weil猜想本身的应用而言(那想必是从实用角度被摆在第一位的,即便不是唯一的位置),根本不需要什么精致的公式与证明——“寻常的”Lefschetz公式(平展版)就足够了\fnote{514}{(*)}{(4月25日)我在此可能出错,因为我还不曾真正了解Deligne对Weil猜想最后一部分的证明,那关乎 Frobenius 特征值的绝对值。看来,Lefschetz束的使用引导他引入比函数$\zeta$(即“寻常”$L$函数)更一般的$L$函数。}。而以下这点当然并非偶然:他偏偏选中了关于“结合于一个闭链的上同调类”的那场报告,把它从 SGA~5 中“借”来,不加任何手续就并入他的摘要——而那场报告恰恰含有(除了“寻常”版的Poincaré对偶,平展版之外)建立“寻常”Lefschetz公式的关键配料——只消几汤匙的工夫。于是人们不禁要说:他本满可以免掉收入这份不三不四的“报告”,它不过为单独的 Frobenius 自同态建立了一条迹公式(同时还执拗地向读者隐瞒:他可以在别处(!)找到明显更一般、而且同样“显式”的公式)。可他既然费心写了这份“报告”,那想必出于两个相互牵连的理由。一方面,六十年代就已十分清楚:Weil猜想一经以“权”为项妥善重新表述,便对奇异簇与非常数“系数”仍保有意义。的确,人们从此可以把它们全然以几何的措辞来表述,无须明确提及\origpage{899}$L$\dmark{}函数的形式主义。我想,Deligne的文章“Weil猜想~II”里所做的正是如此(其中当然丝毫不曾暗示:对于他在其中所证明的那个主要陈述的提炼,我曾出过什么力)。然而,几何-上同调运算的算术诠释(以“带系数的”$L$函数为项)看来仍必定要扮演一个角色,而在那里,以我所发展的形态出现的一般$L$函数公式,将取得一个关键的位置。从长远的眼光看,那就必须在那部被命名为“SGA~4$\frac{1}{2}$”的书卷里提供一份引证。同时,鉴于一般迹公式(Lefschetz--Verdier式)已然构成上同调装备的一项重要配料,这也有助于造成那样的幻觉:这部书卷(如他所宣称的)确乎提供了一套基本完备的上同调装备,以满足“非专业的”$\ell$-adic 上同调使用者的需要。

我还得逐一检视那剩下的三段文字——在“SGA~4$\frac{1}{2}$”中装作就迹公式给出若干历史说明的段落共四段。我将按它们在书卷中出现的次序引用。前两段位于整部书卷的开头(《引言》第1页,以及“阿里阿德涅之线”第2页),其用意显然在于“摆明态度”。它们想必也是被读得最多的两段。第三段是“迹公式报告”一章的简短引言。(第四段前天已经谈到过,它属于同一份报告的正文,想必是所有四段中最少被读的一段。)

在那份置于“SGA~4、SGA~4$\frac{1}{2}$、SGA~5 之阿里阿德涅之线”之后的文献目录里,SGA 这个缩写被解释为“布瓦玛丽的代数几何讨论班”,而不消说丝毫不曾提及我个人。可我毕竟列名于 SGA~4、SGA~5 的主持人之列。况且,主持人的职务想必纯粹是柏拉图式的:把 SGA~4 与 SGA~5 的主要报告过一遍(就再也别谈它们了……),其中谈到Artin、Jouanolou、Houzel、Bucur的报告,却没有一篇是我的。在关于 SGA~4 与 SGA~5 的条目里,没有任何日期——我在整部书卷中不曾找到任何一处暗示,能让未事先知情的读者怀疑:SGA~5(“即将由 Lecture Notes 出版”)如其名义所示,确确实实是\textbf{晚于}那部名为“SGA~4$\frac{1}{2}$”的书卷的\fnote{515}{(*)}{竟没有只言片语能让读者猜到这个“不许阅读的讨论班”讲的是什么,连它的标题(“$\ell$-adic 上同调\textbf{与函数}$L$”)都无人知晓!}。万一对 SGA~5 中某场报告有所提及(通常并不指明是哪一场),反倒会被“明确”\dmark{}\origpage{900}说明那是“零状态”或“原始版本”(言下之意:臃肿而读不通,可想而知……)。这些对 SGA~5 的指涉(是给一位未获知情的读者的,而人们劝他既不要查 SGA~4,尤其不要查 SGA~5),在这位读者的心中便是指向一部\textbf{晚于}他手中正在读的这一部的文本。况且我怀疑,这类未获知情的读者占绝大多数,而其余的人(如我在别处所写)开始老去,并将体面地死去……

我引用《引言》第3段的第一页:

\begin{quote}
“‘迹公式报告’含有对 Frobenius 自同态之迹公式的一份完整证明。该证明是Grothendieck在 SGA~5 中所给出的那一条,剔除了一切无用细节。这份报告当使使用者得以忘掉 SGA~5;人们尽可以把后者看作一系列离题,其中某些很有意思。\textbf{它的存在将使人不久能以原样出版 SGA~5}。”(着重号是我加的。)
\end{quote}

这篇文字有\textbf{两个相反的意思},以精湛的艺术\textbf{同时}奉上。对于了解该公式(对 Frobenius 而言)历史的人,他也许会惊讶于这番满不在乎的呈献(而若他充分了解 SGA~5 讨论班的来龙去脉、以及它在培养那位辉煌而放浪的作者时所扮演的角色,他会加倍惊讶);但他会以为,作者至少指明了其证明的出处。而对于未获知情的读者,他得知:他手中这部书卷里的证明,也见于一部较晚的文本 SGA~5——那文本出自Grothendieck之手,塞满了无用细节,是那位先生在原证之外为取乐而添上的。所引的这段对原证本身语焉不详。如前天所见,在那份“报告”里亲自读那证明,几乎不容人怀疑“SGA~4$\frac{1}{2}$”这部书卷辉煌的作者就是它的生父。当然,人们到处都不屑于指明:写出迹公式这个念头是谁起的;说到底,只要不肯费心去证明,写点什么是不费分文的!也丝毫不提Verdier(是他头一个给出了我当年搁下的“关键情形”的证明)。无疑并非偶然:正当谈到迹公式、正当谈到“那条”猜想的核心时,作者才一味殷勤,满口“无用细节”“离题”(当然很有意思,人要么表现得有风度,要么不\dmark{}那样!)\origpage{901}劝人忘掉,最后又添上一句既含蓄又专断的提醒——“它的存在将使人不久能以原样出版 SGA~5”,仿佛 SGA~5 之所以能“立得住”、之所以可出版,全靠那部名为“SGA~4$\frac{1}{2}$”的文本的“存在”——而它想必正好为那位先生提供了所需的一切,好把在原初文本中以简单方式做出的东西,以复杂的方式呈献出来,如下文所示。

% ---------------- p.712 ----------------
在那份“阿里阿德涅之线”里,我早已(在“操演”那条注(n\textsuperscript{o}~169)的子注“特洛伊木马”(n\textsuperscript{o}~169$_3$)里)指出第2页那两个相连段落(第2、3段)共十七行文字,堪称“浑水摸鱼”之术的“典范”。第二段恰恰关乎那条著名的迹公式。这两段值得在此全文照录:

\begin{quote}
“平展上同调中存在着一种与凝聚对偶类似的对偶形式主义。为建立它,Grothendieck使用了奇点消解与纯粹性猜想(陈述见 [Cycle] 2.1.4),后者在 SGA~4 XVI 中于相对框架下、在 SGA~4 XIX 中于等特征并模去奇点消解的情形下得到确立。关键之点则在 [Th. finitude] 中以另一种方法确立,适用于正则零维或一维基概形上有限型的概形。若干发展见于 SGA~5 I.~在 SGA~5 III 中,人们展示这一形式主义如何蕴含那条极一般的Lefschetz--Verdier公式。

可见,在 SGA~5 的原始版本中,Lefschetz--Verdier公式只是以猜想方式确立的。此外,局部项在那里也不曾计算。为应用于$L$函数,该讨论班含有另一条证明,它在 Frobenius 态射的特殊情形下补足了前者。那就是见于 [Rapport] 的那一条。其他引证:关于陈述与拆卸之概览,见Grothendieck的Bourbaki报告 [5];关于(由Grothendieck做出的)把关键情形归结为Weil已处理过的一类情形的简要描述,见 [2] 第10节;关于对后一情形的$\ell$-adic 处理,见 [Cycle] 第3节。”
\end{quote}

第一段我已经在那条注里评注过了(另见该注第872页的脚注(**),关于那句千金难买的“若干发展见于 SGA I”)。如今且容我追随我的朋友在第二段里的翻云覆雨(或者说其中的某些——太多了,追不过来!)\origpage{902}。那\fnote{516}{(*)}{更确切地说,他让人清楚地以为:单单那份34页的“报告”就包含了(而且是更好地包含)SGA~5 中一切可能有用的东西(后者即便在屠杀版中也还有近500页)。好一堆“离题”,原来都是白费的!}\dmark{}前两句话,一遍又一遍地念叨那条永恒的Lefschetz--Verdier公式,仿佛整个 SGA~5(以及其中某条从未被明白道出姓名的证明,为了某条迹公式……)都死死依赖着它、终生依赖着它,显然属于“乌贼战术”之列:把清楚的事情搞乱,好浑水摸鱼\fnote{517}{(*)}{说Lefschetz--Verdier公式曾是“猜想性的”,这并不恰当。它是建立在“具备一套对偶形式主义(‘六种运算’与双对偶定理)”这一假设之上的,并且于1964年由Verdier在这一形态下确凿地获得了证明。该证明当然已在口头讨论班中给出,而且是完全的。在特征~$p>0$~中仍然“猜想性的”,是双对偶定理的有效性,而它(如我们所说)是在“SGA~4$\frac{1}{2}$”的“有限性”一章中得到确立的。

至于Lefschetz--Verdier公式的局部项,它们被“计算”的程度,与寻常Lefschetz公式(不动点未必“横截”的孤立不动点)不相上下,并且推广了后者中出现的经典“相交重数”。说这些项“不曾被计算”,与说一个向量空间(其定义未被\textbf{指明})的维数、或一个系数未定的多项式的根“不曾被计算”,含义同样多寡。“计算”,在这些场合一如在别处,意味着:在一个被指明的“具体情形”中(例如对Lefschetz--Verdier公式取维数一)确立两项之间的一项\textbf{等式},而两项中没有哪一项比另一项更“被计算”或更已知(例如,在Verdier定义的局部项、与某些结合于 Artin 导子的局部不变量之间……)}。

而那句一语双关的要紧话,却紧接在“淹鱼”之后:

“……该讨论班含有另一条证明,它在 Frobenius 态射的特殊情形下补足了前者。”

那位获知内情却又性急的读者(而哪位读者不性急呢……)有一瞬间被“该讨论班”这个说法的含糊所困扰——究竟是 SGA~5 呢,还是“SGA~4$\frac{1}{2}$”?——而由于他知道在 SGA~5 里确有一份完整的证明,便又一次这样裁定了:作者的确(以颇为含糊的方式)指涉了它。我去年4月头一遍读时(见“白板”一注,n\textsuperscript{o}~67)差点也这样上当,可到底对不上。我明明知道,我所给出的那条显式迹公式的证明,绝不限于“Frobenius 态射的特殊情形”。而且令我震惊的是:人们刚刚还在(用些莫名其妙的“论证”)强调某一种 SGA~5 的呈献(在它的“原始版本”里,拜托!)是“不完整”的——这边猜想性,那边局部项没算出来……有了这个被两个逗号夹住的“补足了”,这种断然的对照便不可抗拒地向未获知情的读者暗示:无须他自问,“该讨论班”显然\textbf{正是}\dmark{}他捧在\origpage{903}手里的“SGA~4$\frac{1}{2}$”这部书卷——而且人们紧接着就在下一句里告诉他去哪儿找:“那就是见于 [Rapport] 的那一条。”而事后去读所引那一章里的该证明,当然也绝不可能在这位读者心中引起丝毫怀疑\fnote{518}{(*)}{见前天子注“真正的数学……”(n\textsuperscript{o}~169$_5$)。}!

% ---------------- p.713 ----------------
要紧处中\textbf{那个}“另一”字被\textbf{着重}了,而这全然不是我朋友的惯常做法。这是两段引言文字中唯一被着重的词,而且如无差错,也是整部书卷中唯一的一个(标题、陈述与新引入的术语除外)。他如此刻意要使这个词跳出来,想必不会没有缘故。(只是到此刻,这件事才引起了我的注意。)这个“另一”的效力——尤其当它被这样推到前台时——在于强调:“那条”公式曾有\textbf{两种}证明:其一是\textbf{不完整}的,而人们刚刚才就这条“Lefschetz--Verdier”公式那绝不出色的处境说了几句(实在拿不出手!);(而在前天看过的著名的《报告》那段更技术的文字里,人们又恪尽职守地就这个凄惨的话题再度施压……)至于那条跛脚的办法究竟有没有多亏那位辉煌作者的有限性结果而终究走得通,谁也永远不会知道这事最后怎么收场。但在这个“反差”效果(归根结底与前天检视的是同一个)之后,驯顺的读者那心理反射便越发专断:与其要那种在某个含糊讨论班 SGA~5 里的\textbf{不完整}办法(它如此不完整,以至于连一份确切的指涉都压根谈不上\fnote{519}{(**)}{在名为“SGA~4$\frac{1}{2}$”的书卷中,我从未找到任何一处指及 SGA~5 中某场报告,其中或含有一条不动点公式的证明,或含有什么$L$函数的著名“上同调理论”。事实上(见下文)已经明确说过:“本着本书卷的精神,将不诉诸 SGA~5……”!}),那种办法人们当然永远乐得避之不及,倒不如在这个\textbf{讨论班}里得到那份好而扎实的证明,那份\textbf{完整}的证明,它已在专门为此而设的报告——“迹公式报告”——里向人们张开双臂,准没错,人们毫不费力就能找着……\fnote{520}{(***)}{最美妙的是:事实上,Deligne的证明正是他1966年在 SGA~5 讨论班上与其他听众一起学到的那一条的忠实复现。}。

那句“该讨论班”简直是神来之笔——我那处于“暂停”状态的朋友在这个词上是攻不破的。然而,无论是在所引那段文字里,还是\dmark{}在“报告”中那段更技术的上下文里(第88页,论及那种注定被遗忘的、所谓“Lefschetz--Verdier”的方法),他都还是又一次冒险\textbf{明明白白}地说出了(或者说,至少半明半暗地说出了)\fnote{521}{(*)}{“又一次”,因为他已在前一段里(这一回甚至说得更明白)“说了假话”,一如我们在子注“特洛伊木马”(n\textsuperscript{o}~169$_3$)中所见。}\textbf{假话}。在两处他都强调(这回倒用得着这个词了)有这么一种办法(人们猜得出,就是 SGA~5 里那条不幸遵循的办法,天知道是在它哪一场“臃肿”的报告里……)用来证明 Frobenius 的迹公式,无非是\textbf{使用Lefschetz--Verdier公式}。然而,除了阿利贝尔1982年的论文(它在维数一情形,对一条任意的、不动点孤立的上同调对应给出局部项的计算)之外,“关键情形”只有两种证明:Verdier的和我的,而其中没有一种(阿利贝尔的也一样)使用Lefschetz--Verdier公式!这原是一桩棘手的事,而且长期未获解决(如今看去也有点旁支):证明 SGA~5 中所提炼的那条显式公式里出现的局部项(对应于远比 Frobenius 对应一般的对应关系),确确实实就是Lefschetz--Verdier公式的那些局部项。Illusie终于把它验明了,照他在 SGA~5 屠杀版导言(第~VI~页)以及他的报告~IIIB“局部项计算”(第~139~页)的导言中所宣布的\fnote{522}{(**)}{至于Illusie这些突然努力的动机,见子注“恭贺——或新风格”(n\textsuperscript{o}~169$_9$),尤其是第916--918页。}。

% ---------------- p.714 ----------------
\origpage{904}如果说Deligne还是不辞辛劳地制造这虚假的印象,那并非没有缘由。因为由此他同时造成这样的印象:SGA~5(那个“技术性离题”的讨论班,“本着本书卷的精神将不指及它”,其用途就是让人“忘掉”它)依赖于那条“猜想性的”公式——况且按其原样也无法使用(局部项没算出来, sic……)——而它之最终得以确立,全赖Deligne在那部名字再明白不过的“SGA~4$\frac{1}{2}$”(读者正捧在手里)里做的事;而(即便只凭这一点)那个较晚的、“含糊的”讨论班 SGA~5 也就依赖于它……

至于所引那段最后一句(以“其他引证”(sic)开头),它同样是这一类把戏的典范,为的是不说那位含糊的Grothendieck先生早已给出过一份\origpage{905}完整的证明——早了整整十一\dmark{}年(在那个“较晚的”、注定被遗忘的讨论班里……)——并且该证明被忠实地复现于“报告”之中。必须造成的印象是:那位先生只做过一些初步的模糊约化,而困难的情形归功于Weil,并由作者以(一次“$\ell$-adic 处理”)辉煌地接手。引到一本读者或许听说过的Weil的权威著作,再加上一条内部引证,这很能增添分量——要么严谨、要么熟悉自己的经典,二者必居其一!偏偏巧得很,对Weil那本书的指涉没有任何日期,也没有章节或页码——那位辉煌作者似乎并不想鼓励读者到辉煌的书卷本身之外去翻检,而一涉及它,引证顿时变得再精确不过(章节、段落全有)。

那条有名的“Weil已处理过的结果”,说白了不过是代数曲线(代数闭域上连通光滑射影)情形下的\textbf{寻常}Lefschetz公式——Weil在四十年代以边缘的办法就能表述并证明它,那时还不曾掌握上同调工具(却用了雅可比簇来定义所缺的$\ell$-adic $H^1$)。在“抽象”代数几何的语境中把这条公式提炼出来,当时原是一个重要的新想法,而且想必正是它把Weil引向了他那著名的猜想。一旦掌握了上同调形式主义,所说的这条Lefschetz公式便大体上是平凡的了。但若是明白说出那位先生的约化就是约化到寻常Lefschetz公式(为此人们骄傲地、不提名字地指到辉煌书卷的“Cycle”一章——那一章恰好是从 SGA~5 抢来的……),那就会给人一个印象:连那条“约化”索性就是那神圣公式的一项\textbf{证明}了。您可别想这样!(*)

我可急着收尾了!剩下的就是“迹公式报告”那一章的引言,前引书第76页,照录如下(略去其最后两行,那指涉书卷作者的一篇概述性文章):

\begin{quote}
“在本文中,我试图尽可能直接地陈述Grothendieck的$L$函数上同调理论。我十分贴近地追随Grothendieck1966年春在 IHES 所给出的若干报告。本着本书卷的精神,将不诉诸 SGA~5——除了两处指及该讨论班报告~XV 中独立于其余部分的段落。”
\end{quote}

乍一看,人会觉得作者毫不隐讳地指出了他的来源,谈到\origpage{906}“Grothendieck\fnote{523}{(*)}{(5月11日)由此可见,这里全部的“暂停!”艺术,就在于在两个相隔遥远的地方(第2页与第88页)指到\textbf{两}处“约化”(!)(容易的,这不用说)都是那位先生做的(只被点名一次,第二次就没有了……),却不让天真的读者有丝毫可能疑心:同一位先生居然\textbf{发现了}并\textbf{证明了}一条迹公式;而且他的证明(注定被遗忘的)被忠实地复现于辉煌的“报告”之中……}的$L$\dmark{}函数上同调理论”,还补充说他“十分贴近地追随”我的若干报告。在\textbf{一部寻常的}书卷里,这没有什么可说的。但同样真切的是,\textbf{上下文}参与构成任何文本的意义。“SGA~4$\frac{1}{2}$”这部不寻常的书卷,深刻地改变了这段话的意义——对于一个天真的读者,他事先已被先前读到的内容提醒过,而事后在读《报告》本身的过程中还会更受教益。到那时他会得到这样的印象:那位慷慨的作者把那名叫Grothendieck的含糊先生记上一笔“$L$函数上同调理论”,实在是天大的善意——而那理论归根结底似乎不过是一个有点神神道道的“诠释”,而且说到底还是\textbf{平凡的}。它只需半页光景就得了证,作为一条“迹公式”的直接\textbf{推论},而那条迹公式可一点也不寒碜,而且除了辉煌书卷的那位太过谦虚的作者之外,谁也配不上说是它的父亲。

% ---------------- p.715 ----------------
诚然,文中说了:在他的“报告”里,作者“十分贴近地追随”那位先生1966年春在 IHES 给出的若干报告。关于这些必定臃肿不堪的报告,再没有多说一个字,除了书卷作者肯为其报告保留的部分之外,它们想必早已彻底失传。究竟是些关于 Frobenius 的呓语(为此人们还会慷慨地指到那个同一位先生“主持”的 SGA~5),还是些关于$\ell$-adic 层的一般议论,抑或是某些下文将谈到的“容易的约化”——全然一团模糊。无论如何,那些想必多半是些“无用细节”,而读《报告》正好为我们免此苦役,谢天谢地——人们要求的也不过多。于是,给那位先生盖上一层贞洁的薄纱,大家且干\textbf{真正的活}去!

而当涉及对那位先生的指涉时,我的朋友喜欢含糊其辞(只要他不是干脆不提),这一回却似乎不能怪他不精确了:在 IHES 给出的报告,1966年春。他若是再精确一丁点,就会添上一句:是 SGA~5 讨论班\textbf{上}的报告。

SGA~5?那不就是那个讨论班吗——它(不着日期)出现在“阿里阿德涅之线”的文献目录里,还注明“即将由 Lecture Notes 出版”?也就是说,人们(想必)已经明白,那个讨论班的所作所为,就是在它前面那个 SGA~4$\frac{1}{2}$(那才真是无可指摘)之上,再添些“离题”(其中某些很有意思,这我们同意)和些“无用细节”。别瞎扯了,SGA~5 哪是什么1966年春的事,您别逗了!而最好的证据就在您眼前,\dmark{}白纸黑字印在“迹公式报告”(Pierre~Deligne著)那篇\origpage{907}刚刚引过的引言里:

\begin{quote}
“本着本书卷的精神,\textbf{将不诉诸} SGA~5”。
\end{quote}

那么,这就清楚了,不是吗?!

\subsubsection*{(e) 变戏法者——或飞走的公式}

\tocsubsub{(e) 变戏法者——或飞走的公式}

\notehead{169$_8$}{}

(3月20日)我开始感到有些疲倦,若不说精疲力竭的话,为我已逾三个星期的这项工作——尤其是(逐细地)最近这几天——孜孜不倦地“拆解”——正是从那些构成一切的“小小无物”入手——我那位最有才华的学生的天才骗局拼凑,好让那些只求被唬弄的人(而他们想必人多势众……)在公共场合上受唬弄。我急着收尾,是的;可我并不后悔为此花去的时间,尽管我年近五十七,而更有意思(或至少更“可喜”)的事情并不缺少。这有点像数学里我称之(想必就在三天前)为“例行工作”的那种活儿——做着的时候牙咬得痒,明明知道这一切不过是粮草,可同时也清楚地知道它非做不可!并非出于什么严峻的“义务”或责任,而是因为一个人不能(至少只有我不能免俗)省掉这一步,倘若我想与所探测的东西建立亲密的接触、在其中“深入”。正是靠了这样的工作,日复一日、周复一周乃至年复一年地与被我们想认识的事物“摩擦”,我们才真正“认识”它们——而唯有出自这份认识(它是一场往往艰苦、其貌不扬的工作的果实),才有\textbf{别的东西}时而迸发,就是我前天谈到的“火花”,它骤然更新我们对事物的领会,以及这引领我们深入其中的工作本身。

正是凭着这般疲倦(它还不是厌倦,而是一份能量耗去的标志),我才得以同样充分地估量:我的朋友皮埃尔在这场微妙的、编排停当的行动(“SGA~4$\frac{1}{2}$”行动,或称“SGA~4$\frac{1}{2}$—SGA~5”行动)中必定耗费了多少惊人的能量。我说不出这艺术家的工作——它比数学家的活儿不知微妙多少,动用的官能全然另一层次——在多大程度上是有意识的,抑或全然是些无意识力量的作为。而这终究是个旁支的问题,只关乎他一个人。无论如何,这能量的分流,以及这投入一项与发现之冲动截然相反的任务——掘墓人-变戏法者的任务——的强度,想必一直庞大,\dmark{}而且(这于我毫无疑问)今天依然如此\fnote{524}{(*)}{这种落在“那条公式”上的侵吞执念,用简单明白的话说,真是疯了。一方面,这种侵吞由于事物的力量,势必大体上(若非完全地)停留在象征层面:一种自我给予的满足,仿佛自己\textbf{果然就是}“那位父亲”,或者仿佛自己\textbf{果然能够让}全世界都相信这一点。它的虚构性、象征性,只要回想一下Deligne本人在那部“SGA~4$\frac{1}{2}$—SGA~5”的骗局出现之前四年发表的文章《Weil猜想I》中所写的(第278页)“Grothendieck证明了Lefschetz公式”(对 Frobenius 对应而言),就已暴露无遗。诚然,仅仅几个月之后,在1974年2月的Bourbaki报告(n\textsuperscript{o}~446)中,Serre呈献Deligne这篇文章时,还(有理地)诧异于Lefschetz公式竟无一份已发表的证明(“人们自1966年以来就等着 SGA~5 的定稿版,它本该比现有的油印讲义更令人信服”),并趁机对 SGA~4 那1583页(它们“带着一切必要的细节、以及许许多多别的”)铺陈平展上同调形式主义的文字冷嘲热讽。Serre当然不会料到,这些针对一位缺席者的讥讽不会落到聋子耳朵里。我深信,它们必定在“让人忘掉”这一天才念头的萌芽中起了作用——忘掉 SGA~4 与 SGA~5 那一“废话矿渣”等等,正如公众借Serre本人的嘴似乎正在要求的……但即便抛开Weil猜想I不谈,仅就已出版的文本(包括 SGA~5 的屠杀版,它虽遭残损却仍是一份确凿的见证……)而言,父权的抹杀根本站不住,以最基本的数学常识也说不通。

此外,如我已强调过的,那条著名公式的制定是一份\textbf{纯粹例行}的工作,只要人知道自己想得到什么。我当年花了几天才把它的主要特征提炼出来——那把我引向一些与 Artin 导子相关的精确可除性问题,而Serre对这些问题早有现成答案,并且能用 Serre-Swan 模的语言优雅地表达。那段略长(但同样例行)的工作,是把受斯塔林斯工作启发(而他这篇东西恰好刚寄到我手里)的非交换迹形式主义仔细地理顺。这一切,都是那种有Deligne那份干劲(或者只有我这份更寒碜些的干劲)的人一年之内成打处理的事情!

诚然,在Deligne笔下,“迹公式”意味着\textbf{任意维数下} Frobenius 对应的迹公式,而他费心(在“SGA~4$\frac{1}{2}$”中)把它同他所谓的$L$函数的“上同调诠释”(“Grothendieck的”,谢谢!)区分开来。他把后者仅仅当作迹公式的一个\textbf{推论}呈献。(事实上,本着1964年我那场Bourbaki报告的精神,两条公式对我曾是\textbf{同义的},是“算术”与“几何”之间同一关系的等价表达式,其一是加性的、其二是乘性的。)

由此可见,在这场围绕“那条公式”的执念背后那真正的动机(固然还只是浮面的),全然不在于上同调装备,而在于尽可能缩小、乃至彻底抹去我个人在“\textbf{那条}”猜想的证明中扮演过角色这一事实。终究是她,在我看来(直到1981年6月的反常层会议为止),作为那场在我昔日的学生围绕被否认的师父身上打成的结的\textbf{重大}纠结点显现。}。这些侵吞-掉包的反射\origpage{908}——至少在与我的工作相关方面\dmark{}——以及对待任何公然打着我的印记的其他工作\origpage{909},终于在(那场漫长的“攀登”,即“亡故师父之葬礼”的过程中)获取了对他整个存在的如此支配,以至于它们成了仿佛第二天性,侵凌并覆盖了他的原初本性,那个“孩子”的本性,那奔向着发现世界的孩子……我不止一次在看似平常的处境中(其规模远不能与我刚刚稍稍细看过的“平展上同调”这样一场“行动”相比)就近看到这些反射无声的效能,它们在和颜悦色的天真外表下以完美的自如运作。还没等你明白过来发生了什么事(倘若你真能明白过来……),他已经占有了你在欢欣中创造的东西,而先是吹上一口含蓄而阴险的轻蔑的气息,使它枯萎。(同样真切的是,有这种气息的绝非他一人,远非如此;而这种气息今天看来很成了时令之气的一部分……)

% ---------------- p.716 ----------------
但是,这气息枯萎了另一个人所创造之物的美、枯萎了他的喜悦,也就枯萎了\textbf{一切}事物之美,以及那创造之力本身——它在他身上一如在我们每个人身上,使我们得以与事物相通、深刻地认识它。诚然,这并不妨碍做出些“困难的”事、受人羡慕、忌恨与畏惧。但他身上所承载的那部作品——我近年还见过它的若干先兆——至今还没有诞生。它将在某一天诞生(如果那天终将破晓),当某种东西已然崩塌,而当那位端坐于上的主-奴,一如他那被否认的师父曾经所做的那样,成了一个\textbf{仆人}。

如今这已是满满六十页的一厚摞了(还不算一大沓脚注!),将近三星期的劳作,我全数奉献给了“平展上同调”这一场行动。它无疑是所有行动中最浩大的一场,即便不是最“肥”的(那一场将在最后检视,就在那道名字当之无愧的注“封神”里……)。我意识到,饶是如此,我还没能把它的里里外外看全。一针引一针,这场预想中的“整理”——在一场“调查”中把“事实更新”——反倒让调查重新启动,引得我去稍稍细看那部不寻常的书卷“SGA~4$\frac{1}{2}$”,而从前我只在跑动中看过它一眼。这也是个机会,让我重看一遍那本阴郁的 SGA~5 Illusie版,而且目光比从前更清醒。如今我看出那两个盗匪之间有一丝不苟的默契:Illusie把自己整个交给Deligne调遣,好呈献一部全然\dmark{}\origpage{910}合乎那位显赫的庇护人与朋友心意的 SGA~5 版本。这番对 SGA~5 的呈献,是以低八度的音调,回应那份在电锯文本中铺展的贬损与轻蔑之气,并为其中设下的骗局提供含蓄而有效的支援。

屠杀版的导言通篇都是为了在未获提醒的读者心中造成这样的印象:这是一部“技术性离题”的书卷,叠在貌似\textbf{中心的}、而且\textbf{较早的}(!)文本“SGA~4$\frac{1}{2}$”之上。而在Illusie执笔的那些报告中,这一印象又因大量指及那部海盗文本而加强——每逢他用到他那位朋友认为宜于收入其摘要的成果,他都慷慨地指向它,哪怕在同一部 SGA~5 里、甚至早在 SGA~4 里就有“量身定制”的指涉\fnote{525}{(*)}{因此,紧支 Künneth 公式(在任意基概形之上)是(对导出范畴而言使用的)真态射基变换定理的直接推论,而该定理曾是平展上同调中第一个重大“突破”,时在1963年2月。它正以此身份列在 SGA~4 的那“废话矿渣”里——人们可不愿让Illusie指到它,既然有那部中心的文本(其用途恰恰是让人忘掉这些含糊的前身……)在向他张臂相迎……}。

我是在去年5月12日那道同名之注“屠杀”(n\textsuperscript{o}~87)以及它的子注中,逐渐发现一场不折不扣的屠戮的现实的。在那组注里,我终于着手对那场拆解做一番详尽的描述(虽说还不是穷尽的)——它在过去两周间渐渐显现在我眼前。由于我当时不曾像近三星期这样,把所谓“SGA~4$\frac{1}{2}$”里围绕“那条”公式精心设置的骗术一一拆开细看,去年我还不曾看清这个侧面:在 SGA~5 的Illusie版的整体呈献中,竟有这样一丝不苟的串通。为给“平展上同调”行动(别名“SGA~4$\frac{1}{2}$—SGA~5”)收尾,我还得说明这场串通在“那条”(不动点)公式的呈献中是何以表现的,给几个细节。

% ---------------- p.718 ----------------
我此前已(在子注“好撒玛利亚人”,n\textsuperscript{o}~169$_2$ 中)指出:Illusie在他的导言里帮着朋友唱和,好让人以为 SGA~5 的出版取决于那条永恒的Lefschetz--Verdier公式的证明。(该证明自1964年起就可获得,而我当然早已在口头讨论班里把它发展出来了;而Illusie虽在1965年承担了整理工作,十二年来却始终不认为值得守他的约……)\dmark{}我同样提醒:去年(在那条已引的注“屠杀”,n\textsuperscript{o}~87 里)我就已经发现了\origpage{911}原始讨论班报告~XI 的若干坎坷。那场报告与紧接其后的报告~XII 不可分割——后者发展了我那版(直到1981年为止所知最好的)一维Lefschetz公式——而它在Illusie版中完全消失了。据Illusie的导言所说,那场报告想必只是“Grothendieck的交换迹理论”(对“\textbf{非}交换迹”而言是天赐的笔误!)“推广了斯塔林斯的(非交换迹)理论”,而且是(同样天赐地)在一次搬家中消失的(!!)。事实上,那场报告发展的是必不可少的代数预备,为的是在下一场报告中描述局部项;在那里,我发展出局部项计算(或不如说\textbf{定义})的一般方法(经由一条“Nielsen--Wecken”式的公式\fnote{526}{(*)}{这条公式被Deligne侵吞了(丝毫不提我个人),连同那种把Nielsen--Wecken公式过渡到常系数(即“寻常的”)系数层的不动点公式的途径——而可构系数层则是任意的。关于这一点,见子注“真正的数学……”(n\textsuperscript{o}~169$_5$,第883--884页)。于是(身价使然……)同一个Deligne小心翼翼地丝毫不提那场“较晚的”讨论班 SGA~5 的报告~XII——那里“Nielsen--Wecken”这个名字正写在报告的标题里(“代数几何中的Nielsen--Wecken公式与Lefschetz公式”)。}),及其在一维情形下的显式应用(若我没记错,是靠了 Serre-Swan 模)。总之,Illusie用一场“新的”报告~IIIB 来“顶替”那场消失的原始报告~XI,并应景地称之为“局部项计算”(这想必也是,如无差错而且巧得出奇,正是被抹掉的那场报告的标题!),\textbf{而他自居其父}。就此他一石二鸟。一方面,这是一项\textbf{残毁}之举,初看仿佛无谓,却以这种\textbf{生硬的割裂}四处播撒混乱\fnote{527}{(**)}{这种残毁与混乱,在我的昔日学生Illusie奉我另一位昔日学生Deligne之命播撒的众多残毁之中,使后者得以居高临下地谈论 SGA~5 的“混乱状态”(“虽说严谨”,因为人到底是好牌手……),而“SGA~4$\frac{1}{2}$”(它再早不过!)据说应当对此“加以补救”……而这一切,都在那群忠诚信徒的慈祥目光下。恭喜!}——把一场报告从它的自然脉络中硬拽出来,在原地留下一个张着的窟窿,只为图个乐子把它塞到别处。这或许是绵密细致的Illusie强加于那部曾经辉煌的讨论班(而他自感骤然成了它的绝对主人……)的一切残毁之中,事后在我看来最暴力、最赤裸裸地招摇的一桩:我能白屠,我就屠——且带着与我的良好教养相称的细腻。恭喜了,Illusie,为这一类活儿,那可不是你跟我学的,而是跟另一个人学的,你把他奉为楷模与导师……

\origpage{912}\dmark{}这是一。而这第二下,同一块石头砸得漂亮,Illusie竟能\textbf{抹掉}我在1965年提炼出的那条不动点公式的\textbf{父权},同时又(尤其是)把\textbf{那条公式本身也抹掉}。自1965/66年起,它就是“\textbf{那条}”\textbf{一维情形下的好的不动点公式},远比Verdier前一年在伍兹霍尔发展的那条更一般(否则我也不必劳那个神),更遑论Deligne那部著名的“报告”里的那条(后者只限于 Frobenius 对应,却步步照搬我在一般情形下提炼出的证明)。直到几年前(晚了将近二十年)它才在阿利贝尔的论文\fnote{528}{(*)}{这篇论文是在Verdier的指导下准备的(没错,永远是那个Verdier),1981或1982年在蒙彼利埃通过(我手头没有确切引证)。它是十年艰苦劳作、显然颇为凄清的结局……}中得到改进,头一次处理了系数层上任意的上同调对应。

% ---------------- p.719 ----------------
Illusie设法把文本编排得使有关公式几乎无处可寻:在两场相互撕裂的报告(带官样文章的 IIIB(sic)与 XII)的技术岩浆中,没有任何东西(在这一篇或那一篇的引言里,也没有在别处)把读者的注意力引向这个构成两场报告之整体的核心成果、引向整个讨论班最重要的成果之一\fnote{529}{(**)}{技术上,正是\textbf{那条}关键公式(“不可约情形”)使得证明那条著名的“$L$函数公式”成为可能,后者等价于(任意维数下)Frobenius 对应的迹公式。这条公式的关键角色,早已由 SGA~5 这个讨论班的名字本身作证(那个名字在那部所谓“较早的”“SGA~4$\frac{1}{2}$”文本里从不曾被提及):“$\ell$-adic 上同调与$L$函数”。}!我甚至承认,我没能绝对确凿地弄清这条公式究竟在不在 SGA~5 里。鉴于文本故意造成的混乱,又鉴于我远离这个题材,我得花上几个小时甚至几天才能找着。真正让我困惑的,是全文丝毫不提 Serre-Swan 模——而(如果我的记忆没骗我)正是它们赋予了我提炼出的那条公式以优雅与概念上的简洁\fnote{530}{(***)}{很可能,在这里以及下一句里,我把Euler--Poincaré公式(见于报告~X)的结构与(报告~XII)那条Lefschetz公式的结构混为一谈了。在Bucur的报告(复述我的口头报告)中所呈现的那种形态的Euler--Poincaré公式里,Serre-Swan 模确实是显式登场的。}。而正是为了这条公式的需要,Serre才做了几场关于结合于 Artin 导子的Galois模的精彩报告——那些报告本来当然应该收入那部出版的讨论班,而它们最终却付诸流水了(连同原始讨论班的另外五六捆\origpage{913}报告——这对于Illusie、Deligne一干人倒不算什么……)。有可能,所说的这条不动点公式就是报告~XII(第431页)中的公式(6.3.1)。匆匆一瞥,无法把它与那几十条被密密麻麻编了号的其他公式区分开来,而它正被埋在其中。显然那位执笔者(Bucur)被那活儿压垮了——而那位才华出众的编辑(sic)Illusie,十五年来惯于做那些清晰完美的文书活儿,绝不会为他的朋友Bucur那些\textbf{正好合他心意}的笨拙\fnote{531}{(*)}{Illusie所作 SGA~5 屠杀版《引言》的最后几行,假装“向1976年死于癌症的 I.~Bucur致敬”——比屠杀版早一年。我不知道其中是否有因果关系——我对Bucur为人的诚实与忠贞毫无怀疑,他是绝不会让像这样一版东西轻易放过而不至少告知我的。无论如何,这场死后致敬所嵌入的那场行动的精神,赋予这致敬一种可疑的滋味。“那是只凭言辞自我交代,而本有一种更符合 Ionel Bucur的意愿与耿直、更能纪念他的方式,那就是缓和它的笨拙,而不是毫无愧色地利用它们。”}抬一抬小指头。恰恰相反,他还设法加剧混乱:使那条本已无处寻觅的关键公式,益发无法与Lefschetz--Verdier公式的那一条区分开来,也无法与它在“报告”中的特殊情形区分开来。在那篇著名的报告~IIIB(sic)的引言里,那位临时凑数的“父亲”Illusie写道:

\begin{quote}
“本报告~IIIB 的第二部分,\textbf{性质上技术性得多}[所以千万别去读它!],是\textbf{受启发}[!]于Grothendieck为建立曲线上某些上同调对应的Lefschetz公式所用的方法[我们可千万别去查是哪些!](见 XII[唯有聪明绝顶的人才知道去哪儿找“那条”公式!]以及(SGA~4$\frac{1}{2}$ 报告)[读者在那里不难找到那条公式,并得知它的\textbf{真正的父亲}是谁……])。”(着重号是我加的。)
\end{quote}

再往下,在同一篇引文中,人们说(即Illusie,不言自明)运用第5号\fnote{532}{(**)}{这回谈的是迹,而且是\textbf{非}交换迹——那些笔误-揶揄,是专为那位亡者保留的,至少只要他还不曾在场反驳……}\dmark{}“技术”:

\begin{quote}
“用以在第6节定义环不必交换的模复形之间上同调对\origpage{914}应的“Lefschetz--Verdier局部项”。”
\end{quote}

而偷偷赐给这些“局部项”的名字,正是我1965年为写那条\textbf{显式的}(“Lefschetz-Grothendieck”)公式而引入的——为的恰恰是不必指及一般Lefschetz--Verdier公式的局部项——这名字显然是为了维持那场由Deligne蓄意制造并维护的混淆,好像所说的那条显式公式在技术上\textbf{依赖于}Lefschetz--Verdier公式似的。再往下几行,我们更是惊喜地得知:“由Grothendieck在(XII 4.5)的Lefschetz公式中定义的局部项”\fnote{533}{(*)}{(5月12日)我被这个关于“我的”公式的异乎寻常的精确(XII 4-5)搞得莫名其妙,刚去查了所引的条目。我发现一条“猜想~4.5”(第415页),它似乎关乎定义局部项的可能性。大家早该料到,这位千金难买的先生又要给我们端出一条他的猜想来了,而不是一个真正的定义……}(偏偏不提这些正是刚刚被慷慨赐名为“Lefschetz--Verdier局部项”的那些)“确确实实就是Lefschetz--Verdier的局部项”(不过这回是\textbf{另一种}意思,当然了:那条\textbf{一般的}、“非显式的”、所谓Lefschetz--Verdier公式的那些)。

% ---------------- p.720 ----------------
要说在污泥里钓鱼的技艺,循着一种我太过熟悉的风格,这真是到家了!

卷首引言里也是同样的混淆手法,其中写道(第 VI 页第5行):

\begin{quote}
“对Lefschetz公式的应用见于报告 XII 与 IIIB。”(着重号是我加的),
\end{quote}

总而言之,为的是让读者毫无希望地迷失其中,连去找、连去试着去找都不必——去找\textbf{那条}一维情形下已知的唯一显式Lefschetz公式(至少直到1981年),那位(既不是Illusie、也不是他的主子Deligne,而是)亡故的前“主持人”(sic)——被他的那位“编辑”-掘墓人Illusie乐呵呵地屠杀了的讨论班的主持人,理所当然地不被提名——的公式\fnote{534}{(**)}{尽管 SGA~5 讨论班的一切本质成果,除了Lefschetz--Verdier公式与 Serre-Swan 模理论(后者不在屠杀版中)之外都出自我手,Illusie却把文本编排得使这些成果中\textbf{没有}一项(不止是那条所谓“Lefschetz”公式失落于某场报告 XII 之中的某处……)显出我卑微的本人与此有过任何干系。由此,他在那场把我这个人从 SGA 中排挤出去的行动中扮演了头等重要的角色——那行动由他的朋友Deligne长期筹备,并在注“葬礼队列——‘为科学尽职’”(n\textsuperscript{o}~175)中找到它的终局。(另见子注“驱逐(2)”,n\textsuperscript{o}~169$_1$。)}。

在原始讨论班中,那场被抹掉并改称为 IIIB(带着一位\dmark{}崭新的父亲)的报告~XI,本插在\origpage{915}六场报告\textbf{VIII至XIII}的序列里,那些报告围绕两个紧密相连的主题——Euler--Poincaré显式公式与Lefschetz公式——并以我在创办该讨论班之际提炼出的共同方法,本着同一精神展开。在那部分讨论班里,如同在其他部分,有一种论旨与视景的显而易见的统一。它被我那位昔日学生一丝不苟地屠戮了,而他正是利用了他作为一艘被他本人和我的其他上同调学生(以对师父的身后谢意)弄沉的讨论班的“编辑”(sic)这一角色。带着堪与缜密的Illusie相称的规律性,六场当中隔一场缺一场——即报告~IX、XI、XIII 三场——从屠杀版中消失了。报告~IX 出自Serre,呈献 Serre-Swan 模理论——眼看事态的走向,Serre宁可收回他的赌注,亲自操心把他那场漂亮的报告交到所有人手中。报告~XIII,据那位“编辑”在卷首引言中的解释,是多余的——显然那位不被提名的主持人不会数到十三——于是就给删掉了!报告~XI,如我们所见,凭一手漂亮的戏法,变成了报告~IIIB,几乎是报告~III 的附录(瞧瞧,瞧瞧——机缘真是巧合得很……),而报告~III 原称“Lefschetz--Verdier公式”,为着混淆的缘故被改称为干巴巴的“Lefschetz公式”。总之,这次“搬场”并非偶然——它始终朝着同一个方向,朝着Deligne-Illusie这对完美搭档从不松懈地维持的那个混淆:一边是Lefschetz--Verdier公式(那条“猜想性的”、“局部项没算出来”的,但终究还是靠Deligne与Illusie的合力得到证明了……),另一边是另一条公式,一条显式的,它必须严格保持隐秘,被小心翼翼地淹没在一大堆编了号、还带四位小数的公式、含沙射影的暗示与精心计算的含糊之中。再次恭喜,亲爱的昔日学生!于是,那场题为“平展上同调中的Euler--Poincaré公式”的报告~X,被掐去了前一场与后一场,悲惨地悬在虚空中\fnote{535}{(*)}{在没有反证的情况下,读者自会猜到:这条著名的所谓“Euler--Poincaré公式”当归于两位以他们的名字为它命名的杰出几何学家。试与上一条脚注比较。}。漂亮活儿,你的时间倒没白费……

\subsubsection*{(f) 恭贺——或新风格}

\tocsubsub{(f) 恭贺——或新风格}

\notehead{169$_9$}{}

(3月22日与4月29日)我想再回到Lefschetz--Verdier公式与那条\textbf{隐秘的}、\dmark{}\textbf{无处寻觅的}公式之间被蓄意维持的混淆。我刚在 SGA~5 里发现了一份相当详尽的\origpage{916}“术语索引”——人要么细心,要么不细心!出于好奇,我去翻“Lefschetz”词条,看我“那条”公式会不会在里面……唯一的指涉是“Lefschetz--Verdier公式(报告~III)”——而那场报告也已(如我们所见)被改称为“Lefschetz公式”了。于是读者得到确切的通知:在这个书卷里不存在(至少不在这里)别的“Lefschetz”公式,只有那条所谓“Lefschetz--Verdier公式”的(正是他刚从别处得知的那条:它是猜想性的,等等;SGA~5 死死依赖着它、终生依赖着它;而“SGA~4$\frac{1}{2}$”如其名义所示在此救急……)。是啊,漂亮活儿!

% ---------------- p.721 ----------------
我继续检视我昔日的学生Illusie——在我另一位昔日学生Deligne的麾下——的种种业绩。我接着引用那部屠杀版导言的引文\fnote{536}{(*)}{引文的开头见前一条子注“变戏法者——或飞走的公式”(n\textsuperscript{o}~169$_8$)。},就在“那条”Lefschetz--Verdier公式(永远是同一条)凭着数学的变戏法之术骤然倍增为“\textbf{若干}Lefschetz公式”的地方,可谁也没能说出是哪些。他接着写道(第 VI 页第6行):

\begin{quote}
“\textbf{报告 XII 的迹公式}[人们满心以为,绝不会有任何读者动念去把它找出来……]\textbf{是独立于报告 III 的一般公式而得到证明的},但人们表明((III B 6))其中出现的局部项确确实实就是一般公式的那些,\textbf{并且后者蕴含前者}。”(着重号是我加的。)
\end{quote}

两手空空,两兜干净——攻不破的Illusie,与那位变戏法的总领班同样攻不破!在先后追随了那一整群一味朝着同一方向的障眼术含混之后,我刚注意到:就在这里,在一个直到此刻都逃过我(一如它逃过这篇四页多导言的任何其他读者)的毫不起眼的迂回\fnote{537}{(**)}{Zoghman~Mebkhout是一位细心却来得稍晚的读者,他告诉我,他自己也曾受骗,深信那条不动点显式公式(对任意维数的 Frobenius,或对一维情形的一般对应关系)确确实实依赖于那条一般的(非显式的)Lefschetz--Verdier公式。由此可见,Illusie那句“暂停!”式的说法,逃过了他的注意一如逃过我的——而这正是所要的效果……

由于我1974年的那场Bourbaki报告,是在人们把一条一维情形下的\textbf{显式}公式明确写出之\textbf{前}写成的,呈献的是“带系数的”函数$L$的公式(可构$\ell$-adic 层中的系数,或与此等价的、Frobenius 对应在这类层中的不动点显式公式)。那时我以为,一维情形下 Frobenius 的显式公式的证明,会是一般Lefschetz--Verdier公式的一个推论——“只差把局部项写出来了”。于是,由于预支了一份尚待完成的工作(而且那活儿照例该归Verdier),我在这场Bourbaki报告里把这条\textbf{显式}公式命名为“Lefschetz--Verdier定理”。后来,无论是Verdier的“伍兹霍尔”证明,还是我的那条覆盖远为一般情形的证明,都不诉诸一般的Lefschetz--Verdier公式。情形对所有 SGA~5 的听众都是再清楚不过的,至少如此。但对于那些只知道我的Bourbaki报告而不知 SGA~5 的人(而 SGA~5 一直被扣到1977年),那里有一个误会,而Deligne与Illusie沆瀣一气,把它彻底利用来搭建那场“SGA~4$\frac{1}{2}$—SGA~5”的骗局(那针脚粗得一目了然)。

况且,就 SGA~5 对那部名字骗人的海盗文本的“逻辑依赖”这一骗术而言,它无论如何也站不住,即便那条显式公式当真依赖于那条“猜想性的”Lefschetz--Verdier公式。因为,正如Deligne本人在那段著名的“方法A”中顺便指出的(对于一个请求宽免的读者——见“真正的数学……”,n\textsuperscript{o}~169$_5$ 第884页),那位不被提名的先生的那些“容易的约化”,把事情归结到一维情形,而在那里“证明的各配料本来就已经齐备”。

所有这些骗术,只要端给一位或睡着、或性急、或正求之不得地要被唬弄的读者,就能成功。对于一位细心而挑剔的读者,整套巧妙的编排就显出它的本来面目:一场不知羞耻的诈骗。但看来,我是八年来这场骗局在数学市场上出现以来,第一位细心而挑剔的读者……}里,

\dmark{}人们半明半暗地说:报告~XII 的某条迹公式(读者只管尽力\origpage{917}去找是哪一条!)是\textbf{独立于}“报告~III 的一般公式”而获得证明的(而后者,为着合宜,同样不配有个名字,合乎所谓“刻意含糊”法……)——以便在同一次呼气、同一句话里接着(仿佛是为了“追回”一句合乎谨慎规则的“暂停!”式说法……)来一个“但人们表明……”。这个“但”指的是那份“柏拉图的”补篇:十二年来,谁也没有为它操过心,首先是Illusie与Verdier,也就是说:“我的”那些局部项——对不起,我是想说“其中出现的”那些(在这条报告~XII 的迹公式里,而那位作者永远不得被明白地提名\fnote{538}{(*)}{对于 SGA~5 的读者,看来该是Illusie——那场关于“局部项”的辉煌报告~IIIB——做那条从未被提名的公式的谦逊父亲。而对于那部自称“SGA~4$\frac{1}{2}$”的书卷的读者,他因此不曾听说过别的公式、只听说过“报告”里的那一条,父亲显然就是那部书卷辉煌的作者。而对于同时读这两者(倘若有人这么做)的读者,他就只好掷硬币猜正面,或者干脆把舌头交给猫……})

% ---------------- p.722 ----------------
\dmark{}——这些项就是那条永恒的“一般公式”的那些项——而关于给\origpage{918}各条公式的名字、以及去哪儿找它们的那种含糊,突然让位于一种堪与缜密的Illusie相称的范本式的精确:这一份“额外赏赐”的证明,却在 III B 6 里——倘若哪位读者想确证它确确实实在那里,他找起来毫不费力,那一条!

那么,对这恒等式为何突然有了兴致,既然整整十一年来,SGA~5 这个讨论班的整个命运,Illusie(一如我的其他上同调学生)都全然无动于衷?为的是能在同一句话里再接再厉地辉煌地(这事儿要么是天意,要么我就再也不知怎么说了)使“那条一般公式”(Lefschetz--Verdier的,就不点名了)\textbf{蕴含}“报告~XII 的那一条”(那位亡故者同样不被提名)。

这真是一手真正天才的戏法!我那位辉煌的昔日学生大汗淋漓(包括靠着数学上的活儿,一点不假),才挣得这句貌不惊人的收尾话的辉煌成果——而在Deligne眼里、以及在他那仆从的眼里,这一句却是性命攸关:Lefschetz--Verdier公式“蕴含”“报告~XII 的那条”(而人们刚刚才说过它是被独立证明的,可为了满足无意识那些全然象征性的快意,这又算得了什么!)。

而这个“\textbf{蕴含}”,就数学而言,性质十分特别——而我敢打赌,我是世上唯一的数学家,除了那个噱头的辉煌发明者(也许还有他的主子Deligne),能品出它的滋味。然而要理解它,并不需要是专家,甚至不必是数学家。两条公式,那条“一般的”(别名Lefschetz--Verdier)与“报告~XII 的那一条”(别名那位不被提名的亡者),分别写成这样的形式

\begin{quote}
$T = L$, \qquad $T = L'$,
\end{quote}

其中项$T$(交错迹和)在两条公式中是同一个,而项$L$、$L'$(局部项之和)则是临时定义的(其一由Verdier本着Lefschetz的精神,另一由那位亡故者本着Nielsen--Wecken-Grothendieck的精神)。

十一年后,Illusie(他的文书热忱一听见主子的招呼便骤然苏醒)做了一场突如其来的努力,本该用在更有价值的事业上,为的是去证明……
% ---------------- p.723 ----------------
$L = L'$\textbf{直接地}(?)(以及对各个局部项逐个地),为的是能说公式$T = L$“蕴含”$T = L'$(并由此隐含地,那部该屠杀的讨论班的公式$T = L'$——对$L$函数理论至关重要的公式——“依赖于”公式$T = L$,而后者在Deligne与他那救命的“SGA~4$\frac{1}{2}$”出现之前始终“属于猜想”——sic……)。

\dmark{}局势对于一个略知底细的人就变得更滑稽了,因为他意识到:若不是$L'$中那些局部项(那位不被提名的亡者的那些)那套稀奇古怪的定义直接被证明公式$T = L'$的运思本身“吹送”而来,\origpage{919}世上任何人都不会想出它来。说实话,我可以说:我在\textbf{甚至还没}定义右端$L'$及其各个局部项之\textbf{前},就找到了公式$T = L'$的一项“证明”——那些局部项是从证明里“冒”出来的,不多不少\fnote{539}{(*)}{我说明一下(这事当然不言而喻):在一切可以想见的应用中(不止于$L$函数公式,那仅关乎 Frobenius 对应),真正\textbf{贴切}的公式确是那条\textbf{显式}公式$T = L'$。从实用的角度看、并且就一维现象而言,Lefschetz--Verdier公式$T = L$只有一种历史(或启发)的兴趣,而Illusie的结果$L = L'$(或者更确切地说,出现在$L$中的那两类局部项与出现在$L'$中的那两类,是同样的)就更是如此了(至少在新命令到来之前)。这些原是再明白不过的事,而那两个同谋却设法(而且,照眼前的时势,还成功地)把它们搅浑。这让人不禁思索:我们眼前这场狂热的科学生产,究竟能有什么意义——而那一切都发生在一些与二十五年以来我们在几何与算术之关系上的认识所取得的关键进步息息相关的问题上,对普通数学常识如此粗暴的歪曲,竟丝毫未引起任何人的注意……}。

恭喜,这是第三次了,Illusie;也恭喜你,Deligne,你给他当了楷模。你们两位合力开创了数学中一种\textbf{新风格}的先河。而这风格已经收了门徒。

如今它已然成了“1980 风格”,显然前程似锦\fnote{540}{(**)}{关于这方面触目的例子,见“黑手党”一注(n\textsuperscript{o}~171$_{2}$)中若干“1980 风格”的样品,出自我们那几位大作者手笔:Brylinski、Kashiwara、Beilinson、Bernstein。显然,什么希望都是允许的!

(5月12日)至于这“新风格”另一些偶然的信徒——他们是在一位从未被提名的默默无闻的身后弟子的工作尾迹里崭露头角的——如今我可以补上 Malgrange、Laumon、Katz。(见“掠夺的空白授权”一注,n\textsuperscript{o}~171$_{4}$。)}。这便是变戏法-风格,别名“掘墓人的风格”,其全部艺术在于不断地\textbf{欺骗读者};不只在主要观念的\textbf{父权}上骗他,而且(顺带着)在它们的师承与相互关系上、在每一项的份量上、在什么是本质、什么是枝节上骗他——为的是让那该受抬举的人受抬举,让那该遭贬损的人遭贬损(或者以一个无精打采的手势、在一句若无其事的话的转折处把他埋葬……);而且\textbf{尤其是},为了一份刺激的\textbf{权力}之感:牵着读者的鼻子走,随心所欲地\dmark{}\origpage{920}炮制并拆解他的科学的历史,并决定那些声称要呈献的数学之物“是”什么、又不是什么。这便是那永远“\textbf{称王}”的艺术:轻轻牵动无形的(?)丝线,而永不、永不纡尊降贵去服侍。而这一切,为的是永远、彻底地享有“暂停!”的豁免:即便有哪位读者比别的更机灵,竟会亲自去看一看,竟会别出心裁地想动用(天知道……)他自己的眼光与官能(这少见,但说到底也不是不可能……),他也永远不能当场捉住你说了某句一经按字面理解、又无含糊或双关可逃的话,那确确实实、无可挽回地是\textbf{假的}。

这种句法上的艺术之艺术,在于那样一种风格条款,它看似一道难题,然而……随着那场记忆古怪的“反常层会议”(那是一场令人叹为观止的“SGA~4$\frac{1}{2}$—SGA~5”行动之后不过四年),人们得以看出这种新的、无害的技艺能走多远——在那场妙不可言的行动的变戏法绝技的种种施展之后,在对一部革新性作品的掉包中、在对那位曾长期孕育并在孤独中耐心成熟这部作品的人的肆意掠夺中……

向师父和学生、向Deligne和Illusie脱帽致敬!艺术家的活儿!你们二位当之无愧地赢得了教团全体的一致感谢。

% ---------------- p.724 ----------------
\subsection{(2) 分赃(“对偶——晶体”)}



\subsubsection*{a.~最后一份——或聋子的耳朵}

\tocsubsub{a.~最后一份——或聋子的耳朵}

\notehead{170(i)}{}

(2月28日)我这就说到围绕我数学著作的“四种运算”中的第三种(第四种且待下一条注,那一条将抹杀Zoghman~Mebkhout的作品)。

\begin{quote}
III “对偶——晶体”行动(或称:“美丽的残骸……”)。
\end{quote}

就我如今看事情的样子,这大体上是对我那部分著作中关于上同调、而尚未被\textbf{Pierre Deligne}以事实(或象征)方式侵吞的部分的一次\textbf{分赃}\fnote{541}{(*)}{(5月1日)尽管如此,仍该把\textbf{凝聚}语境中的对偶形式主义单独挑出来:它(与一个后来显得匆促的印象相反)显然还不曾被我的上同调学生中的任何一位、也不曾被我知识范围内的任何其他人所侵吞。诚然,唯一陈述我在这个主题上的大部分想法与结果的参考文献,是 R.~Hartshorne的《留数与对偶》(Residues and Duality),这使人可以指到它,而无须在任何一个时刻说出一个不受欢迎的名字……}。显然,他给自己留了最大的一份,带着 Motif 与平展上同调,更具体地说,是那套$\ell$-adic 的上同调工具。其余部分的分割(*),则在我的另两位上同调学生\origpage{921}\dmark{}\textbf{J.-L. Verdier}与\textbf{P. Berthelot}\fnote{542}{(*)}{(5月1日)自从写就以来,看来该添上“第四个盗匪”,即Neantro~Saavedra~Rivano,他经由那些应景地受洗为“坦纳克式”的范畴,侵吞 Motif Galois群的哲学。但他只是替Deligne当“稻草人父亲”,而后者在十年后“收回”了父权。事情的始末,见后续诸注“棺材上的第六颗钉子”,n\textsuperscript{o}~176$_1$至176$_7$。}之间进行。那默契是何时、怎样形成的,我说不上来;它似乎是这样:Berthelot得全部晶体上同调,而其余归Verdier——后者基本上并吞了一切围绕对偶瑜伽\fnote{543}{(**)}{见上一页的脚注。}的东西,以及作为其代数前提的导出范畴与三角范畴瑜伽。

关于Berthelot参与分我这具遗骸,我只有一个事实可依据,虽说那事实分量不轻。去年,我在“共同继承人……”(n\textsuperscript{o}~91)那条注的反思中偶然撞上它,并为它写了一篇短小的子注(n\textsuperscript{o}~911)。那就是我在那里引述的Berthelot的那篇综述文章\fnote{544}{(***)}{刚性几何与特征~$p$代数簇的上同调,Pierre Berthelot,载“$p$-adic 分析及其应用”CIRM 吕米尼学术会议(9月6--10日)。},它在一场题为“$p$-adic 分析及其应用”、于1982年9月举行的吕米尼学术会议上,为 Dwork-Monsky-Washnitzer 上同调与晶体上同调的一种“综合”(如他所说)呈献了主要的想法。在那篇引言(b)部分里,他给晶体上同调写了一小段历史,其气量狭隘,毕竟与我对晶体瑜伽所抱的远为宏阔的视野全然不合\fnote{545}{(****)}{关于这一点,见随后那篇子注“聋子的耳朵”(n\textsuperscript{o}~170(i)bis)。}。

我的名字在文章正文与文献目录中都缺席。至于若干评注与细节,读者可参看所引的那篇子注,这里无须重复。我只想补充一句:一旦把我这个人从图景中清除,那晶体上同调之父,不是别人,正是他Berthelot一个人,甚至不消他费心把这话明白说出来——某种侵吞的风格显然已经收了门徒……事实上,正是他那篇论文——他当年照着我的若干出发想法、与我一起筹备的——构成了这个晶体主题上第一部出版的作品(除了我自己\origpage{922}就若干出发想法\dmark{}所做的那份极其简略的勾勒之外\fnote{546}{(*)}{这些想法的唯一一份已发表的勾勒,出自我1966年11月与12月在 IHES 所做的五场报告,由 I.~Coates 与 0.~Jussila 执笔,即《晶体与概形的 De Rham 上同调》(Crystals and the De Rham Cohomology of Schemes),载《论概形上同调的十篇报告》(Dix exposés sur la Cohomologie des Schémas,North Holland, Amsterdam 1968)pp.~306--358。所有本质的出发想法都在那里有了勾勒,包括引入 Monsky-Washnitzer 式局部“加厚”的必要性(pp.~355--356)。})。他的论文为那个我曾向他提出的纲领的至少第一部分(170(i)bis)呈献了一项规模宏大的奠基工作。

这篇可纪念的“综述”写于1982年,即“反常层会议”(1982年6月吕米尼)之后一年;那场会议将在“第四行动”中谈到。我不曾费心去重读我手头Berthelot的抽印本,以便了解这场对我的葬礼的参与,究竟代表他与我和我的作品之关系上的一个晚期转折,还是一种更古老的态度之延续。若是前一种情形,那在某种意义上大可打赌:这个转折是对前一年随那场会议而来的、科学伦理总体败坏中那场骤然且疯狂的自我升级的一种回应。为此我提醒一句:同一年1982年,也因那部“值得纪念的书卷”LN~900 把 Motif 发掘出来\fnote{547}{(**)}{见“沉默”(n\textsuperscript{o}~168),尤其是“……与发掘”(n\textsuperscript{o}~168(iii))。}而出名——在那里,为这场行动买单的不再是一个无名的“当值无名之辈”(如在那场辉煌的会议上),而是一位“亡者”,他的名字毕竟还留在人们的记忆里(哪怕是出于遗憾……)。前一年的那场行动,已经足够清楚地表明:再也不需要任何克制——而“Motif 行动”事实上与“晶体行动”以及此前的所有行动一样,连一点涟漪也不曾激起……

\notehead{170(i)bis}{}

(170(i)bis)(2月28日与4月30日)\fnote{548}{(***)}{本子注源出上一条注“最后一份”的一条脚注。(****)(5月12日)事实上,那已是在1966年,见上文脚注(*)。}我在此所谓晶体理论(特征~$p>0$)的“第一部分”,指的是特征~$p$基概形上\textbf{真而光滑}的概形、以常系数(或“扭曲的常系数”)为项的晶体上同调。那时只需使用我在六十年代末\fnote{549}{(****)}{(5月12日)事实上,那已是在1966年,见上文脚注(*)。}(暂时地)引入的“寻常”或“无穷小”晶体 site 就够了。事实上,与Berthelot乐于赋予“晶体上同调”一词的狭隘含义相反,它对我来说从一开始就有一种远为\dmark{}宏大的含义,而我对谁都不曾\origpage{923}隐瞒,我的学生们显然已经把它忘了——为了在十或十五年之后又“重新发明”了它的一小块……

% ---------------- p.726 ----------------
一方面,我的晶体想法从一开始就绝不限于给定特征~$p>0$~的概形这一情形。我最早的晶体反思——在“带除幂的加厚”这个新念头来到我心中之前——处在\textbf{特征零}的概形上,那里除幂自动在场(正因此,它们往往不被人留意……)。这个研究方向由于Zoghman~Mebkhout的想法而重获生机,其自然的终点,将是在(先从)特征零概形上、针对“De Rham-Mebkhout 晶体系数”的六种运算形式主义——那个形式主义,我在“墓上的旋律——或自负”(n\textsuperscript{o}~167)那条注里已经提到过。早在六十年代,我就预见过一种不分特征的晶体上同调,其形态是在绝对基$\mathbb{Z}$上(例如)有限型概形的语境中的一套晶体“六种运算”形式主义。它本该涵盖$\mathbb{F}_p$($p$元素域)上有限型概形的“寻常”晶体理论(那个理论当年还在寻找——而今也还在寻找)。我深信:忘却并埋葬了亡故师父的这一视景(它却再简单、再富于启发性不过),正是晶体理论那可悲的停滞之因——而在它起步时那股蓬蓬勃勃的生气之后,已将近二十年。

\origpage{924}另一方面,且回到 Monsky-Washnitzer 的进路——它当年曾为把我“触发”到晶体上同调起了作用——我从一开始就清楚:有必要引入(为着一门不只适用于真而光滑的概形的理论)一个比“无穷小”site 更大的晶体 site,在其中,所设想的“加厚”将是适当的\textbf{拓扑}代数(带除幂理想)的谱,或许是 Monsky-Washnitzer 所用的那些(剥去了诸如光滑性这类无用假设的)(*)。提炼出“\textbf{对的}site”\dmark{}\fnote{550}{(*)}{如我在先前一条脚注中所说明的(见第922页),我在那篇关于晶体瑜伽的、唯一一篇已发表报告(1966年末)中谈到的正是这类 Monsky-Washnitzer 式加厚。从那时起我就已清楚:特征~$p>0$~的晶体上同调,将有很大一部分在特征零的刚性解析空间上进行。我当然不曾不把这告诉一切可能与此相关的人,首先肯定是我的学生Berthelot——一旦他选定投身晶体主题。在所引的那篇文章里,循着一种我认得的、而Berthelot未曾发明的风格,人们会以为他刚刚(晚了十五年)发现了与刚性解析几何之间那条未曾疑及的关联(x)。他在那里以一个“共同推广”(Monsky-Washnitzer 理论与晶体理论的)的天才发明者的姿态出现,并大言不惭地把那推广命名为“刚性上同调”(而它很快就会如愿地叫作“Berthelot上同调”)。我还得指出,Berthelot的这项工作,是“与 Ogus 共同进行的一场反思的延续”——而同一位 Ogus,正是在同一年(1982)以 LN~900 共同作者的身份,参与了那场“Motif”骗局而崭露头角。

系统性的埋葬在Berthelot后续的一篇文章(我手头有一份预印本)《刚性上同调与 Dwork 理论:指数和的情形》(未署日期)中继续进行。对$F$-晶体这一关键概念、或紧支上同调的概念(而后者是我1963年2月在代数几何中荣幸地引入的,早了二十年……),文中毫无对亡者的指涉。这些概念既然如此自然,实在不值得为那点区区小事费神……形式概形(在离散赋值环之上)的“泛纤维”的概念,作为一刚性解析空间,则被慷慨地记在我昔日的学生Raynaud的账上。而我早在Berthelot、Raynaud、乃至任何其他人还不曾听说过“刚性解析空间”这个词之前就已知道这一概念——因为能够定义这样一条泛纤维的需要,原是我预想存在一门“刚性解析几何”的两大动机之一,而后来也正是它,成为引导Tate把这样一门几何在形式上搭起来的导线之一:他的定义必须是使得“泛纤维”的概念成为同义反复……

(x)(1985年9月)事实上,头一个预见这样一门理论之存在的是 J.~Tate,时在1959年8月。关于这一点,见注~n\textsuperscript{o}~173$^d$)(“葬礼——或自然的下坡”),尤其是第1132页的脚注。}与“\textbf{对的}系数”,是我遗赠给我的上同调学生们的纲领之一部分(以Berthelot打头)——如今看来,那完全是白费。近来(趁写《收获与播种》之便)“顺带”想了想这件事,又记起一门囊括一切特征于一体的晶体理论之必要,我甚至自问:这些拓扑代数(Monsky-Washnitzer 式的,或任何合理变体)是不是也同样太“粗糙”了(如同受限形式幂级数一样),因为它们太“远离代数的东西”,是不是该用(在适当意义上)是“平展邻域”的“加厚”来代替它们。我想在《收获与播种》之后那部分《反思》(我揣测是第3卷)里回到这些问题,连同六种运算的瑜伽与“系数问题学”的阐述,特别是“De Rham-Mebkhout”类型的晶体系数。

Mebkhout还预见过:他的 D-模哲学该为晶体理论提供一个新的视点。但他这方面的建议——特别是1978年向Berthelot提的——出自一位含糊的无名者、一位不肯悔改的Grothendieck主义者,便落进了聋子的耳朵\fnote{551}{(*)}{耳朵聋并不妨碍同一个Berthelot,在我在上一条脚注里引述的那篇文章里,满不在乎地(在第3A节末尾)指到“复簇上$DX$-模理论的一个类似物”,而人们在刚性解析框架中“暂时”还不曾掌握它。当然谈不到在此提起某位含糊无名者的名字——他曾在四五年前来向他提过一些荒唐的建议,而这就更不必提了:因为前一年的某场学术会议(将在下一条注“封神”,n\textsuperscript{o}~171 中谈到)已经为这位无名者定下了清楚的调门。可以肯定,再过几年,借着那套众所周知的、所谓“Riemann--Hilbert-Deligne”哲学的真正之父的祝福,Berthelot就要以 D-模哲学在“刚性解析上同调”(也称作,尽管他自己避免这样称呼,“Berthelot上同调”)中的天才发明者身份亮相了。由此可见,凭眼下这年景,要远行,耳朵不必有多灵……}……

% ---------------- p.727 ----------------


\subsubsection*{b.~荣耀满天飞——或暧昧}

\tocsubsub{b.~荣耀满天飞——或暧昧}

\notehead{170(ii)}{}

\origpage{925}\dmark{}(2月28日)\fnote{552}{(*)}{这条注的文字曾于5月1日(铃兰日)重新整理,并在若干处作了修正。}为把“对偶行动”定位在 J.L.~Verdier那可疑的受益上,本该先说几句对偶瑜伽的话(所谓“六种运算”的——但那名字已沉没得无影无踪),那是我从五十年代后半期起发展出来的;还有导出范畴的瑜伽,它与前者实则不可分割。我在“我的孤儿”(n\textsuperscript{o}~46,尤其是第177--178页)那条注、以及它的子注 n\textsuperscript{o}~462(第186--187页)里,已就此说得相当周详;最后(作为对Verdier在我那同调代数观点之葬礼中所扮演角色的初步反思)在“本能与时尚——或强者的法则”(n\textsuperscript{o}~48)那条注里也谈过。我看再回头去讲是没用的,只建议读者需要时去参看,然后再继续读“对偶行动”的经过\fnote{553}{(**)}{(5月12日)另见“祖先”(n\textsuperscript{o}~171(i))与“工地的巡礼——或工具与视野”(n\textsuperscript{o}~178)两注,尤其是“六种运算”与“系数”两个工地(n\textsuperscript{o}~3、4)。}。

Verdier在这场分赃行动中的态度,比他那两位朋友的态度显得更暧昧,因为他\textbf{曾}有时同时\textbf{在两个台面}上玩,而那两个台面可能看似互相矛盾。我起先自己也难以在其中认出方向,那局面看去是那样混乱。一方面,在1967年他的论文答辩之后、尤其在我1970年离开之后,他努力(由于一些我所不解的理由)\textbf{埋葬并败坏}那套他得自我的同调代数与对偶瑜伽——而他原本把最充沛的精力、在整个六十年代直到论文答辩,都用于发展这些思想并以他自己的贡献使之丰富。另一方面,至少从1976年起(在他那篇论文(sic)答辩九年之后),并在Deligne的鼓励与有效支持下,他装作\textbf{侵吞}那些出发思想的父权(只要它们没有被抵制),也侵吞\dmark{}我围绕\origpage{926}平展对偶这一主题所发展的全部方法与成果\fnote{554}{(*)}{当然,在所说的那些“其他语境”中,平展语境出发时的困难、即需要一次“突破”以给平展上同调至少一点抓手(在缺乏以奇异单形、收缩法等手段做出的著名超越构造的情况下)——这一困难并不存在。我的学生们都找到了那样的情形:那场“突破”的重活已经由别人替他们干完了——他们只消搬自己的家具来,而往往那“别人”还奉送了家具。机会一到,他们就忙着把那人埋掉,好把自以为可以侵吞的东西据为己有,并嗤笑其余……}——那些方法只需略作变通(mutatis mutandis)便能套用到一切其他语境,诸如拓扑空间或复解析空间的语境。

关于Verdier对单单是导出范畴的态度,我曾在“赊账的学位论文与全险”(n\textsuperscript{o}~81)那注里设法指出这种暧昧的意义\fnote{555}{(**)}{写这条注时,我还不知道Verdier当年是怎样露了一手的——凭着1976年他提供的那“好的参考文献”——见下文“第2步”。}。读者还将在那里找到若干素材,尤其关于他的论文(至今仍未出版)与答辩周围那些古怪的情形。隔了一年回头看,这场反思中渐渐显出的那幅图景,在我看来或许是正确的(至多还需些修订),然而毕竟是皮毛的。我很清楚:Verdier的\textbf{真正}动机,不在任何一种近乎可怜的“回报计算”的层面上,而全然是另一种性质,并且本质地牵涉他对我这个人的矛盾关系。我琢磨,即便对一个皮毛的观察者而言,他的情形也格外触目:他以为埋葬了那位曾是他师父的人,实则日复一日、直到今天所埋葬的,不是别人,正是\textbf{他自己}与他自身之中的创造力。

% ---------------- p.728 ----------------
为把“对偶行动”兜一圈,我现在要对我所知道的这场行动、以及更一般地说Verdier对我的埋葬的参与,做一次简短的回顾。

\textbf{第1步}(1966--1976)。那是在我1970年离开之后——我说不准究竟何时了——Verdier通知我:他不再打算出版他的论文。我提醒一句:那篇论文本该呈献同调代数的新奠基,以导出范畴为着眼点。在我看来,他那篇论文工作的存在理由,就是交到所有人手中,提供一部参考文本,其意义可与 Cartan-Eilenberg 的书相比,并直接顺应\dmark{}五六十年代在我的工作和我的学生们的工作的尾迹中出现的那些新需要。\origpage{927}凭后见之明,我明白:那套新的上同调语言,当年(而我甚至可以说今天也依然……)只被我的上同调学生们完全消化,而Verdier的决定,于是就等于在这幅同调代数的新图景上划了一道大杠。由此,他那“论文”——只有二十五页,仅只呈献了一些想法的令人信服的勾勒,而那些想法他自己就说不该归他——也失去了意义,严格说来成了一份“冒牌论文”。可七十年代初,在(惊讶地)得知Verdier的决定时,我正如此强烈地埋头于与我当年那些数学兴趣截然相反的事务,这些问题对我其时是无穷地遥远。我不曾起念头去细看这件事——它是在匆匆间(我可以想见)夹在关于萨克莱原子废料桶裂缝丑闻的一场公开讨论、与为编写《生存与生活》简报的一次工作会议之间听来的。而当时我更不曾想到要去反应。我头一回终于“细看”Verdier这一举动之意义、并开始隐约看出它是一场刻意破坏的,是在那条已引的注“本能与时尚——或强者的法则”(n\textsuperscript{o}~48)里,几个星期之后,即在发现那场葬礼“以其全部辉煌”之后,又在那一注更为周详深入的“赊账的学位论文与全险”(n\textsuperscript{o}~81)里重提。

回想起来,有一点很清楚:Verdier在他自作自受的那份工作上的分心——那本是他与论文评审委员会订下的“诚意契约”(见所引注 n\textsuperscript{o}~81)的一部分——至少该上溯到1968或1969年;否则他那“论文”的撰写与出版,早在1970年我离开之前就已成事。我提醒一句:我还在1950年就把论文工作的纲领交给了他,而对一位像他当年那样有天份、有干劲的研究者,这个纲领加上一份新奠基的大规模撰写,满打满算也不过三四年工夫,连熟悉入门都算在内。同样真切的是:有某种心态,就是设法预先为一份“预定”的工作支取一笔信用,事后再也不必费神去把它做出来——这样一种心态,在我看来,早在1964年之后就已经显山露水:有那所谓“Lefschetz--Verdier”公式的坎坷,后来有\origpage{928}局部紧空间的对偶(顾名思义,就叫作“Verdier的”对偶),凭着六种运算的精神\fnote{556}{(*)}{关于这种相当特别的心态,见子注“遗产——或骗术与创造”(n\textsuperscript{o}~169$_6$bis),以及去年那两篇子注(n\textsuperscript{o}~812、813)——在那条已屡次引用的“赊账的论文与全保险”里。}\dmark{}(它们至今仍不被提名)。但整个六十年代,我身为囚徒,被自己的任务、被自己通过它们不懈追求的那幅图景所困,如同阿哈那难以捉摸而又无处不在的白鲸,我万万不曾疑心:在那位曾是我在自认为“共同”的任务中的亲密伙伴的人身上,有什么东西“不对劲”——正如我也毫不曾对我的上同调学生中任何别人起过疑。而凭二十年的后见之明,我现在才惊觉:我生命中有十年(若不是十五年、二十年),我活得与周围的现实全然\textbf{错位},而这不仅在我的家庭生活里(我在那里早已有所觉察),也同样在我全身心投入的职业生涯里……

% ---------------- p.729 ----------------
但且回到“第1步”。Verdier对我这个人和我的作品的暧昧关系,无论如何在 SGA~5 讨论班1966年结束之后就已显现:我的上同调学生中没有任何一位(他也不例外)觉得这个讨论班的整理工作与自己有关\fnote{557}{(**)}{回想起来,我不由要问:1964年(那时在我的接触下,他终于投身新的上同调技术的洪流)到1970年之间,Verdier的时间到底都用在什么上面了——因为他竟不肯着手并完成任何一项撰写工作,甚至对那些他即将自居为作者的理论也是如此。至于他那份清单:有分量却没有一项做到头,见那条被屡屡引用之注的子注~n\textsuperscript{o}~811.}——而那整理始终留在那些被任务压垮的、或不屑守信的自愿者(sic)手里。显然,就在那时,我的上同调学生们整体的情形已经腐烂,而我却丝毫不曾觉察,宁愿活在一个一切都只有秩序与美的世界里……直到十八年后,我才开始投去第一道怯懦的目光,看看当年究竟发生了什么——而那些岁月(仅在一年前)在我还显得田园牧歌\fnote{558}{(***)}{尤其见《愚妄与更新》中“一个没有冲突的世界?”一节(n\textsuperscript{o}~20),那里只有节名中的问号,能让人对那“牧歌”起一点疑心。}……

我1970年离开之后,甚至在他向我宣布他那“正式的”破坏基础工作的决定之前,Verdier在那些六十年代的暧昧,已由他与他的朋友Deligne当年那些各色小型骗局的串通得到确认——而他不可能不曾察觉:\origpage{929}\dmark{}我在 Hodge I、II、III 诸文中被抹掉\fnote{559}{(*)}{在那场“权复形”的玩笑(见同名之注,n\textsuperscript{o}~83)里,我几乎以为辨出一句暗讽,以挑衅的口吻,指向我的上同调学生中某一人的最早一桩昭然若揭的骗局——即Deligne1968年那篇论谱序列退化的文章。如果说当年我没看出其中名堂,那么我最出色的学生所树立的榜样,却并没有对所有人都白费!},然后是在我那单值群讨论班的出版版 SGA~7~II 中(它以Deligne与Katz的名义呈献,后者突如其来地占了一个亡者尚且温热的位子……)。同一年(1973年),他也不曾能不读到 Mac~Pherson 的那篇文章——那里解决了一条“Deligne--Grothendieck猜想”,而他心里清楚Deligne与此毫无关系。

直到1976年,在那些默许的并吞行动中,Verdier的角色看来主要还只是被动的。但相反,由于不出版那本该构成他论文的东西(那论文曾是“赊账”给他的\fnote{560}{(**)}{见已引注 n\textsuperscript{o}~81。}),他在我离开之前,就在那场对我的同调代数观点(他曾一度把它当作自己的)及其作为一种“日常”技术在代数几何、拓扑与代数中的使用的葬礼中,扮演了关键角色。正如他的朋友Illusie和Deligne一样,就这样破坏自己双手的工作,为着埋葬那位启发了他的人的乐趣,他当之无愧地赢得了教团毫无保留的拥戴……

这番蓄意的埋葬之意,在他1975年之后对Zoghman~Mebkhout那番令人沮丧的态度中也表现得明白——当时Mebkhout装作要取法我的对偶瑜伽和导出范畴的瑜伽。关于这一点,我仍请读者去参看那些已引的、更周详的注:“我的孤儿”“本能与时尚——或强者的法则”“赊账的学位论文与全险”(n\textsuperscript{o}~46、48、81),以及“当值无名之辈与好上帝定理”一注(n\textsuperscript{o}~48$'$)\fnote{561}{(***)}{(5月1日)另见“封神”一注(n\textsuperscript{o}~171)的子注“一个视景的萌发——或闯入者”(n\textsuperscript{o}~171$_{1}$)。}。

\origpage{930}\textbf{第2步}(1976).1976年出版了Verdier那篇“值得纪念的”发表在 Astérisque 上的文章\dmark{}\fnote{562}{(*)}{J.L.~Verdier,《结合于一个闭链的同调类》,Astérisque n\textsuperscript{o}~36 (SMF) p.~101--151 (1976)。},前面已作为“一场攀登的第3幕”谈过(见“操演”一注,n\textsuperscript{o}~169)。我提醒一句:那篇五十页的文章,除去几页他自己的东西,就是把一些我在 SGA~5 中早十年发展起来的观念与技巧原文照录,而丝毫不提我个人、也不提任何一个谈过这些事的讨论班。这次出版是我一年前在“反常层会议”的尾迹中(在“好的参考文献”一注,n\textsuperscript{o}~82 里)发现的,它当时以一道全新的光,照亮了Verdier本人和我的其他上同调学生们那份把 SGA~5 讨论班(以它这个名字、并带着属于它的父权)交给数学界的不情愿究竟意味着什么。

% ---------------- p.730 ----------------
这里无须回头再谈我昨天在那条已引的注(n\textsuperscript{o}~169)里对这篇文章所做的评论。作为一个可笑的细节,我只想补一句:正是Verdier这份“工作”(sic)的手稿——他上年(1975年)曾好心把它寄给Zoghman~Mebkhout——成了后者踏入簇之上同调的芝麻开门,并成了他对那位从此以“恩人”身份出现的人物的毫无保留之敬仰的根基。而且这份敬仰还很是耐久,我相信,它只是在Zoghman在反常层会议上的那些倒霉遭遇之后,才彻底瓦解。

Deligne告诉我\fnote{563}{(**)}{见“把一个个 i 点清”一注(n\textsuperscript{o}~164),第 IV 1 部分。},说他是在“SGA~4$\frac{1}{2}$”(sic)与 SGA~5 出版之后一年(1977)——才得知Verdier那篇文章的:这与我的信念正好相反,我认为那“好的参考文献”的出版,标志着那场终于导致那场规模空前的“SGA~4$\frac{1}{2}$—SGA~5”行动的骗局攀登中最后一个本质的步伐。细想之后,我难以相信Deligne的说法。他是我认识的消息最灵的数学家之一,而且自始至终与Verdier关系密切;他不可能不曾早已知道Verdier的计划、不可能不曾收到一份预印本(而这甚至早在Mebkhout之前)、也不可能不是在1976年头几个拿到抽印本的人之一。那篇文章填补了文献中一个张着的窟窿(鉴于1966年之后 SGA~5 讨论班一直不曾出版)——这是Deligne本人向我确认过的——\origpage{931}而他也不可能不曾至少费心把它浏览一遍:对一个像他这样“在圈内”的人,那不过是至多一刻钟的事\fnote{564}{(*)}{我甚至可以想见:比数学兴趣更强得多的(而那篇文章对作为 SGA~5 听众的Deligne本来毫无可学之处),必是能第一手、白纸黑字地得知对那位亡故师父的无瑕疵抹杀——循着他自己八年前就已开创的传统!}。无论如何,这一桩明火执仗的抄袭不曾引起 SGA~5 另外六七个明明“在圈内”的昔日听众中任何一个人的反应,这本身就是一切相关者之间一场无瑕疵串通的保证。时机已经成熟:对那母讨论班 SGA~5 的正式屠杀,以及把我的平展上同调作品炸成碎片……

\textbf{第3步}(1977)。在这场1977年发生的“SGA~4$\frac{1}{2}$—SGA~5”行动中,由Deligne发起、并由Illusie踊跃参与,Verdier这一次只扮演一个跑龙套的角色:他为那部名字骗人的“SGA~4$\frac{1}{2}$”这份寒酸的小册子添了一份砖瓦——即他那篇论文(sic)的某种“第0状态”(而论文本身,则彻底消失了……),为着合宜,在沉睡了十四年之后被专门挖了出来!在整部书卷里,无论在那篇把这份额外赏赐的文本(“已无从寻觅”——而且事出有因!)大大标举出来的引言中,还是在这份文本本身中,都没有一处暗示我曾在其间发展的思想中扮演过任何角色;也不曾(顺带地)提到这份文本当年注定要成为一篇论文。与Deligne一样,Verdier也不曾认为值得告知我这次出版(其故也一样),也不曾给我寄一部那部障眼的书卷。至于细节,见“同伙”一注(n\textsuperscript{o}~63$'''$,写在发现这场偷偷摸摸的发掘的激动之下),以及在那条已屡次引用的“赊账的学位论文与全险”(n\textsuperscript{o}~81)里更深入的反思。

% ---------------- p.731 ----------------
于是,在他那场不寻常的论文答辩十年之后,Verdier抓住了Deligne提供给他的机会,实质上是在同调代数中对“导出范畴”这一观点预先占下一份无可争辩且排他的父权,有他那位显赫的朋友的完全担保;而这正是他们两个还在继续对这一同样的观点维持着事实上的\textbf{抵制}\fnote{565}{(**)}{如我在先前一条脚注(注见第~页,底本页码空缺)中所解释的,在那部名为“SGA~4$\frac{1}{2}$”的棺材文本中,Deligne不得已曾在“那条”公式的证明中求助于导出范畴。想必正是这一点向他暗示了:该用一篇沉没的论文的“第0状态”来给他的书卷壮壮声势。事实上,直到1981年,这都没有改变对导出范畴的抵制局面。}的时候。那场抵制沉重地压在Zoghman~Mebkhout的工作上,判了他完全的孤独,并且一直维持到1981年的“反常层会议”。

\dmark{}\origpage{932}于是,在1977年,Verdier以一门上同调瑜伽的储备之父的身份出现——而那瑜伽,眼下还是一股合宜的、心照不宣的轻蔑的对象,可谁也说不准……此外,自上一年起,随着那“好的参考文献”的出版,他又以我发展出的那套对偶形式主义之一部分的父亲身份出现(关于结合于闭链的“离散”同调类与上同调类、双对偶性形式主义、以可构性表述的有限性定理等等)——这还不算局部紧空间的对偶,那也停留在一个暧昧的地位,一个等待的地位——一如赋予它以意义的导出范畴瑜伽。

\textbf{第4步}(\textbf{反常层会议},1981年6月)。那是迄今为止Verdier对我的葬礼之参与的登峰造极。那场会议把对Zoghman~Mebkhout——代数簇上同调中 D-模那套统一而丰产的观点的先驱——的肆意掠夺奉为神圣。作为那场会议的正式组织者(与 B.~Teissier一起),Verdier在其中扮演了头等重要的角色。我将在下一条注里,连同“第四行动”(所谓“反常层会议行动”或“当值无名之辈行动”)回到它。这里我只限于谈它对Verdier的直接后果,以“一份遗产之分赃”(而那遗赠的亡者则被小心翼翼地置之不理……)的名义。

那场会议奉导出范畴与三角范畴凯旋重返数学竞技场。以这些范畴的“父亲”身份(而他十五年来竭尽全力要埋葬它们),Verdier在Deligne之后,是那一事件的主要英雄。至少这是从那场会议的主要文章里透出的印象——那篇文章出自Deligne的笔下,一篇自成第一册、并构成会议《论文集》的扛鼎之作\fnote{566}{(*)}{《论文集》出版于 Astérisque n\textsuperscript{o}~100 (1982)——题为《奇异空间上的分析与拓扑》。事实上,那些标注1982年的《论文集》,直到1983年12月才印完,而Mebkhout是1984年1月才读到的。}。仿佛是巧合,一篇论文的(我当年绝不会接受它作为博士论文的)瘦骨伶仃而又救命的“第0状态”——它当年赶来给那部海盗文本“SGA~4$\frac{1}{2}$”救场——那书卷\dmark{}\origpage{933}一时未免单薄了些——瞧,它摇身一变,成了那件光彩夺目的呈堂证物:靠着对“SGA~4$\frac{1}{2}$”的一大阵引用,好让那位来路不明的父亲Verdier以谦逊的姿态,自诩为那场所谓“反常层”的大冲刺(而反常层与此毫不相干!)以及代数簇上同调一场新的、迟到的再出发的预见先驱(那是在一位没有人肯提名道姓的含糊无名者的足迹上……)。

同一篇文章(署名 Beilinson-Bernstein-Deligne)也奉平展语境中的六种运算形式主义(当然,那名字从来不曾出现过)强力回归,带着我在五十年代引入的、如今已被公认的那些记号。如我在别处所写\fnote{567}{(*)}{见“不义”一注(n\textsuperscript{o}~75),第288页。}:“在所引的那篇文章里……没有一页不是深深扎根于我的作品、并带着它的印记,甚至细致到我所引入的记号、以及每一步出现的那些概念所用的名字——那些名字,是我与那些概念初识、而它们还未被命名时,由我赋予它们的。”

我那平展对偶形式主义——十八年前发展的,而那时我的学生Verdier还在学上同调语言的 ABC——在普遍的狂喜中被改称为“Verdier对偶”\fnote{568}{(**)}{在记号索引里,那个我在1963年于平展语境中引入、并且(顺带说)构成Illusie版 SGA~5 报告~I 之主题的对偶函子,被叫作“Verdier对偶”。这名字在正文中到处重现(例如在第62、103页——只要运气好碰巧翻到……)。我发誓我没有半点编造!}。他那显赫的庇护人,在这样的喜庆日子里,可不会在区区小事上吝惜!亡者的名字在那篇文章\fnote{569}{(***)}{我的名字毕竟还出现在文献目录里,带着 EGA 这个缩写(而它总有一天该被一篇专门的文本取代……)。Mebkhout的名字也在正文和文献目录中缺席,整部书卷里都找不到他的痕迹。}里不出现,在由Teissier-Verdier署名的卷首引言里也不出现。那位含糊的无名者(Zoghman~Mebkhout,姑且不点名)的名字也不出现——而没有他,那篇文章、乃至那场辉煌的会议,本都不会问世……

论干活,那真是干了活!除了 Motif(它们不久就要跟上,就在第二年)之外,或许还有晶体瑜伽,一个从未被提名的亡者的上同调遗产之不成问题的分赃,就此告成,而且在一致同意与\textbf{皆大欢喜}中完成。

% ---------------- p.732 ----------------


\subsubsection*{c.~珍宝}

\tocsubsub{c.~珍宝}

\notehead{170(iii)}{}

\origpage{934}\dmark{}(3月1日)我在前几条注里检视的那三场“行动”,关乎我身后留下的“遗产”的分赃——那遗产,是我关于簇的上同调的、书面的与非书面的作品。这场分赃的直接受益者,是我的五位上同调学生中的三位,即Pierre~Deligne、Jean-Louis~Verdier和Pierre~Berthelot\fnote{570}{(*)}{(5月2日)事实上还该加上第四个“受益者”,那是我最近才发现的,即Neantro~Saavedra——在先前的脚注(第921页脚注(*))里已经谈过。}。但这三场行动(以及随后那第四场)中的每一场,都只能靠一大批或多少“搭在”簇之上同调这根线上的同事的串通(有时还有他们的积极支持)才办得成;而这批人中居首位的就是我的五位上同调学生——除我刚才点名的那三位,还有Luc~Illusie与Jean-Pierre~Jouanolou(*)。

这三场行动,以及即将谈到的第四场,在我看来是不可分解地缠在一起的——无论是在它们深刻的动机上,还是在它们最切实的变故上。最早的隐约征兆可追溯到1966至1968年,但更触目的表现,是在我1970年“离开”之后。这次离开,加上“大世界”数学界某种普遍的世风\fnote{571}{(**)}{(5月2日)那想必有一种双向的作用:某种心态败坏的状态(我在离开之前也曾参与其中),有利于我的一批昔日学生对我的作品进行变本加厉的掠夺与贬损,而他们日益增长的犬儒态度,无疑又反过来造成了今天我所目睹的那种多少带普遍性的腐败状态。},为这样一场规模宏大的行动创造了有利的外部条件——在我们这门科学的编年史上,这样的行动大概是独一无二的。

这场行动首先旨在\textbf{败坏}我引进数学的大部分重大\textbf{核心思想}\fnote{572}{(***)}{(5月2日)关于这方面的细节,见“我的孤儿”(n\textsuperscript{o}~45)一注、尤其是“工地的巡礼——或工具与视野”(n\textsuperscript{o}~178)一注。},并埋葬它们所嵌入的那幅统一的图景;然后,败坏或抹杀\textbf{匠人}在我所锻造的那些工具中的创造性角色——那些工具,是在这些思想的指点下、在那幅全局图景的启发下锻造的,而它们曾作为Deligne和我的其他上同调学生之作品的基础工具;最后,在最后一个阶段,侵吞那些有幸被我的学生们采用、或顶住他们加诸其上的\dmark{}抵制\fnote{573}{(*)}{(5月2日)我在那些我曾引进、后来被埋葬、却又顶住Deligne和我的其他上同调学生们设立的抵制而终于自立的思想与工具中,谨指出以下几项:导出范畴;Motif(诚然是狭隘版本);Galois--Poincaré-Grothendieck范畴的瑜伽(为着葬礼的需要被改称“坦纳克式的”);以及围绕场、gerbe 与 liens 诸概念的非交换上同调形式主义(由Giraud本着我在1955年之后引入的出发思想发展而成)。}而终于自立的\origpage{935}思想与工具的父权。

% ---------------- p.733 ----------------
这场行动在1982年以 Lecture Notes 900 那卷书的出版而告终——那部书卷奉 Motif 重返数学公共广场,形态狭隘(相对于我在六十年代渐渐看清的那幅图景),而父权则(心照不宣而显而易见)归于Deligne。它终于在第二年找到了尾声——在那份为庆祝 IHES 存在二十五周年而编印的小册子里端出的、分三个部分的《葬礼悼词》。

对这几篇文字所构成的那条“矿脉”之发现,是在去年5月12日\fnote{574}{(**)}{而正是在同一天,那场对 SGA~5 原始讨论班的无耻屠杀已经向我显露——它出自Illusie之手,并有我全体上同调学生的积极参与或踊跃串通,在“教团全体”的慈祥目光下……},在“葬礼悼词(1)——或恭维”(n\textsuperscript{o}~104)那条注里。近五个月之后,它在承续它的那条注(n\textsuperscript{o}~105)“葬礼悼词(2)——或力量与光环”\fnote{575}{(***)}{为了把对《葬礼悼词》的反思作一次意想不到的扩展,另见下一条注“肌肉与肚肠(阳葬阴(1))”(n\textsuperscript{o}~106)——它同时开启了“阴阳之钥”那场漫长的反思。}中继续。这里我只想用几句话召回这篇不寻常的“悼词”的精神与全部滋味。

那小册子(除别的之外)呈上一幅“肖像长廊”,由关于这座受庆机构历任与现任教授的一篇篇短文组成。在(出自Deligne笔下)献给我的那篇文字里——那文字本该唤起一部作品——却没有说出“上同调”或“Motif”这个词。“概形”一词也没有,也没有任何别的词能让人想到我曾发展过的一门理论、或我曾证明过的一条定理(哪怕它可能有过用处)。相反,我慷慨地得到一堆空洞的最高级赞辞与种种夸夸其谈的殷勤:“宏大的作品……”“二十卷……”\origpage{936}“最大的自然一般性……”\fnote{576}{(****)}{这法语-小黑人(petit-nègre)真是一项千金难买的妙笔,为的是以滑稽的方式(而且装作不经意地……)唤起一位庞大话痨的那份臃肿而多余的夸夸其谈……}\dmark{}“对术语的高度留意……”“他给自己的问题……沿他自己划定的路线……变得太困难了……”。这便是吹吹打打、在聚光灯下的葬礼,靠那篇送得恰到好处的“恭维”送来——它硕大而臃肿,正如那该被“缅怀”其记忆的亡者;同时又带着一股滑稽的暗示之精细,而那精细实在是那位笨拙的祖先所欠缺的……

在献给Deligne的那篇短文 里(并经他过目),没有任何东西能让人疑心我与“那”Weil诸猜想(“其难度有口皆碑”)的证明有任何关系——而那证明被恰当地大书特书。恰恰相反,文中强调:“这个结果显得越发惊人”,因为它不得不说是在与“一系列”出自我手的猜想作对(这位Grothendieck显然从来只会提这个!),而那些猜想(文中还补了一句,为免留下任何关于该如何看待它们的疑虑)“今天与当年一样高不可攀”(读作:当我竟有那个不幸的念头把它们陈述出来的时候……)。

这两幅微型肖像,再加上出色地补全它们的第三篇(仅用一句三行的凝练警句\fnote{577}{(*)}{这第三篇是我在那条已引的“葬礼悼词(2)——或力量与光环”的反思中发现的——而它立刻显得比另两篇合在一起还更有分量!正是它,启发了我给这条注取的名字“力量与光环”。})——都是些真正的珍宝,在献给一位“亡者”(而这一次他还并未过世!)的、以娴熟手法端出的《葬礼悼词》中,想必同样是同类中独一无二的。它们已在上面引过的那三条相继的注(n\textsuperscript{o}~104--106)里得到了应有的仔细检视,并在(晚几星期的)“阴的葬礼(阳葬阴(4))”(n\textsuperscript{o}~124)那条注里,在“阴阳倒转”这一更具穿透性的动态光照下重加检视。

% ---------------- p.734 ----------------
\subsection{(3) 封神(“De Rham 系数与 D-模”)}



\subsubsection*{a.~祖先}

\tocsubsub{a.~祖先}

\notehead{171(i)}{}

(3月1日与5月2--8日\fnote{578}{(**)}{(5月13日)本注与随后四条注,最初本是单独的一条注“封神”(n\textsuperscript{o}~171),写于3月1日。它还把上一条注“珍宝”(n\textsuperscript{o}~170(iii))包括在内。它曾在5月2日至8日间重新整理并大幅扩充——尤其是在其数学部分——并拆分为四条独立的注:“祖先”“作品……”“……与意外之喜”“光荣之日”(n\textsuperscript{o}~171(i)至(iv),再加上上面已提到的“珍宝”)。此外还有八篇子注(n\textsuperscript{o}~171(v)至(xii))关乎上述四条注,以及4月的四篇子注(n\textsuperscript{o}~171$_{1}$至171$_{4}$),讲述我的朋友Zoghman与“行内的法则”搏斗的种种奇遇——依照他本人向我讲的那样。正是这十六条注(n\textsuperscript{o}~171(i)至(xii)以及171$_{1}$至171$_{4}$)构成了眼下“四种运算”中的“封神”部分(而其中封神是第四场、且——在新命令到来之前——是最后一场……)。})在我那场预期中的葬礼里所区分的四场局部“运算”中的每一场里,明摆着是Deligne在充当指挥的角色(或者不如说,是葬礼上的\origpage{937}大司祭\dmark{}),我的另四位的上同调学生或多或少积极地参与,还有一大群数学家的串通——他们全都深知内情(而这内情显然不是要让他们不快的……)。这个“串通集团”在我正准备着手检视的第四场局部行动中,达到了惊人的、近乎难以置信的比例。

\begin{quote}
IV “当值无名之辈”行动(或称“反常层会议行动”)。
\end{quote}

那是一场对\textbf{Zoghman~Mebkhout作品}的行动——据我所知,他是唯一一位(在我离开数学舞台之后)敢于以“Grothendieck的延续者”身份冒险的数学家。

这场行动持续了十年上下,从1975年直到今天。冒着重复的危险,我先要唤起它的历史语境。

五十年代后半期,我在概形的语境中发展出了一套“凝聚对偶”的形式主义。这些反思——出于理解Serre对偶定理在解析几何中\fnote{579}{(*)}{我最早的对偶反思处在解析空间的框架中,且早于Serre的那些。凭着“平展式”的对偶技术与Poincaré-Grothendieck关于算子$\bar{\partial}$的引理(那是我刚证明的),我曾证明:若$X$是一 Stein 簇,则$H^i(X,\mathcal{O}_X)$(相应地$H^i(X,\omega_X)$)是 Fréchet 核空间,与$H^{n-i}(X,\omega_X)$(相应地$H^{n-i}(X,\mathcal{O}_X)$)处于完美对偶中。那时我不曾想到把同一方法套到向量丛的情形上(我没有看出那个极简单的代数事实:算子$\bar{\partial}$既是$\mathcal{O}_X$-线性的,便能延拓到取值于全纯向量丛的可微分形式),也不曾想到 Stein 之外的复簇(而它们是我当时唯一熟悉的)。Serre在一般情形下证明他的解析对偶定理,其证明与我在特殊情形下找到的几乎是同一个。}、尤其在代数几何中的确切意义与范围的愿望——是在一种近乎完全的孤独中进行的,除了我之外不曾有幸引起任何人的兴趣\origpage{938}\fnote{580}{(**)}{当然,在所有数学家中,我本该最指望Serre对我的凝聚对偶反思感兴趣。他曾感兴趣过,如果我记忆不差,是着迷于把他的对偶结果推广到域$k$上的光滑射影簇$X$上的一个凝聚层$\mathcal{F}$(未必局部自由),把$H^i(X,\mathcal{F})$的对偶等同于$\operatorname{Ext}^{n-i}_{\mathcal{O}_X}(X;\mathcal{F},\omega_X)$。这给 FAC 的一个“计算性”结果赋予了内蕴的几何意义(那个结果当年当然引过我的好奇与启发),而那是当$X$为射影空间时的情形。但除了这个成果——它是我在发现对偶的旅程中最早的成果之一,且与他所熟悉的东西还接近——之外,每当我想跟他谈对偶,Serre总是拒绝听。我想我几乎不曾试图跟任何别人谈过,除了(晚得多的)Hartshorne,他依我的想法在哈佛做了一个非常漂亮的讨论班,1966年出版(《Residues and Duality》,R.~Hartshorne,Lecture Notes in Mathematics,n\textsuperscript{o}~20,Springer Verlag)。}。正是这些反思\dmark{}使我逐步提炼出“导出范畴”这一观念 % ---------------- p.735 ----------------
——那些对象堪称各类空间与簇的同调与上同调形式主义中“天然的\textbf{系数}”,并嵌入一套关于带环空间的“六种运算”之形式主义的最初胚胎(带环的 topos 且留待以后)。这六种运算中有四种,早在1955年我那篇《论同调代数的若干问题》的工作里我就多少熟悉了\fnote{581}{(*)}{载 Tohoku Mathematical Journal, 9 (1957), p.~119--221。},只是还没有导出范畴的语言:凭着以后几年确立下来的记号(同时还有导出范畴的观点),那是些“内部”运算$L\otimes$与$\mathrm{RHom}$(层$\operatorname{Tor}_i$与$\operatorname{Ext}^i$的“总导出函子”形式,即我在《Tohoku》中引入的),以及“外部”运算$Lf^*$与$Rf_*$(逆像,与“Leray式”正像),构成两对(或两个双函子的)伴随函子。在$f$是一态射“浸入”$i:X\to Y$的情形,还要加上伴随函子$Ri_!$、$Ri^!$一对,分别体现“以零延拓”与“在$X$中(以其为支集)的局部上同调”两种运算。我的反思中的导线,是要达到一条\textbf{对偶定理}(整体的,那时还谈不上局部版本……),推广Serre为体上的光滑射影簇上的局部自由凝聚层所证明的那条。当时的问题,是要给出一个适用于任意凝聚层(或这类层的复形)、乃至准凝聚层的表述,对$X$既无光滑性也无射影性假设(只保留真态射,那在当时看来是本质的\fnote{582}{(**)}{关于这一点,见下文第940页的脚注(*)。})。此外,与我对 Riemann-Roch 定理的反思相仿,我感觉到正确的陈述不该关乎体上的一个簇,而该关乎一个随便怎样的概形的\origpage{939}真态射$f:X\to Y$。正是靠了历年的工作\fnote{583}{(*)}{不言自明,在这“几年的工作”中,我心头还压着不少别的炉火,远不止凝聚对偶的问题!那时我熟习着当时已知的代数几何奠基(以Serre的 FAC 观点为主要参考),熟习着Weil诸猜想的问题学,还熟习着在一门Serre的课上学到的相交重数形式主义——他在那里发展他的“Tor 的交错和”的想法。这一切,将要触发我在1957年投入 K-理论形式主义与 Riemann-Roch-Grothendieck 定理,而它们(就其精神而言)与我的对偶反思十分贴近。}、一次又一次的\dmark{}近似,整体对偶定理才逐渐从它那些多余的假设中澄清出来,同时导出范畴这一观念也从被预感的地带里走出,获得具体形态,并赋予那形式主义与那些陈述以一种\textbf{内蕴的意义}——缺了它,我会觉得自己根本无从工作!那正是为了达到一条令我完全满意的整体对偶陈述,我才引入了\textbf{对偶复形}的形式主义,并提炼出\textbf{双对偶定理},又发现(在适当的诺特假设下)存在一个内射的对偶复形、本质典范,我称之为“\textbf{残差复形}”,以及它的一套变差理论。整体对偶定理的一个最初陈述——有一阵我觉得那就是“对的”——是说:$Rf_*$与$X$上和$Y$上的对偶复形函子交换(对于两个互相“对应”的对偶复形)。只是后来我才发现:单就对偶复形而言(经由残差复形)的变差理论,能推广为一个性质全新的函子,即函子$Rf^!$,那个“异常的逆像”,在$X$上是局部性质的。由此,也就出现了对偶定理对真态射$f$的最终陈述:这个新函子是$Rf_*$的\textbf{右伴随},于是嵌进一串三个伴随函子

\begin{quote}
$Lf^*$, \quad $Rf_*$, \quad $Rf^!$。
\end{quote}

为求一套完全成熟的形式主义,只差一个函子$Rf_!$(对任意(分离的)有限型态射的“真支正像”)的描述——它推广浸入情形下已知的函子,对真态射化为$Rf_*$,并与$Rf^!$构成一对伴随函子$Rf_!$、$Rf^!$。我不记得五十年代曾为这一缺陷苦恼——那套形式主义超出凝聚(概形的或解析的)对偶之外的普遍范围,\origpage{940}\dmark{}当时仍为我所不及\fnote{584}{(*)}{当然,我已察觉到:即便在开浸入$f:X\hookrightarrow Y$的情形下(那里函子$Rf^!$与“限制到$X$”的函子$Lf^*$重合),后者在准凝聚层的语境中\textbf{不}容许左伴随。通常的左伴随$Rf_!$(“在$X$之外以零延拓”)不保持准凝聚性。另一方面,我也曾验证:即便抛开准凝聚性假设、甚至对一个以一个点为基的真态射,也不存在“对偶定理”。于是,在一般假设下无法定义一个$Rf_!$,这一不可能性在我看来已成定局、且出于事物之本性。

是Deligne在1965或66年(他才刚落脚!)察觉到:人们可以给$Rf_!$一个意义,并把凝聚对偶定理挽回给非真的、分离有限型态射——条件是改用一些作为准凝聚\textbf{pro-层}的(复形)系数。这个漂亮的想法却不曾得到人们本可期待的运气——正如它使之完善的凝聚对偶最初形式主义一样。

Deligne在特征零代数概形上构造“De Rham 系数”的尝试中,曾重拾这一想法并获成功;但那是个有前途的尝试,而他在我1970年离开之后不久便把它付诸流水了。六年之后,轮到Mebkhout去提炼出“那”一类好的“De Rham 系数”(晶体的)——而那是我已期待了十年之久的……}

% ---------------- p.736 ----------------
这一缺口直到1963年才充分显现在我眼前——那时我发现:在新生的平展上同调(以“离散”系数为项)语境中,存在一套在各方面都与凝聚形式主义类似的形式主义,而且还多一个函子$Rf_!$(“真支正像”),它对\textbf{任何}分离有限型态射皆有定义。况且,那正是因为我亦步亦趋地依循我前几年在凝聚情形下所做的工作(那工作除了我之外不曾引起任何人的兴趣),才得以(在一两个星期之内)从那两条基变换的关键定理出发,建立起所谓的“六种运算”完整形式主义。那是一个比此前在超越语境中所掌握的对偶形式主义(仅对拓扑簇及其上的局部系统)无比完善、无比强大的对偶形式主义,甚至比我在凝聚对偶中达到的那个更令人满意。

我的凝聚对偶工作,陈述在 R.~Hartshorne那场众所周知的讨论班\origpage{941}《Residues and Duality》(直到1966年才出版\fnote{585}{(**)}{所说的那场讨论班(出版于 Lecture Notes in Mathematics,n\textsuperscript{o}~20,Springer Verlag)陈述了我的凝聚对偶形式主义观念之主体,围绕六种运算形式主义、双对偶,以及一套“残差复形”理论(它们是对偶复形的典范内射代表)。这些观念在解析框架中被Verdier、而尤其是Ramis与吕盖重新采用。反之,Hartshorne的讨论班不曾包含若干更精细、与这形式主义密切相连的发展:一套残差理论(对任意基上有限型平坦的概形),以及一套差分的上同调理论——据我所知它们从未出版。

我还在五十年代发展了完美复形的“行列式模”形式主义,它本该最终被纳入 SGA~7,而其执笔者(照着 SGA~5 某些“执笔者”已确立的先例)两年后宣布弃权。

最后我指出:紧接我五十年代那些凝聚对偶反思,当时我也曾引入并略加发展纯代数版本的\textbf{Hodge 上同调}与\textbf{De Rham 上同调},尤其是结合于一个代数闭链(起初假定光滑)的上同调类形式主义,以及一套 Chern 类理论——仿照我在 Chow 理论中发展的那一套。}),而关于平展对偶的工作,则在 SGA~4 的一两章里、而尤其在专门献给它的 SGA~5 讨论班里。而直到写这几行的此刻,我才突然意识到:除了一些零星的先声文本(五十年代的 Cartan 与Bourbaki讨论班),不存在任何\textbf{已出版的}、出自我笔下的系统文本,陈述对偶的形式主义与瑜伽——无论是在凝聚语境中,还是在平展语境中。SGA~4 中专门论述这个主题的各场报告(围绕分离有限型态射唯一的“整体对偶定理”,证明$Rf_!$、$Rf^!$是伴随的),是由Deligne在讨论班之后两三年、依据我的手写笔记写成的\fnote{586}{(*)}{Deligne的这番整理处在 SGA~5 讨论班之\textbf{后}。况且Deligne不曾一味遵循我的笔记,而是遵循Verdier1965年在局部紧空间语境中引入的我的方法的一个变体(他大体上沿用了平展的模型)。那时在任何人心里,对偶的一切主要思想的父权、遑论平展对偶的父权,都毫无含糊;谁也不曾想到(想必连Deligne也不曾!)遵循我最初方法的一个变体,会在随后的二十年间被用来浑水摸鱼,把平展对偶归给Verdier(而Deligne则把“平展上同调……”那一“包”里的其余一口袋进自己的腰包……)}。至于 SGA~5 讨论班,它被我的上同调学生们实际禁闭了十一年,终于出版时(在Deligne1977年那部电锯文本之\textbf{后}),已被大肆掠夺、面目全非,被那位“编辑”(sic)Illusie奉其显赫朋友之命糟蹋得一塌糊涂\fnote{587}{(**)}{关于这一点,见“四种操演”一注(n\textsuperscript{o}~169$_2$),以及随后诸子注。}。

那正是在这片废墟中——它曾是我发展过的、也是\origpage{942}我全部几何学著作中最关键的讨论班之一(与 SGA~4 并列)——还留着我亲手写下、或至少依据我的手写笔记写下的唯一书面痕迹,稍稍唤起平展对偶的形式主义与瑜伽;而越过这套尚且片面的瑜伽,还有那套被它不可抗拒地暗示出的六种运算瑜伽。我的学生们细心抹去了后一套瑜伽的一切痕迹\fnote{588}{(*)}{(5月8日)我刚才匆匆浏览了我为 SGA~5 头三场报告做的手写笔记——去年应我的请求,Illusie好心把它们还给了我。(他是那些昔日执笔者中唯一肯费心把我托付给他们的笔记退还给我的人……)头一场报告是对此前的 SGA~4 讨论班就平展上同调形式主义及其与其他种种语境的关系所取得成就的一次广阔“巡礼”。第二场报告长篇宏论地发展了“六种变换”的“抽象”形式主义。那里有一份基本完备的清单,但还不曾着力去圈定那些典范同构之间的相容性。(那是一件性质更技术的工作,在一个我尤其想把“传递”出去、并清楚感受到其全部力量的对偶瑜伽之当口,它是无用的。)

不言而喻,在Illusie版中,这两场报告都不见踪影。我曾终于相信:(被讨论班那些更技术的方面缠住)我想必是疏于呈献那幅统一的图景了。回想起来,而在几乎恰好在发现 SGA~5 讨论班遭“屠杀”整一年之后,我今天似乎已把手指按在这场“屠杀”行动的要害上。那并不是这样或那样几场报告的消失——被一个Deligne并吞、被一个Verdier掠夺、被一个Serre从灾难中救出、或被一个Illusie为着纯粹的乐趣从一幅和谐的“整体”上撕走。而是这个讨论班的灵魂与要害,是这场由一个人独自完成的巨大工作中那根恒定而无处不在的导线——而Illusie着力把它从 SGA~5 中根除,不留(几乎)一丝痕迹。“六种运算”这个名字本身,在这场讨论班中缺席,正如它在我的学生们的著作中缺席——他们想必订有默契的约:只在极罕见的场合才说出这几个字,即其中某一位还与那位被宣告亡故的匠人(而他虽已亡故,却毕竟还得被蒙在鼓里……)狭路相逢之时。},它在整个六十年代以异常的暗示力量,给我的上同调工作以灵感。那才真正是“系数类型”\fnote{589}{(**)}{这个核心思想同样被我的上同调学生们根除并遗忘了。它也是在我对自己的作品及其坎坷作了“十五年后”的第一次回顾之后,最早重新强加于我的那些思想之一,就在“我的孤儿”(n\textsuperscript{o}~45)那条注里。那条注的名字,比我当年所能梦见的更贴切、更深刻,而它却写在发现“葬礼”(这个词的本义与强义)之前。六种运算与“上同调系数”这同一个核心思想,当《收获与播种》的反思重新让我与我的作品在那些曾是我学生们的人手中所受的命运接触时,几乎作为一个主旋律在处处回响。尤其见“墓上的旋律——或自负”(多少展开“系数类型”的“旋律”或变奏主题)与“工地的巡礼——或工具与视野”两注(n\textsuperscript{o}~167、178)。}这一核心思想的“要害”——而 Motif 的瑜伽,正是它的灵魂……

\origpage{943}这样一种畸形的局势——一门科学中的重要进展,体现于一幅新图景,却被那些曾是它第一批受益者与保管者的人亲手根除——若不是另一种局势,也同样堪称例外,即由我的骤然离开及其周遭条件所造成的那种局势,是不可能确立的。何况,事态的走向,早在我离开之前、在整个六十年代,就已由我所处的那种分裂局面预备好了——一方面,我被无穷无尽的基础工作占着,而那工作只有我能或才愿意承担\fnote{590}{(*)}{我提醒一句:这项维度宏大的基础工作,自我离开那一天起便戛然而止,直到今天依然如此。这是我在“藏宝”一注(n\textsuperscript{o}~169$_5$)里所谈的那个“误会”之一个雄辩的标志。人人都已准备好搬来自己的家具、永久住进我建造的那些房子里——但要再动手、拿起泥刀与铅垂线去继续建造和布置,哪怕只是在种种需要的紧迫压力之下,已经没有一个人了……},另一方面,又被一些主题往往离得很远的问题不断叩问……

% ---------------- p.738 ----------------
那些当时吸引着我的基础工作,以及那些因此往往更强烈、更直接地迷人的东西\fnote{591}{(**)}{若我当时肯听自己的,有多少次我会撂下那些我为众人操劳的没完没了的基础工作,去投入那场不断召唤我的未知冒险——那真正的冒险——而不把由我发现的新土地让别人去游逛的乐趣让给人。今天我看到:那些土地始终是处女地,或差不多,而那些我当年看作先驱的人,早在我离开之前就已选择做安逸的食利者……}即便是那些我曾让自己有空去深究和发展的主题(比如对偶),我也少有那份余暇,把研究结果整理成宜于出版的形式(合乎我那高标准的要求)。于是,我时常索性把整理的工作留给别人(那些我全然信任的人)——就像“对偶”这个主题,无论在凝聚框架还是平展离散框架中都是如此;或者把一些我知道是丰产的出发思想交给别人去发展(比如导出范畴的观念、或晶体上同调,这里只举这两个为例,同类例子还多得很)。在一种“正常”的局面下——即\origpage{944}一片诚意回应着我向一批有干劲的学生发出呼吁时的信任,他们在我的接触中学习他们的手艺、并为他们未来的作品学些有分量的根基——一切本是最好的\dmark{},对一切相关的人都是最大的好处,科学共同体也不例外。但实情是:这种不寻常的局面,把一份可观的\textbf{权力}放到了他们手中(而直到去年之前,这个念头从未在我脑中掠过……),尤其在我离开之后。我离开之后不久(甚至就在之前……),他们中某些人就忙不迭地滥用这份权力,去抹杀作品与图景,去贬损匠人,并把他们以为用得上的一些他亲手打造的工具据为己有。

我的凝聚对偶工作,似乎从未大受欢迎\fnote{592}{(*)}{如我在先前一条脚注中指出的,这些工作曾启发Verdier、Ramis、吕盖在解析空间的凝聚理论方面的工作。“一向以来”就清楚(至少对我如此):同一套形式主义必能在刚性解析语境中重现(而那个语境,据我得到的回声,也还停留在童年阶段)。

另一方面,Mebkhout告诉我,日本的分析学派相当大地从《Residues and Duality》中汲取了灵感——同时却绝口不提那位匠人的名字。凭眼前的时风,反过来的话倒要叫人惊讶了……}。相反,平展对偶的工作却立刻引来注意。但我认为更公允的说法是:引人注意的是,有人“想明白了”,不管用什么办法,在平展语境中证明了Poincaré对偶的类似物——而那对偶,我想人们在熟悉的定向拓扑簇的语境中已经认识了近百年。那对平展上同调来说曾是“一个好得分点”(人们几乎不再怀疑:对Weil诸猜想——“其难度有口皆碑”——来说那才是“对的”)。这就是说,那群在著名的猜想旁逡候着的数学公众,是以一个“消费者”的身份反应的——他不情愿承认并消化一幅新的、深刻的事物图景,只取一个面目熟悉的“结果”。二十多年后,我确认:六种运算与系数类型这个富有力量的图景(以一套简单得令人困惑的形式主义表达),仍然被所有人无视(唯一的例外是那位孤独的匠人);而每当有人敢提一句,它就成了冷嘲热讽或刻薄挖苦的对象\fnote{593}{(**)}{至于细节与评注,见子注“无用的细节”(n\textsuperscript{o}~171(v)),尤其是(a)部分“上千页的一捆……”。}。我那套装备中这样那样的零散成分,在这里那里被人使用,却丝毫不提我个人(而且早已找好的替代父亲),尤其是双对偶形式主义——自从\origpage{945}1981年那场难忘的会议之后、在相交上同调的大冲刺里。但那幅\dmark{}\textbf{图景}——带着孩子般的简洁与完美的优雅,而它曾给出过自身力量的有力见证\fnote{594}{(*)}{至于这些“有力的见证”的细节,见子注“无用的细节”(n\textsuperscript{o}~171(v)),(b)部分“无所事事的机器……”。}——却依然无人识得,只是那些宁可用轻蔑(和掠夺……)代替理解的人轻蔑的对象。

% ---------------- p.739 ----------------
如果说我亲手、并用心所做的东西,比它的时代超前了二十年、或许五十年,那不是因为我三十年前动手时遇到的那门\textbf{数学}还不成熟。那是因为人的不成熟\fnote{595}{(**)}{关于这个问题的一点初步反思,见子注“自由……”(n\textsuperscript{o}~171(vii))。}。而我的身后学生、我唯一的延续者Zoghman~Mebkhout,所面对的也正是同一种不成熟。我很有幸,在1970年离开之前,只以“不理解”的形式遭遇过它——而那不理解从未卸下友好的姿态。Zoghman~Mebkhout,来到一片与他所大胆延续的那份工作时代不同的数学场域,他在不理解与轻蔑之后、在他的某一项成果作为工具的价值终于被承认之后,还领教了前辈们的恶意与一个时代全部的不义——但我抢跑了……

我在数学中最重要的发现之一、而它至今仍几乎被所有人无视的,是那套我在五十年代开始发展的对偶形式主义的\textbf{无所不在}:“六种变换与双对偶的形式主义”同时适用于人们最初设想的“连续”系数(“凝聚”理论),与“离散”系数。这种无所不在,在1963年春显得像一次几乎难以置信的惊喜——正是靠了它、而非别的,我才得以发展出一套平展对偶形式主义,并达到我所谓的对平展上同调的“掌握”。而且早在那个时候,我就被一个问题困扰着(虽说不曾太停下来):一门“共同的”理论的可能性——无论是在概形框架、复解析框架、乃至拓扑框架中;一门能“笼住”两类系数的理论。De Rham 上同调(我的一位老朋友……)给出了这方面的第一个迹象,提示人们到“带可积联络的模”(或“分层模”,或许……)的方向去寻找一个“共同原则”。这些模生出一门“De Rham 上同调”(实质上以离散系数为项),于是被发现与凝聚上同调\origpage{946}处在关联之中。这个进路后来向我提示了“晶体”与“晶体上同调”的念头,然而它似乎还不足以(看来)提供那把钥匙,去描述一套完整的六种变换形式主义——针对那些在适当意义上同时涵盖离散(“可构”)系数与连续系数的“系数”类型。我的学生中,似乎\origpage{947}没有一个懂得感受这个问题\fnote{596}{(*)}{写这几行时,我关于这事的记忆还模糊。它后来变得清晰了,而我在子注“荒唐的问题”(n\textsuperscript{o}~171(vi))里更周详地回到这一点。}\fnote{597}{(**)}{我曾向Verdier谈过这个问题——那时他已(如我所建议的)发展了拓扑空间的(至少是其胚胎的)对偶理论,仿照我在平展语境中发展的那一套(关于这一点,见子注 n\textsuperscript{o}~81、81$'$)。那该是六十年代中期。显然当时没有“对上暗号”——那个问题的意义(或许有点含糊,实情如此)似乎从他嘴边溜过了。可我必定提到过 De Rham 上同调,既可微的、也复解析的,它把Serre对偶与Poincaré对偶联系起来,关乎这两种类型的系数。

(5月14日)况且,早在五十年代我就知道:可以把Serre的对偶定理推广到真光滑相对概形上局部自由层之间的微分算子复形的情形,以便也把 De Rham 上同调涵盖进来(因而,实质上,一门以离散系数为项的上同调)。那是一个与Mebkhout在解析框架中的对偶结果非常接近的对偶结果——而后者将在下一条注里谈到。当时我不曾沿着这条路走下去,尤其因为(我相信)我看不出该怎样用这些微分算子复形做出一个合适的“导出范畴”——因为缺乏一个好的“拟同构”观念。实情也是:我当时工作的孤独——做着一些(凝聚上同调)显然除我之外全世界无人感兴趣的问题——几乎不足以刺激我把又一项推广(用微分算子代替线性映射)堆到我在前几年于自己角落里已提炼出的那些推广之上。然而从Mebkhout的视点看,我本是万事俱备:那里通向相应的 D-模(一个微分算子复形的各个分量)的过渡,给出了一把具有完美简洁性的钥匙,去构造所需要的导出范畴。况且,早在1966年(尽管当时我没有清楚地意识到),我就有了一个对偶的视点,它本可以让我用“分层 pro-模”造出一个导出范畴(这个观念后来由Deligne在他那套 De Rham 系数的草图中发展出来,下文将谈到)。的确,给任意一个凝聚模结合其\textbf{无穷阶主部}所构成的 pro-模,后者带有典范的分层;或者给一个微分算子复形结合一个由这类分层 pro-模组成的复形,其晶体超上同调便与所考虑的微分算子复形的Zariski超上同调相同。(参见我的报告“Crystals and the De Rham Cohomology of schemes”,由 J.~Coates 与 O.~Jussila 记录,收入《Dix exposés sur la cohomologie des schémas》,North Holland,第~306--358~页,尤其是第~6~节。)于是,可以按通常的方法,借助相结合的分层 pro-模复形,为微分算子复形之间的(微分)态射定义“拟同构”的概念。}

% ---------------- p.740 ----------------
\dmark{}只有Deligne例外。他用了一整年讨论班(我想是在 IHES,1969/70年)去发展一套形式主义——那套形式主义至少使他能(对特征零体$k$上有限型的概形$X$)描述一些(所谓“De Rham 的”)上同调空间,它们在$k=\mathbb{C}$的情形下,重现通常的(以超越方式定义的)复“Betti 上同调”。他所用的系数是“分层 pro-模”以及这类 pro-模的复形。然而当时并不清楚这些系数是否能嵌进一套六种运算的形式主义\fnote{598}{(*)}{这里我的记忆仍然是模糊的,而且是个错误——先验地清楚,出于超越性质的启发式理由,这里\textbf{必定}存在一套六种运算形式主义。(细节见子注“……与阻碍”,n\textsuperscript{o}~171(viii)。)

我的错误显然来自一个刻意的(浮在意识表层的)说法:把一件或许看似无法解释的事情合理化、讲得通——即Deligne放弃一条“稳妥”而富于希望的研究方向。那理由绝非数学性质!},而Deligne放弃了继续走这条路。如果我没记错,当时最缺的(*)、令人能放心的一着是:用一种纯代数的术语(借助凝聚或前凝聚模与分层),对$X$上“代数可构”的$\mathbb{C}$-向量空间层之范畴做一个描述——该描述因而对任何特征零的基体皆有效,而这个范畴在基域为复体$\mathbb{C}$时本是以超越方式定义的\fnote{599}{(**)}{我提醒一句:这个可构性观念是我在五十年代引进的,带着众多变体(代数的、实解析的等等),而那正是我严格地独自一人对这些题目感兴趣的时刻。(见我去年的评注,在子注 n\textsuperscript{o}~463 中。)}。



\subsubsection*{b.~作品……}

\tocsubsub{b.~作品……}

\notehead{171(ii)}{}

Mebkhout的工作始于1972年,处在(技术上更艰难的)解析空间的超越语境中。在以后几年里,他在一种几乎完全的孤独中,熟悉我的上同调作品与\origpage{948}导出范畴形式主义\fnote{600}{(***)}{(5月14日)Mebkhout后来向我说明:1972年前后对数学文献的这第一批阅读,是日本Sato学派作者们的著作。他费了很大劲(如他所说)才理清头绪,那在他看来复杂得可怕。而他正是在那里找到了Hartshorne《Residues and Duality》一书的指涉——读那本书对他是真正的放松。诚然,那本书写得极其出色!我为那书写的几行引言,唤起其中所发展之形式主义的无所不在,大大启发了他。正是从那时起,他着手熟悉我的作品,而它后来成为他主要的灵感来源。在这一切工作与报告中,他留心清楚地指明这个来源。}——它们被那些曾经是我学生们的人弃置不顾。\dmark{}一条导线,在他的反思中渐渐取得第一重要的位置,那就是连续对偶与离散对偶之间惊人的平行。后者在这期间取了“Poincaré--Verdier对偶”的名字,而大世界里没有人(尤其是那位新“父亲”Verdier)装作对这个平行的深刻理由有所疑问。那是“实用”而短视的观点当道:使用我早就创造好的现成装备,而不提任何问题——尤其不提那种模糊至极、甚至荒唐的问题。这个问题在任何已出版的文本中都不曾提及,甚至(而我意识到这里是在自责……)在我自己笔下的文本中也不曾\fnote{601}{(*)}{(5月14日)然而我记得:在 SGA~5 讨论班期间,我心中始终挂着我那套形式主义的无所不在,而我不放过一个机会,指出我在离散平展上同调框架中所发展的观念与技术在其他语境中的可能变体。我很难相信,在口头讨论班上我竟没有提过综合这两类系数的问题,哪怕只是在最后那场关于未解决问题的报告中——它也从那部屠杀版中消失得无影无踪。不用说,这部精心删去了所有不符合“技术性离题文集”这一严整标签的内容的版本,绝不会提到这类问题……\par (5月19日)关于此事,另见子注“死去的书页”(n\textsuperscript{o}~171(xii))。}

% ---------------- p.741 ----------------
共同形式主义的表述本身,就离不开导出范畴。Mebkhout始终把它们作为工作工具,顶着时尚的风向与前辈们的轻蔑——首当其冲的,还是那位不知情愿还是勉强地充当这些范畴“父亲”的Verdier。相对于我引入的整套工具,Mebkhout新增的关键成分,是Sato及其学派的微局部分析。确切地说,他借用了他们在光滑复解析簇上的 $\mathcal{D}$-模概念(它等价于我在1965--1966年前后引入的“模的晶体”概念,而后者在更广泛的语境中仍有意义,尤其适用于奇异簇),更借用了 $\mathcal{D}$-凝聚性的概念以及凝聚 $\mathcal{D}$-模上精微的完整性条件。此外,他实质性地运用了Kashiwara1975年的一个定理:与一个\origpage{949}\dmark{}完整 $\mathcal{D}$-模相结合的微分算子复形,其上同调层是解析可构的。在Mebkhout两年前向我讲述之前,这些观点与结果我一无所知;Deligne在1969--1970年间思考德拉姆系数形式主义而没有继续下去时,想必同样不知道这些。正是\textbf{把这两股观念汇合起来},Mebkhout才得以统一把握光滑复解析簇 $X$ 上这两类系数,用微分算子复形来表达,或者更好、更精确地,用较为灵活的 $\mathcal{D}$-模语言,表述为具有凝聚上同调的 $\mathcal{D}$-模复形。\fnote{602}{(*)}{关于 D-模语言的细节、它与微分算子复形语言及晶体语言的关系,见子注“五张照片(D-模与晶体)”(n\textsuperscript{o}~171(ix)),(a)部分。}这正是他对当代数学的伟大贡献。

说得更确切些:若$X$是一光滑复解析空间,记$\mathrm{Cris}_{\mathrm{coh}}^*(X)$为导出范畴$D^*(X,D_X)$中由具有$D_X$-凝聚上同调的$D_X$-模复形组成的全子范畴,$\mathrm{Cons}^*(X,\mathbb{C})$为导出范畴$D^*(X,\mathbb{C}_X)$中由$X$上具有解析可构上同调的$\mathbb{C}$-向量空间层复形组成的全子范畴,而$\mathrm{Coh}^*(X)=D^*_{\mathrm{coh}}(X,\mathcal{O}_X)$为由具有凝聚上同调的$\mathcal{O}_X$-模复形组成的全子范畴。Mebkhout揭示出如下的基本函子图表:
% 原书图内的 coh 下标与正文的记号不完全一致,这里照录原图.
\begin{equation*}
\begin{tikzcd}[row sep=2.5em, column sep=1.7em]
\mathrm{Cons}^{*}_{\mathrm{coh}}(X,\mathbb{C}) \arrow[dr,"M"'] & &
\mathrm{Coh}^{*}(X) \arrow[dl,"N"] \\
& \mathrm{Cris}^{*}(X) &
\end{tikzcd}
\tag{Meb}
\end{equation*}

其中右边的函子$N$是“重言式的”函子——由显然的包含$\mathcal{O}_X\to D_X$所作的标量扩张函子的总导出。左边的函子$M$,或称“Mebkhout函子”,性质要深刻得多\fnote{603}{(**)}{至于一个紧密相关的函子$M_\infty$(在$D^\infty$-模语境中)的“显式”描述,见已引子注 n\textsuperscript{o}~171(ix)的(b)部分:“好上帝公式”。}。它是\textbf{全忠实的},而它的本质像,是由那样一些$D_X$-模复形组成的$\mathrm{Cris}_{\mathrm{coh}}^*(X)$的全子范畴——那些复形的上同调层不仅凝聚,而且还“完整的”且“正则的”。这些是精微的局部条件,前者由Sato学派引入,后者\origpage{950}由Mebkhout临时定义\fnote{604}{(*)}{“正则”这个名字当然是从单复变量函数的微分方程的“正则奇点”的经典术语借来的。若$i:U\hookrightarrow X$是补集$U=X-Y$(对$X$中的一个除子$Y$)的包含,则Mebkhout意义下的(对一个$X$上的 D-模复形$C$)“沿$Y$的”正则性,可以这样写:典范态射
\[
Ri^{\text{mér}}_{*}(C_U)\longrightarrow Ri_{*}(C_U)
\]
从$C$在$U$上的限制$C_U$的“亚纯正像”到通常的正像的那个典范态射,对所结合的 De Rham 复形诱导一个拟同构。

在$F_U$化为一个“局部系统”(即一个带可积联络的$\mathcal{O}_U$-凝聚层)的情形,这个观念等价于Deligne的那一个。后者显然也是受我的比较定理启发(其区别在于:Deligne绝口不提这一点,而Mebkhout则始终留心清楚地指明他的来源)。Mebkhout是在引入他自己的(超越性质的)定义之后,才得知Deligne的观念的。他先前不曾为他的条件寻找一个纯代数的描述。Deligne的工作表明:在所设想的特殊情形中,Deligne的代数条件蕴含Mebkhout的那个,而Mebkhout验证逆命题也成立。这就即刻提供了钥匙,去为Mebkhout的正则性条件作一个纯代数的描述——对任何一个具有凝聚上同调和完整的 D-模复形。

Mebkhout告诉我:日本人有一个“具正则奇点的微微分系统”的观念,他们把它用在全然不同的精神里(为的是分析、而非几何的需要)。在“好上帝定理”的大冲刺之后,那是一个现成的手段(在众多别的之中),为的是打乱牌面并抹杀Mebkhout的先驱工作。看起来两个观念是等价的——而鉴于这个题材中被蓄意造成的混乱状态,很可能谁也不曾费心去验证这一点。Mebkhout始终只用他1976年引入的那个正则性观念工作(它见于他的论文,两年后提交)。},而他尤其(如他所说)是从我的代数 De Rham 上同调与解析 De Rham 上同调之间的比较定理——为了我曾在离散平展上同调(即复 Betti 上同调)的框架中发展的那些观念与技术——得到启发。

% ---------------- p.742 ----------------
这些条件(而在我看来这正是它们的主要趣味所在)事实上是“纯代数的”,尤其在$X$被换成特征零的任意体上有限型概形(要光滑也无妨,但不必需)的情形下仍保持意义。

Mebkhout函子$M$(或称“好上帝函子”\fnote{605}{(**)}{至于“好上帝定理(或函子)”这个名字的来历与意义,见“当值无名之辈与好上帝定理”一注(n\textsuperscript{o}~68$'$)——那注写在我知道反常层会议那套神秘化之前、甚至在我知道“以其全部辉煌的葬礼”之前。})可描述为如下函子的拟逆:

$m : \mathrm{Cris}^{*}(X)_{\mathrm{hol.reg.}}\longrightarrow\mathrm{Cons}^{*}(X,\mathbb{C})$,

由以下式子定义

$m : F \longmapsto \mathrm{DR}(F)\overset{\mathrm{def}}{=}\mathrm{RHom}_{D_X}(\mathcal{O}_X, F)$,

\origpage{951}\dmark{}即(定义在整个$\mathrm{Cris}_{\mathrm{coh}}^*(X)$上的)那个把每个(具有凝聚上同调的)$D_X$-模复形结合到它所结合的微分算子复形(或“De Rham 复形”)的函子的限制\fnote{606}{(*)}{关于这一点,见已引的“五张照片(晶体与 D-模)”一注(n\textsuperscript{o}~171(ix)),(a)部分“‘De Rham 系数’相册”。}。Kashiwara的可构性定理蕴含:当$F$是完整的(从而,当它完整且正则时更是如此),$\mathrm{DR}(F)$确在$\mathrm{Cons}^*(X,\mathbb{C})$中,这使得函子$m$可以定义——那定义固然不言自明、孩子气,然而除了Mebkhout(而且直到五年后的“大冲刺”之际还是如此……)谁也不曾想到\fnote{607}{(**)}{(5月7日)人们不得不把那\textbf{两个}函子$m$、$M$——它们在一个方向和另一个方向建立起关键范畴的等价——称为\textbf{Mebkhout函子};$D^\infty$-模的函子$m_\infty$、$M_\infty$也同样。(关于后者,见所引的“五张照片”一注(n\textsuperscript{o}~171(ix))的(b)部分。)把这些函子与自然的对偶函子复合,人们得到另外两对互为拟逆的函子$(\delta,\Delta)$与$(\delta_\infty,\Delta_\infty)$,它们是反变的,在某些方面更方便(见所引注)。那四个就是“Mebkhout反函子”。}!那要先记得某一种瑜伽——导出范畴的瑜伽,而大家不约而同地决定把它连同引入它的那位亡故者以及他别的同类夸夸其谈一起埋葬\fnote{608}{(***)}{(5月7日)不止一次,Mebkhout被当作一个滑稽人——他以为在导出范畴之间写箭头(您倒是评评这个理!)再加上些$\mathrm{RHom}$,那就是做数学……可他并没有因此动摇,正如我当年(1955年)引入层的整体与局部$\operatorname{Ext}^i$(以及后来带$\mathrm{RHom}$或不带的)时一样——那些让所有的人都晕船,并为他们对我的最直言不讳的保留提供了理由(至少直到1957年,即 Riemann-Roch-Grothendieck 之年……)。

这一切都不曾妨碍Mebkhout相信自己的嗅觉,并跟着它走到它带他去的地方。他赤手空拳地投入工作,没有经验,没有任何人的帮助。他\textbf{确信}他预感到的定理必真——他手上的一切迹象都是相互印证的。稍许有点经验的话,就连他早已握有证明它的一切手段都会是显然的——用如今标准的手段,我的学生中随便哪一个三下两下就能证明。

然而独自一人、只能靠自己,那定理在他看来远得眩目、高不可攀——他几乎不敢指望有朝一日能证明它!

如果说他为证明它辛苦了近两年,那是因为他不曾具备我的学生们曾具备的优势:有一位慈善的长者搀扶,并在我的接触下学到某一种标准技术——可构层的拆卸技术,加上Hironaka的奇点消解。他所提炼出的陈述诚然深刻,而证明也同样深刻,但今天性质上已是标准的。回想起来,他当年要克服的困难主要是心理上的,更甚于技术上的:逆着潮流工作,而且完全只靠自己的光明……}

% ---------------- p.743 ----------------
\origpage{952}\dmark{}此外,那个\textbf{正则性}条件——超出完整性之上的——是由Mebkhout“量身”提炼的,其方式恰恰使得函子$m$受到这样的限制后,可望是全忠实的、甚至是范畴等价。他在1976年就达到了这个信念。他终于在1978年初、在他的论文中以一个非常相近的形式把它证明了出来\fnote{609}{(*)}{(5月5日)在他的论文里,Mebkhout对$D^\infty$-模陈述并证明了相应的等价定理,并给出了拟逆函子$M$的一个卓越的显式表达。关于这一点,见子注 171(ix)((b)部分),以及子注“一个视景的萌发——或闯入者”(n\textsuperscript{o}~171$_{1}$)。Mebkhout早在1976年就达到这样的信念:两个函子 $m$、$m_\infty$(因而还有标量扩张函子 $i$,即上面最后引的那条子注里谈到的那个)是等价,并得到了 $m_\infty$ 的拟逆函子的显式形式。他那论文里关于$m_\infty$的结果,是1978年的。从那时起,他就握有证明(类似的、但带着额外技术困难的)$m$情形的全部配料。

鉴于他的论文(1979年2月通过)所受到的一片冷淡,他那时不曾费力也为$m$的情形写一份像样的证明。那些配料与$m_\infty$的相同,并且受我的复代数簇的 De Rham 上同调比较定理之证明启发(他是1975年得知那个证明的),以及 SGA~5 的拆卸技术(他是从Verdier那“好的参考文献”里学到的,而 SGA~5 讨论班仍在我的好上同调学生们的小心禁闭中)。只是到1980年末,鉴于他的想法对 Kazhdan-Lusztig 猜想的证明日益重要,他才费力为$m$的情形(那里人们并不预先掌握一个拟逆函子)写了一份周详的证明。那证明出版于《另一个范畴等价》,Compositio Mathematica 51 (1984), p.~63--88(手稿收到日期 1981年6月10日)。
\par 我强调:1975年至1980年间,除了Kashiwara在1980年用几行文字提及它(下文“黑手党”子注,n\textsuperscript{o}~171$_2$,将谈到),在Zoghman~Mebkhout自己的著作之外,文献中没有任何地方谈到函子 $m$ 或 $m_\infty$,也没有谈到一种对偶“哲学”,把解析可构离散系数与完整正则 $\mathcal{D}$-模复形,或完整 $\mathcal{D}^{\infty}$-模复形精确地联系起来。正如我们将看到的,当这一联系的重要性终于因Kazhdan-Lusztig 及Deligne主持下的相交上同调热潮而得到承认时,Zoghman~Mebkhout的名字却在一份轻声细语、面带微笑、低调而又有效得冷酷无情的默契中,悄无声息地被抹去了……}。

\origpage{953}\dmark{}那才是Mebkhout带来的伟大的新定理,是八年顽强的、在完全孤独中持续的工作之加冕。它在一个简练的陈述里,包含着一整系列深刻的、一般性递增的结果,那些结果在1972到1980年间被耐心地提炼并逐个证明。至于这场孤独的旅程中发现“簇上同调中一种新哲学”的几座大路标,我请读者参看子注“三座路标——或天真”(n\textsuperscript{o}~171(x))。在这条注里,我的目的主要是用几句话描述那幅新的图景——它呈现在这位孤独匠人劳作的第一个漫长阶段之末——即Zoghman~Mebkhout。

关键的事实(Mebkhout早在1976年就清楚认识到的)是:范畴$\mathrm{Cons}^*(X,\mathbb{C})$(性质上“拓扑的”)可以借助于Mebkhout函子$M$,被解释为范畴$\mathrm{Cris}_{\mathrm{coh}}^*(X,\mathbb{C})$的一个全子范畴——而后者在“抽象”代数几何的框架中仍保持意义;它也可以“实质上”被解释为一种“导出范畴”,由通常意义上的微分算子复形构成\fnote{610}{(*)}{关于这两个视点之间的精确关系,我请读者参看被大量引用的子注“五张照片”(n\textsuperscript{o}~171(ix)),(a)部分。}。这个全子范畴由
% ---------------- p.744 ----------------
完整性(Sato式的)与正则性(Mebkhout式的)条件定义,显然就是我在六十年代就已预见的那种“De Rham 系数”的好范畴——它当年还缺失于我的装备中,在特征零的情形,为的是补全并相互联结(如一幅巨大的扇面)我在1963年提炼出的那些“$\ell$-adic 系数”;那也是Deligne在六十年代末曾试图把握而(似乎)不曾以一种令他满意的方式做到的范畴。这个范畴显然将在代数几何中(尤其是在基概形$X$上 Motif 范畴的描述中……)扮演一个本质的角色。在我看来,至少对我来说,这个范畴当得起的名字是“\textbf{De Rham-Mebkhout 系数范畴}”\fnote{611}{(**)}{对这个范畴的关键角色与意义的普遍不理解,早已显明在这样一个事实中:它至今还不曾得到一个名字、也不曾得到一个简练的记号。代替它的(在我看过的文本中),作者们只是含糊地指到“完整正则的微分系统”(愿有好运气的人在其中找得到门路!),指到这些系统与 $\mathcal{E}$-可构层之间的“构造”“对应”或“(据说众所周知的)关系”——而且当然,照例要严格地沉默掉那位孤独的匠人,是他发动了这场围绕数学俊髦世界的新热门点心“D-模”的大喧闹。},记作$\mathrm{DRM}^*(X)$(或$\mathrm{Meb}^*(X)$),或在概形框架中、当$X$是特征零的体$k$上有限型概形时,\origpage{954}记作$\mathrm{DRM}^*(X/k)$(或$\mathrm{Meb}^*(X/k)$)\dmark{}\fnote{612}{(*)}{在代数情形,除了局部的“正则性”条件之外,还必须加上一个“无穷远处”的正则性条件(在非真簇的情形),才能找到那些“好的”De Rham-Mebkhout 系数——它们在基域为复体时,将对应到$X_{\mathrm{an}}$上具有\textbf{代数}(而不只是解析)可构上同调层的$\mathbb{C}$-向量空间层复形。也正是为着这些系数,人们才有一条“比较定理”,推广我关于 De Rham 上同调的结果,即:晶体“总上同调”$R\Gamma_{\mathrm{cris}}$,无论取代数(Zariski)观点还是取超越意义,都是“同一个”。这个陈述本身又该被看作一个更完备的陈述的特例,即:代数观点的“六种运算”与超越观点的六种运算是“相容的”。

如果我的学生们不曾那么忙着埋葬师父的作品,那该是七十年代最初(若不及六十年代起……)他们就该提炼出那门势所必至的系数理论,以它全部的简洁与全部的力量……}。

正是通过上面那个函子图表(Meb)——它概括了Mebkhout的哲学(可追溯至1976年,并由他在往后几年中确立)——那些\textbf{凝聚晶体系数}(即$\mathrm{Cris}^{*}_{\mathrm{coh}}(X)$的对象)才能被看作是“离散”(可构)系数与“连续”(凝聚)系数的一种“共同推广”。由前者构成的范畴,无论如何通过Mebkhout函子$M$(一个性质深刻的函子),等同于由 De Rham-Mebkhout 系数构成的晶体凝聚范畴的\textbf{全子范畴}。情形对那个重言式的函子$N$就不那么好,它毫无全忠实之处。但为了安慰我们并把图景补全,可以补充说:在所涉及的每一个范畴中,人们都有一个自然的\textbf{对偶函子},它给出一条双对偶定理(对$\mathcal{O}_X$-模与$D_X$-模是“平凡的”,而在可构$\mathbb{C}$-向量空间层的情形要用到Hironaka奇点消解的全部力量)——依我当年在凝聚(交换)框架中、而后在(\origpage{955}1963年的)平展离散框架中所提炼的模型\fnote{613}{(**)}{(5月5日)把我的双对偶结果、以及可构性在$\operatorname{RHom}$之下的稳定性,从平展语境推广到解析语境,那其实是自动的、而且我早在1963年就已知道;那时Verdier已与我共事三年,正投入导出范畴(他已承担了为它做系统理论的工作)与凝聚对偶的瑜伽。正是从我口中,他学到了那些把凝聚对偶形式主义推广到离散系数情形的手法。如我们所见,十三年之后(1976年),他在“好的参考文献”中占有了那套解析复语境中的对偶与双对偶瑜伽,并与Deligne及其他上同调学生串通——他们全都深知内情。

在第二年(1977年)的 SGA~5 屠杀版中,Illusie保留了(在报告~I 中)双对偶定理,因此对同时读过两种文本的读者来说,Verdier的骗局显而易见——但显然所有人都视之为寻常(鉴于眼前这年景……)。相反,Illusie回避把可构性经由$\operatorname{RHom}$的稳定性这一结果收进去——而那当然是我在陈述并证明双对偶定理之\textbf{前}就给了的,而Verdier抄去的我的证明丝毫不依赖它。于是(这种事真做得出来!)Illusie只满足于在第二个论元是对偶复形时确立所说的稳定性!!那是替他的朋友Verdier遮掩的一种方式,使那部从头到尾(除了在他该提的地方提过的那三页)抄自我的 SGA~5 各场报告的文章,不那么显得出这一点。最美妙的是:所说的稳定性早已是对偶形式主义的一个直接推论(这并不妨碍:装作只对$G$是对偶复形时确立可构性经$\operatorname{RHom}$的稳定性,在数学上是荒唐的)。但那位殷勤的Illusie在他的报告里绝口不提这个推论,好保住那副样子:似乎出现在他那同伙的“好的参考文献”里的稳定性结果,确确实实是他的独创。

人们不禁要问:既然如此,Illusie为何还是保留了双对偶定理——要屠杀就屠杀到底,他为这一点不周又算得什么呢!但假如他把它掏空,他就势必也得同时掏空那条永恒的Lefschetz--Verdier公式(而它对那定理的运用是本质的)——那恰是“特洛伊木马的头颅”:那条公式,正是它在 SGA~5 中所谓的关键角色,该为他的另一位同伙那场无礼的“电锯”行动辩护,而那行动炸碎的是我的平展上同调作品的统一性。

再恭喜一次我的昔日学生Illusie,那位机灵的“编辑”-掘墓人……

};
% ---------------- p.745 ----------------
此外,这两个函子$M$、$N$是与自然的对偶函子相容的\fnote{614}{(*)}{对重言式的函子$N$而言,这个相容性本身是重言式的。相反,对Mebkhout函子$M$(或者说,对它的拟逆 $m=(G\mapsto\mathrm{DR}(G)=\mathrm{R\mathcal{H}om}_{\mathcal{D}}(\mathcal{O}_X,G))$)而言,那是一个深刻的结果,由Mebkhout在1976年证明(名为“局部对偶定理”),与 D-模的整体对偶定理同时——后者下文立刻就要谈到。可这并不妨碍“大家”现在把这个结果当作理所当然,而且尤其是(这就更理所当然了呢)丝毫从不提及某位含糊的无名者……}

此外,若$F$、$F'$是在$X$上处于对偶的晶体系数,Mebkhout证明:$F$与$F'$在$X$上的$\mathbb{C}$-向量空间“晶体上同调”复形\fnote{615}{(**)}{我提醒一句(参看“五张照片”,n\textsuperscript{o}~171(ix)):$F$在$X$上的(“绝对的”)晶体上同调定义为
\begin{align*}
R\Gamma_{\mathrm{cris}}(F)
 &\overset{\mathrm{def}}{=}\mathrm{RHom}_{\mathcal{D}}(\mathcal{O}_X,F)\\
 &\simeq R\Gamma\bigl(\mathrm{R\mathcal{H}om}_{\mathcal{D}}(\mathcal{O}_X,F)\bigr)
 =R\Gamma\bigl(\mathrm{DR}(F)\bigr).
\end{align*}
另一方面,下标 $!$ 表示真支撑上同调(这里是晶体上同调),即
\[
R\Gamma_{!}(F)\overset{\mathrm{def}}{=}
R\Gamma_{!}\mathrm{R\mathcal{H}om}_{\mathcal{D}}(\mathcal{O}_X,F).
\]}

\origpage{956}
\begin{quote}
$R\Gamma_{\mathrm{cris}}(F)_{!}$, \quad $R\Gamma_{\mathrm{cris}}(F')$
\end{quote}

作为拓扑向量空间的复形,由一个自然配对“处于对偶”,换句话说,人们有一个配对

\begin{quote}
$H^i_{\mathrm{cris}}(X,F)_{!} \times H^{-i}_{\mathrm{cris}}(X,F') \longrightarrow \mathbb{C}$
\end{quote}

(对一切整数$i$),而它\textbf{本身就是一条(拓扑向量空间的)对偶}。这条对偶定理“笼住”了离散系数情形下已知的(“绝对的”)对偶(Mebkhout称之为“Poincaré--Verdier对偶”),与凝聚系数情形下的对偶(Mebkhout称之为“Serre对偶”),合成为一种我愿称之为“Mebkhout对偶”的对偶——而他自己则称之为“Poincaré--Serre-Verdier对偶”\fnote{616}{(*)}{如我在别处(在“同伙”一注,n\textsuperscript{o}~63$'''$里)已经说到的:Mebkhout“不能不”向他的“恩人”Verdier脱帽致敬(自从后者把那条天赐的“好的参考文献”传给他之后)——凡有机会他都如此。然而两套对偶中任何一套(更遑论笼住它们的那一套)的本质观念,没有一件归功于Verdier。事实上,除了Poincaré对偶定理与Serre对偶定理的初始形态——它们当然曾是我的出发点——之外,一切本质观念都包含在我引入、并在凝聚与离散两个语境中长期发展起来的那套六种变换与双对偶形式主义里,而且是在孤独中发展的。

正是想着这一点,我去年在“受害者——或两种沉默”(n\textsuperscript{o}~78$'$)一注中写道:Mebkhout的那些“保护者”们“肯让他亲手抬着那口装着我的遗骸的棺材的一角”。当时我若还补一句才公道,那就是:Zoghman有那份勇气——尽管他清楚地感觉到俊髦世界里刮的是什么风——在每一篇文章里都明确重申他取法于我的思想,而不像大家那样:一边在(文字中)沉默地掠夺亡者,一边(在言谈中)摆出居高临下的姿态。

至于人们终于给那门我发展了多年、在完全孤独中发展的凝聚对偶理论起的名字“Serre对偶”,它格外有滋味(而Serre——他可没要求这么多——会比任何人都更欣赏!),因为Serre对我的对偶工作曾表现出全然的漠视,从而剥夺了我唯一可以指望的对话者来陪伴我的思索!我相信还可以说:那份漠视一直完好地保持到今天,连导出范畴的观念(以及其他无用细节……)也包括在内。}
% ---------------- p.746 ----------------
在我看来,这些是一个维度广阔的对偶纲领的最初几步,其中尤其包括(除别的之外(171(xi)))为特征零体上有限型概形上的 De Rham-Mebkhout 系数发展一套(双对偶与)六种运算的形式主义(且期待着更好)。鉴于Mebkhout不得不\origpage{957}工作的孤立条件与那种漠不关心的氛围,他当时不可能\textbf{为自己}发展一套完整的形式主义,像我在他所取法的两个语境中发展出的那种(171(xii))。在1972到1980这八年间他所提炼并证明的主要结果中(171(x)),在我看来、以我六十年代那个纲领的眼光看,最重要的当然是把那一类好的晶体系数——所谓“De Rham-Mebkhout”的——揭示出来的那个。巧的是,也正是这个结果,从1980年10月起,见识了最辉煌、甚至惊人的运气——而它却像从前的$\ell$-adic上同调、或特征$p$的晶体热门点心一样,只被当作一件\textbf{工具}侵占,从赋予它全部意义与力量的那幅图景上被撕走。

比Mebkhout的其他结果更甚,也一如我在发展双对偶与六种运算形式主义时那样,导出范畴的语言在这里是本质的——为的是揭示离散系数与凝聚系数之间那条简单而深刻的关系\fnote{617}{(*)}{(5月7日)确切地说,对一个完整 D-模(化归到零次的复形),好上帝函子一般结合出一个可构$\mathbb{C}$-向量空间层复形,它有不止一个非零的上同调层,反之亦然。最简单、最触目的例子是:在$X$上取一个除子$Y$,由此得包含$i:=X\setminus Y\hookrightarrow X$,再取 $i_{*}(\mathcal{O}_U)$ 中由沿 $Y$ 亚纯的函数构成的子层。Mebkhout早在1976年就得到的一个深刻结果(后来被吸入好上帝定理)是:那是一个完整正则的 D-模(在Mebkhout之前,谁也不曾想到把这个层看作一个 D-模,更不曾疑心它至少是凝聚的……)。它经好上帝函子的变换是$Ri_*(C_U)$,它至少在第0次与第1次有非零的上同调层。

那是Mebkhout哲学中为Deligne的进路所缺失的一面——后者得到的是一本词典,在可构 $\mathbb{C}$-向量空间层与$\operatorname{Coh}(\mathcal{O}_X)$(即 $\mathcal{O}_X$ 上凝聚模的范畴)的某些带分层的 pro-对象之间,而不必过渡到复形与导出范畴。(他终究还是留心让那些范畴出了场——那时我还在附近,而谁也想不到有一天那些范畴会被埋葬……)。那(至少乍看)是Deligne进路的一个优势,更贴近离散系数的直接几何直观——但那也是一个迹象,无疑说明他的进路较不深刻。我倾向于相信它还有自己的角色要扮演,但要与Mebkhout的视点“搭档”,而后者(我揣测)在某种意义上是对偶的。

(5月24日)关于这方面的细节,见子注“五张照片(晶体与 D-模)”(n\textsuperscript{o}~171(ix)),(c)部分,尤其是第1009页及以下。},它由好上帝定理所描述(别名:那位从不被点名的Mebkhout……)。于是,在那个\origpage{958}平展同调工具的创造近二十年之后(而今天所有人都理所当然地使用它,同时轻蔑那催生它的图景……),多亏这位隐身的身后弟子那(已成为“热门点心”的)结果,导出范畴的语言才突然获得平反(仿佛它从不曾被埋葬……),在聚光灯下、在人群的欢呼中——那人群是来向昨天的掘墓人欢呼的,而他们(谦逊地)扮演起新父亲的角色。但我又抢跑了……
\subsubsection*{c.~……与意外之喜}
\tocsubsub{c.~……与意外之喜}
\notehead{171(iii)}{}

Verdier多少充当着Mebkhout的“论文导师”;而Mebkhout此前七年的工作,完全是在孤独中进行的。他从来不曾对这个年轻人的工作产生丝毫兴趣——这年轻人显然既倔强又愚钝,不过是一个落伍的Grothendieck主义者,他以居高临下的姿态对待他。自1975年初次见面以来的四年里,他总共只给这个不知从哪里冒出来的无名小卒三次“会面”。我的其他上同调学生,没有一个
% ---------------- p.747 ----------------
肯屈尊对那位无名小卒的工作发生兴趣。他的工作对他们自己的研究的意义,他们全然看不出——而它却那么刺眼,连我这样一个十五年前就已从这一切“脱钩”的老朽都看得出来。他们太封闭在自己的掘墓-狂欢里、封闭在一种闷闷不乐的摇手柄日常里,无从领会一件新事物——它不带名片、不加修饰,只凭那些最简单、太显然不过的事物之力而来。他们早已埋葬自己的创造官能,只限于做名牌时兴货的消费者。然而后来,他们大大地向那位闯入者报复——他竟看到了他们以及所有人未曾看出的东西(而他们其实拥有一切,同他一样、甚至超出,足以看见、足以动手……)。但在这里我又抢跑了……

答辩在1979年2月15日举行,在一片普遍的冷淡中。Mebkhout把他的论文寄给所有他以为(对也罢、错也罢)对解析簇或代数簇的上同调感兴趣的数学家——当然从我的所有学生开始。在收到论文的全体人中,\textbf{没有一个}肯告知收到了寄件,或寄一个谢字。诚然,Mebkhout的论文比他某些文章更(在我看来)带着周围那逆境条件的印记——它在我看来卷帙浩繁、至少可以说不好进入,而那些不在圈内的人有理由不曾立刻咬钩。相\origpage{959}\dmark{}反,我发现Mebkhout对我口头解释他的哲学时说得完全清楚、且立刻令人信服;而没有任何理由说:他对Verdier(1976)、Berthelot(1978)、Illusie(1978)与Deligne(1979)所能作的解释,会比对我作的那些逊色。

那是在1979年6月的Bourbaki讨论班上,Deligne从Mebkhout口中得知那篇未被读过的论文中的“\textbf{Riemann--Hilbert对应}”。(那是Mebkhout给那个范畴等价(或那本“词典”)起的名字,前文刚谈到过。)显然,Verdier在那四年间从不曾想起就他的那位隐身学生的工作跟他多说一句——那工作的兴趣显然直到1981年反常层会议前后都全然逃过他(而在那次会议上,Deligne不得不负责给他解释那是怎么回事……)。而在Deligne那边,那只能立刻“对上暗号”——那正是他十年前自己撂下不管的那个问题的圆满而简练的解答!

在这种情势下,看似理所当然的反射(乃至我此刻都难以想象怎么能换一种做法……),是立刻向这位年轻的陌生人道贺——他终于找到了一个说实在深刻的、人们为之辛辛苦苦干了一整年、最后付诸流水的问题的谜底。世风已经大变……Deligne(当然依旧和和气气)只说了一句含糊的恭维(而那已足以叫天真的Zoghman心里发热——他确实不曾被惯坏,而且远未料到等待他的是什么):是的,他收到了他的论文,甚至读了它的引言,而他觉得那是些“美的数学”。对Zoghman那是黄道吉日!那当然是头一回(也同样是最后一回……)他有权得到一位大先生的恭维——那位人人都认识、人人都引用的先生……\fnote{618}{(*)}{(5月14日)况且那也是Mebkhout与Deligne之间唯一一次谈话的荣幸。

(6月7日)至于另一句恭维,在头一年(1978年6月)、且这次是出自Illusie之口,见“掠夺的空白授权——或高尚行当”一注(n\textsuperscript{o}~171$_{4}$),尤其是第1091页。}我说不出此刻以及随后一年里,Deligne脑子里关于这条他刚从一位无名者口中得知的卓越定理在想些什么。我揣度他必定会跟周围\origpage{960}\dmark{}的人谈起\fnote{619}{(*)}{(5月14日)细想之下,并且凭我从别处对Deligne的了解,我怀疑在他带着一个明确的主意和一个订好的计划去做之前,是否真的曾“跟周围的人谈起”。关于Deligne所扮演的十分特殊的角色、以及他让Beilinson与Bernstein这两个稻草人父亲所扮演的角色,见“父亲们的圆舞曲”一注(n\textsuperscript{o}~176$_{4}$)。(另见“冤大头买卖——或木偶戏”一注,n\textsuperscript{o}~172$^2$(e)。)}

% ---------------- p.748 ----------------
次年的10月,他把它私下通报给苏联数学家Beilinson与Bernstein\fnote{620}{(**)}{(5月14日)这是从Deligne给Mebkhout的一封信(1980年10月10日收到)中透露出来的。关于 Kazhdan-Lusztig 这一段的细节,见子注“黑手党”(n\textsuperscript{o}~171$_{2}$),(d)部分“总彩排”。},准是猜想到他们会用得上。同一年,事实上正是这个“对应”(当人们肯提名道姓时仍叫作“Riemann--Hilbert对应”,而Mebkhout的名字从不曾被说出)成了那场著名猜想\fnote{621}{(***)}{同一条猜想同时(相差几天)由 Brylinski-Kashiwara 独立证明——但其雷同程度却令人吃惊——带着同一主要配料、同一套掉包手法、以及对这个新事实的关键角色和它的作者的名字的同样抹杀。细节见已引的子注“黑手党”(n\textsuperscript{o}~171$_{2}$)(c)与(d)两部分。}——我几乎只知道它的名字,“Kazhdan-Lusztig 猜想”——的证明中那本质的配料、那此前一直缺失的\textbf{新事实}。那同时是代数簇上同调一场突然而壮观的新生的发令枪——它终于走出长达十余年的停滞(除了Deligne关于Weil诸猜想的工作)。这场意外的新生次年即见诸实行,那是1981年6月在主题“奇异空间上的分析与拓扑”下的吕米尼“事件”\fnote{622}{(****)}{会议的《论文集》出版于 Astérisque n\textsuperscript{o}~100 (1982)。而且那些《论文集》直到1983年12月才印出、1984年1月才问世,比卷上标注的日期晚了近两年。}。



\subsubsection*{d.~光荣之日}

\tocsubsub{d.~光荣之日}

\notehead{171(iv)}{}

关于这场“值得纪念的会议”,我请读者参看“不义——或归来的意义”(n\textsuperscript{o}~75)一注,以及随后那几注——它们写于发现之后还带着惊愕(这词不算过分)的热头上。这些注构成我所命名的葬礼第七行列:“会议——或Mebkhout层与反常性”。

我在这里只需提醒一句:在那场由 Bernard\textbf{Teissier}与 Jean-Louis\textbf{Verdier}签名的会议《论文集》引言中,那条著名的“Riemann--Hilbert对应”被作为会议的“Deus\origpage{961}\dmark{}ex machina”来介绍。在主要文章中也是如此——那篇文章(与所引引言一起)构成《论文集》的第一册,署名 A.A.~Beilinson、J.~Bernstein 与 P.~Deligne(而事实上由最后那位撰写、并在会议上报告,另两位共同作者缺席)。况且,头两位署名作者在上一年(1980年11月)就已由Mebkhout亲自(且独立于Deligne)告知了他那定理的来龙去脉——Mebkhout甚至为此特意去了一趟莫斯科\fnote{623}{(*)}{关于这段有教益的插曲,见已引的子注“黑手党”(n\textsuperscript{o}~171$_{2}$),(d)部分“总彩排(封神之前)”。}。Teissier也第一手地、而且早就知情——Verdier更不必说,他主持了Mebkhout的论文答辩……最后,我补充一句:临到末了才“急就章”地决定请Mebkhout作一个关于 D-模理论的报告(那门理论在场的人中除了他没有别人懂),Mebkhout就这样有机会把他谦逊地以 Riemann 与 Hilbert 之名命名的那个定理告知全场\fnote{624}{(**)}{(5月14日)关于这场古怪会议的与会者——那是一场盛大的“Grothendieck数学节”,却对亡故的祖先本人、以及对那位“有那份天赋让如此俊美的世界聚在一处”的隐身学生绝对沉默。作为“1970年前”的学生与会者,只有Deligne与Verdier,但已经足够牢牢占据舞台前排。蹊跷的是,Berthelot与Illusie(他们的工作可以说特别带着在此大张旗鼓地发掘出来的Mebkhout视点之缺席的印记)不在其列。相反,Contou-Carrère(“之后”的学生)却迷路闯了进去,庆幸自己被请来讲述他那套 Schubert 闭链的消解法。

我记得他回来后兴高采烈,完全认同于所有那些杰出而出名的人——他感到同他们平起平坐,而他们是来听他讲的,显然是感兴趣的,一点不假!他带着愧疚的神气向我谈起Mebkhout——后者曾对他打开心扉,却带着苦涩,而他也说不清为什么;可对孔图来说,无论如何日子显然是美的!

那是1981年6月。四个月后(为回应他向佩皮尼昂一个职位的单独应聘),那是结结实实的一记耳光,他当作一种屈辱和冒犯硬生生地受下。(关于这段插曲,见“棺材3——或那些稍稍太相对的雅可比簇”一注,n\textsuperscript{o}~95,尤其是第404--406页。写那注时我还不曾把它与Contou-Carrère参与那场辉煌会议的插曲联系起来——他大概被稍稍甩在了一边。)}
% ---------------- p.749 ----------------
而不对这件事的父权留下丝毫含糊(可想而知)。那件事有那份意想不到的天赋(无论对他还是对所有人)——让如此俊美的世界聚在一处。\origpage{962}\dmark{}而且,读者在会议的《论文集》中寻找Mebkhout那场报告的踪迹,也是白费。Verdier事后和颜悦色地向他解释:只有呈献\textbf{新}结果的文章才会收进《论文集》,而他博士论文里的结果已经两年以上了。读者在那部《论文集》中寻找关于这条著名定理之来历的任何文献引证、或哪怕最含糊一点的提示——同样白费:那条定理当然既不归功于 Riemann 也不归功于 Hilbert。读者也休想在其中找到Zoghman~Mebkhout这个名字。这个名字在第一册中不出现,无论在正文还是在文献目录中。在第二册里,它在文献目录中出现两次,是 Brylinski 与 Malgrange 的“暂停!”式引文(总不能说没人引过他!)——而那些引文与好上帝定理毫不相干——别名Riemann--Hilbert——别名Deligne(而尤其不别名Mebkhout)\fnote{625}{(*)}{(5月14日与26日)除了已提及的与会者,我具名得知的有 Brylinski、Malgrange 与 Laumon。三位都完全了解Mebkhout的工作,而他曾有机会对每一位详加说明,且远不止限于他在会议所作的那场报告。但这并不妨碍 Brylinski 与 Malgrange 在他们收进《论文集》的文章中——那文章本质地使用Mebkhout的思想与好上帝定理——既抹杀这些新思想与新工具之出现所扮演的关键角色,也抹杀它们的作者的名字。

至于 Laumon,他后来在一篇与 Katz 合写的文章中补了过来。而这位 N.~Katz 正是1973年以“SGA~7 行动”崭露头角的那一位——那件事在“一场攀登的若干插曲”一注(n\textsuperscript{o}~169$_3$,插曲2)中谈过。况且,他早在1979年就由Mebkhout亲自告知了他的结果(关于这一点,见“掠夺的空白授权”一注,n\textsuperscript{o}~171$_{4}$)。那是指《傅里叶变换与指数和的估计》一文(它同时也是 Laumon 的博士论文),那文章以预印本形式流通已有两年(我甚至承 Laumon 厚意得到一份)。这些作者们为$\ell$-adic 系数发展一种傅里叶变换,其模型是 Malgrange 1982年在 D-模情形引入的那一种(顺着那位无名者的工作的尾迹,且如所共料不提他的名字)。Mebkhout的工作构成 Malgrange 所发展的理论与 Laumon-Katz 二者的启发基础,一如它们曾是已引的 Beilinson-Bernstein-Deligne 那篇文章(关于他们误称的“反常层”)的启发基础。话虽如此,Laumon 与 Katz 也跟着大伙的调子走(文中与文献目录里都不提那位当值无名之辈——当然更不消提那位祖先……),循着 Deligne、Verdier、Berthelot、Illusie、Teissier、Malgrange、Brylinski、Kashiwara、Beilinson、Bernstein 定下的调门——我为这按字母排的次序道歉,总之这已有十二人直接而积极地卷入反常层会议那场辉煌的神秘化-骗局——还不算 Hotta 也在太平洋彼岸助了一臂之力,就成了十三!

Malgrange 在那篇文章中也并未被多引——显然有些作者小圈子互相狂引,却避开引隔壁的,哪怕他们一个比一个使劲地从人家那儿汲取。总之,一旦事关那位祖先或那位无名者,他们就全体一致。那往往确实是出色的数学——但我这老派人对这心态并非无动于衷,而它败掉了我阅读的胃口,乃至最终连做数学的胃口。至少不做他们那种数学。气味太难忍受……

我也曾对 J.L.~Verdier的文章《特化与温和单调层》投去一瞥,那文章发表在同一批《论文集》中。当然毫不意外,我在那里也看到了“Riemann--Hilbert对应”,却(在正文和文献目录里)丝毫不提那位他曾主持其论文答辩的无名者。他想必是不得已忘记了……那里还谈到一条平展 Riemann-Roch 定理(这名字我仿佛在哪儿听过……)——而我在 Laumon-Katz 的文章里也见过这个。既然双方都不曾提起某位亡者,我便对自己说:那条“定理”想必一定该归 Riemann 与 Roch 两位先生,正如那个在 SGA~5 的“技术性离题”与“废话”之中的特例一样(还不算那篇猜想的报告,被那位深谋远虑、机灵的“编辑”Illusie仿佛天意般地掏空了……)。

Mebkhout况且早在1977年就预感到他的哲学与傅里叶变换之间的联系——那时他正严格地独自一人对一门联结 D-模与离散系数的对偶瑜伽感兴趣(一如我自己当年对凝聚对偶、而后平展对偶形式主义的孤独)。那个“傅里叶变换”的直觉始终停留在含糊的状态——当时的语境对他继续走这条路并不比我1950年前后去拓展我的凝聚对偶理论、扩展到一门涵盖微分算子复形的理论(见第946页脚注(**))更容易。在Mebkhout为巴黎第七大学“奇异性讨论班”(1977--79)所作的报告“Poincaré对偶”之引言的第2页,有一条关于傅里叶变换的暗示。}

% ---------------- p.750 ----------------
\origpage{963}\dmark{}且回到那场有血有肉的会议,该相信,那些聚集在那里的杰出数学家,即便屈尊来听那位当值无名之辈给他们作的报告,也没有一个发觉:那位无名者作为自己的创造介绍给他们的“Riemann--Hilbert对应”,正是他们中最出色的一位早已辉煌地引入的那一个,作为他那辉煌报告的启发性的拱顶石——而那报告(据组织者Teissier与Verdier\fnote{626}{(*)}{那是从所引的会议《引言》中清楚透露出来的“隐含之见”,由 Teissier 与 Verdier 署名。}自己的估计)构成那场关于所谓(人们还真想知道为什么)“反常层”的辉煌会议的“钉子”。总之,该相信他们中没有一个曾惊讶:那位无名者的名字不曾在那场报告中被提起——那报告飞得那样高,固然无须为一个区区小辈费神;而两年半以后,随着《论文集》的问世(1984年初),那位无名者的名字也同样不在其中——既不在(已提及的)引言里,也不在Deligne及其合作者的相关文章里。而那篇文章\origpage{964}\dmark{}对这份对应的真正父权更不曾留下多少疑问,而主要作者兼报告者-变戏法者\fnote{627}{(*)}{关于我的朋友皮埃尔围绕那条从未被提名的定理之父权玩的花招,见去年那注“变戏法者”(n\textsuperscript{o}~75$'$)。},以他惯常的谦逊,也不曾提名道姓,连他那两位显赫先驱的名字也不曾。若说居然有人惊讶过,那他们直到今天也不曾让人知道——至少不曾让我知道,尤其不曾让那位提供笑料的主角知道:那位身后的、照理说完全不为人知的学生——一如既往——Zoghman~Mebkhout\fnote{628}{(**)}{(5月19日)关于我朋友Zoghman的那些坎坷遭遇——他天真地误入一个衣冠楚楚、满面和气的“硬汉”圈子——的详情,见以下几条子注:“一个愿景的萌发——或闯入者”“黑手党”“根”“掠夺的空白授权”(n\textsuperscript{o}~171$_{1}$至171$_{4}$)。}。

\subsubsection*{a1.~无用的细节}

\tocsubsub{a1.~无用的细节}

\notehead{171(v)}{}
\fnote{629}{(***)}{本注(分(a)(b)(c)三部分)源出“祖先”一注(n\textsuperscript{o}~171(i))的两条脚注——见第~944~页的脚注(**)与第~945~页的脚注(*)。}

\subsubsection*{(a) 上千页的一捆……}

\tocsubsub{(a) 上千页的一捆……}

\noindent (5月4日)连Serre也不曾免于这条规矩:他早已(如André~Weil一样)养成一种可憎的癖好,动辄宣判那些不合他心意的数学是“蠢话”。他与Weil那样的规格(人们或以为)本该使他们超越这类幼稚行径。而在此时此地(Deligne那“最后的二十页”姑且置而不论),恰恰是两三千页“Grothendieck式的蠢话”,最终证明了Weil猜想(还顺带证明了Weil与Serre做梦也不曾想到的不少别的东西)。这并没有引得Serre多一份谦逊:就在他陈述Deligne对这些猜想之最后一步的证明的那篇文字里(1974年2月的Bourbaki讨论班,第~446~讲),他恰恰借这个机会(措辞彬彬有礼,这自不待言),就SGA~4那“1583页”必定塞满的“无用细节”大加讥诮。在这种廉价的讥诮里,我辨认出的并非恶意,也并非不老实,而分明是一种无意识与轻率。他倒不辞劳苦,数出三大卷的页数(他自然留意不去读它们,其中的实质他也无从领会),又做了一道加法——好“优雅地”嘲弄一番。

% ---------------- p.751 ----------------
\origpage{965}\dmark{}但一切相互勾连:既有我当年对那些聪颖学生的迁就,也有Serre的“优雅”(而那是在葬礼已经发展了四年之后……)\fnote{630}{(*)}{(5月27日)对谈及Serre处加以接续的一段反思,见本注的(c)部分。}以及随后的一切。仅仅三年之后,人们几乎会在我的非学生Deligne笔下重新见到Serre那些说法本身、或它们的言外之意,还多一点恶意与无礼,连同那些被“剔除”的“无用细节”、那种“混乱状态”与那团“废话硬壳”——那正是这位Deligne学得手艺、觅得主要灵感的地方,而由他笔下的一篇寡淡摘要慈悲地“让人忘掉”。就这样,从迁就到随便、再到无礼,数学界在仅仅十年间,竟走到一种风习状态,连“体面”这简单的感情也仿佛消失了。

创造“那条猜想”所缺的新工具的,既不是Weil、也不是Serre,更不是Deligne;而是那个他们乐于讥讽的人——出于蓄意的无视或算计的恶意,效果并无太大差别。但我自己以万般的小心写了又写、让人写了又写,不倦地度过一月又一月、一年又一年,写出一部文本,它以其应有的全部宏阔,呈献那门语言与若干基础工具——为的是一幅崭新的、丰产的统一图景;我深知,而且明明白白地知道:在那被Serre、我的学生与一致的风尚置于脑后的1583页中,\textbf{没有一页}不曾被匠人称量又称量,没有一页不在其位上、不履行它的职能,而迄今写下的任何一页都不能取代它。这些页不是风尚的产物,也不是虚荣的产物——那种喜欢凌驾于他人之上的虚荣。它们是我的爱情的果实,是那些漫长而晦暗的劳动——它们准备着一场分娩。

至于我作品的这一部分,一如我那些已然进入共同财富的重大数学贡献,到今天为止,\textbf{没有谁}曾能重做我所做的事(仗着“蠢话”“无用细节”与“废话”),除非靠抄我(除了些不关紧要的变体)\fnote{631}{(**)}{(6月7日)我近来读到了 Fulton 的好书《Intersection Theory》(“Ergebnisse”丛书,Springer Verlag,1984年),发觉对 Riemann-Roch-Grothendieck 定理,应当算作例外。}。有些人抄(照抄、或在相近乃至崭新的语境里抄)并且说明(这种事如今已开始变得少之又少……);另一些人则扮演新父亲,并对他们无耻掠夺的作品、\origpage{966}\dmark{}对教会他们手艺的匠人摆出鄙夷的恩赐姿态。这种无礼之所以能滋生并铺展,只因为它找到了一份随时准备接待它的共识,而这首先来自那些(往往以其超常的身量)定调子的人。

% ---------------- p.752 ----------------
\subsubsection*{(b) 无所事事的机器……}

\tocsubsub{(b) 无所事事的机器……}

六种运算瑜伽是这幅在 SGA~4 与 SGA~5 两个讨论班中发展起来的“宏阔统一图景”的组成部分。我甚至要说,那场口头讨论班 SGA~5 的中心主题、或不如说它的“筋”与灵魂,就是这套瑜伽。因此Illusie才着意让它从屠杀版中消失(那部版子按其“照料”本该变成一册“技术性离题……”)。

在“祖先”一注(n\textsuperscript{o}~171(i),第945页)里我(未作进一步说明地)写道:六种运算这个富有力量的图景“给出过它自身力量的有力见证”。对我来说,这个力量最触目的具体标志——如果可以说的话——就在我们对平展上同调的掌握之中。为了在1963年达到这一掌握,那来自凝聚对偶的“六种运算”图景是我恒定的导线。况且我自认是世上唯一有资格就什么在这一工具的发展中具有决定性发言的人。

这里不言自明的是:在发现的进程里,所谓“启发性”的要素几乎总是决定性的。如果我谈一种观点或一幅图景的“力量”(那是与一条定理自身全然不同的另一种东西),那不能用严格技术的措辞衡量它。首要的是它那“暗示性的力量”——在发现的旅途中做谨慎而可靠的向导,在敏感的时刻向我们口授“那条”该引入的正确观念、“那个”该提炼并证明的正确陈述、“那门”仍待发展的理论。正是忘却了这样一幅向导图景(在把它埋葬之后),才使得代数簇上同调理论从六十年代那股强劲的冲劲,在我离开之后的年头里落得个混乱与萎靡。除了那桩“威望大问题”(即 Frobenius 特征值绝对值的问题)之外,一切本质的问题都被冥顽地回避了……

作为这幅图景(或在此情形,不如说这套形式主义)力量的另一标志,我举离散语境与凝聚语境中的Lefschetz--Verdier不动点公式。在这里,“六种运算”形式主义的角色既是\textbf{启发性}的(就这个意义而言:该公式被这\dmark{}形式主义\origpage{967}不可抗拒地暗示出来),也是\textbf{技术性}的(就这个意义而言:该形式主义也给出证明这条公式所必需且充分的手段)。诚然,鉴于葬礼,直到“相交上同调”与被命名为“反常”的层那场大冲刺之前(那时一部分形式主义被发掘出来,却不提匠人……),我所发展的那套上同调形式主义中只有极小一部分被使用过。但我自己清楚:凭着Weil诸猜想与 topos 直觉的无所不在,六种运算的图景是我1955--1970整个年间在上同调反思中的主要灵感来源\fnote{632}{(*)}{(5月15日)不言自明,这幅图景本身是在这一时期中渐渐取得形态的,出发点是我1955年那篇《论同调代数的若干问题》(载 Tohoku Math. Journal)中所含的最初萌芽。它在1963年达到完全成熟,随着平展上同调的突然启动。而这次启动(巧得出奇)几乎就发生在我引入那个“缺失的函子”$Rf_!$(真支正像)的那些日子里。但六种运算作为“富有力量的图景”与无所不在的导线,其角色(我相信)是在 SGA~5 讨论班里才变得充分自觉的。早在1966年,随着晶体上同调的启动,我就已清楚:第一个目标(超出那项“磨合”的有限纲领——它将在Berthelot的论文工作中完成)是达到一套针对“好的”晶体系数的六种运算(与双对偶)形式主义。只可惜须得由一个(被宣告亡故的)老朽从那为他预备的棺材里走出来,才使(在将近二十年之后,并借助一位当值无名之辈的想法,而他也被一同埋葬)这些“好系数”好歹终于\textbf{被定义}出来!至于它们的描述(尤其对$\mathbb{Z}$上有限型的概形),见《数学反思》第3卷(连同《收获与播种》第五亦即最后一部分)。}。这就是说,这幅图景的“力量”对我是显然的事,或不如说是我在数学家生涯的十五年间几乎日日经验过的现实。而且,这种经验在这些最近几周里又引人注目地得到再确认——一旦我重新接触那些“被遗弃的工地”:晶体系数与 De Rham 系数的工地、以及 Motif 的工地\fnote{633}{(**)}{(5月15日)关于“被弃置的工地”(或“荒废的”工地)这一意象,见葬礼仪式的第6幕(注~176$'$、177、178),尤其是三条所引注中的最后一条。在撰写《收获与播种》之余,我不过是在这里那里花上几个小时,处理 De Rham 晶体系数的问题与 Motif 的问题,便足以看到一个针对前者的令人信服的定义、以及至少一个构造后者的原理,在$\mathbb{Z}$上有限型概形这一关键语境中出现。(与前一脚注的评注相比较。)}

\origpage{968}\dmark{}这份关于某种富有力量的图景之力量的全部“主观”经验,也有一个“客观”的意义,那是不可一挥衣袖就抹掉的。只要愿意记住这样一件事,这个意义就显现出来:(除了少数例外)关于“抽象”代数簇与概形的上同调的那些主要观念与概念(今天所有人都使用它们,仿佛它们可以追溯到亚当与夏娃\fnote{634}{(*)}{关于这种数学成品“使用者”(或“消费者”)的心态——他忘记了(若他曾知道过……)创造是什么事——也关于亚当与夏娃以及上帝,我请读者参看去年那些注“一种不义与无力的感觉……”(n\textsuperscript{o}~44$''$)与“当值无名之辈与好上帝定理”(n\textsuperscript{o}~48$'$)。亦见“教学的失败(2)——或创造与浮夸”(注 n\textsuperscript{o}~44$'$)。}),正是由我——而不是别人——在1955--1970这同一时期里提炼出来的。(不言自明,我在此把我出发的 FAC 与Weil诸猜想都排除在外。)

这个伟大的时代,就数学而言(且据我迄今所能看到的),通向一种阴郁、呆滞与贫乏。其深层原因绝不位于技术层面。这种贫乏的标志之一,是一幅为激励并哺养宏大设计而生的有力图景,竟被埋葬、或交由嘲笑——而动手的正是它的保管人与第一批受益者。而另一个标志是:Deligne也罢、Verdier也罢、Berthelot也罢、Illusie也罢,尽管地位与威望带来一切便利、才华横溢、经验老到,没有一个做出一件在 De Rham 系数上非做不可的作品——那本是他们自己的研究(与被否决的图景……)的正统后裔;他们甚至没有本领在面对这件作品时,认出它的发明性与丰饶。而正是本着这同一种精神(因为一切又相互勾连……),一旦那件出自新作品之一的工具的意义终于被承认,他们就赶忙把它据为己有,甚至并不理解它,并把那位锻造了它的无名匠人——连同祖先——一道埋葬……
% ---------------- p.753 ----------------
\subsubsection*{(c) 像什么也不像的东西……——或枯萎}

\tocsubsub{(c) 像什么也不像的东西……——或枯萎}


(5月27日)\fnote{635}{(**)}{本注“无用的细节”的第三部分源出第一部分的一条脚注。见第965页脚注(*)中的转引。}我关于Serre所说的话是自然而然写出来的,并且出自对有关他的事物的一种知觉——那知觉想必是在最近几个星期或几个月里在我心中形成的。\dmark{}\origpage{969}然而在写这几行时,曾有一份残存的犹疑或困惑、或保留,针对我刚写下的东西。我那话里总而言之让人听出:Serre在这一场合上竟缺少“优雅”!

实情是:我认识Serre快三十年了,他一直就是我心目中“\textbf{优雅}”的化身。想必不是只有我这样看他。那是一种优雅,既在他的工作与作品中,也在他对他人的关系中;而那绝不是纯形式的。它包含工作中一丝不苟的诚实,以及对他人同等的诚实要求。不止一次,我注意到他在某个较不谨慎的同事面前判断的锐利——那位同事一心要把一个恼人的困难弄浑(为的是不必承认他不知该如何克服它),或把他自己的某个错误含混过去……因此这种优雅也包含一种\textbf{严谨},对他自己与对他人同样。

正是这一切,对我而言与Serre这个人不可分,它们必定介入了我刚提到的那份“保留”——在一份对事物的知觉不期然地抢在熟悉的知觉之前的时候。我无意为了其中一个的“好处”而排斥另一个。两者都有东西要教我,都是一个复杂现实的不同侧面,而那现实当然绝非静止。该由我把它们相互定位,以达致对一个与我有一段过去、有同情与尊重之情的人的细致入微的领会。

然而我刚谈到的这种“严谨”,并不曾延伸到Serre与数学和数学家的关系上的一切。我前面曾提到一种“不自觉”或“轻率”,我本可以同样称之为一种“\textbf{封闭}”。它与那种“谨慎与谦逊”的态度形成对照——我在大多数长辈身上遇到那种态度,他们同Serre一样在我起步时善待我,而有时(如他那样)还带热情。关于这一点我在下文(在“自由……”一注,n\textsuperscript{o}~171(vii)里)有更多表述,我在那里确认:那种态度曾是“熏陶过接纳我的那个环境的尊重气氛……”的一部分。

% ---------------- p.753（接上段） / p.754 ----------------
我在某些场合在Serre身上见到的“封闭”,并非始于昨日。我在五十年代下半期就感觉到它的最初征兆。我相信,它大大限制了他自六十年代起作品的深度\origpage{970}\dmark{}与广度。我感到这种“封闭”与以下各事之间有着联系:一些不同于他本人的数学进路;以及在他内心逐渐发展起来的一个蓄意之举——把他对数学事物与对数学本身的领会,封闭进一种(或如我真想写的,一副“眼罩”)纯技术的或技术主义的观点之中,从而对一切近似于一种\textbf{视野}的东西关闭门户;对某种超越了有形的、即时的、\textbf{可证明的}陈述(或陈述的集合)的东西关闭门户,或者(退一步说)超越了采取“硬邦邦”的猜想之形式的东西——那些猜想轮廓分明,总而言之是“封闭”的(只不过还剩下要把它们证明……)。回过头看,我觉得他终究是把他创造力的那一面推向了极致:那纯粹“阳”的一面、“超阳”的一面、“\textbf{雄性}”的一面。凭借他在同代数学家中、以及在随后两三代人身上所享有的非凡影响力,在我看来,Serre大大促成了我所见到的、在七十年代与八十年代肆虐的那种过度技术主义精神——那是如今唯一还被容忍的精神,而数学的其他任何进路都已成为众人嘲弄的对象。

借用 C.L.~Siegel的一个说法,我们今天目睹的是一种非同寻常的“扁平化”(Verflachung)\fnote{636}{(*)}{我从Serre刚刚收到的一封信中借来这个(德文)表达式。该表达式取自 C.L.~Siegel为 Hecke 著作集所写的序言。Serre在信的末尾引了 C-L.~Siegel的这一印象,随即补上一句:“那时这不公道,而在我看来,如今就更不公道了。”可是它在我心里“咔哒”作响,并且持续地起着作用。我那段关于Serre与我之间关系的简短反思,大概正出于此。况且我相信:Serre之所以引用Siegel,是因为从某种意义上说,这一出自我们时代一位大数学家之口的印象,必定早已在他心里起作用;那大概是数学的“玫瑰人生”中的一个不谐之音。这样的不谐之音当然不止一个,但这一个看来不那么容易排遣……“Flach”在德文中意为“平的”、“没有深度”;“Verflachung”指通向这样一种“平”的状态的过程,或指这样一个刚刚完成的过程的结局。在正文里,我着力跟随这个极富表现力、可惜无法原样翻译的词在我身上唤起的联想。当然,我完全不知道我感知此事的方式,与Siegel的感知是否有丝毫重合——Serre所引的Siegel原文,我不曾读过。},一种数学思想的“压平”、“收缩”:它被剥夺了一个维度——那个洞见的维度,那个梦境与神秘的维度,那个深度的维度——而在我看来,在此之前它与这个维度从未失去联系。我把这感受为一种\textbf{枯萎},一种思想的\textbf{硬化}:\origpage{971}它失去了活的柔韧、失去了滋养的品质——变成纯粹的\textbf{工具},\dmark{}僵直而冰冷,用来完美地执行那些“费力挣来的”任务、摆上公开拍卖的任务\fnote{637}{(*)}{“公开拍卖”这个意象,大概是从“招标”公告(sic)那里得来的暗示:“法国国家科学研究中心通讯”以及我作为这份可敬机构里一名初出茅庐的研究人员,定期收到的其他纸页上,满是这类公告。这一行话,连同许许多多别的征兆,表明这种发现工作的“压平”绝不限于我昔日熟知的那个圈子,也不限于数学这门科学。我至今还没见到纯数学的招标公告,但那不会太久——我很容易想象我当年的某些朋友或学生,一本正经地坐在包着软垫的门后,坐在某个名称令人望而生畏的委员会里,决定哪些“研究方向”应宣布为优先、哪些“进路策略”应予提倡、哪些“名列胜者”的课题组的“投标”宜予“采纳”进入“预选”,甚或授以头奖:由监护部颁发的官方津贴,经“主管委员会”出具肯定意见后每两年续发一次……},当目的感与方向感、当把这些任务本身视为一个浩大的整全之部分的感觉,都已被所有人遗忘。有一种深层的硬化,被一种狂热的肥大所掩盖。

思想的这种失衡,是另一种更为本质的失衡的征兆之一,也是更为深层的空虚与匮乏的征兆。思想的这种干枯在过去二十年里蔓延并扎根,与此同时,人与人之间交往中那些习见的细腻与敬重的形式也在遭受侵蚀——这不是偶然。同样不是偶然的是:这阵蔑视之风曾经刮起,而我终于感到了它的气息,它与一种或多或少普遍的腐败相伴而来;一年多以来,我一直在绕着这种腐败作巡视,未能穷尽。

% ---------------- p.755 ----------------
Serre直到今天,对这种四面八方环绕着他的腐败依然毫无知觉。我却记得他嗅觉灵敏。然而,光有灵敏的鼻子不够,还得用上它,听取它要告诉我们的东西,即使它向我们报出的气味足以让我们不安;甚至足以让我们忧虑,当它把我们自己牵连在内的时候。我很清楚,Serre同我一样,不会想到与狼同嚎,不会掠夺,不会搞阴谋,不会揭人短处——在“人人”都掠夺、都搞阴谋、都揭短的地方。他确实什么都不做——他只是捂住鼻子(这样一来少了一只手……),装作什么也没闻到。

而他绝非孤身一人——在我隐居中气味仍能传到的那个我们共同的世界里,那些曾是我朋友的人,一个也没有——一个也还没有跟我谈起过——\origpage{972}哪怕只是话里带一点——他\dmark{}闻到的、令他不适的气味。我那些同行中,无疑仍有许多人正直地从事着数学家这一值得敬重的行当。但在那些坐在前排的人中间,我\textbf{一个也不}认识——谁有这份质朴,去相信他自己健全官能(此处即嗅觉)的见证,而不是捂住鼻子,以免对自己说出:这里有什么东西发臭——也许该去看看吧……

但我回到Serre其人与我本人,回到我在他身上感到的那种“封闭”——它不知何时出现,并随着年月日益加重。我相信,他作品中最富多产性的那一部分,那最深刻地影响了其时数学的那一部分,位于早年,位于这种封闭出现之前,或者至少,位于它对Serre与数学、与数学家的关系取得决定性掌控之前。也正是那些年,那些五十年代,他与我的接触对我来说最为多产;也正是在那些年,Serre在我身上扮演了那种“引爆者”的角色,给我的作品某些最具决定性的推动。也是在那些年,我心中诞生并成长起一幅浩大的图景,它在那些年激励并滋养了我的作品,直到今天依然如此。我可以明明白白地、清醒地说:如果在“我”之外曾有人对这幅图景的萌发有所参与,那就是他,Serre,在那些年。而之所以可能如此,只是因为在这些多产而决定性的年头里,他心中有一种对数学事物本身(就其自身而言)的敞开,包括对那些尚未落入即刻之把握的事物,对那些起初似乎不愿被已经成形的语言之网眼所框定的事物——那些也许需要多年晦暗而耐心的劳作,甚至需要整整一生,才能凝聚为有形的实体,才能让一个\textbf{躯体}的四肢、形态与轮廓显现出来,活生生、强健有力,见证着在熟悉的已知者之语境中,一个\textbf{新的存在者}出其不意地出现。

我相信,在我认识Serre的头些年,直到五十年代末前后,他一直保持着一种感受力,去感受“创造”这种摸不着、微妙的事物,去感受那些为一次分娩做准备的、卑微的劳作。我相信,在某个时刻,他曾能感到一幅图景的萌发,以及赋予它形体的语言——如灵魂或精神,如躯体……那时有一种无需言说的温暖,有一种审慎而有效的随时可用性,在他能够协助一场并非他本人的、艰辛而紧张的劳作之处;而通过一份同情与一份期待,他毕竟参与了那场劳作。

\origpage{973}\dmark{}我无法说出,他在我们共同激情的那个层面上,这种鲜活究竟于何时、以何种方式钝化下去,让位于我方才试着描画的另一种东西。大约在六十年代初,甚至更早,他就不复感知森林,而只肯看见他合意的这一棵或那一棵树。其余的东西没有存在的理由。我想,看到我那样专注于不倦地开垦大片无形之地,耐心地在那里栽种这一切尚不像任何东西的事物,一副已经看见一片繁花森林的样子,他仅仅是恼火罢了\fnote{638}{(*)}{(6月17日)在我“巡访工地——或工具与视野”一注(n\textsuperscript{o}~178)里逐一巡视的六个“工地”之中,古时只有一个工地(“Motif”工地)曾使Serre多少发生兴趣——而且……当我最近在附言(P.S.)中不带评述地写信告诉他,自认为已掌握在 $\mathbb{Z}$ 上有限型概形之上构造 Motif 范畴之形式的原理时,他在回信中丝毫不提。看来这些“Grothendieck式的数学”再也无法使他或冷或热了……}。

% ---------------- p.756 ----------------
这并不曾阻止我继续开垦、继续栽种与补种、继续修剪、继续重新开垦与重新栽种——也不曾阻止我们一如既往地成为伙伴,依然花上几个小时又一个小时讨论数学(多半是在电话上)。当我有一个轮廓分明的问题,而且是一个并未被禁绝的问题,我通常尤其向他求教,指望他或许会有所烛照——事实上,他常常确实有所烛照。我继续从他那里学到许多东西,而他想必也从我这里学到那些能够令他感兴趣的东西。这比礼尚往来或互通有无更好——总有一份共同的激情把我们联结在一起,有火,有火花。

可他已经不再是我的灵感的源泉。这源泉如今唯有在我自身之中\fnote{639}{(**)}{(6月12日)关于这段对Serre与我之间关系之反思的续篇,见“家庭相册”一注(n\textsuperscript{o}~173)c.部(“众人之中的那一个——或认同”),6月11日所作,以及 d.、e. 部。}。

\subsubsection*{a2.~古怪的问题}

\tocsubsub{a2.~古怪的问题}

\notehead{171(vi)}{}

(5月5日)\fnote{640}{(***)}{本注源出“祖先”一注(n\textsuperscript{o}~171(i))的一条脚注——见第946页脚注(*)。}我在这里的记忆有些模糊,并在随后几个星期里重新清晰起来,那时我有机会就这些问题重新作些许接触。实际上,我心中有两个截然不同的问题,一个完全确切,另一个相当含糊。

第一个问题关乎这样一个需要:为那些尚待确切定义的“De Rham 系数”,廓清出一套完整的“六种变换”(variances)理论。我的晶体观念,无论在特征 $p>0$ 中\origpage{974}还是在零特征中,都提供一个极为确切开端——人们已经预先知道,\dmark{}什么将取代那些 $\ell$-adic“局部系统”(或“常值扭曲层”,即 Betti 情形,在超越语境中),而必须设法定义“带奇点的系数”,当然是在导出范畴的精神之中\fnote{641}{(*)}{同样清楚的是,当基域为 $\mathbb{C}$ 时,人们想要的是一类等价于如下范畴的东西:其上同调层为代数可构的 $\mathbb{C}$-向量空间层的复形之范畴。这一极为确切的提示暗示出:通过剥丝抽茧,关键的问题是把晶体局部系统(在一个不一定闭的子概形上)附着为一个周围环境概形上的晶体层。这基本上是 Deligne 在1969年所做的事,只是事实证明,人们找到的不是一个晶体层,而是一个\textbf{晶体 pro-层},这在当时便是一个重要的新观念(而且只要你肯一看,它就“显然”)。然而,系统地使用 pro-对象,将要求相当可观的基础性工作,Jouanolou 为其论文(关于 $\ell$-adic 系数)所做的工作给出了一个先声。那里本须再次捋起袖子……\par 从 Deligne 的观点与我的观点来看,Mebkhout 通过 D-模的新进路,从此就归结为把一个晶体 pro-层替换为一个晶体 ind-层(借助\textbf{通常的}凝聚对偶函子 $R\mathcal{H}\mathit{om}_X(-,\mathcal{O}_X)$),并\textbf{取归纳极限},以求得一个通常的晶体层,即(现在假设 $X$ 在零特征域上光滑)一个 $\mathcal{D}_X$-模。于是,Mebkhout 在1972至1976年间(从一个相反的“端点”出发,参见“三个路标”一注 n\textsuperscript{o}~171(x))所确立的那个出乎意料的“奇迹”,是这个 D-模是\textbf{凝聚的}(更确切地说,其同调层是凝聚的)。另一个同样出乎意料的奇迹是:人们能够以一些简单的、相对于Grothendieck的晶体观点而言性质全新的条件,来刻画这样得到的 D-模(或不如说 D-模的复形)——即那些“微观局部”的完整性条件,加上一个由 Mebkhout 引入并在此期间已为人所熟悉的“正则性”条件。\par (5月26日)关于 Mebkhout 的 De Rham 系数与 Deligne 的 De Rham 系数之间的对偶关系之细节,见“五张照片(晶体与 D-模)”一注(n\textsuperscript{o}~171(ix))c. 部。关于必须以“凝聚 pro-模形式的晶体”取代 Deligne 的“pro-凝聚模”观点,以及关于(尚未证明的)以不加修饰的晶体层(通过取射影极限)取代笨重的 pro-对象(晶体式或分层式)观点之可能性,见同一注的 c.、d. 部。}
% ---------------- p.757（本段承上） ----------------
,于是所缺的,便是晶体复形的一条恰当的“有限性”条件。在零特征中,给出答案的是“D-凝聚性”(不论我还是我的学生中谁都不曾想过——多么简单自然的主意!),再加上更微妙的完整性(condition d'holonomie)与正则性条件——正如“好上帝”即Mebkhout的哲学(在晶体瑜伽启动十二年之后)所教会我们的。我怀着好奇等待着我某个前学生终于动起来(不提那位当值无名之辈,也不提那位祖先——这是说好的……),以廓清特征 $p>0$ 中\origpage{975}对应的条件,或不如说,廓清在刚性解析语境的特征\dmark{}零中对应的条件。迟做总比不做好……\fnote{642}{(*)}{(5月26日)自这些行写就以来,作为我一番努力的一份意外成果——我本想为“登峰造极”写一份足以传世的事迹——我竟然(几乎不费吹灰之力)廓清了在我看来正是\textbf{恰当的}De Rham 系数定义,至少对 $\mathrm{Spec}\,\mathbb{Z}$ 上有限型概形如此(在我看来,这是最为至关紧要的情形)。当然,与我1966年的观念相比,本质性的新成分是那“含糊的无名氏”的哲学,我在这里(像所有人一样)将克制自己不去点名。\par 我所设想的针对 $\mathrm{Spec}(\mathbb{Z})$ 上有限型概形的方法,应当同样给出任意域(不论特征零与否)上有限型概形的好 De Rham 系数(Mebkhout式或Deligne式,任择其一)。我打算在《数学反思》第三卷“De Rham 系数”一部分里粗描这条进路,连同其他种种“技术性离题”,我的学生尽可前来抄录……}

我自己没有在六十年代继续追问这个问题,我另有足够多的事情要做,而且想着有Berthelot和Deligne在列,这事便落在可靠的人手中(可见人有时会看错……)。Deligne在1969/70年的工作原则上已提供一个零特征中的答案,倘若Deligne把这项工作做到底,大概会令我满意。

\noindent 可是在我心中,这样一套假想的 De Rham 系数理论,即使它能把“离散的”(以一种晶体上同调的形态)上同调与“凝聚的”上同调联系起来,也并\textbf{不}“统摄”凝聚对偶理论。于是我看不出一个Zariski凝聚层如何定义一个“包络晶体”(cristal enveloppant)\fnote{643}{(**)}{(5月26日)也许最好取那个包络“余晶体”(co-cristal)(见注171(ix)B部,其中对余晶体概念有所暗示)。我大概会在前一页边注所许诺的讲稿里回到这个问题。}(注意:在 D-模的语言里,这就是纯量环的扩张 $\mathcal{O}_X\to\mathcal{D}_X$,至少当 $X$ 光滑时如此)——而即使我看见了,所得的晶体(只取 $\mathcal{F}=\mathcal{O}_X$ 就已如此,它给出晶体 $\mathcal{D}_X$)也\textbf{并非} De Rham 型的。然而我暗自思忖:在一个复解析空间 $X$ 上,凝聚对偶(例如以Serre的形式,设 $X$ 光滑、系数局部自由)难道不能作为“离散”对偶的一个“特殊情形”得到吗——这“离散”对偶是Verdier按照 étale 理论的模型发展出来的。就这么着,它看起来有点古怪,并且立刻掀起一大簇问题:如何“以离散的语言”解释对偶模(最高次微分形式)$\omega_X$ 的作用,又如何处理那些怪异的病理现象,它们可在“离散的”对偶中没有对应物?

\origpage{976}\dmark{}是Mebkhout第一个(而且看来直到今天除我之外也是唯一一个人)领会到两种对偶之间确有一层深切的联系,但这联系并不\textbf{表述为}其中之一“统摄”另一个,而是表述为找出第三种对偶理论\fnote{644}{(*)}{关于这“第三种统摄另两种的对偶理论”的细节,见“……作品”一注(n\textsuperscript{o}~171(ii))。}——即 $\mathcal{D}_X$-模(或 $X$ 上“晶体”)的对偶理论,它“统摄”另两者,并且(除此之外)还在“离散”一边自我限于那些同调层为\textbf{解析可构}的 $\mathbb{C}$-向量空间复形。对我来说毫无疑义:这正是对那个“含糊的问题”(而且有点不着边际……)的“正确答案”,而我却从未有机会把它向我那位身后学生提出……

(5月15日)撰写“登峰造极”同时成了一个意外的机会,使我对Mebkhout的工作、对他引入簇的上同调研究的 D-模瑜伽多少熟悉起来。一路行来,这也打捞起一些沉没的记忆。我尤其注意到:早在五十年代末或六十年代初,我就比我十天以前动笔写本注开篇(“古怪的问题”)时自以为的程度,更接近“Mebkhout的哲学”。在任意基概形上的真光滑概形的框架里,我手中曾有一条对偶陈述(以一个相对微分算子复形与其“伴生”复形表述),它“统摄”凝聚对偶与 De Rham 上同调的对偶。技术上讲,这大略相当于Mebkhout对偶定理之代数版本的等价物(在复解析语境中,“……作品”一注(n\textsuperscript{o}~171(ii))已论及)。然而我那条对偶陈述并不使我满意,我不曾想过要发表它,甚至不曾想过为它做宣传,因为在所述形态下,它在我看来太靠近Serre对偶定理(在任意基概形上相对化——这是理所当然的),它便是后者一个多少直接的推论。要得到一条令我满意的陈述,我便须知道如何用微分算子复形作出一个“导出范畴”,以便仿照此前若干年间廓清的凝聚对偶理论之模型,以这些范畴的对象来表述一条内在的对偶陈述。

% ---------------- p.758 ----------------
\noindent 于是所缺的,便是\origpage{977}微分算子复形之间一个(微分的)态射的恰当“拟同构”概念,\dmark{}以便构成一个导出范畴(将这些拟同构形式地逆反)。显然,惯常的定义(经由相应的上同调层)在代数框架里不可用(在超越框架里恐怕同样不可用\fnote{645}{(*)}{我在这里弄错了。Mebkhout向我保证:对微分算子复形之间的一个(微分的)同态而言,它是拟同构(按相应 $\mathbb{C}$-向量空间层复形的朴素意义)当且仅当相应于关联 D-模复形的同态是一个拟同构。事实上(经由映射柱)这等价于说:一个微分算子复形按朴素意义拟零,当且仅当关联的 D-模复形拟零——这似乎是人人皆知的事(至少Mebkhout知道,他在那篇取之不竭的论文里给出了证明……)。})。转向相应的 D-模复形,如今为我昔日的困惑给出一个美妙简洁的答案!

既然看不出有什么现成的定义可给拟同构概念,我便没有试着去弄清楚这事存在与否,以及那里是否确然有一个非凡的导出范畴。那正是在我是唯一对导出范畴感兴趣的时候——那些由凝聚模及其之间的\textbf{线性}态射构成的导出范畴(远为不那么精微)……我并未清楚地感到:这“须廓清一个拟同构概念”的问题(同样有点含糊,若不说是荒诞)触及一个多产的奥秘,而这个奥秘容许一把孩子般简单的“钥匙”!并且有一类非凡的“系数”范畴,只等人来定义它。为此大概须得我的反思在这样一种气氛中继续推进:这反思在那里至少遇到一点兴趣与回响,哪怕只在\textbf{一个}既是局中人便算的对话者那里!

是 De Rham 上同调把我的注意力引向这个当然不言自明的事实:设 $X$ 为域 $k$ 上(比如说)一个代数簇,凝聚层的全局上同调空间不仅是 $\mathcal{O}_X$-线性同态的“函子”,甚至还是\textbf{所有} $k$-向量空间层同态的函子,特别地对微分算子也是如此。正是这个观察促使我萌芽了一份反思,关于一种“凝聚的”(或“拟凝聚的”)对偶理论,其中层之间的“态射”将是微分算子,而不必是线性的。这份反思夭折了,如上所述,而且到这样的地步:它甚至不曾留在\origpage{978}我记忆一角,作为(诸多事项之中的)一件终须\dmark{}有一天弄清楚的事——它沉没了(我想)在完全的遗忘之中,直到几天以前。甚至我约1966年时对晶体的零星反思,也不曾把它打捞回我的记忆,据我所知。然而,这晶体式的反思,在我当时毫无察觉之下(因为当时连这个问题本身都不记得!),将从1966年起为我提供\textbf{另一}把钥匙,某种意义上与Mebkhout那把“对偶”,为我昔日的困惑而设——经由与一个微分算子复形相关联的无穷阶主部复形。我在昨天写的一条页边注里提到了这一点(第946页注(**)),并且我打算在《数学反思》第三卷发展“系数类型”瑜伽的那一部分里详尽回到这里,尤其给出一个成形的定义,定义我所推测的“那些”
% ---------------- p.759 ----------------
好的德拉姆系数(Mebkhout式的或Deligne式的,任择其一),(比如说,在一个 $\mathbb{Z}$ 上有限型概形上)。

\noindent 技术上如此,甚至“心理上”也如此(就彼时早已提出的问题、以及赋予它们力量与生命的那幅总括图景而言),从六十年代后半期起,一切都已就绪,足以廓清这套德拉姆系数的定义。Deligne在我之后曾与正确的概念只差半步,而他不可能不廓清它——倘若那股他授予其对自己人生与作品之绝对权力的力量,不曾为他在这一道路上的反思早早地、斩钉截铁地画上句号……\fnote{646}{(*)}{关于此事,见“……与羁绊”子注(n\textsuperscript{o}~171(viii))。}

发现,并不是双臂倒挽,以锤或以槌猛砸一颗钉子、一柄凿子或一块钢楔。发现,首先是要懂得\textbf{倾听},以一份敬畏与一份全神贯注,倾听事物的声音。新事物并不如一道炫目的光柱,已然成形地从钻石中迸射而出;它也不出自一台机床,无论那机床如何精密、如何强有力。它并不轰然自报家门,披挂着它的贵族世系文书:我是这个,我是那个……它是一个卑微而脆弱的东西,一个娇嫩而活生生的东西,也许是一枚卑微的橡实,从中将长出一株橡树(倘若节令于它相宜……),或是一粒种子,将生出一茎,茎上生出一朵花。它不在聚光灯下诞生,甚至不在日光之下。它不是已知者的果实。它的母亲是夜与昏暗,是那捉不住、无轮廓的雾气——是那逃逸于想要\origpage{979}框定它的言辞之外的预感,是那仍在寻找自身的大古怪问题,或是\dmark{}那份如此含糊、如此飘忽却又如此真实的不自足,伴着那份说不出的(也赖不掉的……)感觉:有什么东西不对劲,有什么东西歪歪斜斜,石头底下藏着鳗鱼……

当我们懂得谦卑地倾听这些低声向我们诉说的话语,并执着地、狂热地追随它们那飘忽的讯息,那么——在历经晦暗而摸索的劳作,或许浑浊而无外观之后——雾气骤然具形、凝缩,成为\textbf{实体},坚固而可触,成为\textbf{形式},可见而清晰。在这孤独的一刻,在全神贯注与寂静之中,那新事物——夜与雾之女——显现了……

\subsubsection*{a3.~自由……}
\tocsubsub{a3.~自由……}
\notehead{171(vii)}{}

(5月4日)\fnote{647}{(*)}{本子注源出“祖先”一注(n\textsuperscript{o}~171(i))的一条页边注——见第945页注(**)。}我并不自称要摆出“成熟之人”或“智者”的姿态,周围尽是不成熟、不负责任的同类——我想,在《收获与播种》的书页里\fnote{648}{(**)}{(5月26日)我甚至可以说:若说撰写《收获与播种》在某一点上对我有所揭示,那揭示的恰是一种“未成熟”的状态,一份“智慧”的欠缺,而断不是它的反面。这也许是全部发现中最出人意料的一个,也是最关键的一个(就其即刻的意涵而言):我对某段过往、对我作为数学家的作品,依恋竟如此之强。这份依恋当初尚以较为克制的外形,在去年3月末向我显现,那是在《愚妄与更新》最后一条注“一段过往的重量”(n\textsuperscript{o}~50)的反思之中。正是直面“葬礼仪式”那粗暴的现实——尤其在其蓄意轻蔑与暴力的诸面相——在我身上启动了强大的自我中心的防卫反射。它们同时也向我揭示出那些将我系于一份过往的联结之强,而我从前竟曾相信那份过往已经离我而去。在过去一年之中,这些联结似乎获得了一份新的活力,而且常常(尤以近来为甚)我感到它们是一种\textbf{重量},说真的,一份累垮人的重量——正如其他曾经压在我身上、而后又化解了的重量……},从我这个人浮现出的形象大概并非如此。然而,至少在我与数学的关系中,我相信可以这样说:在我一生之中,始终保有一份货真价实的单纯\fnote{649}{(***)}{(5月16日)这里大概须对我“私人猎场”的一种占有姿态破个例外——我在《愚妄与更新》“好斗的数学”一节(n\textsuperscript{o}~40)里恰恰点到了它。这些“好斗的”心性本该促使我贬低他人的想法,每当那些想法我这边早已知晓之时。因此可以说(与我在正文中所断言的相反):在这些情形里,我的虚荣确确实实干扰了“我健全的判断”,并且往往倾向于在某一场合激起一种令人气馁的态度,而那时一句善意的鼓励本才是应有的姿态。不过我觉得,这样的处境在我数学家的一生中是例外的,而且并未曾对我的数学创造力构成羁绊。},\origpage{980}同时保有一份对我\dmark{}原初本性的忠诚。我的虚荣,侵入我生命之广不亚于侵入我任何一个同行的生命,然而(就我所记得的)它几乎不曾干扰我健全的判断,也不曾干扰我数学上的嗅觉\fnote{650}{(*)}{关于此事的若干保留意见,见前一条页边注。}。% ---------------- p.760 ----------------
\noindent 此外,也只是在我1970年离去之后,我才开始逐步地、且每一次都带着惊愕地发现:甚至在那些才具非凡的人身上,这种才具竟有时仿佛被“非理性”的成见所歼灭、所无望地封堵——而且越是……便越是顽固!我在这一方面最初的一些经验发生在1976年\fnote{651}{(**)}{(5月16日)这并不真是我在这方面的最初经验——前几年我另有若干次,尤其与Deligne一起,也在我的离去之前的往岁之中。但那些经验始终是零星的,而围绕拉德盖列叶博士论文的那一幕却令人震撼,因为五位数学家(个个水准极高)在行为与不作为上的完全一致,他们相互之间想必不曾串通一气。在那儿,在我与我的朋友皮埃尔一人之关系的种种变故之外,我第一次接触到了“葬礼仪式”。\par 然而,那些所谓“非理性”的因素在所谓“科学”思想中的惊人分量,远远超出“葬礼仪式”的语境,甚至超出一个时代。无须精通科学史(我全然不精通),便可看出:科学史每一步都铭刻着一种巨大惯性的效应,它抗拒一切新思想的萌发,也抗拒新思想在偶然萌发之后的茁壮成长。关于这一方向的反思,尤其见《愚妄与更新》头两部分(“劳作与发现”与“梦与梦者”),第1至8节。},在“进步是挡不住的!”一注(n\textsuperscript{o}~50)中亦有所追述;而关于此事的第一番书面反思,则在“掘墓人——或全体教众”(n\textsuperscript{o}~97)\fnote{652}{(***)}{这一反思在“阴阳之钥”中大有深化,尤其是在涉及同一“集会”的两条注:“天意攸关的境遇——或登峰造极”与“否决(一)——或召回”(n\textsuperscript{o}s~151、152)中。另可参看“肌肉与肚肠(阳葬阴(1))”一注(n\textsuperscript{o}~106),它开启了关于阴与阳的长程反思。}中继续展开,\origpage{981}其语境正是那\dmark{}“葬礼仪式”的特殊情形。同样,也只是渐次地、并且逆着巨大的惰性力量,我才意识到:这些“非理性的”原因并不因此便全不可解,只要人肯费神驻足、把它们探个明白。正因如此,我终于也“接受”了它们,尽管勉强……

回到我这个人,回到我与数学的关系。以我的工作方式而言,我倾向于靠一份份往往为时过早的臆测前行,而不讲求“审慎”\fnote{653}{(*)}{关于这种工作方式,尤其见“兄弟与配偶——或双重签名”一注(n\textsuperscript{o}~134),以及(《愚妄与更新》中的)“错误与发现”一节(n\textsuperscript{o}~2)。};但我把每一个出现的直觉(或“臆测”)追索到底,于是散布在工作最初阶段的大量错误终归自行剔除,让位给一种坚不可破的理解,而且(多半)确确实实触及事物的核心。当要对别人的工作下判断时,我自发的方式便全然不同,尤其当那工作立于我不熟悉的题目或不熟悉的调门上。我觉得,自己那时一向倾向于表现出审慎与谦逊。事实上,这恰恰是接纳我跻身其中的大多数前辈给我做出的榜样,如Cartan、Dieudonné、Chevalley、Schwartz、Leray——且只说这几位。我不记得听他们哪一位曾斩钉截铁地表过态,不论是为非还是为是,论一件他们摸不着其实质的工作。我现在才意识到:这份审慎原是我曾在他处谈及的那股\textbf{尊重}空气的一部分,它浸润着接纳我的那个圈子\fnote{654}{(**)}{另见“受欢迎的异乡人”一节,n\textsuperscript{o}~9。}。我觉得,在这圈子——我在二十多年的岁月里曾自居其中的圈子——最先退化的,正是这份审慎,而它是一份尊重的表征。也许我的记忆骗了我,也许我在自我欺骗;但我觉得,自己相对而言
并未怎样受那种尊重气氛退化的这一面影响。我想,我一直意识到自己在数学上无知有多广,也意识到自己无力领会他人的某项工作——每当那工作落在我的兴趣中心之外,而我的兴趣又往往是强烈聚焦于一处。

\origpage{982}\dmark{}至于那些我能够领会、因而也便能够赏识或评判的他人工作(只要我肯为此费点神),我同样记不起事后不得不承认有什么粗疏的判断失误,无论在贬在褒。我对自己那些观念与直觉的感觉也是这样,无论它关乎一个“好问题”之存在(或不存在),关乎有待探察的一份厚实材质,关乎某个观念的分量,还是关乎我对某一处境或某一事物那份或多或少的完整领会与深入领会。在所有这些场合,倘若确有差错,那差错每一次都是朝“少”的一边。是的——一个新题目或一个新观念之丰富,它在深度与广度上真正的分量,往往只是点点滴滴才充分显露,随着周,随着月,有时甚至是随着年。于是,一份起初正确的感觉(多半如此)——尽管它起初仍旧模糊而弥散——经由深浅粗细不等的“逐例功夫”而得到的这份渐次确认,于我们便如同一份意外、一抹惊叹,随着时辰与日子不断翻新。这无疑是研究工作(无论是数学的研究还是别的研究)那股非凡迷人之力的缘由:每一步,在我们眼前展现的实在,都在丰富、在精微、在深幽上越过我们最大胆的梦……

\noindent 不过我还是回到我对他人工作的领会,当那工作落在与我熟悉的题目上,甚至落在对我而言“火辣”的题目上。我想我可以说:我对一个观念真正分量的敏锐预感(这分量常常连作者本人都不自知),在我的作品中起过至为重要的作用。我在这里首先想到的,便是Serre所起的非凡作用,以及这么一件事:在1955至1970这十五年我作品异乎寻常地丰饶的岁月里,我的大多数观念、我的大多数重大投入,其出发点都在Serre的某个观念或某条进路上,尽管那观念有时貌似平淡。我打算在《主题纲要》\fnote{655}{(*)}{这些“历史评注”在《引论》之三“罗盘与行囊”中已有预告。}的“历史评注”里更详尽地谈这件事。但这么说,并不因此便是对Serre这个人有什么特别的敞开。(斟酌各比例而言)同样的事也在别的数学家身上发生过,无论在我作泛函分析家的过去,还是在我作几何学家\fnote{656}{(**)}{除却许多别的例子,我提一个为例:有一条原理,它把关于任一基座上“有限呈现”之相对概形情状的陈述,归约到该基座是一个\textbf{有限}局部环(甚或一个有限域)之谱的情形;这是一条意义重大的原理,而我是从 D. Lazard一个出色(且极特殊)的结果的一条引人注目的证法思路里把它提取出来的。关于此事,见“暂停!”一注(n\textsuperscript{o}~77)以及该注在第297页的那条页边注(***)。\par (5月16日)我不敢肯定:每逢我的进路受益于他人的一则观念,我都着意标明过。例如,我不记得在 EGA IV 相关的那一段里,我曾着意把Lazard列为那里所发展的一般归约方法之出处。那是当年看来无关宏旨的一桩疏忽。我想,Dieudonné(与我同为 EGA 撰述人)或Serre那些人,他们想必和我一样知道Lazard的这个结果乃是(无疑)同类中的头一个,也未必就认为非引不可(甚至只是适时)——何况这就不在Bourbaki的文风守则里!诚然,Bourbaki在历史评注里总会补上这笔账,而这类评注在 EGA 里、在我作品的别处都付阙如。今天,经受过数学界科学道德在七十、八十年代那可怕的败坏,我在标明出处上会比当年仔细得多——不仅在技术的意义上,也在启发的意义上标明,而后者常常远为关键。在已引用的那些“历史评注”里,我想我至少会补上我在这方面的一些疏漏。}的过去。

\origpage{983}\dmark{}我可以说,在我整个数学生涯中,我因着这份我刚刚试图稍稍描摹的、对待数学的单纯,得到了绰有余裕的“报偿”。这份单纯,
在我生命的别的圈子中常常令我失落,它本身却是一桩福祉。说实话,我作品的丰饶与力量,都归于这份单纯——它也正是\textbf{孩童}的单纯……

% ---------------- p.762 ----------------
\subsubsection*{a4.~……与羁绊}
\tocsubsub{a4.~……与羁绊}
\notehead{171(viii)}{}

\noindent (5月4日)\fnote{657}{(*)}{本子注源出“祖先”一注(n\textsuperscript{o}~171(i))的一条页边注,见第947页的注(*)。}我在这里说错了;在过去两个月里,我重新与本题稍稍接续,记忆也便弄清了些(纠正了些)。其实,Deligne的主要用意,恰恰就在于给出可构的离散(即 $\mathsf{C}$-向量空间)层的一个“纯粹代数的描述”,连同相宜的导出范畴\fnote{658}{(**)}{这里说的是那个范畴(在“作品……”一注(n\textsuperscript{o}~171(ii))中记作 $\mathrm{Cons}^{*}(X,\mathbb{C})$):它由 $X$ 上 $\mathsf{C}$-向量空间层的复形组成,其上同调层解析可构,并把它看作 $D^{*}(X,\mathsf{C}_{X})$ 的一个全子范畴。}。他所引入的系数(办法是在一个 pro-晶体层上施加一条临时的\origpage{984}“可构性”条件,该条件由存在一个适当的\dmark{}\textbf{“脱卸”}来定义,而那脱卸是照着我当初在平展情形或复解析情形引入的那个仿制的)——这些系数都是“量体裁衣”,专为回应这条希求。于是,六种运算的形式主义对这些系数\textbf{理应}存在(在特征零时),这事(在启发上)便“显而易见”;而且严格说来,它甚至能“粗暴而笨拙地”证出来:只要把“外尔原理”用得恰当,归约到基域是 $\mathsf{C}$ 的那个(已知的)情形。

\noindent 于是,只要人肯驻足一想,Deligne竟会撂下一条显然满含希望的进路,不再朝“系数范畴”的描述走去——而在六十年代中期,这一点已经再清楚不过:这些范畴就要在代数簇的上同调中起关键作用——那么这就成了个谜。于是他把最终提炼出一条某种意义上对偶、且更犀利的进路\fnote{659}{(*)}{我对此毫无疑义:倘若Deligne不曾撂下德拉姆系数这个题目(他是从我这里接过来的),他随即势必会(比那位当值无名之辈早八年)发现 D-模的“对偶”瑜伽,并因而熟悉Sato学派的思想。}的差事交给别人——八年后;而那进路随即\fnote{660}{(**)}{“随即”一词并不完全合乎实情的样子(它合乎的是“本该如此”的情形,假如……)。实则,从新的哲学与新的工具俱已就绪之时,到那些定调的人终于发觉此处有物可用(而且可顺手收入囊中……),中间过了三年。}刷新了几何中的上同调主题。这事先前并不怎样震着我,因为Deligne这一理论的发端就在我离去前不久,而那时还没有任何东西能够预示它将被赋予的命运。另一方面,自我离去之后,而且几乎直到最近这几个月,我与上同调主题全然断了联系。

\noindent 我近来曾想到——有些草草,未曾细究——Deligne对他自己投入了整整一年的理论竟如此疏离,其缘由或许在于:他对自己那条以“脱卸”定义的“可构性”判据并不满意。那判据也许显得过于简单,而确实:它远不如Mebkhout1976年在其“对偶”观点中提炼出的那条关于完整性与正则性的局部代数条件来得深。可深思之后,这份“解释”根本站不住脚!一条通往一个\origpage{985}枢机问题的进路“太过简单”,\dmark{}决不至于叫一位神志清明的数学家连那进路带那问题一并放手!他至多会在某天丢下他最初的进路——当他找到另一条,得以对同一问题得到一份更深、更全的视野时\fnote{661}{(*)}{其实,在这桩具体情形里,在我看来根本无所谓为着“老天爷”的进路(且不说出Mebkhout的名字)而“抛开”Deligne进路之事。两条进路互为补充:Deligne那条的长处在于更贴近几何直觉,Mebkhout那条则在技术上更简单(免得借助 pro-对象),并且在种种方面更深。}!
% ---------------- p.763 ----------------
\noindent 只要我对这桩古怪的情状稍稍驻足,事情便清楚起来:在这个情形里,一如在许多别的场合,我的朋友皮埃尔的动机全无数学的成分,甚至全无“理性”的成分。重新思量之下,我才察觉,围绕德拉姆系数的那整套问题意识——它唯有在六种运算与晶体瑜伽的视野中才获得它全部的意义(那瑜伽是我前几年引入晶体topos时提出的,而其精神恰恰就是六种运算……)\fnote{662}{(**)}{我还记得:Deligne讲他自己的理论时,一贯回避使用晶体语言;而那语言本可给他的理论一份更深的维度,把它嵌进一个业已存在的、topos式的上同调形式主义。同样,我如今也看出:无论是Berthelot,还是我别的上同调学生,都早已失去那种深层\textbf{统一性}的感觉——特征 $p$ 的晶体上同调与特征零的晶体现象(那正是他的讨论班的主题)之间的统一性。这些都是蓄意无视一种根本统一性的迹象;那统一性被任意切割,并因此被摧毁。这种蓄意,正是“封堵”的本性:自我中心的力量的介入,与求知冲动无涉。关于这种封堵在另一位上同调学生身上的例证(我虽曾认他天生具有细腻的直觉),参见“共同继承者们……”一注的子注 $n^{\circ}$912。}——我察觉那整套问题意识在我作品中和在我这个人中扎根之深,而且对人人都是明明白白地摆在眼前的。诚然,霍奇系数的问题意识也出自同一位师长,只是那学生早在其内心深处(也许连他自己也不自知)便已在疏离他。可那师承在外人眼里远不那样显豁(无人——连Serre在内——似乎有所察觉)\fnote{663}{(***)}{从Serre论Deligne工作的报告中,似乎确能看出这一点;该报告引在“为一副模糊的骸骨作安魂曲”一注的子注 $n^{\circ}$1651 中(尤其第813页)。关于这一师承的明白宣示,见“把一个个~i~点清”一注($n^{\circ}$164)第一部分第4节(尤其第793页)及其子注 $n^{\circ}$1641。},而最要紧的是:要做的那番意义重大的工作,其头一段并不落在一种表面上\textbf{大张旗鼓}Grothendieck化的视野(“六种运算”或别的……)之内——至少并不明明白白地摆在人人眼前。

\origpage{986}\dmark{}然而,我早已不止一次地着重指出:Hodge--Deligne的上同调理论,在六十年代末一次惊人的起步之后,至今仍停留在幼年阶段——那里被容许的系数唯有常数系数(或退一步说,“光滑”系数,即局部系统在“Hodge--Deligne”意义下的对应物),而连Leray的高次直接像 $R^{i}f_{*}$(且只说这些)那样要紧的运算,在那里竟也毫无通行的地位!“霍奇系数”的正确概念以及在其上建立恰当运算的问题,在Deligne的作品中(就我所知)甚至不曾被\textbf{提及}——而除非我记错,早在我有幸结识他之前,这问题于我便已熟稔。自我离去之后,多年来我有时曾问起:他究竟在等什么,好把他那非建立起来不可的理论最终建立起来(我终究也就厌倦了……)。他答得始终如一:“那太难了……”\fnote{664}{(*)}{后来,这一回答与Deligne手笔的《葬礼颂词》(或以赞词行葬礼)联在了一起,那篇颂词最近也曾提及(见“珍宝”一注 $n^{\circ}$170(iii))。这篇“颂词”以这样一问收场(此问当真值皮埃尔若干斤两……):“他在1970年离开 IHÉS,那时他对数学的热情正渐消渐隐。该相信他在那条为自己划定的路线上所提的问题,已经\textbf{变得太难了}吗?”(着重号是我加的)。这一殷勤的建议在《颂词》的第二段(专论Pierre~Deligne的那一段)里重述,我们从那里得知:逝者的某些猜想“今日仍如其时一般无从接近”,它们想必(至少这是明白的暗示)就是那位Deligne为证明某条“难点众所周知”的猜想所须克服的主要障碍。这些联想让我明白:在我朋友皮埃尔那句千篇一律的回答“那太难了……”里,含有一种嘲弄的弦外之音;这弦外之音必定给了他一种格外辛辣的满足,因为明摆着,那位高大的呆汉、懵懂的逝者万万想不到这弦外之音(更想不到自己已是逝者……)。}。这话并不能说服我,这是肯定——倘若我不曾抽身投入另一桩全然不同的冒险,我也会当即动手,把这“太难”的理论连同德拉姆系数的理论一并发展起来……

% ---------------- p.764 ----------------
\noindent 回过头来看,Hodge--Deligne理论中的停滞,与Deligne本人对待德拉姆系数主题的那副荒唐态度(\origpage{987}那态度在1981年6月吕米尼那场难忘的座谈会上达到顶峰\dmark{},那桩“反常”的不义之举将永远与它相连……),两者之间的平行使我吃惊。这两桩荒唐事,如今在我看来密切相关,而且关涉的层面全然不是数学的层面。诚然,显而易见,霍奇系数的形式主义之发展\textbf{从属于}德拉姆系数的形式主义之发展(这一点我早在1966年前后就已看得清清楚楚,而人们似乎是在那位从不被提起名字的身后学生的踪迹上,最近一两年才正在发现它……)。这个数学事实使两串事实之间的联系、以及它们各自的反常性质,都更加触目:因为这条“客观的”联系本是一份额外的强有力激励(至少对一个“充分掌握着自己能力”的人而言),促使他把两个理论都发展起来——它们本来只可能相互照亮、相互加强。

\noindent 这两个理论各自的停滞(德拉姆理论直到1981年的反常层会议为止,霍奇理论则直到今天),很大程度上在于上同调主题整体的萎靡不振,关于这种萎靡我不止一次有机会谈到过\fnote{665}{(*)}{关于这种萎靡不振,请特别参看“荒芜的工地”(《葬礼行列》,6.),尤其是“工地巡礼——或工具与视景”一注(n\textsuperscript{o}~178)。}。即便撇开人的精神维度不谈,即便只考虑通过“前沿”科学生产获取“收益”这一仅有的因素,这种停滞在我看来也以惊心动魄的方式同时说明了两件事:一是不为人知的隐秘的自我中心力量能在一个人身上取得怎样的统治——即便是在一门自称“无私”的科学的操持之中也是如此;二是这种统治那(显然)反常的性质,它在这里(至少初看之下)似乎总是在与所追求的目的背道而驰\fnote{666}{(**)}{至少,若把“目的”理解为摆在世人面前的那个(“科学的进步”,姑且说),或者甚至是那个丝毫不假的目的——通过积累迫使他人敬重与赞叹的作品来扩大声望,情形便是如此。然而在我看来,即便是这一种“收益”,比起最强大的那些隐秘力量所追逐的满足来,也只是次要的——而我的朋友正是选择让这些力量统治他的存在。}。

\subsubsection*{b1.~五张照片(晶体与 D-模)}

\tocsubsub{b1.~五张照片(晶体与 D-模)}

\notehead{171(ix)}{}

\fnote{667}{(*)}{本子注是“作品……”一注(n\textsuperscript{o}~171(ii))之下的一条子注,性质纯粹是数学的。读者若无意用数学的措辞,对Zoghman~Mebkhout的工作、以及作为簇的上同调理论中一门新的“系数理论”的“D-模瑜伽”,有哪怕些许的领会,尽可略去本子注。下面各页可以看作这套瑜伽、或“Mebkhout哲学”的一段简短导引,用的是一套晶体式的概念装备与整体视景。这一视景早在1966年就已在我心中成形,却被我的同调学生们Deligne、Berthelot、Illusie、Verdier系统地、几乎完全地遮蔽了,而他们本是它的主要保管人。唯一留下的文字痕迹,是我1966年在 IHÉS 所作报告“Crystals and the De Rham cohomology of schemes”的记录稿(I.~Coates 与 O.~Jussila 记录),载 \emph{Dix exposés sur la cohomologie étale des schémas},North-Holland Pub.~Cie(1968)。然而从技术观点看,这个报告包含着晶体上同调的全部启动观念。除了Mebkhout的工作之外,概念层面(或其他层面)似乎不曾有任何真正关键的进展——恰恰相反,我看到一种相对于我六十年代那些观念的惊人退步。不幸的是,那些观念在所引报告中只以极零碎的方式、或字里行间的方式出现;这里一如别处,最大的空缺,是完全没有明确提到德拉姆系数的问题意识,以及应当为这类系数建立的六种运算(与双对偶)形式主义(x)。我曾能够确认:Mebkhout虽然比任何别人都更熟悉我关于上同调的文字著作(以及我的学生们的那一部分),却全然不知这一原初的问题意识(直到两年前还是如此)——而在我看来,就数学“底层”而言(且不谈非理智性质的心理因素),这直到今天仍是他主要的障碍。今后我将把上述1966年的报告简称[Crystals]。\par(x)(6月16日)关于一条订正,见第990页脚注(**)。}

% ---------------- p.765 ----------------
\subsubsection*{(a) “德拉姆系数”图册}

\tocsubsub{(a) “德拉姆系数”图册}

\origpage{988}\dmark{}(5月4日与5月19--20日)我提醒一句:对一个光滑的复解析空间,记 $D_{X}$(或径直记 $D$)为 $X$ 上复解析微分算子的环层(更确切地说,是 $\mathsf{C}$-代数层)。Sato揭示的第一个关键事实,是这是一个\textbf{凝聚}的环层。第二个事实,性质是重言式的,却也同样关键,那就是:局部自由 $\mathcal{O}_{X}$-模的范畴——其中取作态射的,不只是 $\mathcal{O}_{X}$-线性态射,而是此类模之间的微分算子——经由一个先验上\textbf{反变}的函子,作为一个\textbf{全子范畴}嵌入局部自由 D-模的范畴,这个反变函子是\fnote{668}{(**)}{这里写下的同构是 $u\mapsto\varepsilon\circ u$,其中 $\varepsilon\colon D\to\mathcal{O}_{X}$ 即“增广”$\theta\mapsto\theta(1)$。}
\begin{equation*}
F\longmapsto \mathrm{Hom}_{\mathcal{O}_{X}}(F,D^{d})\xrightarrow{\ \sim\ }\mathrm{Op}_{\mathrm{diff}}(F,\mathcal{O}_{X}),
\tag{1}
\end{equation*}
其中 $D^{d}$ 表示 $D$ 配备其由典范\textbf{右} $D$-模结构诱导的 $D$-模结构,\origpage{989}该结构与 $D$ 在其自身上的左运算\dmark{}交换(这些运算使(1)的第二个成员成为一个 D-模)。这个全忠实函子进而在由自由模组成的全子范畴之间导出一个(反)\textbf{等价}。它不容许一个典范的拟逆函子,“与向开集的限制交换”——这想必就是为什么方才这第一个反变函子一般说来并不是一个等价。如果 $C$($C$ 犹如“cristal〔晶体〕”,见下文)表示一个局部自由的 D-模(甚至是自由的,也无妨),人们确实可以给它结合一个以函子方式依赖于 $C$ 的层:
\begin{equation*}
C\longmapsto \mathrm{Hom}_{D}(C,\mathcal{O}_{X}),
\tag{2}
\end{equation*}
这是一个反变函子,看上去可能为(1)提供“那个”拟逆函子的天然候选。麻烦在于,这个层(2)并非以自然的方式配备 $\mathcal{O}_{X}$-模结构,而只配备 $\mathsf{C}_{X}$-模结构(其中 $\mathsf{C}_{X}$ 是 $X$ 上由复数域 $\mathsf{C}$ 定义的常值层)。当 $C$ 由一个局部自由的 $\mathcal{O}_{X}$-模 $F$ 经由反变函子(1)得出时,(2)便典范同构于 $F$ 底下的那个 $\mathsf{C}$-向量空间层。

\noindent 函子(1)(如任何加法函子一样)延拓到复形的范畴:它把 $X$ 上的一个微分算子复形(通常意义下)变成一个局部自由的 $D_{X}$-模复形,而由此得到的(反)函子当然是全忠实的(对于第一个复形范畴中微分算子复形之间的微分态射)。正是在这个意义上,可以说 D-模复形(其分量为局部自由者)“\textbf{推广}”了 $X$ 上的微分算子复形。

\noindent D-模复形的视点相对于微分算子复形的视点有决定性的优势,就是直接嵌入环空间上模复形的瑜伽(最早发展于我1955年的文章《论同调代数的若干问题》\fnote{669}{(*)}{In Tohoku Mathematical Journal, 9 (1957) p.~121--138。}),并由此尤其嵌入\textbf{导出范畴}的瑜伽(那是我在所引文章之后的几年间提炼出来的)。关键观念“\textbf{拟同构}”,在采取微分复形之间的微分态射这一视点时并不跃然眼前,而一旦过渡到相结合的 D-模复形就变得显然。因此,\origpage{990}Mebkhout所引入的视点\fnote{670}{(**)}{(6月8日)此处应当读作:由Mebkhout为了一种新的系数理论的需要,引入Grothendieck的装备之中。不言而喻,“D-模的视点”应归功于Sato,但他是在一种截然不同的 optique〔取向〕中使用它的。},与其说是微分算子复形视点的一种\textbf{推广},不如说是一个\textbf{关键的放松}:\dmark{}正是借助这个视点、而且只有借助它,微分算子复形从今以后才能被用作一门新上同调理论的“系数”,并带着与之相联系的全部丰富直觉。

% ---------------- p.766 ----------------
\noindent 倘若我在德拉姆的系数理论与 $\ell$-进系数理论(后者事实上是Mebkhout发展其哲学的主要灵感来源之一)之间设立一个平行对照,那么我会说:这概念性质的第一步,这“孩童般的”一步,类似于我在1958年引入 étale 层概念时所迈出的那一步(其中已孕育着“topos”这一关键的统一概念)。在同一类比之中,“好上帝定理”(下文即将重述)类似于 étale 上同调中真态射的基变换定理——它(1963年)是 étale 上同调起步阶段的第一个大定理,短短几周之内就使人对 étale 上同调这件工具达到几乎完全的“驾驭”。而在 D-模的框架中(更一般地,在晶体的框架中),要达到对“晶体上同调”(或广义的“德拉姆上同调”,我早在六十年代就已预见这样一门理论)的同样驾驭,那相应的工作\textbf{仍有待完成}——自从Zoghman~Mebkhout终于完成第一次大突破以来,已经过去七年。

\noindent Mebkhout引入的这个新的系数范畴,“包含”(其确切含义见“作品……”一注,n\textsuperscript{o}~171(ii))既包含解析可构的“离散系数”,也包含Serre引入的凝聚系数(我曾把它们系统化为一套“凝聚系数”的上同调理论\fnote{671}{(*)}{这是我在五十年代后半叶在凝聚框架中发展的六种运算与双对偶的形式主义。}),它就是由上同调层(作为 D-模)\textbf{凝聚}的 D-模复形组成的范畴,视作通常的导出范畴 $D^{*}(X,D_{X})$ 的一个全子范畴
\begin{equation*}
D^{*}_{\mathrm{coh}}(X,D_{X}) \quad \text{或} \quad \mathrm{Cris}^{*}_{\mathrm{coh}}(X),
\tag{3}
\end{equation*}
若只限于上同调有界的复形(它们构成全子范畴 $\mathrm{Cris}^{b}_{\mathrm{coh}}(X)$),那么这样一个“系数”在\textbf{局部上}可由一个各次均为有限型自由 D-模、且次数有界的复形来表示;或者,本质上等价地,由一个次数有界的微分算子复形来表示\fnote{672}{(**)}{(6月16日——参看第988页注(*)之末尾。)Mebkhout刚刚向我指出,此说不尽确切——这一问题意识在 loc.~cit.~1.5~d)(第312页)处已略有提及。Mebkhout在他的工作“Poincaré 对偶”(巴黎第七大学“奇点”讨论班,1977--79)中明确提到了它,见于 §4.4(关于 D-模的相对对偶定理)的最后三行。}。

\dmark{}\origpage{991}使用导出范畴时,当然要把基本函子(1)与(2)替换为总体导出函子
\begin{equation*}
F\longmapsto \mathrm{RHom}_{\mathcal{O}_{X}}(F,D_{X})\ ,\quad C\longmapsto \mathrm{RHom}_{D}(C,\mathcal{O}_{X}).
\tag{4}
\end{equation*}
若要寻找与这两个函子同类型的\textbf{协变}函子,首先碰到的就是“标量扩张”函子(在所引注中记作 $N$):
\begin{equation*}
F\longmapsto D\otimes_{\mathcal{O}_{X}}F,
\tag{5}
\end{equation*}
(总张量积),其中张量积还要用到 $D$ 的 $\mathcal{O}_{X}$-右模结构,即 $D^{d}$\fnote{673}{(*)}{众所周知,$D$ 作为 $\mathcal{O}_{X}$-右模或左模都是平坦的(从 $D$ 的典范滤过以及相伴分次层的已知形状即可立见)。由此可知,(5) 中的“总”张量积其实就是一个通常的张量积。}。这个(就 $F$ 而言的)函子相对于(1)有一个缺点,就是不能延拓到那些变元 $F\to F'$ 仅仅是微分算子(而非线性映射)的态射。第二个函子(4)应当看作一个反函子
\begin{quote}
$\mathrm{Cris}^{*}_{\mathrm{coh}}(X)\longrightarrow D^{*}(X,\mathsf{C}_{X})$,
\end{quote}
它也有一个重要的协变“对应物”,由
\begin{equation*}
C\longmapsto \mathrm{RHom}_{D}(\mathcal{O}_{X},C)\overset{\mathrm{def}}{=}\mathrm{DR}(C) \quad \text{(与 $C$ 相结合的“德拉姆复形”),}
\tag{6}
\end{equation*}
给出,其中第二个成员确实借助 $\mathcal{O}_{X}$ 的、由局部自由的有限型 D-模构成的所谓“Spencer 分解”而明确地写成一个德\,拉姆型的复形。(这一分解取自通常的德\,拉姆复形,经由函子(1)取相结合的 D-模复形而得。)用晶体的话来说(下文将予阐明),函子 $\mathrm{DR}$ 明确写出为函子 $C\mapsto \mathrm{Hom}_{D}(\mathcal{O}_{X},C)$ 的总体导出函子,它给每个“D-模”(或“晶体”)结合由其在可变开集上的“水平”截面构成的 $\mathsf{C}$-向量空间层。这是一项\textbf{局部性质}的操作。对于一个“系数” $C$(即一个 D-模或此类复形)而言,恰当的(整体)“\textbf{积分}”(或整体上同调对象)概念,在这里并非通常的函子
\begin{quote}
$\mathrm{R}\Gamma_{X}(C)\simeq \mathrm{RHom}_{D}(X;D,C)$,
\end{quote}
而是(对我而言作为\textbf{晶体整体上同调}函子十分熟悉的)“水平(整体)截面”\origpage{992}函子 $C\mapsto \mathrm{Hom}_{D}(\mathcal{O}_{X},C)$ 的总体导出函子;我把这一总体导出记作 \dmark{}$\mathrm{R}\Gamma_{\mathrm{cris}}(C)$,使得有那些重言同构
\begin{equation*}
\mathrm{R}\Gamma_{\mathrm{cris}}(C)\overset{\mathrm{dfn}}{=}\mathrm{RHom}_{D}(\mathcal{O}_{X},C)\simeq \mathrm{R}\Gamma_{X}(\mathrm{DR}(C)),
\tag{7}
\end{equation*}
亦即 $C$ 在 $X$ 上的晶体上同调,是由取相结合的德\,拉姆复形的(整体)通常上同调而得到的。

\noindent 在 $\mathrm{Cris}^{*}_{\mathrm{coh}}(X)$ 中可以定义一个\textbf{对偶化函子},它给出一个双对偶定理,其模式就是我先后在(交换的)凝聚语境、随后在离散(étale)语境中提炼出的那些。我把它记作 $D$(如在所引诸语境中一样):
\begin{equation*}
D : \mathrm{Cris}^{*}_{\mathrm{coh}}(X)\overset{\sim}{\longrightarrow} \mathrm{Cris}^{*}_{\mathrm{coh}}(X).
\tag{8}
\end{equation*}
这是一个反等价,本质上是对合的,即有一个关于 $C$ 的函子性双对偶同构:
\begin{equation*}
C\simeq D(D(C)).
\tag{9}
\end{equation*}
\noindent {\small [译按:原书式(9)右端印作 $D(D(X))$;据本段的变元 $C$ 及双对偶的含义,此处改为 $D(D(C))$。]}

这个函子可以(通过复合)把反函子(1)与(2)变成协变函子。要记住的简单事实是:若 $C$ 与 $C'$ 互为“对偶”,则其一的德\,拉姆复形(6)等同于另一个的“余德\,拉姆”(2):
\begin{equation*}
\mathrm{RHom}_{D}(\mathcal{O}_{X},C)\simeq \mathrm{RHom}_{D}(C',\mathcal{O}_{X}).
\tag{10}
\end{equation*}
反之亦然。
在微分算子复形上,这个操作 $D$ 可表达出来(就次数而言相差一个 $n$ 的“平移”),办法是过渡到“相伴的”微分算子复形,其分量为 $\mathrm{Hom}_{\mathcal{O}_{X}}(F,\omega_{X})$,它是逐项取相伴算子而得的。于是,D-模的对偶化函子就与 Serre 对偶中那个熟悉的对偶化函子相容:
\begin{equation*}
F\mapsto \mathrm{Hom}_{\mathcal{O}_{X}}(F,\omega_{X})\simeq F\otimes_{\mathcal{O}_{X}}\omega_{X}.
\tag{11}
\end{equation*}
这里 $F$ 为有限型局部自由 $\mathcal{O}_{X}$-模。
\noindent {\small [译按:式(11)照录原书;右端的 $F$ 疑漏对偶记号。按左端的定义,应为 $F^{\vee}\otimes_{\mathcal{O}_{X}}\omega_{X}$,其中 $F^{\vee}=\mathrm{Hom}_{\mathcal{O}_{X}}(F,\mathcal{O}_{X})$。]}

其中 $\omega_{X}$ 表示 $X$ 上最高次微分形式的“对偶化模”。要注意的是,德\,拉姆函子
\begin{quote}
$\mathrm{DR} : D^{*}_{\mathrm{coh}}(X,D)\longrightarrow D^{*}(X,\mathsf{C})$,
\end{quote}
一般\textbf{不}与对偶化函子交换(在第二个范畴中取函子 $\mathrm{RHom}_{\mathsf{C}}(-,\mathsf{C}_{X})$)。但Mebkhout有一条深刻定理(人人都在用它,当然谁也不注明出处,仿佛那不过是一串简单的连环推理),就是说:对于\textbf{完整}的变元,因而对于诱导出的函子
\begin{quote}
$\mathrm{Cris}^{*}(X)_{\mathrm{hol}}\longrightarrow \mathrm{Cons}^{*}(X,\mathbb{C})\ (,\hookrightarrow D^{*}(X,\mathsf{C}))$
\end{quote}
\origpage{993}而言,与对偶化函子是交换的。我不在此“重述”完整性的条件\dmark{},只愿指出:一个 D-模复形是完整的,当且仅当它的上同调层都是完整 D-模;而这是一条在 $X$ 上局部性质的条件,并且还是“代数的”。另一方面,Kashiwara的可构性定理(Kashiwara陈述这定理时只针对一个完整\textbf{模},那时无论他自己还是任何人——除了Mebkhout——都不在导出范畴下工作……)蕴涵:德\,拉姆函子限制到完整复形上,确实取值于 $\mathrm{Cons}^{*}(X,\mathbb{C})$。引入Mebkhout的\textbf{正则性}概念——它同样是局部性质,并且是“代数的”\fnote{674}{(*)}{我提醒一句:Mebkhout对正则性的原始定义是超越性质的。关于一个“纯粹代数”的转译,请参见关于德拉姆系数的预定报告(“Mebkhout风格”或“Deligne风格”),载《数学反思》第三卷。}——就得到“好上帝函子”(又名Mebkhout函子)
\begin{equation*}
m : \mathrm{Cris}^{*}(X)_{\mathrm{hol.reg.}}\overset{\sim}{\longrightarrow} \mathrm{Cons}^{*}(X,\mathbb{C}),
\tag{12}
\end{equation*}
它这一次是一个\textbf{等价}(如我们在“作品……”一注(n\textsuperscript{o}~171(ii))中所见),因而与自然的对偶化函子相容。正是那个拟逆函子
\begin{equation*}
M : \mathrm{Cons}^{*}(X,\mathbb{C})\overset{\sim}{\longrightarrow} \mathrm{Cris}^{*}(X)_{\mathrm{hol.reg.}}\hookrightarrow \mathrm{Cris}^{*}_{\mathrm{coh}}(X),
\tag{13}
\end{equation*}
使得可以把 $X$ 上“可构离散系数”(即 $\mathsf{C}$-向量空间层)的范畴看作 $D^{*}(X,D)$ 的一个全子范畴、更确切地说是 $D^{*}_{\mathrm{cons}}(X,D)=\mathrm{Cris}^{*}_{\mathrm{coh}}(X)$ 的一个全子范畴,而人们将时而把它解释为一个“晶体”系数的范畴。

\noindent (5月19日)就目前而言,我们可以说,已经用三种不同的“语言”或“视点”,如同用三张不同的“照片”,描绘了同一个实在,或(本质上)同一类“系数”,即所谓“德拉姆系数”:有 $\mathsf{C}$-向量空间层及其复形这一视点(“拓扑”拍摄),配以一条“解析可构性”条件\fnote{675}{(**)}{我提醒一句:解析空间 $X$ 上的一个 $\mathsf{C}$-向量空间层称为“解析可构的”,如果在每一点的邻域内,它有一条合成列,其相继因子形如 $i_{!}(F)$,其中 $i : Y\to X$ 是 $X$ 的一个解析子空间 $Y=Z\backslash T$ 的包含映射(这里 $T\subset Z$ 是 $X$ 的两个解析闭子空间),而 $F$ 是 $Y$ 上一个局部自由的有限型 $\mathsf{C}$-层(或“$\mathsf{C}$-向量空间的局部系统”)。},起着有限性条件的作用(为了能写出 Riemann--Roch\origpage{994} 型的定理,这一点尤为不可或缺,这类定理涉及恰当的“Euler--Poincaré 示性数”与恰当的“\dmark{}Grothendieck 群”)。还有“微分算子复形”的视点,以完整性与正则性两条条件取代可构性条件。还有“D-模复形”的视点,随之而来的是凝聚性、完整性与正则性条件。

% ---------------- p.769 ----------------
\noindent 第二张“照片”(从“分析”的角度拍摄)很诱人,因为它用“经典”的语汇就能为我们所理解,而且它呈现的对象——即微分算子的复形——在我们看来“尺寸”适中,而 D-模,即便是凝聚的(且不说 $D$ 本身!),戴上“$\mathcal{O}_{X}$-模”这副眼镜去看,也显得庞大。然而就技术而言,反倒是它们提供的照片更为完整。事实上,虽然局部上“每个上同调凝聚、次数有界(姑且说)的 D-模复形都可以经由(1)用一个微分算子复形来表示”这件事看着“明白”,整体上还能如此可就难说了,除非对 $X$ 施加苛刻的假设(比如“Stein 簇”,或者——在代数框架中——一条拟射影性假设)\fnote{676}{(*)}{当然,没有什么能阻止人们从 $X$ 上微分算子复形的范畴、以及这些复形之间的微分态射出发,把那些“拟同构”(由过渡到相结合的 D-模复形来定义)形式地“反转”,从而构造出一个“导出范畴”。人们将会得到(我推测)$\mathrm{Cris}^{*}_{\mathrm{coh}}(X)$ 的一个\textbf{全}子范畴,但在缺乏“Stein”或“$X$ 射影”(在代数情形中至少拟射影)一类假设时,恐怕不会是这整个范畴。}。

\noindent “照片”1 的好处在于:当 $X$ 不再被假定为光滑、而是一个一般的复解析空间时,它仍然保有意义。相反,照目前这个样子,“照片”2 和 3 只有在光滑性假设之下才讲得通。人们当然仍可以在 $X$ 上不假定光滑就定义环层 $D_{X}$,也仍然有一部自明的“词典”,在微分算子复形(其分量为局部自由 $\mathcal{O}_{X}$-模)与 D-模复形(其分量局部自由)之间来回对译;只是(据说)$D_{X}$ 不再凝聚了,真可惜!在奇异情形中,恐怕很难照光滑情形中那个已知榜样,提炼出一条“好上帝定理”。而另一方面显而易见的是:在奇异情形中我们同样需要 2 或 3 那一类的照片,因为“照片”1 是\textbf{就其性质而言属于超越}的:天真地照搬它、用代数簇的 Zariski 拓扑或 étale 拓扑来表述,所得的“系数”将特殊得太离谱,不堪使用(因为比起超越拓扑来,\origpage{995}\dmark{}这些拓扑太过粗糙)。相反,“照片”2 和 3一开始把视野限制在“光滑”之内,却在“抽象”代数几何中仍保有意义(比如说,一开始在特征零的域上),这就构成(对我而言)它们主要的魅力。这就是说,非得给它们做\emph{放大},把奇异簇也纳入视野不可。

\noindent 这看来不曾让Mebkhout费神,他忙着别的要事——我向他提出这个问题时,他当即的念头是这样的。设 $X$ 作为一个闭解析子空间嵌入一个光滑簇 $X'$ 之中。于是范畴 $\mathrm{Cons}^{*}(X,\mathbb{C})$ 可解释为 $\mathrm{Cons}^{*}(X',\mathbb{C})$ 的全子范畴,其成员是在 $U=X'\setminus X$ 上限制为零的对象(即支撑含于 $X$ 的对象)。而由好上帝定理,后者也可用“照片”2 或 3 的语言,解释为 $X'$ 上那些在 $U$ 上限制为零的“德拉姆——Mebkhout系数”的范畴。停留在“德拉姆——Mebkhout系数”(即“照片”2、3)的框架内,先验地验证下面这点应当是容易的:这个范畴——在等价的意义下,而该等价本身除差一个唯一同构外即被确定——不依赖于为 $X$ 选定的“光滑化”$X'$。我自己就干过许许多多这样的事,我乐意相信它行得通。而倘若 $X$ 并不“可光滑化”,那也无妨(Mebkhout如是说),我们可以“做上同调下降”,从这些局部的碎片重建出一个整体范畴;或者引入 $X$ 的开集之“光滑化位”,在上面工作。看来确有办法应付,只是与其用一个“光滑化位”(Mebkhout为作出回应的需要而临时搭建的
回答,在一场始终停留在柏拉图式的对话中),这个位在我看来显得强烈冗余,为什么不索性用晶体site呢,它早已受过检验(即便看来它已被那些曾是我弟子的人,带着一种动人的浑然不觉,抛诸脑后……)?更何况,早在1966年我廓清晶体瑜伽的启动观念之时,我就已经看得很清楚:未来的“德拉姆系数”本来就应当用晶体的措辞来表达!

\noindent 这于是引我从抽屉深处取出一张照片,它可有时间积起灰尘了,可怜见的——然而一旦把它吹净,它在我看来恍若全新,且清晰得无以复加。况且,这是我去年写“我的孤儿们”一注(n\textsuperscript{o}~46)时(甚至还没遇见“葬礼”之前)最早想到的事情之一,当时我隐约感到,是时候有人怀着敬意谈一谈那些值得尊重的事物了……此外,自从Mebkhout\origpage{996}\dmark{}跟我谈起 D-模以来(1980年——天晓得,那时我可并不“在行”!),我就禁不住把它们想成是“晶体”,并且把“D-模”与“晶体”(即 $\mathcal{O}_{X}$-模的晶体)当同义词使用,当然(明显地)偏爱后者。

\noindent 我于是来到许诺过的第四张照片,即“晶体”照片。先设 $X$ 光滑。在 $X$ 上给定一个 D-模 $F$,等同于给定一个 $\mathcal{O}_{X}$-模,外加一道附加结构,而这道结构可以用若干等价的方式来表达。其中重言式的一种,是说把 $\mathcal{O}_{X}$ 在阿贝尔层 $F$ 上的运算“延拓”为环层 $D_{X}$(它包含 $\mathcal{O}_{X}$)的一次运算。由于 $D_{X}$ 由 $\mathcal{O}_{X}$ 与导子的加法子层生成,可见在 $F$ 上给出所谓“\textbf{可积联络}”,也就是说给出这样一个法则:对 $X$ 的每个开集 $U$ 上的每个导子 $\xi$,结合 $F$ 的一个“$\xi$-导子” $\theta_{\xi}$,关于 $\xi$ 是线性的,并且与导子的“括积”运算相容\fnote{677}{(*)}{当然,还须有一条与限制到开集相容的条件。}——二者是一回事。可以说,这是 $F$ 上一道“微分”性质的结构,阶为 1。

\noindent 由于我们处在特征零之中\fnote{678}{(**)}{在下面,我们可以免除一切特征假设(比如说,在光滑相对概形的框架中),办法是把沿对角线的 $X\times_{S}X$ 形式完备化,换成“S-除幂”的形式完备化。这同样导致,对 $X$ 上的 $\mathcal{O}_{X}$-模层 $F$,要把它的“无穷阶主部”的 pro-层 $P^{\infty}(F)$,换成“(无穷阶)S-除幂主部”。在对偶方面,这就等于把相对微分算子的环层 $D_{X/S}$(即使 $S$ 是诺特的,它也毫无凝聚性可言),换成 $\mathcal{O}_{X}$ 在 $\mathcal{O}_{S}$ 上的相对导子的“包络”环层(据Mebkhout向我保证,它会是凝聚的!)。事实上,这就是德拉姆系数的概念语境,它将推广Mebkhout的 D-模语境,尤其是为了发展一套 $\mathrm{Spec}\,\mathbb{Z}$ 上有限型概形的德拉姆系数理论。},这道结构也可以被诠释为一道更丰富的结构,即一道无穷阶的微分结构,我曾把它称为 $F$ 上的“\textbf{分层}”(此时 $F$ 便得名“\textbf{分层模}”)。表达一个分层的一种方式,是把它写成 $F$ 上(相对于态射 $X\to\mathrm{pt}$ 的)一道“\origpage{997}\textbf{无穷阶无穷小下降资料}”,或者更确切地说,写成\dmark{}沿对角线取 $X\times X$ 的形式完备化之上、$F$ 的两个拉回像(由两个典范投影 $\mathrm{pr}_{1}$ 与 $\mathrm{pr}_{2}$ 给出)之间的一个同构,它延拓对角线上的恒等,并且此外还满足一条适当的“传递性条件”。

\noindent 从一道可积联络过渡到一道“无穷小下降资料”(即分层结构),这是一个新观念——而且“平凡”,一如我有幸发现的所有新观念!然而,只有一旦用模的晶体这一概念对它重新诠释,它才取得全部的力量。事实上可以看出,$F$ 上所说的这道结构,也就等同于给出:对 $X$ 的每个开集 $U$ 的每个“无穷小邻域” $U'$,给出 $F|_{U}$ 到 $U'$ 的一个\textbf{延拓} $F_{U'}$(总之,$F$ 在无穷小邻域上方“生长”,犹如一颗“晶体”——晶体有种种品类,此处是模的晶体……)——这个延拓,就限制到 $U$ 的一个开子集 $V$ 的概念、以及同一个 $U$ 的无穷小邻域(或称“加厚”)$U'$、$U''$ 之间的态射(当然是诱导 $U$ 上恒等的态射)而言,表现得一如人们预料。

% ---------------- p.771 ----------------
\noindent 晶体视点的好处,在于所要研究的对象(即 D-模)可以诠释为一个适宜位\fnote{679}{(*)}{要注意,在晶体site上并不能找到所有的模层,而只能找到那些满足一条简单的附加条件的模层(在[Crystals]中称为“特殊的”层)。}上的“寻常”模层,这个位以交换局部环为环层,即由 $X$ 的各个开集 $U$ 的加厚 $U'$ 构成的“\textbf{晶体site}”(晶体结构层简单地就是 $U'\mapsto\Gamma(U',\mathcal{O}_{U'})$)。由此,人们就拥有一部与这种局面相联系的全部几何直觉武库。我在1966年发现并当时为之震惊的一条卓越关系,就是:晶体site(或它所对应的晶体topos)以结构层为系数的上同调(或者更一般地,以 $F$ 为系数,至少当 $F$ 在 $\mathcal{O}_{X}$ 上凝聚时),等同于 $X$ 的\textbf{德拉姆上同调}(此处是以 $F$ 为系数,即以 $\mathrm{DR}(F)$ 为系数的 $X$ 的寻常超上同调)。这就是晶体上同调的起点\fnote{680}{(**)}{在这里又是如此——那些启动观念“平凡”到这种地步,实在不必为这一点点东西而局促:一个人后来花了十五年的生命,去发展它的一个小角落(而把其余忘掉……)。}。

\origpage{998}\dmark{}于是我们有一部完美的词典,在我已引的1966年那些讲演\fnote{681}{(*)}{见[Crystals]报告,引于本子注的第一条脚注(第988页注(*))。}中有详尽的阐述,它在 $X$ 上四类对象、或一个 $\mathcal{O}_{X}$-模上四种类型的结构之间建立对应:
\begin{quote}
(Cr)\quad D-模\\
\qquad 带可积联络的 $\mathcal{O}_{X}$-模\\
\qquad 分层模(无穷阶无穷小下降资料)\\
\qquad $\mathcal{O}_{X}$-模的晶体
\end{quote}
这部词典有效,对 $F$ 不附加任何凝聚性或拟凝聚性类型的限制。然而要注意,若比较两个极端项
\begin{quote}
D-模 $\Longleftrightarrow$ $\mathcal{O}_{X}$-模的晶体,
\end{quote}
那么在这两个语境中“凝聚性”的自然概念\textbf{并不相对应}。晶体结构层是凝聚的,但晶体topos上的凝聚模,恰好对应于作为 $\mathcal{O}_{X}$-模凝聚的 D-模,此时它们甚至是有限型局部自由的。它们构成的范畴,经由相对于 $\mathsf{C}_{X}\to\mathcal{O}_{X}$ 的“标量扩张”函子,典范等价于局部自由的 $\mathsf{C}_{X}$-模层范畴,即 $X$ 上“$\mathsf{C}$-向量空间局部系统”的范畴。因此对这类对象,就有五种可能的描述(若算上表(Cr)先前的四项,就是五张“照片”)!但这只是些性质极其特殊的“系数”\fnote{682}{(**)}{事实上,在此是 D-凝聚性(我六十年代未曾察觉它)才是重要的有限性概念。},在(de~Rham--Mebkhout)那些我们感兴趣的系数之中。

\noindent 我们还是回到上面表(Cr)的四张照片,看看不再假定 $X$ 光滑时会怎样。所考虑的四类对象仍保有意义。另一方面,看来头两类并不构成重要的范畴——更确切地说,人们自然地遇到的、作为“具有几何意义”的所有 $D_{X}$-模与所有带可积联络的 $\mathcal{O}_{X}$-模,都(在一个显然的意义下)“来源”于分层模,而这些分层模仍可诠释为 $\mathcal{O}_{X}$-模的晶体,与光滑情形一样\fnote{683}{(***)}{这个论断下得仓促,而且照原样是错的。要使它为真,必须把由 $X$ 的开集的一切无穷小加厚构成的“晶体site”,换成由那些在 $X$ 上局部容许一个收缩的加厚构成的子位(称为“分层位”)(当 $X$ 光滑时,这个条件自动满足)。若给定 $X$ 上一个分层模 $F$,它经过这样一个收缩的拉回,除了差一个唯一同构之外,不依赖于所选定的收缩,由此得到 $F$ 到所考虑加厚上方的一个“典范延拓”。因此可见,当 $X$ 不光滑时,$F$ 上一道晶体结构比一个单纯的分层“更丰富”,因为它使我们能把 $F$(即把它“养大”)延拓到 $X$ 的开集的\emph{任意}无穷小邻域上方,尤其是(而这是一件特别重要的事)延拓到浸没在一个光滑环境空间中的 $X$ 的各阶无穷小邻域上方。事实上,在分层模概念与模的晶体概念之间,最新颖、最富生机的是后者。正是它注定要主宰德拉姆系数的理论。关于这一点,我“提醒”一句:对 $X$ 上真且光滑的相对概形 $Z$,$Z$ 在 $X$ 上的相对德拉姆上同调(无论在超越语境中还是在代数语境中……)都“不仅”配备一个分层,而且实实在在地配备一道晶体结构,使它在任何无穷小邻域上方“生长”。这是一个至关重要的数学事实,而Deligne早在1969年我离去之前就把它忘了,当时他用分层 pro-凝聚模、而不用更强的晶体版本、即 pro-凝聚模的晶体,来描述德拉姆类型的系数。应该说,我的名字与分层模概念的联系,不如我与模的晶体概念的联系那样显赫(这一概念如此自然,人们简直要赌咒说它必定上溯到上个世纪;而模的晶体概念则外观远不那么“传统”)。关于这一点,见“……与羁绊”(子注 n\textsuperscript{o}~171(viii))中的思考。}。

% ---------------- p.772 ----------------
\origpage{999}\dmark{}我承认,由于未曾就此深思,我还看不太清:对浸没在一个(姑且说)光滑的 $X'$ 中的 $X$,$X$ 上的晶体与 $X'$ 上的晶体之间的确切关系(即便 $X$ 本身光滑,也是如此)\fnote{684}{(*)}{(5月26日)随着余晶体(co-cristal)概念的引入,局面为我大为澄清,下文 D) 中将提及这一概念。}。

\noindent 可以肯定的是,晶体site、或不如说带环层结构的晶体topos $X_{\mathrm{cris}}$,以协变的方式依赖于解析空间 $X$,即:若 $f : X\to X'$ 是解析空间之间的一个态射,就由此导出
\begin{quote}
$f_{\mathrm{cris}} : X_{\mathrm{cris}}\longrightarrow X'_{\mathrm{cris}}$ ;
\end{quote}
由此特别有这些带环topos上模层的一个“直接像”函子。我们想要理解这个运算(尤其在闭浸入 $X\hookrightarrow X'$ 的情形),并理解在什么条件下一个晶体被变成晶体。在闭浸入的情形,我们还想要这个函子正合。这里的思想是这样的:若 $F$ 是导出范畴 $D^{*}(X_{\mathrm{cris}},\mathcal{O}_{X_{\mathrm{cris}}})$ 的一个对象,$F'$ 是它经由 $f_{\mathrm{cris}*}$ 的总体导出函子的像,并且再假定 $X'$ 光滑,那么 $F'$ 为\origpage{1000}完整正则这一条件\textbf{本不应}依赖于把 $X$ 浸没到一个\dmark{}(光滑)空间 $X'$ 中的所选浸入。若情形果真如此,那么就将把 $X$ 上de~Rham--Mebkhout的晶体系数的范畴,定义为(导出范畴的)由前述条件定义的全子范畴(显然在 $X$ 上是局部的)。

\noindent 因此,且不论一项二十年前就该做、而显然至今仍未做的基础工作(关于晶体的模上的基本运算),可以说在 $X$ 是任意解析空间(未必光滑)的情形,只剩下\textbf{两}张照片(而非四张)来描述我们所得的“德拉姆系数”:一是 $\mathrm{Cons}^{*}(X,\mathbb{C})$ 原封不动,二是(目前仍是假设性的、而且照原样我还看不真切的\fnote{685}{(*)}{我在下面提及一张“第五张照片”,它对我而言眼下就清晰得多,用以一种纯代数的晶体措辞捕捉“好的”德拉姆系数,且不加任何光滑性假设仍保意义。这张照片是从某种意义下与de~Rham--Mebkhout照片“对偶”的角度拍摄的。})de~Rham--Mebkhout系数 $\mathrm{DRM}^{*}(X)$ 范畴,我刚为它冒险提出一条定义原则。范畴 $\mathrm{Cons}^{*}(X,\mathbb{C})$,其描述从超越观点看毫无问题,然而一旦进入代数语境就\textbf{消失}了。这就使得廓清 $\mathrm{DRM}^{*}(X)$ 的一个好定义、使其在该语境中仍保意义的需要,变得显而易见。而对我来说同样清楚的是,这张照片的恰当“框架”——它于是看来(至少初看之下)是唯一剩下的一张——是由晶体的模形成的\fnote{686}{(**)}{我把 $X$ 上一个“\textbf{晶体的模}”称为晶体topos $X_{\mathrm{cris}}$ 上的一个模层。因此可以把模的晶体视为晶体的模的特例。}

\noindent 我况且承认,即便在 $X$ 光滑的情形,在Mebkhout对“德拉姆”系数的描述中,我也并不怎么认得过来:用好上帝函子来说,它不尊重自然的乘法结构;而据说,与之相容的是Mebkhout的反函子(下文(b)中将予讨论)。\fnote{687}{(***)}{这个“据说”是一种略显轻慢的说法(近乎“新风格”的腔调……),为的是一笔略过一条漂亮的定理,它照例出自同一位当值无名之辈(但我似乎得知,其年份比好上帝那一枚要新)。例如,对 $K$ 的两个解析闭子空间 $Y$ 与 $Z$,它蕴含局部上同调上下述公式,它漂亮得显然不可能是真的(然而……):
\begin{quote}
$R\Gamma^{\mathrm{alg}}_{Y\backslash Z}(\mathcal{O}_{X})\simeq R\Gamma^{\mathrm{alg}}_{Y}(\mathcal{O}_{X})\otimes^{L}R\Gamma^{\mathrm{alg}}_{Z}(\mathcal{O}_{X})$,
\end{quote}
我敢打赌,某些漂亮的先生们近日就会把它揣进兜里,俨然“仿佛他们一直就知道”——同时等着把它拍卖给其中最漂亮的那一位……}。更何况,这个函子不与“六种运算”交换。\origpage{1001}与Mebkhout系数相联系的直觉\dmark{},初看之下与离散系数相联系的直觉,性质似乎截然不同。从某种观点看,这是一个好处——我们拥有两张从截然不同角度拍摄的照片!这只是使习惯于从其中一个角度观看的人,更难在从另一角度拍摄的照片中辨认自己。

\noindent 事实上,除了已经浏览过的四张照片(我指的是就“德拉姆系数”而言)之外,还有\textbf{第五张}\fnote{688}{(*)}{这样,我就不只是兑现了本注“五张照片”这一标题的许诺:我实际上揭示了\textbf{两系}五张照片,第一系只描述“德拉姆系数”,第二系则描述一般的晶体系数。},我把它留作储备:那是Deligne的那一张,用分层 pro-模拍摄\fnote{689}{(**)}{如先前一条脚注(第998页注(***))所指出,这张德拉姆-Deligne照片是用一个略有扭曲的“镜头”拍摄的(且说其原因,也超出制造工匠的能力所及)。必须把它修版,还要把它放大,让它走出特征零的框框。这将在《数学反思》第三卷中完成,在那里我亲爱的前弟子们可以安心地来汲取一切“无用细节”和别的“技术性离题”,而那些是他们近二十年来一直未曾有闲情自己去找出来的——我放任他们凭一个华丽的主题自己摸爬滚打……}。它的优点,是极其贴近地“贴合”可构离散层的直觉:一个“零次”对象对应于一个同类型的对象,张量积与拉回像的概念通过Deligne的等价相互对应;因此六种运算的整体也将如此(它们确实可以用这两个概念来描述)。另一方面,从“德拉姆-Deligne”系数 $\mathrm{DRD}^{*}(X)$ 到“de~Rham--Mebkhout”系数 $\mathrm{DRM}^{*}(X)$ 的过渡运算,在我看来原则上特别好理解,用 $\mathcal{O}_{X}$-模上的运算(“$\mathcal{O}_{X}$-对偶”)来说(至少一开始,对光滑的 $X$ 是如此)——我在先前一条脚注\fnote{690}{(***)}{这条“先前的脚注”在此期间已变成本注“五张照片”的(c)部分。}中已经提及。因此我有这样一种感觉:在此我处在一片既坚实又熟悉的土地上,\origpage{1002}一旦有暇,它定能使我认明方向。我甚至本想在本注中\dmark{}勾勒Deligne的观点,并把它与Mebkhout的观点、与我在已引的1966年那些讲演中勾勒的形式主义联系起来。但这个子注开始变长,而且越来越变成一处离题支节!因此我宁愿把此事推迟到《数学反思》第三卷,在那里我还想给出 $\mathrm{Spec}\,\mathbb{Z}$ 上有限型概形上“好的”德拉姆系数(Deligne风格或Mebkhout风格,任选其一)的描述。

\subsubsection*{(b) 好上帝公式}

\tocsubsub{(b) 好上帝公式}

\noindent (5月5日与5月21日)我想在此回到对Mebkhout函子(亦称“好上帝函子”)的描述:
\begin{equation*}
M : \mathrm{Cons}^{*}(X,\mathbb{C})\longrightarrow \mathrm{Cris}^{*}_{\mathrm{coh}}(X)\overset{\mathrm{dfn}}{=}D^{*}_{\mathrm{coh}}(X,D_{X})
\tag{1}
\end{equation*}
,其中 $X$ 为一光滑复解析空间。

% ---------------- p.774 ----------------
如“作品……”一注(n\textsuperscript{o}~171(ii))中所说,这是一个性质深邃的函子,它定义为德拉姆函子 $\mathrm{DR}$ 限制到 $\mathrm{Cris}^{*}_{\mathrm{coh}}(X)$ 的全子范畴 $\mathrm{DRM}^{*}(X)$(即 $X$ 上“de~Rham--Mebkhout系数”的范畴)上的函子的拟逆:
\begin{equation*}
m\overset{\mathrm{dfn}}{=}\mathrm{DR}|_{\mathrm{DRM}^{*}(X)} : \mathrm{DRM}^{*}(X)=\mathrm{Cris}^{*}(X)_{\mathrm{hol\,reg}}\longrightarrow \mathrm{Cons}^{*}(X,\mathbb{C}),
\tag{2}
\end{equation*}
而这个函子结果是一个等价(“好上帝定理”)。事实上,Mebkhout还得到了函子 $M_{\infty}$ 的一个卓越的直接描述,它是把函子 $M$ 经过环同态
\begin{equation*}
D_{X}\longrightarrow D^{\infty}_{X}
\tag{3}
\end{equation*}
的“标量扩张”函子 $i$ 导出的,其中 $D^{\infty}_{X}$(或 $D_{X}^{\infty}$)表示“$X$ 上无穷阶微分算子”的环,即(由定义)把层 $\mathcal{O}_{X}$ 看作复拓扑向量空间层时它的($\mathsf{C}$-自同态)环。众所周知,$D^{\infty}_{X}$ 在 $D$ 左、右都忠实平坦,因此环扩张函子
\begin{equation*}
i : \mathrm{Cris}^{*}(X)=D(X,\mathcal{D})\longrightarrow D(X,\mathcal{D})\longrightarrow D(X,\mathcal{D}^{\infty})\overset{\mathrm{dfn}}{=}\mathrm{Cris}_{\infty}(X)
\tag{4}
\end{equation*}
\noindent {\small [译按:原书式(4)连续两次写出 $D(X,\mathcal{D})$,两者之间有箭头;此处照录,其重复疑为排印之误。]}

的总体导出函子由一个寻常的张量积来阐明。且提一句:人们不知道环 $D^{\infty}_{X}$ 是否凝聚,但显然没有它也过得去。人们把由 D-模复形中那些“完整”的组成的全子范畴
\begin{quote}
$\mathrm{Cris}^{*}_{\infty}(X)_{\mathrm{hol}}\hookrightarrow \mathrm{Cris}^{*}_{\infty}(X)$
\end{quote}
定义为:局部地(经由函子 $i$)可从一个完整的 D-模复形 $C$ 导出的条件。\origpage{1003}(由下文重述的双重好上帝定理将可得:那时甚至可以取 $C$ 同时完整且正则,即一个“de~Rham--Mebkhout”系数,并且这将在整个 $X$ 上除唯一同构外确定 $C$……)考虑函子 $M_{\infty}=iM$,它嵌入交换图表
\begin{equation*}
\begin{tikzcd}[row sep=2.5em, column sep=1.2em]
& \mathrm{Cons}^{*}(X,\mathbb{C}) \arrow[dl,"M"'] \arrow[dr,"M_{\infty}"] & \\
\begin{gathered}\mathrm{DRM}^{*}(X)\\[3pt]={}\mathrm{Cris}^{*}(X)_{\mathrm{hol\,reg}}\end{gathered}
  \arrow[rr,"i"] & & \mathrm{Cris}^{*}_{\infty}(X)_{\mathrm{hol}}
\end{tikzcd}
\tag{5}
\end{equation*}
结果发现(或者不如说,是那位无名工匠证明的……)函子 $M_{\infty}$ 本身也是一个范畴等价(因此 $i$ 也是)。人们也可以把它作为“德拉姆”型的函子 $m_{\infty}$(与 $m$ 类似,定义在 $\mathrm{Cris}^{*}_{\infty}(X)_{\mathrm{hol}}$ 上)的拟逆而得到。为了描述函子 $M_{\infty}$,更方便的是描述反函子
\begin{equation*}
\Delta_{\infty}\overset{\mathrm{dfn}}{=}M_{\infty}D=D_{\infty}M_{\infty}=i(MD)=i(DM),
\tag{6}
\end{equation*}
其中 $D$ 表示已提及的对偶化函子,在 $\mathrm{Cons}^{*}$ 或 $\mathrm{DRM}^{*}$ 中;而 $D_{\infty}$ 表示在 $\mathrm{Cris}^{*}_{\infty}(X)_{\mathrm{hol}}$(乃至 $\mathrm{Cris}^{*}_{\infty\mathrm{coh}}(X)$)中也存在的类似的对偶化函子。(注意:(5) 中出现的三个函子都与对偶化函子交换。)于是 $\Delta_{\infty}$ 的拟逆 $\delta_{\infty}$ 由公式

% ---------------- p.775 ----------------
\begin{equation*}
\delta_{\infty}\overset{\mathrm{dfn}}{=}Dm_{\infty}=m_{\infty}D_{\infty}
\tag{7}
\end{equation*}
给出,它与(6)类似。于是人们得到Mebkhout对 $\Delta_{\infty}$、$\delta_{\infty}$ 的表达式,由下面两个对称得惊人的公式给出:
\begin{equation*}
\left\{\begin{aligned}
\Delta_{\infty}(F)&=\mathrm{RHom}_{\mathsf{C}}(F,\mathcal{O}_{X}),\\
\delta_{\infty}(C)&=\mathrm{RHom}_{D}(C,\mathcal{O}_{X}).
\end{aligned}\right.
\tag{8}
\end{equation*}
要注意,在第一个公式中,第二个成员借助于 $D^{\infty}_{X}$ 在第二个参数 $\mathcal{O}_{X}$ 上的运算而继承了一个 $D^{\infty}_{X}$-结构;而在第二个公式中,第二个成员只是被诠释为一个 $\mathsf{C}$-向量空间层复形。这两个公式中的第二个,放在那里“以备查阅”,本质上重言式,它只是说:函子 $\delta_{\infty}$ 把一个 $D^{\infty}_{X}$-模复形 $C$ 结合到(由德拉姆函子 $\mathrm{DR}_{\infty}$ 与 $C$ 相结合的)那个微分算子(无穷阶)复形的“相伴”复形——这个复形被诠释为 $\mathsf{C}$-向量空间层复形。(人们由此确实得到一个上同调层可构的复形,这件事等价于Kashiwara的可构性定理。)

\origpage{1004}\dmark{}(相反,第一个函子 $\Delta_{\infty}$ 把可构层变成完整的(复形的)$D^{\infty}_{X}$-模,却是一条深刻的定理。仅由这一结果蕴含的有限性定理\fnote{691}{(*)}{例如,这个有限性结果蕴含:在 $X$ 上局部地,复形 $\mathrm{RHom}_{\infty}(F,\mathcal{O}_{X})$(在导出范畴中)同构于一个 $D^{\infty}_{X}$-模复形,它在每个次数上都是局部自由有限型的,并且它的上同调模(局部地)经标量扩张来源于凝聚的 D-模。事实上,甚至可以假定后者是完整且正则的。}(且不说完整性)本身就已经是一个引人注目的新结果。然而更加非凡的事情,\textbf{是这两个函子互为拟逆}。形式上,这件事使人想起双对偶关系,它既可以在范畴 $\mathrm{Cons}^{*}$ 中表达,也可以在范畴 $\mathrm{Cris}^{*}_{\infty}(X)_{\mathrm{hol}}$ 中表达——只不过(在两种情形中都表示为一个 $\mathrm{RHom}_{\infty}(-,\mathcal{O}_{X})$ 的)“对偶化”反函子把两个\textbf{不同的}范畴彼此相连。正是这种形式上的类比,使Mebkhout把断言同构
\begin{equation*}
\Delta_{\infty}\delta_{\infty}\simeq \mathrm{id}\qquad \text{在 }\mathrm{Cris}^{*}_{\infty}(X)_{\mathrm{coh}}\text{ 中}.
\tag{9}
\end{equation*}
的定理称为 $D^{\infty}_{X}$-模复形的“\textbf{双对偶定理}”(这个术语反正有引起混淆的危险)。这个关系,加上函子 $\delta_{\infty}$ 全忠实这一事实(或更确切地说,$\Delta_{\infty}$ 是它的一个伴随,这一点他把它纳入自己的双对偶定理的陈述中),Mebkhout早在1977年、在完整的好上帝定理之前就已得到。所谓的“双对偶”定理因此本质上意味着(正如它所启发的“我的”双对偶定理一样):一个\textbf{完整的} $D^{\infty}_{X}$-模复形,可以凭借对相结合的微分算子(无穷阶)复形的知识——即把该复形仅仅看成一个 $\mathsf{C}$-向量空间层复形(在适宜的导出范畴中)——把它\textbf{重构}为导出范畴的一个对象;更确切地说,可以凭借显式的反演公式(8)(第一个公式)来重构。更何况,完整 $D^{\infty}_{X}$-模复形之间的一个态射是拟同构,当且仅当无穷阶微分算子复形之间相应的态射在朴素意义下是拟同构(即在同调层上诱导一个同构)\fnote{692}{(**)}{(5月26日)事实上(如我在下文(c)开头所指出)Mebkhout证明这后一个结果时,甚至不借助任何完整性条件,其等价形式为:若与一个 $D^{\infty}_{X}$-模复形相结合的微分算子复形是拟零的,则该复形也是拟零的(D-模也是如此)。}。

% ---------------- p.776 ----------------
\dmark{}\origpage{1005}(Mebkhout的双对偶定理,在某种意义下构成好上帝定理的“一半”(就 $D^{\infty}_{X}$-模而言),当后者取其最强形式、即断言函子(8)互为拟逆时,这一点如此。)这是Mebkhout论文的核心结果,该论文于1980年1月提交。但这“一半”单独而论,已经是一个新结果,而且(据我所知)完全出乎意料。它是一个典型的结果,在Sato的观念与我的观念之间架起桥梁,但却是在我由来已久的纲领的视野中:以“连续”或“微分”的路子(并在导出范畴的视野中)来表述“离散系数”。就此而言,我觉得这个结果以其精神与灵感而论,完全逃出于日本分析学派的问题意识之外。Kashiwara的可构性定理似乎在这一点上只是一个“侧面”,而绝非一门新系数理论的出发点。正如1976至1980年间的出版物不容置疑地证明的那样,Mebkhout是当时唯一发展这样一种哲学的人。

\noindent 1978年1月,Mebkhout在巴黎过境时曾向Kashiwara讲述他的结果,当时他刚刚结束论文的撰写。应Kashiwara之请,天真的Mebkhout——他为终于找到一个似乎对他要说的东西感兴趣的人而满心欢喜——把还热乎乎的第三章寄到了普林斯顿,那一章里除别的之外就有所谓“双对偶”定理。那是1978年2月。三年之后,同一个结果(连同一个假装是证明的东西)出现在Kashiwara--Kawai的一篇著名文章中\fnote{693}{(*)}{M.~Kashiwara, T.~Kawai,\emph{On holonomic Systems of micro-differential équations}, III \emph{Systems with regular singularities},Publ. RIMS 17, 813--979 (1981). 从Mebkhout那里剽窃来的“重构定理”,见于这部长篇著作(1980年11月收到)的第4节。该著作的主要结果,是(5)中函子 $i$ 是范畴等价这一事实的一个弱化变体。因此这是Mebkhout(几何)理论的一个直接推论,是这些作者用解析的路子(独立于Mebkhout)得到的结论。详见“黑手党”子注 n\textsuperscript{o}~171(ii),(b)部分:“最初的麻烦——或太平洋对岸的卡伊德”。译按:此处171(ii)照录原书;所引标题“黑手党”实际对应171$_{2}$,疑为原书自引误。}。它为此番场合被重新命名为“reconstruction theorem”,而对某一位Zoghman~Mebkhout只字不提。那年也正是反常层会议值得纪念的一年——那光荣的一年,某种“新风格”\fnote{694}{(**)}{关于这种“新风格”(Kashiwara与Hotta是它在太平洋对岸的杰出信徒),见“祝贺——或:新风格”一注(n\textsuperscript{o}~169\textsubscript{9})。}强行征服了(且未遇丝毫抵抗……)数学的这一个部分,而在所有部分中,这里正是我当年惯常感到如鱼得水的地方……

\origpage{1006}
\subsubsection*{(c) 第五张照片(用“pro”)}

\tocsubsub{(c) 第五张照片(用“pro”)}

\dmark{}(5月21日)“双对偶定理”(9)属于1977年。为了证明 $D^{\infty}_{X}$-模“好上帝定理”的另一半——那时归结为证明函子 $\delta_{\infty}$ 本质上满射,第一个困难是证明:对 $\mathrm{Cons}^{*}$ 中的 $F$,用第一个公式(8)定义 $D^{\infty}_{X}$-模复形 $C=\Delta_{\infty}(F)$,则它至少局部地可以经由函子 $i$、借助于一个(完整、正则的)D-模复形得到。先验地,按照Mebkhout的观念(即遵循双重好上帝定理,它蕴涵(5)中的函子 $i$ 是一个等价),后者除唯一拟同构外必定唯一。

\noindent 我不曾试图理解Mebkhout最终在他的论文中是如何设法构造出这个 D-模的。在我看来,这里的情形应当用Deligne的观念来澄清,即与一个可构 $\mathsf{C}$-向量空间层 $F$ 相结合的 pro-凝聚层\fnote{695}{(*)}{这是他在1969--70年 IHÉS 讨论班上发展过、随后弃置不顾的观念。关于这一点,见“……与羁绊”子注(n\textsuperscript{o}~171(viii))。}。这个观念是他在 $X$ 上\textbf{代数簇}的语境中发展的,但必须能够 mutatis mutandis(作必要的变更)适用于解析情形,条件也许是只能在 $X$ 上“局部地”工作,或在 $X$ 的每个紧集上工作。与 $F$ 相结合的 pro-凝聚层,因此(至少在每个紧集 $K\subset X$ 上)是凝聚层的一个投射系统 $(F_{i})$(定义在 $K$ 的邻域上),它可以很简单地定义为 pro-表示以下函子的层
%JOIN:para
% ---------------- p.777 ----------------
\begin{quote}
$G\mapsto \mathrm{Hom}_{\mathsf{C}}(F,G)$
\end{quote}
(在 $X$ 上(在 $K$ 的邻域……)凝聚 $\mathcal{O}_{X}$-模 $G$ 的范畴上),这个函子正合于左,因此确实是 pro-可表示的。例如,若 $F$ 是 $X$ 的一个解析闭子空间 $Y$ 上的常值层 $\mathsf{C}_{Y}$,“以零延拓”到整个 $X$ 上,就得到由那些 $\mathcal{O}_{X_{n}}$ 组成的 pro-层,其中 $X_{n}$ 是 $Y$ 在 $X$ 中的无穷小邻域。(注意:这个投射系统的投射极限,是 $\mathcal{O}_{X}$ 沿 $Y$ 的形式完备化。)人们注意到(回到一般情形)pro-层 $(F_{i})$ 配备一个典范分层\fnote{696}{(**)}{pro-模的分层概念,与模的分层概念定义方式相同——即在前一日的注中(a)部分给出的描述,原则上每当有一个“相对”概念(如模、pro-模、相对概形等等)容许一个“拉回”概念、即给出一张以“簇”为基的“纤维范畴”时,就适用……要注意,若 $(F_{i})$ 是一个 pro-模,它的一个分层一般说来不能用一个 $F_{i}$ 的分层“相容”系统来描述——所考虑的对象,其性质远比分层模范畴的 pro-对象更为一般。}。\origpage{1007}Deligne的观念是:从 $X$ 上可构 $\mathsf{C}$-向量空间层的范畴、到分层 pro-凝聚层范畴的“\textbf{Deligne函子}”\dmark{},是\textbf{全忠实}的,从而使人能把前一个范畴(它性质上是超越的)诠释为分层 pro-凝聚层范畴的一个全子范畴。后者有一个纯代数的意义,而所说的全子范畴也能(以多少是重言式的方式\fnote{697}{(*)}{“重言式”,至少就这部已知的词典(最初由Deligne廓清)而言是如此:它在解析空间 $Y$ 中除子 $Z$ 的补集 $Y-Z$ 上局部常值的 $\mathsf{C}$-向量空间层(或“局部系统”),与 $Y-Z$ 上沿 $Z$ “正则”(Deligne意义下)的凝聚分层模之间。})用纯代数的术语来定义。这就是我将记作
\begin{quote}
(10)\quad $\mathrm{DRD}^{*}(X)$ 或 $\mathrm{Del}^{*}(X)$
\end{quote}
的范畴,它构成“\textbf{第五张照片}”,而我昨天不愿把它讲明白\fnote{698}{(**)}{最终,这一阐明(被称作“重言式”的!)在此也没有给出,至少没有当场给出。不过它将在更后面(第1011页)给出。要注意,记号(10)指的是“导出范畴”变体。}。况且我记得,Deligne当初费心把他的诠释(以及前述的全忠实性陈述)发展到能过渡到导出范畴(那时我的同调学生们、Deligne打头,还没有一致决定把这些范畴甩卖掉),而我用记号(10)所指的,当然正是“导出范畴”的版本。

\noindent 话虽如此,$\mathrm{RHom}_{\mathsf{C}}(F,\mathcal{O}_{X})$ 中的“代数部分”必须能够以一种十分自然的方式定义为一个(在适宜意义下的)归纳极限,即 $\mathrm{RHom}_{\mathcal{O}_{X}}(F_{i},\mathcal{O}_{X})$ 的归纳极限——并且特别地(过渡到上同调层),人们描述出典范箭头
\begin{quote}
(11)\quad $\varinjlim_{i}\mathrm{Ext}^{d}_{\mathcal{O}_{X}}(F_{i},\mathcal{O}_{X})\longrightarrow\mathrm{Ext}^{d}_{\mathsf{C}_{X}}(F,\mathcal{O}_{X})$ \quad ($\forall d\in\mathbb{Z}$)
\end{quote}
利用 pro-对象 $(F_{i})$ 上的分层与第二个参数 $\mathcal{O}_{X}$ 的重言式分层\dmark{},人们必须能够在(11)的第一个成员上定义一个分层、即一个 D-模结构,使得(11)与(相应的算子环的同态\origpage{1008}$D\to D^{\infty}_{X}$)相容。话虽如此,Mebkhout的好上帝定理必须能够精确化,即说(11)把第二个成员等同于由第一个成员经标量扩张而得到的 $D^{\infty}_{X}$-模\fnote{699}{(*)}{此外,当然,(11)的第一个成员(与Mebkhout的哲学一致)必须是一个\textbf{凝聚}、\textbf{完整}且正则的 D-模。}——这特别蕴含该箭头是一个\textbf{包含}。因此,左边一个成员应当被看成是右边一个成员中的某种\textbf{“代数部分”}(或“\textbf{亚纯部分}”),而后者性质上是
“超越的”。

\noindent 一般情形由前面这个特例 $F=i_{*}(\mathsf{C}_{Y})$(其中 $i : Y\to X$ 是 $X$ 的一个解析闭子空间的包含映射)而得到大量澄清。于是(11)的第二个成员是一个支撑在 $Y$ 中的局部上同调层——一个\textbf{超越的}不变量,而第一个成员
\begin{quote}
$\varinjlim_{n}\mathrm{Ext}^{d}_{\mathcal{O}_{X}}(\mathcal{O}_{X_{n}},\mathcal{O}_{X})$,
\end{quote}
却是我当年在概形框架中为局部上同调引入的那个众所周知的表达式。这个层在一点 $x\in Y$ 上的纤维,不是别的,正是结构层在 $Y$ 于 $X_{x}$ 上的“迹” $Y_{x}$ 中的支撑上、在 $\mathcal{O}_{X,x}$ 的谱 $X_{x}$ 上的局部上同调。

\noindent 从这个例子可以看出,Deligne的观念与我六十年代初就局部上同调主题所发展的那些观念何等接近\fnote{700}{(**)}{下文将显出,Deligne的观念也与我在1966年(在[Crystals]中)引入的观念密切相关:对任何微分算子复形,我把它的“形式化” $P^{\infty}(L^{\bullet})$ 看成一个分层 pro-模复形,或者更好地说,看成一个\textbf{晶体复形}的定义,其(整体)晶体上同调等同于 $L^{\bullet}$ 的(整体)上同调。}。无论如何,Mebkhout1972至1976年间工作的主题,正是在这个关键情形下研究箭头(11)
\begin{equation*}
\varinjlim_{n}\mathrm{Ext}^{d}_{\mathcal{O}_{X}}(\mathcal{O}_{X_{n}},\mathcal{O}_{X})=H^{d}_{Y}(\mathcal{O}_{X})^{\mathrm{alg}}\longrightarrow H^{d}_{Y}(\mathcal{O}_{X})
\tag{12}
\end{equation*}
他在这种情形下证明了上文宣告的关系,此外还证明了(这是我方才一度忘记纳入陈述的)(12)的第一个成员是一个\textbf{凝聚的} D-模,乃至完整且正则。\origpage{1009}由此起,(11)的类似陈述必须是经 dévissage(拆卸)而来的直接结果\fnote{701}{(***)}{(5月22日)我在这里有点急躁!可构 $\mathsf{C}$-层的“基本类型”比单纯的 $\mathsf{C}_{Y}$ 性质更为一般,(但一般定理的证明确实用了与1976年那个特例相同的技术。)},\dmark{}包括在 $F$ 不是可构 $\mathsf{C}$-向量空间层、而是 $\mathrm{Cons}^{*}(X,\mathbb{C})$ 中一个复形的情形下也是如此。唯一的一粒盐,除了Deligne函子形式的构造之外,是在一个分层 pro-模复形、取值于一个分层模复形(即取值于一个 D-模复形,此处是 $\mathcal{O}_{X}$)的 $\mathrm{RHom}_{\mathcal{O}_{X}}$ 的定义中,是作为一个 D-模复形(且作为导出范畴的一个对象)来取的。

\noindent 且不论这一粒盐,人们于是得到好上帝函子 $M$“代数的”(相对于好上帝函子 $M_{\infty}$“超越的”而言)、或者不如说相联系的反函子 $\Delta$ 及其拟逆 $\delta$
\begin{quote}
(13)\quad $\Delta=MD=DM$ ,\quad $\delta=mD=Dm$,
\end{quote}
的一个至为简单而概念化的描述,用的是一条双重公式,它把(8)复述一遍。但为了把它写出来,要利用Deligne的等价
\begin{equation*}
\mathrm{Del} : \mathrm{Cons}^{*}(X,\mathbb{C})\xrightarrow{\ \simeq\ }\mathrm{DRD}^{*}(X),
\tag{14}
\end{equation*}
人们宁愿看相应的函子 $\hat{\Delta}$、$\hat{\delta}$,它们在 $\mathrm{DRD}^{*}(X)$ 与 $\mathrm{DRM}^{*}(X)$ 之间,其中的帽子符号意在提醒:人们将在(“可构”那一侧)用 \textbf{pro-对象}来工作。于是人们得到那些卓越的公式(实质上包含在(8)中,但这一次联系的是“代数性质”的系数双方,并且用的也是“代数性质”的公式):
\begin{quote}
(15)\quad $\hat{\Delta}(C')=\mathrm{RHom}_{\mathcal{O}_{X}}(C',\mathcal{O}_{X})$ ,\quad $\hat{\delta}(C)=\mathrm{RHom}_{\mathcal{O}_{X}}(C,\mathcal{O}_{X})$。
\end{quote}
这里因此两次出现“同一个”公式,唯一的区别在于 $C'$ 此处是一个分层 pro-凝聚层复形(或者说\fnote{702}{(*)}{关于这一译法,见第1006页脚注(**)。},一个 pro-凝聚模的晶体复形),而 $C$ 是一个 D-模复形(本质上可以把它看成一个分层 ind-凝聚 $\mathcal{O}_{X}$-模复形,或者再看成一个 ind-凝聚模的晶体)。本质上正是“同一个”函子使人们从一类走到另一类,即“寻常的对偶化”(凝聚的)函子,我五十年代的老友……诚然,“显然”它必须交换 pro-对象与 ind-对象(不妨在后者中取归纳极限……)。

\origpage{1010}\dmark{}(当然,有一项基础工作要做,为的是给这些公式一个精确的意义——一项Deligne在他那个著名的沉没讨论班、或儒瓦诺卢在他那篇同样著名的沉没论文中所做的那类工作(自从反常层会议以来人人都引它,却没有一个人把它捧在手里……)。这是一项我确信或许会有点长、但本质上“连环推理”的工作。“硬”的部分包含在Mebkhout的好上帝定理中,由Mebkhout的公式(8)(也许不恰当地)称作“双对偶”公式者补全。而它们的代数翻译,即断言两个函子(15)互为拟逆,却实实在在是(实质上)凝聚 $\mathcal{O}_{X}$-模系数的“那个”寻常双对偶定理,只是佐以 ind-pro 的风味、并以分层为关键罢了(它们必须在对偶化函子中“通行”无阻)。

% ---------------- p.779 ----------------
\noindent 两类对偶对象之间的对应,用微分算子复形来观照则完美无遗(且不须任何基础工作作关键!)。(在这场对偶中,况且完整性条件(进而正则性条件)不扮演任何角色。)对这样一个复形 $L^{\bullet}$,把昨天(在(a)、(1)中)考虑的那个(反变的)函子
\begin{quote}
$F\mapsto \mathrm{Hom}_{\mathcal{O}_{X}}(F,\mathcal{D}_{d})$
\end{quote}
结合到一个分量局部自由有限型的 D-模复形,即 $C$。另一方面,把这个复形 $L^{\bullet}$“形式化”,过渡到无穷阶主部 $P^{\infty}(L^{i})$(被看作分层的 pro-模),就提供一个分层 pro-模复形 $C'=P^{\infty}(L^{\bullet})$。话虽如此,人们看到这两个复形由公式(15)相互对应,在其中此处显然 $\mathrm{RHom}$ 归结为 $\mathrm{Hom}$。(只须对各分量 $L^{i}$ 逐项验证这一对偶性,它就归结为下述多少是重言式的事实:从 $P^{\infty}(L^{i})$ 到 $\mathcal{O}_{X}$ 的“连续”线性同态,与从 $L^{i}$ 到 $D$ 的线性同态一样,恰好对应于从 $L^{i}$ 到 $\mathcal{O}_{X}$ 的微分算子,分别利用“万有”的(无穷阶)微分算子 $L^{i}\to P^{\infty}(L^{i})$,与由 $\theta\mapsto\theta(1)$ 给出的“增广” $D\to\mathcal{O}_{X}$。)由于至少局部地在 $X$ 上,$\mathrm{Cris}^{*}_{\mathrm{coh}}(X)$ 的每个对象(即每个上同调凝聚的 D-模复形)都可以借助一个微分算子复形 $L^{\bullet}$ 来描述,因此可以认为,只要对 $C$ 与 $C'$ 作适宜的 D-凝聚性与“D-pro-凝聚性”假设、彼此“对偶”,这个特例就在一切实际用途上完美地掌握了两类系数之间的对偶性(15)。\origpage{1011}于是只须发展我所提到的“连环推理”,在 $C'$ 或“pro”一侧局限于分层 pro-凝聚层复形\dmark{}(它们局部地可以(除拟同构外)描述为一个 $P^{\infty}(L^{\bullet})$)。

\noindent 相对于Deligne原本的进路,他所引入的 pro-凝聚模及其复形能够局部地用一个微分算子复形来实现,这件事本身是一个\textbf{全然出乎意料的现象},由Mebkhout的理论所带来。在我看来它本质上是
等同于\fnote{703}{(*)}{(5月26日)在这里我又“有点急躁”了,1976年的那个结果并不足够。请与第1008页脚注(***)的评注相比较。}上面提到的那个Mebkhout定理(可追溯至1976年,早在好上帝定理的证明之前),该定理关乎出现在上方(12)中的层 $H^{d}_{Y}(\mathcal{O}_{X})^{\mathrm{alg}}$ 的 D-凝聚性。那是一条深刻的定理,是四年工作的结晶,用尽了广中奇点消解的全部力量(更不消说那位廓清并证明了它的工匠的勇气——顶着普遍的冷漠)。我刚才指出的那个后果\textsuperscript{703(*)},是德拉姆系数(如我自1966年起所窥见的那样)与微分算子复形之间一层深刻的关系,这层关系我全然未曾预见(Deligne发展他通往德拉姆系数的第一条进路时,也同样未曾预见)。至于所考虑的微分算子复形上的完整性与正则性条件,它(事后而论,多亏那救场的好上帝定理)应当等同于Deligne的“有限性”(加“正则性”)条件(我在引入范畴 $\mathrm{DRD}^{*}(X)=\mathrm{Del}^{*}(X)$ 时一度未曾把它讲明)。条件如下:$P^{\infty}(L^{\bullet})$ 的上同调 pro-层在局部上由合成列“拆卸”,使得相继因子能(经由Deligne函子)由 $X$ 的子空间 $Y-Z$ 上的 $\mathsf{C}$-向量空间局部系统来描述(其中 $Z\subset Y\subset X$ 是 $X$ 的解析闭子空间)。为了最终给这条判据一副“代数”的面貌,只需把 $\mathsf{C}$-向量空间局部系统换成 $Y-Z$ 上一个\textbf{凝聚}分层,并附加条件:表达该分化的联络(注意可设 $Y-Z$ 光滑)在 $Z$ 的邻域上是“正则的”,取Deligne的意义\fnote{704}{(**)}{考虑到Deligne所廓清的范畴等价,即在 $Y-Z$ 上的 $\mathsf{C}$-向量空间局部系统、与 $Y-Z$ 上配备沿 $Z$ 的“亚纯结构”、且沿 $Z$ 联络正则的可积联络丛之间的等价,这条正则性条件在此以一种自然的方式引入。这种亚纯结构(蕴含把 $Y-Z$ 上的凝聚模延拓为 $Y$ 上的凝聚模的可能性,至少在 $Z$ 的每点邻域上如此)在方才给出的描述中是暗含的。除非我弄错,当在前一条条件中去掉正则性条件时(只需假定在 $Z$ 的邻域上给定 $E$ 的一个亚纯结构,以便能借Deligne的办法给它结合一个整个 $X$ 上的 pro-凝聚模),人们就得到完整性条件的一个“上同调式”描述。Sato的定义是借“微局部”的路子作出的——我承认,我还从未真正领会过它……}。(注意:相联系的 pro-层可以这样得到:\origpage{1012}把 $Y-Z=T$ 上拥有的\dmark{}晶体在 $T$ 的无穷小邻域上方养大,并沿 $Z$“压碎”,以便处处都有凝聚的层,而不只在 $Z$ 的补集上……)

\subsubsection*{(d) 晶体与余晶体——全忠实?}

\tocsubsub{(d) 晶体与余晶体——全忠实?}

\noindent 当不再假定 $X$ 光滑时,因此,除了性质上超越的“照片” $\mathrm{Cons}^{*}(X,\mathbb{C})$ 之外,还剩下两张“照片”(两者性质上都是晶体的)$\mathrm{DRM}^{*}(X)$ 与 $\mathrm{Del}^{*}(X)$,而它们却具有纯代数的意义。我昨天(在(a)中)为 $\mathrm{DRM}^{*}(X)$ 草拟了一条定义原则,今天又为范畴 $\mathrm{DRD}^{*}(X)$ 拟了一条。后者从眼前起对我来说就完全可理解。如我昨天所指出(见(a),第998页脚注(***)),这里应当把分层 pro-模的观点,用 pro-模(pro-凝聚的)\textbf{晶体}的观点来细化\fnote{705}{(*)}{(5月27日)细想起来,我甚至难以相信Deligne的定理 $\mathrm{Cons}^{*}(X,\mathbb{C})\cong\mathrm{Del}^{*}(X)$ 对非光滑的 $X$ 也成立,如果 $\mathrm{Del}^{*}(X)$ 像Deligne所做的那样、不借助晶体site来定义的话。也许正是为了意识到这一点,他最终宁愿把整套理论沉没掉,而不肯同意重新引入那个犯忌的位……(请与“……与羁绊”一注(n\textsuperscript{o}~171(viii))相比较。)}。这个观点眼下仅存的那点麻烦,就是它将要迫使人们发展的那套“pro”连环推理,而它(据我在此类事情上微薄的经验)有变成骇人规模的风险!这些 pro-模的晶体,给 $X$ 的每个开集 $U$ 的每个无穷小加厚 $U'$ 结合一个 $U'$ 上的 pro-凝聚模,“以与拉回像相容的方式”对应于加厚之间的态射 $U''\to U'$,它们甚至不能被诠释为晶体site(或同样地说,晶体
%JOIN:para
% ---------------- p.781 ----------------
topos $X_{\mathrm{cris}}$)上的 pro-层!因此人们先验地不能把诸如 $X_{\mathrm{cris}}$ 这样的(交换地)带环topos上模层的已知上同调形式主义应用到它们身上。

\noindent 这里的诱惑很大,就是取人们在每个加厚上拥有的那个 pro-层的投射极限。这样得到的是晶体的模(而非模的晶体),其“值”在每个 $U'$ 上既非凝聚也非拟凝聚。希望在于,至少对\origpage{1013}我们感兴趣的那类 pro-模晶体(尤其是那些由Deligne函子得到的)而言,这样一个 pro-模晶体能够从经过取极限而得到的晶体模 $C$\textbf{重构}出来,办法是在每个加厚 $U'$ 上取 Zariskian 层 $C_{U'}$($C$ 在 $U'$ 的 Zariski 开集上的限制)的“pro-凝聚包络”\fnote{706}{(*)}{我在这里谈到“Zariskian 层”(与“晶体的”相对)时,已悄悄滑回了概形的语境。偏爱解析语境的读者自会加以校正。}。在我看来这确实如此,至少对如上所说联系于 $Y-Z$ 上一个凝聚分层模的 pro-模晶体是这样,例如在取 $\mathcal{O}_{X}$ 沿 $Y-Z$ 的形式完备化、并在别处以零延拓(在加厚上亦然)的那个典型情形中就是这样。

\noindent 如果我的“希望”得到证明,那么 $X$ 上德拉姆-Deligne系数的范畴 $\mathrm{DRD}^{*}(X)$ 就可以诠释为寻常导出范畴 $D^{*}(X_{\mathrm{cris}},\mathcal{O}_{X_{\mathrm{cris}}})$ 的一个全子范畴,以上同调层上“有限性”与“正则性”类型的条件(它们本身如上所述以拆卸来描述)来定义。那将是一个令人困惑的简单描述,我早在1966年就可以把它给出,如果当时我有暇继续我的晶体思考的话……

\noindent 这个“基础的”问题(取极限是否合法)显然不依赖于 $X$ 是否光滑的问题——如果它不光滑,就把它浸没在一个光滑的 $X'$ 中而归结到光滑情形。如果这个观点(简直美得不像真的!)确实行得通,那么(现在在光滑情形下)顺便就该(我认为)把“双对偶”公式(代数版本)(15)诠释为 $\mathrm{RHom}_{\mathcal{O}_{X}}$ 的\textbf{寻常}公式,不必为 pro-问题费神(而只需小心搬运分层……)。这方面的第一个检验如下:若 $u : C_{1}\to C_{2}$ 是上同调凝聚的 D-模复形之间的一个态射,它经由朴素的对偶化函子 $\mathrm{RHom}_{\mathcal{O}_{X}}(-,\mathcal{O}_{X})$ 的像是一个拟同构,那么 $u$ 是否也是拟同构?但这一点(借一个映射柱的论证)归结为问:一个上同调凝聚的 D-模复形,若其“朴素对偶”为零(取导出范畴的意义,即上同调层为零),它本身是否也为零(取同样的意义)?或者再说一遍,如果有微分算子复形 $L^{\bullet}$,说相结合的 D-模复形上同调层为零,与说“形式化”的 $P^{\infty}(L^{\bullet})$(这一次不看作一个 pro-层复形,而看作一个寻常层复形(经取 $\varprojlim$))上同调层为零,是否一回事?Mebkhout肯定能告诉我……

\dmark{}\origpage{1014}(5月23日)昨晚我又给Mebkhout打了电话——况且数学的或历史的问题,我几乎每晚都给他打电话,已经有一两个星期了,加起来可要创下一笔天文数字的电话账单了!但《封神》,我这三个紧巴巴的星期以来孜孜不倦而又擦拭不休的那篇,值这个价……

\noindent 无论如何,Zoghman向我保证了一个结果,它看上去同我昨夜收尾时写到的那个“检验问题”相近:若 $C$ 属于 $\mathrm{Cris}^{*}_{\mathrm{coh}}$ 且使得相结合的算子复形 $L^{\bullet}=\mathrm{DR}(C)$ 是拟零的,则 $C$ 本身是拟零的(解析情形)。人们有一个层(的 $\mathsf{C}$-向量空间)复形之间的同态,由“无穷阶主部”给出
%JOIN:para
% ---------------- p.782 ----------------
\begin{quote}
$L^{\bullet}\longrightarrow P^{\infty}(L^{\bullet})$,
\end{quote}
由此有上同调层上的同态
\begin{quote}
(16)\quad $H^{i}(L^{\bullet})\longrightarrow H^{i}(P^{\infty}(L^{\bullet}))$\quad ($i\in\mathbb{Z}$)
\end{quote}
人们很想说这个同态(16)总是单的,并把第一个成员等同于第二个的“水平”截面子层(这将是某种“水平截面层”函子在适宜的分层 pro-模范畴上的正合性……)。单性便将蕴含:若第二个成员为零,则第一个成员也为零,因此若这一切对每个 $i$ 都成立(且据Mebkhout向我保证的那样),那么与 $L^{\bullet}$ 相结合的 D-模复形是拟零的——这正是我想要的。

\noindent (16)中的单性也意味着:对一个微分算子 $E\overset{d}{\longrightarrow}F$,以及 $F$ 的一个截面 $f$,若 $f$ 在每点 $x\in X$ 上“形式上”属于像中(过渡到该点的完备局部环),此外“形式解”(即方程 $d(g)=f$ 在 $g$ 中的解)可以取成对变动的 $x$ 解析地依赖,那么该方程局部地有解。Mebkhout告诉我,他不知道有这样的结果;然而这个问题如此自然,答案想必应该是已知的!

\noindent 为了结束“五张照片”,我想在此再回到那两张“晶体照片”上,一张对应于Mebkhout的 D-模观点,另一张对应于与之对偶的观点。不言而喻,应当在导出范畴的精神中工作——因此一种名副其实的“晶体”诠释必须顾及这一点。所以那两张晶体照片只有在\origpage{1015}那相应的函子\dmark{}(从范畴 $D^{b}_{\mathrm{coh}}(X,D)$(姑且说)出发、到一个适宜的晶体范畴、例如 $D^{b}(X_{\mathrm{cris}},\mathcal{O}_{X_{\mathrm{cris}}})$)本身是全忠实的时候,才是“全忠实”的。我希望情形正是如此,\textbf{甚至不必}给所考虑的 D-模复形\textbf{添上完整性与正则性条件}。

\noindent 最简单的情形无疑是第四张照片,它在于把 D-模的范畴诠释为模的晶体的范畴,由此有一个总体导出函子(不妨说是“Grothendieck的”——以抢在那些“无用细节”与“技术性离题”的爱好者之前……):
\begin{quote}
(17)\quad $G : D^{*}_{\mathrm{coh}}(X,D)\longrightarrow D^{*}(X_{\mathrm{cris}},\mathcal{O}_{X_{\mathrm{cris}}})$。
\end{quote}
这里关键的问题,是这个函子是否全忠实。只有在此时,记号 $\mathrm{Cris}^{*}_{\mathrm{coh}}(X)$ 用于第一个成员才是完全有道理的——并且同时,德拉姆上同调(至少在此情形中,在复解析框架或特征零的域上代数概形的框架中)的晶体观点也才有道理。为了证明全忠实性,比如说在代数几何中,经标准论证可归结到 $X$ 是仿射的情形(或在解析情形中,归结到一个多圆盘的情形),以及归结到第一个成员中所考虑的两个对象 $C$、$C'$(问题在于比较它们在一个方向与另一个方向上的 $\mathrm{Hom}$)都等于 $D$ 本身、只差一个次数平移的情形。(这个归结毫无问题地作出,至少假定 $C$、$C'$ 次数有界,从而局限于 $D^{b}_{\mathrm{coh}}(X,D)$,而这对应用而言看来是大为足够的。)于是最终被引导去验证公式
\begin{quote}
(18)\quad $\Gamma(X,D_{X})\xrightarrow{\ \sim\ }\mathrm{Hom}(G(D),G(D))$ ,\quad $\mathrm{Ext}^{i}_{\mathcal{O}_{X_{\mathrm{cris}}}}(X_{\mathrm{cris}};G(D),G(D))=0$\quad 对 $i>0$
\end{quote}
($X$ 仿射,相应地 Stein)。我不曾有暇去验证它\fnote{707}{(*)}{我有借口:一年多以来,我最清醒的时间都用来追踪我昔日学生中某几位的业绩了……},但毫不怀疑它是真的。某个非常接近的东西,我似乎在[Crystals](1966年)中已经证明过了\fnote{708}{(**)}{这是我已在别处提到过的结果:对概形 $X$ 上的一个微分算子复形 $L^{\bullet}$,(在 $X$ 对 $S$ 光滑相对的框架中,或在解析框架中无疑也同样) $L^{\bullet}$ 的“Zariskian”超上同调等同于它的形式化 $P^{\infty}(L^{\bullet})$ 的晶体超上同调。诚然,这个陈述更直接地关乎(17)的“对偶”箭头(19),也可以这样表达:对 $C$、$C'$ 上同调凝聚的 D-模复形,箭头 $\mathrm{Hom}(C,C')\to\mathrm{Hom}(G^{\circ}(C),G^{\circ}(C'))$ 在 $C=\mathcal{O}_{X}$ 的情形下是双射(这已经不错,且容许一切希望……)。}。

% ---------------- p.783 ----------------
\dmark{}\origpage{1016}(至于第五张照片,它有几种不同的印片。Deligne原初的印片是用分层 pro-凝聚模作出的。为了推广到 $X$ 非光滑的情形,第一次重要的修版,在于把所说的那些动物诠释为 pro-模的\textbf{晶体}。但那样就卷入了(毫无吸引力的!)无休止的 pro-上同调代数 pro-基础的齿轮——并且失去直接依附于 $X_{\mathrm{cris}}$ 的topos式直觉的好处。所以我宁愿(如果可能的话)干脆另取一张照片,视角几乎相同,经由一个\textbf{反变}函子(同样不妨说是“Grothendieck的”,识者自会领会……)
\begin{quote}
(19)\quad $G^{\circ} : D^{*}_{\mathrm{coh}}(X,D)^{\mathrm{opp}}\longrightarrow D^{*}(X_{\mathrm{cris}},\mathcal{O}_{X_{\mathrm{cris}}})$。
\end{quote}
可以说,这是从Deligne的照片出发、在 $X$ 的每个开集 $U$ 的每个无穷小加厚上粗暴地取投射极限层而得到的那一个。若第一个成员中的 $C$(以反变的方式,如(a)中公式(1)那样)结合到一个微分算子复形 $L^{\bullet}$,它经(19)的像可以这样得到:把 $P^{\infty}(L^{\bullet})$(复形 $L^{\bullet}$ 的“形式化”)看作一个分层 pro-模复形(在[Crystal]中引入的观念),或者再看成一个 pro-模的晶体复形,并在任何加厚上取投射极限。另一种说法是:对 $X$ 上任何局部自由的 $\mathcal{O}_{X}$-模(例如)$L$,结合一个晶体模(它\textbf{不是}一个模的晶体,除非我弄错),我把它记作 $P^{\infty}(L)_{\mathrm{cris}}$,其方式固然是“显然的”(而我的学生们早已把它忘掉了),该模以微分算子为中介函子性地依赖于 $L$,从而过渡到微分算子复形。

\noindent 前述对函子(19)的两种描述都还是不完整的,尤其是因为第一个成员的一个对象未必在整个 $X$ 上来自一个微分算子复形。我推测可以给这一启发性描述一个内在诠释,用公式
\begin{quote}
(20)\quad $G^{\circ}(C)\xrightarrow{\ \sim\ }\mathrm{RHom}_{\mathcal{O}_{X_{\mathrm{cris}}}}(G(C),\mathcal{O}_{X_{\mathrm{cris}}})$\quad (其中 $G$ 定义于(17))
\end{quote}
但不曾验证它正确。经标准论证,在这里人们又归结到(为了证明当 $C$ 如上所说结合到 $L^{\bullet}$ 时自然箭头(20)确实是一个同构)$C=D$ 的情形,而那时(20)归结为公式\dmark{}\origpage{1017}
\begin{quote}
(21)\quad $\mathrm{Ext}^{i}_{\mathcal{O}_{X_{\mathrm{cris}}}}(X_{\mathrm{cris}};G(D),\mathcal{O}_{X_{\mathrm{cris}}})=0$\quad 对 $i>0$,
\end{quote}
它们与(18)相当神似。

\noindent (19)的全忠实性之意,无论如何是相当清楚的,并且(如(17)那样)又经拆卸归结到 $C=D$、$C'=D[i]$(次数平移 $i$)的情形,于是归结为公式
\begin{quote}
(18.1)\quad $\Gamma(X,D)\simeq \mathrm{Hom}(P,P)$ ,\quad $\mathrm{Ext}^{i}_{\mathcal{O}_{X_{\mathrm{cris}}}}(X_{\mathrm{cris}};P,P)=0$\quad 对 $i>0$,
\end{quote}
其中记
\begin{quote}
$P=P^{\infty}(\mathcal{O}_{X})_{\mathrm{cris}}$,
\end{quote}
%JOIN:para
% ---------------- p.784 ----------------
它是 $X$ 上一个无比卓越的晶体代数。这里(就晶体 $\mathrm{Ext}^{i}$ 的零性而言)假定 $X$ 是仿射的(相应地 Stein)。

\noindent 最终,昨天在我看来还“简直美得不像真的”、而那时我还在透过Deligne的照片看事物的那个东西,一旦把事情写成不必添上完整性条件(更不必说正则性条件),就突然显出一切合情合理的模样。若上帝假我以寿命,且没有别人抢先替我做这工作,我希望赶在年底以前把那件事弄清楚(连同(21)与(18)的有效性),借《数学反思》第三卷中专门讨论德拉姆系数的那一部分。

\noindent 如我所说,是第五张照片——那张与离散系数相联系的拓扑学直觉“贴合”得最紧的——得我偏爱。若是有一天我心灰意冷地得知公式(22)是假的(而若是公式(18)如此,我还不至于那么烦恼,虽说它们技术上看着更没把握),那将说明必须回到(修过版的Deligne照片的)pro-观点——一个不怎么乐观的前景!无论如何,我毫不怀疑,除技术性调整之外,人们手里确实握着一张出色的照片,它在代数几何中尤其有效(甚至在异于特征零之域的基上),且不加任何光滑性假设。

\noindent 至于第四张照片,其忠实性取决于(18)的有效性,我再度承认,在光滑情形之外(甚至在光滑情形之中)“我还看不明白”它,且不能确定对非光滑的 $X$,我所提出的晶体诠释是否原样行得通。不过我觉得,我那些关于 D-模观点(尤其是我对这一观点的晶体诠\dmark{}\origpage{1018}释)的根深蒂固的变差困惑,正借助于引入一个与晶体概念相对偶的概念——我称之为\textbf{余晶体}(co-cristal)——而即将解决。就在不久前的昨天,那种(就闭浸入下 D-模的“变差”而言)一直存在的不适的弥散感觉,终于分娩出一个“好概念”(据我所知,还不曾真正写下什么)。它在“ind”一侧看着很合身,犹如晶体概念(为我所熟悉)在“pro”一侧那样。在光滑簇上,两个范畴(晶体与余晶体)典范等价(而这正是我必然倾向于把它们混同的原因——情有可原……),但对任意的 $X$ 就不再如此。局面与上同调环 $H^{\bullet}(X)$ 和同调群 $H_{\bullet}(X)$、Chow 环 $\mathrm{Ch}^{\bullet}(X)$ 和 Chow 群 $\mathrm{Ch}_{\bullet}(X)$、或Grothendieck环(请原谅我这处不规整……)$K^{\bullet}(X)$ 和Grothendieck群 $K_{\bullet}(X)$(再次恕罪)之间发生的事完全类似。在那里,当 $X$ 是(拓扑的、或代数的等等——视情形而定)光滑簇时,人们也曾长期混淆这两类对象。事后“可以解释”:第二项无论如何配备了一个第一项上的模结构(cap-积——在后两种情形中它由一位我不敢在此指名的祖先引入),而在光滑情形中,人们发现这个模是秩 1 自由的且配备一个典范基,这使它不幸与那个环(无疑美得多)混同。

\noindent 那么,对于 $X$ 上模的晶体的范畴 $\mathrm{Cris}^{\bullet}(X)$(以张量积配备“环”结构)与模的余晶体的范畴 $\mathrm{Cris}_{\bullet}(X)$(前者以一个 cap-积完美地在其上“作用”)之间,情形也是这样。

\noindent 但现在是停下这一长段数学离题的时候了,我承认它在一场漂亮的葬礼的编排中全然不合时宜。有兴趣知道下文(自然纷繁)的读者,只好去买《数学反思》第三卷(若他不吝惜钱财),在那里一位不知悔改的逝者打算继续他那些混乱的“技术性离题”\fnote{709}{(*)}{这一次,不言自明,是作为我另一位学生的“合作者”,而这位学生早已高升为“晶体”之父……}。

% ---------------- p.785 ----------------
\subsubsection*{(e) 好上帝的无处不在}

\tocsubsub{(e) 好上帝的无处不在}

\noindent (5月27日)一篇“最后的”脚注,昨天(在把《封神》头十二篇注交给打字之前)临急之际补进“五张照片”,结果又涨到“骇人的规模”,而我\dmark{}\origpage{1019}最终要把“这一长段数学离题”用最后一个(简短的)章节继续下去。这样一来,《五张照片》就将由(a)至(e)\textbf{五}个章节组成——可见一切都会圆全和完善……

\noindent 这是对Mebkhout“好上帝定理”(设想的)真正有效性域的一番评论,它(据我看来)远远超出复解析空间这一初始框架——不仅靠它带来的新\textbf{哲学}(而这哲学眼下已经刷新了上同调主题),而且也在一个技术意义上超出。

\noindent 一旦把 $X$(光滑)上可构的 $\mathsf{C}$-向量空间层,或者诠释为分层 pro-凝聚模(Deligne式),或者(通过在 $X$ 的开集的无穷小加厚上取投射极限)诠释为晶体层(Grothendieck式),那么“好上帝定理”(别名Mebkhout)就断言两个范畴的等价,而这两个范畴这一次\textbf{两者都是}“纯代数”性质的。换句话说,这定理现在有一个精确的意义,是在复解析语境以外的其他语境中:既在体上光滑概形的语境中(甚至无须假定特征零——关于这一点见上面第996页脚注(**);在特征 $p>0$ 中“除幂晶体”的观点在此是本质的),也在任意特征的刚性解析簇的语境中,又在 $\mathbb{Z}$ 上有限型光滑概形的语境中(此外还有许多……)。

\noindent 好上帝定理的“形式”部分关乎\textbf{所有}凝聚的 D-模复形,而不只是那些完整的,它说:好上帝函子经祖先修订校正之后(即本质上相对于结构层 $\mathcal{O}_{X}$ 的对偶)是范畴 $D_{\mathrm{coh}}(X,D_{X})=\mathrm{Cris}^{*}_{\mathrm{coh}}(X)$ 到所考虑的系数范畴 $\mathrm{Coeff}^{*}$(随感兴趣者的口味任选)的一个\textbf{全忠实}函子。把事情看透了,这应当或多或少是“连环推理”的。

\noindent 但在到达范畴中,人们“以拆卸”定义两个值得注意的全子范畴,即“完整系数”的与“完整正则系数”的(如(c)末尾以及第1011页脚注(**)中那样)。话虽如此,“Mebkhout广义定理”(就所考虑的语境而言)肯定毫无连环推理之嫌,但必定深刻,它将说两件事:
\begin{quote}
1.~\dmark{}\origpage{1020}完整“系数”的范畴 $\mathrm{Coeff}^{*}_{\mathrm{hol}}$ 在范畴 $\mathrm{Cris}^{*}_{\mathrm{coh}}(X)$ 经“Mebkhout-Grothendieck”(全忠实)函子的像之中。(注意:实质上,这个函子就是Mebkhout函子,但是着眼于 $\mathrm{Cris}^{*}_{\mathrm{coh}}(X)$ 全体看的,并且还经“祖先修订校正”,使目的落在 $\mathrm{Coeff}^{*}$ 中,而后者有纯代数的意义……)

2.~以“微局部”的“完整性”与“正则性”条件,用微分算子复形的术语,来刻画 $\mathrm{Coeff}^{*}_{\mathrm{hol}}$ 与 $\mathrm{Coeff}^{*}_{\mathrm{hol\,reg}}$ 的逆像。
\end{quote}
关于最后这一点(它对于我六十年代的纲领或许相对次要),在特征零中人们已经有一个现成的完整性条件。至于正则性条件,现在正好该看看日本人手里是否恰好有那个恰当的概念——但这可不是Mebkhout能告诉我的,因为他见得太多,不想再听人提起。

\noindent 至于我,我没有他那样的见识,我觉得正则性有\textbf{三个不同的侧面},它们相互补充:
\begin{quote}
1.~Deligne以拆卸在 $\mathrm{Coeff}^{*}_{\mathrm{hol}}$ 中廓清的“几何”侧面,归结为一个“局部系统”(例如带可积联络的丛)在一个除子邻域上的正则性条件。

2.~“微局部”或“日本式”的侧面,直接用微分算子复形的术语来表述(?)

3.~Mebkhout引入的“上同调式”侧面,这个侧面目前(据我看)只在复解析情形中被很好理解。我丝毫不知道它有没有机会推广到刚性解析情形。
\end{quote}
侧面 $3^{\circ}$ 当然将是关键的,每当需要为定义在一个完备赋值体上的代数簇、以完整系数,在“Zariskian”上同调与“刚性”上同调之间建立一个\textbf{比较定理}的时候。

\noindent 对于我六十年代宏大的“变差纲领”,最重要的恰恰是“几何”侧面。要紧的是为 $\mathrm{Coeff}^{*}_{\mathrm{hol\,reg}}$ 定义一套六种运算的形式主义。如果甚至能为 $\mathrm{Coeff}^{*}_{\mathrm{hol}}$ 找到一套,如Mebkhout似乎所信,那就更好。

\noindent 但是(如果我没弄错)Motif(我对此早有预感\dmark{}\origpage{1021}超过别的一切)只会生出同时完整且正则的系数来。

\noindent 我回到问题1,它有一个显而易见的变体,即一个“问题 $1'$”(更为谦虚),把 $\mathrm{Coeff}^{*}_{\mathrm{hol}}$ 换成 $\mathrm{Coeff}^{*}_{\mathrm{hol\,reg}}$。一旦证明了Mebkhout-Grothendieck函子的全忠实性,人们显然就归结到下面这件事:在 $X$ 的一个光滑(未必闭的)子簇 $Y$ 上,给定一个带可积联络的丛(或一个 $\mathsf{C}$-凝聚的 $F$-晶体 $C$,视所选的语境而定……),必要时给它附加一条Deligne的补充正则性条件(在 $\overline{Y}-Y$ 的各点上)。Deligne的办法(或由祖先修订以便过渡到晶体语境)使我们能给它们结合 $\mathrm{Coeff}^{*}$ 的一个对象(按定义它甚至会是“完整的”,乃至“完整正则的”)。这个对象在Mebkhout-Grothendieck函子的像中吗?或者,同样地说,在 $X$ 上局部地,所说的这个 $\mathrm{Coeff}^{*}$ 对象能用 $X$ 上的一个微分算子复形来描述吗,借祖先专利的办法,即过渡到该复形的“形式化”,把它或者诠释为Deligne式的复形,或者诠释为一个晶体复形?

\noindent 这个问题的答案,无论如何在复解析情形中、以及在特征零的体上光滑相对概形的情形中是肯定的(除非我弄错),甚至无须引入正则性条件。那就是我此前已经特意强调过的(在(c)中,第1011页)“由Mebkhout的理论带来的全然出乎意料的现象”\fnote{710}{(*)}{在当今,强调这样的事实(至少就在此所论的数学那一部分而言)已变成一桩真正的\textbf{公共卫生事业},在一个上同调主题的几乎全部出版物、而且(我怕)所有那些以今天享誉的署名发表的出版物,都被写成恰好\textbf{略去}使所有这些文本活生生的那些强有力的伟大观念、并\textbf{抹杀或根除}如此关键的工具(旧的或新出现的)、如此要害的概念、如此多产的观念的作用与出处。有一种智识上的\textbf{腐败}(一种更深\textbf{腐败}的标志……)当今在我们的科学中人人眼皮底下蔓延,对于任何别的科学、在任何别的历史时刻,我都未曾有所耳闻。}。在正则情形中(包括“在无穷远”处),这本质上就是好上帝定理。在一般情形中,如果我没弄错,这应当毫不费力地归结于我称之为“完整性的上同调判据”(或“Kashiwara可构性定理之逆”)者,它归功于Mebkhout,在下一篇注“三座路标——或天真”(n\textsuperscript{o}~171(x),见第1028页)中有所讨论。

% ---------------- p.786 ----------------
\subsubsection*{b2.~三座路标——或天真}

\tocsubsub{b2.~三座路标——或天真}

\origpage{1022}\notehead{171(x)}{}

\dmark{}\noindent (5月5日与5月23日)\fnote{711}{(*)}{本子注“三座路标”出自“作品……”一注(n\textsuperscript{o}~171(ii))的一条脚注。请见置于该注末尾附近的参照符号。}Mebkhout在1972至1980年间发展的哲学,可以归结为\textbf{三大定理},三者都紧密联系着我在五十与六十年代发展的那些观念,但其中没有一项是我(或任何人)当年预见到的\fnote{712}{(**)}{如我在“荒唐的问题”一注(n\textsuperscript{o}~171(vi))中所指出,然而我早已知道Mebkhout整体对偶定理的一个变体,是关于 $X/S$ 真光滑相对概形的,用相对微分算子复形表述。确切地说,若 $L^{\bullet}$ 与 $L'^{\bullet}$ 是这样的复形,彼此“相伴”,则 $Rf_{*}(L^{\bullet})$ 与 $Rf_{*}(L'^{\bullet})$,作为导出范畴 $D(S,\mathcal{O}_{S})$ 的对象,是“完美”复形(局部上可由次数有界的有限型自由模复形表示),并且按完美复形的通常意义互为对偶。在 $S=\mathrm{Spec}(\mathsf{C})$ 的情形,这个定理或多或少等价于Mebkhout的定理(限制在代数且真的解析簇情形),然而有一个重要区别:我当年缺少一个“导出范畴”的观点来处理微分算子复形。另一方面、而且尤其是,我丝毫没有想到过这些复形(附加以Mebkhout廓清的适宜条件)构成“离散系数”(即德拉姆系数)的完美替代物。另一方面,至少从1966年起我就清楚,必定存在这样一种代数可构 $\mathsf{C}$-向量空间系数的替代物,对任意特征的相对概形有意义,而我的晶体观念正是朝此方向的第一步。如将在[Crystals]中看到的(即前一篇注“五张照片(晶体与 D-模)”(n\textsuperscript{o}~171(ix))所引的那些讲演),我的晶体反思的内在逻辑却已然把我重新带到微分算子复形的门前。那时我已经离Mebkhout的哲学很近。需要我的同调学生们(尤其是Deligne、Berthelot、Illusie)被“葬礼”综合征堵住,才没有在此后几年间廓清这一哲学。(我本人当时全神贯注于别的基础工作,而把晶体主题交给我的学生们照料。)}

% ---------------- p.787 ----------------
\noindent 第一座大定理是Mebkhout1972至1976年间工作的主要成果。它关乎光滑复解析簇 $X$ 的结构层在解析闭子空间 $Y$ 中的支撑上的\textbf{局部上同调}层 $H^{i}_{Y}(\mathcal{O}_{X})$(此概念由Sato与我各自独立引入)。这里Mebkhout之前没有人想到去作的关键观察,是 $X$ 上无穷阶微分算子环 $D^{\infty}$ 的运算\fnote{713}{(***)}{关于这些算子的定义——其名字初看吓人,但它们引出一套在每一点上都与寻常微分算子形式主义平行的形式主义——见前一篇注“五张照片(晶体与 D-模)”(n\textsuperscript{o}~171(ix))的(b)部分。},由于它们\dmark{}\origpage{1023}作用于参数 $\mathcal{O}_{X}$,也作用于这些上同调层。另一方面,在代数几何的“Zariskian”框架中,我(在五十年代末?)曾把这些层描述为 $\mathrm{Ext}^{i}$ 层的归纳极限。这经类比引导Mebkhout引入局部上同调的一个“代数部分”,以及一个典范箭头
\begin{equation*}
H^{i}_{Y}(\mathcal{O}_{X})^{\mathrm{alg}}\overset{\mathrm{def}}{=}\varinjlim_{n}\mathrm{Ext}^{i}_{\mathcal{O}_{X}}(\mathcal{O}_{X_{n}},\mathcal{O}_{X})\longrightarrow H^{i}_{Y}(\mathcal{O}_{X})\overset{\mathrm{def}}{=}\mathrm{Ext}^{i}_{\mathsf{C}_{X}}(\mathsf{C}_{Y},\mathcal{O}_{X}),
\tag{1}
\end{equation*}
其中 $X_{n}$ 表示 $Y$ 在 $X$ 中的第 $n$ 个无穷小邻域,而 $\mathsf{C}_{X}$、$\mathsf{C}_{Y}$ 是 $X$ 相应地 $Y$ 上的常值层 $\mathsf{C}$(后者以零延拓到 $X-Y$ 上)。第二个本质观察,是这一次 $X$ 上寻常微分算子的环 $D$ 作用于第一个成员。人们早已知道,所得到的这类层,无论右边那个性质上超越的成员,还是左边那个性质上“代数”的成员,作为 $\mathcal{O}_{X}$-模都相当庞大——毫无凝聚性可言,这是肯定的。诚然,人们也曾有一种感觉(至少在代数那一侧):总还是有某类“有限性”或“余有限性”,其意义Mebkhout之前没有人想过去精确化。

\noindent Mebkhout那条卓越的定理,就是第一个成员是一个 D-模,而且此外第二个成员(它看着更加难以处理) 只是由第一个成员经环的更换 $D\longrightarrow D^{\infty}$ 得出。由于第二个环已知在第一个上平坦,这也顺带蕴含(1)是单的。\noindent 同时,鉴于凝聚性结果——关于第二个成员(Mebkhout之前没有人对它有任何理解)的这个结果可以被视为一条非常强的有限性定理——特别地,它作为一个 $D^{\infty}$-模是有限呈现的(但也许不凝聚,因为人们不知道 $D^{\infty}$ 本身是否凝聚)。

\noindent Mebkhout处理的头一种情形,即具正规交叉的除子的情形,是他1974年通过的第三阶段论文的课题。即便这种情形已非平凡,而且当然是全新的——Mebkhout所解决的那个问题本身从来没有人看见过。况且这种情形结果正是关键的情形,Mebkhout(以一般性递增的逐次近似)借奇点消解能够归结到它\fnote{714}{(*)}{关于Mebkhout的局部上同调定理,请特别见:\emph{La cohomologie locale d'une hypersurface},载 \emph{Fonctions de plusieurs variables complexes III},Lecture Notes in Mathematics n\textsuperscript{o}~670,p.~89--119,Springer-Verlag(1977),以及 \emph{Local Cohomology of analytic spaces},Publ. R.I.M.S. Kyoto Univers. 12,p.~247--256(1977)。}。

\origpage{1024}\dmark{}(我方才陈述的这个结果,单凭自身在我看来就有这样的分量:在稍稍正常一点的条件下,它本会为它的作者赢得国际声誉。同样,他所处理的头一个关键情形,已然显示了一种视觉上的独创性,“正常地”本会为他赢得他的长辈中那些有资格品鉴其妙处的人(例如我的每一位前学生,无一例外)的热烈鼓励。罢了……

\noindent 事实上,在这四年间,Mebkhout得出一个比我方才陈述的还要详备的结果。他证明他所研究的 D-模不仅凝聚,而且此外是\textbf{完整的}(这是他在日本学派中找到的概念),而且此外是\textbf{正则的}\fnote{715}{(*)}{Mebkhout对正则性的原始(超越的)定义,在“作品……”一注(n\textsuperscript{o}~171(ii))中有所回顾,见第950页脚注(*)。}(取他 ad hoc 定义的意义,其灵感来自我的代数-解析德拉姆上同调比较定理)。更进一步,他证明出发的那个可构 $\mathsf{C}$-向量空间层 $\mathsf{C}_{Y}$(它进入(1)第二个成员的定义)从 $D^{\infty}$-模复形 $\mathrm{RHom}_{D^{\infty}}(\mathsf{C}_{Y},\mathcal{O}_{X})=C$,凭那条非凡的反演公式
\begin{quote}
(2)\quad $\mathsf{C}_{Y}=\mathrm{RHom}_{D^{\infty}}(C,\mathcal{O}_{X})$。
\end{quote}
\textbf{重新构造出来}。从来没有人梦想过这样的公式——而在五年后那决定性的一天之前,也不会有人梦想它,那时这哲学的力量显露出来,同时给葬礼发出信号,在祖先身边、连同那位把它带来的人……为了梦想它,本不该已经埋葬了祖先的哲学(借导出范畴、带不带下划线的 $\mathrm{RHom}$ 以及别的“无用细节”之名);此外,还该懂得欣赏一个极平淡却充满神秘的几何局面(支撑在具正规交叉的除子中的局部上同调),并走到神秘\textbf{尽头}。这个“尽头”,它还不在我刚才描述的那座1976年的辉煌定理中——但早在那一刻,Mebkhout就对它有清晰的愿景:那就是双重的“好上帝定理”,一个用于完整正则 D-模,另一个用于完整 $D^{\infty}$-模,以及前面说过的双重反演(或“双对偶”)公式\fnote{716}{(**)}{在前一篇注“五张照片(晶体与 D-模)”(n\textsuperscript{o}~171(ix))的(b)部分。}。这同时也是离散(解析可构)系数与“连续”系数之间关系问题的一个美妙简洁的解。

\origpage{1025}\dmark{}(但我在预支。当他证明了构成他著作与哲学的第一座大路标的那个定理时,那个“尽头”,虽已被清晰洞察,在他看来仍远得令人眩晕。如果他在身边找到一位能干而仁慈的长辈,且有最低限度的经验与数学嗅觉,那人会纠正他:显然他已经离它很近,而所要克服的困难,一如发现工作中常见的那样(甚至不妨说,向来如此……),更多是心理上的,而非技术上的。但在着手追逐那无穷远之前,他先攻整体对偶定理——那个注定要“统罩”已知对偶定理的定理,无论对凝聚系数还是对离散系数。

% ---------------- p.789 ----------------
\noindent Mebkhout著作中无处不在的那种深刻动机,把局部上同调问题与整体对偶问题这两个问题联系起来,那就是对离散系数与连续系数之间一种本质统一的预感。那也正是我1966年晶体进路的指导线索,它努力以“连续”的术语把握“德拉姆系数”(性质上本质是离散的)……

\noindent 这里不是重温Mebkhout对偶定理陈述的地方\fnote{717}{(*)}{这个陈述在“作品……”一注(n\textsuperscript{o}~171(ii))中有所回顾。}。它的证明曾撞上严重的技术性困难,源于超越语境,而他以上同调下降与核 EVT(拓扑向量空间)的技术克服了它们(这些技术我本人也并不陌生,即便Mebkhout是至今唯一还坚持引用祖先的人……)。就他的对偶哲学而言,这个定理是一个本质的路标。如果同Mebkhout一起牢牢记住:把它应用到完整 D-模复形上,它除凝聚对偶之外还包含解析可构离散系数的整体对偶\fnote{718}{(**)}{在Mebkhout确立他的整体对偶定理(1976年)时,此外他还不曾证明每个解析可构的 $\mathsf{C}$-向量空间层都来自一个 D-模复形。但他对此深信不疑。},那么可以说,它本身已经孕育着Mebkhout式 D-模的全部哲学。它的分量,早在1980年(他的论文答辩后一年\fnote{719}{(***)}{见“墓外相会”一注(n\textsuperscript{o}~78)。})他第一次跟我谈起时,在我看来就是一件明摆的事。我不认为我曾有幸启发过任何一位\origpage{1026}与我接触的学生\dmark{}作出分量相当的工作\fnote{720}{(*)}{我在这里尤其想到那些跟我准备论文的学生。Deligne的情形是另一回事,因为他的论文在我离去之后进行,而且通篇不提我的名字,尽管他工作(关于Hodge--Deligne上同调)的初始灵感来自我对种种“系数”的问题意识,而后者也预见了“霍奇系数”的一套形式主义。Deligne的工作是朝此方向的第一步,远比Mebkhout所完成的那一步片面,而后者是在(与霍奇密切相关的)“德拉姆系数”方向上。诚然,Mebkhout因长辈的漠视与鄙薄而面对重重障碍,却不受那麻痹了我的学生们的葬礼综合征之苦。(关于这一点,见“……与羁绊”一注(n\textsuperscript{o}~171(viii))。)}。

\noindent Mebkhout当时还费了很大的劲才让这个定理得以发表,它浑身散发着“Grothendieck味儿”。(《数学年刊》把它退回了,让他明白这类东西够不上所需水平。结果它终究还是发表了,在 \emph{Mathematica Scandinavica} 上,1982年\fnote{721}{(**)}{整体对偶定理(用于凝聚 D-模),\emph{Mathematica Scandinavica} 50 (1982) p.~25--53。另见“Poincaré 对偶”,载巴黎第七大学奇点讨论班(Pub.~n\textsuperscript{o}~7),1977--1979,以及尤其“The Poincaré-Serre-Verdier duality”,载 \emph{Proceedings of the Conf. of Algebraic Geometry},Copenhague(1978),Lecture Notes in Mathematics n\textsuperscript{o}~732,p.~398--418,Springer Verlag(1979)。这两篇讲演的引言、尤其是第二篇的引言,代表了Mebkhout所带来的哲学的一幅草图,那是在他还是其唯一保管人与辩护人的时候。}。)我相信那正是他讲授 D-模哲学时最爱的主题,但其精神与日本人的截然不同。他对我说,这个定理确实有让听众或偶然的对话者惊叹的天赋,唯独那些属于建制圈子的人\fnote{722}{(***)}{(5月24日)这与我自己的观察若合符节。看来名人的处境容易使人养成这样一种自满,对“再美也不配她屈尊欣赏”的人。我不知道这些习性在当今整个科学界、乃至自来如此,是否通例。我起初受到一个不存在这种自满精神的圈子的接纳,曾是我极大的幸运。它想必是蹑手蹑脚地到来的,经年累月,在这些人那些人身上安家落户,渐渐无人(只有Chevalley除外……)察觉。一切看似如旧——然而一切已经不同。它已经像一层细薄的灰尘,覆盖着事物原初的新鲜。我也曾触及这灰尘,与别人一样。而今天,当我重新面对从前的学生或朋友时,我常有这种感觉:这灰尘已积成厚密的一层,形成一副密不透风的铠甲,透过他们向我发问……})例外,而且无一次例外。这是一件让我感到宽慰的事。它
%JOIN:para

表明这种餍足的自满精神(它使一切事物之美黯然失色,无论它多美)在数学界并未成为普遍。它主要(即便不说是排他地)肆虐于高层圈子,近十年来我确实有充分机会与它相识……

\dmark{}\origpage{1027}(应当用一个已提及的局部性质的结果来补全这个整体对偶定理,它同样深刻,说的是:上同调层凝聚的 D-模复形的自然对偶化函子(它把完整复形变成完整复形,完整正则复形亦然),在这些复形上此外还与德拉姆函子 $\mathrm{DR}$(“相结合的微分算子复形”,被看作上同调可构的 $\mathsf{C}$-向量空间层复形)相容,相容于我当年在这些复形上引入的那个自然对偶化函子\fnote{723}{(*)}{这就是那个对偶,经我的学生与旧友们的一致同意,它在此期间成了“Verdier对偶”(无论在复解析情形还是平展情形)……(关于这一点,例如见“好的参考文献”一注,n\textsuperscript{o}~82。)}。这种相容性显然是Mebkhout对偶形式主义的一个本质成分,为的是理解他的整体对偶定理的意义。出于一种我不明的原因,他把它称作“局部对偶定理”\fnote{724}{(**)}{这个结果就是以这个名称见于Mebkhout论文第三章的。他对我说,他这个名称(以及“双对偶定理”之名)的灵感来自我引入的术语——然而在我看来,“局部对偶定理”只是我所廓清的“双对偶定理”的另一个名字,它代表着后者的一个重要侧面,即“几何”侧面。这个相容性结果(Mebkhout向我解释)是他证明他在同一章中称作“双对偶定理”者的重要一步。(关于后者,见前一篇注“五张照片”的(b)部分。)证明的问题且搁在一边,从一种“哲学”或一套“瑜伽”的观点看,好上帝函子应当与对偶化函子交换,这当然是一件“显然”的事(因为既有好上帝!)。一则滑稽的细节:Kashiwara(Mebkhout1978年1月曾有机会与他面谈)\textbf{不相信}这个定理为真!这足以说明他多么不在行,而他恰缺少“六种运算”风格的几何视觉。这并不妨碍他后来在Mebkhout把自己的第三章寄给他(1978年2月)之后,把这结果据为己有(当然不提作者),收进他与Kawai那篇已引的大文章(见第1005页脚注(*))(loc.~cit.~第4节命题1.4.6)。正是在这部著作中,“双对偶定理”(loc.~cit.~第4节1.4--9)也被不经任何手续地据为己有(以“重构定理”之名)。这足以说明,生于巴黎的“新风格”诸位大师的太平洋对岸信徒们(他们出现在一个“Grothendieck学派”蒸发了、未留痕迹的地方),比起他们的法国同行来毫不逊色。我的(离散系数)双对偶定理也同样见于Kashiwara--Kawai这同一著作的同一不竭的第4节(命题1.4.2),然而当人们恬不知耻、不假思索地劫掠那位身后无名的学生(老板们公然弃之不顾)时,人们却按规矩向对面那位显赫的同行脱帽致意,恰如其分地引用Verdier提供的“好的参考文献”(而他本人也在劫掠一位从不被点名的逝者……)。这些骗局在消息灵通人士中固然臭名昭著,Mebkhout在这方面也有过若干反响。但显而易见,只要事关消灭那位不可引用的祖先及其倒霉的后继者,它们就被视为得体而受欢迎。}。这个深刻的定理,连同那著名的“对应”(人们肯赏脸称它“Riemann--Hilbert”时),被“所有人”(Verdier与Deligne打头)当作一件“众所周知”、不言自明的事,而且尤其从不提名某一位无名者(而“所有人”都清楚,尤其不要引用他)……

\dmark{}\origpage{1028}我终于来到Mebkhout著作中的第三座大路标。从技术上可以说它由三个(或至少两个)不同的定理构成,但它们联系得如此紧密,以致在Mebkhout的心中它们显得不可分离。早在1978年1月,他就证明了“$D^{\infty}$-模”那一侧面:即“相结合的德拉姆复形”函子限制到完整 $D^{\infty}$-模复形上的限制 $m_{\infty}$(即“Mebkhout函子”)是一个范畴等价(与上同调可构的 $\mathsf{C}$-向量空间层复形之间)。已知这个函子与对偶化函子交换,把这个定理改经相联系的反变函子 $\delta_{\infty}$ 来表述便是自然的,它由
\begin{equation*}
C\longmapsto \mathrm{RHom}_{D}(C,\mathcal{O}_{X})
\tag{3}
\end{equation*}
给出,
%JOIN:para
% ---------------- p.791 ----------------
而断言这个函子是一个(反)等价也是一回事。这个定理于是可以由Mebkhout那条绝妙的\textbf{反演公式}(或“重构”公式,或“双对偶”公式)来精确化,它把拟逆函子表达为
\begin{equation*}
F\longmapsto \mathrm{RHom}_{\mathsf{C}_{X}}(F,\mathcal{O}_{X})
\tag{4}
\end{equation*}
与此同时,Mebkhout还证明了Kashiwara可构性定理的一个\textbf{逆定理},即:若一个上同调凝聚的 $D^{\infty}$-模(或 D-模)复形使得相结合的德拉姆复形(作为 $\mathsf{C}$-向量空间层复形)上同调可构,则它是完整的(\textbf{完整性的上同调判据})。在 $D^{\infty}$-模复形的情形\origpage{1029}(其中不存在正则性问题),这因此蕴含:在那个\dmark{}导出范畴中(1978年直至1981年早就没人在里面工作了……),该复形(或不如说它的对偶)凭反演公式“重构”,除唯一同构外唯一。

\noindent 如我在别处所解释\fnote{725}{(*)}{见“作品……”一注(n\textsuperscript{o}~171(ii))第952页脚注(*)。},从那时起,Mebkhout手里就拥有证明好上帝定理(对 D-模亦然)所需的一切:即德拉姆函子限制到完整正则 D-模复形上的函子 $m$ 是一个范畴等价。这个结果给他的鼓舞较小,因为显然没有配套的反演公式\fnote{726}{(**)}{前面已经看到毕竟有一条——而我稍后还要回到这一点。}。无论如何,即便他那条绝妙的反演公式,对谁都不冷不热——从他论文的准导师Verdier开始(而Verdier后来还赏脸担任答辩委员会主席)。为了重新作技术努力去证明一件他反正确信、且感觉自己已握有全部证明手段的事,这可不是一种很令人鼓舞的氛围。他要到由一条据说高不可攀的猜想(这回不是Weil的,而是Kazhdan--Lusztig的那条)的证明所触发的“冲刺”启动之后,才去管它。

\noindent 这恰好像有意安排,正是人们突然急需的另一侧面。无论如何,当时“所有人”都那么急着使用这件崭新出现在市场上的“崭新工具”,而且所有人都心照不宣:尤其不要提出证明的问题——万一显出那工作早已由一个不可引用的人做完了呢。似乎除了当事人自己,没有人想到去把已经写就的 $D^{\infty}$-理论的各块抄录拼合起来,以证明 D-理论中该证明的那条定理。迄今唯一公开发表的证明\fnote{727}{(***)}{参考文献:\emph{Une autre équivalence de catégories},Compositio Mathematicae 51 (1984), 63--88。}看来确实是Mebkhout的那一篇,去年发表(而于1981年6月收到,正是那场难忘的反常层会议的同一个月……)。

\noindent 我在前一篇注(b)部分中解释过一个简单的原则,受Deligne通往德拉姆系数的进路启发,用以在(完整正则的)D-模框架中收复一条“反演公式”(或\origpage{1030}借Mebkhout的说法\dmark{}说,“双对偶”公式)。自从世界各地的人们开始就 D-模这门新的“时髦点心”办讨论班以来,我不知道这种十分自然的进路是否已被廓清——Mebkhout无论如何不曾有所耳闻。可以肯定的是,如果Deligne有那种“我那个时代”被视为理所当然的本能反应,那么正是他本人,在1979年6月得知一位无名者的美妙观念之后,就会鼓励他把其关键结果的 D-模那一侧面(更接近代数)的证明也写出来,并会向他建议他那条美妙反演公式的“pro”变体——说到底相当显然。同样,从那时起,对已经为此付过学费的Deligne来说,Mebkhout的观念将

% ---------------- p.792 ----------------
带来所缺的德拉姆系数(至少在特征零的体上代数几何中);理所当然的事,是鼓励他作出必需的调整,以陈述一个(不如说这一回是Mebkhout的)好上帝定理,用于复代数簇\fnote{728}{(*)}{如我已曾有机会指出的,在代数框架中,当需要以代数方式复述代数可构的离散系数时,除了局部的完整性与正则性条件之外,还应给所考虑的 D-模复形加上一条“Deligne-Mebkhout”的在无穷远处的正则性条件。}。

\noindent 但此一时彼一时。将不说一场代数簇上同调的新出发,是由一位无名的泛泛之辈孤独而顽强的努力完成的,他自承师法一位逝者,而风雅世界中好久以来没有人敢贸然说出他的名字\fnote{729}{(**)}{诚然,人们还没有找到办法为 EGA 与 SGA 找到替代的参考文献。但这些天降的缩略语中不含任何对一个必须保持沉默的名字的暗示。众所周知,SGA 这个缩略语指一个由布瓦玛丽主持、在一批十分体面且完全可具名的数学家推动下举办的代数几何讨论班,诸如 M.~Artin、J.L.~Verdier、P.~Deligne、L.~Illusie、P.~Berthelot、N.~Katz、P.~Jouanolou,甚至另一些较不出名但同样可引的人物。显然那里曾有过一个繁荣的代数几何“学派”,号称“布瓦玛丽学派”,其核心与灵魂是所引姓名中最光辉的那一位。关于这个“布瓦玛丽学派”以及作为其表现的 SGA 缩略语的更多细节,请特别见“驱逐(2)”与“殡仪馆——为科学服务”两注(n\textsuperscript{o}~169$_1$ 与 175)。(另见“双重含义——或骗术”一注(n\textsuperscript{o}~169$_7$)第899页第3段。)}。将不说复兴恰是通过十年来\origpage{1031}逝者的继承人们一边瓜分掠获物、一边埋葬的那种数学。天真的Mebkhout,若他想“活下来”\dmark{}并且“出人头地”,只须走“新风格”那条现成的路\fnote{730}{(*)}{关于这种风格(它取代了一个消失得无影无踪的“Grothendieck学派”……),见“恭贺——或新风格”一注(n\textsuperscript{o}~169$_9$)的结尾。},如其他杰出的年轻人(甚至有不太年轻的)急忙所做的那样。引用自己思想的(不可引用的)源头,这是怎样的怪癖,而把鱼搅浑、只引那些\textbf{该}引的人是多么简单。Mebkhout,我相信你的账到头了!

\noindent 你落进了一个不是为你而设的世界——而我仍然为此为你高兴,你不是为\textbf{那个世界}而造的。你做了你感觉自己该做的工作,不赶时髦,不盘算回报,只是信任你自己的本能——哪怕因此孤独地走你的路。你做了\textbf{你的}工作,而不是窥伺那些决定何为善、何为体面、何为不善的人发出的隐晦(和不那么隐晦的)信号。你没有为讨人喜欢而见风使舵,你没有看到黑却说“白”,或反过来——你是用\textbf{你的}眼睛看的。我不必为此向你道贺——你不曾寻求祝贺,不论是我们的还是任何人的。而为此我为你、为所有人感到高兴。

\subsubsection*{b3.~师父的角色(2)——或掘墓人们}

\tocsubsub{b3.~师父的角色(2)——或掘墓人们}

\notehead{171(xi)}{}

\noindent (5月5日)\fnote{731}{(**)}{本子注出自“作品……”一注(n\textsuperscript{o}~171(ii))的一条脚注。关于指向本子注的参照,见所引注末尾附近(第956页)。}这里自然的问题,当然是在代数几何中是否存在一套“六种运算”的形式主义,用于未必属于 DRM 类型的 D-模(或“晶体”),它将“统罩”我在凝聚与离散两种情形下引入的那两套——先假定,为明确思想,我们处在体 $\mathsf{C}$ 上。第一个困难源于 D-凝聚性概念对晶体的自然张量积概念、或类似的拉回像运算都不稳定\fnote{732}{(***)}{(5月22日)Mebkhout向我指出,他证明了完整性与正则性条件对(在 $\mathcal{O}_{X}$ 上的)总张量积运算以及对拉回像概念是稳定的,并且好上帝\textbf{反变}函子 $\delta$ 在其上交换。(相反,好上帝协变函子 $m$ 在其上并不交换,它把寻常拉回像变成非凡拉回像。)利用这个结果可以证明,对de~Rham--Mebkhout系数不存在一套“六种运算”的形式主义,能“延扩”张量积与拉回像这两个已知的基本运算。尤其,范畴 $\mathrm{DRM}^{b}(X)$ 不容许“内 $\mathrm{Hom}$”运算(扮演 $\mathrm{RHom}$ 的角色),且对 $f : X\to Y$,函子 $f^{*}$ 一般不容许右伴随 $Rf_{*}$。Mebkhout目下引入的函子 $Rf_{!}$(对光滑的 $X$、$Y$ 且 $f$ 固有)是 $f^{*}$ 的一个\textbf{左}伴随。(注意:de~Rham--Mebkhout系数上的 $Rf_{!}$ 运算的定义方式,使好上帝\textbf{协变}函子在其上交换,$Rf_{*}$ 亦然——无论对错……)因此可见,就用de~Rham--Mebkhout语境中现有的“自然”运算而论,它们本身\textbf{不}构成一套“六种运算”理论,而构成一种对偶理论。于是问题就在于看这种理论在多大程度上能延扩到不再假定完整且正则的(姑且说拟凝聚的)D-模上(例如仅仅完整——这个条件经张量积与拉回像保持)。似乎特别是,整体对偶公式可以为上同调凝聚(乃至仅仅拟凝聚)的 D-模复形、以及特征零的体 $K$ 上有限型分离概形的任意态射 $f : X\to Y$ 写出,以便同时统罩凝聚对偶定理与离散对偶定理,至少取如下形式:对偶化函子“交换”函子 $Rf_{*}$ 与 $Rf_{!}$。}。\origpage{1032}因此,为了希望有一套六种运算的形式主义,就必须\dmark{}用一个比 $\mathrm{Cris}^{*}_{\mathrm{coh}}(X)$ 更大的范畴来工作,

% ---------------- p.793 ----------------
也许是“拟凝聚”晶体的范畴(取一个显而易见的意义)——但那样就几乎无望收复一条双对偶定理!此外,经 $\mathcal{O}_{X}\to D_{X}$ 的自然标量扩张函子显然不与张量积交换——因此,即便对晶体有一套六种运算理论,它延扩那套(多亏Mebkhout,目下在实质上已为人知)de~Rham--Mebkhout晶体理论(借“结构搬运”从“离散”理论经好上帝函子得到),它也不会延扩凝聚 $\mathcal{O}_{X}$-模的那一套\fnote{733}{(*)}{这个陈述更应借一套“六种运算的对偶理论”来重新表述,见前一条脚注。}。这也许并不排除:仍然可能存在一条“整体对偶定理”(拟凝聚晶体版本),用于特征零的体上有限型概形的一个(姑且说)真态射,它(在一个显而易见的意义下)“统罩”de~Rham--Mebkhout晶体的“已知”对偶定理(实质上,仍是借结构搬运),以及凝聚情形中已知(不带引号)的类似对偶定理\fnote{734}{(**)}{可以设想这样一个对偶定理有三种不同形式。或说上与下的对偶化函子“交换”函子 $Rf_{!}$ 与 $Rf_{*}$,或说两个恰当定义的函子 $Rf_{!}$ 与 $Rf^{!}$ 互为伴随,或写一条“投影公式”(它统罩两种陈述): $Rf_{*}(\mathrm{RHom}(F,Rf^{!}(G)))\approx \mathrm{RHom}(Rf_{!}(F),G)$。}。

\dmark{}\origpage{1033}(令我相当震惊的是,Mebkhout本人连这最后一个问题都没有提出过,早在他达到自己的“绝对”对偶定理(对应于靶簇退化为一点的情形)的表述那一刻——就在最近他还显得不太“感受”到它\fnote{735}{(***)}{(6月8日)然而Mebkhout向我保证,他早就切实提出过这个问题。如果说我有过相反的印象,那必定是因为这个问题对他来说始终完全停留在柏拉图式的层面上。}。这使我深切地感到,某一种“哲学”——它在六十年代上半叶对我已成为第二天性,而且(我曾觉得……)对我的学生们也是——已被所有人忘却到何种程度,首先那些自命为其掘墓人而非传承者的人。而我看到,这也正是我离去后一门理论(概形的上同调)所经历的那场惊人停滞的主要原因,而我当年把它留在全盛之中。

\noindent 应该说,Mebkhout置身于复解析的超越语境中,而非概形语境。这就引入了可观的技术困难,在某种意义上说是“寄生的”,当需要达到对本质的变差现象的理解时更是如此。在这方面,他的长辈们又一次辜负了他们的任务,那就是把他们在我接触下获得的经验提供给新来者(正如我当年曾把我的经验提供给他们一样……),由此引导他(或至少给他一些启发),尤其是在选择投身的方向时。

% ---------------- p.794 ----------------
\noindent 但照亮与引导,这也是一种\textbf{服务},而他们早已毫无保留地选择了师父的角色。

\subsubsection*{b4.~死去的书页}

\tocsubsub{b4.~死去的书页}

\notehead{171(xii)}{}

\noindent (5月5日)\fnote{736}{(*)}{(5月22日)本子注与上一篇一样,出自“作品……”一注(n\textsuperscript{o}~171(ii))的一条脚注。见该注末尾的参照符号,第957页。}Mebkhout对我说,在我两年前我们相遇时跟他谈起之前\fnote{737}{(**)}{这次相遇在“墓外相会”一注(n\textsuperscript{o}~78)中有所讨论。},他还从未听人说出过“六种运算”这几个字——他很纳闷我想说的是哪些“运算”!显然,他从来没有想到(除我之外,似乎也没有任何别人想到)去回顾一套十分简单的上同调形式主义的主要成分,发现其中有六个基本函子或双函子,分成三对伴随函子,带着这样的箭头与相容性等等。这些事在我看来是如此显然,\origpage{1034}以致我原以为任何读者——无论\dmark{}是读《留数与对偶》阐述凝聚对偶要素的,还是读 SGA~4 或 SGA~5 阐述离散对偶要素的(况且基本用同一套公式)——都会(像我早在五十年代所做的那样,虽则我承认没有做到底)为自己立起一套多少系统、多少完备的主要同构与主要相容性的公式,因为只有这样、别无他法,才能领悟一门新语言的精神,把它同化为己有,使它成为“自己的”。毫无疑问,微积分的先驱们正是这样、也正是别样地做的,为的是在对无穷小获得一种精致而可靠的直觉,而那时他们尚缺概念工具去按后来(或重新)出现的严格性准则把握它们……

\noindent 隔了二十年回头来看,我意识到在所引的那些“参考文本”中——它们都以极大的细心乃至才华写成,而一切“真正的工作”(按时下的要求)都已做出,其顶峰是“那条”主对偶公式,即 $Rf_{!}$ 与 $Rf^{!}$ 之间的伴随公式(几乎是唯一被认为值得注意与努力的一条,哪怕第二天就忘掉它,如同没见过森林就忘了树……)——然而在所有这些文本中,\textbf{主要的}东西没有说出口,也没有从作者传给读者(假定它被作者本人很好地看到与感受到)。“主要的东西”是一套“瑜伽”,一种“哲学”,是一条穿越(在此情形下)代数几何中上同调丛林(以及别处)的万试万灵的指导线索。人们可以在“一切都已做完”(据说)之后,用五十页或一百页把它大加发挥;也可以只用几页把它唤起,而把发挥的工作留给读者,由他按自己的需要或自己的满足,以为己用的程度去发展它。

\noindent 正是那么几页——无论关于“六种运算”、关于Motif,还是关于许多别的东西\fnote{738}{(*)}{这些行写就之后,我曾能够确认:关于六种运算,我在此弄错了——事实上我是被 SGA~5 的屠杀版欺骗了,Illusie细心地把我在口头讨论班中大加发挥的“六种运算瑜伽”的一切痕迹根除了,连同一套被充分注释的完整公式。}——那些我强烈感受到、却不知道把它们写出来何其重要的书页——恰恰是它们,首先在我的文字著作中缺失了。我当时全神贯注于\origpage{1035}无休止的细心任务,\dmark{}为大家服务,做那沉重的“计件工作”,而那是唯一被认为该出版的东西——我不曾感觉到还有更本质的书页,而写那些书页我是\textbf{唯一}能做的人。\textbf{我要说的本质}东西没有进入写就的书页,只是口耳相传而已
%JOIN:para
% ---------------- p.795 ----------------
——在它传得过去的时候!至多,它也许是在无尽的奠基卷帙的字里行间——可当今有谁懂得读字里行间呢?

\noindent 因此,本质的东西,就是日复一日托付给那些我作为数学家的生活中堪称“亲近”的人、首先是我的学生们的东西。那是一件理所当然的事,毫无刻意。我不曾想到,以某种方式我也是在这样授予他们一份可观的\textbf{权力}。并不是我不曾感受到我所构思与传递之物的力量,但那种力量,它也是理所当然的。对我而言,肯定至少在数学中,“力量”与“美”过去是、现在仍然是同一回事。我不曾想到人们会滥用这些东西,这些对我充满平静而炽烈的生命、为生存与生育而造的东西。当我离去时,以着实再突然不过的方式,我对它们没有一丝忧虑。这些我从未想过要写的书页——我心中毫无疑问,它们的讯息早已被接纳并铭记,而这些“亲近”者将成为许多活生生的书页,说出那讯息,并以他们所能献上的最好东西丰富它。

\noindent 那些我曾以信任与尊重相待、视作弟弟并在他们身上认出我自己的人,选择了埋没与沉默。而当那一位到来时——忠于他自己,而他们在其中认出了我——这些饱有一切的人,却选择把他留在他们紧闭的门前——一个异乡人与一个闯入者。我不认识你!而这些未写的书页,这些白说的话,变成了死去的书页,死在这些门第高傲而紧闭的殷实人家中,而被拒认的弟弟只好尽力在自己心中把它们重新找回,经漫长而摸索的劳作。他独自不得不在那以千百卷计的无从解脱的丛林中开辟道路。经历过这一切的人,即便他曾有幸(如我当年)得到经验丰富而仁慈的向导们兄弟般的扶助,也深知我所说的是什么……

\noindent 他在日子与岁月中艰难地走出了自己的路——一条坎坷的路,事后看来有时没有罗盘,或至少没有别的罗盘,只有一种还在寻找自己的嗅觉,穿过\origpage{1036}一段\dmark{}经痛苦与艰辛获得的经验。他没有为自己的用途重写那些现成的书页、那些罗盘书页(它们在高门大户中变成了死去的书页)——除了零星碎片。他写了\textbf{另一些}书页,\textbf{他的}书页,痛苦地属于他自己。他坎坷地、顽固地写下它们,在所有人的漠然中。然而这些往往笨拙、配上个莽汉也相称的书页,我那些精明殷实的前学生(若他们肯屈尊读一读)定会以怜悯视之而一无所见——它们是\textbf{本该}写的书页,是那些我连想都不曾想过要写的书页的自然的、“显然的”续篇,因为它们在我看来太理所当然了……

\subsubsection*{c1.~一个愿景的萌发,或闯入者}

\tocsubsub{c1.~一个愿景的萌发,或闯入者}

\notehead{171$_{1}$}{}

\noindent (4月15日)\fnote{739}{(*)}{(5月30日)接下来这三篇注(n\textsuperscript{o}~171$_{1}$ 至 171$_{3}$)写于4月15日至18日(1985年),那时《封神》还只有十页光景。由于Zoghman~Mebkhout的来访引发的对“四种运算”反思的重新启动,这些注在5月份大为扩充。十页变成了一百多页,因而其中绝大部分比接下来这三篇注要晚。由此有若干局部的重复,某些事实或插曲在前面的注与接下来的注中都有提及或描述,只是角度不同。为保持文字的自发,我不愿作调整以消除这些重复。}我想借我的共同下葬者Zoghman~Mebkhout本人最近一次在我处的做客,提供一些关于他的奇特遭遇的最新鲜细节,如他在我们谈话的当口这儿那儿零星吝惜地告诉我的。

\noindent Zoghman曾有幸与他的“老板”\fnote{740}{(**)}{(5月24日)Mebkhout坚持认为“老板”一词(即便带引号)在此用得不妥。自1972年入门以来直到今天,他做工作都没有老板,只靠自己想办法。Verdier只不过是他论文答辩委员会的主席。除此之外,他的作用仅限于把“好的参考文献”传给Mebkhout,而这份文献大有用处,那时 SGA~5 仍在我的同调学生们联手扣押之下(而这正是为了诸如“好的参考文献”那种操作的需要……)。}J.L.~Verdier有三次“会谈”。第一次发生在1975年——他当时需要一条技术性的结果,而这条结果(事后看来)就包含在关于解析可构离散系数的双对偶定理之中——那时Zoghman连可构性这个概念都还不知道。(这是我很早——五十年代——就引入的一个概念,后来在平展拓扑的框架中又得到重提,见SGA~4。)在那个时刻,\origpage{1037}这个概念在分析中还根本谈不上“尽人皆知”,不像\dmark{}今天这样。事情恰恰是:这正是他的工作所需要的那个概念。Houzel(他与Verdier同期听过SGA~5,只是把我当时在那里讲的东西多少忘了些)建议他去找Verdier。这便是与那位大人物的第一次“会谈”。Verdier于是让他明白:他所要的东西(即两个具有同构“对偶”的离散复形必定同构)在某些技术性条件(恰恰就是“可构性”)之下成立,这些条件他会在Verdier正要交给他的那份手稿里找到现成的阐述。那份手稿正是“好的参考文献”\fnote{741}{(*)}{即 J.L.~Verdier的文章《结合于一个闭链的同调类》,Astérisque n\textsuperscript{o}~36 (SMF),第~101--151~页(1976年)。在相继的两条注“好的参考文献”与“玩笑——或‘权复形’”(n\textsuperscript{o}~82、83)中有详尽的谈及,在“攀登的若干插曲”一注(n\textsuperscript{o}~169~(iii))的第3插曲中也有较简略的谈及。},在那里面(在性质相同的种种壮举之间)他还装出发明了可构层、发现了双对偶定理(连同它的证明)的样子——而这些东西,他早在十二年之前(1963年)就从我口中得知了\fnote{742}{(**)}{早在五十年代下半期,我就对离散层的种种“可构性”概念(代数意义的、复解析意义的、实解析意义的、分片线性的——温和拓扑的语境则尚待到来……)发生了兴趣,与凝聚性概念并列,把它们当作在层的框架中表达有限性条件的自然概念;我还提出了这些概念在“六种运算”下的稳定性问题。平展上同调此后的进展(1963年及其后)促使我回到平展框架中的这些问题,并发展出使之能以统一方法处理的技术(拆卸与消解),它们同样适用于复代数簇与复解析簇的超越语境。双对偶定理——在平展框架中(以纯粹性与消解为前提)与在超越语境中同样成立(且证明相同)——我早在1963年就已提炼出来。它而且见于SGA~5的最早一讲(1965年),并在1977年Illusie版的大屠杀中幸存了下来。}。关于我的这个人,他在这个话题上一个字也没有提,无论在那次会谈里,还是在次年刊行的那份手稿里。反正Zoghman是心满意足地离去的,对那位大人物满怀感激——正是大人物,在那一刻、以及此后几年里(可构性概念将在他所有的工作中扮演关键角色),恰好供给了他所需的东西。

% ---------------- p.796 ----------------

\noindent 到了1976年初,他开始对对偶发生兴趣,并为我曾在凝聚情形与“平展”离散情形中发展、而后由Verdier在拓扑离散情形中接手的那两套对偶形式主义之间的相似所吸引。那正是这套形式主义多年以来\origpage{1038}已陷于废弃、而我的学生们对导出范畴实行着一场\dmark{}默契而严格的联合抵制的时刻——而导出范畴正是它的自然语言。“六种运算的形式主义”这一概念连同这个名字,自五十年代以来、贯穿整个六十年代,一直是我最重要的主导观念之一,却在我离去之后即刻变得(并且直到今天依然)是严格意义上的禁忌。(Zoghman两年前来见我的时候\fnote{743}{(*)}{这次来访在“墓外相会”一注(n\textsuperscript{o}~78)中有所谈及。关于针对“六种运算”实行的联合抵制,另见“死去的书页”一注(n\textsuperscript{o}~171(xii))。},还不曾听人说出过“六种运算”这几个字,起初也不知道我说的“运算”是些什么——而我还以为这是二十年来人人都熟悉的概念呢!)由此可见,投身这一方向的处境是何等不利——在那里,他注定要在全然的孤独中工作。但这并没有妨碍他早在1976年就提炼出一条对偶定理——关于非奇异复簇的——一举“统罩”Serre对偶定理与离散对偶定理(他称之为“Poincaré--Verdier对偶”),用的则是关于D-模复形的一条对偶陈述(其中还包含关于微分算子复形的一条整体对偶陈述)。他取的“系数”,其一般性也远远超出了Serre的情形(限于局部自由层)与Poincaré的情形(限于局部常值的离散层)——忠实于我以那套当时备受排斥的“六种运算”形式主义把这些主题引入时的精神。

% ---------------- p.797 ----------------
\noindent 两年前当Zoghman向我讲解这一定理时,我既感到它的趣味(这对我是一望而知的),也感到它的局限;因为以“六种运算”的精神来看,对我同样一望而知的是:“对的”陈述应当是关于解析空间的态射 $f : X\to Y$ 的一个陈述,取(比如说)两个函子 $Rf_{!}$ 与 $Rf^{!}$ 之间伴随关系的形式。诚然,置身超越语境带来了相当可观的额外困难,它们(在我看来)有力地遮住了Mebkhout的眼睛,使他看不清对偶之中那些本质代数机制的单纯性——而他身边竟无一人——尤其是昔日我的学生们中间——能够(或乐于……)让他体会到这一点。尽管如此,他毕竟把手指按在了一个重要的“原理”上——那就是:D-模理论(我自己更愿意称之为“晶体模”\fnote{744}{(**)}{这个“晶体”术语的缘由(显而易见)——它反映了一种更内在的D-模视野(我的学生们曾从我这里学到,却早已忘却)——见“我的孤儿”一注(n\textsuperscript{o}~46)中的评注(尤其第~179~页),以及子注~n\textsuperscript{o}~46$_{4}$(第~188~页)(x)。至于Mebkhout的哲学与我六十年代末提炼的晶体瑜伽之间那些显而易见的联系所遭遇的“健康官能的堵塞”,见“装神弄鬼”一注(n\textsuperscript{o}~85$'$,第~350--351~页)。\par(x)(5月24日)另见“五张照片(晶体与D-模)”一注(n\textsuperscript{o}~171(ix))。})提供着一个\origpage{1039}“公分母”,足以“统罩”起\dmark{}那些现象(尤其是对偶现象)——它们出现于离散上同调与凝聚上同调之中。沿这条轨道,倘若当时有一位“身在局中”的人给予鼓励,再加上最低限度的数学本能\fnote{745}{(*)}{这并不是说我的上同调前学生们缺乏“最低限度的数学本能”——否则他们中间当初谁也不可能随我做出那份好工作。但这种本能被埋葬师父的综合征带偏了道,或者堵死了。}与善意,那么毫无疑问,他会在随后三四年里,为(至少)特征零的代数几何发展出一套完整的六种运算形式主义,给出那同一形式主义(诚然备受排斥)在超越语境中那种形态的一个纯代数的、忠实的“范式”,用于代数可构的$\mathsf{C}$-向量空间层。

\noindent 他深深感到自己刚刚发现了某种重要的东西,于是满心欢喜地请求并获准与他的恩人作一次会谈,向他陈述自己的结果。这恰恰是\textbf{对}我十年或十二年前向Verdier提出、而他似乎毫无所动\fnote{746}{(**)}{(6月5日)关于这一点,见“祖先”一注(n\textsuperscript{o}~171(i)),尤其第~946~页的脚注(*)。}的那个问题的\textbf{回答}——他甚至很可能已把它忘得一干二净。无论如何,他对这个从天而降、正在做着他——Verdier——早已一笔勾销的事情的年轻人的善意,已经耗尽了。他连听也不想听Zoghman对定理的来龙去脉及其证明所作的解说。他(礼貌地)让他领会的要旨是:他Verdier早已不再相信圣诞老人,这个年轻人最好还是把东西收起来吧。

\noindent 奇怪得很,Zoghman周围竟\textbf{没有一个人}“接得住”这一结果\fnote{747}{(***)}{(6月3日)这里有过一场误会。如“三座路标——或天真”一注(n\textsuperscript{o}~171(x),第~1026~页)所说,这定理常常能让一位偶然的交谈者惊叹不已。但到目前为止,这惊叹似乎始终停留在柏拉图式的层面——定理并没有变成一件工具,一种人们熟知、随手使用乃至不加思索的东西。这肯定与下面的事实有关:为这结果显而易见的美而欣喜的人,从来不是那些“定调子”的人,不是那些决定什么“重要”、什么算“瞎拼凑”的人。(而在眼下这年头,“昨天”的瞎拼凑摇身变成“今天”的奶油甜点,可不是什么稀罕事……)——在他4月22日的评注里,Zoghman给我写道:“……面对这一定理,人们心里有一种别扭。有人暗地里嫉妒它。但鼓励过它的人实在寥寥,恰恰相反。”}——\origpage{1040}无疑这太像六十年代的“Grothendieck做派”了\dmark{},谢天谢地,如今人们早已超越了那一套!也许,两年前,我正是他所遇到的头一个体会到这一结果的重要性、以及它所孕育的那门新“哲学”的人——那就是在形形色色簇的上同调(先从代数的与解析的做起)之中,“离散”侧面与“微分”(或“解析”)侧面之间一场宏大综合的哲学。这个构成他论文其中一章的定理,最后于1982年刊于 Mathématica Scandinavica(第~50~卷,第~25--43~页)。同一篇文章曾投给 Annals of Mathematics,后者让这个年轻的妄自尊大者明白:他不够在这么一份有身份的刊物上发表的档次。

% ---------------- p.798 ----------------
\noindent 甚至直到今天,这定理在上流社会里通常仍遭无视或蔑视,尽管它已然孕育着那门新哲学——经由好上帝定理(别名Mebkhout),它提供了代数簇上同调一场壮观更新的手段。可是“人人”——包括我的上同调前学生们(尽管有一天我曾认为他们天赋一种健全的数学本能)——都一窝蜂扑向了那道新的“奶油甜点”,即某一件强大的工具(而“人人”却又装模作样,只肯以影射或迂回说法提名道姓\origpage{1041},比如“可构层与完整微分系统之间的关系”,或者“那本该正常地在这一套注里占有一席之地的东西”\fnote{748}{(*)}{这是凭记忆引自Beilinson--Bernstein-Deligne那篇“值得纪念的文章”(由Deligne执笔)——关于它,见“光荣之日”一注(n\textsuperscript{o}~171(iv))。关于这个不愧流传后世(作为一份提醒、也作为一纸警告……)的迂回说法及其来龙去脉,见“牵线人”一注(n\textsuperscript{o}~75$'$)。前面那处引文(“可构层与完整微分系统之间的关系”)摘自Beilinson--Bernstein(同在1981年)的一篇文章,下一子注(“黑手党”,n\textsuperscript{o}~171$_{2}$)将谈到它;在那里,读者还有幸结识Brylinski-Kashiwara为这门风格的开花结果所作的贡献——它服务于同一桩骗局。}……),扑向那“最新潮的玩意儿”(相交上同调),而\dmark{}那幅\textbf{视野}——它使这件工具有所由出——却依旧像先前一样无人理会,而这两者的父亲则被当作听差使唤。

\noindent 这里的情形,与我关于topos、导出范畴、六种运算、上同调系数、以及更进一步关于Motif的那幅宏大统摄视野的情形如出一辙。平展上同调与晶体上同调这类工具,正是从这幅视野中诞生的;而这幅视野本身——今天“人人”转动曲柄般使用的那些工具正是从中诞生的——在我离去那一天还强旺地活着,次日便被埋葬了。而我清楚地看到:我离去时正值蓬勃的这个辉煌主题\fnote{749}{(*)}{关于这种“阴郁的停滞”的印象,我第一次谈起是在“对一份遗产的拒绝——或一个矛盾的价格”一注(承“我的孤儿”之后,n\textsuperscript{o}~47,第~195~页)的末尾。自写那条注以来的一年里,这份印象只有加深的份儿——带着一条与我在所引注的子注~n\textsuperscript{o}~47$_{3}$中所表达的根本相同的保留:即Deligne关于Weil诸猜想(WeilI与II)的工作,以及“好上帝定理”(它把好上帝连同他的仆人Zoghman一并清洗出去)与相交上同调的“大冲刺”之后的新出发。但这些局部性的成功,在我看来,与那些我熟知的、此后“入驻”这个“辉煌主题”的人们的卓越——乃至超群——的手段,全然不相称;自离去至今十五年已经过去。也与我遗赠的那些主导观念的丰富与强劲全然不相称——而今天我重又见到它们时,它们已气息奄奄……},十五年后所呈现的惊人停滞,并非源于智力手段或天赋的匮乏(在我认识得那么好又那么坏的那些人中间,不止一人天赋出众),而是源于掘墓人的心性,或肆无忌惮的裙带之风,或两者兼有——那是一些与那种使人认出、并使人找到简单而本质之事物的天真无邪截然对立的心性。

\noindent 为发展他那门新哲学,Mebkhout汲取了导出范畴与六种运算的精神——那时导出范畴正被当作Grothendieck式的花架子,而他连听人说出“六种运算”这个名字的机会都不曾有过。今天,随着新工具的问世——它与导出范畴密不可分——人们大张旗鼓地把后者重新发掘出来,却对两个人的名字讳莫如深:一个是多年孤独劳作、把它们从虚无中提炼出来的人,另一个是从中汲取灵感、终于使一门把拓扑、复分析与代数几何连为一体的新系数理论开花结果的人——后者同样形单影只。

% ---------------- p.799 ----------------
\dmark{}Deligne们、Verdier及其同伙一窝蜂扑向这些簇新的时兴货,一面嚷嚷(那份\origpage{1042}照例必得、冠冕堂皇的谨慎,这还用说):“是我,是我!”。他们中间,还没有一个能在自己身上找到那份勇气与对自己的忠诚,把一份视野在孤独中养熟,沉重地怀上它几个月、几年,远离掌声——那时只有他们自己看得见,而他们所见之物在这世上无法与任何第二个人分享。

\noindent 可我扯远了,该回到我对\textbf{一个愿景之萌发}的讲述了。就在1976年——Mebkhout证明那条“统罩”Poincaré对偶与Serre对偶的对偶定理的同一年——他便得出了三个范畴等价的观念;这三个范畴分别体现同一实在、同一类型对象的“拓扑”侧面、“代数”侧面与(超越的)“解析”侧面。以一门“上同调系数”的普遍理论的眼光\fnote{750}{(*)}{“\textbf{系数类型}”各各不同的这一观念——其中每一个在我面前都显现为六种运算形式主义(与双对偶性)的一个特殊化身,或多或少切近于一切系数类型中最精细的那个、“绝对的”或“普遍的”类型,即“\textbf{Motif}”——或许正是引导我贯穿整个六十年代、尤其自1963年以降,在发展我对代数簇及其他对象的上同调视野时的那个最主要的主导观念。这一观念在我身上的力量,在我专门回顾我的著作及其在时风手中的坎坷遭遇的第一条注(“我的孤儿”,n\textsuperscript{o}~46)里,已清晰可见。在葬礼反思的各处,我以执着的口吻回到它,尤其在“墓上的旋律——或自负”与“工地的巡礼——或工具与视野”(n\textsuperscript{o}~167、178)两条注中。这也是一个首要的数学主题——在被我的上同调前学生们连同一场时风埋葬的那些主题之中——我打算在《收获与播种》之后加以发展,为它在数学思考中取得它应得的位置。\par 奇怪的是,我上同调著作的这一核心主导观念,以及表达它的那个(归根结底十分简单的)代数-范畴结构,在文献中从来不曾被明确陈述过,连我自己在六十年代也不曾(x)。它在我笔下的著作中见于字里行间,而它的传布主要通过口头的交流。在我心里,这原是理所当然:我的学生中间总会有一个,舍得花上所需的几天或几个星期,把这组观念系统地陈述出来——那时我正全力投入EGA与SGA的基础工作。\par 事后回望,我更清楚地意识到:一些非形式的文字(在此情形哪怕只有几页,也不必追求精确而系统的表述)——它们恰好把那些难得提名道姓、却隐藏在往往面目技术性的文字背后的“主导观念”传达出来——这样的文字对于给研究者指引方向、对于不时给一部容易被自己的技术性闷得透不过气的文献送进一口新鲜空气,是何等重要。说到这里,Zoghman对我说过:他在我的笔下找到的寥寥几段这样的文字,曾给他很大的帮助。其中,他近来还特别向我举出我附在Hartshorne那卷《留数与对偶》(原书拼作 \emph{Résignes and duality}——那卷书主要阐述我五十年代下半期在凝聚框架中发展的六种运算形式主义)卷首的几行引言。如今我掂量得出:倘若当年我肯费神在里面加进哪怕一两页非形式的文字,讲一讲“六种运算的瑜伽”、并强调它作为一条无处不在的引线在建造那些尚待诞生的上同调理论时的意义,这几行引言本会对他有用得多……\par(x)(5月24日与6月1日)这些行文写成之后,事情显明:早在SGA~5口头讨论班的最开头(在我的第二讲里),我就曾花大力气,长篇阔论地展开六种运算的“抽象”公式库——它后来主宰了讨论班的全部进程。(关于这一点,见“祖先”一注(n\textsuperscript{o}~171(i))第~942~页5月8日的脚注(*)。)此外,在整个口头讨论班期间,我不曾放过一个机会,不断指及我所发展的那套上同调形式主义的无所不在——它在原则上对“$\ell$-adic系数”之外的形形色色其他“系数”类型一概有效。而Illusie,却精心地把谈论六种运算形式主义的那一整讲、以及任何一丝超出讨论班主要对象这一特定语境的“上同调系数”视野的暗示,统统从屠杀版里剜了出去。\par 关于此事另见“死去的书页”一注(n\textsuperscript{o}~171(xii)),以及“无用的细节”一注(n\textsuperscript{o}~171(v)),(b)部分(“无所事事的机器……”)。},\origpage{1043}我把这些对象称作“德拉姆—\dmark{}Mebkhout系数”\fnote{751}{(*)}{(5月30日)在(后来写成的)“五张照片(晶体与D-模)”一注(n\textsuperscript{o}~171(ix))中,我采用了一套略有不同的术语,把“同一类型对象”径直称作“德拉姆系数”——在这里我们将给出它的三幅\textbf{描述}(或三张不同的“照片”)。其中两幅将分别冠以“de~Rham--Mebkhout系数”(或简称“Mebkhout系数”)“无穷阶”与“有限阶”之名。}。设$X$是一个光滑解析空间\fnote{752}{(**)}{(5月30日)在这几条注的初稿里,我因偏爱“代数几何”观点而忘乎所以,曾假定$X$是$\mathsf{C}$上的代数\textbf{簇}。这并不符合Mebkhout当初置身其中的框架;更何况,这还害我陈述了一个关于$D^{\infty}$-模复形的“好上帝定理”的变体——它只有当假定$X$为真(固有)时才如其所是地为真。于是我的心里出过一些误会,Mebkhout不得不温言提醒我注意。在把这些页重新誊清时,我作了应当作的订正。},那么一方面有 % ---------------- p.800 ----------------
由$X$上“可构”$\mathsf{C}$-向量空间复形组成的(“导出”)范畴 $\mathrm{Cons}^{*}(X,\mathbb{C})$,或径记作 $\mathrm{Cons}^{*}(X)$(\textbf{“拓扑”侧面});再是由那些具有凝聚上同调层的 $D^{\infty}_{X}$-模组成的复形构成的范畴\fnote{753}{(***)}{关于模理论与D-模理论的定义及最初那些连锁的小事实,读者可参看“五张照片(晶体与D-模)”一注(n\textsuperscript{o}~171(ix)),尤其是(a)与(b)两部分(“‘德拉姆系数’图册”与“好上帝公式”)。},\origpage{1044}\dmark{}它是无穷阶微分算子复形的推广,我记作 $\mathrm{DRM}^{*}_{\infty}(X)$(超越的\textbf{“解析”}侧面);最后是由那些具有凝聚上同调层的 $D_{X}$-模组成的复形构成的范畴——它是寻常微分算子复形(有限阶)的推广,我记作 $\mathrm{DRM}^{*}(X)$(\textbf{“代数”}侧面)。有一个重言式的标量扩张函子,从凝聚环 $D_{X}$ 到环 $D^{\infty}_{X}$:

\begin{quote}
$i\;:\;\mathrm{DRM}^{*}(X)\longrightarrow \mathrm{DRM}^{*}_{\infty}(X)$
\end{quote}

\noindent 它嵌入一个(本质上交换的)函子图表:

\begin{equation*}
\begin{tikzcd}[row sep=2.5em, column sep=2em]
\mathrm{DRM}^{*}(X) \arrow[rr,"i"] \arrow[dr,"m"'] & &
\mathrm{DRM}^{*}_{\infty}(X) \arrow[dl,"m_{\infty}"] \\
& \mathrm{Cons}^{*}(X) &
\end{tikzcd}
\tag{1}
\end{equation*}

\noindent 其中两条斜箭头是“相结合的 De Rham 复形”箭头\fnote{754}{(*)}{(5月24日)见已引的“五张照片……”一注(n\textsuperscript{o}~171(ix)),(a)部分。}——它不是别的,正是 $\mathrm{RHom}_{D}(\mathrm{Sp}_{*},\,\cdot\,)$,这里 $D=D_{X}$ 或 $D^{\infty}_{X}$,而 $\mathrm{Sp}_{*}$ 是 $\mathcal{O}_{X}$ 经由局部自由 D-模的“Spencer 消解”。

\noindent 那些竖直箭头的存在,来自“Kashiwara可构性定理”——它蕴含:与一个完整D-模复形相结合的 De Rham 复形,其上同调层是解析可构的。Kashiwara早在1975年就证明了这个重要的定理\fnote{755}{(**)}{Masaki Kashiwara,\emph{On the maximally overdetermined System of linear differential équations}, I,Publ. RIMS, Kyoto university 10 (1975),第~563--579~页。},但却是从一个全然不同的着眼点。他处理的是单独一个完整D-模,取它的 De Rham 复形,证明其上同调是可构的。直到1979年9月、直到好上帝定理引发的“大冲刺”为止,他同上流社会里任何别人一样,从不曾以导出范畴的精神工作;把那些竖直箭头写进图表(1)这个念头,压根不曾浮上任何人的心头!

\noindent 一旦(1)的三条箭头作为导出范畴之间的箭头写下\fnote{756}{(***)}{严格说来,更正确的说法或许是:这是寻常意义下导出范畴的(由“可构性”或凝聚性、完整性、正则性诸条件定义的)全子范畴。},问题便随之而来:它们果真是范畴等价吗?Mebkhout自1976年起就对这一点深信不疑。他的信念来自\dmark{}\origpage{1045}编列的一张约莫十个“典型的”例子的清单(复现于他与 Lê Dung Trang 合写的那篇综述文章里\fnote{757}{(*)}{Lê Dung Trang 与 Zoghman Mebkhout,\emph{Introduction to linear differential Systems},Proc. of Symposia in Pure Mathematics, Vol.~40 (1983), part 2,第~31--63~页。Zoghman向我推荐了这篇短文,说它是文献中对他自1976年以来所发展的那门哲学的最好导引。在该文书目里还可以找到一份(完备的?)Mebkhout关于这一主题的出版物清单,至少截至1983年。}),都是那些可以称作“初等的”可构$\mathsf{C}$-向量空间层——它们同时也是平展上同调理论中人人熟悉的、层的“拆卸”作业里不断出现的那类对象。就在这决定性的1976年,对其中每一个层,他都成功地在$D_{X}$(“\textbf{代数}”)上、也在$D^{\infty}_{X}$(“\textbf{解析}”)上构造出一个非同寻常的完整复形——以六种运算的眼光看,它具有非常简单的代数或解析的上同调意义,而它的 De Rham 复形恰是那个层。值得注意:他从出发时手里拿的是一个可构层、而非一个层复形,可在相当一批情形里,生出它的那个完整复形却绝非仅有一个上同调层。这向他清楚地表明:循着“六种运算”的精神(他并不知道这个名字……),倘若真有等价,它绝不可能从模层范畴(在$\mathsf{C}$上或在$D$上)彼此之间的等价推演出来;它唯有在过渡到导出范畴时才取得意义。

% ---------------- p.801 ----------------
\noindent 在我看来一清二楚的是:此时的\textbf{创造行为},就在于看出并写下那两条\textbf{显而易见的}箭头 $m$ 与 $m_{\infty}$——人人却都不屑一写,就在于提出那个“再幼稚不过”的问题:它们会不会碰巧正是范畴等价,从而给出可构$\mathsf{C}$-向量空间层(或层复形)这一拓扑观念的一个微分代数的、以及另一个微分解析的解释。\textbf{问题在那里,对这一问题之关键地位的清晰意识也在那里}——意识到它的分量;而由此,作为理所当然之事,一种把这问题\textbf{担在自己肩上}、要把它坚持到底的内在姿态。以那些“典型的”或“初等的”例子所作的预备性“实验”,便是朝这个方向迈出的第一步。

\noindent 那是婴孩般的一步,也是本质的一步——唯有知道如何独处的人才迈得出这一步。一旦迈出,我的上同调学生们随便哪一个,凭着在我身边、在SGA~4与SGA~5中学到的拆卸与消解技术\origpage{1046},都能在短短几天\dmark{}之内把它证明出来,或者几个星期——只消他“接得住”,当然,只要他能(像Mebkhout那样用整个内脏)感到这个问题的意义、它的\textbf{实质}。可他们中间没有一个人,连那个为提炼一套超越“六种运算”主导观念、\textbf{超越}其上——而它正是把连续系数与离散系数联结起来所仍欠缺的——的统摄视野而\textbf{宣告弃权}的Deligne也包括在内\fnote{758}{(*)}{(6月5日)重读时,这个提法在我看来显得仓促,而且有点“不着边际”。实际上,我的“六种运算这一主导观念”与一门“系数哲学”密不可分——后者(至少从1966年起就十分明确地)预见了一套“德拉姆系数理论”(与我的晶体观念血肉相连),它具有与$\ell$-adic系数理论相同的本质形式性质,并与后者(当$\ell$变动时)共同构成同一个终极对象类型——“Motif”——的种种不同的“实现”。Mebkhout完成于1972至1980年间的工作,在我看来是朝这一直觉的实现迈出的第一大步——迈进这一步的条件,实际上至少从1966年起、随着晶体瑜伽的启动就已完全成熟;那一年,一套德拉姆系数理论的问题(至少在我的心中)已被清晰地提出。假如我的上同调学生们中间——早自六十年代——无一人迈出这一步,在我看来这主要应归咎于种种堵塞自发创造力的机制;而那种创造力,他们中间没有一个人是缺乏的。关于这一点,见“……与羁绊”一注(n\textsuperscript{o}~171(viii))。},没有一个能看出Mebkhout的思想、这位依旧浑身透着我那股气息的无名之辈的思想的分量——尽管那分量一望即知……

\noindent 至于那位“无名之辈”,凭着他自己的手段与他自己的阅读,向自己提出范畴等价的问题,在他看来(而且也确实)是天底下最显而易见、最孩子气的事,或者由此达到确信:这些箭头货真价实地就是范畴等价。相反,由于缺乏经验、缺乏比他更有经验的长者的鼓励,那证明在他眼里成了一座大山,长久以来显得完全高不可攀。

\noindent 然而仅仅一年半之后,他便找到了一个证明——起初是对箭头$m_{\infty}$,在1978年3月。他对我说,在心理上,我的代数与超越 De Rham 上同调之间的比较定理曾给他很大帮助,把他引上了证明之路。出于一个我未曾很好理解的理由,他还把自己的定理(即:函子$m$——人们称它“好上帝”函子,恰恰是为了不提Mebkhout……——是一个等价)\origpage{1047}视为我的比较定理的一个“推广”。从这一刻起,他也知道\dmark{}自己手头已有必需的工具(连同Hironaka的消解技术),足以处理$m$的情形——对一个像我这样的代数几何学家来说,那才是远为有趣的情形。而他,身为分析家,起初倾注心力的本是函子$m_{\infty}$的情形——那原是他的偏好\fnote{759}{(*)}{(5月24日)另一个也许更强劲的缘由是:在$D^{\infty}$-模的情形,他手头有一条绝妙的反演公式——关于这一点,见“五张照片”一注(n\textsuperscript{o}~171(ix)),(b)部分“好上帝公式”。}。对那个在他看来多少有些次要的问题,他直到论文答辩之后才重新拾起,并在一个月之后(1979年3月)证明:函子$m$(就是那个今天人人以迂回说法使用、却从不肯写出来——免得提名道姓一位不可提名道姓的作者的……函子)确实是一个范畴等价\fnote{760}{(**)}{Mebkhout把“$m$是一个等价”这一事实的正式成文的证明(其原理与“解析的”好上帝函子$m_{\infty}$的那个证明相同),是两年之后、1980年末才写出的。这份证明陈述于两篇相继的文章中的第二篇(第一篇处理解析的好上帝函子$m_{\infty}$,重述其论文),即“一个范畴等价”与“另一个范畴等价”,载 Compositio Mathematica 51 (1984),第~51--62~页与第~63--88~页。(手稿1981年6月10日收到。)但早在1969年3月[原书如此,当为1979年]起、在其后的岁月里,凡有机会,他便传播这一结果(连同关于函子$m_{\infty}$的那个),尤其是同年6月就告知了Deligne。\par 我想,由于他极度的孤立,加上他作为分析家的“眼镜”,他不曾意识到:真正会引起Deligne之流兴趣的,恰恰是代数的好上帝函子——因为它在拓扑与代数几何之间搭起一座“桥”(算术尚待到来,而看来我是第一个、并且是唯一一个预见到它的人……),其分量堪比平展上同调工具所提供的那一座。否则,他定会立刻把它写成正式的文字并火速发表——尤其考虑到他误入的那个古怪圈子(他当时还一无所知……)的种种行规。不过,他在1980年3月的头一场噩运(与Kashiwara打交道),本该让他心里起疑(x)。\par 恰恰也是在同一个3月,Mebkhout在法国科学院《通报》上发表了关于“Riemann--Hilbert问题”的一篇短注(第~290~卷,1980年3月3日,Série A,第~415~页),其中他陈述了论文里(关于$m_{\infty}$)的等价定理,并谨慎地声明:“人们希望证明——以对偶定理[7]所用的上同调下降法为法——函子$S$[即我所谓的$m$]乃至$T$[即我所谓的$i$]也是范畴等价”。实际上,他的证明表明:它们是“在$X$上局部地”的等价;单是这一点,已经蕴含了著名的Kawai-Kashiwara定理(下一子注将谈到它),即:标量扩张函子$i$在完整正则$D_{X}$-模范畴与完整$D^{\infty}_{X}$-模范畴之间诱导一个等价。我顺带指出:Mebkhout的最终结果,即便应用于单个的模(而非模的复形),也远为有力,因为他同时断言:由“标量扩张”函子导出的那些典范箭头

$\mathrm{Ext}^{n}_{D_{X}}(M,N)\longrightarrow \mathrm{Ext}^{n}_{D^{\infty}_{X}}(M_{\infty},N_{\infty})$

也是同构(而不止于$n=0$)。\par(x)(5月25日)在4月24日的一封信里,Mebkhout还向我说明:“要知道,论文答辩之后我松了一口气。整整四年,我一直处在极大的紧张之中。”}。由此立刻得出:那个从“代数的”(他当时对它还不曾有多大兴趣)通往“解析的”(超越的)的“换环”函子$i$,同样是一个等价。

\asterisms

% ---------------- p.802 ----------------
\dmark{}1978年3月,Mebkhout与他的“恩人”Verdier有了第三次会谈——此前他们已两年未见。他向他讲解了(未来的)“好上帝定理”的来龙去脉;\origpage{1048}他谦逊地(这一谦逊害了他!)称之为“Riemann--Hilbert等价”。回想起来,Mebkhout确信:自己的解说多半从Verdier的头顶上飘了过去。可以肯定的是:Verdier压根没有意识到,他这位“受保护者”刚向他提交的是一些值得驻足的观念。他没有向身边任何人提起,连Deligne也没有——Deligne直到一年多以后,才从Mebkhout口中得知好上帝定理(连同那条号称“Poincaré--Serre-Verdier”的对偶定理——三年前这位Verdier本还死活不肯相信它……),那是1979年6月的Bourbaki讨论班(答辩之后四个月)。无论如何,Verdier开了绿灯,让Mebkhout把他的结果作为国家博士学位论文提出,并且同意组建并主持答辩委员会。答辩之所以推迟了一年才举行,是因为那个有名的“巴黎地区各大学论文委员会”办事的拖拉(对这个机构,Verdier可是珍爱如眼珠……)。

\noindent 正如我在一条先前的注里说过的\fnote{761}{(*)}{见“……与意外之喜”一注(n\textsuperscript{o}~171(iii))。},答辩在一种人人漠然的气氛中进行。Mebkhout固然可以把论文左一份右一份地寄出去,它依然默默无闻——连那块大部头的一份收执,都没有人屈尊回一张。

\noindent Mebkhout却没有被击倒。尽管种种迹象指向反面,他仍感到自己属于一个“家庭”——毕竟是一群做同一类数学的人:那类数学,他很大程度上是\origpage{1049}凭读我的著作学来的,而更多的是凭对这几部著作中的某种\textbf{精神}持开放、\dmark{}倾听的姿态学来的\fnote{762}{(*)}{人们不禁要问(或者要问我):我的著作里那种如此独特的著名“精神”究竟是什么,它竟启发了我这位“身后的学生”Zoghman~Mebkhout,却遭我其余所有学生——Deligne打头——以及一支亦步亦趋的时风的“排斥”。倘若我试着为这种精神找一份谱系(在我对数学史的那点支离破碎的知识允许的范围内),我会说:它属于\textbf{Galois、Riemann、Hilbert}的谱系。倘若我试着以心灵中种种作用力的动力学来界定它,我会说:这是一种通过“阴”“阳”两种创造力之和谐平衡而显自身的精神,其“基调”或“主导音”是\textbf{阴},是“女性的”。对数学中、乃至一般地对世界之发现中这一进路的更详尽的描述,可在“上涨的海”“九个月与五分钟”“阴的葬礼(阳葬阴(4))”(n\textsuperscript{o}~122、123、124)诸注的反思中逐步见出;这一反思又在“兄弟与配偶——或双重签名”“阴,作为仆人——与新主人”“阴,作为仆人(2)——与慷慨”(n\textsuperscript{o}~134、135、136)诸注中得到接续。至于当代世界对这种“精神”的某些“本能的”拒斥机制,见“天赐的情势——或封神”与“否认(1)——或提醒”两注(n\textsuperscript{o}~151、152)。}。他显然还不曾意识到——至少在意识的层面上——这种精神恰恰早已被那些构成他自信已然加入的“家庭”的人们所排斥;而对那些踩着厚绒地毯步入数学的翩翩君子来说,他不过是个拖拖拉拉的落魄鬼、一个闯入者。

% ---------------- p.803 ----------------
\subsubsection*{c2.~黑手党}

\tocsubsub{c2.~黑手党}

\notehead{171$_{2}$}{}

\noindent (4月15--17日)

\subsubsection*{(a) 家庭合影上的阴影}

\tocsubsub{(a) 家庭合影上的阴影}

\noindent 可我们这位朋友Zoghman,对这一切还蒙在鼓里;他虽则孤立,却并不不幸。自1973年起,他有幸在奥尔良谋得一个助教的职位,这让他得以安安静静做他感兴趣的数学——哪怕眼下感兴趣的只有他一个。他继续住在巴黎地区,继续出入讨论班,继续熟悉文献……

\noindent 倘若他稍稍琢磨过这件事,他本会发觉:在这个佯装不知有他、而他自信身属其中的“家庭”里,一切并不那么美满。凭着阅读我的著作,他终于多少觉察到:曾对他如天降甘露的那份“好的参考文献”,至少有相当一部分,根本不出自他“恩人”Verdier的自酿。可构性概念早在1963年就在SGA~4里得到长篇展开,比Verdier在那篇文章里装作发明它早了十二年。随着SGA~5于1977年问世——哪怕是以Illusie屠杀版的形式——事情变得清楚了:那著名的、关于解析或代数可构的$\mathsf{C}$-向量空间层复形的“Verdier双对偶”,被原封不动地\dmark{}\origpage{1050}抄自SGA~5的第一讲——那个名字古怪的书卷“SGA~4$\frac{1}{2}$”所指涉的正是这一讲,其说法是:“若干补充见 SGA~5 I”\fnote{763}{(*)}{关于这个叫人哭笑不得的委婉说法——它意在(这一回由他Deligne)把那条倒霉的双对偶定理据为己有——见“特洛伊木马”子注(n\textsuperscript{o}~169$_{3}$)第~872~页的脚注(**)。}!在同一部古怪的书卷里——它的作者乐于以十足的高傲俯视环绕着它的姊妹卷SGA~4与SGA~5——他还可能读到关于结合于一个闭链的上同调类的一讲;那一讲(不知何故)被人从那部“技术性离题”的书卷SGA~5(号称更晚……)中卸了下来;他由此同时可以看出:他那恩人的文章得名的那个主题的上同调侧面(同调侧面的对偶),同样是抄自SGA~5。然而,对“好的参考文献”里的这三个主题\fnote{764}{(**)}{这里说的是那“三个主题”:可构性,可构层的双对偶性,结合于一个闭链的上同调类(与同调类)。},却没有任何一处提到我的这个人或SGA~5……

% ---------------- p.804 ----------------
\noindent 他固然还无从知道:Verdier那篇文章剩下的部分(五十页里除去三页)是从我关于平展同调形式主义、关于结合于代数闭链的同调类的报告中“抄”来的——那些报告(仿佛有约在先)已从那部声名可哀的Illusie版中消失,连一丝提到它们存在过的影子都没有。但他手头的那点事实,对一个知情而警醒的人来说,确实已经绰绰有余。总而言之,这情形同我十年前翻阅Deligne关于谱序列退化那篇文章时的处境一模一样:在那篇文章里,他既把最初的动机、整个权的瑜伽(连同我这卑微之人的角色)掉包藏匿,也把Blanchard思想的贡献藏匿起来——那贡献恰恰是把“蛮横的”Lefschetz定理用于纤维之上\fnote{765}{(***)}{详情见“驱逐”一注(n\textsuperscript{o}~63)的开头,以及该注第~233~页的脚注(**)。}。像当年的我一样,Zoghman那时只得把对一件不快之事的清明感知压下去,对自己说(在此情形):这想必是师父与学生们之间惯常的一种“默契”——当学生们把他直接传授的观念、技术、结果当作自己的东西来呈现时,师父便睁一只眼闭一只眼\fnote{766}{(****)}{(5月30日)而且顺带还温言称他是个江湖骗子……}。正如在\dmark{}这类情形中常有的那样,这种解释(它正合Zoghman的心意)本就不乏真实的成分——而且还不仅如此。\origpage{1051}我不止一次确实曾是这类暧昧局面的当事人。(但也属实:在我离去之前,事情从不曾发展到这个地步——师父的著作成了遗骸,人们厚着脸皮瓜分它的碎片……)

\noindent 再说,在那个由所有对簇上同调感兴趣的人组成的扩大的家庭里——包括Sato学派的日本人——情形也不怎么美妙。那位Kashiwara,他1975年的可构造性定理曾为定义“好上帝函子”立下汗马功劳,却也装模作样,要把那些倒霉的可构造层——忽而间人人竞相争抢的东西!——的发明权据为己有。为着这份事业,他在所引文章的第2节里给它们重新施洗为“finitistic sheaves”,在该节里或多或少逐字照搬SGA~4中关于此主题的展开。据我从多方听到的,Sato学派熟悉我的上同调著作,哪怕他们引用我时吝啬得很\fnote{767}{(*)}{Mebkhout就此给我写道(1985年4月24日):“我在Sato的日本学派那里见到的唯一提到你的地方,涉及EGA~III第0章——而他们却在毫无顾忌地从你的著作里汲取灵感。”};很难相信Kashiwara不知道可构造性概念——至少在平展语境中,在那里它是全部理论的中心有限性概念。不言而喻,第二年Verdier对“finitiste”(原文如此)概念照样不提Kashiwara,正如他对某位亡故者、某个讨论班只字不提\fnote{768}{(**)}{仿佛有约在先:这个讨论班(SGA~5)恰恰就是(与SGA~4一道)经我的上同调学生们共同商定、并按他们头目Deligne的说法,注定要被“忘掉”的那一个(多亏了由他亲笔操刀的那份“电锯摘要”的出版……)。}。两人都是上流社会中人,说不定还同属一个“家庭”——可一到作者虚荣这牛排面前,各人只顾各人的……\fnote{769}{(***)}{(5月24日)Mebkhout提醒我,我在这里把画面涂得略黑了些。Verdier对Kashiwara的文章一无所知,对完整性(holonomie)概念同样一无所知——那是Mebkhout在1976年与Verdier的那次“会谈”中教给他的。(这发生在“好的参考文献”刊行之前——它似乎刊于1976年末;不过按情理,既然他知道自己与他的同行都在从同一处未提名的源头“抄”,也就不能指望他会去引Kashiwara……)反过来,Kashiwara不知道“好的参考文献”、也不知道我的双对偶定理(它在那里顶着Verdier的父名)——是Mebkhout在1978年1月把这两样东西连同他论文第三章的结果一并告知他的。而后,那些结果就在已引的Kashiwara--Kawai的文章里被毫无顾忌地(而且几乎不带证明地)据为己有——关于这一点,见“五张照片(晶体与D-模)”一注(n\textsuperscript{o}~171(ix)),尤其第~1005~页。Kashiwara不知道离散系数的双对偶定理这件事,连同这里或那里记录的许多别的迹象,表明他离Mebkhout那门直接源于我著作的对偶哲学有多远。}

% ---------------- p.805 ----------------
\dmark{}我看,对Zoghman来说\origpage{1052},宁可对自己说,某个他从未谋面的日本人\fnote{770}{(*)}{(5月24日)他终究见过他们一面,这些大名鼎鼎的日本人!Mebkhout就此给我写道(1985年4月22日):“1972年,Sato学派全员出动,来开一个关于超函数的会议。他们把自己的方法藏得很好。很长时间里,他们的成果都无法企及。围绕这个学派有某种神话,如今Kashiwara之所以能随心所欲,正赖于此。”\par(6月4日)应当说,倘若果真(如Mebkhout在此所暗示的)是Sato学派首创了那种为求支配而故示幽暗的手法,那么这手法在太平洋此岸也不乏徒众——如今他们比起师父来毫不逊色!而把反常层会议那场令人难以置信的蒙骗局操办起来的,恰恰是他们,而绝非Kashiwara之流:在这场蒙骗局里,Kashiwara被当作一枚方便的“小卒”用来铺路——随后便遭丢弃……}准是个“骗子”,也不愿不得不在下述事实面前认账:那是几位显赫的长辈——其中一位,于他犹如一位强有力而高不可攀的父亲、一位恩人;是他有机会在讨论班里晤对的长辈,而且他还荣幸地同他们以“你”相称(自Bourbaki时代以来,法国数学圈的惯常做派)。

\subsubsection*{(b) 最初的麻烦——或太平洋对岸的卡伊德}

\tocsubsub{(b) 最初的麻烦——或太平洋对岸的卡伊德}

\noindent 说来吊诡,Zoghman的麻烦,恰恰是从某个圈子开始察觉他随一整门哲学(诚然是一种看来分明已过时的哲学……)带来的某一件工具之威力那一天开始的。1979年6月,他已把这一点告知Deligne;Deligne聚精会神地听了他关于对偶定理的解说,而(可想而知)更加聚精会神地听了关于好上帝定理的解说。他甚至十分和蔼地对他说,他读过论文的引言,并且觉得这项工作里必有精彩的数学\fnote{771}{(**)}{(6月3日)前一年,在雷恩的$p$-adic分析讨论会上,Mebkhout已经从Illusie口中领受过一个同样无谓的恭维。关于这一点,见“掠夺的空白授权”一注(n\textsuperscript{o}~171$_{4}$),第~1091~页(尤其同页脚注(**))。}。那一天,Zoghman的日子真美——可惜好景不长。

\dmark{}同年,1979年9月,他参加了莱苏什讨论会\fnote{772}{(*)}{莱苏什讨论会(1979年9月1--13日)的《论文集》刊于 Lecture Notes in Physics n\textsuperscript{o}~126 (1980),Springer Verlag。这部《论文集》里,既有Mebkhout的报告《论Hilbert--Riemann问题》——它以完美清晰的方式陈述了他的全部哲学(我会称之为“德拉姆系数”的哲学),证明都有文献支撑——也有Kashiwara与Kawai提出的报告。任何善意的读者,把两篇文章一比即可验证:在这两位作者的文章里,既没有一丝这类哲学的萌芽,也没有半句提到好上帝定理之类的暗示。\par(6月4日)在其4月22日的评注信里,Mebkhout就前一年(1978年8月)举行的赫尔辛基国际数学家大会表达了同样的意思:“我必须说,我听了Kashiwara在赫尔辛基大会(1978年8月)上的报告,他是大会的主讲人之一。无论远远近近,那里没有任何一种哲学,可以同离散系数与连续系数之间的比较沾边。我特意把我一周前在哥本哈根所作的报告立即整理成文,交给那公认有裁判之责的数学共同体。而这位Kashiwara先生的报告,则刊于[赫尔辛基]大会《论文集》。”},在会上作了一个报告\origpage{1053} % ---------------- p.806 ----------------
《论Hilbert--Riemann问题》,陈述他的等价定理。他的报告看来是完全无人理睬地过去了。相反,讨论会的一个“压轴节目”,却是几天之前Kawai的一场报告,宣布了一个他与 M.~Kashiwara合作取得的、不寻常而出人意料的结果。以一种多少拐弯抹角、刻意叫人看不懂的形式(循Sato学派特有的风格\fnote{773}{(**)}{(6月4日)关于这一点,见一处先前的脚注(第~1052~页的注(*))。在我看来,恰恰是在反常层会议的尾迹里,这种蓄意晦暗的风格在太平洋此岸被精益求精,变成一门系统蒙骗与趁乱窃取的方法。}),这定理断言:在一个(光滑的)复解析簇上,从 $D$ 到 $D^{\infty}$ 的“标量扩张”函子在“具正则奇点的”完整D-模范畴与完整D-模范畴之间诱导一个\textbf{等价}。他们的证明后来写成一篇一百五十多页的长文,业已刊行\fnote{774}{(***)}{M.~Kashiwara, T.~Kawai,\emph{On holonomic Systems of microdifferential équations III}, System with regular singularities,Pub. RIMS 15,第~813--979~页(1981年)。}。

\noindent Mebkhout当场同所有听众一样,有些懵了。这个被炒得耸人听闻、却没有一个人太弄得清究竟是怎么回事的定理,对他却有一种说不上来的\dmark{}熟悉劲儿。随后几天里,他缓慢然而稳当地、按他惯常的方式\origpage{1054}反复琢磨这件事。我可以想象:在讨论会的喧嚣之中,单是把这定理整理成非日本人能看懂的形式,他少说也得花上一两天。从那一刻起,便稳操胜券了!

\noindent 我敢打赌,在场的西方人里,没有一个对“正则奇点”是什么有过半点概念。可Mebkhout,他几年前就为一门尚在寻找自身的“系数哲学”之需,定义过一个\textbf{正则}完整D-模的概念\fnote{775}{(*)}{关于Mebkhout对一个完整D-模复形(沿一个除子$Y$)的正则性的定义,见“作品……”一注(n\textsuperscript{o}~171(ii))第~950~页的脚注(*)。光说“正则”,意即:沿(每个开集上的)\textbf{一切}除子都正则。}。至少对他来说,这概念有一个十分精确的意义——而且,取合适的\textbf{导出范畴}、并且“走到镜子另一面去”,他知道如何以“可构造离散系数”的相应导出范畴来解释这个范畴。至少,他已在论文中长篇地证明了:就在“镜子另一面”的这同一离散系数范畴的术语之下,$D^{\infty}$-完整模范畴的类似解释——而且他清楚,自己手头已有证明“完整正则D-模”情形的类似物所需的一切。这他在论文里实际上已经做到,其形态是关于$X$的一个\textbf{局部}结果——单是这一条,已足以蕴含Kashiwara--Kawai那个“耸人听闻的结果”。由此,导出范畴的观点,连同连续系数与离散系数之间的游戏,给出了一个Kashiwara--Kawai型的结果,但原则上远为有力:它同时给出诸高阶 $\mathrm{Ext}^{i}$ 之间的同构,而不止于 $\mathrm{Hom}$ 的层面(那是单用D-模、而不用由这些模构造的导出范畴所能得到的全部)。有见于此,倘若这个日本人的“正则奇点”概念竟不等价于他自己的概念,那才是见了鬼——这样一来,那个声誉显赫的结果实际上是他的系数哲学的一个纯然推论;而对这门哲学,此前还没有一个人屈尊理睬过。

\noindent 当讨论会全员出动,来给一个无名之辈的报告捧场——那报告在日程上占着一行,谁也说不清为什么——而报告的末尾\fnote{776}{(**)}{(6月4日)实际上,Mebkhout早在报告开篇就特意暗示过这一点,天真地以为这准能抓住他的听众。},那人面不改色地一连串箭头加图表地摆开(那类玩意儿\dmark{}在六十年代还时兴,在严肃人士之间早已过时\origpage{1055}),宣布说:讨论会那个著名的“压轴节目”——没人能把它的大意复述出来,这只会叫它愈发唬人——那个“压轴节目”原来是一个范畴等价定理的直接推论(说得倒轻巧!),一个他在相应的\textbf{导出范畴}(这又是些什么动物?)之间得到的等价定理,外加一个看上去同它们不大相干的范畴;而这定理,他赌咒发誓说,早已写在一篇\textbf{论文}里(这可真是锦上添花!),而且早就寄给了Kashiwara先生以及满堂显赫同行中的许多人——这情形,怎么看怎么像一个拙劣的玩笑。一阵尴尬的沉默,几声会心的窃笑。那位Kashiwara先生亲自提出例行的那个问题,想必(正是)为了驱散这个不懂规矩的年轻人造成的难堪。不过说实在的,他看上去也有点被震住了;他准在纳闷自己是不是在做梦\fnote{777}{(*)}{(6月4日)Mebkhout就此给我写道(4月22日):“莱苏什讨论会之后,有人告诉我,这位Kashiwara先生觉得他与Kawai合写的文章空空如也。可他为了不老实地补上自己落下的功课,可没少下力气。打从[1975年证明其可构造性定理的]那篇文章算起,已有五年他没碰过离散系数了。[因那篇文章而得来的]骤然的名声——它全亏了另一个全然不同的问题——使他得以去忙一些更‘严肃’的事情——千万别是瞎拼凑!1975到1980年间,在一片普遍的敌意之中(这是我后来才明白的),我是\textbf{唯一}一个在发展我从你的著作里学来的这门孩子气的哲学的人。”}……而那位无名之辈呢,却没有被吓退——他差一点要在第一场报告之上再作第二场——真是无奇不有!

\noindent 一分钟之后,我们这位无名之辈Zoghman发现,自己独自一人立在黑板前,守着他那些漂亮的图表,面对一屋子走空了的座位……那一天,以及随后几天,没有一个人屈尊问一问那个不懂规矩的家伙那些所谓的“结果”的来龙去脉——那人错就错在被邀来这么一场身价显赫的讨论会。

\noindent 应当说,大场面的吹吹打打一旦过去,那位Kashiwara先生的脑袋里总归还是起了波澜的。事实是:不过短短几个月之后,在1979--80年度的古卢阿伊克-Schwartz讨论班上,在4月22日的一次口头报告\fnote{778}{(**)}{(6月4日)古卢阿伊克-Schwartz讨论班1979--80年度,M.~Kashiwara1980年4月22日的报告,“具正则奇点的线性偏微分方程的完整系统与可构造层”。关于这场值得纪念的报告会——\textbf{Mebkhout当时在场}——的详情,见“掠夺的空白授权”一注(n\textsuperscript{o}~171$_{4}$)。}里,他宣称\textbf{那是他自己的创造}——正是那个\dmark{}\origpage{1056}曾在某次讨论会上善于投下一阵冷场的定理!不过他倒还有那份“好意”,在第2页补了一句:

\begin{quote}
“请注意:该定理\textbf{也}由Mebkhout\textbf{以另一条途径}作了证明”(着重是我加的)\fnote{779}{(*)}{我在此引用的是那份书面讲稿的文字;它由Kashiwara在口头报告一年之后整理而成。详情见前一脚注所引的那条注。}
\end{quote}

\noindent 这个“也证明了”真可谓字字千金——要知道,这是一个他连同任何别人此前连想都不曾想到的定理,是他(不过几个月前)刚从当事人本人口中听来的,而他连人家差不多一年前寄给他的论文都懒得一读!倘若他先前就知道这个定理,那他肯定不会费神去写那167页密密麻麻的证明,来证明一个“蛮横的”分析结果——而那结果不过是它的一个直接推论,甚至是推论的推论。

% ---------------- p.807 ----------------
\noindent “以另一条途径”同样叫人绝倒。那份报告里没有一丝一毫证明的痕迹,在Kashiwara其后的著作、或其任何一位日本同行的著作里同样没有。Zoghman向我担保:文献里不存在他的定理的第二个证明,只有他自己的那一个;而我(鉴于那种证明的类型——我对它再熟悉不过,理由不言自明)极疑心人们永远也找不到第二个。那是一个对应于一种几何进路的证明,使用的是Hironaka式的奇点消解——这个工具对我(和我的学生们)已成了第二天性,而分析家们(尤其是Sato学派的那些)对它一无所知。以至于Kashiwara显然觉得自己连\textbf{照抄}Mebkhout的证明都力有不逮……

\noindent 这种白线缝得极粗的骗局,\textbf{只要有一种普遍的一致同意罩着它}就能行得通——而受损的(在此)是一个无名之辈。这满堂的高雅人物\fnote{780}{(**)}{(6月4日)若要看一看直接、积极参与围绕Zoghman~Mebkhout著作之蒙骗-骗局的那帮演员的“列队”(至少是我知情的那一些),见“黑手党”一注(n\textsuperscript{o}~171$_{2}$),(f)部分“演员列队——或黑手党”。不过这份列队并不齐全——要看一份更全的名单(一口气排出十三位国际知名数学家的名字),见“光荣之日”一注(n\textsuperscript{o}~171(iv))第~962~页的脚注(*)。其中还缺一个名字:R.~Remmert——他在此期间浮出水面(见已引的“黑手党”一注,(c1)部分“记忆的失灵——或新编历史”)——十三变成十四!(还没算那位隐名的审稿人——十四变成十五……)}大可不必客气:须知上述那位\dmark{}无名者,分明正被那样一些人弃之不顾\origpage{1057}——他们恰恰是最有资格第一手知情、且对当事人负有直接人身责任的人:J.L.~Verdier(答辩委员会主席)与 P.~Deligne(头一个感受到那结果之分量的人——他是前一年从Mebkhout口中得知它的)。

\noindent 既然说到了太平洋对岸的卡伊德Kashiwara,不妨就在这一章上收尾,讲一讲那位当值无名之辈被彻底清洗的尾声——顺着三年前1981年6月反常层会议上那道响亮的先例。这说的是 R.~Hotta与 M.~Kashiwara的一篇文章“半单李代数上的不变完整系统”(Inventiones Mathematicae 75,第~327--358~页),刊于1984年(1983年3月2日收到)。从引言第6行起便可看清:这篇文章是那道没完没了的、所谓好上帝(或当值无名之辈)的“Riemann--Hilbert对应”的又一应用。在这篇文章里,\textbf{那位无名者的名字不再被提起},书目里也没有他。Zoghman已知第二作者的心性,却无从预断第一作者是否也不老实,于是给他写信,告诉他自己是文中以举足轻重的方式用到的那条定理的作者,并就自己未曾以此被引用一事提出异议。结果呢?引文给了已引的Kawai-Kashiwara那篇文章(167页那篇)——那条定理在那里恰恰根本不见踪影\fnote{781}{(*)}{(5月25日)如他处已解释过的(在“五张照片(晶体与D-模)”一注,n\textsuperscript{o}~171(ix),尤其第~1005~页),那部著作只含好上帝定理的“一半”,而那一半是从Mebkhout论文第三章劫掠来的。}。Hotta回复他说:他觉得没有必要引用他,因为\textbf{反正人人都知道,那条对应是Kashiwara与Mebkhout的}。幕落……

% ---------------- p.808 ----------------
\subsubsection*{(c) 入门费——或有为青年}

\tocsubsub{(c) 入门费——或有为青年}

\noindent 但日本毕竟路远;我的朋友Zoghman经年累月地同遥远的日本人交锋,想必是因为:对他来说,直面一桩绝不限于天涯海角大陆的黑手党的现实,要痛苦得多——这黑手党无论在巴黎那些气派的讨论班里,还是在莫斯科或\origpage{1058}东京,都高踞\dmark{}首座。该回到这可爱的法兰西国土了,回到我亲爱的上同调前学生们组成的“小家庭”——以及(那稍大一点的)自我“亡故”的遥远时日以来围绕他们结成的那个圈子。

\noindent 消息有时跑得飞快。1979、1980年间,借着Deligne与莱苏什讨论会的东风,“人家”总该渐渐察觉:数学市场上刚刚冒出一个定理,说实在的颇有苗头,可惜出自一个落伍的Grothendieck派余孽;不过,这个叫人提不起劲的父权,有一个现成的替身——人所共知的日本分析家Kashiwara,他正求之不得,要出演那著名的“Riemann--Hilbert对应”的父亲。

\noindent 1980年1月,Mebkhout就他那条来得不是时候的定理,在巴黎第七大学 Lê Dung Trang 主持的“奇点讨论班”上作了一个报告。Jean-Louis~Brylinski没有到场听讲,但 Lê Dung Trang 向他谈起了它,并让他读了自己的笔记。据Brylinski自己后来向Mebkhout坦率讲出的:他一得知Mebkhout的定理,便脱口而出——有了这个,Kazhdan--Lusztig猜想就有法子证明了!(猜想嘛,照占卜师们的说法,是“无从下手的”。)

% ---------------- p.809 ----------------
\noindent 人们或许以为,Brylinski会去找当事人,请他把自己所需的那条定理所系的“完整性”与“正则性”条件的种种奥妙讲个明白。可照他自己向Mebkhout坦白讲出的:是“人家”劝他别去找当事人,而去找那位大人物Kashiwara。他没有点明那“人家”是谁。不过显然,他耳朵尖(脑子也快),而且那时候他也还同Mebkhout至今一样默默无闻。这种话无须说第二遍,他便去找Kashiwara打听去了——Kashiwara那时想必还在近旁。那是他最正当不过的权利。结果便是与Kashiwara合写的一篇文章,刊于 Inventiones Mathematicae 64(第~387--410~页),1981年(1980年12月19日收到),题为“Kazhdan-Lusztig 猜想与完整系统”。Brylinski一夜之间成了明星——这原是他应得的;而Kashiwara呢,又往那已够辉煌的战绩簿上添了一根花翎\fnote{782}{(*)}{把Kashiwara的名声同他刚刚找到的证明——一个Kashiwara毫无寸功的证明——撮合在一起,同时只字不提那位无名的年轻同行所扮演的关键角色:这就是Brylinski不加推辞地付出的“入门费”,为的是跻身某个“圈子”——名流们的圈子,本注“黑手党”之名即得自它……}。

\dmark{}一切本可以在这最好的世界里尽善尽美,但是……想必还是那个“人家”又出了主意\origpage{1059}:某个无名的泛泛之辈,提得越少越好。反正,寄给 Inventiones 的那部手稿里,\textbf{Mebkhout的名字不见踪影},正文里没有,书目里也没有。

\noindent Mebkhout看到了文章的预印本,就此向Brylinski提出了抗议,并致信 Inventiones 的编辑 R.~Remmert。Brylinski的反应很“柔顺”(那笔调如今对我已再熟悉不过……):在校样上于书目末尾(字母次序之外)添了三条“暂停!”式的Mebkhout引注(顺手为之!),但正文里对这位Mebkhout先生却始终一字不提\fnote{783}{(*)}{Brylinski-Kashiwara文章的引言,以对几位作者的鸣谢收尾,其中包括Jean-Louis~Verdier(而那位当值无名之辈,不必说,只字未提)。随后是第1节,概述“具正则奇点的完整微分系统”(这是日本人对完整正则D-模的叫法)。在该节开头几行,人们读到:“至于细节与证明,请读者参看[6, 15--17]。”文献[6]即Kashiwara1975年确立其可构造性定理的文章,而[15--17](校样上追加的)就是对Mebkhout的“暂停!”式引注。无论发生什么,“有为青年”Jean-Louis~Brylinski的荣誉总是无恙的……}。读这篇文章的人,倘若居然在书目末尾天晓得为什么多出一个“鼎鼎大名的无名之辈”的名字,只会想:准是为了照顾哪位朋友的情面,才把他放在那儿……

\noindent \textbf{Brylinski是靠一场骗局跨入名流的}。实情是:只要那件新工具不曾问世,他所证明的那个猜想便无从下手。且不说那件工具的\textbf{发明权},这篇文章里没有任何东西把这件新工具亮出来——它的角色从第6--8行起就被那句不咸不淡的“解释”(原文如此)掉了包:

\begin{quote}
“本文采用的方法,是把旗流形上具正则奇点(R.S.)的完整线性微分系统同 Verma 模联系起来,\textbf{并使用完整系统与可构造层之间的对应}。”
\end{quote}

\noindent (着重是我加的。)对这条未加限定的著名“对应”,没有片言只字的指涉或解释。“人家”想必让这位年轻后进心领神会:这条“对应”从此便属于人人皆知之事,根本无须为此\dmark{}\origpage{1060}援引一条专门的定理——那会掀起发明权这类枝节的、而且(尤其)不合时宜的问题。而Brylinski,这个有为青年,自然无须人说第二遍……

\noindent 至于Remmert,他把那位无名控告人的信转给了Brylinski-Kashiwara文章的审稿人。审稿人驳回了控诉,表示的意见是:“该结果\textbf{独立地为已知,而且大概更早,系由Kawai与Kashiwara所得}”,指涉的是他归于这两位作者的“重构定理”(指涉上引作者文章第~116~页,载 Guillemin 主编《微局部分析讨论班》,Annals of Math. Studies 丛书,n\textsuperscript{o}~93)。

% ---------------- p.810 ----------------
\noindent 审稿人的这一评价——他可是被认为知其所言——在两个层面上都骇人听闻,并且表明:他本人就是同一场骗局的当事人,与(此刻的)Kashiwara和Brylinski沆瀣一气。仅仅凭一个关于\dmark{}种种独立所得结果\origpage{1061}(据审稿人本人所表达的意见)之在先性的单纯\textbf{假定}\fnote{784}{(*)}{(6月4日)我在这里甚至撇开了一点:这个假定根本没有根据。而且,转达审稿人答复的Remmert的信(1981年1月26日),对(信中援引的)Guillemin 讨论班与Kashiwara报告的日期,也只字未提。我刚刚火速打电话追到意大利的Mebkhout,问他这条指涉及其日期的细节。我得知:Kashiwara的报告在1978年——就在Mebkhout把自己论文的第三章告知他(1978年1月)之后几个月——Kashiwara先生可没有浪费时间!由于论文答辩直到1979年2月才举行(这要归功于巴黎各大学论文委员会——那个 J.L.~Verdier珍爱有加的机构……——的拖拉),这本来可以为审稿人的在先性“假定”提供一个说得过去的依据——至少就“重构定理”而言。可倘若那位审稿人(姑且不说他老实不老实——而这一点,明摆着已经谈不上)尽心尽职地干了他的活,他本会注意到:在Kashiwara那份所引的报告里,没有任何东西像一个“重构定理”的\textbf{证明}。\par Mebkhout后来还再度出击,在1981年3月25日的一封信里,他强调:1\textsuperscript{o})审稿人所援引的定理是“他博士论文最重要的成果之一”,而且他已把这一结果连同其证明告知了Kashiwara(可他忘了说是\textbf{何时}——Zoghman从来不肯说这个!);2\textsuperscript{o})这条定理“对于确立所论的范畴等价,是远远不够的”。R.~Remmert没有屈尊回复这封信——它来自一个无名无援的控告人。\par Zoghman方才向我说明(我这样追问下去,终会查明一切的……):他是第二年、即1979年——他答辩的那一年——才在 Guillemin 讨论班上得悉Kashiwara那场骗局的。因此,那正是他头一回同“黑手党”圈子里通行的这类手法打照面。等到这年9月的莱苏什讨论会,他对大明星Kashiwara的底细,心里已经有数了。不过,既然他的哲学与结果连同证明都已白纸黑字地写出并刊行,他原不会想到:一旦其重要性得到承认,他的著作竟会被干脆利落地整个掉包。而他的进路之力量的第一个征兆,恰恰就在莱苏什讨论会上、关于Kashiwara--Kawai定理一事中显现。\par 当然,1978年1月,Mebkhout(那时他还没有任何理由起疑)向Kashiwara谈起的,不只是他所谓的“双对偶定理”(后来为一场骗局之需被重新施洗为“重构定理”),还有完整的好上帝定理——前者总而言之是它的“一半”(两半中较浅的那“一半”)。他对我说:对于双对偶定理,Kashiwara倒确实“接住了”,看来他此前已自己提出过这类问题;但显然,他全然不知如何证明它。(然而,Mebkhout的证明并不使用奇点消解。)至于好上帝定理,那可是彻底地从他头顶上飘过去了——以至于到莱苏什讨论会时,他把这事忘得干干净净。然而Mebkhout像对所有人一样,同年(1979年)年初就把论文全文寄给了他(也就是说,那时他还不曾察觉前一年 Guillemin 讨论班上的那场骗局)。还有一件事表明好上帝定理彻底溜出了这位卡伊德的掌握:就在同一份 Guillemin 讨论班报告里,他甚至没有想到——为了良心安稳,不妨这么说——把它也一并收进腰包(哪怕他根本不明白那是怎么回事……)。\par 我至今无缘把Kashiwara这份报告(*)拿到手里,于是不免自问:它会不会让一个不知情的读者得到这样的印象——Mebkhout发展的那门哲学,Kashiwara(而且靠他自己的手段,如他所说)至少从1978年起就已知晓。Zoghman答应给我寄一份这份报告的拷贝;他向我担保,它会让我打消这种误会。那里头(他说)堆砌着一批技术性的陈述,或多或少(不)可理解(Kashiwara不可能做得更少了……),没有证明,也没有一条显见的引线,更没有任何东西(同他同年赫尔辛基的报告、以及次年莱苏什讨论会的报告一样)像一个联结连续系数与离散系数的“系数哲学”。\par(x)(6月16日)Mebkhout向我说明:那份报告实际上是由\textbf{Kawai}宣读的,作为与Kashiwara的一项合作成果。},就认定那位假定的后作者(偏偏碰巧正是那位无名者……)可以全然不被引用,这已经够骇人听闻的了;这类做法,明摆着为最严重的滥用敞开大门(而且早已敞开多时……)\fnote{785}{(*)}{这与 R.~Hotta三年后以同样犬儒口吻表示的态度(见上文所引他对Mebkhout的答复)如出一辙:新的“规矩”,或者更恰当地说“圈子的法则”,就是引用身居权位者(哪怕名不副实),不引用无名之辈——哪怕后者的贡献是决定性的、并有无可抵赖的出版物为证。\par 我并不怀疑 R.~Remmert在此事上的善意。但我还是要指出:身为 Inventiones 的编辑,他在这场骗局中的责任无可推卸、是直接牵涉在内的——即便撇开下述事实不谈(这一点他当时本无从料想):他是被一位不诚实的审稿人引入了歧途。那位审稿人曾表示过一份“希望”(就当时的情形而言,这希望堪称犬儒):“出于礼貌,Brylinski与Kashiwara会提及Mebkhout的结果”。而身为编辑,R.~Remmert的职责恰恰在于:确保Mebkhout的结果在正文中得到应有的提及——不是作为“礼貌”,而是出于对数学家这一行当基本伦理规范的尊重。\par(5月30日)自写下这些文字以来,我又得知了一桩“新事实”:它以意想不到的方式照亮了 R.~Remmert在围绕Zoghman~Mebkhout著作那场骗局中所扮演的角色,表明他也积极参预了围绕我本人著作的那场骗局。于是我对他所抱的那份善意推定便就此消散(那是出于老习惯,也因为此前没有无可辩驳的反证)。感兴趣的读者可以在下文(注“黑手党”的)(c1)部分、在“记忆的失灵——或新编历史”这个标题下,找到关于这桩“新事实”的详情。}。可还有更甚的。他所引的那条“重构 % ---------------- p.811 ----------------
定理”(它\dmark{}同样是从Mebkhout的论文里劫掠来的\fnote{786}{(*)}{关于这一劫掠,见“五张照片(晶体与D-模)”一注(n\textsuperscript{o}~171(ix)),(b)部分(“好上帝公式”)末尾,第~1005~页。},在那里它顶着\origpage{1062}(不确切的)“双对偶定理”之名)离Brylinski-Kashiwara那篇问罪文章的证明中所用的范畴等价(号称“Riemann--Hilbert”等价)还远得很——那等价出自Mebkhout一人,而那条定理对它丝毫无所蕴涵\fnote{787}{(**)}{Mebkhout的“双对偶定理”与“好上帝”定理——前者构成后者的“一半”,是两半中较浅的那一半——两者之间的关系,见已引的那条注(同样在(b)部分)。前者无须诉诸消解,而完整的定理则动用Hironaka奇点消解的全部力量(那是一件典型的“几何”工具,日本学派至少直到八十年代初还对它一无所知)。}。在我这里,审稿人的不老实——他倚仗上同调建制的一致默契,去封杀一个无名之辈的名字与著作,以成全名流们的“好处”——不容丝毫怀疑。凡略具上同调-分析素养、且对这一迷人主题有起码兴趣的人,都能亲自核实事实,并且看出这是一场粗劣的戏法——而那位匿名审稿人竟趋奉其间\fnote{788}{(***)}{(5月30日)而 R.~Remmert,以 Inventiones 编辑的身份,对它作了毫无保留的襄助……}。这局面的暧昧余地就更小了:因为直到1981年,在Kashiwara那里、也在任何其他日本或其他微分系统的行家那里,“导出范畴”这个词还无人说出口\fnote{789}{(****)}{(5月25日)Mebkhout提醒我,这句斩钉截铁的话应当稍作限定。尽管导出范畴自我离去之后在法国实同禁忌,日本学派却一直对它作节制的使用。那是一件方便的技术手段(尤其为免于诉诸谱序列),却绝非一套为形形色色簇与空间的上同调中“系数”的内在几何视野“量身裁制”的语言。},更谈不上有\dmark{}丝毫反思朝着一种联结离散系数与连续系数的“哲学”的方向——\origpage{1063}而老实说,这门哲学在后来的那些含混其词的指涉里同样无影无踪——什么“(原文如此)完整系统与(原文如此)可构造层(再来一个原文如此)之间的对应”——这些翩翩君子,直到今天,也没有一个人拿出这份诚实:\textbf{把所涉的范畴、以及确立其等价的那些箭头,白纸黑字地(像我不久前所做的那样)明写出来}。相反,自1977年以来,Mebkhout的一系列讨论班报告、笔记与文章,为他自1972年起在全然孤独中完成的开拓工作作证\fnote{790}{(*)}{这些文章的清单,我在这里既不评述、也不开列,请看已引的Mebkhout与 Lê Dung Trang 的文章(载 Proceedings of Symposia in Pure Mathematics, 40 (1983), part 2)。(5月25日)另见“三座路标——或天真”一注(n\textsuperscript{o}~171(x))中随文给出的书目。}。

\noindent 我必须承认:在直面这件事、并久久地、从方方面面把它端详与检视之前\fnote{791}{(**)}{(6月1日)我先是在去年,5月2日至9日那一周(写“送葬行列VII”——题为“讨论会——或Mebkhout层与反常性”)做过这件事,而后近两个月以来,在写《登峰造极》时,又做了一遍,……},我连做梦也不曾想到:在学者的世界里,竟会发生这样一场恬不知耻的集体掠夺。而不得不对自己说:这场不义的蒙骗,主要是由我最亲近的两位昔日学生合手搬演上场的;更有甚者,那开场锣竟是由\textbf{我的著作的一个接棒人的出现}敲响的——那份著作,我曾满怀激情投入其中,把我所能奉献的最好的东西都放了进去\fnote{792}{(***)}{在把这一页(涂改得相当厉害)誊清时,我心中浮起一个念头:倘若我对这份著作的投入居然结出了(别的)这样的果实——出乎意料而不受欢迎的果实——那无疑是因为,在这份投入本身、在我当时的精神状态里,并不只有我在这里乐于强调的那个“我之最优”,同样还有“我之最劣”。这一点,《愚妄与更新》(《收获与播种》的第一部分)里固然已相当清楚地显明,可它也是这样一件事:种种强大的自我机制不断地逼着我把它忘掉!我开始意识到,这个“最劣”,去年的反思中只是匆匆一瞥;我并没有对它作过一次真正深入的检视,或作一次把它形形色色的面相真正细致地揭示出来的“巡礼”。这就是为什么我对它的认识仍然浮于表面,一如这份认识的作用(尤其在我同葬礼的关系中)仍然浮于表面。\par《收获与播种》这第四部分“四种运算”,主要是一份细致的清点工作——把与葬礼相连的种种原始事实重新清点一遍。可这份“后勤”式的活计,毕竟让我感到:对葬礼的一种更深入的理解,与其说将来自我近三个月来所做的这类工作,不如说将来自对《愚妄与更新》中那份工作的深化——也就是说:也来自对“我离去之前”那些遥远时日里的那个我的认识的深化……}。我离去之后,这份著作竟成了那些人\dmark{}觊觎的靶与猎物\origpage{1064}——而他们恰恰曾是与我最亲近的人;还有一种隐秘的暴力,越过了我本人与我的著作,击中的还有那些公开从中汲取灵感的人……

% ---------------- p.812 ----------------
\subsubsection*{(c1) 记忆的失灵——或新编历史}

\tocsubsub{(c1) 记忆的失灵——或新编历史}

\noindent (5月30日)写下前几页六周之后,我在讲述我朋友Zoghman坎坷遭遇的叙述里,在此打开一个括号,稍稍停留在一个“新事实”上——它在先前一处脚注(第~1061~页的注~793[原书如此,当为注~785])里已有暗示。接下去的几页,读者可以读作一份有趣的补遗:关于“新风格”的开花结果——那门风格别处已谈过\fnote{793}{(*)}{见“恭贺——或新风格”一注,n\textsuperscript{o}~169$_9$。},它擅长(让所有人满意地……)为当代数学的(在此情形下的某一)主题写出一部“新编历史”。急着知道我朋友Zoghman后续遭遇(他迷失在一场他无从预料的马戏里)的读者,可以径直跳到“总彩排(封神之前)”(下文(d)部分,日期为4月16日)。

\noindent 我读到了 S.~Bosch、U.~Güntzer 与 R.~Remmert 的书《Non Archimedian Analysis》的引言与书目\fnote{794}{(**)}{即 Grundlehren der Mathematik 丛书第~261~种(1984年)。}。这本书阐述刚性解析空间的理论,并(理所当然地)把 J.~Tate1962年的(“私人”)笔记《Rigid-analytic spaces》呈现为这门理论的出发点。引言里说明:R.~Remmert“曾得以获得一份”这份稀世文献的拷贝;而这份文献,在\dmark{}“簇”概念的名流集会中,多多少少算是一位新来者的出生证(解析的“簇”,在此情形)。

\noindent \origpage{1065}Remmert想必忘了:把这份文献交由 IHES(那时刚刚起步)油印、并给他以及另外几位复解析空间的专家寄送拷贝的,正是我——为的是让他们注意他们偏爱的主题这一出人意料的扩展。那是在他们中间还没有一个人装作对实数域或复数域之外的基域感兴趣的时节——不过世事难料……

\noindent Remmert想必还忘了:我当年之所以那么起劲地在朋友们中间传布这份宣告一个几何“宇宙”之诞生的文献,(除别的之外)是因为我同这次诞生有密切的干系。“刚性解析空间”这个名字本身,就是我想出来的——那时不仅Remmert,乃至Tate本人,都不曾听说过这个名字,连这名字所应表达的东西,也无人梦想到过。我是头一个看出:Tate的椭圆曲线“loxodromique”理论,应当成为一种当时尚不存在的“解析”簇的“取商”,并且应当引出Serre“GAGA”型的代数-解析比较定理。还有另一个动机给我指明了通往这同一类新对象的道路:需要能够对离散赋值环上有限型的形式概形定义“泛纤维”。

\let\partdivider\mtOrigPartdivider

% ================================================================
%  18.5.4 注 (c1) 之尾 — (d) — (e) — (f) — c3 — c4 — Epilogue outre-tombe
%  对应原书 p.813–836（PDF 物理页与印刷页同）
%  源文: _work/chunks/chunk01–07.md
% ================================================================

% ---------------- p.813 ----------------
作为指向同一方向的第三个迹象:我曾听说,克拉斯纳(五六十年代的巴黎
数学圈里人尽皆知的一位怪才:他家里收容着一支猫大军,裹着一件厚重的俄式大衣,永远一副
乐呵呵的表情,在各讨论班之间走来走去……)——就是这样一位克拉斯纳,据说在非阿基米德
赋值域上“做解析延拓”。我所知仅此而已;我甚至不能确定可曾遇见过哪一位读过克拉斯纳这
方面工作的人——但这件事本身已够引人好奇。应当说,“解析延拓”这个词本身可没有让我心
跳加快的魔力(毋宁说相反,它叫我想起学生时代一些索然无味的记忆……);可一旦隐约察觉
到需要一类新的几何对象,它就难免让人心头一震……

回到Remmert——既然他的记性糟糕到这般地步,Tate的原文本(他还自夸握有一份)总该能替
他唤回记忆。在笔记里,Tate对我在这套理论的构想中\dmark\origpage{1066}所扮演的角色没有丝毫讳言\fnote{795}{(*)}{自Tate与我、他的家人与我的家人亲密往来的那些遥远日子以来,二十多年已经过去。多少年了,再没有他的只字音讯。这本书印出来,我的前学生与旧友当中没有人会看不到这篇引言里对我本人的这番抹除——可我也从未听说他们中哪一个为此动过声色。此一时,彼一时……},其中写道(凭记忆
引述):他“完全忠实地”遵循一位“总施工”(maître d'oeuvre——这是就通过“拼
合碎片”来构造该概念的方法而言),而这位总施工正是他从我这里得到的。此外我还为他提供
了某种“砌块”(或者说:在受限形式幂级数代数中进行“局部化”的“方法”),以应形式概形
纤维之需。他又以第二类“碎片”(或“方法”)补足这最初的一批,两者在某种意义上互为补充。

若不是我手里有“广义空间”这条主线(它后来成为\textbf{topos}),这个新概念恐怕根本不会问世
——平展上同调不会,晶体上同调不会,尾随其后涌现的许多东西也不会,甚至连那道最时兴的
“奶油甜点”、大名鼎鼎的 D-模也在内;这条主线的理论当时尚待建立,但四年来已被人隐隐预感。
正是这条直觉,为我指明了通往一类“簇”的道路——这类“簇”恰恰\textbf{越出了}寻常拓扑空间
(局部赋环空间)的框架。

自从John~Tate启动\textbf{刚性解析空间的局部理论}那一刻起,同样是我提出并推广
了关于这些新簇的首批关键“整体”定理的表述;这些表述早在第一批基础工作完成之前,就已
经在我心中成形:关于刚性解析空间上相对真概形的代数-解析比较定理;关于刚性解析空间的
真态射 $f$ 的 $R^i f_*$ 的有限性定理——这些问题由Kiehl在随后几年里解决\fnote{796}{(**)}{这里要指出,早在Tate打下刚性解析空间理论最初的地基时,我就清楚:他所处的那个框架仍是暂定的,丝毫不足以穷尽我试图用“刚性解析空间”这个名字表达的直觉内容——正如域上有限型概形穷尽不了“概形”一词所系的直觉。有一条线索指向Tate框架的实质扩充(凡愿意听的,我都讲过……),而提供它的是Tate自己:他在某个拓扑环上写出了一条“Tate泛椭圆曲线”(这个环是形式幂级数环 $\mathbb{Z}[[t]]$ 的一个子环,由那些在复平面开单位圆盘内收敛的级数组成——如果我记忆无误),而这个环显然应当被看作某一个刚性解析空间的“仿射坐标环”,其类型并不在Tate开出的清单之内。自我离去之后,一切基础问题都跌入人人不屑的境地;1962年Tate搭起的这台概念机器从那以后纹丝未动,也就不足为怪了。}\fnote{797}{(***)}{(6月4日)我还曾最先力主:对刚性解析空间,必须引入比Tate所设想的更一般的“点”(Tate的点只取基域的\textbf{有限}扩张为值)。这一必要性,既受代数几何之类比的提示,也出于一个愿望:为与所论刚性解析空间相伴的 topos 的“点”找到一个具体的解释。}。不过
话说回来,照眼下盛行的风气,\dmark\origpage{1067}预想新概念、物色总施工、提出供“真正的数学家”去解决
的问题,都被当作无关紧要之事——往极端里说,简直就是唬人的花招……

% ---------------- p.814 ----------------
不管怎样,这篇引言里没有一字提到我的名字与刚性解析空间有什么相干。克拉斯纳的名字同
样不见踪影——非但如此,Tate的理论还被说成引入了“一种足够丰富的结构,使不可能成为可
能:全不连通域上的解析延拓”——可早在1962年,这项(“不可能”的)解析延拓,十年也罢,
二十年三十年也罢(我说不准),早已是克拉斯纳“注册在案的正业”(若可以这样说)了。汗
牛充栋的参考文献里,克拉斯纳和我都无迹可寻。不过快到引言末尾,我的名字倒是顺带闪现了
一回,在“Grothendieck 拓扑”这个名字里;这个概念所引的是Artin1962年的笔记,
而对 SGA~4 里那份精益求精的打磨功夫(自1963年起,贯穿整个六十年代——却顶着一件显然
不受欢迎的父名……)则傲然无视,完全步了我那班前学生的后尘……;同样只字未提的,是当
年(1966年)我曾为刚性解析空间在晶体上同调的发展中指派的那个角色——那时Remmert(连同
他那些显赫的复分析家同事们)还全无迹象要理会这类所谓的“刚性解析”的怪“簇”(这名字
您给评评理……),它们是几位代数几何学家在自己角落里捣鼓出来的——仿佛复解析空间还
不够那些严肃的分析家与几何学家消遣似的……

\pmark\origpage{1068}关于书中所述这套理论的诞生史,只要是从第一手知情,就能看出这篇引言里招摇着同一种犬
儒做派——它同样流露在那份匿名审稿人对一位无名申诉者的答复里(并蒙这位Remmert先生祝福):显然,在作者们的心目中,这不过又是一桩“礼数”问题,说到底是一种“客
气”,施与不施全凭他们高兴:在新理论的孕育过程中起过关键作用的某某某之名,要不要写进
他们的“历史概述”(原文如此),悉听尊便。在他们看来(恐怕在数学界几乎所有建制中人看
来亦然——他们对此类歪曲照单全收,连眉头都不皱一下……),“历史”并不是\textbf{实际发生过
的事},而是一件可以由僭取书写之权者\textbf{乾纲独断}的东西,或者由一小撮人合议而定:什么有资格
存在,什么有资格算作发生过,全由他们说啦。

这班人就爱对苏联发生过、还在发生的事情津津乐道,而且一桩也不会放过(我知道自己在说
什么),忙着签署什么“捍卫自由”(思想自由之类……)的宣言——那是在\textbf{别人}家里;而在
\textbf{他们自己}掌权的地方,同样的谎言专制一天也没停过。

(6月3日)几天前在前几页里提到克拉斯纳那副独特而可亲的形象时,我心头冒出一个问题:
他还在世吗?他比我年长一两辈,我已有很久(少说十五年,也许二十年)没听人提起他的名字。
人物我还记得清清楚楚,他的名字却要过几秒钟才浮上来。(说实话,这一类事如今常在我身上
发生,年纪不饶人……)克拉斯纳以好客闻名,他的俄国血统本是另一个本可让我们接近的共同
点。可我那时一头扎在数学里,腾不出闲工夫只为“有趣”去交朋友。我们俩对待数学的路数,
想必正处在对跖点上。我们大概在Bourbaki讨论班的间隙闲聊过一两回,数学是肯定没有谈的。
那时真正让我着迷的,除了数学几乎没有别的……

% ---------------- p.815 ----------------
\pmark\origpage{1069}话虽如此,今天我收到Deligne的一张便条,不过寥寥数行,谈一桩无关紧要的实际琐事,
或许是想让自己还留在我的好记忆里(我们之间没有书信往来,算来已有几个月了);又或许是
想找个地方安插一条附言——我冒昧把它抄在这里(想必他会同意):

\begin{quote}
“又及:惊闻克拉斯纳已于半月前去世,不胜哀恸。我一直记得二十来年前他在布鲁塞尔所
做的一场滔滔不绝的长篇报告,我自然是云里雾里,却硬是撑到最后,成了所剩无几的听众之
一。令我吃惊的是,你那张五十年代人物表\fnote{798}{(*)}{这里对我在《收获与播种》第一部分(“Fatuité et Renouvellement”)中的说法有一个明显的误会。那一部分的主旨从来不是勾画一幅数学“五十年代群像”——哪怕只限巴黎圈子,或者Bourbaki周围那一圈。我的根本意图,是去发现我作为数学家的过去。正因如此,我才谈到我与这样那样同事或学生的关系——每当这层关系在我的生命中显得重要,或者能够照见我自己的时候。}里竟没有他——他在那个年代做了不少漂亮事,虽然与Bourbaki的精神格格不入,
而且有一种专把定义写得潦草马虎的天才。”
\end{quote}

这样一来,这里又多出一篇悼词,这一次是为我的一位同穴之人而作。这一篇里,我自信看出
一丝同情,或者毋宁说是某种情感的余光,那种情感往日曾经鲜活。可是,同在我的悼词里一样,
我的朋友皮埃尔决不会松一松牙关,说出——这一次是为了纪念一位一去不返的逝者——他津津
乐道却从不点名的那些“漂亮事”究竟\textbf{是什么}。然而他和我一样清楚:这些“事”为一个
如今正开得繁盛的理论的问世铺了路;而且,出于他或许心知肚明的原因,新主子们乐于抢先把
这位和善、潦草、“马虎”的先行者提前下葬(就葬在我的旁边)——他刚刚离世。他确凿无疑
是一位在超度量域上“做解析延拓”的人物——那时Tate、Remmert或者我还在“做”三角形全等
的种种情形和毕达哥拉斯定理,而我们的朋友皮埃尔还得让妈妈替他擤鼻涕(顺便擦洗……)呢!

% ---------------- (d) La Répétition Générale (avant Apothéose) — p.815–818 ----------------
\subsubsection*{(d) 总彩排(封神之前)}
\tocsubsub{(d) 总彩排(封神之前)}

(4月16日)可我必须回到我身后学生Zoghman~Mebkhout那一连串够味儿的“倒霉遭遇”上来。
\dmark\origpage{1070}1979年6月,当一个籍籍无名的陌生人——口口声声以Grothendieck的思想为宗——亲口告诉他,
那个让Deligne十年前整整一年绞尽脑汁、终未得到满意解答的关键问题\fnote{799}{(*)}{(5月25日)也可能,Deligne早就失去了对“关键”二字的分寸。关于这一点,见注“……以及羁绊”(n\textsuperscript{o}~171(viii)$'$)。}有了一个漂亮的解法时,Deligne的脑子里究竟发生了什么,我毫无概
念。以他一贯的心性,可想而知,他绝不会去祝贺这个年轻人在他Deligne栽过跟头的地方得了
手。可我深深感到,他那掘墓人的心性使他的嗅觉(我曾领教过的惊人嗅觉)失灵到这般地步,
以至于直到今天(六年之后),连他自己也没能把握住那个无名之辈的思想与眼光的真正分量。
他和所有人一样,最终只看见了那道“奶油甜点”:那件人人翘首以待的意外工具,那把砸开
“出了名的难题”的铁钎。可他也曾有一天,把别人传递给他的一整套广阔眼光据为己有——为
的是把那眼光连同生下它的人一并埋葬,再顺手抢走又一件工具,把它也改造成一柄“铁钎”……

我所知的Deligne对Mebkhout定理的最初一次反应,是一封未注明日期的短信(手写),写给Mebkhout,收到于1980年10月10日\fnote{800}{(**)}{这就是那份“以保密为条件提供、我在此不再多置一词……”的文件,“受害者”一注(打字稿第309页)里提到过它。事隔一年,Zoghman慨然允准我在此照录。}:

\begin{quote}
亲爱的Mebkhout:

我已把我那本你的论文寄给Bernstein和Beilinson:他们证明 Kazhdan--Lusztig 猜想用得
着你的结果(我有一份他们工作的俄文摘要,你要的话我可以转寄)。能不能再给我寄
一本?

谢谢。

P. Deligne
\end{quote}

从这封信我推想,Deligne必定已把“好上帝定理”的事告知了那两位苏联数学家,或许还暗示它
可以用来证明那个猜想;要么这是他自己悟出来的,要么当时已有风声,说Brylinski在这个
题目上有想法。Mebkhout那场“点着了”Brylinski的报告,早在1980年1月。Brylinski-Kashiwara
一边,Beilinson--Bernstein一边,两篇文章——用那个没有名字的定理(它的作者更没有名字)证
明了那条著名猜想——\dmark\origpage{1071}收到日期一为1980年12月19日,一为1980年12月8日,前后只差十一天。
巧合?

我甚至冒出过一个念头:Deligne1979年6月就比谁都先知道这件新工具(反正谁也不曾费神去
读那个无名之辈的厚砖头,Deligne也没有)——他为什么不自己想到把它用在那个猜想上,收
获一顶新桂冠,反倒去帮苏联同行摘桂冠?论头脑的敏捷,他不见得输给Brylinski吧?说不定
从那一刻起,他就已经瞥见一种可能:从旁门把“好上帝定理”本身的发明权捞回来——这发明权
(他心里必是这么感觉的)十年前就该是他的;说不定在他看来,正是由于一场不可饶恕的发牌
错误,那个不自量、欠打磨的狂妄后生才僭取了权利,去证明他Deligne早已久久钻研而终无决定
性成果的东西。他最终离成功只差那么一丝丝;换了别人在他流汗无功的地方收割,这可不公道
……但若想收回那件归根结底“理所应当”归于他的东西(依照某个人上人圈子里终于通行起来
的不成文法——他自认是那个圈子的中心和大佬……),那他就得以全然不同的细腻手法来操盘,
不能妄图一口吞下太多\fnote{801}{(*)}{这当然只是揣测:关于那条著名“对应”的据为己有之念,是否早在Deligne初闻其事时就已存在。就我自己而言,我深信不疑。诚然,上引那封信看上去会给人相反的揣测。而在我看来,这仍是那种挑衅姿态的又一个信号:他Deligne根本不必有半点顾忌——反正对方不过是个无名之辈,孤身一人对全世界,\textbf{横竖动弹不得};而他Deligne尽可以“自甘出丑”,正如他尽可以借“反常层”(\emph{faisceaux pervers})这个挑衅性的命名,以一种象征性却刺眼的方式,宣告自己真实的心性。关于这一点,见注“反常”(n\textsuperscript{o}~76),以及(在心理情境十分相近、却不那么极端的场合)注“玩笑——或‘权复形’”(n\textsuperscript{o}~83)。}。

话虽如此,Zoghman吃过与Kashiwara、与Brylinski打交道的怪插曲的亏,已有戒心,他审慎行事,决定亲自向Beilinson、Bernstein两位先生通报那个定理——Deligne说他们用得着它——以防Deligne这么大的
人物在谈起定理时,忘了提一句谁是它谦卑的作者。时机正好:下个月,1980年11月24日至28
日,莫斯科有一场“广义函数及其在数学物理中的应用会议”。Mebkhout在会上报告了他的定理,
已刊行的\dmark\origpage{1072}题目是“高维的Riemann--Hilbert问题”;他特意与Beilinson、Bernstein当面交谈,把自己
结果的来龙去脉原原本本讲给他们听。

他来的正是时候。会后不过十天,两位作者就寄出了他们关于 Kazhdan--Lusztig 的工作,以《法
国科学院院报》快报的形式(第292卷,1981年1月5日,第I辑--15):“群论——$g$-模的局部
化”,Alexander~Beilinson与Joseph~Bernstein署名,\textbf{由Pierre~Deligne转交}。理所当然,他们的
手稿上只字未提Mebkhout之名——显然,Deligne把那个无名之辈忘了个一干二净,尽管他明明把
此人的论文转寄给了他们,其目的正是……?谁懂,谁就懂吧!Mebkhout费了九牛二虎之力,才
让Beilinson(“两人中更老实的那一个”,Mebkhout一本正经地向我担保)相信:他们书目里引的
Kashiwara--Kawai那篇文章里,唯独没有那件“构造”(此处替代那句说了八百遍的“对应”)——而这件
“构造”,他们同所有人一样,也只敢用暗语提及。(多半是Deligne,一边把那篇无名之辈的论文
转寄给他们——想要的结果清清楚楚就在里面\fnote{802}{(*)}{(4月17日)论文里至少有一个非常接近的结果,尽管Beilinson--Bernstein(以及Brylinski-Kashiwara)所用的那套表述在其中并未逐字出现。细节见同日一条脚注(打字稿第1047页注(**))。}——一边必定向他们递了话:若非要给个参考文献,引Kashiwara的某一篇文章或许更稳妥些,哪一篇倒在
其次,反正没人会去查证到底。)对于那位亲自到场的无名之人,他们倒也当面许诺:会想着这
事的,会在Kashiwara那一处把引用改正过来。抱歉抱歉——我朋友Zoghman的倒霉史,确实老是那一套!

在这几位杰出作者的呈文里——由Deligne转交(他的信我刚才才照录,写于仅仅一个月之前)
——\textbf{Mebkhout的名字只字不见}。顺便说,Kashiwara的名字也不见(我已经瞥见那儿支起了一只耳朵
……)。不过,在呈文的最后一部分(证明 Kazhdan--Lusztig 的那一段)里,倒有一处偷带的两重
引注,引一件“\textbf{见于}[4]、[5]\textbf{的构造}”\fnote{803}{(**)}{读者不妨细品“\textbf{见于}……的构造”这一说法的含糊其辞:这件“构造”(或“对应”、或“关系”……)究竟出自谁手,问题被完全悬置;而仅仅六个月之后,在那场著名的讨论会上(见注“变戏法者”,n\textsuperscript{o}~75$'$),这个问题将以人们熟知的精湛手法得到“解决”:从Beilinson--Bernstein-Deligne的文章里,读者将获悉,那电报式的引注[4][5](在两处地方——那两处,想必那件构造恰好(托天之福)“有所陈述”)纯属客套,而“对应”的光辉父亲,果然是大家猜想的那一位……不过,撇开我刚提到的这套魔术手法不谈,引用一个新定理——又深又难——却只写成“见于……的构造”,仿佛那不过是件寻常“构造”,碰巧被丢在那里,作者们又碰巧选中它,拿来支撑他们辉煌的证明——这本身就已是一桩骗局。我在这里认出的,正是“SGA~4$\frac{1}{2}$--SGA~5”那场手术的同一副心肠:那场手术的要领,是(顺带)提一句“平展上同调形式‘有所陈述’于 SGA~4 与 SGA~5”(以及那位光辉作者不得不从中剔除的“废话矿渣”),随后挽起袖子,做作地开始“干\textbf{正经}数学……”。(5月25日)关于这种“新风格”,另见注“恭贺——或新风格”(n\textsuperscript{o}~169$_9$)。},\dmark\origpage{1073}这件“构造”(您准猜着了!)不是别
的,正是那位当值无名之辈那个从未被命名、其人更不曾被叫出名字的函子。参考文献[4]是Kashiwara的
一篇文章(临时的替身爸爸)。那篇文章里(Kawai-Kashiwara那篇也一样——它被扔进了盈亏账)自
然没有任何东西,哪怕远一点儿,像这几位作者所说的“构造”;况且那篇文章是1975年的\fnote{804}{(*)}{核实过了:就是前文引过的Kashiwara那篇文章,他在其中证明了他的可构造性定理——这个定理当然对定义那些“好上帝函子”起着关键作用(不过在1980年那阵抢购潮之前,除了Mebkhout,谁也不曾梦想过这些函子)。佯装把Kashiwara这个定理(无人打算跟他争)与“好上帝定理”混为一谈——后者无比深刻,层次完全不同——这是一桩粗劣的骗局。就证明而论,“好上帝定理”动用了Heisuke~Hironaka式奇点消解的全部威力。就“哲学”而论——这还要重要得多——它在同调形式主义内部架起了先前缺失的桥梁,沟通拓扑、代数与分析(算术尚待来日——前提是,我眼看着正在掘墓的某些先生,总有一天能找回他们健全的心智……)。},比莱苏什讨论会上某位无名之辈的报告早了差不多五年——正是那场报告让这位
Kashiwara先生开了窍:把“导出范畴”这个词说出口,毕竟不算傻事,而且可以依仗“强者即有理”这条再简单不过的法
则,把别人流汗挣来的功劳据为己有。至于参考文献[5],那是Mebkhout1979年9月在莱苏什讨
论会上的报告——恰恰就是在那场报告上,Kashiwara先生学到了:导出范畴原来是有用的——而且
除了用来坑害一位被他的大佬和前辈们晾在一边的无名之辈,还能派上别的用场……

\pmark\origpage{1074}同Brylinski-Kashiwara那篇文章一样,这里也没有任何东西,能让一个真正“不入局”的读
者起半点疑心:若不是那件新兴而天赐的工具现世——它被“见于……的构造”这句托词悄悄
藏起——这篇光辉的呈文根本不会问世。我还认出了那套屡试不爽的“溺毙鱼儿”术\fnote{805}{(*)}{关于这种号称“同化稀释法”的招数的其他例子,见子注“\textbf{正经}数学……”(n\textsuperscript{o}~169$_5$),即打字稿第885页脚注(*)。},号称“稀释法”:把要抹掉的那个人(却又想
给自己留“叫停”的余地,好在必要时说自己引过他……)\textbf{配对}到另一个与此事毫不相干、
或作用微不足道的人身上——仿佛在此处(在字里行间,却又明白无误地)说:我们把那个无名
小卒放在这儿(纯粹出于客套,也架不住他一再坚持),他和那件大名鼎鼎的“构造”(按新近确
立的共识,只能用暗语谈它,把它当作人人皆知的东西……)的干系,就跟1975年那篇文章的干
系一样多——那一年,“上流世界”还没有谁肯屈尊说出“导出范畴”这个词(顶多当笑话讲讲
……)。

\subsubsection*{(e) 冤大头买卖——或木偶戏}
\tocsubsub{(e) 冤大头买卖——或木偶戏}

我不后悔花这番功夫——既为我自己的备忘,也为哪位对这桩事感兴趣的数学读者——在这里
逐一清点围绕当值无名之辈那条定理的三桩前奏骗局。这三桩骗局,一桩是Kashiwara家的土产,一桩是
Brylinski-Kashiwara的土产(还劳了一位始终匿名的审稿人从旁襄助),一桩是Beilinson--Bernstein
的土产——而Deligne,则在幕后\fnote{806}{(**)}{(6月5日)Deligne“在幕后”扮演的角色,至少在第三段插曲里是明摆着的;第二段插曲,也有种种强烈的迹象指向同一方向。但看起来,Kashiwara早在1978年就已经在为他自己“打响第一枪”(就围绕Mebkhout著作的那些骗局而言)——当时(看来)Deligne还什么都蒙在鼓里。关于这一点,见本注 c) 部分(“入门费——或有为青年”),即打字稿第1050页脚注(*)。}。三桩在风
格上如出一辙,齐整得惊人,无须我再多着墨。这种风格,在我对\emph{葬礼}\fnote{807}{(***)}{关于这种风格,特别可看前已引过的注“恭贺——或新风格”(n\textsuperscript{o}~169$_9$)的结尾。}所
作的漫长调查一路之上,我早已领教到倒尽胃口;而在我那位天资最出众的学生——正是那位Pierre~Deligne——1968年的一篇文章里,它已经以刺眼的方式预先现出了形\fnote{808}{(****)}{见注“驱逐”(n\textsuperscript{o}~63)的开头。}。而这层情形也足以\dmark\origpage{1075}提醒我的记性:当年面对Deligne以及另一些在我看来天资出众的人,我抱着一种模棱两可、曲意逢迎的态度,
对今日铺展在我眼前的这股败坏,我也不是没有添上我自己的一份。

现在也清楚了:1981年6月那场“反常层会议”(Colloque~Pervers)的\emph{封神},离我们方
才清点的第三段插曲不过半年,绝非从天而降。说也奇怪,据我所知,这场讨论会是我离去之后头
一场、也是唯一一场专门(嘴上自然绝口不提,实际上却毫不含糊)用来把“Grothendieck数学”的
某个侧面重新发掘出土的讨论会——起因是一件猝然降临的新工具,谁料它竟无可替代。这件工
具,唯有在一种看待事物的路数里才活得下去——而时尚的种种“共识”,老早以前就把这种路数
归档为过时,还搭上几分隐隐的滑稽\fnote{809}{(*)}{至于这些“时尚共识”背后起作用的心理机制——它们与面对某种数学研究路数时某种“出自本能的”排斥反应彼此叠合——见前已引注“天赐的机缘——或封神”与“拒认——或召回”(n\textsuperscript{o}~151、152)。}。而仗着我那位天资夺
目的前学生的特殊天才,事情起了一场古怪的回旋:这条被他本人、被所有人弃绝的路数,竟在需
求的推动之下于事实中得到辉煌的证实;而恰恰又是借着同一场讨论会,这场证实成了机缘,把那
位已故而未被提名的大师——连同那位身后的学生(同样无人提名,他曾有幸(或不幸……)让这一大帮上流人物为他挪步前来)——彻底而最终地下葬。

% ---------------- p.819 ----------------
这场讨论会,可不是从天上掉下来的,决不是。我的朋友Pierre~Deligne有个特点(别的且不
论):他懂得等待,也懂得抓住有利的时机。围绕那道“奶油甜点”的三段插曲,加上对当值无名
之辈的提及早已被删削得几乎一干二净,分明向他表明:时机已经熟透,该神不知鬼不觉地、带着他那副
招牌式的笑容可掬与温文尔雅,把那件横竖都理应归他名下的东西收进兜里了。我猜他跟Verdier做过一番周密的串通:得让Verdier明白,如今时候已到,该锣鼓喧天把导出范畴和那份早已被
拒认的“父亲名分”从土里重新刨出来;与此同时,又把那个模糊的无名之辈、连同那位谢世多年的大
师,一并埋进舞台脚灯的光焰底下(就怕哪位不识相,偏偏记起:那些俨然“最新时尚”的美好
玩意儿,那位大师当年原也出过一份大力……)。

\pmark\origpage{1076}让Kashiwara给某条从未被点名的“好上帝定理”当个偷偷摸摸的爹,这一套一度还对付得过去
;而且恰恰也只在这样一个期间对付得过去:大家讲好了,那个定理既不会被点
名,也不会见诸文字。说到底,Kashiwara本人对这条定理也未必真有多么上心——这条定理,说不定他懂
的比Verdier本人还少;八成是机会凑手、积习相帮,他顺手就把它
揣进了腰包,仿佛出于无心。至于Deligne,那位懂得等待的人,他清楚得很:这条定理不会永远是
一条既无地址、也无姓名的定理。总而言之一句话:这是一条\textbf{在给自己寻找父亲的}定理,
要找还得找\textbf{配得上它的}父亲;而只有等到那“真正的”父名——按常理本就该属于他的那
一份(而且十二年前就该……)——成为举世公认、不容置疑的共识,它才够资格到光天化日之下
露面。那篇“反常”文章,同名讨论会的一颗明珠,正是朝这个方向立下的第一块界石,由头号当
事人以他惯常的娴熟手法亲手砌下。

我的感觉是:Beilinson和Bernstein多半还挺受用——想不到竟平白沾上了对那些层的父亲名分
(那些层也被称作“反常”——同样是叫错了的),而且绑在一起的还有一位至今这般显赫的黑帮
老大——可实际上,他们是被Deligne当了枪使,好在“万一……”之际给他充当“不在场证明”。
文章既然是这样写的,任何一位消息不够灵通的读者都只能认定:那件天赐“对应”的作者,不是
别人,当然是Deligne——尽管它从头到尾没被点过名,也没被明白说出过(反正人人都该早已知道
它……)。

剩下的,只有那一点暧昧的暗影(精心算计过的),藏在一句妙笔里:那条未被点名的“关系”,
“本该在这几篇笔记里找到自己的位置……”\fnote{810}{(*)}{关于这一点,见前已引注“变戏法者”(n\textsuperscript{o}~75$'$)。}。这正是“叫停”式的笔法:含蓄然而清楚地点拨一句——却又不
肯一字一句挑明——上述那条“关系”(既然通篇并无相反的说明),至少要归到那篇光辉文章的
\textbf{三位作者之一}名下,或者(退一万步讲)由三人共享。但同样明白的是:等到时机一到
(对那位懂得等待的人而言……),无论Beilinson还是Bernstein,都不会去向Deligne争一顶实际上已
然到手的父亲名分。想必是做成过一笔交易\fnote{811}{(**)}{会想到这样一桩“交易”,是联想到了两个类似的局面。其一是Deligne与Verdier之间的交易(或许是心照不宣的,但痕迹显然可见):后者“牺牲”掉 Lefschetz--Verdier 公式,把它划进盈亏账,以成全“SGA~4$\frac{1}{2}$--SGA~5”那场手术;作为交换,他却“收进”了过世者全部的“对偶”遗产,还外带导出范畴(一篇打折的文章)作添头。(详细的来龙去脉,见“公式”那组子注,n\textsuperscript{o}~169$_5$--169$_9$。)其二是Deligne跟一位被宣告已故的大师做成的那桩“交易”——反正此人早已从流通中销声匿迹,不愁他起而反应——事涉1967/69两个年头里合开的 SGA~7 讨论班:三年之后,这个讨论班被“对半分”掉,一半归那位已故者,另一半归Deligne和一位临时拉来的搭伙人。(详情例如见注“步步升级的插曲”(n\textsuperscript{o}~169~(iii)),插曲2.)这又让人联想到跟同一位已故者(他还蒙在鼓里)就那条(据MacPherson(Mac~Pherson)说)称作“Deligne--Grothendieck”的猜想做成的“交易”(见同一条前引注的插曲1):头一半归“邮差”Deligne——是他把一条猜想(此前一直由我那班上同调学生悉心保密)通风报信给了MacPherson;后一半归那位已故者,身份则是前者的“合作者”……}——\dmark\origpage{1077}即便不曾形诸明文,至少也彼此心照:Beilinson和Bernstein分得 Kazhdan--Lusztig 猜想,外加(为
了凑足分量——反正Brylinski-Kashiwara已经
% ---------------- p.820 ----------------
趴在上面了)对那些层的共同父亲名分——那些层,我猜想是出于众人的合议,才被叫作“反常”
的\fnote{812}{(*)}{见注“反常”(n\textsuperscript{o}~76)。};至于Deligne,则分得那条没有名字的著名“关
系”,只待那指日可待的一天到来——届时他的谦逊无须劳驾挪动一步——全世界都会管它叫“Deligne定理”。而这位未来的“父亲”鼻子也着实够灵,至少把这一点嗅了个分明:就这个孩子而论
(当初他宁可拒之门外,也不肯亲自动手把他生下来……),自己做成的是一笔“好买卖”
\fnote{813}{(**)}{这笔“好买卖”,我看同时又是一桩糟透的买卖;而且即便(甚至尤其……)在事事都遂了当事人心愿的情形下,也是如此:把宝贵的天赋与创造力,挥霍在扮黑帮的营生上。}。至于
Kashiwara,他的戏份儿已经演完:在那篇光辉文章里,就那件天赐的“关系”而论,再没有人提起他,正
如再没有人提起那位当值无名之辈。人多势众围攻一个,只要那是个无名小卒,好说;可地盘上一
旦清走了闯入者,那就各人顾各人……

% ---- 以下“(f)”节, 仍属原书 p.820 ----
\subsubsection*{(f) 演员列队——或黑手党}
\tocsubsub{(f) 演员列队——或黑手党}

那份才翻开三个星期的“家庭相册”\fnote{814}{(***)}{见3月22日那条同名注(n\textsuperscript{o}~173)。},想不到又添了几
张新面孔。看得出来,“这个家”丁口大为兴旺,我这副老朽之躯都快认不过来了,更何况时世也
已大不相同。这一回,按\dmark\origpage{1078}出场的先后,是\textbf{Kashiwara}、\textbf{Hotta}
\fnote{815}{(*)}{细心的读者也许会奇怪,在这份“演员列队”(在围绕Zoghman~Mebkhout著作的那场骗局兼蒙骗之中)里,怎么见不到Kawai的名字——他是Kashiwara那篇屡被引用的文章的合作者,其第4节肆无忌惮地劫掠了Mebkhout博士论文的第三章。(关于这一点,见注“五张照片(晶体与D-模)”(n\textsuperscript{o}~171(ix)),特别是第1005页。)Mebkhout一口咬定,不能把Kawai同Kashiwara装进“同一个口袋”(Kawai嘛,闭着眼睛跟着Kashiwara走也就是了……)。他向我形容Kawai是个有点儿跟不上趟的家伙,我有个印象,他几乎已经喜欢上这个Kawai了——总之一句话,这可是他的“好日本人”,在他那里,我去碰是万万使不得的!大概也正因如此,他才没有给Kawai写信(就像他给Kashiwara的另一位搭档Hotta写过信那样),去点破他与Kashiwara合写文章里的那些骗局,从而逼他公开表态,同他的搭档兼老板团结一致。}、
\textbf{Brylinski},还有Brylinski-Kashiwara投给《数学创造》的那篇文章的
\textbf{匿名审稿人}。一帮“硬茬儿”,这准保没错,反应招数训练有素;更何况只要事关坑害一
个无名路人,只消幕后大头目一个隐而不显的暗号,便个个心领神会、行动一致(甚至等不及暗号
……)。

我又一次认出了那副\textbf{黑手党}的做派\fnote{816}{(**)}{这种离奇的感觉,去年在注“讨论会”(n\textsuperscript{o}~75$'$)(猜得出是哪一场……)里就已攫住过我:那里一派收保护费的气氛,叫人恍然如梦,或者以为正“在看一部讲黑手党称霸某座遥远大都会阴暗底层的电影……”。这一回,在这趟穿行于当值无名之辈种种倒霉遭遇的漫游里,这种感觉又寸步不离地陪伴着我……}:在他们无可争
议的领地上称王称霸,而领地的核心,便是代数簇以及诸如此类的上同调理论。一帮才华耀眼而心
狠手辣的人物,脑瓜子个个无懈可击;在所谓“当值无名之辈”那场行动接连四段的插曲里,我一
路重新见识了他们的身手,直到“反常层会议”登上顶峰。除我方才点到的四位黑帮老大(其中一
位匿名)之外,还请诸君记起“硬核”的另外五名成员;这样一共九位,为埋葬\textbf{那位闯入
者,那个不是他们一伙的人}而全体出动。

有大头目\textbf{Pierre~Deligne}——那位永远懂得“下水”下得最少、进账却捞得最多的主儿。
有他的副手\textbf{Jean-Louis~Verdier},人称“大善人”——正是此公,主持过某位无名之辈某
篇论文的答辩委员会;又正是此公,当过那场令人难忘的讨论会的两位组织者之一,那场会肆无忌
惮地掠夺了同一位无名之辈。另一位主其事的是\textbf{Teissier},他同Verdier联名为那场令
人难忘的讨论会的令人难忘的文集,签下了那篇令人难忘的引言。跟旁人不同,他\dmark\origpage{1079}似乎不过
是跑了个龙套、当了个挂名替身,自己从中一无所得——唯一的乐趣,不过是讨得
他知道既有名望又毫无顾忌的那班人的欢心。最后\fnote{817}{(*)}{(5月25日)这个“最后”终究是言之过早——打那以后,又有别的帮派成员自行冒了头,进入了我的视线。关于这一点,见“光荣之日”注(n\textsuperscript{o}~171(iv))里打字稿第962页的脚注(*)。\par(5月30日)最新消息:又一名成员——R.~Remmert——刚刚被认了出来。关于这一点,见本注 c) 部分(“记忆的失灵——或新编历史”)。}还有\textbf{Beilinson}和\textbf{Bernstein}(对这两
位,就在这几页里
% ---------------- p.821 ----------------
我刚刚有过一番更充分的领教),被一根根看不见的提线轻轻地牵动着……

我还会等着,不急不躁,也不抱幻想,看未来还会给我们预备下哪些别的“反常层会议”,有整
个教团毫无保留的首肯作保,好为“科学”增最大的光荣,为“人类精神的荣誉”添彩。

\subsubsection*{c3.~根与孤独}
\tocsubsub{c3.~根与孤独}

\notehead{171$_{3}$}{}

(4月18日)一步一步跟完我朋友Zoghman那串倒霉遭遇的这第四天行将结束之际,我比去年更懂得了一些态度与心性——尤
其是冲着我来的那些——而直到去年,它们在我眼里还显得古怪。总而言之,凭着他那份自知分量的工作,他原以为进了“一个大家庭”,
有点像那位从不被人提起的已故宗师的家——的确,无人提起;可即便不提,他分明在场。谁承想,他
发现自己落进了一个鲨鱼的世界:个个彬彬有礼、甚至和蔼可亲,牙齿却毫不留情——转眼之间,他
带来的东西被剥夺得一干二净:八年孤独劳作的果实;然后有人让他明白:你,我们已经看够了——一个讨厌
的、不请自来的外人。换了他那个位置,没几个人能不留下创伤。我不知道他可曾对什么人敞开过他
的挫败,除了几声苦涩的暗示,又含混得反倒像在作证对他不利:一个怨气冲天的人,边缘上还有点不合群。

尽管没有人点我的名,我却俨然是这伙毫无顾忌、毫不留情的世界的“爸爸”;他实在没有什么理由
信任我。诚然,我们的初次见面——1980年,那时他还远在千里之外,猜不到等着他的是什么——曾
经打下过信任的底子;我分明感到,这个底子,顶着一切,一直保存到今天。说到底,他心里明白:我
纵然是一切鲨鱼的“爸爸”,也不会像他们那样行事。可\textbf{怨怼}是有的,这确定无疑;而且它乐于\dmark{}\origpage{1080}
摆出一副打心底不信任的架势——虽然(至少我感觉如此)那是“贴上去的”。

为自己认定的正当权益“斗争”,这容易——只要你属于一个团体,多小都行,同它声气相通。可那个
孤身对抗全世界的人、那个被排斥者、那个不请自来的外人,就像一棵被夺走了沃土的树。他身上的
力量帮不了他;那力量变成苦涩,反过来咬他自己,仿佛要同那个排斥他的全世界合唱一个调子。

当我手里捧起那本书——它为 Motif 掘墓复生,同时埋葬了让它们现形的工匠;那本书,署着一代
光辉人物(其中有我参与培养的)里最光彩的四位作者的大名——当我终于读到它(纯属偶然,因为
此前没有任何人觉得有什么值得特别提醒我的……),那一刻,在我结识数学家世界三十六年之后,
我第一次知道:\textbf{我孤身一人,对抗所有人}。过去八年里发生的许多事,忽然拼合起来,显出全部
意义。当一个人蓦地重新发现这份孤独,那滋味很是古怪。那一天我不得不缓一口气;此后一连几个
星期,一天一天地认清\emph{葬礼}的全部规模——一场与那部著作相称的\emph{葬礼}。

可这同Zoghman毫无共同之处——他的“自己人”在他还没真正扎下根之前,就把他“挂账”了。而命
运对我,是笑过的。多亏长者们接待了我(他们后来故去也罢、退休也罢、早就不再搞数学也罢,统
统不要紧)——多亏我年轻岁月里寻得的兄弟般的接待,我,得以在我自己选定的土壤里,“扎下根
来”。这些根,扎下去,长起来,随着岁月,变得深沉而有力。这些根,牢牢栽在一片不是“共识”、也
不是任何时尚的土壤里——扎得想必比任何以制造时尚、追随时尚为乐的人都深\fnote{818}{(*)}{如果说我在数学里也好、别处也好,从来不曾操心去追随或制造时尚,我知道,这正是我有幸在幼年便得以发展起来的那种深根的表现之一。既然我从一开始就在自身之中有了深根,投进我的大事业里的那份精力,就不会被补偿性的饥渴——比如领风气之先的饥渴,或保持并显得合乎必备之“调门”的饥渴——所分散。
关于我的童年和这些“根”(我想,我没有说出这个词),我在注“天真(阴阳的联姻)”(n\textsuperscript{o}~107)里说得具体。}。

\pmark\origpage{1081}说到底,我还“孤身对抗所有人”得起——想说什么就说什么,走我自己的路。

(5月25日)\fnote{819}{(*)}{下面两页,原本是作为注“……与意外之喜”(n\textsuperscript{o}~171(iii))的一条脚注写的。在哪里安插它们,我犹豫了一阵,最终决定收进眼下这条“根与孤独”里。因为这是“封神”全部注中唯一的例外:我试着从自己的经历出发,尽力去体会Zoghman本人如何经历了那些我为之作编年史的事件与处境。}不消多少想象力,就能体会Mebkhout的沮丧:他的核心结果之力一旦得到承认,他
就感到自己像一根麦秸似的被“扫掉”了\fnote{820}{(*)}{“扫掉”一语,借自Mebkhout的一封信(比正文所引那封早两天),相关段落我在这里照录:
“的确,[Kashiwara的]那个可构造性定理……让我爆发了。而且从那一刻起,像Deligne这样的人,一眨眼就
能找到我全部的结果,连‘好上帝定理’的各种形态在内,证明嘛,照你的说法,四个小菜的功夫。这就
说明了为什么这一切几天之内就被扫光了。”
我想,Mebkhout在这里把一个Deligne式的、默而不宣的“推理”,说得再确切不过:据他人\textbf{本来可
以}(而且\textbf{本应})凭他自己的手段、行囊与一切,“四个小菜的功夫”找到这些果实,所以就可以据为己有。这类推理——在类似的情境里,人们常常禁不住要做——唯一的别扭之处在于:\textbf{难的全是想
到}。而真正“想到”了的,是Mebkhout,绝不是Deligne,更不是别的什么人。创造不归“技术”管:技术,
在终于有人看见了无人曾看见的东西之后,顷刻间便可“扫光”一个局面,快过把它写下来。创造不在
“扫光”里,而在\textbf{看见}无人曾看见的东西之举里;用自己的眼睛看,不“追随”任何人。把这两者区分
开来——把创造之举,同摇一圈就转圆的那根曲柄区分开来——正是数学家操守的一部分。}。他给我写信(4月24日,在他新近来我家之后)说:“我花了八年,搭起
Kazhdan--Lusztig 证明所用的那些结果。他们用一个星期就把它证明了。”这一回,还是有一种羞
怯拉住了他,不肯把心里真真切切感到的那些话说到底;\dmark{}\origpage{1082}我在这里替他补上那句“没说出的”:而事情一旦做
成,“他们”便拿着那件别人在孤独里打造出来的崭新工具,彼此之间昂然招摇,同时让工匠明白:你,
我们已经看够了……

可这件事实在大得离谱,以至于当场,Zoghman还不大相信他自己那健全官能的见证——正如去年5
月2日,读到吕米尼讨论会文集时,我自己也迟迟难以相信自己官能的见证\fnote{821}{(*)}{关于这场讨论会(1981年6月),见注“不义——或归来的意义”或“光荣之日”(n\textsuperscript{o}s~75、171(iv))。说实话,去年5月头一个星期写下“送葬行列之七:讨论会——或Mebkhout层与反常”(n\textsuperscript{o}s~75--80),还不足以克服那近乎不可克服的惰性——在一个严格说来只有你一人使用这种见识的处境里,“相信自己健全官能的见证”何其难。直到五个月之后,终于“有血有肉地”面对了现实——我的朋友皮埃尔(Deligne)来到我隐居处看我——那股隐秘而顽强的怀疑,才终于消散。关于这一点,见注“义务已尽——或真话时刻”(n\textsuperscript{o}~163),尤其打字稿第782--784页。}。Zoghman是在去年1月读到这同一部
文集时——在 Kazhdan--Lusztig“总彩排”三年之后——才终于多少弄明白,实际发生的是什么事。

打击是可怕的,我自以为懂得——Zoghman当场以为他要把命搭上。幸好人这身子很结实——Zoghman
今天还活着,其间还结了婚,做了父亲……可我想,即便手捧“文集”的那一刻,他仍旧不能完全相
信。总有什么东西“卡着”。说不定,在我写这些行的时候,他仍旧没有全信。应当说,即便只按“理性”
或“客观”的字面讲,这件事也实在难以置信、实在巨大:直到今天,除了我(或许除了他,也未必
……),竟没有\textbf{一个人}敢相信自己眼睛、看见它——虽然它比一座主教堂还大!

可对一个被当头击中的人——被他的那些敬重的、什么都不缺的长辈,施以犬儒而\textbf{无缘无故}的
不义——这种事,只怕永远没法全信,只怕永远“\textbf{超出理解}”……而恰恰是这种事,能够毁掉一个
人的一生。它们之所以有这份\dmark{}\origpage{1083}毁灭之力,是因为一种幽暗的、绝望地被压下、却无可否认的知觉:
\textbf{毁掉}的\textbf{意图}——就这样无缘无故地,为了\textbf{取乐}:为了取乐于漫不经心的一挥手,碾碎那
对你珍贵的东西,碾碎(如果办得到)那构成你生命的实质与盐分的东西。正是这种“无缘无故”的
恶意里的反常快乐,才真正“超出理解”……

我想,Zoghman从来不曾真正对什么人谈过这些,大事之前也罢,之后也罢——至多几个单音节的
词,除了他自己,谁也解不开。单是 Kazhdan--Lusztig 那一幕,已经太大、太不像话,他不敢指望谁
会信。根深蒂固的共识,把最明显、最明白、最无可抵赖的事实,像麦秸一样扫掉。而这里说的,
是一件疼得那么切近、在他的存在里那么“一碰就疼”的事:只消想到,他若向谁袒露,那人可能拒斥这
份不速的信息,他面对“超出理解之物”的痛苦可能不被接纳——这个风险、这个概率,便膨胀成
\textbf{不堪忍受}之物;这种事,付任何代价也不去碰——哪怕为此当场憋死,要憋就憋……

两年前,他倒是对我“用单音节的词”说过。说不定,在他心底,他盼着我会听懂这些单音节——不只
按字面,而是同时听出他不敢高声说出的一切(也许连对他自己都不敢说……)。这盼望确实疯得
可以(在一个一切都疯得没边的处境里!);对我后来确凿得知的那些,我当年连想象的边都没沾
到。只能如此,谁叫我那时没有一份详尽的、原原本本的信息呢\fnote{822}{(*)}{(6月1日)更公平的说法或许是:以我这种有限的敞开与在场的状态,除非在极罕见的时刻,事情“不可能是别样”。然而我相信,我们人人都有“耳中之耳”,完全听得见没有说出的那部分——只不过,我们通常留神把这只耳朵捕到的讯息,排除在自觉注意的领域之外……}。而Zoghman那一边,也远不敢把这份信息
给我。这很疯;可这并没有妨碍他怨我。他总得怨一个人,怨一个足够亲近的、摸得着的人,好把“超
出理解之物”在他心里搅起的东西,至少卸下一部分在他身上,好让自己从那啃噬着他的东西里,解
脱那么一点点。

\subsubsection*{c4.~掠夺的空白授权——或“高尚行当”}
\tocsubsub{c4.~掠夺的空白授权——或“高尚行当”}

\notehead{171$_{4}$}{}
\dmark{}\origpage{1084}(6月2日)我给\emph{葬礼}画上“句号”、连同末尾那条注“自深处”(4月7日)一道的那份心满意足,到今天已满两个月——而这两个月来,我也一直在拼着老命,给\emph{葬礼}的最后一部分“做最后的润色”!这几乎是去年同一时节那出戏的原样重演——当时,我给那即将成为\emph{葬礼}第一部分的东西“收尾”,也是收个没完。
% ---------------- p.823 ----------------
和眼下一样,拖个没完的正是“最后一刻”——拖到我连吃喝都顾不上,最顾不上的还是睡觉。就这么一直拖下去,直到我的身体油尽灯枯,宣布弃权。那正是整整一年之前(差不了几天),我不得不撂下一切,足足三个多月,一门心思只顾从一场急性衰竭里爬出来\fnote{823}{(*)}{关于这一段,见注“事件——或身体与精神”(n\textsuperscript{o}~98)。}。可这一回我多了个心眼,格外当心,不再走同一条老路。这条命,我还想留着呢……

这一回,又是那场没完没了、一弹再弹的“调查”。我原先估计,撑死也就一条十来页的注,取名“四种运算”,把去年那阵旋风式调查的结果“归置归置”、做个总结。可如今,调查卷土重来、干得热火朝天,眼看要满四个月了;十页变成了三百页上下,而且这还不算(完全)完事!我再不敢作什么预言——从9月底重新动工算起,已经到第九个月了,我一直在“马上就要完工”!只有等到最后一批笔记全部打成清稿、通读修改完毕、交去复制的那一天,我才会知道它\textbf{真的}完成了。(到那之后,剩下的就不再是我的活了。)我唯一知道的,是我巴不得赶紧到那一步,就像巴不得看到一场漫长而磨人的疾病有个尽头;而我必须尽我所能走到终点,不被种种想象出来的期限推搡着走。只有等该\textbf{在此时}看见和说出的都已看见、都已说出,我才会停下来喘一口气。

害我吃了最多苦头的,正是那该死的“\emph{封神}”——我说不上为什么。这几样“四种运算”,是《收获与播种》全书里唯一一处写得磕磕绊绊、零敲碎打、一路费力挣扎的部分——按说本该是现成的活儿:不过是一次简单的“归置”,没错;没有任何东西会以触到\dmark{}\origpage{1085}要害的方式,把我这个人牵连进去、拿我这个人问罪,从而动员起抵抗的力量,一个“\textbf{摩擦}”。可天知道这摩擦有多少,而且在\emph{封神}这里,比其余一切加起来还多!它打哪儿来的?

早在“操演”那条注里,写起来就已经够费劲了。没完没了地拉长,就是从那里开始的。光这一道“运算”,最后就足足写了八十页——而如今,一个月之后,\emph{封神}已经写出了整整两倍的篇幅。话虽如此,除了“操演”里兴许有那么几页(边上带点儿太浓的“侦探”味儿……)(在那里,我兴许一头扎进一桩离谱透顶的“骗局”的絮叨细节里,扎得比必要的还深……)——除了这件写得眉目详尽、对“不在局中”的读者来说多半还有点磨人的“凭卷宗办案”的活儿之外,我并不觉得我最终码出来的那一摞摞上百页的东西是画蛇添足,是回锅的冷饭,是头发丝里切四瓣的功夫。让我一直提着气撑下来的,恰恰是\textbf{新的实质内容}大量地、出人意料地涌到我身上,它们我无论如何都得安置妥当,由不得我愿不愿意——其中甚至有,可不是嘛,数学的实质内容!有好些时刻我几乎招架不住:一时要落到白纸黑字上的东西实在太多,而且刻不容缓——这些事情一个个都还热乎着,甚至烫手,可你偏偏只能一件挨一件地去料理……

这样的丰富,本身就自带着一股强劲的工作推动力,丝毫不是那种会惹出“摩擦”的东西,恰恰相反。这摩擦,可以肯定,并非来自实质内容本身,而是来自我投进这份工作里的那股唯我的执念之力。说来像是悖论:恰恰是我那股“赶紧收场”、“一把甩到桌面上”的不耐——急着说出我该说的这样那样的事,而它们就发生在此时此刻,同我息息相关、触我至深——正是这种不耐(我想),制造了摩擦,制造了能量的涣散。摩擦是分裂的征兆,是几股力量朝相反的方向对拉,每一股都因对方的抵抗而愈激愈怒:
% ---------------- p.824 ----------------
一头是“赶紧收场”的急切——趁我还在打磨这块料,就想把它“甩手”交出去;另一头,是要求把眼下这一刻让我隐约望见的东西走到底,不满足于差不离,不让人推搡,也不把自己关进一个赶着“结项”的“计划”、一份预先定死的“日程”里。我很清楚:一旦把“意外”排除在外——这个防着我原地打转的搅局鬼——我的工作就失去了品质和意义。它\dmark{}\origpage{1086}就沦为“文牍差事”。这些年下来,我对这处看似不算什么、实则是一切的“小小差别”,已经变得非常敏感。这样的大转弯,如今偶尔还会启动,总在沉重已极的时刻——但从来没有维持多久。一旦走上这条路,那孩子把一切都扫到一边去——再想撑下去也是白搭。工作本身的兴致,那份\textbf{渴望}——它全然不同于攒页数的饿劲儿、不同于抢着画句号的急色——霎时间无影无踪,你便傻乎乎地守着涂黑纸张的活计。到那时,真就再没意思了——剩下的只有一件事:把航向拨正,而且马上就拨!

工作里总有一种\textbf{不耐}(我的一个老相识……),不停地把我往前拽。我觉得,这和自我跟这几样“四种运算”较上劲以来沉沉压在我身上的那种不耐,不是同一种。那另一种不耐,不是压下来的重量,而分明是一种往前拽的力量。它是胃口好的征兆,不是倦怠、疲乏或餍足的征兆。它不是累积的不耐,不是收尾的不耐,不是赶着把一份“计划”“结项”的不耐,而是要认识我面前那个即将向我把底交出的未知者的不耐。是一个赤身的孩子,独自站在无边的大海面前,急于纵身潜入其中去认识它……\fnote{824}{(*)}{这一形象,已见于注“孩子与海——或信与疑”(n\textsuperscript{o}~103)。}

不过,现在该回到我朋友Zoghman那串倒霉遭遇的叙述上来了——就在这条预定作为\emph{封神}最末一笔的注里。前面说过,这段叙述,Zoghman本人从来只是东一块西一块地零星递到我手里:随着信件、电话、会面,碰上什么算什么;可以肯定,我这边思考的推进和\emph{葬礼}的写作都因此受了影响,至少在专写我朋友那番坎坷的那一部分里是如此。如今我更懂得这种欲言又止的意味了,因为对“受害者”角色的一切执着(去年我还自以为从中看出几分)早已烟消云散(就算它当初确实在过)。我这边,想必也有过那么一些时刻,多少有些饱和了,态度上是那么一副“求您别再倒了,行行好!”的架势。这八成也没给他鼓劲儿。说实在的,那句“这儿是日本人、那儿是Kashiwara”的老调,Zoghman想必唱了总有四五年了,叫我不胜其烦;他跟他们之间确实领教过一些事,这不假。可我心里明白:他之所以领教过那些事,他的成果之所以就这样被交出去任人掠夺——近乎官方式地:“来吧,好人们,放量地拿,千万别客气……!”——那\textbf{不是}因为几个遥远\dmark{}\origpage{1087}的日本人。这\textbf{是因为“自己人”}:是“小家庭”\fnote{825}{(*)}{(6月16日)关于这一点,Mebkhout特意要我强调:他早已完全不再认同于这里所说的那个“小家庭”。}里的人——都是我们自家阵营里的体面人物,而他从不点名,除非引用他们的著作,并带着配得上他们崇高声望的全部敬意。

至于Kashiwara之流,我可再也不想听人提起了!眼看事情卡住了,Zoghman这时拿出他的智慧和耐心,索性撂开不谈,却并没有因此放下对我工作的关切,也没有停止在这里那里默默地、实实在在地帮我一把。

直到4月初他最后一次来我这儿,我才终于看了看那个“日本包裹”。起初还有点勉强。我原以为,自己准会陷在一堆解不开的超技术纠葛和读不懂的文章里腻烦得要死(没准儿还是日文的……),反正我是永远不会去读的——结果不然!事情简单得像打声招呼——有点像巴黎地铁(或者该说东京地铁)里“扒手的故事”。说真的,还挺有趣(至少,被摸走钱包的是别人的时候……)。
% ---------------- p.825 ----------------
这一来,Zoghman和我之间的僵局也松动了,他的倒霉史,这里那里、一道道闪电似的,也让我听到了些零头碎块。原先我多少按“技术情报卡片”的样式记下的几段插曲,如今被鲜活的临场追忆充实起来;恰恰是这一类东西,在科学文本那副纹丝不动的“立正”姿态里,甚至在同事之间的通信里,似乎永远被禁绝——您总不会想看吧!在“四种运算”里,我甚至不得不使劲摇晃自己,免得恰恰又落回那种腔调——“侦查结论”式的腔调(甚至,“诉苦清单”式的……)。Zoghman交出来的这些“碎块”,帮我从里面走了出来,帮我同一种活生生的实质保持接触。

Zoghman离我家的当天,我就重新拾起\emph{封神}的活儿,趁他讲给我听的东西还没凉,再写它一条两条子注。这就有了那几条注(或者说子注——写得多了,我自己也数不清了……):“眼光的绽放——或闯入者”、“黑手党”(后来我把它又分成七部分,每部分起了名字),还有“根与孤独”。我把整个儿火速寄给了他,请他在我送去打字之前提提意见。到这一步,我觉得自己多少有点在替他说话,我要确证:凡是我根据他讲的转述出来的东西,都得到他毫无保留的认可。他随即寄回了详细的意见(4月22日与24日的信)。

\pmark\origpage{1088}这些意见里,不乏这样的“碎块”,为我那些注里偶尔显得有点瘦骨嶙峋的事实骨架,敷上了鲜活的血肉。

也正是这样,我才知道Zoghman当时在场——就在那个令人难忘的1980年4月22日,Goulaouic--Schwartz讨论班上。就在那一天,Kashiwara把“好上帝定理”当成自家出产的定理当众宣布——而这条定理,是他几个月前在莱苏什讨论会上从Mebkhout亲口听来的\fnote{826}{(*)}{关于莱苏什讨论会与Goulaouic--Schwartz讨论班上的这段插曲,见注“黑手党”(n\textsuperscript{o}~171),部分(b)“最初的麻烦——或太平洋彼岸的黑帮老大”。}!这胆子实在大得没边,何况Mebkhout人就坐在底下,难怪听起来像天方夜谭。Mebkhout当场没有爆发(我直纳闷他是怎么忍住的……)。他客客气气地等报告讲完,“为的是公开抗议这种做派,向他提起莱苏什的那场会议、提起他的那个问题\fnote{827}{(**)}{指的是1979年9月莱苏什讨论会上,Mebkhout报告结束时Kashiwara所提的那个问题。关于此事,见上一条脚注中所引的那条注。}。Goulaouic请我私下解决我们之间的事。大厅几秒钟之内就空了”。

这便是那些“碎块”中的一块,由这段惜墨如金的叙述交了出来。后来我又在电话里听到些细节。这件事值得停下来细看。关于八十年代数学界的风气民情,它透露的东西可不少。这里涉及的,已不是某位心狠牙长的“黑帮老大”的心态——科学界传统价值败坏的一种极端症候——甚至也不是那班体面风光、处处周全的“建制中人”——在他们身上起作用的,是护着“自己人”的阶级本能。在这里,是整个大厅眨眼之间一空——转眼间一个人都不剩\fnote{828}{(***)}{这幅景象在我心头不可抗拒地勾起一连串联想,同三年前我亲历的一个场面如出一辙:那是一场Bourbaki讨论班的末尾,承人们好意赏我十分钟,让我谈谈某条坑害外国人的凶恶法令。关于此事,见“我的告别,或:外国人”一节,n\textsuperscript{o}~24。}!你们自己私了吧——我们嘛,什么也不想打听……

我真想知道,Goulaouic和那场讨论班上其他安安稳稳的听众,脑子里当时到底转过了什么念头——在那儿做报告的,可是一位声望斐然的外国讲演人(讲的题目,我相信他们当中没有谁真正熟悉)。这场插曲,说到底,本来够人好好想一想。但我\dmark{}\origpage{1089}恐怕他们谁也没费这份心,我更愿意猜想:全体一致,大家争先恐后把这段难堪的插曲忘了个干净。不过说到底,只要肯费点心思想一想,而不是拔腿就逃,这桩阴暗的事情里,毕竟还是有\textbf{一件}事是清楚的。Mebkhout的语气和措辞(何况他们认识此人,至少在各处讨论班里同他打过照面),使人几乎不容怀疑:
% ---------------- p.826 ----------------
这件事里必有\textbf{一个骗子}——不是Mebkhout,就是Kashiwara。

当然可能,他们在心底里早已提前断好了案:Mebkhout在编故事——那位尊贵的来访客人会去掠夺一个无名听众,这叫人怎么想象!这便等于说:面对一个无名之辈,名人无论做什么,都在一切嫌疑之上。这便是\textbf{掠夺的空白授权},交给有名望的人,用来对付求告无门的人。他要说的话,不会有人听:“你们自己私了吧!”

要不,他们就是一头埋进了怀疑里:谁知道谁说真话、谁说假话?(更何况,只要捂起耳朵,就更不知道了!)不错,一个Kashiwara,当着事主的面公然掠夺一个不知名的陌生人,这份蛮横的厚颜,看上去几乎难以置信。可说到底,还有更难置信的事:一个无名的陌生人(他们全都认识他,而他也从未靠骗子的把戏、靠那份胆量在他们眼前露过一手……)竟敢当众指控一位Kashiwara犯下粗暴的剽窃——假如他要说的纯属凭空编造……再退一步,就算他讲的或许有根有据,拿一句“你们自己私了吧!”把他打发到九霄云外,这依然是掠夺的空白授权。这就好比,一个人在大街上被一群穿燕尾服的恶棍洗劫,喊一声“抓贼呀!”——人们却冲他喊:“你们俩自己私了吧!”

况且据说,长期以来,纽约和其他美国大城市的下等街区,事情就是这样办的——在那里,谁也不想跟当地称王称霸的黑手党找麻烦。反正如今,数学界里事情就是这样办的(从什么时候起,我说不上来),在那些被当作“上等街区”的地方就是这样办的——比如Goulaouic--Schwartz讨论班\fnote{829}{(*)}{在这里我可以高兴地指明:那桩令人难忘的插曲发生的那一天,“他的”讨论班上,Laurent~Schwartz本人并不在场。后来他是否听人说起过,我就不得而知了。},或者在那班“搞”代数簇上同调的名流当中。

\pmark\origpage{1090}若按理性的说法、照字面来理解,“你们自己私了吧”这句话简直近于低能——何况无论怎样,两造之中必有一方是不老实的,这一点再清楚不过。在心理层面上,这句低能的话所表达的,是在一种被感到“碍事”的处境面前的一种\textbf{卸责}。这也是对下述明摆着事实的蓄意无视:遵守数学家这一行当基本伦理规则的问题,绝不是一件纯属“私下”的事,可以由僭取了蔑视这些规则之权的一方,同吃了亏的另一方,关起门来了结。这是一件\textbf{公共事务},一件关系到\textbf{每一个}数学家的事。

正是靠着普遍的冷漠,靠着每个人在承担个人责任时的惊慌失措,黑帮心态和“反常层会议”(Colloque~Pervers)那样无耻透顶的勾当,才能在科学界肆无忌惮地兴旺繁盛。一些人的惊慌,同另一些人的无耻,是\textbf{同一种腐败}的正反两面。1980年4月22日那一天拔腿就逃、捂着耳朵的那些人,同样促成了那场令人难忘的讨论会在次年的\emph{封神},其功劳丝毫不亚于那班凭空搭起这场辉煌大骗局、又神气活现跑去里面招摇过市的黑帮老大。

(6月3日)也正是在Mebkhout最后一次来我这里的时候,我从他那里听到了一些很有教益的细节,涉及同一场辉煌讨论会的某几位与会者,以及他们各自身上竞相盛开、一朵赛过一朵的“新风格”。我有机会翻了翻那次会议的报告集第二卷,里面有Verdier的文章、Brylinski--Malgrange的文章,又扫了一眼Laumon的博士论文(用一双比初次收到它那天更有经验、也更不分心的眼睛)。这篇论文其实是同N.~Katz合作的工作。关于这些工作里遵循的“新风格”,我在注“光荣之日”(上帝知道,这个名字它受之无愧……)的那条长长的脚注(打字稿第962页)里写了几句评论;在那条注里,为了更细的情节,我还指向了眼下这条注
% ---------------- p.827 ----------------
(当时它还没写出来)。说过的话,就要兑现!

Mebkhout给我讲过,他有怎样的荣幸和便宜,得以两次向N.~Katz谈起自己关于对偶的想法、关于连续系数与离散系数之间联系的想法。第一次是在雷恩的$p$-进分析讨论会上,1978年7月。当时他\dmark{}\origpage{1091}在“小圈子”里讲解了他的定理:复解析空间上D-模的整体对偶定理——那条为Serre对偶和Poincaré对偶加冕的定理\fnote{830}{(*)}{关于这条定理,注“作品……”与“三座路标——或天真”(n\textsuperscript{o}~171(ii)、(x))里都有谈到。}。在场的有Katz,有Illusie——正是在\emph{葬礼}里已经不止一次谈到过的那两位。Illusie呢,像他惯常那样和和气气、客客气气,觉得这真是漂亮极了——大致是这么句话\fnote{831}{(**)}{况且这本身还是一种分文不费的“客气”。两人的反应路数彼此不同(Illusie是“阴”,Katz是“阳”),实质却并无二致:只要是“无名先生”(Monsieur~Personne)说出来的,那就从这只耳朵进、打那只耳朵出!关于这一点,见注“蒙骗”(n\textsuperscript{o}~85$'$),尤其是打字稿第351页上我对Illusie的那几句观察。}。至于Katz——我猜,他在那儿是平生头一回听人谈起D-模(那阵子,这东西还远不是那场令人难忘的讨论会之后的那种大时髦)——他只干巴巴地撂下一句“这早就有人知道了!”,随即转身就走。既然跟他说话的,不过是个不知打哪儿冒出来的“无名先生”——而他N.~Katz,就在同一年,还要在几千位尊贵同行面前发表演说,向新科菲尔兹奖得主Pierre~Deligne致敬呢……——那当然只能“早就有人知道了”。

第二次,是1979年9月莱苏什讨论会之后不久\fnote{832}{(***)}{关于莱苏什讨论会,以及Kashiwara在Goulaouic--Schwartz讨论班上的那桩骗术,见注“黑手党”(n\textsuperscript{o}~171),部分(b),“最初的麻烦——或太平洋彼岸的黑帮老大”。}。Katz那时在法国高等科学研究所(IHÉS)。Katz在$p$-进微分系统方面的功力尽人皆知,而Mebkhout心里明白,这同他刚在莱苏什讲过的“好上帝定理”多少有些干系;于是他专程赶到法国高等科学研究所,把他在莱苏什的文章送上门去,同Katz谈自己的想法和结果。有了雷恩那场接待打底,可以说他这人是有点韧劲的,竟不知疲倦!反正,这一回的戏码跟上一回差不多。Katz又一次居高临下地接待了这个不知打哪儿来的无名之辈——此人竟敢第二次跑来催问,而且没准儿连招呼都没事先打。人一旦成了大人物,有时就不知道该怎么躲开这些缠磨人的家伙了……

一年之后,只消这同一批想法——由一个无名的陌生人,在孤独中久久揣着、慢慢养熟——跟着那场Katz先生遗憾未能赏光的辉煌讨论会的余威,被四处吹得震天响,说成是某位Deligne(或者某位Kashiwara,已经说不清了……)最新挖到的宝贝,它们便\dmark{}\origpage{1092}在这位大人物眼里陡然兼得了分量与斤两。替他把来龙去脉讲解清楚的,准是Laumon——Deligne最出众的弟子之一。而这位Laumon,对这些想法的来历也一清二楚,而且是第一手:因为正是那个无名的陌生人亲口告诉过他。可弟子以踏着老师的足迹为荣;而老师呢,早已明明白白、毫不含糊地示范过:对那个被判给沉默与湮没的人,该采取什么样的做派。

台前灯火的荣耀,归Deligne们,归Verdier们;Brylinski、Katz、Laumon诸公,也都赶在好时辰跑来,好从中分得自己的一份!乐声与吹吹打打归他们,感恩戴德的群众的欢呼也归他们——这群人欢天喜地蜂拥而至,庆祝这落入“新主子”们之手的“高尚行当”。

\subsubsection*{墓后尾声——或洗劫}
\tocsubsub{墓后尾声——或洗劫}

\notehead{171$'$}{}

\pmark\origpage{1093}(6月14日)直到一个月以前,我始终觉得,“葬礼”的精神所及,仅限于我间或称作数学“上流社会”或“大世界”的那个天地,尤其是那个天地里我素来出入、自己也身列其中的那些圈子。在朗格多克科学技术大学(USTL,蒙彼利埃)——它十二年来一直是我的挂靠机构——我不曾察觉到任何排斥的迹象,任何故示轻慢、失礼乃至粗暴的姿态,与那场十五年来方兴未艾的“葬礼”同一走向\fnote{833}{(*)}{在注93(第396页第3段)里,我尤其表达的就是这层意思。}。一桩新的事实刚刚闯入这幅安宁的图画,以一种严酷的方式改写了这幅图画,也改写了我同我的挂靠机构之间的关系。

% ---------------- p.828 ----------------
依照积习已深的机制,我起先并没有想到要把这桩新近的插曲收进我的见证《收获与播种》——乍一看,它出现在这里,仿佛“汤里落进了头发”,好不是地方。是顶住了几分切实的抗拒之后,我才终于承认:把这段插曲按下不表,便是有负于我的见证的精神。诚然,这段插曲余温犹在,何况我这一次“承受”得相当吃痛——这倒又给我方才影射的那些“积习已深的机制”平添了一分力气。可这一次,我领教这桩事故那些雄辩而不合时宜的教训时,感受如此鲜活,这本身也是一个征兆:它触及我至深——触及我的职业活动,触及我同我所隶属的职业圈子的纽带。所以,这恰恰属于《收获与播种》立志要作深入见证的那类事情:不允许有任何“保留的角落”,是我禁绝自己去触碰的——无论那“谨慎”是对我自己,还是对任何人,都用错了地方。

此外,在我对“葬礼”的思考这个更专门的框架之内,我感到一目了然:两者之间有着直接的关联。有可能,这类关联并不是单纯的因果关系:当地的某些同事,或许是终于把“葬礼”当真看在眼里,于是认定,从今往后他们也不妨“照样来上一场”。即便果真有这样一层因果,依我看,它触及的也只是事态一个附带的、偶然的侧面。然而另有一个更其本质的侧面,也是最令我震动的,无论在大写的“科学大世界”里,还是在一所外省的平凡大学里,皆无二致:那就是学术与大学环境中某种\pmark\origpage{1094}\textbf{退化}——也许是空前的退化:人际关系的品质在退化,礼数与尊重他人这些最起码的形式在退化,科学伦理亦然;而科学伦理,本就与尊重他人、尊重自己牢不可分。因此,下面这些页,读者不妨看作对“一个时代的风俗画”——在数学环境里,想必是一个时代的末日——的一份贡献(在围绕“葬礼”的思考一路提供的众多贡献之外)。

与其在这里把事件的经过或详或略地重述一遍,我更愿意转录四份\textbf{文件}——它们同样能把事件描述清楚。即:
\begin{enumerate}
\item 一封“致USTL数学任教同事书”,日期为5月28日,我在信中向他们通报了某种情况,并表示希望在全体大会上加以讨论;
\item 一份“答复”,出自Charles夫人——USTL数学楼的房舍负责人——形式为5月30日的一封通函,抬头具我的名,而实际上寄给全体数学教师;
\item 一项决议,由UER~5(第5教学与研究单位)全体大会于6月6日表决通过,大会议程为:“关于Grothendieck教授办公室搬迁的情况通报与讨论”;最后则是
\item 一封“致数学楼旧日同事书”,日期为次日6月7日。
\end{enumerate}

% ---------------- p.829 ----------------
我没有把两封信列入这些文件:一封是我5月21日致Charles夫人的信(第一份文件里提到它),另一封是我致USTL临时行政负责人Cano先生的信(这份文件里提到它,第四份文件——即“一场误会的尾声”——里也提到它);比起下面转录的文件所含的内容,这两封信在我看来没有提供任何新的信息。

对Charles夫人的信,我只作一条评注(“事实上很难联系上他”——这里的“他”,指的是区区在下,而这封信名义上正是写给我的):蒙彼利埃寄到我住处的信,一天即可到达;而多年以来,我离开住处,唯有去USTL的时候。

% ---------------- p.830 ----------------
\begin{quote}
\pmark\origpage{1095}朗格多克科学技术大学\\
数学研究所\\
数学楼里的一场洗劫\\
致USTL数学任教同事书\\
Alexandre~Grothendieck 撰

\hfill 蒙彼利埃,1985年5月28日

亲爱的同事:

上星期,我托UER的一位女秘书去我四楼的办公室取一份文稿,才得知那间办公室已被搬空,我的东西荡然无存——这一点我就在今天亲自核实过:只剩下一片光地板。从没有人通知我:我的办公室竟可以不经任何手续就被征用。因此我也就无从对这一举动表示同意,更谈不上准许任何人在我不在时进入我的办公室、触碰我的东西。当天我就给UER主任Lefranc先生(Lefranc)打了电话,把情况告诉他;此事(看来)出自Charles夫人的主动,这通电话似乎也证实了这一点。我向Lefranc先生言明:这种做法令我震惊;以如此粗暴的形式调换办公室,要我同意,绝无可能;我期待我的东西在最短期限内物归原位。他向我保证,会去办妥该办的事。

也就在同一天——5月21日,星期二——我写信给Charles夫人,告诉她:我把我办公室这场来得不是时候的“清空”视作滥用权力,并且把它感受为一种暴力;我期待她作出原原本本的解释,期待她毫无保留地道歉。否则,我将把问题提交大学理事会,由它裁定:对USTL的一位教师使出这类手段,是否应当被当作理所当然之事。

今天来到USTL,我得以证实:Charles夫人认为无须回复我的信(信的副本我也送交了Cano与Lefranc两位先生)。Lefranc先生同样认为无须就下述事实向我作任何解释:自从他向我保证会办妥该办的事、让我的东西回到办公室以来,一个星期已经过去,我的办公室依然空空如也。那些被席卷而去的东西如今安在,他也好,Charles夫人也好,都觉得无须告诉我。我辗转从几位秘书那里得知:这些东西存放在她们之中某一位的办公室里。此外,我曾在会议室里与Charles夫人打了个照面,她向我担保,她不过是执行UER主任Lefranc先生的指示,并请我就此事去找他,说这事与她不相干。在事态了结之前,Nguiffo~Boyom先生慨然愿意与我合用他的办公室。

% ---------------- p.831 ----------------
\pmark\origpage{1096}也许,只有我一个人觉得这里面有什么不对劲——一种暴力,一种蔑视;诚然,被人摆出这般架势、不经任何手续扫地出门的,也的确只有我一个。(倘若除我之外还有谁认为,这可不是他愿意在USTL工作的那种气氛,那就请他站出来让我知道,我会由衷高兴……\fnote{834}{(*)}{不言而喻,这样一个举动,对我来说只有一种情况下才有意义:大家明白,它使签署者本人担责,并准许我公开提及此事。})。至于我,我认为,在这场“误会”(姑且借用我一位同事那个可爱的委婉说法)之后,召开一次UER会议,让主任Lefranc先生和Charles夫人有机会就他们的意图与动机作出交代,也让UER的教师们说说,他们是否认为这类做法属正常(当它们落到别人头上的时候……)——这算不得什么奢侈。

我在USTL十二年,常常有机会领教Lefranc先生助人时的善意、热忱与高效——对此我心怀感谢。正因如此,眼见他甘当他人手中的工具,纵容一种专断与蔑视的气氛确立起来,我要收回对他的信任,遗憾也就更添一层。从现在起,我请求他要么切实承担UER主任之责,要么辞去职务。我也请求Charles夫人辞去她那UER“房舍负责人”的职务——这套职务,她滥用得津津有味。

敬候您的(或你的)答复

\hfill Alexandre~Grothendieck

又及:我天性乐于助人。去年,应Lefranc先生之请,我曾同意与Lapscher先生(Lapscher)对换办公室;而不久之后有人告诉我,他随后改变了自己的计划。不言而喻,我的同意,并不意味着我准许对我办公室来一场洗劫——无论在当时,还是在任何别的时候。
\end{quote}

% ---------------- p.832 ----------------
\begin{quote}
朗格多克科学技术大学

\hfill 1985年5月30日,星期四

\underline{数学}

\pmark\origpage{1097}J.~Charles夫人 “数学研究所房舍负责人”\\
致A.~Grothendieck先生,数学教授。

\begin{center}亲爱的同事:\end{center}

1. “数学研究所房舍负责人”的“工作”,从何处开始,到何处为止?这位“负责人”受理数学教师们提出的申请:\\
——或是为一位新任的教师(或研究人员)安排办公室\\
——或是为一位已有办公室的教师(或研究人员)另行安排办公室。

后一类申请,一般都举有工作上的理由:把同一研究组的成员归拢到一处。

这位“负责人”于是着手研究各种可能,首先是同UER~5的主任一道研究——数学研究楼的房舍,其正式委派的管理人,正是由USTL校长指定的那位主任。随后他同有关人士一起寻找可行的方案;调整须待各方谈妥之后方才实施。

2. 近年来照此办成的事情:\\
——几何研究组成员的归拢\\
——力学研究组成员的归拢

3. 这项“工作”中遇到的困难:\\
——实际接触到的每一个人,几乎都觉得自己的办公室归自己“所有”\\
——要强迫什么人“更换”办公室,看来绝无可能。

4. 我接到的最近一份申请,以及为所提问题寻求“方案”的经过评估:\\
——Lapscher(LAPSCHER)教授提出的申请:把Lapscher先生本人、他秘书的办公室、Micali先生归并到同一层。\\
——设想的第一种方案:在三楼与四楼之间交换办公室,好让“申请人”们归并到四楼。这次交换尤其涉及Grothendieck先生与Therond(THEROND)先生。Grothendieck先生经UER~5主任亲自接洽,\dmark{}\origpage{1098}向后者言明:\textbf{办公室在哪儿对他全然无所谓,只要有一间就行}。相反,Therond先生一度表示同意,随后却拒绝了任何搬迁。\\
——设想的第二种方案:我随后请Lapscher先生自己去联络他的同事,以提出另一种方案;这一点UER~5主任向他确认过。他的奔走情形,他不时向我们通报:五间办公室的“占用者”都同意作一次轮换;Grothendieck先生的同意,则出自他同UER~5主任的那次谈话。\\
——第二种方案的实施:UER~5主任得知这一致同意之后,为拟议中的办公室调整开了“绿灯”。

% ---------------- p.833 ----------------
Lapscher先生曾向我谈起钥匙的问题——搬迁启动而未完结的那段期间怎么办;我向他指出:\\
——新钥匙八成一把也没有,\\
——这场搬迁本可由全体有关者一齐动手,几个钟头做完,拖长它,在我看来并不可取。

Lapscher先生随后告知我:Grothendieck先生办公室里的物件,已被搬进他未来的办公室;这件事办成之前,未能先联系上Grothendieck先生。

应当指出:Grothendieck先生住处远在蒙彼利埃之外,目前又借调在法国国家科学研究中心(CNRS);事实上,要同他取得联系非常困难。

5. 以“负责人”的身份,我对这件或许可以称作“冲突”的事情的印象:\\
——我曾有机会向Grothendieck先生言明:我以UER~5的名义行事,不能亲自答复他的信;因此他应当向UER~5的主任索求答复。继这封发给大家的第二封信之后,我认为自己必须走出曾自我施加的“克制义务”。\\
——至少在搬动有关人士的东西之前先通报一声,这在我看来本是可取的\\
——把搬迁压缩到至多半个工作日之内,这在我看来本也是可取的。\\
——所设想的方案在我看来是妥当的:它丝毫改变不了有关各位的办公室占用率。

我不等待答复。

先生、亲爱的同事,请接受我最诚挚的问候。

附记:此信副本送下列各方知悉\\
——收讫Grothendieck先生1985年5月28日来信的全体数学教师;\\
——UER~5主任,他另外还收到Grothendieck先生1985年5月21日写给我的那封信的副本;\\
——USTL临时行政负责人,他已收到5月21日那封信的副本,我再随函附上5月28日那封信的副本。
\end{quote}

% ---------------- p.834 ----------------
\begin{quote}
朗格多克科学技术大学\\
数学研究所\\
\pmark\origpage{1099}数学研究所

1985年6月6日(星期四)18时会议纪要。

出席者:Auberson先生、Charles夫人、Ciulli先生、Contou-Carrère先生、Cuer先生、De~Lima先生、Delobel先生、De~Robert先生、Grothendieck先生、Hocquemiller先生、Escamilla先生、Hubert~Coulin小姐、Lefranc(LEFRANC)先生、Loupias先生、Meden夫人、Molino先生、Pierrot夫人、Pinchard先生、Saint~Pierre先生、Voisin小姐

经讨论后,到会者(19人)以16票赞成、3票弃权通过下列文本:

“数学教师们就Grothendieck先生的东西被搬动时所处的不可接受的情形,向Grothendieck先生致歉。他们共同承诺,警惕这类令人遗憾的事件不再发生。尤其必须明确:一间办公室的钥匙,未经占用者明确同意,任何人不得动用。”

\begin{center}
Lefranc先生\\
主任
\end{center}
\end{quote}

% ---------------- p.835 ----------------
\begin{quote}
\pmark\origpage{1100}朗格多克科学技术大学\\
数学研究所\\
一场“误会”的尾声\\
致数学楼旧日共事的同事们(教学人员与技术人员、第三阶段学生)书\\
Alexandre~Grothendieck

\hfill ……1985年6月7日

亲爱的同事:我在这里为我的办公室遭洗劫一事写下尾声——此事已在我5月28日的信里提起。那封信当时只寄给了数学教师,而它同样关乎、并以同一名义关乎在数学楼里占有一间办公室的每一个人,无论男女。我当初没有把信同时寄给技术人员和第三阶段的学生,是出于疏忽,也失之欠考虑——我(仓促地)以为,那样做等于给这桩插曲一个它并不应得的扩展。我在此向有关各位诚挚致歉——更何况,他们中有好几位(照理本该是“不知情者”……)曾向我表示同情,令我感动。大概也是因为这个疏忽,昨天专门讨论这桩插曲的UER全体大会,只限于“UER~5成员”参加。

在许许多多别的事情之外,这桩插曲还让我领教到:这类事在UER~5并非头一回发生——只是头一回,被瞄准的是一位“A级教师”。昨天表决通过的那份虔诚的决议,能否阻止这类事件重演——在普遍的冷漠之中(一如从前),尤其是对非在编教师、对第三阶段学生的冷漠——我不知道。我会留意向Mori夫人与Moure 夫人核实:她们是否确实收到UER主任的指示——今后无论以什么借口,都不得把任何一间办公室的钥匙交给任何人,也不得为任何人动用,除非得到其中一位占用者的明确许可。

我上一封信,结尾用的是这样几个字:“敬候您的(或你的)答复”。作为对这一期待的回应,我收到三份同情与声援的表示。它们来自Louis~Pinchard、Pierre~Molino与Christine~Voisin。同样,我还收到一份同一方向的表示,出自第三阶段学生Philippe~Delobel——他(与Christine~Voisin一样)曾随我读过DEA(深入研修文凭)。昨天有几位第三阶段学生出席全体大会,正是出于他的倡议。对他,也对我方才提到的所有那些人——他们(不含糊、不闪避)向我表示了声援——我乐于在此表达我的敬重与感谢。像这类的“严酷”经历,果实之一便是让人认出自己的朋友——倘若运气好,还有朋友可认……

我还收到另一封回应我的信,出自一位同事,他分明为发生的事情而喜不自胜,还借此机会温文尔雅地拿我开涮。这是我收到的这一方向的唯一回声。其余诸人那里,不少是这一位的全然冷漠、那一位的局促不安(其间不止一次,我触到一种没有说出口的恐惧:怕给人留下坏印象,从而断送晋升的机会,或危及本已朝不保夕的位置)。在那些人当中——凡是激动到肯为此劳神出席那场全体大会的人(大会是最后一刻偷偷摸摸才召集的,虽然一星期前就已排定……)——我感受最深的,是一种处心积虑、打定主意要把鱼搅浑的心计,口口声声“人人都和善,个个都可爱”。末了(在足足三刻钟的空谈之后),大家转而拿那个众望所归的“坏蛋”问罪,那位缺席者(多么凑巧),Lapscher先生——(据人们方才刚刚暗示出来的意思)\dmark{}\origpage{1101}发起这场偷袭的,正是他。至于走到点名道姓追究他那一步,那是谁也没有打算的,可怜的人——当然,对其他任何一个人,也同样如此,这不必说。

% ---------------- p.836 ----------------
在这桩洗劫插曲中以这种那种名分牵连在内的各“负责人”那里:一位Lapscher的无耻蛮横,一位Charles夫人“以取乐为务”的粗鄙(她在既成事实面前把这场偷袭遮掩过去,还添上她自酿的那一份放肆),以及USTL临时行政负责人Cano先生的失礼——我写信告知他事态原委,请他把问题提交大学理事会,他竟免作任何答复——这一切都令我震惊。然而最甚于一切的,是UER~5主任Lefranc先生那模棱两可的态度,令我错愕,令我痛心。从5月20日星期一(那一天,我把刚发现的情况和我对此事的感受告知了他)直到昨天,他始终不认为适宜之举,或者是向我说明究竟发生了什么,或者是毫不含糊地同一位Lapscher的强盗行径、一位Charles夫人的粗鄙划清界限。自始至终,他竭力维持那场倒霉“误会”的虚构,硬是给那些行径——在我看来无法容忍的行径——装出一副无关痛痒、甚至堪可敬重的模样。为的是不叫任何人难受,可以断言,他选择了“又护山羊又护白菜”:(多多地)照顾山羊,(稍稍地)照顾白菜。

在其他种种迹象之列,我也好好记下了一笔:我原以为可以列入朋友之数的许多人(其中有三位曾是我的学生)所守住的沉默;这一位招摇的漠然,那一位的窘迫不安,还有一位甜腻腻的欢欣雀跃。还有一位Micali的沉默(他是这场偷袭的共同受益人;而几年前,他就已有充分的机会确信:招得Charles先生与夫人不快,会带来多少麻烦……),以及Brun小姐的殷勤卖力——她领一位Lapscher的吩咐,客串了一回雇佣锁匠兼搬家工(等到这场行动的性质再无可疑之时,连一句遗憾的话也没有)。

在这一切的底子上——昨天我又见到那间十二年来一直是我的办公室,这一回它成了战场:我的东西(连家具一并)被七手八脚重新堆作一团(距那场“闪电”偷袭,足足过去了十五天……)——如今我已没有心绪再在里面重新安顿下来。人们向我担保,同样的插曲在我身上重演的可能性微乎其微;何况我尽可以先发制人,把那把迄今托付给Mori夫人与Moure 夫人的第二把钥匙收归己有。但只要物质上办得到——尤其在我借调法国国家科学研究中心的整个期间——我宁愿从今往后放弃在USTL使用办公室,不加争斗地把地方让给Lapscher们、Charles们以及各色同好。

如果能够避免,我不会再在USTL执教。在那儿,我自始至终——可以肯定——不过是个异乡人,一个故乡在别处的人:无论我对待数学的路数,还是我对待教学的路数,抑或我的生活方式,皆是如此。大学这个小天地所能教给我的,我想我已经学到了;而最后一“卷”,便是这桩插曲教给我的一切——它刚刚在一片皆大欢喜中落幕。我刚刚参与的这场UER~5会议,很可能是我的最后一场;这封信,也很可能是我有机会写给你们(或写给你)的最后一封。而这一回,我不等待答复。

\hfill Alexandre~Grothendieck
\end{quote}

% ================================================================
%  18.5.5 Le seuil + 18.5.6 L'album de famille (原书 p.836–864)
%  源文: _work/chunks/chunk08.md, chunk091.md, chunk092.md,
%        chunk093a.md, chunk093b.md, chunk094a1.md, chunk094a2.md, chunk094b.md
% ================================================================
\subsection{门槛}   % 18.5.5. Le seuil

\notehead{172}{}

\pmark\origpage{1102}(3月22日)我原以为顶多一两天、撑死十来页,就能把这些自从10月起我就打算重新巡视的
出名的“四种运算”走完一遍。可如今一鼓作气干了三个多星期,页数排到了一百上
下——居然还没有全然写完!初稿从2月26日写到3月1日,就花了四天。它给我的,不过是一块底布,
好(终究)在上头绣出一桩“事情”来,而不只是一纸侦查结论。3月1日的第二天重读初稿,那印象可
不妙:一张没完没了的“诉苦状”,而且,除了三四位真正在行的专家(还得假定他们有耐心读完
……),只怕谁也读不懂。我明白:至少得约略交代一下所为何事,也就是说,至少得铺一个语境——
否则就白搭\fnote{835}{(*)}{整部《收获与播种》的思考里,除这里之外,我曾对“即兴”写作方式作过这种(规模确实小些)让步的,只有“那张条子——或新伦理”一节(n\textsuperscript{o}~33)与“不义——或归来的意义”那条注(n\textsuperscript{o}~75)。}。

这就不可避免地带来一些重复,同\emph{葬礼}第一部分相比——可有些情形,重复不但有益,甚至
必不可少(在数学里,一如在别处)。而且在这些情形里,人们很快发现:那些所谓的“重复”其实并
非重复,因为被“重说”的东西,实际上是\textbf{重看了一遍、以新的眼光再看},并且照明已经不同。在
为“四种运算”铺设“语境”、交代我著作的某些侧面时,我觉得自己对这份著作学到了些东西,把它
摆得更正了。这么做,我也许并没有就我自己或他人学到什么真正新的东西,可我不后悔费了好几
天的功夫,把那份“诉苦状”式的初稿这样重写。这份著作,我放入过我所能给出的最好的东西;它
值得我怀着成熟给我的那份回顾,以不同的光线重新去认识它。正当我准备原原本本记录这份著作
自我离去后(交在了可靠的人手里,对此我从不怀疑……)遭受的一切之时,好事在于:我在它身上
\textbf{停一停},在它的位置上、在构成它之美的那个统一体上,哪怕只停几页——以此再一次,向那
些我亲眼看见被践踏的东西,表达我的敬意。

\pmark\origpage{1103}可这还不算完,差得远呢!放弃了“诉苦状”的写法——那种夹着编号、指向\emph{葬礼}第一部分
那些血肉更丰满的注的写法——我明白了:这些我重新拿起的笔记,同《收获与播种》里所有其他
章节与笔记一样,必须自身可读,把它们要说的实质,完整地传达出来,独立于那些指向思考\textbf{另一
时刻}之注的引用。这又带来许多“并非重复的重复”——也就是,以新的照明,重看将近一年前、在发
现的新鲜震动里逐日记下的东西。何况当时,那么多意外的、有时难以置信的事实劈头盖脸,压根谈
不上一场稍具方法的“侦查”。那时,我不过竭力去接住劈头砸下来的东西,勉力把它们“安放”进去,
细节顾不上。我的精力,大部分被\textbf{招架}这些我发现的原形毕露之物的那种\textbf{癫狂}、那种难以置
信(就像那篇皇帝新衣的故事里一样……\fnote{836}{(*)}{见同名之注,n\textsuperscript{o}~77$'$。})所吞掉;尤其是,要承受那股突然扑面而来的暴戾、犬儒与轻蔑之气——“顶着那副光鲜的面孔……”我太熟悉的面孔;那是另一个时代的气息,我经历过它,我没有忘……

而这最近三个星期,反倒成了补完去年那场风暴式侦查的机会:把某些文本再往细里翻一翻(SGA~5,
尤其是那卷所谓的“SGA~4½”)。于是生出一串(有一阵眼看不象话地收不了尾的!)多少原原本本
的脚注;其中一些成了子注,而其中一个(预定的名字叫“那个公式”)占用了我整整四天,又裂成
另外四个\fnote{837}{(*)}{(6月1日)如今已经变成六个了……}……有几次,我觉得永远写不完了——可到头来,它还是收敛了
\fnote{838}{(*)}{(6月1日)当然,只是一场临时的“收敛”,因为“封神”那条注,最后炸成了三十来条各自独立的注、子注等等,单是它们就有足足一百五十页!}。眼下撂下的,还有十来页涂改得太厉害、须得重写的页,以及“分赃”与“封
神”两条注的脚注,留待日后补上。眼下这样就可以了!哪怕日后还要回头管\dmark{}\origpage{1104}“后勤”,我急于收场,
急于不再耽搁,把“四种运算”这一章里我看到的还有实质可说的东西说出来。

我在\emph{葬礼}里区分出两个“面向”或“层次”:它们紧紧相系,却毕竟不同。它们之间(至少在
我眼中)相当分明地隔着一道\textbf{门槛}。

一方面,是“时尚之风”的面向(有时直到那股“嘲弄之气”——我在《收获与播种》里已不止一次
有机会说到它)。它主要表现在我在别处\fnote{839}{(*)}{见注“掘墓人——或教团全体”,n\textsuperscript{o}~97。}称作“自动排斥的态度”的东西上:这种态度常常拦腰截断数
学常识的简单反射,冲着某些人及其数学贡献而去——这里说的是我,以及另外几个人;他们被归
类(尽管当事人有时竭力想同我划清界限)为同我“串通一气”。在我的情形,要“排斥”(或“埋葬”)
我带来的\textbf{一切},当然不可能:好大一部分,早在我1970年告别数学舞台之前,就已进入日常通用
的公共领域\fnote{840}{(*)}{不过,确实,连某些早已进入“日常”通用(至少在我学生与近旁合作者的小圈子里)的观念与技术,自我离去后也被埋葬了。不妨说,尤其如此的是$\ell$-进上同调那件工具:我在 SGA~5 里(以 SGA~4 的关键结果为起点)极详尽地发展过它。它被我那班上同调学生压在斗底下,Deligne首当其冲,直到1977年,才以我所知的那个形状与精神被掘出土来。}。不过,
确实(而且我在一年前的注“我的孤儿们”(注 n\textsuperscript{o}~46)里第一次确认了这一点):我的著作——
写下的与未写下的——在\textbf{上同调主题}上,绝大部分自我离去的次日起就被埋葬了,头一份经手
的正是我的学生。(我引入的某些主题,四年后、七年后或十二年后被掘出土来,不提我这个人——
但那我们已经触到“第二层次”了……)

这样的自动排斥,人们尽可以惋惜:它有时违背对他人的最起码的体贴与敬意,而且在任何情形下
都与数学的常识与鉴别力无关。而当它打击的是那些手段有时十分出色的年轻数学家时,就更值得
惋惜:“轻蔑之齿”咬灭了一种欢乐,糟蹋了一份原本美好的激情,让它化为苦涩——投入眼看着被
糟蹋了(按那些说一不二的共识……)。还值得惋惜的是,\dmark{}\origpage{1105}当这种排斥打击的是那些简单而丰
饶的观念——它们早已充分证明自己,从无中生有地催生出如今“人人”看也不看就使用的强大工
具。在前一种情形(一个被毁掉的志向),伤害多半不可逆转;在后一种情形则不然——
因为迟早,那些简单而根本的观念,那些“在路上”的观念,终会显形或重新显形,终会进入公共遗
产。不管怎样,谁也不能合情合理地强求任何人,对一个人、一部著作、一个观念\textbf{往好里}想——
只要他(出于只有他自己知道的理由)一心想\textbf{往坏里}想,或干脆忘了它。这类问题,当然属于“个
人伦理”,而且属于其中精细而要紧的部分;但依我看,把它弄成一桩集体的“科学伦理”问题,恐怕
不妥;倘若硬要试,只怕药比病还凶……

我方才影射的第二个“面向”或“层次”,恰恰是这种集体伦理被破坏的层次。我说的那道\textbf{门槛},
是一条\textbf{共识}——据我所知,自有科学见诸文字记载以来,它在一切科学里都得到普遍接受。这条
共识规定:谁也不得把他从\dmark{}\origpage{1106}别人那里取来的观念\fnote{841}{(*)}{我在这里说“观念”,当然,在数学里绝不单指“结果”。常常,一个提得好的\textbf{问题},它触及一个此前无人看出的关键点,就比一个“结果”更重要,哪怕更难。即便这个问题还没有凝结成一个精确的\textbf{命题}——它本可以构成一个假设回答的胚胎,或者一个(尚属猜测的)多少完整的回答——情形也是如此。当然,从一个起初含混的问题里提炼出这样一个命题,是数学工作中本质而具创造性的一部分。把一个(也许深刻的)问题加工后的版本端出来,却隐去它的出处(即便加工是这位魔术师自己做的),同隐去一个形状完整的深刻命题的出处——借口是自己给出了证明——一样,都是剽窃,正如把别人给的证明当成自己的拿出来一样。
对丰饶\textbf{观念}的引入,道理相同——它们往往比好的命题更关键:因为“好命题”的问题,只有在已经提炼出好观念之后才谈得上。这里的道理相同:借口修改过、甚至改进过取自他人的观念,就隐去其出处,其不老实,同原封不动地“借用”这观念,毫无二致。最常见的情形是:提出一个问题(哪怕还含混),拿出一个命题或一个观念(哪怕不完善、临时),才是关键的一步,而不是人们后来加上的那些改进(更精确、更广、更深)。可即便情形不是这样,对那个在改进所受之物上做出原创工作的人,这仍旧不能当作隐匿其所受之物的“理由”(或者,同一回事:把它“拆掉”……)。
如我在别处已强调过的(在注“肌肉与肚肠(阳葬阴(1))”(n\textsuperscript{o}~106)的子注 n\textsuperscript{o}~106$^{1}$里),一条猜测性命题的“价值”,既不取决于其设想的难度,也不取决于它多“可信”,也不取决于它最终被证为真或为假。无论如何,人们愿意给一个数学观念(它表达于一个问题、一个命题、一个观念或一个证明之中)或一组观念的“价值”,很大程度上是主观的,难以成为科学伦理之共识的对象。所以,一个诚实的科学家,会标明他所使用的(显式地或默会地)、不属于“众所周知”的一切观念的出处,而不肯顺从那条斜坡:对他已在内心决断(也许是为了某桩可疑的案子……)反正“显而易见”、“不足道”、“不要紧”(或同类评语)的某个观念,隐去其出处。}

\pmark 这条共识责令我们:凡我们所陈述、使用或发展的观念,每次——至少当这些观念并非出自
我们自己的手笔、也还不是公共遗产、尚未被(不是三四个入室弟子,而是)“\textbf{所有人}”知晓的
时候——都要标明来历。

我不记得曾听人质疑过这条共识。在我身属数学圈子的岁月里——从1948年(二十岁的年轻新手,
去高等师范学院听Cartan的课)到1970年(我告别数学舞台)——我极少有机会目睹、甚至极少听说
有人公然违背这条共识或原则;只在一位同事兼朋友那里,他在这一章上稍欠检点\fnote{842}{(*)}{这位同事的情形,顺带谈过,在《收获与播种》第一部分,就在紧接着下一句所引的那一节里。经过一年多的回顾,这个“案子”如今所显出的分量,是我从前没有给它的。}。正如我在《收获与播种》第一部分里强调过的
(在“一个严守的波利钦内尔秘密”一节,n\textsuperscript{o}~21),遵守这条原则,对一个具备起码的诚实与自尊
的人,绝非自然而然之事。相反,它要求极大的警觉:童年便根深蒂固的反射,自然而然地推着我们
高估自己的功劳,把我们在自己之内消化他人观念的工作,同这些观念本身的构思混为一谈——两
者实属完全不同的层次。写上面所引那一节——一年多了——的时候,我显然对自己也没弄清:该给
这条共识多大的分量。那时我心里有一片含混(在思考的这个阶段,我自己也没有清楚意识到\dmark{}\origpage{1107}),连着一种弥漫的感受:对\textbf{他人}(比方说对我自己的学生)提出严格要求——在他们
同我的关系上遵守这条原则——是小气,是一种配不上我的卑微。所以那时,我之内有一道\textbf{暧昧};
我只是在6月1日那条同名之注(n\textsuperscript{o}~63$''$)的思考里,才把它看清。这场思考把这道暧昧驱散净尽;
而我当时意识到,这道暧昧,从最初(六十年代初)直到去年,一直沉沉地压着我同学生的关系。我
懂得了:数学家这一行(或更一般地,科学家这一行)的严谨,首先意味着对自己怀有极大的警觉,
遵守这条众共识中之首要者;同时也意味着对他人同样严格的要求,尤其对那些我们负有领入门墙
之责的人。

年复一年,我越来越懂得:这一行\textbf{绝不只是}一门技术上的窍门,甚至不只是发挥想象、解
决出了名的难题的能力。从某种意义上,这一点我从来就知道——可我低估了“伦理”的一面,或者
说\textbf{集体}的一面\fnote{843}{(*)}{我在这里无意说:一个处境的“伦理”面向,总是同时就是“集体”的面向,牵涉一个人同群体(这里即“同事”或“同侪”之群)的关系。不过,在我正在检视的这条“共识”上,情形确实如此。按照自童年起塑造了我的世界观的那些特殊条件作用,直到去年,我总倾向于低估(乃至无视)集体之物,而偏重个人之物。我的“数学历险”中“集体历险”的那一面,去年才清楚地向我显现,先是在“Galois的遗产”一节(n\textsuperscript{o}~7),但尤其《收获与播种》第一部分末尾的几节:“孤独的历险”与“一段过去的重量”(n\textsuperscript{o}s~47、50)。},以为它在善意的人们、体面的人们之间,“理所当然”。就这样,我为我刚
才说的那道“暧昧”备好了土壤;那道暧昧,也是(顶着一副虚假的“慷慨”的门面)对我的学生与
类似者的\textbf{纵容},而更隐蔽的,是对我自己的一种纵容。

我离开了这个“善意的人们、体面的人们”的圈子——它也曾是\textbf{我的}世界,我曾乐于自认同属
的世界。在去年4月19日之后的几个星期里,我稍微仔细地看了它一眼:离开不到十五年,我
在里面看到一种\textbf{腐败},做梦我也想象不出它的样子。

\pmark\origpage{1108}对我而言,这是一个谜:身为这个世界的一员,“搞数学”还能有什么\textbf{意义}——如果只
是把它当作一种\textbf{权力}的手段,或者(对身份低微的)当作一种谋一碗饭的手段,物质条件嘛,
说实在的还挺舒服(只要运气好,早早就“安顿”下来了……)。

\subsection{(5) 家庭相册}   % 18.5.6. (5) L'album de famille
\notehead{173}{}

\fnote{844}{(*)}{眼下这条“家庭相册”之注,最初本是前一道注“门槛”(同日,3月22日所写)的直接续篇。这一部分,如今构成其 a.部分(“一位备受簇拥的死者”);6月10日、11日又添了 b.(“新面孔——或吊嗓”)与 c.(“众人之上的那一位——或首肯”)两部分。下一道注“升级(2)”(n\textsuperscript{o}~174)——也是3月22日——紧接本注 a.部分(同日所写)而来。至于 b.、c. 两部分的脚注,日期是6月13日与14日。最后,6月18日又添了 d.(“最后一刻——或一个禁忌的终结”)这一部分。}

\subsubsection*{a.\ 一位备受簇拥的死者}
\tocsubsub{a.\ 一位备受簇拥的死者}

\pmark (3月22日)说得粗鲁些:\emph{葬礼}里有“时尚”的层次,也有“诈骗”的层次。也许不过是
我落伍了;我那个年代被看作诈骗的行径,如今只要干它的人属于上流社会,就早已成了完全正当、
体面的行当。也许那道“门槛”,早就荡然无存了呢?“第二层次”是\textbf{一场唯一的、规模浩大的诈骗
行动}:目标是我在上同调主题上的全部著作;继它之后,又是Zoghman~Mebkhout——那位鲁莽的继承
人、被埋葬的宗师的身后学生、固执而无名之辈——的著作。这场行动的总指挥,是另一位学
生:绝非身后的,倒是隐秘的,没错;他玩着一个心照不宣的角色——我著作的“继承人”——同时对
著作、对工匠,一并拒认、一并拆毁。他是我的朋友\textbf{Pierre~Deligne}。他的干练副官,不是别
人,正是随他一道选了“上同调”这条线的四位学生:\textbf{J.L.~Verdier、L.~Illusie、P.~Berthelot、
J.P.~Jouanolou}。死者想必被簇拥得妥妥帖帖:既有同享\emph{葬礼}哀荣的\textbf{共殓诸友}
\fnote{845}{(*)}{说实话,据我所知,“共殓者”不是一位,而是\textbf{四位}——即四道棺材注(棺材一至四,n\textsuperscript{o}s~93--96)所写的那四位。},又有生前曾是他的亲近者的那些人。至于临时凑数的殡仪工——由总指挥导演
的那场双重埋葬里,前来帮衬的——我看见另外七位“世界闻名”的数学家(借某张广告招贴的\dmark{}\origpage{1109}措辞
\fnote{846}{(*)}{指 IHÉS 1983年为建所二十五周年出的纪念小册子。关于它,见“悼词(一)(二)”两条注(n\textsuperscript{o}s~104、105),尤其第454页。}),在葬礼仪式进行之中不时露面;仪式已在家庭相册(又名“四种运算”)里
巡视过一遍。他们是(按其在仪式中的分量)\textbf{Teissier、Beilinson、Bernstein、Milne、Ogus、K.Y.~Shih、Katz}
\textbf{(N.~Katz)}。

\pmark 到此,我巡视了据我所知曾以这样那样名分\textbf{积极参与}“葬礼行动”的数学家。共十二位
\fnote{847}{(*)}{与(《收获与播种》第一部分)“耶稣与十二使徒”一节里相同的“十二”——那一节巡视了所有跟我做到国家博士论文层级的学生。确实,在我葬礼的积极参与者之中——这一回,是Springer殡仪股份有限公司那一边,而非信众教团那一边——还有 K.F.~Springer博士(那家可敬公司的共同老板)与Peters、Byrne两位博士;后文(注 n\textsuperscript{o}~175)还会谈到他们。于是是十五位!}。至于后四位,就我所知的事实,我无从断定他们存心不良。但我认为,他们的责任并不因此就轻了分毫。因为,他们乐于对他们所做的事视而不见,这本身是一种选择;选择,豁免不了他们对自己行为的责任。

至于那些直接串通的参与者,要我开一张哪怕不完全的名单,或估一估他们的人数——那人数肯定
是另一个量级——我实在无能为力。只需提醒一句:这里面有1981年6月吕米尼那场“难忘讨论会”
(人称反常层会议)的全体参加者;还有那卷顶着“SGA~4½”之名的文本的读者们之中,多少知道
SGA 这个符号意味着什么、却“听之任之”的所有人。

\pmark 我看见两份文字,它们见证着七十年代、八十年代数学里的一场\textbf{凋败},在我们这门科学
的历史上,恐怕还从来没有堪与比拟的先例。其中一份,凋败在它给自己起的名字里就绽开了——那
名字本身就是一次(天才的……)冒名:那份数字标着“SGA~4½”(作为常用的参考符号)、又称
“平展上同调”的文本——P.~Deligne署名,“协作”(除 L.~Illusie与 J.L.~Verdier等人之外,还)有
A.\textbf{Grothendieck}\fnote{848}{(*)}{关于这份“协作”的含义——它是Deligne所设骗局的一部分——见注“对调”(n\textsuperscript{o}~68$'$)。}。第二份文字,是1981年6月吕米尼\dmark{}\origpage{1110}讨论会的文集;尤其并且首要的,是其第一卷:由讨
论会引论(Teissier与Verdier署名)与讨论会的主报告(Beilinson、Bernstein、Deligne署名)构成。

倘若直接参与过这场凋败的、以这样那样名分参与的人之中,有一位能在自己身上找到那份朴素与
勇气,公开道个歉——或者更好,公开就\textbf{他自己}经历的那部分原委说个明白——这对所有
人都是一桩善事,也是那一代容忍了这等凋败的数学家的光荣。可这多半是奢望。

同样是奢望的,想必还有:让 J.L.~Verdier不再在高等师范学院占着Henri~Cartan曾占的那个位置。
那无疑是法国培养数学“接班人”的关键岗位。很久以前,当我得知Verdier——他曾是我的学生,
我曾钟爱于他——升到了这个职位,我自己觉得脸上有光(同时暗地里也有几分得意)。那时,我丝毫
没有怀疑过:Verdier会完美地承担曾属于Cartan的角色,让那些最有志于数学的年轻人在与他的接触中学好这门手艺。可今天我看见(多年前已在看,却从来没有以今天这般粗暴的清晰看见)我
错了;我在这里明白地说出来,不是要给他或任何人抹黑。我认为,他已丧失了指导研究的资格。这样
说,我并不推卸我自己的一份责任:是我把他教坏了(他同我所有其他学生一样)——那门我热爱并
且依然热爱的手艺。

\subsubsection*{b.\ 新面孔——或吊嗓}
\tocsubsub{b.\ 新面孔——或吊嗓}

(6月10日)“家庭相册”这条注动笔以来,两个半月过去了。我当然没料到,还须回到它上面来——
因为对\emph{葬礼}的调查又有了新的反复。主要是那场微薄的“封神”——我当时刚写了五六页——
炸开成一场大写的、足足一百五十页密密实实的\textbf{封神},让我随之发现了一些“新面孔”:它们理应
在家庭相册里占一席之地。也有一些早已熟识的面孔,如今显出他们同样属于那支大军——那些
在“诈骗”的层次上积极参与“葬礼运算”的人。我在这里“备而勿忘”地把他们再巡视一遍\dmark{}\origpage{1111},也好让
每一位当事人确知自己与名流为伴(不过,这事想必早就办妥了……)。新到的照片,按它们
引起我注意的先后,插入相册。

先说Springer出版公司那一边:\textbf{K.F.~Springer}(公司共同老板之一)、\textbf{Peters}、与
\textbf{Byrne}夫人(C.M.~Byrne)。细节见下文“殡仪馆——‘为科学服务’”那条注(n\textsuperscript{o}~175)。写本注开头、
3月22日的那几天,我刚收到 K.F.~Springer的信(日期3月15日)没几天;那封信,把我对那家可敬
的殡仪公司——忠于“为科学服务”的箴言——的精神的最后一点疑虑,一扫而空。

封神那一边(经由对当值无名之辈的埋葬),我知晓了\textbf{Kashiwara、Hotta、Brylinski、Malgrange(B.~Malgrange)、Laumon、Remmert}诸位的贡献;还不算一位\textbf{匿名审稿人}——其存心不良,
毋庸置疑;不过说真的,倘若要把与\emph{葬礼}沾亲带故的那些可疑文章与书籍的审稿人(reviewer) 也一一
点验,那准得另备一册相册。此外,我的老朋友Katz重新露面了,这一回的语境,让我对他本来(至少)
存着的那份善意推定,化为乌有。这样一来,据我所知曾以这样那样名分积极参与那场所谓“反常讨
论会”的骗局行动的、国际知名数学家,人数到了十四位(算上那位出名的匿名审稿人,十五位)。详
情备有文件,见“封神”,尤其“……以及横财”、“光荣之日”、“黑手党”、“掠夺的空白授权——或‘高
尚行当’”几条注(n\textsuperscript{o}s~171(iii)(iv)、171$_{2}$、171$_{4}$)。

最后,在“Motif 运算”那一边,出现了(迟做总比不做好)我的又一位学生,稍稍离大队人马有点距
离。事后,我几乎是被逼着把他算作我的“上同调”学生的第六位,尽管“我那个年代”,他对上同调是
什么,连半点概念都没有。他是Neantro~Saavedra~Rivano;显然,他被人当了“棋子”(当然,出
于他自己的自愿),而非为自己行动。他与Monsieur~Verdoux(化装成“献殷勤的骑士”)\dmark{}\origpage{1112}周旋的历险,已
在“第六根钉子(钉进棺材)”那组注(n\textsuperscript{o}s~176$_{1}$至176$_{7}$)里逐页复原——那是4月19、20日写的
(最后一道除外,尚待写出)。于是,在我“从前的”学生之中,积极参与宗师\emph{葬礼}
的,人数也到了六(十二之中)。Saavedra在\emph{葬礼}中的份额,其特别之处在于:他参与的那场
“塔纳卡范畴”(原文如此)运算,是第一场大规模行动,目标是变走我著作的一个重要侧面、以及我
发展的那门哲学(就其情形而言,是在 Motif 理论的带动之下、乘其便而发展的)的父亲身份。

把新到的面孔算进相册,并撇开Springer-殡仪馆一方的贡献、只算信众教团一方的,那么据我所
知曾积极参与\emph{葬礼}——按我那个年代人们管它叫“诈骗行动”的那种水平——的知名数学家,
人数到了十九位\fnote{849}{(*)}{算上那位出名的匿名审稿人,二十位。}。这些参与者中,只有三位——即与 P.~Deligne合署“难忘文集”《数学讲义》
第900种的那三位——在我看来,存心不良尚未坐实。

这份名单,当然远没有穷尽我的同事、旧学生或朋友——他们以这样那样名分、或主动或被动地参
与了我的葬仪,却没有走到参与一场铁板钉钉的诈骗的地步。我点出过三十来位;其中大多数,已在
对\emph{葬礼}的思考中提到;算上前面那些,凑足了结结实实的五十来位——而这还只是:最近八九
年,我在遥远的隐居处,身不由己地得知的那些;或是一场刻意保持得极有限的调查之中,自己撞上
我注意的那些。

这些数字本身已经雄辩,并且出人意料地印证了去年就已强加于我的那个印象:对我的著作、对我
这个人的\emph{葬礼},不是一个人的事业,也不是一个严格有限的集团(比如我离去前的学生,或我
那些“上同调学生”)的事业,而确实是一项集体事业——其水平,是“教团全体”的水平;或者至少,
是数学建制中曾经\dmark{}\origpage{1113}\textbf{见证}并参与了我1955至1970年间作为几何学家的著作之兴起与繁盛的那一
部分的水平。我1970年的离去,至少在这一片数学里,成了一个信号:对“Grothendieck式”
数学,一场即刻而严厉的\textbf{排斥反应}——人们把它感受为“女性气质的数学”的象征与化身
\fnote{850}{(*)}{关于对某种对待数学的进路风格的这类排斥反应,见“肌肉与肚肠(阳葬阴(一))”、“阴的葬礼(阳葬阴(四))”、“天赐的机缘——或封神”、“拒认(一)——或召回”、“拒认(二)——或变形”诸注(n\textsuperscript{o}s~106、124、151、152、153)。在“上涨的海……”、“九个月与五分钟”、“箭与浪”、“兄弟与配偶——或双重签名”、“阴,作为仆人——与新主人”、“阴,这仆人——或慷慨”诸注(n\textsuperscript{o}s~122、123、130、134、135、136)里,我试着勾勒“女性气质的数学”的某些显著特征,与互补的“男性”特征并列。}:在那
里,视野永远先于并启发着技术的层面;困难永远被化解,而不是被砍掉;与事物在表面参差之下的
深层统一的持续接触,让人每时每刻都能从偶然与次要的无定形团块里,认出什么是要紧的。与此同时,我的离去也是一个信号:一切概念性的工作骤然停摆——或者更确切地说,一切此类工作被
\textbf{褫夺公权},以“深化”为名,骤然打上嘲弄的烙印。

就这样,砍掉数学创造工作的一个本质“侧面”——“阴”的或“阴柔”的侧面——人们达到的是一种
惊人的“扁平化”、“削平”、“干涸”\fnote{851}{(*)}{关于这种“削平”的初步确认,见注“无用的细节”部分(c)“什么都不像的东西——或干涸”(n\textsuperscript{o}~171(v))。}。这件事(在我看来)是通过一个骤然而严厉的急转弯完成的,几乎一
夜之间。这件事如此奇怪、如此闻所未闻,以至于显得难以置信。我花了一年多对\emph{葬礼}的高强
度思考,才终于把握住发生了什么,才认了这份明摆着的事实。我不知道,近年或近几十年,或任何
别的时代,在科学的一支里、或任何别的动用我们创造之力的人类活动里,可曾有过可相比拟的转
向。

可我回到我的相册上来。除了方才点到的诸位,把另一些人对\emph{葬礼}的参与——在我看来毫无疑
义的——的名字收进来,我想是有用的。况且,我并不确信他们中哪一位对我怀有恶意;他们之中肯
定不止一位,\dmark{}\origpage{1114}甚至对我怀着同情乃至好感(回应着我自己心里相似的情感)。他们中也许
没有一个人,会不真心错愕地听说:竟然有一场“葬礼”在对我这个人、对我的著作举行;更不用说
听说自己据说还以这样那样的方式参与其中。在这里指名道姓,本身就会产生那个(对我而言受欢迎
的)效果:把这事告知他们;并且(如果他们自己关心)由此给我们之间一个说明的机会。当然,关于
我所感知到的——无论对错——他们对我的葬礼的直接参与、或经由“共殓者”的参与,我随时听候
诸位垂询,提供一切细节。我绝无意质疑他们中任何一位的善意与职业操守
\fnote{852}{(*)}{(6月16日)随着新近传到我这里的信息,这份善意推定,在Borel的情形里化为乌有了。照他去年与Zoghman~Mebkhout的通信——事由是Borel在苏黎世主持的一个 D-模理论讨论班——我早已知道:Mebkhout曾告知他,自己是那门理论(所谓“Riemann--Hilbert”)中心的范畴等价的作者,并给出了确切的参考文献,还寄去了他的全部著作,Borel本可轻易从中确证事实。这并没有妨碍Borel以照例的居高临下(甚至失礼)对待他。在刚于奥伯沃尔法赫结束的一场同题讨论会上(Algebraic~theory~of~Systems~of~partial~differential~equations,~Oberwolfach,~1985年6月9--15日),Borel做了头三个引言报告(题为“代数 D-模理论”),为“好上帝定理”铺路;\textbf{Mebkhout的名字,在这些报告里一次也没有被提起}——在随后的报告里也没有(Brylinski的报告里,除了一句顺带的“喊停式引用”之外)。照Mebkhout刚寄给我的会议纪要,这场讨论会——Borel在其中指挥乐队(顶替没有出席盛会的Deligne)——是四年前那场反常层会议的名副其实的\textbf{重演}。“黑手党”几乎全员到场:Verdier、Brylinski、Laumon、Malgrange,甚至(这一回)Kashiwara(此人在苏黎世那场讨论班里,早已扮演头等角色,尽管Mebkhout已就其为人向Borel作过周详的通报)。不用说(同苏黎世一样),没人觉得有必要请Mebkhout做个报告;而且(除了这位Mebkhout偶尔的插话——它们都掉进了冰窟)先祖的名字一次也没被提起(除了那个不合时宜的“Grothendieck群”里到底有他的在场)。双对偶定理在那里依旧叫“Verdier对偶”,连Borel的报告里也是。Mebkhout早在去年就再三提请注意:这份双对偶,是从 SGA~5 的第一讲抄来的——看来Borel对某种风格、对某位缺席者,发展出了一种过敏症,禁止他考虑这一类参考文献……他自己在《交叉同调》(Intersection~Homology,Birkhäuser,1984)里,也成了同一骗局的参与者——而那本书,是在Mebkhout向他点破Verdier的把戏\textbf{之后}出版的。
对Borel,我曾把善意推定保持到极限:五十年代我们同在Bourbaki,一起共事,我对他知之甚深。在真正可称为我在数学世界的“出身圈子”的成员里,他是第一位,我不得不在今日、以无可怀疑的方式,确认其直接参与——而且是在“诈骗”的层次上——葬礼。},其中不止一位,我甚至可以补一句:
他们的全然的善意与诚实,于我是不容置疑的。\dmark{}\origpage{1115}与其愚蠢地按字母序开一张清单(那种
事计算机比我强),我宁可按大致的时间顺序,列出为我的葬仪齐声合唱的信众们的名字:不是按他
们在葬礼仪式上露面的时刻(那多半我并不知道),而是按我清楚得知其参与的时刻。另一方面,我
要把我全体学生单独列出\fnote{853}{(*)}{我在这里说“我的学生”,指的是跟我做到博士论文层级、并且(除Deligne外)跟我做了博士论文的那些人。共十四位(其中两位在我“离去之后”),见“耶稣与十二使徒”一注(n\textsuperscript{o}~19)的巡视。}。除了Hoang~Xuan~Sinh夫人——她在越南工作,分明显然太远,
帮不上我的葬礼——我的学生之中,没有一位,不曾以这样那样的方式参与其中。此事我已在“沉默”
那条注(n\textsuperscript{o}~84)与“棺材1——或知恩的D-模”(n\textsuperscript{o}~93)的开头说明过;这里不是重谈的地方。对
于我的每一个学生,我认为最值得做的,是一份就发生了什么而作的深入的说明。

“为我的葬仪合唱”,分布在极不相同的调门上。我辨出了四种主要的调门;它们为我说成了一场第一
流的、气派十足的多声部葬礼!有对一切带着Grothendieck气味的数学之发展的“沉默而有效”的\textbf{抵制};
有\textbf{失礼}与不体贴——我\dmark{}\origpage{1116}离去之前,在数学世界里我从未遇到过的;在一两个极端的情形,它们采取
了几乎不加掩饰的\textbf{嘲弄}的形状。有\textbf{刻意无视}或贬低我的思想与观点在其个人著作中的影响、
或在当代数学某一部分中的影响——而在那些情形,这影响分明而关键;或把无可置疑地属于我的
结果与观念,安到一个第三者头上。最后,还有那种(“\textbf{鸵鸟}”)姿态:碰巧面对一场刺眼的骗局
的人,把头埋进沙子,装作什么也没看见、什么也没感觉到。

不用说,信众的唱诗班里,不止一位在同时吊着好几个调门的嗓子。

话说到这里,终于可以兑现那份答应过的清单\fnote{854}{(*)}{我没有把那八位“非上同调”学生的名字收进这份名单;连同上面已经巡视过的上同调学生的名字,它们可在已引之注(n\textsuperscript{o}~19)里找到。
同样,在我的“家庭相册”里,也理应列出我的同事或旧友之中那些我确知为“不予埋葬者”的名字——他们以绝无含糊的同情与敬重之表示证明了这一点。首先,关于我1983年所做的那份工作《追寻栈》(我还打算回头谈这份工作),我收到过热忱的鼓励,来自 J.~Benabou、N.J.~Baues、A.~Jouanolou,尤其来自Ronnie~Brown与Tim~Porter——他们(不止以一两种方式)在我做那份工作的整个期间给了我切实的帮助。
确实,这几位同事所属的环境,与我惯常自认同属的那个环境相当不同;而我那场气派十足的\emph{葬礼},也正是自然而然地发生在后一个环境里。作为属于或接近这一环境、并且最近(在过去一两年间)给了我同一方向之表示的数学家,我乐于在这里点出 B.~Lawvere、J.~Murre、D.~Mumford、I.M.~Gelfand,以及(并非最末!) J.P.~Serre。这最后一位,独享同时出现在两份“名单”上的殊荣——“埋葬者”的名单,与忠实朋友的名单!}了——它来充实我们的家庭相册:Eckmann、Dold、Campo(N.~A.~Campo)、Mazur、Poenaru、Epstein、Cartier、Quillen、Kuiper、MacPherson(R.~D.~Mac~Pherson)、Heisuke~Hironaka、Hirzebruch、Tits、陈省身(S.~S.~Chern)、Artin、Langlands、Rota、Goulaouic、Fulton、Borel、Tate、Serre。

\subsubsection*{c.\ 众人之上的那一位——或首肯}
\tocsubsub{c.\ 众人之上的那一位——或首肯}

(6月11日)昨夜敲出那份名单时,我觉得自己有点傻:那里老老实实排开的每一个名字,单是它本
身,就唤起一整团丰富的联想,而在这里全无踪影。可这里绝不打算一一停在那些名字上、停在它们
所唤起的东西上——那需要再写一卷,而我正急着结束\dmark{}\origpage{1117}这一卷!我向当事的人们致歉:如此草率地把
他们“贴”进一幅不太能启人深思的“(出席我葬礼的)合影”里。确实,他们中的大多数,已在《收获与
播种》的这里那里以这样那样的名分出现过,即便未必是作为我的葬仪的参与者。其中有四位,属于
我Bourbaki小组的朋友圈子;三十年乃至更久以前,我便与他们因工作而紧密相系(其中两位更有友
谊的纽带)。在这份石碑式的名单里,另有九位,我自觉以温暖的友谊之情与他们相连;那份情谊,到
我写这几行时还没有熄灭。可是,逝去的岁月里,不止一次,当我又面对这些昔日朋友中的某一位、或
我昔日学生中的某一位,一种奇怪的印象攫住了我:我心头那份依然完好的、向对方奔涌的同情之
流,所奔向的那个人,已经不在了——或者说,与那个人的联系已经断了,也许无可挽回地断了;另一
个人顶替了我从前认识的那个人——那曾是个充满热烈而颤动的生命的人——而新来的这位,似乎
已把他的一切痕迹抹得干干净净。那像是一场\textbf{干涸},一场枯竭;在一度有敏感而活的血肉之处,
生出了一层又硬又不怕水的壳……

在这本我才刚刚半开的家庭相册合上之前,我还想在匆匆插入相册的众人之中,单挑一位,稍稍多
停留片刻。他在这本相册里来得最晚。比之于我最终收进来的任何别人,在我这里都曾有(当然是不
自觉地)严肃的抗拒:不肯抛开一些关于我们关系的陈旧的现成图像,不肯认一份谦卑的明摆着的事
实。我说的是Jean-Pierre~Serre。

在《收获与播种》里,我不止一次有机会谈到Serre,通常指名道姓\fnote{855}{(*)}{在《愚妄与更新》里,有两三次我忍住没有点Serre的名;在那些时候,对我多少持批评态度的人,通常并不需要点名。《收获与播种》中我谈到Serre、谈到我们之间关系的最周详的段落,在“九个月与五分钟”、“兄弟与配偶——或双重签名”、“无用的细节”三条注里(n\textsuperscript{o}s~123、134、171(v))。}。这里那里说的那一点,想必已足以让人感到:他
在我的数学家过去里,扮演过一个无人可以相比的角色。这件事,在写\dmark{}\origpage{1118}《收获与播种》之前,我自己
也从未在上面停留过;它是我一路写下去才发现的。二十年间——从五十年代初,直到我告
别数学舞台的那一刻——他为我扮演了“得天独厚的对话者”的角色\fnote{856}{(*)}{1965到1969年间——Serre与我的关系依旧紧密——充当得天独厚的对话者的,更多是Deligne。缘由,想必在于性情上的强烈投合,尤其在于Deligne的一种开放态度(对我感到是我所当贡献者之精华的那种开放),而Serre常常缺这个。我下文还要回到这两种关系的截然不同的性质上——它们是我数学生涯中最紧密的两种关系。另见下一条脚注所引之注。};我的许多大的执念观念、大的投入,都是直接被Serre的思想(有时“貌似平淡”)激发
起来的。有某些时刻——我印象里尤其是五十年代下半叶,也许还有六十年代初——他与我之间有过
一种强烈的数学“共生”;我们的数学性情是互补的\fnote{857}{(*)}{关于这种互补性,以及Deligne与我之间的投合,见已引注“兄弟与配偶——或双重签名”(n\textsuperscript{o}~134)。};这种共生,每次都证明极其丰饶。Serre与我之间的关系,并不是“对称”的:比如,
Serre完全不像我这样,倾向于托付给一个或几个“得天独厚的对话者”,来了解他可能感兴趣、或他自
认为需要的东西。尽管如此(至少我揣想),我在他的数学家过去里,想必同样扮演过一个非同寻常
的角色;我不妨想象,我1970年那次不辞而别,在他的数学生活里,是一个断点(也许是某种平衡的断
点——在那平衡里,我代表着“阴”的一极),是一次骤然的转向,凭着一股突然出现的“空缺”。我不知
道……

不管怎样,Serre同我这个人、我的著作之间的这层紧密关系,数学界想必是知觉到的,哪怕它一直
在“不宣之列”。可以肯定:除Deligne之外,Serre被(有理由地)看作与我的著作“最亲近”的数学家。Deligne同我的著作、我的关系则全然不同——那是一种学生与“继承人”的关系。Deligne以我的思想、
我写下与未写下的著作为养分;而我的任何大的执念观念、大的投入,没有一样是由他唤起或激发
的。在某种意义上,他比Serre“更亲近”我:在他与我接触的那些年(1965--69)里,他身上没有Serre身
上的那种对我的著作、对我对待数学的进路的某些侧面的排斥反应;正是这一点,让他得以在仅仅\dmark{}\origpage{1119}
三四年之间(凭他出众的才能,也凭格外有利的种种机缘),把我此前数年诞生、发展的那套宏大的
统一视野,由内而外、整个地吸收进去。但他同我的关系,根子上是深深的暧昧——而且
他系统地玩弄着这层心照不宣的学生与继承人的关系:那对他是一种\textbf{权力}的手段;同时他拒认
这层关系,忙于把宗师、把宗师的视野一并埋葬……

Serre与我之间,没有这一类的暧昧——无论哪一方,任何时候,都不曾有过要从对方身上攫取“权
力”的念头、或把关系用作权力工具的念头。我甚至可以说,在我四十年代末加入的那个“Bourbaki
圈子”里,这类权力游戏并不存在;直到我1970年离去,我不曾目睹、更不曾(哪怕身不由己地)参演
过这类游戏\fnote{858}{(*)}{不过,我应当作一点保留,顾及某个完全在我不知不觉之中、在我某些学生之间、围绕我这个人我的著作进行的游戏。那游戏至少从1966年(SGA~5 讨论班结束之年)就开始了;第一个清晰可见的插曲,是Deligne1968年关于谱序列退化的文章(关于这一点,见“驱逐”一注,n\textsuperscript{o}~63)。我开始知悉这些把戏——它们确实是权力游戏——是去年,近二十年之后。诚然,积极的演员,并非当初接待我、我也自认同属的那个圈子的成员(在那个圈子里,即便凭着成熟给我的回顾,我依然看不出这类把戏)。他们组成了“接班人”。同样确实的是:我在这些接班人身上察觉的、相对母体圈子的品质败坏,与那个初始圈子的几乎每一位成员身上(几乎无一例外)悄悄发生的类似败坏,必然紧密相联。关于这一点,见“Bourbaki,或我的大机遇——及其反面”与“自深处”两节(n\textsuperscript{o}s~22、23)。}。另一种说法(关于Serre与我的关系,或关于我在Bourbaki圈子里
观察到的关系)也许是:任何时候,我在那里都没有察觉到丝毫\textbf{对抗}的成分——哪一方都没有
\fnote{859}{(*)}{不过,七十年代初“生存与生活”的插曲,应当算作例外。那一插曲让事情大白于天下:在许多我以为重要(今天依然如此)的问题上,我自己的伦理与意识形态抉择,同我的数学建制朋友们的抉择——Serre也在内——几乎全体处于对跖点。正是它,骤然终结了我同那个“建制”的认同感——我那时倾向于把它混同于一个理想(且田园诗般)的“数学共同体”。(关于这一点,见“‘数学共同体’:虚构与现实”一节,n\textsuperscript{o}~10.)这场意外的揭示,以及数月之内由此而来的“改换阵营”,当时使我采取了对某些旧友的敌对姿态;我从此把他们归入“反动派”等等。后来,我从这些斩钉截铁而肤浅的分类里退了出来。话虽如此,经过一个并不奇怪的翻转,Serre属于那样一些人:有一阵子,我把他们感知为“对手”,甚至“讨厌鬼”。我高兴地确认:这段插曲没有在他身上留下怨恨或敌意的痕迹——在我身上,也丝毫没有,这还用说么!}。偶有摩擦\dmark{}\origpage{1120},这确定无疑;前面谈到过,或许还要回头谈;但那是全然另一回事。Serre与我的关系之力
量,在我印象里,恰恰来自我们对同一位情人——数学——的共同激情,其间不掺任何唯我的“寄生”
成分:在其中,对方不是手段、不是工具、不是靶子。所以,最近当我同Serre恢复中断了十来年的通信,
我从他给我的两三封信的字里行间,重新找到完好无损的友谊与体贴的迹象,仿佛我们昨天才刚分
手。

再说,即便十多年没有机会通信,Serre偶一传来的音讯,全都朝着同一个方向:不变的友谊——而不
是葬礼的调门,后者是我许多旧友的情形。所以,直到最近的这几个星期,我都未曾想过:Serre会在
我的葬仪上扮演角色。凡从他那里传到我这里的一切、凡我知道的他的一切,似乎都朝着相反的方
向。而且确定无疑:单是他在数学舞台上的在场,就曾给\emph{葬礼}划定过某些界限(不得不承认,那
界限质朴之至……)。翻阅 J.S.~Milne那本《平展上同调》\fnote{860}{(*)}{普林斯顿大学出版社,普林斯顿,新泽西。正是这位 J.S.~Milne,两年之后,参与“难忘文集”《数学讲义》第900种的骗局(见“……以及掘墓”一注,n\textsuperscript{o}~168~(iii))。}(1980年出版,故而在那场难以置信的
“SGA~4½--SGA~5 运算”\textbf{之后}),我惊讶地看到:Milne“放心地”、几乎逐字逐句地沿用了Serre在
某次Bourbaki讨论班(1974年2月,n\textsuperscript{o}~446)上关于平展上同调之父的说法——理论“由Grothendieck发展,\dmark{}\origpage{1121}得 M.~Artin之助”\fnote{861}{(*)}{两年前,1978年赫尔辛基国际数学家大会上, N.~Katz(还是那个Katz)颂扬新科菲尔兹奖得主Pierre~Deligne的演说里,平展上同调理论被说成“由 M.~Artin与 A.~Grothendieck发展,方向为Grothendieck所预定”——可见字母序有时候还真管用……Milne在诸说之中选择沿袭Serre的说法而不是Katz,在我看来,是其善意的又一个征兆。}。有种种迹象表明,Milne只在 SGA~4 与 SGA~5 里零星翻过
\fnote{862}{(*)}{尤其令我惊讶的是:Milne(Mebkhout也一样——他可是我的著作的一位细心读者!)没有察觉 SGA~5 里有一条显式的Lefschetz公式,关于代数曲线上一般的上同调对应;那公式,被魔术师-骗子二人组Deligne与吕西耶妙手变走——说这是艺术家的活儿,再贴切不过!关于这一点,见“变戏法者——或飞走的公式”与“恭贺——或新风格”两条子注(n\textsuperscript{o}s~169$_8$、169$_9$)。};他一面沿袭Serre(在同一篇Bourbaki讲义里对 SGA~4 与 SGA~5 满不在乎地发议论),一面又沿袭Deligne(在其笔名“SGA~4½”的锯倒之作里,对这些讨论班肆无忌惮地拆台),为的是在他的引言里,把原
始文本 SGA~4 与 SGA~5 呈现为“难以进入”\fnote{863}{(*)}{至于 SGA~5 的出版版——托“编者”(原文如此)Illusie的“福”,它只是原讨论班一座毁容的废墟——Milne觉得它“难以进入”,情有可原。好撒玛利亚人Illusie竭尽全力(顺着好撒玛利亚人Deligne的心意),把它做成了一本难以消化的“技术性枝蔓”大杂烩……}——而这,恰恰是他的书(继Deligne三年前那本——说实话,边上
有点单薄——之后)理应来补救的情形;说得明白些:免得用户去费那无用而烦人的工夫,去读原始文
本。最高权威们的意见(这里首先是Serre,继之以Deligne,外加一位待在软壁棺材里一动不动、一声不
吭的死者……)——意见,一个Milne如每一个平常人一样闭着眼睛跟从(在殡葬的语境下,甚至抢着跟从……)——斩钉截铁地排除:这些文本会是别的什么,只会是“无用的细节”(或“废话矿
渣”……),而绝非一门新的“广义拓扑”(topos 版)的基础——这门新拓扑曾使(除其他理论外)Milne
书中所论理论的绽放成为可能;长远看,人们没法省掉这门新拓扑,正如没法省掉寻常的广义拓扑
——那种拓扑,Milne、Deligne、Serre(同我一样)有福气在学堂板凳上学过,所以他们顺从地(当作
理所当然)承认:那门拓扑的戏,值得那根蜡烛……

\pmark\origpage{1122}我想是去年,我才第一次匆匆看了Serre那篇Bourbaki讲义;我最近在“无用的细节”(n\textsuperscript{o}~171(v))
部分(a)“一摞摞千页的大书……”里谈到过它。Serre讥讽 SGA~4 共1583页的那一段,当时给我留下的
印象如此之淡,以至于一两个月前,为写《四种运算》重新拿起同一篇讲义时,我竟把它忘了个干
净。应当说,Serre对我那些出名的“千页大书”的这种拉开距离的态度,我早已深知——早在 SGA~4
系列出现之前——所以它毫无让我意外之处。我记得,这种“发自本能的排斥”第一次在Serre身上发
作、冲着我对待数学的某种进路,是围绕凝聚对偶理论——我在五十年代下半叶发展的。那才是潜
在的“千页大书”(至少),尤其若算上随附的一整套新的上同调代数——导出范畴版;但“大书”也罢,
潜在也罢,明摆着的一点是:Serre听都不愿意听,正如韦尔受不了看人白纸黑字写一个上同调群、
听人说出“拓扑向量空间”这几个字一样。

\pmark 然而这一回\fnote{864}{(*)}{其实,这已经只是我第三次手里捧起这份文本,它才“终于让我心头一动”。},当我重读Serre1974年的这份文字——底色则是一整年关于某场\emph{葬礼}的思考(那
葬礼,1974年已经“大行其道”四年了……)——那一段终于让我心头一动。它在我心中不动声色地起作
用,一天天、一星期星期。我意识到:Serre的那种态度——我最终习惯了它、它在我离去之前“并不
牵扯后果”——竟像是对那场葬礼的\textbf{一盏绿灯}。这一方向的头一件事,以明摆着的力量向我显形:
Serre的措辞本身(但“添上了恶意与无耻”),仅仅三年之后,就被一位Deligne急不可待地接了过去(或
更确切地说,怀着隐秘的快意),当作他那场难忘“操演”的“背景噪声”。

我在已引的那道注——5月4日——里第一次这样表态;这场思考在同注部分(c)(5月27日)“什么
都不像的东西——或干涸”里得到深化。那里,也是我第一次开启关于Serre与我之间关系的思考——\dmark{}\origpage{1123}
以\emph{葬礼}所提供的特殊照明\fnote{865}{(*)}{在前面一条脚注(打字稿第1117页注(*))里,我还点过另外两道注,其中我谈到过Serre与我的关系,但照明相当不同——那是“葬礼之前”的照明。}。写那几页时,我心中想必已经有一种弥漫的知觉:Serre在\emph{葬
礼}里扮演的关键角色。此后两星期,一场对整整一批事实与印象的整合与消化,想必在无声中进行;
抗拒对事物直接而细腻知觉的惯性之力,我相信,已不经战斗、不费力气地消解了。时机,在我看来
已经成熟,可以把知觉到的东西表述出来,尽我所能。

人们也许会想:Serre由来已久的、对我著作的某些侧面与部分拉开距离的倾向,不过起了一种阴差
阳错的偶然作用——可惜,恰好促成了一场同样阴差阳错的葬礼。然而那是一种肤浅的看法,碰不到
事情的深处。径直说到问题的核心吧:鉴于Serre同我这个人、我的著作之间独一无二的关系,鉴于他
对其同辈、以及其后各辈数学家的非凡影响,有一件事在我看来已经清楚:\textbf{若没有他对我葬礼的
心照不宣的首肯,这场葬礼原是举行不起来的}。

除了一位分明缺席的“死者”,这场\emph{葬礼}还有\textbf{两位主角}:他们的作为与不作为互相衔接、互
相补足,似乎没有丝毫摩擦与涂痕(但在我这里,并不打算谈什么串通——两位主角的调门实在太不
相同):他们是Pierre~Deligne,与Jean-Pierre~Serre。

关于前者,在这场关于\emph{葬礼}的长考的一开始,就已有长篇的交代:他代表着葬礼图景的“前景”,
作为葬仪大主祭,同时是那些由他倡导的运算的隐秘继承人与主要“受益人”(而且,早在死者的象征
性“死亡”之前就如此……)。Serre——这里第一次作为葬礼仪式的前排人物谈到他——则代表着图景
的“第三平面”,由“信众教团”构成的那一平面。

\pmark\origpage{1124}早在去年——或者更确切地说,早在我以它最粗粝、最悖谬的形状(以及这个名字)发现\emph{葬
礼}之前——我就清楚地知道:在一个我从不曾知道自己有仇敌的世界里,以这样的热忱来埋葬我的,
首先是\textbf{我昔日的}\textbf{朋友们};他们中的某些人,却并没有因此停止(哪怕只从唇边……)把自己算作
我的朋友——如今,我同样看清:在这些朋友之中,那些同时(而且尤其)是我的学生的人之外\fnote{866}{(*)}{在《收获与播种》的整个思考里,越来越清楚地显出:做过某个人的学生(比如做过我的学生)这个单纯的事实,就给一层关系打下烙印,赋予它一种特殊的质地,使它接近与父亲或母亲的关系。},真正
构成那场仪式的\textbf{台柱}的——作为教团的代表、作为全体信众首肯的担保人——正是他们中间
那位,在我们共同激情的层面上,曾与我最亲近的人。

Serre首肯的最耀眼的标记,当然不在我所熟知的那句随口的俏皮话里——那句话差点从我眼皮底
下溜走(虽说并没有逃过所有人的眼睛……)。真正令我驻足之后惊骇欲绝的标记,是\textbf{他所乐于
保持的蒙昧}:对那场就在他眼皮底下进行的葬礼——这样说毫不夸张——蒙昧到最后一刻
\fnote{867}{(*)}{可以说,在他那篇已引的1974年Bourbaki讲义里——他陈述Deligne对Weil猜想最后一角的证明——Serre的鼻子正埋在葬礼的正当中,却天真得连想都没想过记下一笔。我隐约感到他的不安:他发现自己面对一种表面荒谬的处境:我作报告(也在Bourbaki讨论班)给出某条 $\ell$-进上同调的 L 函数公式之证明的大纲,十年之后,我在其中承认(当作已知)的那条关键的“不动点公式”,居然仍未在文献中得到证明。Serre于是选择用一阵火气来\textbf{排遣}这份不安:讥讽 SGA~4 那出名的“1583页”(言下之意:而且\textbf{连}人们需要的那条公式\textbf{都没}提供)。这是一条省事的路:把一件不愉快的现实一笔带过(x)。然而他心里非常清楚(也许他乐于忘记……):在 SGA~5 讨论班里,我曾原原本本证明过一条不动点公式,其范围远远超出对 Frobenius 对应的那一条;他也知道,我的讲稿的成文,拖在那些所谓的义务“编者”手里,已经八年了。他即便乐于忘掉 SGA~5 的主题(“L 函数与 $\ell$-进上同调”——这标题好歹还是说了它该说的)及其内容,可他认识我二十多年,看我做数学看了二十多年,深知我把事情做半截不是我的习惯——恰恰相反(而且我做事情“不半截”到这样的程度,常令他不胜其烦,乃至忍无可忍……)。这本可以帮他刷新对 SGA~5 讨论班上所发生之事的记忆——他的脚,反正常踏进那扇门,足以大致知道我在那里做什么、进展如何。
显然,他既不想刷新记忆,也不想提问题。这只是众多例子中的一例:我的朋友宁可闭眼、塞住鼻子,也不肯去知悉一种他若接受就必然深深牵连他自己的真实。
(x)(6月22日)这几行写下之后,我有机会确认:这一类“不愉快的真实”,如今竟受到热烈的欢迎,如同天降的横财!关于这一点,见“家庭相册”的 d.、e. 两部分。}——那是一场对一部著作的葬礼;\dmark{}\origpage{1125}著作自其源头起就与他相系,比世上任何人都系得紧。而读《收获与播种》(假如他肯读这“一大摞”上千页的书的话……)会不会终于促使他动一动自己的鼻子(十五年来,馊味早就很大了……),以及其余的一切——对我而言,这是一个彻底的谜。可我深知:对他来说,同对葬仪的任何其他参与者一样(甚至更多),接纳我的讯息、动用他健全的官能,也意味着接受让他自己被深深牵连进去。

Serre的角色——率领那支前来观礼、为我葬仪合唱的信众教团——在我看来既典型,又例外。说它
例外,一是其极端性:他是与我最亲近的人,比教团任何成员都近;二是其非凡的身量(器局)
\fnote{868}{(*)}{还有第三种情况,赋予Serre在葬礼中的角色这种“例外”或“极端”的性质。他属于在我初次接触数学家世界时接待我的那群“善意的长者”。(关于这个群体,我有生以来第一次在“受欢迎的异乡人”(n\textsuperscript{o}~9)一节、继而在《收获与播种》引论(I~5,“一笔情愿的债”)里谈到过。)除了我们之间的友谊与好感之外,这也许是最主要的缘由,使得我花了一年多,才认了那份明摆着的事实,确认了Serre在我的数学葬礼中所扮演的关键角色。}:这
身量(器局)排除了深层动机中“补偿性的对抗”这类惯常的“寄生”成分
\fnote{869}{(*)}{这种“(表面上的)无因对抗”,我在这里那里已影射过两三次,尤其在(4月3日)“信使(2)”(n\textsuperscript{o}~182)一注(见下文)里。对我葬仪的绝大多数参与者而言——也许甚至是除Serre之外的全体——这种“原型式”的对抗肯定在起作用,这一点,我毫不怀疑。这股力量,在我看来,有别于那股通过“被自己否认的女人”之压抑(或“埋葬”)过程表达的力量。但这两种力量又紧密相联;在葬礼中它们结成一伙,以某种合金出现,常常难以分开。不过,我相信我已在其中辨认出葬礼中起作用的\textbf{两大力量}。可眼下,我很难说其中哪一方占先,是哪一方。我倾向于认为,是我先辨认出的那一方,即:对自己存在中阴柔侧面的压抑之力。
Serre的情形之所以在我眼前显得“典型”(同时又例外),想必是因为,在场两股力量中,正是后一种(我倾向于视为根本的那一种)在其中显出其全部力量,而丝毫没有另一方的痕迹(我称之为“寄生”的那一种——其义为:它会遮蔽对我所相信知觉到的\textbf{要害之物}的清晰把握)。我揣想(只要对原始事实与知觉的整合消化工作继续下去),未来几个月会带来更细腻的理解:在葬礼中、也在我以这样那样名分参与的其他冲突处境中,两股力量各占几分。}。如我方才强调的,我在\dmark{}\origpage{1126}Serre同我这个人、
我的著作的关系里,察觉不出丝毫对抗的痕迹;很清楚,在他的首肯之下起作用的深层力量里,也没
有对抗的痕迹。据我所知,除了那句出名的俏皮话,这份首肯纯然被动地表达:不过是一些\textbf{遗漏}
罢了。可是,这道给一场规模浩大的葬礼的心照不宣的“绿灯”——伴随的那些行动,有时庞大到仿佛
既定义了常识、又定义了体面——如今在我看来,恰是Deligne对同一场葬礼那积极而热切的参与的
不可或缺的\textbf{对应物}、某种意义上的“底片”\fnote{870}{(*)}{这里,Serre与Deligne在葬礼中的角色分配,有一处相当显著的\textbf{颠倒}:Serre的角色几乎全然被动,Deligne的角色则极度主动(虽说这位“操盘手”的角色,为事势所需、也合我朋友皮埃尔特有的风格,始终隐而不显)。然而,事实上,正是Serre其人,具有强烈得鲜明的“阳刚”主导;而Deligne其人,具有同样鲜明的“阴”(或“阴柔”)主导;并且(对两者皆然)无论在唯我机制的层面、“自我”及其条件作用的层面(亦即“\textbf{老板}”的层面),还是在发现之冲动的层面——那本源的、就其内在本性而言逸出条件作用的层面(即“\textbf{孩子}”的层面)——皆然。在Serre与Deligne这对截然相反的极端性情之间——葬礼的两大“台柱”——死者本人,则代表一种中间项:在“老板”一侧,阳刚主导强烈;在“工匠”(或“孩子”)一侧,阴柔主导同样强烈。(这一“基调”的分布,出现在注“兄弟与配偶——或双重签名”(n\textsuperscript{o}~134)里。)
阴阳角色之间的“对调”之力与机制,本是\textbf{主要}的思考主题:它催生了“阴阳之钥”这场长考,并作为底纹贯穿始终。它从《钥》的第一道注“肌肉与肚肠(阳葬阴(1))”(n\textsuperscript{o}~106)起就隐含在场,并在后续十一道注(n\textsuperscript{o}s~124、127、132、133、138、140、145、148、151、153、154)中或多或少走到注意的前台。而在这里,我意外地撞上了一个多少类似的“对调”处境——其动因,是葬礼中起作用的深层力量的内在逻辑。
近来,另一处似乎悖谬的阴阳角色“对调”,也打动了我——在这场不乏表面悖论的葬礼里!这一回,说的是“过早的死者”一方,与葬礼全体参与者一方,各自的角色。在集体无意识的意图层面上,这场对一个死者(按理,他只应安于与其身份相称的完全被动)的葬礼,首先是一场“女性气质的数学”的葬礼——一种带着强烈“阴柔”意涵的数学风格与进路的葬礼;而埋葬的教团,按理体现着“纯粹而坚硬”的阳刚价值,把那软塌塌的阴柔酥软之物交付给应得的蔑视。(比如见“阴的葬礼(阳葬阴(四))”与“天赐的机缘——或封神”两注,n\textsuperscript{o}s~124、151.)然而,处境的内在逻辑,却逼着这些“纯粹而坚硬”的参与者,在其中演一出典型的“阴”的或“阴柔”的戏:一场“天鹅绒爪子”的戏——半明半暗、沉默、遗漏、不动声色安插的暗示,或不断地暗示一事而佯言其反——总而言之,“喊停!”的那一套;我朋友皮埃尔是其中的至圣至师,而每一位埋葬者都不得不多少把它据为己有,事势使然。(关于这套风格,见“喊停!”一注,尤其“天鹅绒的脚掌——或微笑”与“对调(四)——或夫妻的马戏”两注,n\textsuperscript{o}s~77、137、138.)相反,是那位“死者”——那丰沛的阴柔酥软之物的化身——在人们最不提防的时刻,从舒适的棺材里坐起来,骤然接过一个他熟悉的“男子汉”角色:摊牌亮底,把好奇的鼻子和不逊的舌头、连同手电筒,伸进那些最精微暧昧的昏暗角落,粗声粗气地逐个点名,是猫说猫,是贼说贼——总而言之,一个真正的没规矩的家伙,一个不折不扣的搅局者,搅的正是那一团和气、细语呢喃的安稳局……}。而我朋友身上发生的这场变化,分量惊人。

\pmark\origpage{1127}我似乎确实深切地感到了在Serre身上起作用的那股力量。它所处的层次,比一种个人之间的对抗、或比谋求某种“好处”(按这个词的通常意义)更深。近来与他交换的那几封信,在这一点上很说明问题。说来,十五年来——自我离去以来的十五年里——我的朋友身上完成了一场\textbf{转变}
\fnote{871}{(*)}{“转变”(\textbf{transformation})这个词,立刻把我带到我朋友皮埃尔身上的“\textbf{变形}”(métamorphose)——那是去年10月他来访时,我第一次清楚知觉到的。(关于这一点,我在“拒认(二)——或变形”一注(n\textsuperscript{o}~153)里谈到。)“变形”一词更有力,对应着这一事实:在我朋友皮埃尔身上,发生了一次真正的翻转——从一种“阴”主导鲜明的本来性情,翻转到一副借来的、“彻头彻尾男子汉”的姿态。撇开这一点,我在两位朋友身上感到的转变,方向相同,动力相同:那股压抑“阴柔”特质的力量。},方向恰恰是那种对我对待数学的进路中的某些主导侧面的“发自本能的排斥”
\dmark{}\origpage{1128}。这些侧面,当年在Serre自己的进路里也曾存在,只是程度较轻——那是在他
作为数学家的过去里最丰饶的那些年:开放的、创造力澎湃的年月,尚未有一道\textbf{压抑}其创造个
性中这些侧面、压抑其内在“孩子”的工序立起来之前。这些,是创造力的“阴”的或“阴柔”的侧面与特
质。我在朋友身上以惊人的力度感到的转变,就是:从阴与阳两股创造之力和谐协作(而“阳”或“男
性”主导鲜明)的状态,转为一种“彻头彻尾男子汉”的失衡状态——“阴”的或“阴柔”的品质被毫不留
情地连根拔去。

说实话,正如两周前我已暗示过的(前引之注),这是一场演化的终点;五十年代,我已辨出它的最初
迹象,六十年代里它愈演愈烈。那时起,平衡就逐渐破裂,表现为视野的\textbf{窄化},表现为获准
上场的那份创造才能的清单的\textbf{缩水}。对我对待数学的某些主要侧面的排斥反应,以及渐渐地,
对一切真正构成我的著作的生命、深度与力量的东西的排斥——那不过是向外投射:是在他与我的
关系的层面上,一份远为深重的排斥的可触显形;被排斥的,是他自己存在的一个本质侧面,是他自
己的创造才能的一个本质侧面。

也许(如我方才暗示的),只要我还在跟前,同我的关系就对Serre的这一演化起着刹车的作用;在
他的生命里——五十年代,尤其六十年代——它代表着某种配重,由此成为一个相对平衡的因子。果真如
此,我的骤然离去,必定给那股压抑阴柔品质的力量放了行——那类力量,我并不陌生:它
也是作用于我自己生命的主导唯我之力之一;只是有一个显著的差别:在我这里,这份压抑之力只局
限于唯我机制的层面、我与他者的关系的层面,\dmark{}\origpage{1129}从不干涉我同数学夫人的恋情,也从不干涉(更一般
地)我在发现历险中的自发举步——无论数学的还是别的
\fnote{872}{(*)}{关于这股压抑之力在我自己生命中的角色,我在“超级爸爸(阳葬阴(二))”(n\textsuperscript{o}~108)一注里谈到过。我是在1976年开始察觉这股力量的;那一年,是我的心灵历险的一个关键转折。那转折,在“重逢(阴的苏醒(一))”与“接纳(阴的苏醒(二))”两条注(n\textsuperscript{o}s~109、110)里谈到。至于“阴柔”特质在我的数学工作中的主导地位(这些特质似乎在那里避了难,躲开了一切怀疑!),见“上涨的海……”一注,n\textsuperscript{o}~122。}。

回到\emph{葬礼}:眼下我做得最好的,莫过于引录11月10日那场思考的结尾几行——“阴的葬礼
(阳葬阴(四))”一注(n\textsuperscript{o}~124,打字稿第564页):
\begin{quote}
“……而这葬礼,骤然间以一种新的、意想不到的光呈现于我:连我这个人也变得次要了,我成
了那‘理应交付蔑视之物’的\textbf{象征}。这不再是某个人的葬礼,不再是一部著作的葬礼,甚至
不是一场不可容忍的异端的葬礼;而是‘数学中的阴柔’的葬礼——而或许更深一层,是在那众
多前来为悼词鼓掌的与礼者每一个人身上,\textbf{那活在他自己之内的、被自己拒认的女人的
葬礼}。”
\end{quote}
这最后一条直觉,是那一天、写下这最后两行的当口,以一道闪电突现的,像是一份意想不到的启
示,加在前几行所说的那份启示之上。此后几个星期,这道直觉一直作为底纹留在我的思考里,最终
在12月23日至26日连续的三道注里被重拾、深化:“天赐的机缘——或封神”、“拒认(一)——或召回”、
“拒认(二)——或变形”。

在这道直觉初次显现的那一天,以及在所引三道注的前两道——我在其中对它稍作探测——里,我
心里并没有一个具体的个案,除了(稍稍地)我的朋友皮埃尔(在第三道注里得到更周详的检视)。此
外我也深知:这一个案,对构成我\emph{葬礼}上出名的“第三平面”的信众教团之全体,绝无典型性。
所以,由于没有一个具体个案来作例证,我对那场突然瞥见的真实的把握,仍带着\dmark{}\origpage{1130}一层含混——被预感之物的那层含混:在某个层面上“知道”,却未曾充分而清晰地“看见”。我隐约
记得,我曾为这层含混稍感不安;曾有一阵冲动,想在我的朋友们——我知道他们参与了\emph{葬礼}
——找一个“有代表性”的,好把这份弥漫的认识“挂”上去,看它化身为一个可触的真实。

那时,Serre的念头,一次也没有掠过我的心头——他属于我旧友中那寥寥数人:我曾断然认定(至少
在意识的层面上),他至少不在参与我葬礼之列!而我那摸索着的思绪,当时(以及此前……)没有找
到那个在葬仪上须以一身体现“教团全体”的人,想必是因为,在我之内的某个地方,一清二楚:世上只
有\textbf{一个人}能演这个角色——而恰恰是这个人,我之内的一种沉重,让我以某种心照不宣而斩钉
截铁的禁忌,一开始就把他排除了……

如今,随着一场缓慢而幽暗的地下功课,那份沉重已经消散;现在,全光照亮了:恰恰是他——众人
之中恰是此人——那份寻找化身的直觉,在他身上以如此完美的方式落实,以至于人们几乎要相信:
正是他本人,让它在我之内浮现;正是他,从它显现的那一刻起,就赋予了它那种斩钉截铁、无可辩驳
的“已知之物”的力量\fnote{873}{(*)}{我甚至满心倾向于认为:这个“人们几乎要相信”,恰与事情的真实相合。这将再一次证明:我们的认识官能,远远超出那片可怜而微末的映像之域——我们只容许它们进入意识目光那片划得极窄的田地。(6月14日)前天的思考收尾的那道念头、那道骤然的直觉,也是在落笔之时以“闪电”方式显现的,没有明显的准备,也没有检视的打算。它带着某种“明摆着的力量”呈现在我面前。只是事后,我才记起:在紧挨着11月10日那段引文所出自之注的前一道注里,我恰好有机会相当周详地谈到Serre其人、谈到他与我的关系(而且,这还是《收获与播种》中的第一次)。}



\subsubsection*{d.\ 入殓——或自然的坡道}
\tocsubsub{d.\ 入殓——或自然的坡道}

\pmark\origpage{1131}(6月17日)现实,每一次都超出一切预感(无论多敏锐),甚至超出人们对它可能有的“知
识”;只有当我亲身与之摩挲——在路上的转弯处,而且往往全然出其不意——我才渐渐浸透它的滋味
与气味。即便这种接触似乎只是\textbf{印证}了已被预感或已被“知道”的东西,它却往往出人
意料,毫不客气地撼动一种几乎拔不掉的\textbf{怀疑}:对那确凿已知、已说、已写、已重述、又已重写的
东西的怀疑——尽管在某个层面上(那是巨大的沉重所在的层面),它依旧是一纸空文。我不止一次察
觉这份沉重\fnote{874}{(*)}{关于这份“沉重”、这份“不信自己健全官能之见证”,另见“义务已尽——或真话时刻”一注(n\textsuperscript{o}~163)第782至784页,尤其第782页脚注(**)。},我的不耐为之火起——一股顽强的沉重,执意要把我拽回熟
悉的观念与图像的车辙里,或拽回那些多少得到普遍首肯的观念与图像里;而我分明“知道”(或者说,
我之内的\textbf{别的}什么人、什么东西知道……)这些安顿妥帖的观念与图像是一场障眼戏,往往是明
摆着的障眼戏,根本站不住……思想,哪怕怀着要知道底细(那件既被“知道”又被否弃之事的底细)的
强烈渴望,单凭思想本身,无力抹去这份深深锚在自我结构之中的沉重。唯有与真实直接接触的斩钉
截铁之力,才有能力有时撼动这份沉重,侵蚀它、挪动它分毫——若谈不到真正抹去的话。

昨天我给Serre打了电话。只为一件求实的小事:关于Tate的“刚性解析空间”笔记——近日谈到的
\fnote{875}{(*)}{见“黑手党”一注(n\textsuperscript{o}~171)部分(c),“记忆的失灵——或新编历史”。}。我模模糊糊记得,那份文本原有一篇短引,交代这项工作的来历;我印象里,这篇引言在1971年
《数学创造》做的那一版里“丢”了。其实,Serre向我确认:Tate的笔记里,本来就没有这样的引言。那
不过是一些日志式的笔记,Tate把他关于刚性解析的思量寄给Serre,几乎像写信一样,而且(当然)
绝无发表的念头。我记得是我安排它们经法国高等科学研究所散发的(作者署名之下,加了一个小
标题“私人笔记,经(未经)本人许可刊布”\dmark{}\origpage{1132}),可我忘了Serre是中间人。反正,1962年刚性
解析空间诞生之时,除Tate与我,最“在局中”的就是Serre了。是他,或许早一两年,向我讲解过所谓
“Tate”椭圆曲线的理论——取一个完备离散赋值环的分式域 $K$。我被那——我记忆中是一场显式公
式的洪水(据说还是“经典的”)——砸得有点发懵;它们从我头顶掠过,没有“挂住”。但有一幅惊人
的几何图像留了下来,想必是Serre一句评点所致:归根到底,Tate椭圆曲线(至少它的“点”),是在
乘法群 $K^*$ 里“商化”、被一个同构于 $\mathbb{Z}$ 的离散子群相除而得。这正是复情形的类比:先把
$\mathbb{C}$ 除以一个因子 $\mathbb{Z}$ 得 $\mathbb{C}^*$,再除以一个因子 $\mathbb{Z}$,这次得到的就是一条椭
圆曲线。在复情形,商化在复解析的领域里有确切的意义,而Riemann-Serre型(即 GAGA 型)的定理保
证:最终的商(一条紧复曲线)具有典范的\textbf{代数}曲线结构。而在Tate的情形,唉,在 $K$ 这个完备赋
值域上那套让人不惯的解析空间里,得到的商是一个\textbf{全不连通}的紧解析空间,根本没有可能从
中取出一条椭圆曲线。然而(Serre当时必是这样对我说的)一切却仿佛……总之,Tate硬是靠着 $K^*$
和它的离散子群,凭着一手显式公式,造出了一条名副其实的椭圆曲线。

我记得清清楚楚:无论Serre还是Tate,都不相信Tate那套计算性的构造,会真有一个“解释”——以
$K$ 上一种“解析簇”的新概念来解释
\fnote{876}{(*)}{(1985年9月)正如去年7月我与Serre的通信所显现的,这里有一处我记忆的变形(Serre也有一处)。Tate(8月4日、10月16日)与我(8月18日、10月1日、10月19日)的几封信——都是写给Serre的——使事件的影片得以重排。是Tate(既非Serre,也非我)最先有了直觉与确信:应当存在一种“解析簇的新概念”,以便简单地解释“Tate椭圆曲线”的形式主义——那是在1959年8月前后。在我这里,它并没有“立刻让人心头一动”(照我记得的样子);我对Tate这一建议的第一反应相当怀疑,然后才开始思考这个问题。想必不久之后我就被说服了——一旦我察觉:既有的观念(尤其形式概形)无法解释与Tate椭圆曲线相联系的现象。此后两年,我相信,独自思考新概念的界定原理的,只有我一人;Tate与Serre都毫无头绪。这种情形持续到1961年10月:我向Tate提供了一门理论的总施工图。这立刻触发他去发展必需的基础,以便把握那些局部碎片(这项工作,若没有一个清楚的想法——日后如何把这些碎片拼装成整体的对象——本不会有太大意义)。更详细的评说与相关书信的引用,见《数学反思》第3卷预定的“历史评注”。}。

至于我,它当年是立刻让我心头一动的(*),我压根不曾想过,可以不把Tate曲线“看”作商化的结果
——为一个尚待提炼的合宜“簇”概念商化的结果——那正是让我着迷的那类工作!很可能也是Serre,
尽管满腹疑虑,向我提过:有一些人,至少克拉斯纳,在完备超度量赋值域——因而全不连通的域——
上“做解析延拓”。这似乎为我那个(有点怪诞的)希望添了水:不管怎样,应当会有一种“好的”解析簇
概念,比已知的更聪明,并且(凭“连通”一类的性质)接近实或复的、甚至代数的\dmark{}\origpage{1133}解析簇。可
再一次,在那个三人小组里,真正相信的只有我一个——至少我当时有这个印象。

这事在我脑子里转了好几个月,也许一年——说不清了。这处境让我想起一桩旧日的困惑:在当
时可用的概念框架里(靠赋环空间,概形与形式概形之类),人们无力赋一个形式概形——取所论的
离散赋值环 $A$——的\textbf{一般纤维}以意义。很快清楚:这本质上是同一桩困惑——我为给Tate的构造
以几何意义而寻找的那类“簇”,应当正是那种能给这出名的、尚不存在的“一般纤维”以意义的簇。
我终于有了第三条主线(克拉斯纳的风声之外):那是1968年出现的“广义拓扑空间”的直觉——当
时还没有得到“site”或“topos”这样的名字,因为我还没有开始边写边琢磨的概念工作——它将使我
得以定义那出名的、(隐含地)进入Weil猜想的“$\ell$-进Weil上同调”。这提示我:正如Weil上同调
的情形,我寻找的那种新“结构类”,不应到那老一套的寻常“赋环空间”里去找,或许要到那些“广义
空间”里去找——配以合宜的环层。

这些零零星星的直觉,何时才强大到足以推我开一个括号、从手头的任务里(主要是 EGA 与
SGA)抽出时间,动手做一点边写边琢磨的雏形工作——我说不清了。我知道的是:这项工作,一如既往
地,是在孤独中完成的——只有我一人“看见”那里有什么;也因此,只有我一人当时有能力做那把它
带到光天化日之下的最初的工作。我记得,我这里想几个钟头、那里想几个钟头,甚至有时一整天,
像逃学一般(虽然“日常”的工作并不缺!)。终于有一天,我咬紧缰绳,要弄个水落石出,真刀真枪地
干上了——我至少连续花了好几天,也许一两个星期。最难的是超越那些根深蒂固的思维习惯:它们
不停地要把我拽回已知的常规——“寻常”解析空间的车辙(我想,如今叫“松弛的”——德文作
“welk”)。我不得不重来三四次——一发觉自己又回到老辙里(马回马厩了!)就再爬出来\dmark{}\origpage{1134}。但分明,这\textbf{不是}那匹老马能顶用的时候……


这项工作做完,我心里有了底:再加一份我那时没有兴致去做的技术性工作,我已经立起了一个
“刚性解析空间”的观念(这是我给它的名字;“刚性”一词,为的是表达连通一类的性质——贴近代数
簇,而与所谓“松弛的”解析簇的性质恰成对跖)——它至少足以回应当时我心中的两个愿望:以这些
空间解释Tate的构造,以及解释形式概形的一般纤维。


我当时没有想得更远——我最急着回到那些被我暂时撂下的任务。倘若我再多玩一玩,我很快就会
发现:连闭环 $r\le r\le R$ 这样简单的空间(它们也配得一种“刚性解析”结构)都逃出了我的构造。是
Tate——我当然让他知道了我的思量——做了必要的调整,好把它们包括进来。除了我做的大半的
概念工作,还有一份更技术性的活儿要做:对所用的“砌块”(它们扮演仿射概形的角色)要有好的掌
握。这活儿,恰以其特有的优雅与细致,在Tate1962年的笔记里完成了\fnote{877}{(*)}{为求公允,我想可以说:我的工作与Tate的工作,代表刚性解析空间理论之绽放的两个阶段,同等不可或缺。我的那份,在于最初的视野(Tate与Serre都不曾有)与一份以概念为主的劳作——它也绝非没有必须迎面处理的技术侧面。Tate的那份,以技术为主,但其中同样有一份概念工作。我的工作,基调是“阴”的、“阴柔”的(所以,除了我已不在场,它还成了人人轻蔑的对象);Tate的工作,基调是“阳”,合乎好品味与好规矩的法度。\par 倘若我不曾把自己的工作推进到足以不再有任何怀疑的地步——怀疑一种好的“刚性解析”观念的\textbf{存在},以及对一门理论的总施工图的清晰视野——那么这门概念,恐怕今天还不会问世。因为,虽说人们“早晚”必然会发现并发展这个概念——它当然不是什么“发明”——其需要,在此后二十三年里,却从未以足够紧迫的方式出现,逼人“跨出那一步”。显然,是我第一个(1966年)预见到了刚性解析理论在最初的两个动机之外的应用领域:晶体上同调的发展。\par 此外,除我预见的三个用途——包括把Tate理论推广到一般的阿贝尔概形——我没有听闻别的几何应用。看来,后来“在这一行里工作”的人们,只从中看到关起门来发展理论的材料(既然理论已存在,而共识把它列入“严肃研究课题”),却不把它嵌进更宏大的几何视野。这是一个鲜明的案例:数学思想的原子化、碎块化——它与一切奠基工作的失宠相连,与一切不能归结为某种解“竞赛题”的技术绝活的失宠相连。一个特别雄辩的征兆是:无人尝试发展一个更一般的刚性解析空间观念——使Tate发展的那种,与“概形”之对“域上的代数簇”相仿——好把“变动的”完备赋值域上的诸刚性解析几何联系起来(尤其特征变动的,并包括实、复与“超度量”诸情形)。这一缺席,是近十五年数学在一切奠基工作层面上惊人停滞的众多征兆之一(这层面的关键,在此显然)。\par 回到Tate与我:同样可能的是,倘若我最初的“突破”没有在Tate那里“让人心头一动”、没有触发他的“第二回合”,刚性解析空间也不会存在。这事我原会在周围这里那里谈起的;可诱人的题目(包括看似更“紧迫”的题目)从来不缺,恐怕没有人会留意它——尤其今天,肯定不会:如今,引入如此怪诞之物的念头本身,就会让人想起某位,还是不点名的好……}

\pmark\origpage{1135}话说回来,还要过一阵,我才认了这份明摆着的事实:我所用的砌块,边边上短了点。它们对激
发我的那两个初始问题,已然够用——何必再往远里找!我很难改口。最终是Tate说服了我,以他那种
既平和又细致的方式:毕竟,不止这两个例子;就算我此生看来没遇到过圆环,也不构成排除它
们的理由。而显然,用我那套砌块,没有办法把它们“收编”(除非用无穷多块——那又差不多落回“松
弛”的老辙)。

我做好了我那份工作——在我自己的角落里,这很自然,反正我是唯一相信的人——不过当然,一旦走到(暂时的)终点,我也没有放过机会,把这事告诉两位主要当事人\dmark{}\origpage{1136}(实际上也是仅有的两位知情人):Serre与Tate。Tate那边分明心头一动;而Serre,我想,在我告诉他我的进展时,想必也被说服了。对此我没有确切的记忆,但倘若事情偏偏不是这样,我肯定会记得的。

\pmark 所以,昨天我给Serre打电话时,在我看来天经地义:他和我一样清楚,在这种新簇观念的诞生中,我扮演了什么角色。我没料到会有机会影射这件事;可当我跟他谈起Tate的笔记时,倒是他提醒我:那些笔记原封未动地刊在《数学创造》上,而且Remmert与另外两位作者刚出了一本专讲那些出名的刚性解析簇的书。这本书,就是我最近在“黑手党”一注部分(c1)(“记忆的失灵——或新编历史”)里谈到过的那本——我在其中痛斥Remmert“记忆的失灵”(可Tate的笔记本可以替他唤回记忆),服务于一种在我看来明摆着的蓄意刁难。我顺带向Serre提了一句——在给他写的上一封信里,我已有机会影射过某场葬礼\fnote{878}{(*)}{这是对那封信的复信里的话——那是我收到的Serre的最后一封信——其中Serre引了Siegel关于当代数学“扁平化”(Verflachung)的评语;我对它作了评说,并在“无用的细节”(n\textsuperscript{o}~171(v))部分(c)“什么都不像的东西——或干涸”里继续讨论。如我在该注所说,Serre曾把Siegel的这份印象斥为“不公正”——可我感到,Siegel竟这样想,多少让他心绪不宁。而这同一个词,他(肯定不是故意地)又拿来,打发我对一场葬礼的影射。不用说,他压根没有想到要问我:\textbf{什么}使我说有葬礼(我在信里一个字没提,宁可等他来问)。显然,案子早已审理完毕……},而这毫无疑问是一幅相当触目的写照。

头一件相当疯的事,是Serre(天知道他当年可是在最前排的!)——唉,连他也不记得了,而且是压根儿不记得了:我在那些出名的刚性解析簇里曾经有过一份!我简直瞠目结舌!这真叫人发疯——当我向他影射我以为自己有过的那份微薄份额、从激发我的那两个例子说起时,Serre记得的恰恰相反:差不多是我压根儿不想知道这些新簇,说(照他讲)有了形式概形,就已经有了所需的一切\dmark{}\origpage{1137}!当下我简直不敢相信自己的耳朵\fnote{879}{(*)}{事后回想,我明白了朋友的记忆(边上略有些失灵)发生了怎样的变形。既然我曾把形式概形当作主要的、几乎是唯一的向导,来提炼刚性解析空间的定义(以便给形式概形配一个刚性解析的一般纤维),他便(二十三年后)记成了:我曾斩钉截铁地坚称,不需要簇的新概念,因为“我的”形式概形已够用一切。(记忆的失灵,倒常常把事情“办”得挺好……)然而,连 $K^*$(我的第二条主线)也并非来自形式概形。反正,这里依旧是那句老话:案子早已审理完毕!}——然而不过几天前,我还刚刚以世上最平静的心情写下几页,谈的正是Serre在某场\emph{葬礼}里将要扮演的某个关键角色、某个“支柱”的角色。唉,这一回可好,我正正地站在\emph{葬礼}之中,就在我鼻子跟前、电话线的另一头,而且恰恰就是这位Serre本人,一如既往地从容自在,分明怀着世上最好的信念!(话说回来,我怎么也想象不出Serre会不怀好意,尤其在谈到数学的时候……)

我自然没有心思争辩,Serre更没有,不过到底还是有一场东拉西扯的闲谈,五分钟或十分钟。这十分钟真是用得其所:让我亲身摩挲到一场\emph{葬礼}的可触真实——色、香、味,一应俱全;那场葬礼,因为我一味只盯着纸看,已经变得有点遥远了。

我不得不说出的头一件事,就是“刚性解析空间”这个名字本身,是我给起的——还暗示,如果我未曾明说的话,那是在我还独自梦想着这些我如此命名的东西的时候。Serre有点愣住——显然,他同样不记得;但同样明显,我不是在编故事。不过不要紧,名字终究只是名字,而且毕竟“自然得很”……这句“自然得很”明白地暗示:正因为太自然了,它已不再意味着什么;任何人只要鼻子对着那件东西,都免不了就叫它这个名:“刚性解析”。总而言之,我的朋友就这个名字,给了我一句他并无此意的恭维——却配着“就这点小事……!”的调门。再说,关于这个,我什么也没发表过,不是么?那就无话可说了……

我越来越瞠目。发表也罢,不发表也罢,对我而言,现实分毫不改。一个女人怀胎九月,把孩子生到世\dmark{}\origpage{1138}上,孩子已经又蹦又跳、好好的,有人对她说:这不是你的孩子,因为什么也没发表,而且她也没本事
出示出生证——对说这种话的人,她必定当面嗤笑。说实话,我没有冲着Serre笑——那不是我的
做派;何况我还没回过神来。我也没有想争辩:连Tate本人在他的笔记里,都对我在理论起步中扮演
的角色没有丝毫讳言(此事,Serre看来同Remmert忘得一样干净
\fnote{880}{(*)}{我再一次感到:“无论如何,案子早已审理完毕”。既然Tate说他“完全忠实地”遵从我提供的一位总施工,那好——不过是个总施工罢了,一张潦草的图,随便哪个小把戏都能在沙子上画出来;一股淡淡的Grothendieck味儿,没错——Tate还算客气,真是够朋友,居然费神提了一笔……});而1972年,当我写“主题概要”\fnote{881}{(*)}{即那份日期为1972年的文本,对我当时数学贡献的一份略嫌干巴巴(也不太启人)的概览,为申请法兰西公学院教席而写(该席位后来给了 J.~Tits)。这份文本,加上更周详的历史评注,将刊于《数学反思》第3卷。《引论》之三(“罗盘与行囊”)里谈到过它。“主题概要”5 e)中,我写道:\textbf{刚性解析空间。}受“Tate椭圆曲线”之例、以及完备离散赋值环上“形式几何”之需的启发,我达成了完备赋值域上刚性解析簇观念的部分表述;它在其后 J.~Tate对这一观念的首次系统研究里发挥了作用。此外,我在特征 $>0$ 域上的代数簇上引入的“晶体”,有时可以解释为特征零域上某些类型的刚性解析空间上带可积联络的向量丛;这预示着特征 $p>0$ 的晶体上同调、与特征零的刚性解析簇上局部系统上同调之间,存在着深刻的关系。}——这还是“如假包换的Grothendieck”!——其中影射此事时,Serre也未曾有过一丝
皱眉(他的记性,想必从那时起就开工了)。反正都是白费劲——分明只要没有发表,我说什么都会被
当作刮过耳边的风……

可“没发表”三个字让我心头一动,我顺势接下去:恰恰,我的著作的很大一部分,就是没有发表的东
西,口耳相传的东西。我感到Serre又是一愣——在他看来,这想法似乎有点荒诞,像一种用词上的自相
矛盾:“著作——没发表……”,在他那里,这两者仿佛凑不到一块儿\dmark{}\origpage{1139}。我说出了“Motif”这个词;他
立刻扑了上去:这一回,他要纠正我那套葬礼观念,兴冲冲地向我宣布:两三年前,恰好出了一整本关
于 Motif 的书——真的,在“Motif”这一章上,我没有什么可抱怨的!

“那么,那本出名的书,你亲手捧过吗?”我问他(问得正巧——我琢磨着问他这个有趣的问题,已有
一阵了)。

捧过——但我也许是想笑吧,Serre回敬道;他当然熟悉那本书;他甚至像通读过全书的人那样谈起
它;想必他确实读过。倘若里面没有什么特别的东西,我本可以免了这一问——明摆着没有;可是(人
就是这样,我毫无办法!)我还是问了!见他似乎没听懂问题的意思,我告诉他:去年我捧起它时,我
几乎不敢相信自己的眼睛。

\pmark 我不得不说出了“诈骗”这个词,但我感到,那是一个纯粹委婉的说法。照我真正的感受——写
这几行时依然如此——那是一场\textbf{无耻之行};可我忍住没有把它说出口。说到底,我心里明白,用哪个
词都无所谓:十五年来,一直“馊得厉害”,而Serre选择什么也感觉不到(那正是我几天前刚写下的),
无论我说什么,都不会“通行”。

Serre这下打开了话匣子,仿佛他等的正是这个。诈骗?你做梦吧,可怜的朋友;可这本书正是Deligne
本人写的,还写得漂漂亮亮呢,没错——行,人人都很清楚是你引入的 Motif,可这不构成每次一说
到“Motif”这个词就要重念一遍经的理由,对吧?更不必说你从来没发表过一行字,而你的瑜伽靠的是
未证明的猜想(我觉得听见另有某人借Serre的嘴对我说话……);而那本书的妙处恰恰是:它不
用任何猜想,实际上,它不用你当年做过的\textbf{任何}东西……

语气斩钉截铁、不容回嘴,出自一位深知自己在谈什么、再无可学之人——还带一丝恼火:一个有
点赶时间的人,被一个死心眼的大胖子缠着,偏\dmark{}\origpage{1140}不肯明白最明摆着的事情。那气氛,容不
得人对任何东西稍作停留——一切早已定案、拍卖成交,分明如此。Serre的公理——关于这一行的伦
理、何为要紧、何为次要——分明变了;对此无可奈何。我只能连他的新公理一起,照单全收。

于是我无可奈何,只在“猜想”上接了话!我本可以对他说:Weil猜想也是猜想——然而他也好,谁也
好,从未想过把它们不当回事——但确实,那些猜想,Weil费心发表了!可我恰好在写“第六根钉子”
(钉向我的棺材)
\fnote{882}{(*)}{指的是那组注(n\textsuperscript{o}s~176$_{1}$至176$_{7}$)——我正在给它做最后润色;在那组注里,我拆解的恰恰是围绕 Motif Galois群观念、围绕Galois--Poincaré-Grothendieck范畴(应景地被封为“塔纳卡”的)的骗局——一场由Deligne策动、(最初)由“棋子”Saavedra经手的骗局……},于是我把话岔到“Motif Galois群”上:它可没有任何
“猜想”的成分;我发展过一整套高度精确的理论,关于Galois--Poincaré型的范畴——那正是那本出
名的书里用到的基础观念之一,而书里似乎认为无须对我这个人做任何影射。

Serre又扑向了这个影射;这一回,他又能纠正我的葬礼幻觉了——那套理论,白纸黑字发表在一本
书里,而且还是我的一个学生写的:Saavedra
\fnote{883}{(*)}{即同一位Neantro~Saavedra~Rivano的那本出名的“塔纳卡范畴”(原文如此),《数学讲义》第265种(1972),Springer出版公司。}
——难道不正是我让他做的那篇论文吗?这一回,分明又是一本他烂熟于心的书;他不止一次引用过
它
\fnote{884}{(*)}{我还知道:每当Serre有机会引用这本(几乎)不提我名字的书——而他本人并不觉得有什么不对——他总(出于我不知道的什么顾忌)小心地同时提一提我这个人。肯费这番神的,他该是最后一个了……}。“那么,那本书里,就再没有什么让你吃惊的了”——我又问他(而这一回,分明我早已知道
答案)。

没有,显然不曾让他吃惊:我的名字不在书里出现,无论是作为主题的作者,还是作为其中引入的附
属观念(诸如 Motif、晶体等等诸如此类)的作者——那些观念\dmark{}\origpage{1141}从头出现、作为例子展开,出自我
的又一位学生之笔,依旧没有我的名字。而这一回,Serre看来并没有记忆失灵——他还记
得(至少眼下……)那些观念归谁。倘若我这位朋友确有“失灵”,那总归不在“记性”这一层……

我们又就“塔纳卡范畴”这个名字聊了几分钟:我暗示,我认为它是一场骗局;而Serre有根有据,觉
得它再贴切不过。这一点,其实早在挑起这只新野兔之前,我心里就明白;我也知道\textbf{为什么}这名字,
对那位曾是我的朋友的人,如此贴切——而我这个孕育并分娩了这件东西的人,却觉得有话要说。

一如我们之间的常态,是Serre斩断了谈话——事实上,谈话确实拖得够久了。自始至终没有一刻“沟
通”;所以它留给我那份不满与不谐,想来正因为此。然而,正像最近收到的那两三封短信——以更强
的、无可辩驳的力量——这场简短的谈话,教了我许多。都是“知道”的事,没错,却被打了对折;知道,
却不信!而那份挫折感(至今没有消散),正是我在抗拒接纳那个讯息的征兆。

一个不速的讯息,当然。几个月前,我还毫不怀疑:Serre(一如我对他鲜活的记忆——犀利优雅、方正
不苟的化身),当他(迟做总比不做好……)凭借阅读“收获与播种”这份天赐的文本,得知某场\emph{葬
礼}的种种丑行(他当然万万料不到,可怜的人……)之时,必将拍案而起,投身这场战斗
\fnote{885}{(*)}{写“这一回”时,我想到了另外两次:我两次下本钱,试图向那个出名的“数学共同体”传递讯息——甚至动员它。第一次是1970年,我告别数学舞台之际,为了科学建制与军事机器的串通。第二次,层次低些,只在法国同事之间,为了一项打击在法外国人的不公正条文。(关于这一点,见“我的告别——或:外国人”一节,n\textsuperscript{o}~24.)两次,我的努力都遇到普遍的冷漠;Serre并不比我刚离开的那个圈子里的任何别的朋友(Chevalley与Samuel这两位除外)更例外。人们不妨就“收获与播种”这块大砖头,在那个建制里——Serre首当其冲——将产生(或不产生)何种效应,打个赌……}。这幅埃皮纳勒版画式的想象,在\dmark{}\origpage{1142}最近几个星期里消散了——多亏一次平淡的书信往
还。而昨天,我有幸亲眼看清,毫无半点疑义:Serre置身\emph{葬礼}之中,早已既久且深,而且他在其中
颇为受用。这一点——无须说(我也绝无半点讽刺)——是抱着世界上最好的信念的!

其实,我早已明白:“善意”绝非我大半生所觉得的那样,是一件如此简单、如此非黑即白的事。有一
种最常见的“善意”,不过是给自己递烟幕:一面体面的旗帜,遮着有时可疑的货物。我们的心灵由层
层叠叠的膜构成;目光越细,越看见某一层的“善意”,如何有时给底下一层的骗局充当掩护与托辞。

至于Serre的善意,我继续这样信任他:他绝不会写一本书、在实质上使用他人的思想而不明白说出
——哪怕那些思想从未发表、除了传达给他的人(倘若还在世)与他本人无人知晓。也就是说,我相
信我知道:Serre绝不会写昨天我们谈到的那一类书。我甚至可以说:单是\textbf{像Serre、像我}这样的人\fnote{886}{(*)}{在这里说“Serre或我”,我实际想的是:五十年代我们同属的那个圈子里的任何一位成员——那个圈子,我在《愚妄与更新》第三、四部分里试图略加勾勒,尤其在“Bourbaki,或我的大机遇——及其反面”一节里。不过,确实,我发觉即便在这个小圈子里,我也知道有两位成员“变坏了”(《收获与播种》里在适当的地方谈到过)。},
写一份面向公众的文本(这里是数学文本),就牵动根深蒂固的职业良知反射;它们(我相信)会消除、
或至少纠正某些“记忆的失灵”——这类失灵,在昨天那样一场随意的闲谈里,本算不得什么
\fnote{887}{(*)}{因此我毫不怀疑:倘若Serre是一本讲刚性解析空间的书的作者或合著者(Remmert正是),他不会由着自己滑上那条“自然的坡道”——把那位理当被略过的人略过;他会越过那些对那条坡道颇为迁就的“记忆失灵”——而他在一次私人谈话里,恰恰乐得滑上那条坡道。同样确实的是:若在十五年前,凭我当时认识于他的那份严谨,他不会任自己滑上这样一条坡道——我以为如此,哪怕只是在一次私人谈话里……}。而这一切,正合三星期前\dmark{}\origpage{1143}我在“什么都不像的东西——或干涸”一注(n\textsuperscript{o}~171(v),部分(c))里写下的话的方向:“我知道,Serre同我一样,绝不会想到跟着狼嚎,在‘人人’抢、骗、拆的地方去抢、去骗、去拆。”

话虽如此,我发现:这一切并不妨碍Serre——至少在某些情形里——实实在在受用于\textbf{别人的}抢、
骗、拆;而且是公开的、明摆着的,“在大庭广众之间”、“在聚光灯下”。他当然可以“怀着最好的信念”
如此——他不必脏自己的手:他只消对别人的掠夺、操纵与拆毁毫无保留地祝福;更何况他分文不取:
他不夸耀他人劳动的果实,却觉得别人(我几乎想说:那些领了专营执照的经销商)当众玩这种把戏
挺好。他收下的“收益”,比(边上有点不干不净的)出版物、比另一些人垂涎的银行账户,更精微。可
见它们毕竟举足轻重,足以造成我所认识的那个人令人瞠目的变形:今天(说不出从何时起),他闭
着眼睛、捂着鼻子,参与着普遍的败坏\fnote{888}{(*)}{这一对参与集体败坏的确认,与“掠夺的空白授权——或‘高尚行当’”一注(n\textsuperscript{o}~171$_{4}$,其名与其所指相符)——里所作的确认(针对1980年3月某讨论班的听众)相合;尤其打字稿第1090页第二段。}。

\subsubsection*{e.\ 最后一刻——或一个禁忌的终结}
\tocsubsub{e.\ 最后一刻——或一个禁忌的终结}

(6月18日)昨天我犹豫过:要不要给“家庭相册”那条注(n\textsuperscript{o}~173)再添第四部分,好“趁热”报
道前天与Serre的那通电话。诚然,那通电话留给我的是一种“不满之感、不谐之感”(如我昨天所写)
——而且这还是委婉语:那份刺痛逼近焦虑。这份不安催促我回到那段插曲,像对待一颗已经熟透的
脓疮:早该切开引流了\dmark{}\origpage{1144}。还有惯常的迁延。几个星期以来,朗格多克科学技术大学的复制室,就等着人
把《收获与播种》那册出名的分册IV的续篇送去——它迟迟不肯分娩;眼下时间已紧,就得抢印抢装,
赶在学院7月15日年度闭馆之前,何况不只我一家:学年将尽,各色论文蜂拥而至,人家有优先权。总
之,我对自己说:一本书总得知道收尾;若我继续往里塞“最后一刻”,明年这时候也完不了;够久了
……

可我终究还是动笔了——哪怕《收获与播种》的印刷要等到开学!它已经等了十五年(三十年也差
不多),再等两三个月无妨;只求容我从容看一看我该看的,说一说我该说的,不被种种“截止日期”推
搡……

那是艰苦的一天劳作——不如说一夜加大半个上午——我要让这份“加料”文本赶上今天的邮班。
事已成矣。

此刻,我有一种走到某项\textbf{必做}之工尽头的感觉。我忽然觉得轻了,仿佛卸下了一副我背了不知
多久、却从未察觉的重负。那想必是一种顽固幻觉的重量;它早在四十年代末,就在我之内扎下根来:
当时,一种领来的身份在我之内绽放——某个“共同体”(数学的)成员的身份、某个圈子的身份;那个
圈子,在我眼里满盛着温暖与生命。这种新身份的绽放,我在《愚妄与更新》里——“受欢迎的陌生
人”与“‘数学共同体’:虚构与现实”两节(n\textsuperscript{o}s~9、10),还有“Bourbaki,或我的大机遇——及其
反面”(n\textsuperscript{o}~22)一节——谈到过。诚然,这份认同,被1970年我离去前后的事件一扫而空,无影无
踪——那是在我投入战斗行动之后。可是,如今回顾,我意识到:对那个我已离弃、不再自认同属的圈
子,毕竟还系着一条\textbf{纽带}——一条也许无形、却力量巨大的纽带;它属于那份“过去的重量”(去年,
同名一节,n\textsuperscript{o}~50,我开始瞥见它)。我离开了那个圈子,没有回头的念头;可那个“家”\dmark{}\origpage{1145}——我为了
另一场历险而离弃的家——的一幅\textbf{图像},却在我之内活着,维系着那条纽带。这
幅图像,想必自我离去(以及更早,当然)直到《收获与播种》的思考,大体静止。《收获与播种》开始
修饰我对某段过去的图像,并勉力把现下的元素——常常令人错愕、常常不请自来——纳入其中。我
最终认了那份明摆着的事实:在我曾自认同属的圈子的后续、以及(据说)整个数学世界里,人心与
风俗发生了一场惊人的\textbf{败坏}。我意识到:这场败坏并非始于昨日;早在我离去之前,我就有足够的
时间,贡献过我的一份(至少,《愚妄与更新》的思考里瞥见过一份)。可是我有印象:我离去之
后,这场败坏有一场疯狂的升级;我的一些前学生,在其中扮演了头等的催化剂角色。

不管怎样——在对\emph{葬礼}的调查接踵而至的层层揭示之中,我心里始终维持着一种心照不宣
的“禁忌”,围着我旧日朋友中那些属于我青年时代那个圈子的人。我简直无法设想:他们中会有谁,
被这场我不断确认的深层败坏严重侵蚀、或“伤到”。当我说起“教团全体”对那些(至少在我眼里)超
出想象的行动的纵容时,我之内想必设有一道内在的“除外条款”,把那些在我看来应“立于一切嫌疑
之上”的人排除在外。他们分明毫不知情——他们想必在忙别的,肯定如此——可别怪他们!大致是这
样的调门。至于我的长者中最年长的几位,这种看法,我愿意相信,符合真实——或至少符合真实的
一个侧面。但对Serre、Cartier、Borel、Tate、Kuiper、Tits以及我熟知的其他人,却断然不然:他们与我
同代,正当盛年,完全融入我此刻检视的这个圈子,直到今天仍在其中行使着不小的权力、定着调
门,正如某些新来乍到者——他们在长者们的祝福之下,终于在那个圈子里结成了一伙毫无顾忌的
“黑手党”。

于是,在我据以看现实的那幅图像里——现实透过\emph{葬礼}这台头等的“显影剂”呈现——有一
道顽固而触目的矛盾。正是这道矛盾,无疑,\dmark{}\origpage{1146}在某个层面上被知觉、又在另一个层面上被否弃,造成了方才
说的那种逼近焦虑的“不安”——那焦虑,泄露了一场\textbf{分裂}。而那个人——比任何人都更能在我眼中体现这个圈子,体现这群被我内心的某个角色执意视为“亲近者”的人;而他本人,也曾是他们之中“最亲近”的一个——就是Jean-Pierre~Serre。正因此,被回避的矛盾之结,主要系在他身上。

\pmark 我开始怯怯地触及这道矛盾,不过六星期前,在“无用的细节”(n\textsuperscript{o}~171(v))一注的第一
部分(5月4日)。这场思考在同一注的第三部分(5月27日,即三周之后)——“什么都不像的东西——或
干涸”——里大为深化。一周之后(6月11日),我顶着顽强的内心抗拒,在本注的 c.部分(“众人之上的那
一位——或首肯”)里,重新回到Serre其人。这一回,Serre在\emph{葬礼}中的关键角色,终于大白于光
天之下。这是我对\emph{葬礼}的理解中的新的一大步——可矛盾的结,依旧没有触及!Serre其人,在我
这里(仿佛什么也不曾发生过)依然是无畏无瑕的“优雅”与“方正”的化身。“禁忌”完好无损!

是前天的那个电话,炸开了这道矛盾,把我的鼻子“直接摁进”(\emph{葬礼})里——不管我乐意与否。
随之,当然,是强大的抵抗之力(方才提到过的)的立即动员,以维持现状,而不是承认矛盾:以这样那
样的方式记录在案,由此解决它。我可以这样做,也可以不这样做。

\pmark 我跨出了那一步——而且我为此高兴。报偿是即时的:一场\textbf{解放},表现为轻快与如释重
负之感;一场内在张力的松解,诚然,但更是一场重负的卸除。

在《收获与播种》里,唯一可与之相比的解放时刻,同样标志了思考的第一次大转折:那是《愚妄与
更新》里的“竞技数学”,紧随其后的是“旋转木马收摊了!”(n\textsuperscript{o}s~40、41)。我还有印象:我方才“跨出”的
这新的一步,接续着去年那一步。何以如此\dmark{}\origpage{1147}、在何种意义上如此,我一时说不好。当时那声凯旋的欢呼
“旋转木马收摊了!”,肯定是过早了(次月我就察觉了)。可我刚迈的这新的一步,至少是又一步,把我带离那
台出名的旋转木马更远一步。将来会告诉我,实际情形是否如此。

经过昨天与6月11日的思考,我觉得已抵达一幅对\emph{葬礼}不那么含混的图景。此前,尤其
是这个“第三平面”模糊不清。6月11日的思考,让它“以可触的方式”化身于Serre其人;而这具化身,又
在昨天的思考里,取得了(这样说毫不夸张)一切具体的轮廓。

\pmark 说到底,在《收获与播种》这整个第四部分里,关于我与Serre关系的思考,在我看来,是对我理
解\emph{葬礼}最关键的部分——在“侦查补遗”与那幅数学大都会阴沟的浓墨重彩的图景之外。诚然,
倘若我不曾出于对我交给自己的题目的敬意,以我所能的全部细心,硬着头皮做这份“侦查的归置”,
并且尽力照亮沿途出现的每一个稍暗的角落,关于Serre的这场思考,恐怕也不会问世;而我对\emph{葬
礼}(以及我自身对它的牵连)的理解,将依旧如前含混。研究的工作里,一切互相维系!

\pmark \emph{葬礼}最后部分里那场思考的最有分量的部分,实际上是“最后一刻”才出现的。照理,这一
部分的“句号”,两个半月前(4月7日)就已画下。当时只剩十来页要重抄清稿、补几条脚注(一年前五
月末,情形也一样……)。意外之事,随即接踵而至:Zoghman来访——来读这(原则上已完成的)最后一
部分、向我提意见。它们凝结成了大约三百页的追加文字——其中,就有这几页:在葬礼(至今被回避)
的光照之下,我重新审视我与Serre的关系。


% ================================================================
%  18.5.7 -- 18.5.9 (原书 p.864–893)
%  源文: _work/chunks/chunk10.md, chunk11.md, chunk121.md, chunk122.md,
%        chunk123.md, chunk124.md
% ================================================================
\subsection{升级(2)}   % 18.5.7. L'escalade (2)

\notehead{174}{}

\dmark{}\origpage{1148}(3月22日)\fnote{889}{(*)}{(6月14日)眼下这道注,接续前一道注(同日所写)的 a.部分(“一位备受簇拥的死者”)而来。}我在别处已经强调过:其实并没有四种运算(对\textbf{一场}葬礼而言),只有
独一无二的一场“\textbf{葬礼运算}”。把它分成四大块,不过是为了叙述方便,其实是人为
的,而且(若太按字面当真)容易误导。确凿无疑,在那位导演——乐队指挥——葬礼主祭的身
体里,从来没有\textbf{四}只小魔鬼蹲在脑袋的四个角落里给他递台词,只有独一个!在关于阴与阳
的那场漫长冥想里\fnote{890}{(*)}{这就是构成《收获与播种》第三部分主体、Notes 104--162$''$的那场思考。},我曾试着与这只小魔鬼比从前更多地相识——从前,我不过时时察觉
它还在那儿折腾,转眼又去管别的。不过我并不声称已经充分认识了它;说到底,认识它或许根本
不是我的分内之事。确定无疑的是,它还像从前一样折腾个不停,而且未必会在我朋友咽气之前收
场。话虽如此,那场著名的“葬礼运算”至今仍在进行——就在我写这几行的此刻仍在进行。我不知
道,眼下这本“家庭相册”的散发,能不能至少终止全部局部运算中最庞大(也最不义)的那一项:
活埋一位年轻数学家——Zoghman~Mebkhout,“所有”在代数簇或复簇上同调里做工的人,四五年来
一直在用他的思想和结果……

既然分明只有一场运算,我就放下“四种”运算的虚构,不妨按时间先后,把我知道的主要插曲
与阶段勾一个轮廓。这里我不做了——我自认为已经做够了:在前四条主注(“沉默”、“操演”、
“分赃”、“封神”,n\textsuperscript{o}s~168、169、170、171)里,我把所知的全部插曲都收拢在一处,好奇
的读者自会把它在时间的标尺上排定次序。说来奇怪,若从“第二层次”或“运算”的观点(容我
用用这些委婉语)看,我1970年告别数学舞台的那一年,在这串插曲的接续之中竟不构成什么断
裂:在我印象里,它们自1966年 SGA~5 讨论班结束起,一直以相当均匀的节奏延续到1977年——
那一年,“SGA~4½--SGA~5”双双出版,SGA~5 又出了Illusie的\dmark{}\origpage{1149}编订版\fnote{891}{(*)}{(6月3日)这个印象应当修正,须把那场声势浩大的“Tannaka 范畴”(原文如此!)运算考虑进来:其第一插曲(与“稻草爸爸”N.~Saavedra搭档)在1972年(尾声则在1982年,由“真爸爸”Pierre~Deligne接力)。关于这一点,见那组注“第六根钉子(钉进棺材)”(n\textsuperscript{o}s~176$_{1}$--176$_{7}$)。}。这场运算,我看作一个骤然而惊人的\textbf{质变}。之前,是一次不动
声色的“收割”。从这里,我感到一阵暴烈与蔑视突然发作,死死缠住一个已被宣布“亡故”的缺席
者的作品。

在我那班上同调学生的这场上上下下齐心的集体发泄之后(“教团全体”在旁边看得心满意足),
似乎有过四年的平静。而从1966到1977的十一年间,我隔一两年就能辨认出一个类型分明的“插
曲”;1977到1981(“反常层会议”之年)之间,我却一无所知。相反,Deligne的长文《父亲猜想,
II》(La conjecture de père,~II),1980年发表于《数学公刊》——就在那场令人难以置信的讨论
会的前一年——照眼下的世道,几乎可以算作正常了……\fnote{892}{(*)}{当然,这项工作的主要结果没有一字提到我这个人;而那个结果的表述,本就是Deligne从我这里拿去的 Motif 瑜伽的一部分。反倒有一件事让我吃了一惊:我的名字,与Miller的名字一道,出现在文中一个段落里,说起那个除幂德拉姆复形——它大约在1976年由Miller和我各自独立地引入。1976年我曾在法国高等科学研究所就这个题目做过一次报告(顺便说,那是我一生中最后一次公开演讲),但明摆着我什么也不会发表。换了旁人,把这份纯属非正式的共同父亲身份略去不表,想必谁也不会察觉,更不会有人有半句闲话……}。也是在这一年,Deligne在一次Bourbaki讨论班上,从作者本
人嘴里得知了“好上帝定理”(又名Mebkhout定理)\fnote{893}{(*)}{(6月1日)其实,这一幕发生在前一年,即1979年6月,就在Bourbaki讨论班上。}。这是长期停滞的上同调主题冰雪骤消的开端。也是信号:第二年,葬礼
运算的第二场也是最后一场(?)高潮就此发动,这一回的调门是不义的,一切克制,连同最起码
的审慎,都被轻快地抛到了九霄云外。

翌年“难忘文集”LN~900那一幕(它为 Motif 掘墓复生,只字不提我这个人——去年某个4月19
日,这一幕曾让我百感交集……),还有同年Berthelot的报告(它把卑不足道的我从晶体上同调的
“历史”——原文如此——中一举清除),事后看来,都像是\dmark{}\origpage{1150}那场讨论会所发生之事的自然延伸,总
的说来也算得平淡无奇;那场讨论会的名字(或它的残余)也许会进历史,作为一次\textbf{警告}。
而翌年的“悼词”——无论在稍稍“摆一摆”它的人看来多么难以置信——同样像是这样一种延伸,
或者(如我先前所写\fnote{894}{(*)}{见注“宝石们”(n\textsuperscript{o}~170(iii))。})像一篇“尾声”。至于此后又过去的两年,它们不过是把一场光辉讨论会及其种种延伸
的“既得成果”,在文字里、在人心,一一盖棺论定……

这是一个显著的巧合——或者说,这分明\textbf{不是}“巧合”的效应——早在去年,在我还未见识
“SGA~4½--SGA~5”那场运算、也未见识反常层会议之前,我已经在我朋友皮埃尔与我的私人关系
中,辨认出恰好落在1977与1981这两个年头的两个“转折”。我第一次把它们纳入同一道目光、试
图探一探它们的意义,是在4月25日的注“两个转折”里——就在我发现\emph{葬礼}(读到那本难忘
的 LN~900)之后六天。两个转折发生的当时,早在多年之前,我远没有料到有一场\emph{葬礼}正在
暗中操办(至少在意识的层面上一无所知);要把其中任何一个同我所知的某件能照亮它的事联系
起来,我也会一筹莫展。

\subsection{殡仪馆——“为科学服务”}   % 18.5.8. Les pompes funèbres - “Im dienste der wissenschaft”

\notehead{175}{}

(3月23日)要把这场“葬礼运算”巡视完毕,还剩最后一位卖力的积极参与者的角色要看一看——
在这场关于葬礼的长考里,我已有机会屡屡“顺带”谈到他。这里说的是那家令人尊敬的Springer
出版公司(Springer~Verlag~GmbH,海德堡),科学书籍与期刊的出版者,大名鼎鼎——而且
还以“Im~Dienste~der~Wissenschaft”(“为科学服务”)的箴言而自豪\fnote{895}{(*)}{(6月1日)向Heinze博士打听之后得知,这其实算不上一句“箴言”,倒不如说是一条广告用语,英文作“Springer~for~Science”。}。这家出版社的数学书里,最负盛名的,大概要数“数学讲义
LECTURE NOTES IN MATHEMATICS”丛书。它恐怕也是世界上运气\dmark{}\origpage{1151}最惊人的科学丛书:二十来
年间出了一千多种。我想,我也为这份空前的成功出过一份力:这套丛书草创之初,正是靠我在六
十年代直至七十年代初,陆续放入的一大批我自己的或我的学生的著作给它撑的场子。我也曾以
“数学基础”(Grundlehren~der~Mathematik~und~ihrer~Grenzgebiete)丛书联合编辑的身份,与
Springer有过交往;三本书(包括 EGA~I 的重印本)经我的手在那里出版\fnote{896}{(*)}{另两本是Jean~Giraud与Monique~Hakim的学位论文(一是“场”与非交换 $1$-上同调的形式主义,一是任意赋环 topos 上的相对概形)。}。

1970年我告别数学舞台之后,我停止了一切编辑活动。仅凭惯性,我一直挂着那套丛书的编辑之
名,直到去年,才终于“正式”从Springer的一切编辑职务上退了下来。促使我这样做的是两个一
致的理由。其一,正当我重归“正统”的数学活动、重新发表数学著作之际,我要给这个“回归”划
出明确的界限:它对我绝不意味着重回某个“powerstructure”(权力与影响的架构),而只是回到
一份以发表为目的的\textbf{个人}数学工作。其二,自1976年起(经由Yves~Ladegaillerie
学位论文的那段插曲),我早有机会嗅到某种“葬礼的气味”——那时,对去年才发现的那场大规
模行动,我还毫无所觉。(关于这篇学位论文的插曲——它是我有幸启发过的最出色的论文之一——
见注“进步是挡不住的”(n\textsuperscript{o}~50),尤其是原委更详的注“二号棺材:或电锯式切块”,n\textsuperscript{o}~94.)这让我懂得:
“我所喜爱、我所愿意鼓励的那一类数学,在Springer出版公司里已经没有位置了”\fnote{897}{(*)}{这句引文(经翻译)出自去年2月18日致Peters博士的短信,信中我告知了他我退出“数学基础”编辑委员会的决定。Peters博士其实当时已经离开Springer(他如今在 Birkhäuser 出版社),通信便改由主管“数学基础”的Heinze博士接续。我曾要求把我那封信的副本寄给“数学基础”的每一位共同编辑(共十八位),并两次(1984年4月与1985年1月)向Heinze博士重申此请,他却始终认为无须告诉我副本究竟寄了没有(看来是\textbf{没有}寄)。我于是费神自己把信的副本一一寄给十八位编辑,并附了几句说明。他们之中我亲自结识七位,五位算得上我的朋友。只有一位(M.~Artin)费心回了信;其余诸位看来都没觉得有什么不对头——哪怕仅仅就\textbf{他自己}而言——:Springer公司居然没有费心把那封信(早在1984年2月)寄给他。};而且更重要的也许是,我在那里感到的气息\dmark{}\origpage{1152},并不鼓励我延续或重续与这家
公司哪怕稍微密切一点的往来。自去年2月我递交“数学基础”编辑委员会的辞呈以来,
过去的一年,只是把这份感受一再证实、一再加深。

不过这些都还是这场“葬礼运算”本体的边角——我昨天所说的那个“第二层次”,如今该回到它
了。据我所知,有\textbf{五本书}直接与这场运算相关\fnote{898}{(*)}{(6月1日)这几行写下之后,事情表明,还应在下面的清单上添入第六本书,它的书名本身就是一个骗局:“塔纳卡范畴”(Catégories~tannakiennes),Neantro~Saavedra~Rivano著。耐人寻味的是,这本书也出了Springer同一种“数学讲义”丛书。不过在那一桩运算里,Springer公司的责任在我看来并未牵连在内,不像其余五卷那样。至于“塔纳卡范畴”那场运算的细节,见那组注“第六根钉子(钉进棺材)”(n\textsuperscript{o}s~176$_{1}$--176$_{7}$)。}。按出版先后,它们是:SGA~7~I(1972年以我的名义出版)与 SGA~7~II(1973年
以Deligne-Katz的名义出版),整理的是1967/69年的单值群讨论班;题为“SGA~4½”(Deligne
编)的一卷,以及 SGA~5 的Illusie编订版(以我的名义出版,1977年);最后,是1982年以
Deligne--Milne--Ogus--Shih四人的共同署名问世、专为 Motif 掘墓复生立碑的那部“难
忘文集”。耐人寻味的是:这\textbf{五}卷皆出自\textbf{同}一家公司,皆在同一套“数学讲义”丛书\fnote{899}{(*)}{即第 288、340、569、589、900 卷。};前四卷出版时,主管“数学讲义”的是Peters博士\fnote{900}{(*)}{如我在倒数第二条脚注里指明的,Peters博士后来已离开Springer,转投 Birkhäuser 出版社。},最
后一卷出版时,主管丛书的是Byrne夫人。

这五次出版,其条件在我看来粗糙地不合规矩。我在别处已经指出:SGA~7~I 与 SGA~5 两卷
\textbf{以我的名义}于1972与1977年出版(LN~288与589),Springer却自始至终认为无须同我
打一声招呼,来征求我的同意,或仅仅把出版计划通知我。题为 SGA~7~II 与“SGA~4½”的两卷
的出版——它们打的也是 SGA 的旗号,而我以为这个旗号绝非人人可用,它与我的著作、我\dmark{}\origpage{1153}这
个人的干系,人所共知——事先同样没有征得我对使用这一旗号的同意,尽管(按理人们满可以期
待)我本该作为作者、或作为该卷(或它所整理的那个讨论班)的主编(或主编之一)列名其中
。最后,LN~900 一卷,在没有我的名分之下,呈现了那些概念、思想与构造;而凡消息
稍稍灵通的数学家都知道,这些是我引入的东西。在这一桩,明摆着(根本无须是 SGA~5 或 SGA~7
讨论班的那几个入室弟子之一)该卷就是人们通常所说的\textbf{剽窃}。我当然不指望Byrne夫人——出错
的话请指正——在该卷出版时主管“数学讲义”的那位——能凭自己的本事,看着稿子识破这场骗
局。但我想,一家严肃的出版社的本分,本该包括保证其出版物的严肃性,办法就是给自己配上胜
任的顾问。

这些顾问若诚实地做他们(想必)拿了报酬要做的事,也完全来得及向有关方面指出:SGA 这个
符号不是什么人都可以打的旗号,它自有\textbf{含义},而要尊重这含义,就该请教唯一有资格决定这个
旗号用法的那个人——也就是我自己。最后,就那卷顶着障眼法旗号“SGA~4½”的出版而言,作为
加重情节,只消翻一翻该卷的引论,或其后的那根“阿里阿德涅之线”,或第一章的引论,就能看
到:其中对待 SGA~4 与 SGA~5 两个讨论班,是何等漫不经心的轻蔑;而在消息稍稍灵通的人当中,
人尽皆知:那两个讨论班是六十年代中期的事,而顶着“SGA~4½”名目的一卷,装的却是七十年
代的伪作。依我看,对一个消息大体灵通、心智健全的人来说,这场骗局只能是一目了然的。因此,
更有十足的理由,在出版这样一卷、用上这样一个名头之前,先按正规手续征得我的同意。

依我看,在这五卷的每一卷的出版里,Springer出版公司的责任都牵连无疑——这五卷,正是围
绕我的著作、打着上同调主题所做的那场可歌可泣的诈骗大业中,一幕幕醒目的插曲。经由这些出
版,Springer公司成了这场非常行动的帮手与\textbf{贩运者}。我当然不能断言这是明知故犯。但我可以
说,这家公司在与我的关系上自\dmark{}\origpage{1154}1976年以来对我屡施的种种失礼(1970到1976年间,我想我没
有同它打过交道),也正朝着这场行动的方向,并且嵌入某种\textbf{精神},与它密不可分。

在注“操演”的子注“驱逐”(n\textsuperscript{o}~169$^{1}$)里,我曾提到我致Byrne夫人论 SGA~5 出版一事的信,和
她的复信——那复信,我必须说,真叫我目瞪口呆。(在这场“为科学服务”的光辉行动里,我“目
瞪口呆”既不是第一回,也不会是最后一回……)我从她的信里(日期是1985年2月15日)得
知:按照“一部著作包含多位作者的贡献时的惯常做法”(原文如此),无须专门知会我——我不
过是讨论班的\textbf{主编}罢了……SGA~5 的五位“作者”是Bucur、Houzel、伊
利耶、Jouanolou、Serre,独独不提我这卑不足道的人:我只以“主编”的身份列名——
想必纯属荣誉虚衔——为了这台光辉的讨论班,我说得未免太多\fnote{901}{(*)}{在那卷题为“SGA~4½”的书里那根著名的“阿里阿德涅之线”(贯穿 SGA~4 等等)之中,没有任何东西能让读者想到,我竟有幸在 SGA~4 与 SGA~5 里做过报告(相反,“有幸”为“SGA~4½”“合作”,倒是有的……)。关于这一点,见我在注“双重含义——或行骗的艺术”(n\textsuperscript{o}~169$_7$)里的观察,打字稿第899页。}。

早在收到这封富有教益的信之前,等得心里发闷(整整一个月杳无音信),我就拿出了我最漂亮
的笔(用德文),写信给 K.F.~Springer博士本人——公司主管负责人之一。那是一封漂亮的两页打
字信,我向他解释:公司与我的关系中绵延已久的一长串不快令我深为痛心;此外,还有若干起
对我不合规矩的粗暴做法,眼下我姑且奉上两桩,在我看来格外触目:“数学讲义”两卷(LN\textsuperscript{o}s~288、589)以我的名义出版,却认为无须征求我的意见;这两部书里,我在讨论班口头讲过的
思想、方法与结果,被删削,被残害,有时竟至于面目全非;这一事实,与环绕这两卷出版的种
种非常情形,凑在一起,对我而言绝不可能出于纯粹的偶然;我期待Springer公司公开地、毫无
保留地道歉,形式容当后议,一俟原则达成一致再行商定;我盼着他同我一样,把了结这一令人
不快、令人无法容忍的局面放在心上,找一个\dmark{}\origpage{1155}与事体相称的解决办法(“eine~dem~Fall~geziemende~Lösung~zu~finden”——这话说得还要更
雅些),落款则是规矩十足的“hochachtungsvoll”(谨上),出自我最漂亮的手笔。

开诚布公地说,我想我算是开诚布公了!这一回,Springer先生,总没法说,
没有人把情形亲自通报给他吧——而且是第一手的,通报者不是别人,正是当事的主要利害人
本人!

无巧不成书:回复(又过了一个来月)偏偏就在昨天到了我手上。它短得可怜,我禁不住要在此
全文照录(译文)。何况我还愣了一会儿,才明白这确实是上个月我那封漂亮信件的回复。回复如
下:

\begin{quote}
海德堡,1985年3月15日

亲爱的Grothendieck教授:

仍要感谢您2月9日的来信。Byrne博士夫人2月15日的来信想必已答复了您的问题。

此致,等等

K.F.~Springer
\end{quote}

这一回,至少我算彻底明白了!那班“消息灵通的人士”(前文已经说过的)想必跟他解释过:
何必为一个给他写信的、有点激动过头的先生费神——此人分明不属于上流社会。而且,倒也
不假……

还在等着收这封来自Springer殡仪股份有限公司管理层的、如此富有教益的回信的时候(承蒙
公司还肯屈尊,由大老板亲自署名作复,真是有礼),我有的是时间省察我自己的打算。这家可
敬的公司扮演的角色,在我看来实在太大;我一度设想过一场好看的大官司:以一位被冒犯的
“体面先生”的名义,作为无可名状的专横行径的受害者,索取天文数字的赔偿。可我又想:这样
一场官司,得吞掉多少精力。就算我赢了,并且拿到一笔叫人头晕的赔偿(往好处想!),那也是
若干年之后——我又落着什么好呢?我不缺钱,也不需要比我已有的更多;一场骗局,并不因为
某场官司打赢了或打输了,就多一分或少一分是骗局。我不会因此把世界变得更好,不会把我自己
变得更好,也不会把 K.F.~Springer先生和他麾下\dmark{}\origpage{1156}某些雇员的做派变好——当然更变不了他们对本
行的看法——即便我动员起我的律师,再让他们动员他们的\fnote{902}{(*)}{我还有过另一种念头:局面兴许会倒过来——这可敬的公司反过来告我一状,说我损害了它的名誉。这些“为科学服务”的人士,一碰到自己的名誉,想必是要较真的(只要涉及的是他们自己的名声……)。}。我也不会让某个上流社会里某种精神变得更好——那种精神,使得像Springer博士及其可敬公司(十三年来)甘为仆役的那类行动成为可能。我还剩下(我由衷希望)几年可
活——光阴似箭,而剩下的年月里,我满眼是叫人兴奋的事情。收集证据,好让法官们相信我与
SGA 有点干系——这类事想必不大叫人兴奋。我当年写它们,不是为了他们,也不是为了 K.F.~Springer先生,才费的力气……

至于那些(除我之外)我为之写作 SGA 的人:他们同这份东西——它(至少对我而言)仍是我
自身的一部分——所保持的关系,绝非与我无关。它是他们同我这个人的关系的一部分。说来奇怪,
这层关系我真正了解(或者说多少了解一点)的,只有我的五位上同调学生:正是托他们的福,
如今一位 K.F.~Springer博士才可以把我像块破布一样打发走,仿佛我对顶着 SGA 旗号的文字被如
何处置、或不如何处置,连发表意见的资格都没有——那位无名之辈的名字在不在封面上,都是
一样。

哪位数学读者若有耐心陪我读到此处,又若某一天也曾在 SGA(我说的是\textbf{真正的} SGA)里
出没,兴许愿意跟我说说他自己怎么看这件事。要是能收到一个字,来自某一位——他发觉,那部
我曾经独自一人整个儿扑在上面、花了十年生命、而\textbf{世上没有一个人}在工匠离场之后还肯费心
把它继续下去的著作——的确带着它的构思者与孕育者的印记,带着它成形之前那必需的全部岁
月的印记,并终于长成一间\textbf{人人可进的房子}\fnote{903}{(*)}{“盖房子”这个执念——而且是“对一切都好”的房子——在我的数学著作里扮演过可观的角色,早在五十年代初就是如此。这是我在劳动中的那股被我称作“服务的冲动”的具体表达;它属于那些赋予我的数学工作以生气的深层力量(尽管在“阴阳之钥”的思考之前,我自己都未曾察觉)。“房子”的原型第一次出现在我的思考里,是在11月26日的注“阴,作为仆人——与新主人”(n\textsuperscript{o}~135)里,来势之猛,出我意料。}——
收到这样一个字,我一定会高兴。况且,一间人人可进的房子,不是贫民区里一座公共小便池,谁
都觉得可以随意进去方便,在又破又黏的墙上涂写脏话……

\pmark\origpage{1157}而倘若读到这里的人,曾经是我的学生,或者曾经是我的朋友,却始终不觉得自己被
激发起写信的冲动、或当面谈一谈的冲动——哪怕只是为这件事,而不是为别的——那么让他知
道:他的沉默同样雄辩,而且,它是会被听见的。

\subsection{第六根钉子(钉进棺材)}   % 18.5.9. Le sixième clou (au cercueil)

\subsubsection{a.\ 预发掘}

\notehead{176$_{1}$}{}

\fnote{904}{(*)}{(6月16日)下面这一组注(n\textsuperscript{o}s~176$_{1}$至176$_{7}$),以“第六根钉子(钉进棺材)”为名,应当被看作“沉默”那组注(n\textsuperscript{o}s~168~(i)至(iv))的自然延续;那一组处理的是“Motif”运算,尤其是其中最后一道“预发掘”(n\textsuperscript{o}~168~(iv)),日期为4月8日。下面这些注——除最后一道(n\textsuperscript{o}~176$_{7}$)外——写于4月19日与20日。我之所以宁愿把它们撂在这里、“四种运算”的末尾,而不并入“Motif”运算,是因为此前几个星期里围绕其余三种运算——尤其是那场(所谓“反常层会议”或“当值无名之辈”的)运算——继续进行的思考,而那场运算正是“封神”那组注的题目——给刚刚出现的那个(同样出人意料的)“新事实”,投下了意想不到的光。这里提醒一句:写下面这些注的时候,原则上,我已经在\emph{葬礼}底下画过“句号”(其终极之注“朋友”(n\textsuperscript{o}~188)日期是4月7日),并且打算一两天内就把葬礼(III)的全部手稿送去打字。也就是说,这些笔记是在“最后一刻补遗”的心境里写下的……}

(4月19日)我终于有机会(4月10日)读到了“预发掘”一注(n\textsuperscript{o}~168$^{1}$)所引的Langlands的文章。照Deligne去年8月寄给我的那份 Motif“带评注书目”的说法,Langlands这篇文章,与Deligne刊在同一卷里的那篇(即前引之注的对象)一道,是自1970年我离去以来头一批用到 Motif 的文章\fnote{905}{(*)}{不过,1972年Kleiman与Saavedra的那些报告除外——它们承接了围绕 Motif 范畴描述的那几份微薄的“练手音阶”(可参照注“把一个个 i 点清”(n\textsuperscript{o}~164)里打字稿第794页脚注(**))。}。直到去年我还不知道Langlands的文章(Deligne的那篇也不知道),这情有可原:作者(同我的前学生一样)不认为有必要给我寄一份抽印本。人们也不禁要问:他何苦费这个神呢——翻一翻他的文章就清楚,我这卑不足道的人,同他文中所谓“自守表示、Shimura簇与 Motif”那个题目,严格说来毫无关系\dmark{}\origpage{1158}。我的名字(借用一句我的打字机恰好一年前就背熟的套语)在这篇文章里、在书目里,一处也不出现。不过,有几个我大约在1964年提炼的观念(或者是我梦见自己提炼过——看来我还在重蹈覆辙……),我自以为认出来了;我甚至把这个梦的追忆(或者也许是一个关于梦的追忆的梦……)白纸黑字写了出来,恰恰就在1984年4月19日\fnote{906}{(*)}{关于这一点,见注“一个梦的回忆——或Motif的诞生……”,n\textsuperscript{o}~51。}。我恍然觉得自己回到了整整一年前的同一天。

诚然,隔了这一年,我已有时间见怪不惊。纵有不快,也几乎算不上意外(人家会说,本来也没多
少东西……),更谈不上震动。何况,那本为三年后那部难忘文集 LN~900 鸣锣开道的先导文章,与
后者之间有一个不小的差别:我没有荣幸亲遇Langlands;他的“Motif Galois群”(或“基本群”)的瑜伽,
不是从我嘴里听来的(Deligne1965或1966年前后可是听过的)。不过,整个六十年代下半叶,我在
周围、在凡愿听者那里,把它讲了个够(Langlands总不是刚下船的外乡人……),所以我有一份揣测:
Langlands心里明镜似的,知道这门关于各种Galois群与基本群的新“几何”哲学——把它们看作合宜的
仿射 pro-代数群——是从哪儿来的。我猜他心知肚明:这门哲学,并非1972年从某位Neantro~Saavedra~Rivano的脑子里生出来的——此人此后从流通中销声匿迹,没有留下任何痕迹\fnote{907}{(*)}{据Deligne去年10月来访时向我暗示的,Saavedra实际上已经改了行(如今“搞经济”了),自1972年论文答辩之后,就再没碰过数学。}。我觉得,Langlands倘若觉得有必要,就此表个态,并不算什么奢侈。当然,照眼下的世
道,盼他费这份神,也许是我过分乐观了……

\subsubsection{b.\ 意外的惊喜}

\notehead{176$_{2}$}{}

\pmark\origpage{1159}好事从不单行。读到Langlands那篇文章的次日,我又有机会翻了一翻Neantro~Saavedra~Rivano的那本Langlands大量援引的书,书名叫“塔纳卡范畴”(《数学讲义》265卷,1972)
。

我离去之前有过九个(男)学生,Saavedra是其中唯一一个我后来再无音讯的;也因此,没有任何
风声向我提示,他多少沾上了某场\emph{葬礼}的“颜色”或“气味”。我于是草草下了结论——带着我惯
有的那份天真:他(哪怕只因为没了机会——听说他后来离开了数学的水域……)是所有学生中唯一
全然置身于“葬礼运算”精神之外的人。可是,同Jouanolou的情形一样,音讯稀少到这个地步,本该给
我提个醒。我当然知道:他跟我共事时预定要成为论文的那份东西,最终以“塔纳卡范畴”之名,于
1972年在“数学讲义”出了书;而那本书,上星期之前,我不记得自己曾费神看过一眼\fnote{908}{(*)}{(6月16日)Saavedra想必不认为有必要寄我这本书;我没有它的本子。不过,七十年代里,我也许亲手捧过它。我留着的印象,仅此而已:他做了一份细致周到、拿来就能用的活儿;至于这印象从何而来,我说不准。这印象尤其在场,是在写“清场”那注(n\textsuperscript{o}~67,尤其打字稿第252--253页)的时候;那里我评说了Deligne的这个“谜”:他几乎逐字“誊抄”了Saavedra跟我一起做的那篇论文。}。我满脑子别的任务,压根没想到:Saavedra从此杳无音信,连他的论文答辩也不
通知我一声、不请我进答辩委员会——按说我是最清楚底细的人——这有点奇怪。正是读了这本书,
才明白他为什么宁可不打扰我,而把论文“偷偷”过了关,交给一个我完全不知其组成的委员会
\fnote{909}{(*)}{这个答辩委员会组成的谜,在“第六根钉子”的第七道即最后一道注(n\textsuperscript{o}~176$_{7}$)里,以全然出人意料的方式解开了;此处不再多置一词……}。那时\emph{葬礼}已经在加速进行:因为委员会里没有一个人认为有必要通
知我答辩一事,更不用说请我参加委员会了(Jouanolou的论文答辩,大约也在同期,情形如出一辙)
\fnote{910}{(*)}{关于一次更正,见前一条脚注所引之注。}。

\pmark\origpage{1160}这本书,呈现的是那门“算术几何”的一个关键侧面——它的视野(总体眼光)六十年代里在我心中诞生
并成长(当时还没有名字),而 Motif 的瑜伽曾是(并且始终是\fnote{911}{(*)}{不过,其间,这份“灵魂”又添了“非阿贝尔瑜伽”的滋养;《纲领概要》里略微谈到过它。(关于这份文本,见《引论》之三“罗盘与行囊”。它反正将收入《数学反思》第3卷。)})它的灵魂。可以说,Saavedra这本书的实质,是对我的某些
思想的一份形状完整、细致周到的陈述——关于某些范畴的一门“Galois--Poincaré理论”(我做梦也
不会想到叫它们“塔纳卡的”……);那些思想,是我在Saavedra还拿不准肯不肯下功夫去熟悉、去
消化,以便能把它们收进其论文的“陈述性”部分的时候,长久的、不厌其烦地讲给他听的。我曾交给他
一批原原本本的手写笔记,命题钉是钉铆是铆,证明概略俱全;至今我还等着他肯把它们还给我
\fnote{912}{(*)}{把我的手写笔记这里那里散发给学生,按需取用,是我的老习惯——而他们要学的头一
课,就是辨认我的笔迹。向来讲好:他们用完,就把笔记还我——但据我所知,这个愿望很少被遵
守。这不过是众多征兆之一,表明我的学生并不怕我,而宁可把我看作一块“好脾气”的面团:干起活
来要求是严的,除此之外,随和得没边……}。当然,论文本身的题目,并不是去陈述一个其动机他完
全无从把握的人的思想。要做的,是显式地给出我在这里称为“Galois--Poincaré”范畴的那些“张量”范畴\fnote{913}{(*)}{为了不叫它们“Grothendieck范畴”!然而,在我有幸引入并命名(因而并不冠我之名)的众多范畴与新观念之中,若有一类非用这称呼不可——单凭体面,我就想写——正是这一类!(topos 除外,但它的名字我觉得本就好……)至于“塔纳卡范畴”这个名字——由一位天才的前学生悄悄塞进来的,又被全体一致的教团欣然采纳——这不多不少,正是一场骗局,如下文原原本本所述。(见“会等待的人……”之后的注,n\textsuperscript{o}~176$_{3}$。)}的一种“便利”的内蕴刻画,亦即:允许对它们作“Galois--Poincaré--Grothendieck式”的描述,把它们表达为基底环 $k=\mathrm{End}(1)$ 上一个“仿射(pro)代数 gerbe”的线性表示。当 $k$ 是域时,我曾以(我引入的术语)所谓“刚性”性质给出过这样一个条件;我记得自己在1964/65年、思考 Motif Galois群之始,就写出过一个完整的证明\dmark{}\origpage{1161}\fnote{914}{(*)}{我没有费时去我的笔记里核实这件事——关于 Motif Galois群的笔记(或者说,其中没交给Saavedra的残存部分)。反正第3卷里我要回去谈,多半在“Motif,吾爱”那一章里。}。我得把原理指给他;至于我在这题目上写下的笔记,我克制着没有交给他——因为
学他未来的行当,是他自己的事,不是我的事,活儿得他自己做。如果我记得不错,当时对我而言悬
而未决的,只有一件事:这样的Galois--Poincaré式理论,就基底环 $k$ 须满足的假设而论,其自然有
效域何在——我尤其关心 $k$ 取 $\mathbb{Z}$ 这样的环的情形(因为牵涉到 Motif 理论的应用)。

在我离去前的所有学生里,Saavedra是来得最晚的\fnote{915}{(*)}{如果我记得不错,Saavedra是1968或1969年要求跟我工作的,在我不期然告别数学舞台之前一两年。},也是准备最差的,而且(至少起初)最没
有干劲想“露一手”。所以我不抱希望:他会超出我给他出的那个非常有限的技术问题——它只要求最
微薄的知识(一点概形语言、线性代数、平坦下降、gerbe 的语言,如此而已)。而他书里第四
至第六章所处理的那些更精细的问题(纤维函子的滤套、实域上Galois--Poincaré范畴上的极化结构,
以及“可极化”的这类范畴的清单,及其对 Motif 范畴和许多变体的应用),要求的却是“四面八方”
的知识\fnote{916}{(*)}{这尤其要求对约化代数群结构理论、及其在实域上的分类的深入了解,还要熟悉整整一批观念:Motif、晶体、$F$-晶体、分层模、局部系统(而此人当初连拓扑空间的奇异基本群也只略知皮毛),再加上 Hodge 理论,以及那些从未在文献里显式陈述、只“藏在行间”的精微的“极化”性质。},从而要求可观的入门功夫;我不认为Saavedra供得起。我至多盼着他也许在论文后
面附一个摘要(多多少少由我口授),把那套理论中不会收进形状完整的陈述性工作里的要点列一
列。直到上星期我才被纠正过来,并且意识到:Saavedra交出了一份真正惊人、而且是以破纪录
的速度完成的活儿\dmark{}\origpage{1162}\fnote{917}{(*)}{关于这个“纪录”的更深的思考,以及它(明摆着的)解释,见注“Monsieur~Verdoux——或献殷勤的骑士”(n\textsuperscript{o}~176$_{5}$)。}。这活儿结晶为一本书:陈述细致周到、一丝不苟,分毫不差地
摹拟着、拿来就能用,以(在我看来)近乎穷尽的方式,呈现了我在六十年代发展的那套几何-代数
形式主义。从这个角度看,我认为他做了一份有用的、处处都无愧于人的工作;我方才说的那份“惊
讶”,确实是一份“意外的惊喜”。

这份工作,精确地说,在于把我提供的一整套观念、命题与证明,以“典范的”、可径直发表的形状
整理出来(按照我当年那些严格的标准)。这样的陈述性工作,当然属于数学家这一行的本分——无
论处理的是自己的观念与结果,还是他人的。与我许多同事不同,我不认为:在评估一篇论文或任
何出版的质量时——及至授予做这种工作的人数学“博士”的头衔、把他当作不折不扣的数学家
时——这样的工作可以被当作可忽略的量。不过,有一点在我看来至关重要:这一行的某种起码的
伦理必须得到尊重;凡工作在于陈述与发展他人观念之处,这一点必须清清楚楚地写明,不留丝毫
暧昧。

然而,在眼下这个案子里,全书除了一篇光辉引言末尾附着的、三行纯客套的、含混其辞的“致
谢”\fnote{918}{(*)}{这篇引言的实质,是逐字照抄四条主要命题——那是我向Saavedra指出的、有待发展的Galois--Poincaré瑜伽的“四根台柱”(至于纤维函子上的滤套所牵涉的问题,它们很难被概括成一条石碑式的命题,故未列入);但其中一条——本该构成其论文“最低纲领”的那条——被添加了一个显赫而明摆着的错误,使得它平凡地成为假的!下一道注(“会等待的人……”,n\textsuperscript{o}~176$_{3}$)就会谈到,尤其见已引的注“Monsieur~Verdoux——或献殷勤的骑士”(n\textsuperscript{o}~176$_{5}$)及其后一道注“肮脏的差事”(n\textsuperscript{o}~176$_{6}$)。}之外,没有任何东西能让读者猜到:书中展开的任何一个主题——从全书的主题本
身开始——跟我这卑不足道的人有什么相干。我恍然又回到读到那本纪念性的 Motif 掘墓文集的
那一天(到今天恰好整一年,一天不差)!我的名字在全书里几乎无处可见,只有两三处例外:那里
需要形状完整的参考文献,而除我之笔,别无可用者\dmark{}\origpage{1163}。

而且,这绝不仅是——怕被人看清作者“只是”在陈述别人的观念与结果而感到的那点尴尬(这——
尤其在这个案子里——若工作做得聪明,已经很不少了)。我从许多不会骗人的“小细节”里看清:这
里绝不止是稍微“收割”一把、给自己的纹章镀点金、然后退回幕后的意思。这真是\textbf{为葬礼而举
行的葬礼}。只举\textbf{一}个例子——上帝知道,我曾花了多少天、多少个星期,向Saavedra原原本本、
翻来覆去地讲解那些观念:他初来乍到,一无所知;晶体、$F$-晶体(在特征 $p>0$ 的情形顶替缺失
的 $p$-进“系数”,好定义 $L$-函数……)、分层模(及其与局部系统的关系),最后还有一点 Motif
的瑜伽(暂时以标准猜想作启发的底座);这一切,为的是借一大把例子让他明白,我拿这些Galois
--Poincaré范畴要到哪里去;也为的是——以防(谁说得准呢……)他有勇气有毅力,在预定的“最低纲
领”之外,至少收进一章典型的例子。他心里一清二楚,不必我多费唇舌:这些都是关键的几何观念,
而且并非开天辟地就有;把它们讲给他听、一遍遍讲给他听而不知疲倦的,不是别人,正是我——
它们是我在此前五到十年间引入的,用来伺候某一视野(总体眼光)的工具(即使那视野从他头顶掠过,
正如它掠过我所有学生的头顶,只除了一位\fnote{919}{(*)}{此人在师父一转身之后,就按捺不住把它埋葬了……})。可是,在他引入并稍稍发展这些观念的地方(在专述例子的第六章里),也如同在他佯装以作者自居的那部分理论正文里,我的名字照样不见踪
影。然而我很难设想,Saavedra会真的以为:读者无论多么孤陋寡闻,哪怕满心愿意相信他是这些
范畴(他慷慨地称之为“塔纳卡范畴”)的父亲,也不至于进一步以为:连 $F$-晶体、Motif 与“塔纳卡
装备”(原文如此)里的其他小玩意,也是这位Saavedra为应一时之需发明出来的。倘若这些观念被
当作信手拈来的寻常货色、或从最近的\dmark{}\origpage{1164}孤儿院里捡来的弃儿,那我在其中认出了一种\textbf{风格}——
一年来我为\emph{葬礼}兜圈子兜个没完,这种风格我再熟悉不过了……

是Mebkhout把那本书带给我的,他满心欢喜,总算能给我看一个我的学生的例子:这位至少还是
“诚实的”\fnote{920}{(*)}{(6月16日)再说,这事没成,他可是难过极了,还尽力来安抚我——这让我想起Kawai的情形(见打字稿第1078页脚注(*)),也想起了Beilinson:Mebkhout觉得Beilinson“比Bernstein更诚实”(见第1072页)——真像第欧根尼提着灯笼,只不过这一回,是在那伙一头扎进上同调这个名声实在不佳的题目的人当中,找一位诚实的数学家……}。他显然被引言末尾那三行致谢打动了——也难怪:1972年,感
谢某位已故者的风气,本来就已不大流行;从那以后,在我不止一位前学生那里,取而代之的调门,
不是冷嘲热讽,就是打趣,再不然就是彻底的沉默。话虽如此,这一回我到底得到一份“深深的感
激”:感谢我“把[作者]引入了这个题目”,感谢我的“劝告与鼓励……对完成这项工作必不可少”
\fnote{921}{(*)}{这些“致谢”,就情形而论,不过是一句玩笑话:人们会以为,我把作者“引入”了单复变函数这一“题目”,或任何同类的经典题目。其实,当我在一位正愁没论文的Saavedra面前谈起这桩事时,这个“题目”还并不存在——除非是在我之内、与 Motif 共生而发展出的那一份视野里,以及在我赋予它形状的手写笔记里。关于这份视野的诞生与发展,我在“一个梦的回忆——或Motif的诞生……”一注(n\textsuperscript{o}~51)里谈到;至于我的一位前学生(在众人纵容的目光下)对这根源所作的、漫不经心的彻底抹除,见其后一注“葬礼——或新爸爸”(n\textsuperscript{o}s~51、52)。(6月16日)Saavedra的这些致谢之所以更是一句“玩笑话”:其作者从来不曾费心,哪怕只寄我一本他的书、连同这些冒牌的致谢。既然我已经把“塔纳卡范畴”(原文如此)这场运算巡视完毕,如今我更明白:我这位前学生没有理由为他的“工作”(原文如此)感到骄傲,也不急于让我看到它。而照当时、乃至直到两年前事情的走势看,那位工匠似乎根本不大可能有机会看到它……}。这就是所谓的拿话头搪塞人;而在我看来,在\emph{葬礼}分明大行其道的时刻,在陈述
自己的工作时保有起码的诚实,才是表达“感激”的更有力的方式。
\subsubsection{c.\ 会等待的人}

\notehead{176$_{3}$}{}

\dmark{}\origpage{1165}说实在的,这本书一到手,我就看清:在那场难忘的“SGA~4½--SGA~5 运算”之前,整个
\emph{葬礼}里,还没有哪一幕的分量,赶得上这本名叫“塔纳卡范畴”——一个再平淡不过的书名——的
LN~265。此前的插曲\fnote{922}{(*)}{这里说的“插曲”,在“埋葬……”那条注(n\textsuperscript{o}~168~(ii))里匆匆勾勒过;它属于“Motif”运算那组注。}不过是或隐或显的“收割”,把某些重要思想的承续隐瞒起来。而这
里,我的“算术几何”视野的一个整个的关键侧面,被人不动声色地“挪用”了——经手的,竟是我的
所有学生当中看上去最“不起眼”的一位!

\pmark 的确,在这位学生背后,我凭一种不会骗人的笔法,认出了那个牵线的人;而且此公在Saavedra极尽殷勤的致谢名单里,赫然排在前列\fnote{923}{(*)}{就“数学”本身而论,这几位依次是(按出场顺序)我自己(不按字母序,算他有礼貌)、Berthelot与Deligne。}。Saavedra给这部书、给我引入的那个关键观念所起的名字,是一件精巧
的\textbf{褫夺之作}。论其斩钉截铁的效力,五年之后才被另一个名字超过——这回,那名字出自Deligne
本人之笔,落在另一部书上\fnote{924}{(*)}{正如下文(已引的注“Monsieur~Verdoux——或献殷勤的骑士”)将显现的,至少有强有力的揣测:与其读作“但这回出自Deligne本人之笔”,不如读作“并且同样出自Deligne本人之笔”。}。

\pmark 如果说,顶着一卷“锯倒之作”的“SGA~4½”之名,是一桩天才的冒名顶替,那么“塔纳卡范畴”
这个名字,就是一桩同样天才的\textbf{骗局}。哪怕对一个“平凡”或“中性”的Galois--Poincaré范畴——
等价于域 $k$ 上一个仿射群概形 $G$ 的有限维线性表示的范畴——我发展出的那套瑜伽,也十足是
“Grothendieck式”的:它由我早先在拓扑空间、概形或(更一般的)topos 的基本群情形发展的类似
瑜伽启发而来。把基本群定义为“覆盖层范畴上一个纤维函子”的自同构群,以及(同样古怪,因为
新,故而不惯)系统地同“不必\dmark{}\origpage{1166}\textbf{连通}”的平展覆盖的范畴打交道的想法,当年没少招来讥讽。我从
来不放在心上,深知那班爱开玩笑的家伙——自以为在学堂板凳上学过Galois理论或Poincaré理论就
算懂了——没有一个真正懂过;而且直到今天,他们之中谁也不能,不逐字照抄我做过的东西、不照
搬我以范畴等价给出的覆盖的Galois--Poincaré理论的表述,独自迈出(比方说)稍一般些的概形的覆
盖的Galois理论的\textbf{最初几步}\fnote{925}{(*)}{“稍一般些”在此可以精确理解为“一个非正规的概形”。在我之前,代数簇的基本群(由朗与Serre)只在正规簇的情形引入过,办法是把它描述为其函数域的“绝对”投射有限Galois群 $\mathrm{Gal}(\overline{K}/K)$ 的一个合宜的商。}\fnote{926}{(*)}{今天,这种表述基本群与覆盖之关系的方式——哪怕在普通拓扑空间(局部弧单连通)这个可以说是“教科书式”的特例里——也开始到处流通,而对我们共同的先祖只字不提,这还用说么……}。

同样,从一个“抽象”的有限维线性表示范畴出发——带着它自然的乘法结构与它自然的“忘却 $G$
之作用”的纤维函子——重建一个(为确定起见,取域上的)仿射群概形,把它作为\textbf{这个函子的自
同构的群概形}——这个想法,既不归功于塔纳卡(他可从未要求这么多),也不归我那微末的前
学生Saavedra,也不归我最出色的学生Deligne(令我遗憾——不过那时他还没在跟前),而是典型
“Grothendieck式”的。由此得到 $k$ 上仿射群概形与配有 $k$ 值纤维函子的刚 $k$-张量范畴之间的完
美对应,道理亦然。再者,倘若碰巧(特征非零域上的 Motif 范畴,就往往如此)一个刚张量范畴
(不幸,或者幸运加倍……)偏偏\textbf{没有}纤维函子可用,那么“代数群”就该换成一个代数
“gerbe”——这个想法,早在年轻的Deligne在数学里还没听人说过“gerbe”这个词、做梦也不曾想
到会有类似东西的年代,就已经被阐发得淋漓尽致。同样,当Giraud在六十年代自告奋勇,发展一整套
二维以内的非交换上同调代数的武器库\dmark{}\origpage{1167}——栈、gerbe 与 lien 的语言\fnote{927}{(*)}{这套形象化的术语是Giraud引入的,代替我从1955年起使用的临时术语(有点随手抓来,如“局部性质的纤维范畴”之类不成体统的名字;那些观念的根本性质,要求的是斩钉截铁、掷地有声的名字)。
(6月16日)在他那本书的引言第一页,Saavedra说到了Giraud“引入”的非交换同调代数“形式主义”。这是众多地方中的一处:在其中,我感觉到有一个比本书作者更机灵的人,在“搀着他的手”……正是同一位,一提“导出范畴”,就乐于紧接着补一句“由Verdier引入”(而这两处的底细,他心里一清二楚……)。}——讥笑声同样不缺。这类东西,如今、而且早就被Deligne之流叫作
“废话矿渣”。这些讥笑不曾让我却步\fnote{928}{(*)}{Giraud却不然——他在刚与我一同开辟这主题不久,就与之决裂,再无回头。关于这一点,见“共同继承人……”一注(尤其打字稿第386--387页)及其后一道注“……以及那把电锯”(n\textsuperscript{o}s~91、92)。};我知道自己在往哪里走——我看着这块“矿渣”以
完美精细的手法,抓住那些微妙而深刻的关系(我深知,任何别的“语言”都无力把握它们),心中是
“复苏”(如我在别处所写),却谈不上惊讶。

话虽如此,当这班讥笑者某一天发现一道被自己漏掉的“奶油甜点”——无论是那被人抢着命名为
“塔纳卡的”(且听下回分解……)范畴,还是某种“对应”或“关系”或“构造”(边上带点新Grothendieck味儿),被人用委婉语一笔带过、或被命名为“Riemann--Hilbert”(同样且听下回分解……)
\fnote{929}{(*)}{关于这后一个“且听下回分解”,见“反常层会议”的全套,尤其注“魔术师”与“冤大头买卖——或木偶戏”(n\textsuperscript{o}s~75$''$与 171$_{2}$~(e),后者属于长注“黑手党”(n\textsuperscript{o}~171$_{2}$)的一部分)。}——这时人人蜂拥而至,争当光辉的发明者。这就是本世纪七十年代、八十年代数学的
“时代精神”……可以确定的是:给这些范畴(我曾长久讲给他听的)起“塔纳卡范畴”这个真正光辉
的名字,这样的主意,不是一个Saavedra想得出来的。放任他自己,他绝不敢改动从我这里承袭的
术语,至少得先征求我的同意——这原是最起码的规矩!非得有上头的榜样与鼓励,他才敢\dmark{}\origpage{1168}这样把我
当作可忽略的量打发掉。再说,这可怜人为了熟悉必不可少的东西——若他想完成我交给他的那个雄心勃
勃的写作计划的哪怕一部分——手头的活儿已经够多了\fnote{930}{(*)}{他以仅仅两年的创纪录时间完成了这个计划——从我离去(那时计划几乎尚未动工,除了对概形基本技术的初步熟悉)算起。即便有Deligne在一旁搀扶(在我离去之前,他对这个学生毫无兴趣),这份成绩也简直堪称奇迹;这个“奇迹”,在注“Monsieur~Verdoux——或献殷勤的骑士”(n\textsuperscript{o}~176$_{5}$)里得到较细的检视。},哪有余暇再去翻文献、去读塔纳卡以及天知道什么的东西——他跟我共事的时候,肯定
从没听说过这些\fnote{931}{(*)}{提醒一句:Saavedra跟我工作,恰好在我离去前一两年(约1968、1969年),此后我几乎完全失去了他的踪迹。他那时的行囊,同任何一个来自第三世界(或我们外省院校)的三年级研究生相比,不厚也不薄。}。

这个名字之所以“光辉”,在于两种看似矛盾的品质的精巧结合。其一,在肤浅的观察者眼里,这
名字不显得全然荒唐。“人人”都约略记得存在一个“塔纳卡对偶”,乘法结构在其中扮演角色——
这看起来同那著名的 $\otimes$-范畴(某位Saavedra——这是谁呀?——称之为“塔纳卡的”)身上发生
的事,还真有几分像;那就“塔纳卡的”吧,为什么不呢!可是,对会等待的人,事情自己会成熟。十三
年过去;如今,顶着谁也没见过的无名作者之书的,换了一部三年前问世的、显赫得多的参考:光辉
的 LN~900,出自非别人而是Deligne之笔,外加一位名叫Milne的搭档。这两位大名鼎鼎的作者,从
头发展了他们同样称作“塔纳卡的”那类范畴的全部形式主义。显然,这是一个根本性的
观念,多年来已被Langlands、Deligne、Serre等人使用,前程似锦。当然,谁也不会相信:在这篇文章里被
顺带引了两三次的某位Saavedra,竟是这一关键观念、以及由它生出的那套精微形式主义的主人。
两位光辉作者的这篇文章,以其主笔人所熟知的全部娴熟,重拾这个题目,其笔调本身,就把这层意\dmark{}\origpage{1169}
思交代得没有半点含糊\fnote{932}{(*)}{关于这篇文章,尤其见“下葬——或新父亲”(n\textsuperscript{o}~52,尤其打字稿第214页)与“清场”(n\textsuperscript{o}~67,尤其打字稿第252--253页)两注。}。更不必说,他们还从Saavedra书里陈述的理论中,挑出一处粗糙已极的错误(以致
他们不得不从一个全然不同的定义出发——那个定义,终于像是正确的\fnote{933}{(*)}{关于Deligne的这桩绝技(Milne作陪衬),见屡次所引注“Monsieur~Verdoux——或献殷勤的骑士”(打字稿第1176页)的开头。}),叫人不免要问:这可怜的
Saavedra(当年必有某人——猜猜是谁……——费过心给他解释是怎么回事)到底懂不懂自己在
说什么。而Milne纵然光辉,且有幸与显赫的Deligne合署一篇文章、发展一个分明根本性的观念,
也决不会起念头认为自己能算它的父亲或哪怕共同父亲;Beilinson与Bernstein同样不会声称,他们发
明了(或哪怕共同发明了……)那句有名的“本该在这几篇笔记里找到自己位置的关系……”——他们
有幸与同一位显赫的Deligne合署了它,而在那之前,正是Deligne以他无上的好意,给他们指点了
Kazhdan--Lusztig 证明的门径……最后,又有谁会\textbf{认真}相信:那位把自己的名字(想必未经商
量)借给这一根本观念的、大名鼎鼎的塔纳卡,果真与此\textbf{有}几分干系?等到哪天世人都看清了这
个观念、以及随之而来的整套精微理论的\textbf{真正的父亲},他也不至于出面认领——倘若他还在世
的话。对此若还有半点疑虑,只消翻翻塔纳卡的著作——嫌多的话,就翻他论“塔纳卡对偶”的那
篇——就会发现:说到底,风马牛不相及……

这里,仍然是:几块界石一立,剩下交给时间。显然,这门理论——它将越来越多地显形为一种新的
\textbf{哲学}的技术手段,用以沟通几何与算术——在未来的年月里,注定要越来越走到数学舞台的中
央。五年、十年之内,再不会有人想到,就这个题目去援引某个无名作者的书;而那位想必搀过他手
的人,倒是不辞辛劳,写出了一份势在必行的陈述,由一位光辉的合作者协助,构成那部光辉的文集
的核心——Motif 的观念,终于在\dmark{}\origpage{1170}那片坚实的土地上得到发展。(至于那位含混而潦
草的先行者——他早已被遗忘——以及那道分明超出他能力的、惯常的猜测性“废话矿渣”,文集里
只字不提,倒也显得仁慈……)引用“P.~Deligne 与 J.~S.~Milne 的塔纳卡范畴”,将如同引用
FAC 或 GAGA(Serre的)或 SGA(法国高等科学研究所那部著名的匿名讨论班,人称“布瓦马
里”)一样,成为第二天性。而这么做的时候,任何人的心里,对这些新观念的父亲是谁,都不会有
半点含糊——那父亲断然不是合著者Milne,更不是塔纳卡,甚至也不是某位严格无名的作者(一个
叫Saavedra的):他在人家文章里被顺带提了两三次,为的是他(在他那卷书引言里)写过关于这
个题目的一份“出色的摘要”(带若干保留)。

不过,人们不会指望这门理论的父亲,让他去强暴自己那份出了名的谦虚,以至于把(按“消息灵
通”人士——他们在这类事情上说了算——的全体共识,理应如此称呼)本该如此称呼的东西,
叫作“Deligne范畴”(或者,在另一个天地里,“Deligne对应”)……

\subsubsection{d.\ 父亲们的圆舞曲}

\notehead{176$_{4}$}{}

(4月20日)昨天的思考,让我以新的眼睛看清一件去年——我刚一头栽进\emph{葬礼}的时候——曾让
我瞠目的事:“……这件看似荒唐的事:Deligne十年之后‘重做’Saavedra的论文*!”从去年4月19
日——我读到那本“难忘文集”LN~900 的日子——起,这件事就被谈到了:文集里(除其他美好东西之
外),Saavedra的论文被几乎逐字重现\fnote{934}{(*)}{见倒数第二条脚注所引诸注。}。一周之后,我在“清场”那条注里又回到它。当时我已得出“内
心确信”:这堆废话背后的意思,是光辉的Deligne(自任Saavedra的文书)的如下愿望:
\begin{quote}
“给自己一种幻觉中的解放感,解脱他肯定感到的一桩痛苦的义务:不得不时时援引那个他正要
取而代之、加以否定的本人,或哪怕任何一个援引了他的人。”
\end{quote}
可上星期,我第一次费心翻阅了这位“别人”的工作,惊讶地发现:他压根\dmark{}\origpage{1171}没想“援引我”(除那三行
冒牌的“深深感激”之外——明摆着是障眼的烟幕)。于是,一年前我的“内心确信”就站
不住脚了——里面肯定有对的成分,可仍剩下谜:总不至于就是那三行字(没有哪个读者会想到跑去
引言末尾把它们挖出来),竟能促使一个Deligne,去当一位早已作古的宗师门下最无名学生的抄写员
吧!更不必说,在这段引言的末尾,我几乎与他、与Berthelot并列出现:他们同我一样(人们会说
\fnote{935}{(*)}{只是有这样一点不同:是我“把他引入了这个题目”(原文如此),而且“他至今在数学家的养成上,还大大受惠于我”(这真是受之有愧)。}),有资格得到“他们在这一工作中慷慨给予的帮助与建议”的
致谢……

这个“谜”在昨天的思考里完全解开了;我不必去找,甚至不必刻意去想。写完之后重新回味,种种联
想浮现出来——写作之时它们必定已经在场,只是我毫无察觉,还在暗中引导我的笔。打动我的,不
仅是风格的相似,更是一整套\textbf{专利的据为己有之术},贯穿\emph{葬礼}的三大“运算”(四者之
中,Deligne本人是主要的——即便不是唯一的——“受益人”)。不妨称之为“临时替身爸爸”的手法:
悄悄把它摆上数学勒索的棋盘,好变走一份真实的父亲身份,而我朋友皮埃尔其人,则暂时隐在幕
后。一旦生身父亲被体体面面地彻底清出舞台,替身爸爸也就被变走,仿佛从来不曾存在;于是\textbf{真
正的父亲},谦和地微笑着登台,甚至无须亲口说破——因为那个不动声色牵线、又懂得等待的人,
诸事自会水到渠成,毫无阻拦:教团全体的同心一致,早已把他按他所应得的名分加了冕。

这套手法,直到前几天,我循着我朋友Zoghman在第四场运算(所谓“当值无名之辈”)各插曲中的
倒霉遭遇一路追索,才开始看清。这一场里的“替身爸爸”(为某件“对应”……)是\textbf{Kashiwara}——我
说不好,他是就这么从天上掉下来的,\dmark{}\origpage{1172}似有天助且纯属巧合,还是未来的真爸爸巧妙地让他明白:
一个无名之辈的那个没有像样父亲的结果,横在这里,弃之可惜\fnote{936}{(*)}{(6月16日)看来,对Mebkhout著作施行扒手手术的主意,正是企业家Kashiwara出的,而且早在1978年——其时Mebkhout刚把刚杀青的论文第三章拿给他看不过几个月。关于这一点,见注“黑手党”部分(b)(“最初的麻烦——或太平洋彼岸的黑帮老大”),打字稿第1060页脚注(*)。}\fnote{937}{(*)}{(6月16日)其实,两年前他就开始露头了——见前一条脚注。1980年3月的那一幕,是Goulaouic--Schwartz讨论班;已引之注以及“掠夺的空白授权——或‘高尚行当’”一注(n\textsuperscript{o}~171$_{4}$,尤其打字稿第1088--1090页)都谈到它。}。反正,皮埃尔朋友深谙
此道,在一个人为的、由“行家们”斩钉截铁的共识拼凑出来的父亲身份的暧昧上,把戏做得
滴水不漏——而这早在那件新东西的分量得到公认之前就动工了。替身爸爸Kashiwara,1980年3月即已
登场,若不是早在六个月前的莱苏什讨论会上就已露头;十五个月后,在1981年6月那场难忘的
讨论会上,他被变走,不留痕迹(而且看来他自己也不太计较)。在这里,变走的手法炉火纯青:再引
进两位——这次不妨叫作“名义上的共同爸爸”——Beilinson与Bernstein;他们登台,纯粹是一道讲究
格式的保险条款——“喊停!”——因为当然,谁也不会想象孩子是他们俩谁做的(虽说两位都沾了
光……)。

与“Motif 运算”的类比,着实惊人!当人们还能拿出来的、关于 Motif 的“外行废话”的父亲身份,
知名度还太高(尤其七十年代初),无从下手时,Motif 瑜伽有两个\textbf{关键侧面},还不曾有一行
字公开发表过,哪怕以影射的方式。其一,“权的瑜伽”,早在1970年就被\textbf{超级爸爸}据为己有,
纹丝不动——反正被变走的那部分,不过是“猜想”,跟一句纯形式的影射差不多。另一侧面却毫无猜
测成分,六十年代下半叶就已打磨停当。按安排,由一个有点跟不上趟的无名学生,来陈述这套瑜伽
的启动机制——技术上不难,但(至少直到那位不受欢迎的生身爸爸“死亡”前后)这件活儿看来超
出了那可怜人的能力。

\pmark\origpage{1173}正是这位学生,Saavedra,是现成的替身爸爸:有幕后人暂时的担保撑腰,足以让一个一心
只想忘掉该忘之人的教团点头称是;但与此同时(妙就妙在这里),这位“爸爸”分明“没有分量”。
时候一到,谁也不会——Saavedra想必比任何人都更不会——起念头,宣称他可能是某种新哲学的父亲;只要
肯稍微想一想,这种假设简直滑稽……这里,功成身退的替身爸爸被请出去,要晚到十年:1982年,
难忘的 LN~900 出版。应当说,从1972年(第一场所谓“Motif 运算”中抬出“替身爸爸”)到1980年
(第四场“当值无名之辈”运算中同样应运而生的替身爸爸登场),多少水从桥下流过,再用不着绕
弯子了!耐人寻味的是,这里又插进来一位“名义上、装装门面的共同爸爸”,好让替代父亲身份(说
到底,就是木偶戏里的父亲身份……)与\textbf{真正的}父亲身份之间“平稳”过渡(而且没人显得在
抢镜头)。我确信,Milne对那双按他人意志摆布他的无形之手,看得不会比Beilinson与Bernstein更清
楚。人人都分到了碎屑;人人(至少是有发言权的人……)都有一切理由心满意足
\fnote{938}{(*)}{(6月16日)在他与Saavedra、与一位(暂时)在幕后的Deligne之间必然达成的“交易”里——这位Deligne,只等时候成熟便重新露面……——Saavedra的“份额”是:一张国家博士学位文凭,以及作为“数学讲义”这一显赫丛书作者而获得的相对知名度——这足以为他启程回国、远离他只远远瞥见过一眼的枯燥数学追求,铺好道路……}。

\pmark 这一切,昨夜还让我想起那第三场大运算——那场直接惠及“未来无所不在的爸爸”的运算:
“平展上同调”。此前我已使自己相信:这场运算的初始动机\fnote{939}{(*)}{关于这一点,见“那个公式”那组注(n\textsuperscript{o}s~169$_5$--169$_9$)。这最初的用意,后来大大扩展开来——尤其见“悼词(一)——或恭维”(n\textsuperscript{o}~104)与“宝石们”(n\textsuperscript{o}~170(iii))两条注。},是据为己有某一\textbf{不动点公式}的
打算:人们可以把某个父亲身份不受欢迎的“L 函数公式”呈现为该公式的一个平凡推论。麻烦在于,
所论的迹\dmark{}\origpage{1174}公式偏偏带着\textbf{同样}不受欢迎的父亲身份。幸好,还有另一位可能的父亲——一位好朋友(不点名,Verdier),他做
出了两份公式:一份太一般(但启发式地至关重要),另一份略显局促,却足以“罩住”人们想要的东
西。可朋友归朋友,这里的“父亲”既不能是朋友,也不能是那个不受欢迎的死者——因为事关“那
个”著名猜想的那把钥匙公式\fnote{940}{(*)}{当然是“那个”Weil猜想。关于这一点,见注“那个猜想”(n\textsuperscript{o}~169$_4$)。}。糟糕的是,L 函数公式与它那不合时宜的父亲身份,偏偏尽人皆知;所以
这里的难点不在朋友(朋友之间,最后总能商量……),而在死者。更倒霉的是,他对“推论”的证明,
1964年就白纸黑字发表在一场Bourbaki讨论班上——只是(万幸!)该公式的那桩例行情形(呃,
对不起,我是说关键情形!——这公式或迹公式,反正一码事,但这话千万说不得……\fnote{941}{(*)}{这两条公式,实际上互为直接推论。既然我对其中之一(L 函数公式)的父亲身份尽人皆知,Deligne便设法(在那份名叫“SGA~4½”的难忘文本里)把它呈现为另一条的推论,并且使尽浑身解数,以远比我这微薄的证明(以及随附的命题)艰难得多的魔术-行骗手法,把自己装扮成后者的父亲。这套戏法见已引的那组公式注(n\textsuperscript{o}s~169$_5$--169$_9$)——在我们这门可敬科学的年鉴里,这大概是绝无仅有的一场大变活人。}),当时还来不及被验证。

这里的手法,是拿那位朋友充当他的超一般公式的父亲(这原是分毫不差的事实——只除了一点:
他从来没费神去证明它……),同时悄悄把它与那累赘的死者所\textbf{显式}证明的公式(对那一条,自
然讳莫如深,一字不提)混作一团,再把那条超一般公式\textbf{拆毁}(说它是猜想的、不完全的、说
白了,没用的)。这叫溺毙一条鱼:叫读者压根不去想,翻一翻某套 SGA~5 讨论班(人们本来就责成
他“忘掉”它)会就这个问题说些什么。至于朋友那条(边上略局促、却完全有效的)显式公式,大家
一致同意也不再提了,只在一个文本的末尾留下一处暧昧而装门面的引用——那文本拖沓冗长、令
人生畏,世上没有哪个读者有勇气读到底。总而言之,可以说:“替身爸爸”(这里是Verdier)确实出
了场,但与其说是凭他对一项结果的默许的“父亲身份”——这里要\textbf{彻底变走}的,恰恰是\dmark{}\origpage{1175}死者的那
项结果——不如说是凭他的默许,参与了一场搅浑水、拆两条“孩子”的把戏:那两个孩子
确凿是他的,搅浑水是为了在混战中变走第三个孩子——那孩子的父亲不可宣之于口;如今是一个
孤儿,谁也找不回他,谁也不上心去找他\fnote{942}{(*)}{关于这一点,见注“变戏法者——或飞走的公式”(n\textsuperscript{o}~169$_8$)——以及“家庭相册”一注打字稿第1121页的脚注(**):它表明Deligne与Illusie这对好撒玛利亚人的变走-蒸发之计,是如何大功告成的。}。在这套手法里,Illusie扮演一个辅助角色,同方才
的“名义共同爸爸”相仿——只是他的父亲身份,同Verdier的一样,任何时候都不许落在那神圣的
\textbf{对 Frobenius} 的迹公式头上——那是唯一算数的,理当(当然,以应有的细腻)留给Deligne一
人;它落在那个要被变走的、不可宣之于口的孩子身上——Illusie以他特有的、堪称楷模的虔诚,参
与了这件事。

\subsubsection{e.\ Monsieur~Verdoux——或献殷勤的骑士}

\notehead{176$_{5}$}{}

可我还想回到Saavedra的“论文”上来。大约在我告别数学舞台之际——1970年初(如果我记得
不错)——Saavedra才装作终于真正“上钩”了:此前一两年,他一直显得三心二意。他对
我说,他已提炼出我最初给他出的那个命题的一种表述与证明,使它适用于基底环 $k$ 为\textbf{任意}环
的情形。他甚至给我画了一张证明的草图,我只半心半意地听了听——我那时几乎全部精力,都被我
正经历的人生转折占着。没有细核Saavedra的话,我有一种印象:他总算起步了,往后可以自食其
力。也许,在一个我对真正的科研指导已经几乎无暇顾及的时刻,我有些性急,把愿望当成了现实
\fnote{943}{(*)}{至少,同我离去之前的投入相比;但并不比我后来在大多数同事——那些担任科研指导的同事——身上看到的更少。}。此后,据我记忆,再没有他的音讯\fnote{944}{(*)}{我的记忆在此稍稍背叛了我——见注 n\textsuperscript{o}~176$_{7}$,其中对此有出乎意料的揭示。}。直到上星期,我还一直以为:我给他出的那个最低纲领,他总该完
成了,也许还稍微\dmark{}\origpage{1176}多做了些——把 Motif 的情形也处理了(照Deligne去年8月随其 Motif 带评注书目
写给我的信所言)。

\textbf{我这才刚刚发现:压根没有这回事}。这可怜人花了三四年功夫,竟在他称之为“塔纳卡范畴”
的东西的\textbf{定义}本身里(我指的是那个内蕴性质的定义),犯了一个粗糙已极的错误\fnote{945}{(*)}{错误的来由,是Saavedra心里对我所说的张量范畴的“基底环”有混淆:它绝不是使该范畴“线性”、使张量积“双线性”的随便什么环,而是典范的环 $\mathrm{End}(1)$(1 为范畴的单位对象)。我给Saavedra讲理论入门常识的时候,他想必“不在局中”到这般地步,这些话全从他头顶掠过,沉入遗忘。Deligne——他似乎多少接了我的班,当上了Saavedra的指导(显然另有一本账……)——却绝不肯替他纠正航向。这让他得以(十年之后)悄悄拆掉Saavedra的纸牌楼,以救世天使、以及(这一回又是)人人翘首以待的真爸爸的面目出现……}:而这个定义,正是要用来证明
它蕴含“Galois式”的、以一个合宜 gerbe 的表示写出的描述。他在引言里陈述的定理三(那篇引言,
按理至少应当\textbf{陈述}我交给他的那理论的四条根本定理)因此\textbf{平凡地为假}。Deligne与Milne以指出这个显赫的错误为乐事,提议把所论范畴的“新”定义取作 gerbe 式的描述(先天地明摆
着,这才是对的,必要时改一改内蕴描述就是……),并一本正经地追问:Saavedra的“定义”(把那
愚蠢的错误剔除之后)是否蕴含“他们的”(原文如此)\fnote{946}{(*)}{前引书第160页(我可不编瞎话!)。}——而这恰恰原本就是预定要构成Saavedra博士论文的那桩题目!

这一处境,是十足的于布爸爸的翻版!而且一上来就是三十六个花样。譬如,交给Saavedra的那
桩题目——唯一需要一点原创贡献的部分,哪怕多么微薄(为一个尽可能一般的基底环上的Galois-Poincaré范畴,提炼好的内蕴条件)——即使在(我相信我在遇到Saavedra之前很久就已处理过
\fnote{947}{(*)}{那是1964或1965年,即Saavedra那篇出名的“论文”(原文如此)之前七八年;比Deligne--Milne二人组赶来救驾——也\textbf{不做}这份微薄的活儿——还早十七八年:那是我对最微薄的学生所期待的唯一“原创”工作……})基底环\dmark{}\origpage{1177} $k=\mathrm{End}(1)$ 为\textbf{域}的情形
——也根本没有被处理。Saavedra的“论文”因此精确地在于:虔诚地照抄理论中(Grothendieck式瑜伽的启动之外)早已由我全部做完的那一部分——基底域上的那部分——并且,以一个
骗人的定义和一个假定理的“证明”顶替那先于一切的必做之事;那“证明”(Deligne以指出这一点为
己任——前引书第160页)归结起来,不过是一个简单的循环论证!

还不止这些。论文站不住——而答辩委员会毫无察觉!想必没有一个委员真正弄懂了所为何事。可
这并没有促使他们中任何一位向我示意:至少他们中该有一位,有能力为那项他们一本正经摆出架
势来评审的工作的严肃性,提供一份有效的担保\fnote{948}{(*)}{这支可悲的委员会的组成,终究要(在坚持读到那里的读者面前)在“第六根钉子”的终极之注 176$_{7}$ 里揭晓——在我棺材上的第六根钉子……}。答辩居然举行了,而且没有让我参与——这只能拜Deligne的担保所赐;
(如Saavedra的致谢所暗示的)在我实际上从舞台上消失之后,Deligne必定稍稍跟进过他的工作
\fnote{949}{(*)}{Deligne对一个默默无闻、想做论文的学生的这份突如其来的关切,是在该理论的自然父亲(且不受欢迎……)死后才出现的——何苦来哉;而那位学生(显然被任务压垮了……)本该陈述的,正是这门理论。}。

再说,打那时起,Deligne居然没有察觉这个错误,我简直无法想象——我深知他的敏锐与犀利,直至
最小的细节——而这绝非什么“小细节”!当然,我曾把已达成的瑜伽的全部精妙之处讲给他听;我
给他解释的最初几桩里,若不包含这个反例——他同Milne如今佯装当作最新发现拿出来的那个反
例,它早在我思考瑜伽的最初就为我所知(我最终打算把这瑜伽就叫“Grothendieck式”,不必去麻烦
Galois与Poincaré的大名……)——那是根本不可能的。他若听任“门生”(原文如此)的“论文”(再如
此)里留着这样一个粗糙的错误——只等时机一到,便可以纯然干脆地搞垮那位(暂时的)“替身爸
爸”——那想必不是没有充分理由的。昨天的思考,把那些理由已经摆得再明白不过。

\pmark\origpage{1178}也许有人会说,我在编故事:Saavedra所说的“帮助与建议”,并不意味着Deligne当真费神
稍稍认真读过引言里概括理论之实质的那四条命题\fnote{950}{(*)}{纤维函子上的滤套的那些结果除外——它们更技术性,更难压缩进一条醒目的命题。}。这些命题,他早在认识当事人之前就熟悉了;那么,那至多是一时轻率:为一项工作作
保,却连一刻钟的功夫都不肯花,验一验引言里预告的主要命题是否成立。可在我心里毫无疑义:Deligne确实费了这个神。因为这项工作\textbf{不是随便什么工作}——它出自一个多少有点找不着北、正
愁没论文的学生之手。继我之后(而且还在Serre之前),Deligne是最有资格感受那份形式主义的全部
分量的人:那是已故宗师留下的、未写出(至少未发表)的遗产的一个关键侧面。他尽管对这一侧面
摆出他惯常的满不在乎的架势\fnote{951}{(*)}{关于这些架势,以及它们所伺候的据为己有之术,见注“据为己有与轻蔑”(n\textsuperscript{o}~59$'$)。},骨子里他比谁都清楚所为何事。他,光辉的Deligne,极端的精英主义
者,居然肯费神跟进一个分明天资平庸的人的工作——这断然不是为了当事人漂亮的眼睛、为了帮
他拿到(按时下的共识,更按他所自诩奉行的、推到极致的严格标准)一纸\textbf{冒牌文凭}。

这话一出口,人们立刻面对一个奇怪的矛盾。一方面,一个据称多年全职投入本题目的人,犯了如
此显赫的错误,很难不把它解读为根本上无能的征兆——那个摆在他面前的问题,哪怕单看其技术
面(其实并不深奥),在答辩之时、在书出版之时,似乎根本没有被把握。另一方面,同一个学生,跟
我一两年没做出什么名堂之后,不到两年,忽然习得一门堪称令人叹服的数学素养:代数群结构理论,
一般域与\dmark{}\origpage{1179}实域上的、硬邦邦的概形理论、Hodge 理论、Motif……不止如此——我不记得读过任何
出自他手笔的数学文字,哪怕几页;而且我深知(对手段微薄的学生尤其如此)学会给数学
“成文”,绝非易事——可翻阅那本以他之名问世的书,它的“成色”之精良令我震动。我曾想,单就技
术而论,这份显然以与 EGA、SGA 同等的标准参考书自任的文本,本可以出自我的手笔,或Deligne
的手笔,或我那四五个才具出众、娴熟于把一套交织复杂的思想与事实精确、完整而优雅地陈述出
来的学生中任何一位的手笔。我很清楚,这样的行文功力,比数学素养更甚,绝非可以临时抱佛脚(除
非是天纵之才,如那位Deligne,及寥寥数人);它(倘若终究能获得)只有经年累月的练笔方能获得。
我自己就花了十年,而那时我同所要表达的实质的接触是极强的。这份接触,同Saavedra对他论
文题目的接触,当然不可同日而语——他写了关于这个题目的文字,却始终没有懂它;而这文字,后
来(至少到1982年……)竟然\textbf{是}“那一份”好的参考文献,为那套精微而关键的形式主义。分明
是两件干脆“对不上榫”的事……

昨夜起就掠过我心头的一个念头,如今,当我把处境白纸黑字对自己讲述一遍之后,带着明摆着的
力量回来了:Saavedra——我深知他的可能,尤其他的局限——不可能是那本光辉著作的作者;这
书以纯然技术的面貌(诚然,在这方面却又是穷尽的、针脚细密的),陈述一门完全超出他的“哲学”
的基础:也许头三章——其中两章主要是人人早已熟知的函子式泛论,第三章呈现的是书中中心观
念的、Saavedra版的十足的骗局——这几章,本该构成他从未完成的“最低纲领”;也许这几章全
然出自Saavedra的手笔。第三章尽管是十足的骗局,却仍足以\dmark{}\origpage{1180}让人窥见要到哪里去——即某些 $\otimes$-范畴的“Grothendieck式”眼光(姑且不点名),或“gerbe 式”眼光;它给后面的第四至
六章赋予意义。一旦接受了 gerbe 式的描述(在Deligne与Milne那双料海盗式的文本里,它被明智地
径直取作那些所谓“塔纳卡”范畴的\textbf{定义}),后三章便是要被据为己有的那套形式主义的
心脏。我猜,这几章是Deligne统统亲笔写的,或部分由他、部分由Berthelot写的;而且比我交给
Saavedra的笔记详尽得多,以致他几乎只消逐字照抄——如果曾有人要他费这番形式主义的神
的话。他准觉得“赢了”:人家“送”他一篇论文,外带一枚文凭——虽然他心里分明感到,他自己做的
那点东西(甚至是在自欺欺人地以为它站得住的情形下),对一篇国家博士论文来说,恐怕太单薄
了。而Deligne(又化装成撒玛利亚人……)赢了:这便是所需要的参考文献——纵然不为眼下,也为
“日后”(对会等待的人……),而且在那份参考文献里,那不受欢迎的名字,至少在一切实际的用处
上,不再出现。

为让欢乐登峰造极,我再补一句:那位名叫Saavedra的,似乎已从流通中消失,再无踪迹。去年,预
备寄发(我眼看就要寄出)《收获与播种》的印装本之时,我曾翻检《国际数学家人名录》——那册
子厚得很,“人人”都在(名录就是干这个的);唯独当事人不在:Saavedra之下没有,里瓦诺之下没
有,连Neantro之名下也没有(我凭良心翻过)。于是,这段历史有了黑色侦探小说的模样。想到那位
满面春风、和蔼可亲的Deligne——俨然又一位Monsieur~Verdoux(别名朗吕),以那本称心如意的“恰好的参考文献”遂
了他那番折磨人的心愿之后(比他朋友Verdier的还早四年!\fnote{952}{(*)}{关于它,见那道(理应如此)取名的注“好的参考文献”,n\textsuperscript{o}~82。})——人们不寒而栗:眼看着他让那桩恶魔般的阴谋的“证物”消失:那可怜的Neantro~Saavedra~Rivano本人,被放进奥尔马伊(Les~Ormails)一座精致的小壁炉里,长时间地焚化——那壁炉,
本就是特地为这类用途设计的\fnote{953}{(*)}{“奥尔马伊”(Les~Ormails)是法国高等科学研究所(IHÉS)住宅区的名字;我朋友皮埃尔——别名Monsieur~Verdoux,别名朗吕(还化装成献殷勤的骑士)——在那里及时接班,顶替了某位已故者:那人被排挤出这块地盘,又被我朋友所擅长的那种不动声色的暗算,抛进了暗无天日的地牢。住宅区有十来幢家庭小楼,和一栋大些的房子,带舒适的单人工作室;那些工作室想必不久也会各自配上一座小巧的全能壁炉……}。

\pmark\origpage{1181}我宽慰自己:没听说Kashiwara或Verdier也从这个世界消失——说真的,就在前天,我还拨通了后
者的电话,问他(没什么信心,似乎也没成功)能不能给我带来另一位“失踪者”的消息:人人谈论、
却看来无一人见过的——Jouanolou的论文。那篇论文,我仍然一无所知;但至少,Verdier似乎
还活着,尽管他本人就是一份“证物”——而且我怀着良好的希望:Neantro~Saavedra~Rivano也
一样。

\subsubsection{f.\ 肮脏的差事}

\notehead{176$_{6}$}{}

话虽如此,Saavedra论文史的那些于布式面向,我甚至还没有巡视完——看来那些不同寻常的论文与论文人,我是
收集上了瘾!到这里,我已有了揣测(不说“内心确信”吧):Deligne(由一位殷勤而义务的助手协助)
之所以装模作样,在Saavedra答辩十年之后一本正经“誊抄”其论文,想必只是“取回”他当年乐意
“借”给Saavedra一段时间的东西(这段时间,恰够Saavedra把论文过了、然后消失);所以这
不过是天理昭昭——只除了一点:他“借”来一段时间的,是从那位从不被点名的死者那里“借”来的。
可是,借死者的东西不还(那还了得!),风纪倒也无亏:这一头,也是万事大吉。

最妙的还在后头:即便在第二位前学生也过了这道手之后(而且还是我有过的学生中最光
辉的一位),我交给Saavedra的那个谦卑(微末)的问题——那曾是二十多年前我的出发点,
也是我自认当时就解决了的第一桩事,在所论 $\otimes$-范畴的定义环是域的情形——这个谦卑(微末)的问题,
直到今天,即便在那个情形里,居然仍然没有“被解决”!Deligne满足于点破Saavedra那个粗糙的错
误(那错误肯定十多年前就被看破了,他只是等待时机……)。他一边在128页上誊抄先前那本参考文
献\dmark{}\origpage{1182},却无心去修复这个错误。何必费这个神呢——追求的目的分明已经达到!要做到那一步,他
心里就必须在这场行动之中,除了据为己有的饥渴之外,另有别的东西:一份醒着的关切,一份对他
所处理的数学实质的\textbf{敬意},一种超出眼前“收益”的视野。

如果我在六四、六五年前后,曾费神为那些可以用“代数 gerbe”表述的 $\otimes$-范畴提炼一套“Grothendieck式”的瑜伽,而不满足于可用群概形描述的那些,那是因为在“打动”我最深的那个例子——一个域
上的 Motif——里,众所周知(凭Serre一个非常简单的论证):当该域特征 $p>0$ 时,不存在“有理
于 $\mathbb{Q}$”(甚至有理于 $\mathbb{R}$)的纤维函子。这在当时\textbf{逼着我的手},把理论表达在像 gerbe
与 lien 的形式主义这样“不正经”的东西里;同时当然,也要找到简单代数性质的内蕴判据,保证这种
“Galois式”或“Grothendieck式”的视野几乎“总是”行得通,而且花费极小。我所提炼的判据(如果
我没弄错,还给出了证明)——由 $k$ 的一个扩域 $k'$ 上纤维函子的存在给出——二十年后,文献里仍
然没有建立!直到今天,照Saavedra、Deligne之流写下的文字,纵然人们对(比方说)有限域上
“Motif 上同调类”的形式主义承认我们想承认的一切,半单 Motif 范畴(比方说)是“Grothendieck式”
的(或照这些先生的叫法,“塔纳卡的”)这一点,仍然没有建立(至少文献里没有)。前后共 418+128
=546 页的文字,出自Saavedra(由一位Deligne与一位Berthelot协助)继而Deligne与Milne之笔——
到头来却连二十年前我的出发点都没有抵达:那出发点,是让我自己相信,“Motif Galois群”\textbf{确
实存在}。

是啊,一个Deligne何必费这个神呢——他早已忘了那视野;他要的声望反正已到手;他用来做
\textbf{他的} Motif 理论(它当然同某位死者的理论毫无干系……)的那些域,全是特征零的域——于
是他那些所谓的“塔纳卡”范畴,统统是“中性”的(或“平凡”的)。照这个账,自然犯不上为 gerbe 之类
大做文章\dmark{}\origpage{1183}——那从此不过是障眼的烟幕。犯不上,除非是为了\textbf{据为己有一个}
\textbf{连灵魂与精神都已被遗忘的东西的字面}。而这场令人瞠目、可悲的故事的收梢,我看到,是这
样的:正如那套被埋葬了十五年的 Motif 视野的入门常识,还得由那个老朽来做——他刚巡礼完那
场光辉的葬礼及其绝技,就得自己把这份小活儿干了:他的学生们在他“死后”,还没有谁把它放在
心上。因为他们早就忙于扮演主子,忙得没有几天工夫,哪怕片刻,再当一回\textbf{仆人}
\fnote{954}{(*)}{我在这里有点性急,装作把我所有学生同他们当中最光辉的一位装进同一个口袋。我先向他们中那些不乐意与如此光辉的人士为伍的诸位致歉!无论如何,我高兴地记得Giraud:他硬着头皮接下了(不期而至的)读Contou-Carrère论文的活儿;那心境,对Contou-Carrère、对我,无疑是一种“服务”的心境,也许对数学共同体也是如此;关于这一点,见“耶稣与十二使徒”一注(n\textsuperscript{o}~19)最末一段,打字稿第151页。}。

\subsubsection{g.\ 一场屠杀中的五篇论文——或孝道}

\notehead{176$_{7}$}{}

(6月19日)到今天恰好两个月,我乘着一阵旋风,动手写前面那些注(4月19、20日的),带着那
个现成的名字“第六根钉子(钉进棺材)”(n\textsuperscript{o}s~176$_{1}$至176$_{6}$,不计本注——它也在其列)。此前一星期,Zoghman~Mebkhout刚把Saavedra那本书带给我——只消看上一眼,我就已经明白是怎么回事了。

我必须承认,这一发现带来的震动,几乎不亚于一年前(整整一年前)那本 Motif 掘墓“难忘文集”
(《数学讲义》第900种)带来的震动。更确切地说:去年的那份震动,某种意义上重新出现了,被这
个与那场掘墓紧密相联的“运算”的发现猝然重新点着;这场运算(一开始就明摆着)曾为它作了准
备,规模也相仿。我再一次被那种安详的无耻之感攫住(简直要说被噎得透不过气)——\textbf{同样的}无耻(这一点,凭种种不会
骗人的迹象,一开始也明摆着)——冲着一件与我血肉相系的东西,一件世上除我之外无人长久怀抱
与滋养的东西……那感受如此强烈,近乎焦虑,连我自己都吃惊。

\pmark\origpage{1184}自发的反应、自然的出口,本该像去年那样:趁感受还新鲜就说出来,由此径直进入这新的侧
面——那场由我的亲近者对我施行的活埋的新侧面。可我忍住了\fnote{955}{(*)}{我到底还是乘着当时的激动写了四五页,但九天之后(4月19日)写下正文时,那些痕迹已所剩无几。};因为Mebkhout的来访,需要我至少有一些闲暇;何况他
有些话要对我说——即便不那么触动我的要害,我感觉得到,它们对他而言同样是“要害”,对\emph{葬
礼}同样意味深长。再说,趁新知还热,把它们记下来,在我看来也重要;而那本出名的“葬礼之书”的
来龙去脉,即便我日后再来清理,也不会溜掉。所以,朋友刚走,我(4月15日至18日)就去写他的倒
霉遭遇,成如今“封神”末尾的那组注(n\textsuperscript{o}s~171$_{1}$至171$_{4}$)。

也就是说,在来到那出名的“第六根钉子”之前,我已经有时间恢复镇定。说实话,刚才重读最初几
页,在我对那桩新揭发的丑闻所作的挖苦(而稍带距离)的描述里,我找不出最初袭我的心绪的痕迹
——那心绪曾让我失眠一夜,而我当时多么需要睡眠。这一回,我是真切地感到了那份“过去的重量”!

那是6月10日,在给\emph{葬礼}画上那个出名的“句号”之后三天——葬礼随之重新开动!当然,我远
没有料到它开动到了什么地步——还有(差不多)三百页要写!写完“第六根钉子”的第六道注(“肮
脏的差事”)时,我满以为连“四种运算”也一并走完了——只除十来页(第三、四场运算的)要重抄
清稿、补上预定的脚注。再过几天,我想就可以把葬礼(III)的全部手稿送去打字了。

然而,随后几天里(也许就在我以为钉完最后一根“钉子”的次日或隔日),发生了一场出人意料的
剧情突变,对此我还得回头再谈。这里,我自发的冲动本该是立刻扑上去。我又等了两个月才动手,
当然不是不情愿\dmark{}\origpage{1185}——是有更急的事要备好送去打字。从头重读《四种运算》,发觉这里那里大有充实的必
要——后事如何,大家都知道了!

不管怎样,今天(只要不再有新的意外——摸摸木头!),终于到了那个黄道吉日:我给\emph{葬礼}
画上真正的句号——照实说了:写下最后几页的日子;它们理应构成我在《收获与播种》之中关于
葬礼的思考的尾声。此后,只还剩写那封充当《收获与播种》前言的“信”——然后我想休息几天,
受之无愧,而且也实在需要……

写完前面六道注几天之后,我得知了Saavedra答辩委员会的组成——那支我在“Monsieur~Verdoux——
或献殷勤的骑士”倒数第二道注里痛加讥讽(罪有应得)的委员会。论文于1972年2月25日在奥赛的
理学院答辩,委员会由\textbf{Demazure}(J.~Demazure,报告人)、\textbf{卡斯泰勒}(Castelle)与
\textbf{Grothendieck}组成。

要说“剧情突变”,这才是剧情突变!于布戏的加冕礼!而且,面对这份官方来源的信息,我好不容
易才相信——我连出席过这样一场答辩的一丝记忆的影子都没留下。Verdoux-朗吕先生的故事,愈发精彩了!我姑且给
Demazure打了个电话,问他可记得同我一起,给一个叫Saavedra的人当过论文答辩委员会。Demazure也记不得多少了,但好歹足以向我保证:答辩确凿举行过(何时如何,他自然说不太清),我们
两位都在场,还有卡斯泰勒(连这名字我都不记得了……)。他知道的也不多,只记得自己是论文报
告人。还是我告诉他:论文官方说来,是一份25页的文本(想必这让他的报告人差事轻松了不少)。
这一来,吃惊的反倒是他。他答应想法给我寄一份论文来。我本来很想知道它长什么样,可至今没有
等到——看来(照Demazure几个星期后终于对我说的)这份论文已无处可寻;也许他也没怎么费神
去找。反正,同我一样,他的文件里也毫无踪迹。不过,这只是个细节……

\pmark\origpage{1186}这一来,我可够瞧的了!我对那支分明无能的委员会——“一本正经地摆出架势评审一项他
们‘未必弄懂了所为何事’的工作”——所发的那些冷嘲热讽,如今都泼回我自己身上!可以想见,我多么想收
回那些讥讽,保住体面,撑住场面——可是不,那是作假。整场\emph{葬礼}里的作假已经够多了,犯不着
再添上我的一份。再说一遍:那些讥讽完全是有道理的。如今知道了委员会的组成,我甚至可以指明:
正是我,先于任何人,最应当受这些讥讽。毕竟,Demazure与卡斯泰勒当时记住的,主要是:这篇论文,
Saavedra是跟我准备的,或至少是跟我开的头,题目是我给的。按说“在局中”的是我,而他们信任
我。说不定,那出名的25页——Demazure据称做过报告人的——本来是站得住的;即便那桩显
赫的蠢事也在其中,它也不过是一份理论概述,不在局中、信任我的Demazure,绝无机会察觉。

至于我——除了上课,两年来实际已脱离数学——这场答辩,我草草了结,想必不过是在奥赛的两堂
课之间、某场“生存与生活”的集会、或某场关于附近(萨克莱)核废料的公开讨论之间,纯然一道
行政手续。确定无疑的是:两年来,Saavedra的工作我早已没有跟进,任何人的也没有;而且我对
Saavedra的工作站得住,毫不怀疑。这份确信从何而来,如今说不清了。同此前我所有的学生不同,
对Saavedra,我没有任何直接的凭据——凭他跟我一起已完成的工作——来判定他的严肃。难道那
些年,我竟把我的学术职责轻率到这般地步,凭他的面色,就信任了他?倘若那本书(同年出版,那
25页的论文想必是它的摘要)当时已经杀青,并让我形成了印象,那“看一眼”确实印象极佳,说
不定我连想都没想过要去核实:那本该构成Saavedra个人贡献的部分。也可能,甚至更可能(此事
我毫无记忆了),我信赖了Deligne的判断——我离去之后,是他跟进了这项工作
\fnote{956}{(*)}{我连Deligne曾跟进Saavedra的工作这件事,都没有记忆。这是4月里,看Saavedra那本书的引言时才知道的。}。

\pmark\origpage{1187}无论哪种情形,我都必须承认:我负有同样的责任——我把科学博士的头衔,授予了一篇二十
三年后显出是\textbf{冒牌货}的论文(借已引之注的说法)。然而,我自己在不知不觉中成了这场骗局的
工具、并因轻率作保而负有责任,这一事实,并不改变骗局之为骗局。它反倒显得愈发天才。因为说到
底,牵线者的真正动机,当然不是让一个落魄的论文新手廉价地得一顶头衔,然后改行退场——而是
让一个绝不糊涂的人,细腻地、不动声色地,把某种在我之内诞生、并且在他还没有听到(在数学
里)说出“gerbe”或“Motif”这类字眼之前,就被我孕育至成熟的那份视野的父亲身份,据为己有。正是趁
我骤然全力投入物种生存与其他种种最紧迫的美好事业之时(那位前学生兼朋友曾对我说,他不得
不同这些保持距离,因为他全身心绝对奉献给唯一的数学
\fnote{957}{(*)}{关于这一点,见注“兄弟与配偶——或双重签名”(n\textsuperscript{o}~134),尤其打字稿第614--615页。})——趁我的精力完全被别处占尽的时刻,我那位光辉的学生兼朋友,完成了这桩
真正独一无二的障眼大戏:让我成了我自己被褫夺的工具!我当时的心境,完全切断了对昔日数学兴
趣的牵挂,对我那些学生——Deligne首当其冲,他们自 SGA~5 讨论班结束起,就已开始按自己的路
数玩起了小把戏——报以盲目的信任:随便什么名字,只要安到那些我只剩遥远记忆的出名范畴上,
我都会说:是,阿门!正如我对Verdier宣布不会再有新式同调代数的书时说了是、阿门;正如Deligne向我宣布,我们合搞的 SGA~7 讨论班的一半,将突然更换父亲身份时,我说了是、阿门……然而,
受骗者乖乖地点头、毫无所觉地同意,并不改变诈骗的性质——只给它添上一重背信。而当Serre之
流的其他占卜师也从中各得其所、毫无保留地祝福其事
\fnote{958}{(*)}{这份再明白不过的祝福,见“家庭相册”一注部分 d.(“入殓——或自然的坡道”)。},这件事便有了非同寻常的维度\dmark{}\origpage{1188}——它成了一个圈子、一个时代的集体败坏——但它并不因此就光彩起来,无论多么天才;它的无耻,一分也不少。

正如剧情突变式的好事从不单行:得知我前学生Saavedra答辩委员会的组成才几天,关于Jouanolou论文的相应信息也到了——那也是一篇有点特别的论文,我在思考中这里那里谈到过它
\fnote{959}{(*)}{这篇论文,在注“团结”(n\textsuperscript{o}~85$^{1}$,打字稿第349页)的子注,以及“共同继承人……”(n\textsuperscript{o}~91,第387--388页)一注里谈到过。另见“学生与纲领”一节(n\textsuperscript{o}~25)。}。同Saavedra一样,他也从未想过给我寄一份他那出名的论文(“人人都引(自反常层会议以来),却无人见过”)——于是我最终给他写了封略嫌干巴的信(4月25日),就这
篇论文的离奇遭遇问了他几个问题。他几乎当即回了信(5月1日):对实质性的问题闪烁其词(因为
“重提往事总归非常痛苦”),但行政细节却精确之至:论文于1969年7月3日在法国高等科学研究
所(巴黎)答辩,委员会主席为\textbf{P.~Samuel},考试官为\textbf{迪斯米耶}(J.~Dixmier)、
\textbf{Grothendieck}与\textbf{Verdier}。我的通信人还添了一句,不无恶意:“据我所知,委员会全
体成员都到了!”(此事 J.L.~Verdier不久后也在电话里向我证实。)

\pmark 又是如此:对这场答辩,我连一丝记忆的影子都没有——它分明也是“偷偷”过的(不得不这
样糟蹋我自己的形象,抱歉!)\fnote{960}{(*)}{答辩的日期,我记得,应在我已经“脱离”数学、转而关心生物学(尤其分子生物学)之后不久。}。我以为答辩是在斯特拉斯堡(因而把它放在七十年代初——知道Jouanolou那几年在斯特
拉斯堡有职位),多半是因为Verdier在这篇论文上的一条隐晦的引文(1975年2月Bourbaki讨论班
讲义 n\textsuperscript{o}~464),引作“J.P.~Jouanolou,~Thèse,~Fac.~Sc.~Strasbourg”(无日期,无题目)。可他是答辩委
员会的一员,同我一样——他的记性也同我一样失灵,或者说得更准,任性,把举行答辩的 IHP\dmark{}\origpage{1189}(Henri~Poincaré研究所)安到了斯特拉斯堡?谁懂,谁就懂吧!

还是这位Verdier,好意把自己持有的那份论文寄给了我。起先,看着这包208张散页\fnote{961}{(*)}{索邦大学论文管理处存有一份215页的存档本——Verdier的这份看来缺了六页。说不定,存放在该管理处的那份,是世上唯一完整的一份——而且,人家还向我保证,是装订成册的。那里想必有一个专为弃儿论文装订的部门,它们都是散着零件到的……},我以为那是一份草
稿的复印件——我记得当年Jouanolou跟我做这篇迟迟不完的论文时,我亲手捧过并细细评注过的正
是它。可Verdier向我确认:这确凿是那篇倒霉论文的定稿本;它看来从未荣幸地印过三四份以上
(我那本,连同我的批注,想必回到了Jouanolou手里,从此再没见过……),也从未装订过。

后来,Jouanolou惠赐的稍详细的说明(6月3日的信),加上给Verdier的电话,让我重新入了一点
水。显然,Jouanolou对他的论文到达了一个“饱和点”:他从一开始就毫无信念地做(而我也懒得把处
境看清楚\fnote{962}{(*)}{见已引的“学生与纲领”一节,n\textsuperscript{o}~25。})。到1969年,他想必到了一个堵点,以致无力哪怕稍作修改,来照顾我的
许多批注。我只好认了,随它去。反正,重读一遍之后,我觉得这份文本代表了一份认真的、可用的
成形工作,虽远非完美——分明好于“聊胜于无”,在没有更合我意的文本的情况下,可以充当一份
必不可少的参考文献\fnote{963}{(*)}{不管怎样,它至今仍是世上唯一的、以导出范畴的版本陈述$\ell$-进系数理论的文本——而且偏偏无处可购,喜上加喜。电锯从那里经过了……}。

\pmark 当然,我做梦也想不到:Jouanolou会以他自己的方式,报复他跟我做这项工作时的三心二意
\dmark{}\origpage{1190}——亲手凿沉它,几乎抹掉那份我如此珍视的出名的“参考文献”的一切痕迹!这又是
一次“往事回头”,我实在没什么可抱怨的(尽管怨气并不是没有!)。在同Jouanolou的关系里,我要
的是在他身上找到“臂膀”,来推一辆规模可观的大车的轮子。我预先认定,他Jouanolou是我事业的参
与者,却从来没有停下来想一想,那些执拗的迹象分明告诉我:根本不是这么回事。诚然,是他自己选
择来跟我工作的(同一位名声显赫的“老板”共事,想必他觉得划算,却没料到自己陷进了什么……);
也是他自己,从我不吝扶持的一大片题目里,自由选定了他的工作题目(这些题目,当然全连着那辆
“大车”——而那车,骨子里多半引不起他半点兴致)。换句话说:同每个人一样,Jouanolou也在他自己
之内、在他自己的愿望与抉择的层面上,同某些矛盾搏斗——这里,是在他的工作里。

我自己的矛盾,不在我同工作的关系里,而在一种对任务的如此的偏执,以致我在学生身上只看得
见受欢迎的臂膀,无法想象他们中任何一位会在我与之共事的工作中心怀分裂。有了对\emph{葬礼}
的长考所给的这一份额外的回顾,我如今发觉:在我的学生里,“分裂”的远不止Jouanolou一个。但他
是一个极端的个案:他是其中唯一没能认同自己所选任务的人;他的工作是在没有信念、没有欢乐
之中完成的。我在这一处境中的责任,在于我不肯真正去知悉它,宁可把本该次要的东西(完成我的
任务)\textbf{摆在前头},而把要紧的东西(学生所“选”的任务真正\textbf{属于}他、并且怀着欢乐去
做)撂在后头。

这想必就是为什么,Jouanolou也是我唯一曾在其身上察觉到怨怼的前学生(当然,那怨怼从不报姓
名)。怀怨是一个出口、一剂解药;它当然一无用处,只除了回避自己的问题(而人们很少往深里追
究)。可它毕竟是事出有因的;今天(二十年之后)我收割到它的某些果实,无权抱怨。

\pmark\origpage{1191}不到两个月里,接连面对Saavedra论文、继而Jouanolou论文的这些不寻常的插曲,让我切
身感到一件在《收获与播种》第一部分里才刚刚瞥见的事:早在我离去之前,以及紧接着的那些年
里,我的学生们和我之间,并不像我想当然的那样一切安好。譬如,在跟我做到国家博士论文层级的
学生们所通过的十二篇论文中,\textbf{四}篇赫然是针对宗师的“葬礼论文”!它们相继出现在1967至
1972年的五年之间,其中两场“葬礼论文”在我离去之前就已举行。第一场是Verdier1967年的论文,
缩减成一份28页的摘要,是我引入的新同调代数被埋葬的前奏;那套代数,Verdier曾受命去发展。
此事已有相当周详的交代\fnote{964}{(*)}{关于这一点,尤其见“赊账的学位论文与全险”与“随心所欲的荣耀——或暧昧”两条注(n\textsuperscript{o}s~81、170(ii))。},不必重复。第二场是Jouanolou1969年的论文,它为$\ell$-进上同调的形式主义——从导出范畴
与三角范畴的观点看(这对六种运算分明至关紧要;Verdier本该提供基本参考)——的埋葬盖了
章。第三场是Deligne1970年的(?)论文:一篇十足光辉的论文,也深深扎根于他从我这里得来的思
想\fnote{965}{(*)}{这是Deligne的“Hodge 理论 II”。关于这项工作在 Motif 瑜伽、在我的“系数理论”视野(包括一种“霍奇系数”的理论)之中的扎根,我在注“把一个个 i 点清”(n\textsuperscript{o}~164,尤其打字稿第739--740页)以及子注 n\textsuperscript{o}~164$_{1}$(第805--806页)里给出细节。同 M.~Raynaud与 C.~Contou-Carrère一样,Deligne自选了他的工作题目——包括论文题目——并没有等我给他出题,而且完全独立地做了这项工作,甚至在几乎完成之前都没跟我提起。尽管如此,他的工作(论混合 Hodge 结构)扎根于我的思想,比Raynaud与Contou-Carrère更深:后两位主要使用我带来的语言与技术,而他们各自追究的问题域,则是全然原创的。
的确(照今天刮的风),思想随风而去——尤其当它们没有发表的话(几天前,Serre又斩钉截铁地向
我解释过一遍……)。}——而我的名字在文中连提也不提!第四场是Saavedra的论文——方才长
篇谈过:不是那位推定的作者\fnote{966}{(*)}{这至少是我在“Monsieur~Verdoux——或献殷勤的骑士”(n\textsuperscript{o}~176$_{5}$)那一注里达到的确信。},而是别人,以他为人熟知的娴熟手法,陈述了第三个人关于 Motif Galois群的思想与\dmark{}\origpage{1192}结果(经由一套完整的所谓“塔纳卡”范畴理论——而且是整整四种!),却对我这卑不足道、亡故之人只字不提!

这四场葬礼式论文(它们正是“四种运算”——大写的那种!——的序曲)分明以许多方式互相勾连
\fnote{967}{(*)}{往深里探测这些联系,当然有趣——但如我在几行之下所说,现在不是时候。}。它们相续于不到五年之间,始于紧接 SGA~5 讨论班结束的那一年。那场讨论班,看来确是
我前学生们身上掘墓心性的起点与集结之地——而且早在我离去之前!这些心性先于我的离去而存
在,这是一个耐人寻味的情况,关乎\emph{葬礼}的那个“第二平面”——由“从前的”全体前学生组
成;这个情况,我还没有真正整合进一个总体的理解。眼下,这“第二平面”似乎是三个平面中最不被
理解的一个。但现在不是重新开启思考的时候;接下来的几个月,想必不会缺少新的材料,它们会由
我的前学生们亲手送来。到那时,再把它们拼成“第二平面”的一幅活的总体图景。

还有第五篇论文\fnote{968}{(*)}{在共十四篇论文之中——十四位(无论“从前的”还是“后来的”)跟我做到国家博士论文层级的学生。也就是说,按数目算,\textbf{三篇之中就有一篇以上}是葬礼论文——已经不算糟了!},对我而言,它也插进葬礼论文的序列,不过是一篇“之后”的论文——比前一序列晚了十
年。那是Contou-Carrère的论文,1982年12月通过;它也在不止一个方面很特别。它不同于前四篇之处在
于:Contou-Carrère为讨好要人们、为求得原谅自己曾多多少少当过我的学生而作出的英勇的掘墓努
力,并没有能避免Verdier(他以为选他做论文导师是聪明的
\fnote{969}{(*)}{而且那时候,我还一直(照Contou-Carrère本人向我保证的)以为自己是他的官方论文导师。我得知还有一位“平行”导师的存在——在这一对里,本该是我充当“备位”导师,以防万一……——是在Contou-Carrère不得不回头找我的时候;与此同时(局势已经变得有点太那个了),他向我揭示了Verdier扮演的角色。经年累月地玩着这种没谱的把戏,Contou-Carrère最后几乎不再做数学,难怪。可以说,他不是唯一的一个……})突然装模作样地要把他\textbf{搞沉}\dmark{}\origpage{1193},不宣而战——于是,不得已,他又回头找我。我本不必
出场充当论文导师:Contou-Carrère自己找到了题目、自己发展了方法,我也没有跟进他的工作,何况
他的工作所在的语境(约化群概形)我已有些生疏。尽管如此,他工作的初始想法——为 Schubert 闭
链的闭包搞一种“等变”的奇点消解方法——直接受启发于我(约1975或76年)向他原原本本解释过的
一个想法:关于约化群在自身上的伴随表示,其轨道闭包的典范的、同步的奇点消解
\fnote{970}{(*)}{六十年代末,我曾被Brieskorn关于所谓“有理”(曲面)奇点的漂亮工作、及其与某些单根系(根全等长的那些)的联系所吸引;我给自己出了一个问题(不用说,有点想入非非):用与根图相应的单代数群,直接描述一个有理奇点。就这样,我得到了所论奇点消解的一份非常简单(说白了,甚至是显然)的几何描述,靠的是 Killing 对;外加一整套精彩的猜想作为钥匙,后来我有些忘了,当年讲给过凡愿听的人。可我什么也没发表;而照Serre刚向我惠赐的新公理——归第一个捡到的人所有——我发觉,有人捡了不少,难免。有时候,换换公理,真方便……}。不用说,Contou-Carrère早就嗅准了上流社会里风往哪儿吹——那上流社会,他正合法地盼着跻
身;对这层渊源,他守口如瓶。倘若大家又都开始提这类不可捉摸的东西,那还了得:一个\textbf{观念}
(还没发表呢),据说是为了\textbf{激发}另一个(您给评评理……)——当然,除非被光荣引用的,正
是那些让所呈工作增光的名字之一(到那时,自然也无须解释为什么致谢滔滔,那致谢想必是名正
言顺的……)。

\asterisms

\centerline{\headfont “四种运算(对一具尸体)”终}


% ================================================================
%  18.6 Les chantiers désolés (原书 p.894–908)
%  源文: _work/chunks/chunk131.md, chunk132.md, chunk133a.md, chunk133b.md
% ================================================================
\section{荒废的工地}

\subsection{(1) 悬而未决之事}   % 18.6.1

\notehead{176$'$}{}

\pmark\origpage{1194}(3月25日)昨夜,我在床上辗转了几个钟头,让自己重新浸入“Motif 的瑜伽”里,而不是像
本该的那样安然入睡。稍后,我也没有接着写笔记,而是又花了一两个钟头,潦潦草
草画出一堆蕴涵图,涉及我所知的、一个德拉姆上同调类(比方说,取特征零域上一个非奇异射影簇的上同调类)
要成为“代数的”所需的那些内蕴条件。我一共找到\textbf{十二}种路数,统统是 Hodge 猜想和
Tate 猜想的变体\fnote{971}{(*)}{(3月27日)这十二种变体,应当每一种都能对任意一个特征零的基本概形 $X$ 给出 $X$ 上一类相应的“系数范畴”(“系数类型”一词,取注“墓上的旋律——或自负”(n\textsuperscript{o}~167)中所说的含义)。倘若所论猜想为真,这条系数范畴应当把 $X$ 上的 Motif 范畴作为一个全子范畴(三角的)包含进来(这个猜想,在 $X$ 取一个域的谱这一特殊情形下,说的恰好就是这件事……)。细节我留给《数学反思》第3卷中专论 Motif 理论的部分(“Motif,吾爱”)。这也等于说,这十二种人人皆知的猜想的变体,给出(至少先天地)同样多种特征零域上“Motif”的不同观念。这就让将来会有十一位我朋友皮埃尔的摹仿者,各自“发现”属于自己的那种 Motif 观念,同时装作不知道别人的——尤其是(十五年来照例如此……)装作不知道某位已故者的,他主要正是以偏爱无用细节而著称……}。与此同时,我也说服了自己:要在一个 $\mathbb{Z}$ 上有限型的概形上定义“那个”(三角的)Motif 范畴——或者至少给出它的一个极紧的近似(假如这仍不是“那一个”对的)——所需的材料,应当差不多已经在手,只消再有一套“神秘函子”的理论即可;这个函子,是我在六十年代末曾公设过的\fnote{972}{(*)}{这个“神秘函子”的问题——它在特征 $p$ 的晶体上同调(经由“滤套 $F$-晶体”观念,$F$ 即“Frobenius”)与特征零的 $p$-进上同调之间架起“缺失的一环”——对我们的代数簇上同调认识而言显然至关重要,可将近二十年过去,自我以再明白不过的方式把它提出以来,这个问题至今无人认真问津……}。

在这里就此铺开,当然不合时宜。可我现在已经看清:正是在此时此刻,眼看 Motif 这个主题在我
把它留在一伙靠不住的人手里十五年之后,落到这般可怜的荒废地步,把当年我已抵达的那些思想
的几条主线勾画出来,恰是其时。我没有心思再等下去,等哪天有了闲暇(等《追寻栈》写完之
后)再去写那部本该写的“\textbf{那}部”系统著作;那份对一个梦的原原本本的讲述,本该作为让这
个\textbf{梦}终于扎根于精心酝酿成熟(并且发表出来……)的表述沃土之中的第一个\dmark{}\origpage{1195}大步骤,好让它
依着自己的本性开花结果。除了已为这部“数学虚构”之书预告并宣布过的第一块界石——即所谓
“六种运算”的代数对偶形式主义的一个概要——我还要在《数学反思》第3卷里附上一篇短文,提出
几个与 Motif 和代数闭链相关联的关键问题\fnote{973}{(*)}{照我目前的出版计划,《收获与播种》的前四部分(以\emph{葬礼}的第三部分亦即最后一部分作结)应构成《数学反思》的第1、2卷。第3卷将由《收获与播种》的第五部分(读 C.G.~Jung自传的笔记)和若干较短的文本组成,其中多数已在《引论》中预告。这样,《追寻栈》的第一卷就预定为《数学反思》的第四卷。}。眼看它们在坟墓里凋敝,我于心不忍;我巴不得看见它们重见天日,重新加入季
节的轮转……

思考回到\emph{葬礼}上来,已经五个多星期,而且再不肯离开。想必正因如此,“孤儿们”的念头——
他们被抛在一个病了的世界里,无人认领——近日又执拗地回到我心头。上一次有根有据地谈到其
中一位孤儿的笔记,是“墓上的旋律——或自负”(n\textsuperscript{o}~167),主题与昨夜以及稍前的这段
Motif 思考正相毗邻(我方才提到过它)。那正好是一个月前,就在我动笔写一条预定(这早已预先
定好)称作“四种运算”的笔记的前一天,我还不知道等着我的会是什么!到头来,一条笔记变成了
十六\fnote{974}{(*)}{(5月9日)写下这几行才不过两个星期,临头冒出的新事实,又把“四种运算”这场调查重新推动起来;它已经添了足足二十来条新注和子注。},我还以为永远做不完了——可到头来,我终究还是把这几件拖泥带水的“运算”走完了一遍!

而在这里,我最想做的,首先是回到那些孤儿身边,至少逐个叫出他们的名字——这也许对他们
有点好处,对我则一定有。第一次谈到他们,是一年前,就在那道同名的注“我的孤儿们”里,去年
3月底,与紧随其后的那道注“拒受遗产——或矛盾的代价”(注 n\textsuperscript{o}s~46、47)一口气写
成。写下这些笔记、给它们起这些名字时,仿佛有一股冥冥中的先知引路,我还没有料到,我留下
的那些东西,究竟“孤”到了何等地步\dmark{}\origpage{1196}——其含义之沉重、之切肤,连做梦我也想象不到;也
没有料到,我当时初次怯怯地指出的那道“矛盾”,会延伸到多远。这段记忆立刻又勾起另一段,是
前一个月的:我看见自己写下——仿佛是另一个比我更透彻的人借我的手在写——“\textbf{谁也斗不
过腐败}”。那是在写“无爱的世界”一节(n\textsuperscript{o}~19)的时候。我还记得,看见“\textbf{腐
败}”这个词白纸黑字地立在眼前,我起初愣住了。我心里某个“讲道理”的声音在训斥我:真是的,
你下手未免太重——“腐败”可是个大字眼,可不能信口开河!你还是换一副腔调吧!

我着实自我盘问了好一阵子,也许有几分钟。随后我明白,我既不会改掉这个“大”字,也不会加
一条注,解释说这个词是笔下一篇酣畅之际溜出来的,不必太当真。从那个世界里,这里那里,陆
陆续续回到我心头的这些“呛人的气息”,我心底有个人——比那个负责挑选“讲道理”标签的“我”
更有眼光——早就一清二楚它们意味着什么,根本不等我费神替自己编个说法\fnote{975}{(*)}{这段经过,我先是在去年3月“那张条子——或新伦理”一节(n\textsuperscript{o}~33)里讲过;两个月后,在发现\emph{葬礼}之后,又在 Notes 93--97 那一组原原本本得多的笔记里讲了,它们构成送葬行列之十“灵车”(与掘墓人同行)。}……

我也清楚记得那一天,那天的思考突然变了质地、我心中那\textbf{另一个}接手执笔的确切时刻。
那紧接在一段追忆之后:我数学生涯最初几年,数学圈子里那团裹着我的温情——多亏长者们给予
的接待,一直深到他们家里:Schwartz家、Dieudonné家、戈德芒家……转折就发生在我接着写下“显
然,对今天许多年轻数学家来说,要命的正是被切断了……与一切情感的潮流、与温暖的联系
……正是这一点折断了工作的翅膀,又夺去了它比一份郁郁寡欢、朝不保夕的营生更深的意义……”——就在
此刻,这个\textbf{无爱的世界}陡然在我眼前浮现,活了起来,再一次向我发问……

去年,“我的孤儿们”这个名字,是不请自来地落到我离去时留下的一切头上的(它们被我托付给
的那些近亲宣告了“死亡”……)。想必是因为,这个名字\dmark{}\origpage{1197}道出了一种\textbf{实实在在}、触手
可及的真实:我所“留下”或“托付”的,既不是“物件”,也不是“财产”,而是一些\textbf{活
的东西}。如今想起它们,我依旧把它们当作活的东西来想:茁壮、丰饶、生来就要生长、
绽放、又要孕育和催生别的茁壮而丰饶的活的东西。我确实感到自己留下了一份“财富”,可那不
是银行家的财富,而是园丁的、砌墙匠的财富:他们用双手让这些繁花似锦的园子、这些宽敞好
客的房子从地里长出来。这种贵重(乃至脆弱)之物的感觉,最紧地系于那些\textbf{观念}、那些
\textbf{问题}、那些我知道生机勃勃的\textbf{大主题}——我把它们留在了更年轻的人手里;这些东
西还需要劳作与呵护——远甚于那些我打磨到位的工具,或者那些我一路盖到头、布置停当的“房
子”\fnote{976}{(*)}{关于我心里那股催我“盖房子”(数学的房子)的冲动,见注“阴,作为仆人——与新主人”(n\textsuperscript{o}~135)。}。自有别人照自己的意思去那里忙碌,做他们自己的饭,享他们自己的清福;哪间房子嫌小
了,他们会照自己的需要扩建——就像我自己也常常不得不一扩再扩,虽然当年看上去,我“眼看
得很远”了。可是,我的数学家过去之所以仍旧攥得住我,靠的恰恰是\textbf{那些悬而未决之事},
是那些刚刚在绝好的地基上、用那么美的石料开了工的工地(而工人们已经走了,拿走了他们中
意的,糟蹋了剩下的……)。正是这些被撂荒的\textbf{工地}——今天我重访时,只见它们横遭劫掠、
残破不堪——我现在要把它们逐一巡视。

\subsection{(2) 吝啬鬼与老朽}   % 18.6.2. L'avare et le croulant

\notehead{177}{}

(3月27日)昨天一整天都交给了庶务。为了送去打字,我不得不把\emph{葬礼}第三部分亦即最后
部分的开头五十页重读一遍——连这里那里的小小措辞修整,外加又添上的几条脚注,足足花了五
个钟头。“阴阳之钥”的打字眼看就要完成。自从为这部分打字吃了那些没完没了的苦头之后\fnote{977}{(*)}{关于这些“苦头”(这还是委婉的说法),见注“祈祷与冲突”(n\textsuperscript{o}~161)的开头,以及《收获与播种》第五部分即最后一部分的开卷之注“Jung——或‘恶’与‘善’的循环”的开头。},我最终退而求其次\dmark{}\origpage{1198},找
了一位系里的秘书,在正职之外替我做。苦头到头了,谢天谢地——她做得又认真又麻利,一
个星期三十来页,无可挑剔。总能到终点的。也该到终点了!

除此之外,在一个 $\mathbb{Z}$ 上有限型的概形上,以三角范畴的形状构造 Motif 范畴这个问题,
继续在我脑子里打转——我又在床上熬了大半夜想它,没有睡觉——当心啊!起初似乎我的想法只
对特征零的概形行得通(比方说 $\mathbb{Q}$ 上有限型的),在基底 $\mathrm{Spec}(\mathbb{Z})$
本身上就不大灵;随后我记起,六十年代我就原则上确定过有限域上 Motif 范畴的结构。假定把我
当年做的工作明确化,我终于看到,一般情形下至少有一个完整描述的原理显现出来——说实话,
够紧凑的,但在我看来绝非不可企及。比起我六十年代的想法,唯一的新配料,是Mebkhout的哲学,
它凝结在那条说来奇特的“好上帝定理”里。除此之外,我拿来作假设性配料的,是“神秘函子”的理
论。眼下它还拿不出来,这肯定不是因为它“无从下手”(借一个我早领教过的说法\fnote{978}{(*)}{这正是我那位杰出的前学生Deligne用以埋葬“标准猜想”的斩钉截铁的考语——于是,在将近二十年里,我那些大胆的同时代人竟没有一个敢向它们下手。完整引文见注“悼词(1)——或恭维”(n\textsuperscript{o}~104)。}),而
是因为,我所认识的那些在代数簇上同调里做工的人,连在数学里也丢掉了对本质之物的分寸——
当然,他们为一场需要悉心照料的葬礼忙得不可开交……

平心而论,Deligne围绕Weil猜想做的工作——“Weil~I”,尤其是“Weil~II”——终归会在紧要关
头派上用场,到要为那些被指望表达 Motif 的系数范畴搭建六种运算的时候。尽管如此,还是得
有一位“糊涂”的、如人所愿地老朽不堪的已故者,在十五年之后灵光一现,从他那些亲爱的学生
与继承人乐于把他安置其中的软壁棺材里坐起来——而他本人对这一切毫不知情,而且差不多把生
前懂得的那一点点也忘光了——这样,在基底概形 $S$ 上描述 Motif 范畴的问题才算\textbf{一字一句地提
出}来,而搭造的原理(把与一个 Motif 相关联的一切已知结构要素都考虑在内的那种)也才顺带\dmark{}\origpage{1199}——纯属巧合地——终于得到明白的表述\fnote{979}{(*)}{正如我昨天在思考里预告的,我打算把这一描述收进《数学反思》的下一卷,附上一幅“Motif 广景图”的(极简略的)总轮廓——暗中的 Motif 把戏,闹得够久了。这里就预先声明:所设想的构造原理,不依赖代数闭链上任何一种猜想,诸如“Hodge”或“Tate”(或者昨天谈到的那十二种变体之一)。}。

1982年那部关于 Motif 的“难忘文集”之后,这份“Motif 藏宝”——十来年间一直是一个人的专
属禁地、一个秘密——似乎变成了三四个人共有的宝藏;这几位彼此往来,一派合谋者的神情,或
者像某个秘密教派的入会大佬。可是,把几个简单的问题白纸黑字写下来、提交给公众注意,不过
需要几天;要是想稍加用心地把它们界定清楚,明白交代手头有哪些配料、还须发展哪些配料,也
不过几个星期。从1970年起十五年,从“难忘文集”起三年,先是一个人,后来是那么几位,竟没
有一个肯从自己固然宝贵的时间里拿出这么几天——更不必说几个星期——这必定有极好的理由,
尽管他们当中没有谁打算去探个究竟。可他们煞费经营的那种气氛、他们赖以自持的那种心气,本
身就已经是一场发现之旅的堕落:发现之旅沦落为凌驾于他人之上、乃至蔑视他人的手段。这样的
气氛天然传播腐败,与创造背道而驰——哪怕乐在其中者是天才里最耀眼的几位。守着这样的心境
——一头孵着金币的守财奴的心境——他们切断了自身之中的创造之力,正如他们乐于扼杀他人身
上的这种力量。

\subsection{(3) 巡视工地——或工具与视野}   % 18.6.3. Le tour des chantiers - ou outils et vision

\notehead{178}{}

(3月30日)前天我满五十七岁,稍稍歇了歇工。只为“阴阳之钥”的结尾做了一点打字校对,昨天
又继续了一点。这是一桩让人歇息而愉快的活儿——至少,当打字的人自己也肯上心、一份我全身
心投入的稿子不至于被弄得面目全非回到我手里的时候,是如此。这回,我给自己放了两天假,\dmark{}\origpage{1200}细
细重读了五十来页清稿,就为了这里那里揪出一个还没放对地方的逗号……

不过,工作的劲头并不在顶点。几个星期以来,我心里有一种悲哀在提醒我:比起把手头这些笔记
一路写到它们自然的终点,还有更要紧的事在等我。我写着,像在逆水行舟;然而我知道,只要不出
意外、不遇不可抗之力,我只有给\emph{葬礼}终于画上那个句号,才会停笔。可把这股悲哀压下去、
放逐出去——它于是重得像一块石头——不让它在这些笔记里发声(除了眼下这一刻,顺带地、影
影绰绰地),这已经相当清楚地表明:我的思考,早在不知多久之前,就不再有资格称作“冥想”了。
它落在一种分裂之中:一边是写的人(却提防着不肯整个人扑进去\fnote{980}{(*)}{可是,就在上一段里,我刚写过(心里并无半点保留):托付去打字的那些稿子,我是“全身心投入”的。可见同样的话(或差不多的话……),随语境不同,可以有不同含义,或指示不同的细微差别。}),一边是活着、感受着的人(却也不停下来
“摆弄”他所经历的、浸透其中的意义)。在那里我感到,到达那个“句号”(当然,也不是说就把尚
待看见、尚待说出的这一点点草草了结……)然后返回我自身,已是刻不容缓……

除了笔记的活儿,这几天还有另一件事在给我调剂。那是一场数学思考,不以我意志为转移地重新
开了头。几天来我已明白:一套 Motif 理论形状的构造,带着二十年前我所望见的那全部气魄,离
“地平线”其实并不像我想的那么远。甚至有可能,一套“完全成熟”的理论——连六种运算的完整
形式主义(外加双对偶)在内——对肯把自己整个儿扑上去的人(不以掘墓的心性糟蹋自己的创造
之力)来说,只是几年功夫的问题。我还看出,要显式描述“那个”Motif 范畴\dmark{}\origpage{1201}——比方说,取绝对基底 $\mathbb{Z}$ 上有限型的概形(人们总应当能归约到这种情形)——有\textbf{两把}“钥匙”\fnote{981}{(*)}{其实还有第三把“钥匙”,这里不提,因为在我看来(不管对错)它没那么棘手。这就是“de~Rham--Mebkhout系数”的恰切定义(先不管滤套与 $F$-结构),比方说,取绝对基底 $\mathbb{Z}$ 上光滑概形。这一定义应当同时提供特征 $p>0$ 的一般晶体系数“那一个”恰切定义的钥匙——而我那些亲爱的前学生们(这一回首当其冲是Berthelot),始终没能、或者始终不愿把它提炼出来。
1983年6月(快两年了),Mebkhout向我讲解他围绕“好上帝定理”的“哲学”时,我的印象是:他对复数域 $\mathbb{C}$ 上光滑概形的(可构造的)离散系数范畴给出的“纯代数”(“德拉姆型”)描述,恰是Deligne在(别处已提过的)他1969/70年间在法国高等科学研究所所给讨论班(若无记错)中采用、却从未发表的那条进路的对偶;那条进路,靠的是带联络的 pro-模。我猜,从一个观点到另一个观点的过渡,经由对偶化函子 $\mathrm{R}\mathcal{H}om(-,\mathcal{O}_X)$ 完成——相对于所论概形的结构层——它把有限型的 $\mathcal{D}_X$-模(可以看作配备了一个可积联络的“ind-凝聚的 $\mathcal{O}_X$-模”)变成“pro-凝聚”的模(同样配备可积联络)。Mebkhout观点的长处,在于它为“好系数”给出一个简单而深刻的代数表达($M$-凝聚、完整性、正则性),这正是Deligne所缺的。Deligne观点的长处,在于它给出的是(与所要表达的超越性质的系数之间的)一个等价,而非反等价,并且更宜于表达所论系数范畴的乘法结构(张量积)。我猜,实践中人们会常常愿意在两张互为对偶的图式上同时工作。Deligne的解读,经由可积联络的模(或 pro-模)这一直觉,在我看来更贴近直接的几何直观。这一点的表现,尤其是(若基域为 $\mathbb{C}$):一个可构造的 $\mathbb{C}$-向量空间层,对应的是唯一一个带联络的 pro-模,而不是这种 pro-模的一个复形。正因如此(读者想必猜得到,这让我多么遗憾……),我预料最终最宜于用来发展六种运算的形式主义、并充当 Motif 范畴构造中第三种配料——钥匙的,恰恰会是他的观点——尽管正是他,毫无惋惜地把它埋葬了,仿佛借此把那位被拒不认的宗师所遗留的系数问题也一并埋掉……(5月9日)关于这一点,另见子注“……以及羁绊”(n\textsuperscript{o}~171(viii)),以及“五张照片”(n\textsuperscript{o}~171(ix))。}:一把,是“神秘函子”的理论,要有足够的普遍性与灵活性,好过渡到相应的三角范畴,把de~Rham--Mebkhout系数与寻常的 $p$-进系数(特征零)连起来;另一把,是在一个\textbf{有限}域 $k$ 上显式构造 Motif 范畴的问题(最好经由一种“纯代数”的构造,不援引 $k$ 上的代数几何),此外还有从 $k$ 上分离有限型概形(先从射影光滑概形开始)到这一范畴的“Motif 上同调”函子。后者我在当年已经构造到了差一个等价的程度,用的是Weil猜想与 Tate 猜想作启发\fnote{982}{(*)}{如果我记得不错,我当时只限于描述半单的 Motif 范畴。按同一原理,这一构造稍作变通,就能给出一个颇可信的候选者,充当不必半单的 Motif 范畴。我在这里说“Motif”,实际上说的是“iso-Motif”,即差一个同源不计的 Motif。但若对每个素数 $\ell$ 使用“$\ell$-进实现”函子,便能从那里重建不经同源的 Motif 范畴(其中诸 Hom 将是 $\mathbb{Z}$ 上、而非 $\mathbb{Q}$ 上有限型的模)。
说到我的构造“以启发的方式”使用 Tate 猜想,不可按字面理解。倘若确实存在(比如说,在有限域上)射影光滑概形上的上同调类,它们是“Motif 的”(这含义正待提炼)而非“代数的”(即并非来自代数闭链),那么 Tate 猜想(同理还有 Hodge 猜想,这回是在 $\mathbb{C}$ 上)就应当改述,把“代数类”换成“Motif 类”。假定人们果真能(如我下文所提示)在有限域 $k$ 上射影光滑概形的范畴上,定义出典范的(并且在适当意义上应当是“普遍的”)上同调函子,取值于已经构造好的范畴(所谓“$k$ 上半单 Motif 范畴”),那么这就顺理成章地给出一种形状完整的定义,界定人们将要称作“Motif 类”的上同调类:即 $\mathrm{Hom}(T^i,H^{\bullet}_{mot}(X))$(2i 维情形)的元素,其中 $T$ 是 Tate 对象,而 $H^{\bullet}_{mot}$ 是所拟的假设函子。这正是为什么,这个函子的构造,在我看来眼下是形塑 Motif 理论(不再是六十年代那种假设性的,而是形状完整的)的\textbf{头号}关键问题。}。

我对这一\dmark{}\origpage{1202}构造的正确性没有半点怀疑。尚待完成的工作——无疑明显更精细——在于:用给定的
有限域 $k$ 把这条范畴“钉死”,尤其是定义“Motif 上同调”函子,哪怕先只在 $k$ 上阿贝尔概形
的范畴上定义(这应当就足以把所求的范畴“钉死”……)。这第二个问题,在我看来,比起神秘函
子,较少技术性,而更直接是“几何的”。而且,它向我显形为\textbf{那把}能解开标准猜想\fnote{983}{(*)}{“标准猜想”一词在这里不可按字面取义,正如前一条脚注里的“Tate 猜想”。毋宁说,这些猜想的表述里,应当把所论的闭链类扩大(起初只限于代数闭链)。在“重新调整过”的标准猜想的“最终”表述里(即便它们按原样已然成立),“代数的”上同调类还要再换成“Motif 的”类。关于标准猜想,我将在“Motif,吾爱”(《数学反思》第3卷)里更详细地回头来谈。}的钥匙,并由此也解开特征 $p>0$ 的上同调理论里那些恼人已极的整性问题。凡此种种,使得这个问题对我有一股强大的吸引力!今晚已是第三个晚上,我重新拿起笔记,打算快速巡视一番那些在我看来最灼热的主题——十五年前我告别数学舞台时,被我的学生们、也被所有人抛下的那些主题\fnote{984}{(*)}{这些主题的一个极简略的初步“巡视”,见去年那道注“我的孤儿们”(n\textsuperscript{o}~45)。}。这一回,我总算要成功了!

\asterisms

\pmark\origpage{1203}\textbf{工地一:Topos}。我在这里提到它,主要是备而勿忘;在注“我的孤儿们”(n\textsuperscript{o}~46)
里,我已经相当周详地说过它。由于我的某些前学生——Deligne首当其冲——乐于糟蹋这条
关键性的统一性观念,自我离去之后,它被打入冷宫,只落得在边缘上苟活。如我在所引注里提醒
的,topos 与形形色色的“多重叠”其实在几何里俯拾皆是——不过,当然,尽可以看不见它们,正如
人们几千年里看不见对称群、看不见集合、看不见数零一样。

一套贴着拓扑直觉、柔软而精微的关于 topos 的语言,已在 SGA~4 的前两卷里被精心发展完备
(就是那块著名的“废话矿渣”——光辉的“SGA~4½”一书第一讲的引言里,Deligne如此称呼它)。这是
Leray所创“层”的观念周遭的语言与直觉的自然结晶;层直觉与层工具发展中的这第二步(或这“第
二口气”),在我看来,分量不下于第一步(其临时表达即戈德芒那本名著)。正是这一视野,让
$\ell$-进与晶体的上同调工具得以问世——然后,恰恰是那班佯装把工具据为己有的人,把它无限期
地埋葬了。

SGA~4 中关于 topos 的种种展开,并不自命完备与终极,但我认为,对 topos 式视野的大多数
直接的几何应用而言,它们绰绰有余。正如一般拓扑或寻常的层论,“topos 式的一般拓扑”本身,
在我看来并不提出什么真正深刻的问题。它是一套精心打造的语言,服务于对形状的拓扑-几何直
觉的某种扩展;这种扩展,是事物本身口授给我们的。这一视野长期蒙受的污名,以及落在它头
上的\textbf{嘲弄},在我眼里属于七十年代、八十年代数学世界的几大凋败之列。

这里说的,不是一间需要有人来救活的“破败工地”,而是一所全然竣工、设备齐全的房子:那些
住过它、又本被唤来在其中工作与生活的人,选择离它而去,还顺手撵走了盖房子的工匠。房子宽
敞、健全,一切各安其位,一如工匠出门做别的营生的那一天。它若还缺什么,缺的\dmark{}\origpage{1204}不是他这双手
的劳作,也不是任何人的劳作。也许,工匠本人对这些东西的敬意——那双手怀着爱做
出来的、他知道是美的东西——会让这股暴戾与轻蔑的气息消散,让这间本为迎接而造的房子重新
学会迎接。

\textbf{工地二:上同调的语言}。这首先是指导出范畴的语言;再者,是我为非交换上同调引入
的诸观点——两者都在五十年代下半叶。

头一条路线,本该由Verdier那篇著名的“论文”来承载;而Verdier本人对自己论文的埋葬
\fnote{985}{(*)}{关于这一点,见注“赊账的学位论文与全险”(n\textsuperscript{o}~81),以及“随心所欲的荣耀——或暧昧”(n\textsuperscript{o}~170(ii))。},同时也就是代数同调中导出范畴这一观点的埋葬。这一观点,曾在六十年代代数几何中上同调主题的繁盛期——尤其在对偶形式主义、不动点公式(Lefschetz--Verdier型)的发展中——扮演过关键角色。实践的需要,早已显出Verdier六十年代初发展的三角范畴框架之不足;这一框架,至今没有像本该的那样得到更新。

“非交换”一路,有Giraud的论文打下的良好基础工作,但它止步于一种 1-叠的形式主义,只能把
上同调对象直接几何化到二维为止。以 $n$-叠与 $n$-gerbe 的语言发展非交换上同调形式主义——众
多例子急切地提示了这一点——一路撞上严重的概念困难。鉴于在某个上流社会里,基础
问题跌入了人人不屑——说得更直白些,人人轻蔑——的境地,这些困难从来无人问津,直到两年
多以前,我亲自去啃\fnote{986}{(*)}{这指的是1983年2月我致Daniel~Quillen的信中的思考;在那封信里,我发现可以“一跃而过”那道豁开的“炼狱”——随着 $n$ 增大,($n$ 不严格的)$n$-范畴(或按我如今的说法,$n$-叠)的形状化描述中,似乎总要冒出来的、越收越紧的相容性关系。$n=2$ 的情形已然不是一桩轻松活,而 $n=3$ 的全部相容性,我想还没人有勇气把它们一一写出。这封信,如我下文还要提到的,成了长途跋涉《追寻栈》的发令枪;次月,乘着已经启动的思考的余势,我动身了。
这封信没有被视为值得一读,也没有得到回复。一年多之后,我才收到当事人的一点评语;对此我在“一段过去的重量”一节(n\textsuperscript{o}~50)里有交代。(参见打字稿第136页第二段。)}。

\pmark\origpage{1205}如今我看见两条路线在一门新学科里汇合——我在别处\fnote{987}{(*)}{见注“阴,这仆人——或慷慨”的子注 n\textsuperscript{o}~136$_{1}$(尤其打字稿第638页)。}提议把它叫作\textbf{拓扑代数}(l'algèbre~topologique):它是传统同调代数(当然是导出范畴风格)、同伦代数、(尚在娘胎里的)$n$-范畴、$n$-广群及相应的叠与 gerbe 的形式主义,以及 topos 的视野——它如今提供了已知最广的“纯代数”框架,供安放拓扑直觉——这几家的合成。这样一种合成的初始观念,六十年代已经齐备,其中包括\textbf{dérivateur}的观念——它注定要取三角范畴这个不够用的观念而代之,并且同样适用于“非加性”的语境。同伦代数中的某些重要进展,比如Bousfield与Kan七十年代初发展的同伦极限与同伦余极限诸观念——他们并不知道我的想法(在我亲爱的学生们那里,这些想法被当作“Grothendieck式的老糊涂涂鸦”打发了)——正落在这些想法的延长线上。

两年前,我以致Quillen的那封信\fnote{988}{(*)}{关于这封信,尤其见“一段过去的重量”一节(n\textsuperscript{o}~50,打字稿第136页第二段)。}起手,为我看见该做的工作起草了一份总施工图。那封信是《追寻栈》一书的发令枪;其第一卷(“模型的历史”)业已基本完成,无疑将作为《数学反思》第4卷出版。我预计还需要一两卷、一两年功夫,才能完成这一前期勘探——勘探的是一片极其丰饶的矿脉;二十年后,我似乎依然是唯一看见它的人。所以,这确实是一处撂荒了十五年的工地——但近一年来,它已在我手中复活。《概要》的写作,继之《收获与播种》的写作,曾把它打断;然而一旦《收获与播种》和那些(篇幅都有限的)文本——它们将与《收获与播种》最后一部分一起构成《数学反思》第3卷——写完,我就打算回去,把它做到底。

\pmark\origpage{1206}\textbf{工地三:六种运算,双对偶}。这指的是我引入Poincaré式或Serre式对偶形式主义(取离散或连
续系数)的那个观点。我当年起的“六种运算”这个名字,被我那班上同调学生细心地铲除
干净了。他们不过这里那里挑几样合用的来使,却把这几样合成一体所成的那种结构(连同双对偶
形式主义)扔进了盈亏账,尤其是,把这一观点所提供的那条无可替代的主线扔了(这条主线,对提
炼好的“系数范畴”至关紧要,见下文)。二十多年来,这套形式主义早已存在并经受了检验;可在“局
中”的人里,没有谁(除非是在注定保密、我无从知晓的文字里)费过神,去提炼那“代数总目”——
凡有这种“万能”对偶、以六种运算形式主义表达的地方,通用的那一套代数公式\fnote{989}{(*)}{(5月9日)在 SGA~5 最早的一次报告里,我曾煞费苦心,把这套公式说得详详细细;它将像发动机的神经一样,驱动整个即将展开的讨论班。这一讲是 SGA~5 里最关键的一讲,却在“屠杀版”里不见了踪影。全卷之中,竟无一字暗示它曾存在!见注“先祖”(n\textsuperscript{o}~171(i))里打字稿第942页脚注(*)。}。

可见,这里严格说不是一处“撂荒的工地”(要补的形式化工作微不足道),而是一个被系统性回避
的丰饶观点(正如 topos 的遭遇)。这场抛弃,无疑在很大程度上,造成了我在代数簇上同调主题上
看到的可悲停滞(少数例外\fnote{990}{(*)}{这些“少数例外”,主要是(1981年以前)Deligne那两项重要工作 Weil~I、II,以及晶体上同调、Barsotti--Tate 群的 Dieudonné 理论在特征 $p>0$ 的一般基底上的若干零星结果(后者是我在1969年前后开创的)。如我在别处指出的,在“好上帝-Mebkhout定理”的余波里,有过一阵回潮(定理与作者照旧一个被无视……),尤其是Mebkhout层理论的发展(它们被错误地叫作“反常层”,顶了本该属于别人的名分……),由Deligne等人展开。}除外)——尤其与我1955到1970年间赋予这一主题的强劲起飞相比。

如我已在《引论》(I~8,“一个秘密的终结”)里预告的,继《收获与播种》之后\fnote{991}{(*)}{提醒一下:这指的是《数学反思》第3卷,原则上它也收《收获与播种》的最后一部分。},我打算收入一篇短
文,勾勒“六种运算”形式主义的要点。多亏我那些\dmark{}\origpage{1207}学生的照料,今天它的存在本身已无人知晓——
除了曾直接参与 SGA~4(1963/64)与 SGA~5(1965/66)两个讨论班的人,而他们显然忘了\fnote{992}{(*)}{巧了,这也正是那篇以“核心”自居、名为(何其讽刺!)“SGA~4½”的文本劝人千万别去读的两个讨论班……(5月29日)六种运算视野的分量,见注“无用的细节……”(n\textsuperscript{o}~170(v)),部分(b)(“一些什么也干不了的机器……”)。}。这样,凡我力所能及的,我都做了:为一件效力完美的工具、一个在
上同调主题里不断把我们径直引向关键问题的丰饶观点,恢复名誉(只要还找得到守着好工具的工
匠)。

我刚巡视过的这三处“撂荒的工地”(或房子,或工具……),牵涉到的,更是一种\textbf{共同的代数
语言}——用来表达最形形色色的几何情形——而不是某个个别的几何情形,比如代数簇的上同调。
第二处工地,我称之为“拓扑代数”的,有时会让我跟一些无疑深刻的问题擦肩(比如球面同伦群
的某些问题),那是偶然,不是本意。我在那里的主要动机,一如既往,是发展出普遍性、灵活性
都足够的代数工具,来服务于那门尚在襁褓中的\textbf{算术几何}:从Weil猜想那颗胚胎起,我花了
一生中十五年的好时光孕育它、接生它、喂养它。几何实质的精华正在这门几何里;这些年里,它
才是我同数学的那场恋爱真正的核心——今天依然如此。接下来还要巡视的三处“最灼热”的主题,
说的正是这精华。

\textbf{工地四:“系数问题”}。这个问题,在Weil猜想的表述本身里,已然萌动\fnote{993}{(*)}{关于这一点,见注“操演”(n\textsuperscript{o}~169)的开头,那里我评说了Weil猜想的初始问题域。(5月29日)这一开头已独立成注“‘Weil猜想’的语境”(n\textsuperscript{o}~169~(i))。}。整个六十年代,它是我在上同调里的兴趣中心。对当时已约略望见的几种主要系数类
型,它已被提得清清楚楚,带着全部必要的普遍性与精确性\fnote{994}{(*)}{似乎并没有出现什么新式的“系数类型”,超出我六十年代下半叶所预见的那些。}。这一问题域——对理解代数簇上同调显然至关重
要\dmark{}\origpage{1208}——我早在第一次回顾自己的作品、那道以敬意命名的注“我的孤儿们”(n\textsuperscript{o}~46)里就说
过,又在注“墓上的旋律——或自负”(n\textsuperscript{o}~167)里回到它。两条根本的主线:一条,是刚才说
过的六种运算与双对偶的形式主义;另一条,是需要在一个或普或泛的基底概形上,给出基底域上
已知的诸种“系数”的合宜推广——它们(哪怕只是默而不宣)进入了对该域上射影光滑概形范畴已
知的那些上同调函子的描述:$\ell$-进上同调、晶体上同调、德拉姆上同调,乃至($k=\mathbb{C}$ 时)
Betti 上同调或 Hodge 上同调。

我不认为这样说有什么过分:这一问题域已经以萌芽的形式\fnote{995}{(*)}{作这一说明,我丝毫无意贬低Deligne与Mebkhout上述贡献的原创性与重要性,正如我不会因为说算术几何“早已在Weil猜想中萌芽”,就去贬低我自己的贡献在这门几何的诞生与最初起飞中的原创性与重要性。}包含着那套至今仍在等待问世的“完全成熟”的“Hodge--Deligne理论”,也包含着同样仍在等待的“德拉姆--Mebkhout系数理论”\fnote{996}{(*)}{大体可以说,先是Deligne的贡献(1969年前后),后是Mebkhout的贡献(1975年以后),分别朝着两个很不一样的方向,回应了定义合宜的“德拉姆系数”的问题——这种系数将能把光滑概形的寻常德拉姆上同调嵌进一个六种运算的形式主义。Deligne只在概形 $\mathrm{Spec}(\mathbb{C})$ 上定义了一类“好”系数,并且对 $\mathbb{C}$ 上分离有限型概形的结构态射 $X\to\mathrm{Spec}(\mathbb{C})$、对 $X$ 上的常值系数(可惜!)定义了函子 $Rf_!$ 与 $Rf_*$。Mebkhout定义了一类“好”系数,原则上对特征零域 $K$ 上任意分离有限型的 $X$ 都成立——但他没有进而对这类概形之间的态射 $f:X\to Y$ 定义 $Rf_!$ 与 $Rf_*$,也没有对 $Rf_!$ 与 $Lf^!$ 发展对偶定理(除了 $Y=\mathrm{Spec}(K)$ 的情形——而且即便那时,也只在复解析簇的超越语境里,那无疑明显更难)。Mebkhout迄今发展的理论(应当说,是在一种令人丧气到极点的氛围里发展的)还有一个局限:它目前只对光滑的 $X$ 做成(我猜,是因为没有系统地使用晶体观点;晶体观点为微分算子层提供了一个满意的替代物,而后者在光滑情形是那样便利)。};而这两套理论之所以至今仍在襁褓之中,没有像 $\ell$-进系数理论(素数 $\ell$ 与特征互素)那样,一两年间便臻于成年,原因是同一个:我那班上同调学生,Deligne首当其冲,忙不迭地把宗师遗留的问题域连同宗师本人一道埋葬。

\pmark\origpage{1209}Deligne这边(他享尽了科学界宠儿的一切便利),与Mebkhout那边(他身陷的孤立再彻底不过——
而这孤立,恰恰是最该接纳他的那些人压在他身上的),所迈出的步子固然零碎,却仍不失为宝贵的
引线,让人得以提炼出某些关键的系数范畴。写那道已引的注“墓畔的旋律”时,这些重要的贡献就
在我心里。此后,我又更深地浸入了那套六十年代已然成形的“系数与 Motif 的瑜伽”,如今有了一幅
更精确、也更完整的图景。因此,我打算在《数学反思》第3卷里,接着六种运算形式主义的概要,回头
再谈系数问题(同时也谈 Motif 问题)。

现在只需说:我分明看见,在一个或普或泛的基底概形 $X$ 上,有本质上的三种“基本”系数类型\fnote{997}{(*)}{我说“基本”的系数类型,是想借这称呼提示:我眼下能窥见的其余各种重要的系数类型,都应当能借这几种来描述——或以合宜的方式把它们“组合”起来,或对它们添加合宜的结构增益,或两者并举。在为de~Rham--Mebkhout系数所设想的诸种结构增益之中,有(除了似乎“内在于”所论系数范畴的“权滤过”之外)一种“德拉姆滤套”,它在 Motif 式的应用中扮演头等角色。有可能,这层附加结构(从六种运算形式主义的眼光看),唯有当它与一种贝蒂式的“离散”结构相结合时,才有意义;这应当使我们得以表述这一滤套所须满足的好性质。关于这些问题,我打算在“Motif,吾爱”(《数学反思》第3卷)里更周详地回头来谈。}:$\ell$-进系数($\ell$ 为\textbf{任意}素数);\textbf{de~Rham--Mebkhout系数}\fnote{998}{(*)}{提醒一句:对这种de~Rham--Mebkhout型系数,我如今看见互为对偶的两种变体:Mebkhout的那一种,与我迟疑着不愿称作“Deligne的”那一种——尽管后者是一个被他拒认的孩子!(5月29日)关于这个被拒认的孩子,见注“……以及羁绊”(n\textsuperscript{o}~171(viii))。关于德拉姆系数的细节,见注“五张照片(晶体与D-模)”(n\textsuperscript{o}~171(ix))。}(主要对在基底概形 $S$ 上有限型的 $X$ 有意义,最\dmark{}\origpage{1210}重要的情形是 $S$ 为环 $\mathbb{Z}$、$\mathbb{Q}$或 $\mathbb{C}$ 的谱);最后,是 Betti 系数(对 $\mathbb{C}$ 上有限型的 $X$)。这三种之中,唯有第三种,在我看来眼下已经确定下来,不带任何假设成分。要定义第一种(哪怕只对绝对基底 $\mathbb{Z}$ 上有限型的 $X$),或描述它与第二种
的关系,关键配料在我看来是那套“神秘函子”的理论——我六十年代末就公设过它,这个问题,看
来也同其余的一切一道沉了下去……到时候,我自当在合适的地方更周详地回到它。

\textbf{工地五:Motif}。关于我朋友Pierre~Deligne对 Motif 的埋葬——蒙教团全体祝福——我已
经说得相当周详,此处不必再度铺陈。这里我倒想强调一件刚刚才向我显形的新事,它本该在十五
年、二十年前就显形的。一个月前,在一个或普或泛的基底概形(比方说 $\mathbb{Z}$ 上有限型,或
只取一个代数闭域的谱……)上“形状完整地”构造 Motif 范畴,在我看来还是一件分明“在地平
线上”的事,淹没在遥远将来的迷雾里。这种心境,无疑是那些已然遥远的日子的顽固遗产:当年
Motif 的思考起步时,基地再假设不过,那时连 $\ell$-进上同调的形式主义都还没有。对我而
言还有一条“从轻发落的情节”:1958到1970年间,那些近在眼前的奠基性写作任务,把我的精力吸
干了;我的 Motif 思考(以及关于其他在当时那些紧迫任务面前显得“奢侈”的主题的思考),一直
被压缩到一份可怜的口粮,那份口粮,我几乎怀着逃学孩子的负咎才拨给自己!不管怎样,我留在
脑子里的印象是:系数问题,是马上就该做的成熟的事(不过得靠别人,因为我已经忙在别处
……);而 Motif,眼下顶多配上一本“数学虚构”的书,要是哪天我抽得出闲,去写一写。说实在的,
倘若我真的一头扎进去写它,而不是耗在那些后来世上没有一个人肯接手的任务上——尽管人人都
乐于享用我做的东西……——事情的步调,想必很快就会改观……

\pmark\origpage{1211}话虽如此,我终于看清了这件事——它本身,一旦摆在眼前,原是明摆着的:只要你肯花功夫去描
述足够“细”的系数,也就是说,照顾到一个 Motif 相关联的全部已知结构,你最终描述的,就是\textbf{Motif
本身}。或者说得对些:你最终描述的是一个范畴,它将把(三角的)Motif 范畴作为一
个\textbf{全子范畴}包含在内(这已经很可观了)——正如复数域上的 Motif 范畴(若承认 Hodge 猜
想一个相当强的版本)作为 Hodge--Deligne结构范畴的全子范畴而出现。至于要确切刻画——用与手
头系数直接相配的“代数”语言——这个全子范畴\textbf{究竟}是什么,也就是说,\textbf{哪些}系数恰
好“是 Motif”,那就落进了恐怕要棘手得多的问题。它们关乎一个 Motif 相关联的种种几何-算术
结构之间的\textbf{相容性}(这些相容性,我想我在前面引过的注“墓畔的旋律”里已经影射过)。解决
这些问题的路数(它们在我看来,与“Motif 理论”的实际构造毫不相干),兴许真的“要等一百年”。
不管怎样,经验一次又一次让我们看到:这类关于一个问题多么“无从下手”的预言,没有多大意
思,顶多是在勇气本没有挂稳的地方,劝人泄气……

(4月1日)关于“\textbf{Motif 的Galois群}(或\textbf{基本群})”的形式主义,再添几句评注。这个
概念(1964年,在我有幸结识我未来的前学生Pierre~Deligne之前,我已经提炼出来并开始发展)带
来一些极其精确、极其精微的直觉和形式主义。它的存在及其要点,不依赖于人们为一个域上
的 Motif(或任意概形上的“光滑”Motif)的观念所选定的特定构造——只要那个构造满足几条合
理的条件。我曾把这样一项任务交给Neantro~Saavedra:把我大约1964年提炼的辞典整理成可
发表的形状,放进尽可能一般的语境——辞典的一边,是我当年称作“刚张量范畴”的那些范畴里
的几何($k$-线性范畴,带满足若干合宜条件的“张量积”运算,$k$ 此处是一个\textbf{域});另一
边,是 $k$ 上 pro-代数群的线性表示理论(更确切、\dmark{}\origpage{1212}更一般地说,是 $k$ 上的“pro-代数 gerbe”的表
示理论)。他把这项任务在他的论文里做到底,1972年由“数学讲义”出版(LN~265)
\fnote{999}{(*)}{(5月10日)这几行写下之后,我有机会读到了那本书——作者可没觉得有必要寄我一本。我看清了:书中,Saavedra俨然是那里所展现的新哲学的光辉发明人;他一丝不苟地遵照我交给他的那些笔记,而几乎不提我的名字(无论是书中引入的观念与关键结果,还是一些早已为人所知的观念,比如晶体、分层模或 Motif)。他为那个主要观念重新起的“塔纳卡范畴”这个名字,是一场天才绝伦的骗局——那名字,断不会比他所自命的这套理论,更是他自己的发明。何况这桩“父亲身份”本来就是暂时的:我的朋友皮埃尔已经在丛书问世十年之后,把必要的事情料理妥当,好叫它(如众人所期待的那样)回到那位早已内定的父亲名下。这场对一具尸身的辉煌手术的细节(在“SGA~4½--SGA~5”那场以同样不可模仿的手法进行的手术之前,它是第一场、也是唯一一场这种规模的),见那组注“第六根钉子(钉进棺材)”(n\textsuperscript{o}s~176$_{1}$至176$_{7}$)。}。这本辞典我当年还推进得更远(尤其涉及把滤套结构或分次结构等,翻译到某些纤维函子上,或涉及与塔纳卡范畴相关联的一个“极化”观念),这些在Saavedra的论文里都没有做\fnote{1000}{(*)}{(5月10日)这原来是一句错误的揣测。它出自我的一个信念:Saavedra断无能力“收尾”我给他指的那个计划——单是“pro-代数 gerbe 的线性表示”这一观点的驾驭,很长时间里似乎就超出了他;他的数学行囊又极其单薄。以Saavedra那点并不出众的本钱,从我离去(那时他对上同调、对代数群的构造一无所知)到书出版,不足两年,他竟有机会消化(而且消化得尽善尽美——书的成色即是明证)书中随手把玩的那一大堆四面八方的观念——这在我看来是不可设想的。关于这一点,见那组已引注“第六根钉子(钉进棺材)”里的注“Monsieur~Verdoux——或献殷勤的骑士”。},或者在“难忘文集”LN~900里都没有(在那里,Saavedra
的论文被重做了一遍,而 Motif Galois群的观念居于问题的中心——只是我的名字,在这个话题
上,同在所有其他关于 Motif 的话题上一样,一次也没有被提起)。

我还想指出:前面谈到过的那第一步——确定有限域上 Motif 范畴(差一个等价)——原本是对
该有限域的 Motif\dmark{}\origpage{1213} Galois群的确定\fnote{1001}{(*)}{(5月10日)我注意到,这一确定,也见于Saavedra那本取之不尽的书里(对我这卑不足道的人只字不提,这还用说么)。它用的是整体类域论的上同调理论(确定群 $H^2(\mathbb{Q},T)$,其中 $T$ 是 $\mathbb{Q}$ 上一个乘法型群)——所以,这也算我的前学生(其本钱俨然超乎常人)不到两年里就消化掉的东西之一……},该群应当是交换的(它被 Frobenius 元素拓扑地生成),实际
上是一个经某些 $\mathbb{Q}$ 上的 pro-代数环面扩张的 $\widehat{\mathbb{Z}}$(由 Frobenius 生成)\fnote{1002}{(*)}{这里说的是对\textbf{半单}Motif 分类的那一个 Motif Galois群。要得到一般的 Motif,还须用 $\mathbb{Q}$ 上的加法群 $\mathbb{G}_a$ 作它的乘积。}。第二步,是 $H^2(\mathbb{Q},T)$ 中那个元素的描述——按照Giraud的理论,它恰把纤维函
子的 $G$-gerbe 加以分类\fnote{1003}{(*)}{要点在于:这个类在所有 $\ell\neq p=\mathrm{car.}\,k$ 的位置上都变为零(多亏“$\ell$-进上同调”纤维函子的存在),而晶体纤维函子的存在,让我们对这一类在缺失的那个位置 $p$ 上的命运,有了足够的信息。}。

正如我在注“一个梦的回忆——或Motif的诞生……”(n\textsuperscript{o}~51)里所说,我是在寻找下述联系时撞
见 Motif Galois群的:对一个投射有限Galois群 $\mathrm{Gal}(\overline{K}/K)$,取变动的素数 $\ell$ 的各
个 $\ell$-进表示——例如取 $H^i(\overline{X},\mathbb{Q}_\ell)$ 得到的那些,其中 $X$ 是射影光滑概形,$i$ 为整
数(或它的某个合宜的子模)。Serre看的是Galois群在 $\mathrm{Aut}(V(\ell))$ 中的像——那是一个约化的
$\ell$-进李群——而这个像的结构(按李理论的意义)似乎与 $\ell$ 无关。正是在探寻这一现象(直到今
天,它本身仍属假设)的深层原因、把它与 Tate 猜想相联系的过程中,我发现了 Motif Galois群的
观念——紧接着“Motif”与“Motif 上同调”的观念之后。

如果说我曾把一件既简单又深刻的东西带到世上,如果说我的数学家生涯里有过创造之举,那正
是这条沟通几何与算术的关键观念的诞生。也正因为如此,去年那个难忘的4月19日,当看到这件东
西被人以那样一副气定神闲的做派据为己有——仿佛它是某人在一个技术段落的拐弯处,即兴哼出
的最末一句小玩意:瞧,简单得像剥白菜,只消把我们的那篇小文章里陈述塔纳卡范畴理论的命题
4.7.3 拿来用在这里就行——我给一股难以想象的\textbf{无耻}之气噎住了\fnote{1004}{(*)}{写这几行时,又逼着我想起一件十分相似的事:把一个 $\ell$-进层系数的 $L$-函数的定义引入,同样不提任何人,同样仿佛是信手拈来的最末一条老生常谈,由同一位光辉的作者即兴完成。关于这一点,见注“操演”的子注“……以及‘废话’”(n\textsuperscript{o}~169$_6$),打字稿第891页。}。七十年代里那些光辉的\dmark{}\origpage{1214}先例之后\fnote{1005}{(*)}{甚至早在六十年代——见“驱逐”一注(n\textsuperscript{o}~63)。},八十年代的数学,就是这样做的。

不过我扯远了,好吧——我本该是领人参观一处工地的,不是来抒情的。那就说明一下:同投射有限
基本群的情形一样,若 $X$ 是域 $k$ 上几何连通的概形,应当区分概形 $X$ 本身的 Motif 基本群,与
“\textbf{几何的}”Motif 基本群。两者并不重合,\textbf{哪怕} $k$ 代数闭——因为 $k$ 的 Motif 基本群
并不因此就是平凡的(它连通,仅此而已!)。所以要引入 $X$ 的“几何的”Motif 基本群;它的职责
(除了别的之外)是在与几何投射有限基本群 $\pi_1(\overline{X})$ 相关联的诸 $\ell$-进李群(作为其商)之间
建立联系。它的定义,是自然同态
$$\pi_1^{mot}(X)\to\pi_1^{mot}(\mathrm{Spec}(k))$$
的核(相对于 $X$ 上光滑 Motif 范畴上一个纤维函子的选取)。

我想要抵达的地方正在这里:这个核——不妨记作 $\pi_1^{mot}(X/k)$——应当是构造“$X$ 在 $k$ 上的
(几何的)Motif 同伦型”的第一步;这一同伦型,我前文已经顺带影射过\fnote{1006}{(*)}{见注“为一副模糊的骸骨作安魂曲”(n\textsuperscript{o}~165)。}。“这个同伦型”的形状完整的描述\fnote{1007}{(*)}{作为对象的类型,我预料它将是一个相对同伦型(按Illusie的意义),放在 $\mathrm{Spec}(K)$ 的 fpqc topos 的“扩张”topos(按Giraud的意义)里——该扩张对应着 $X$ 上光滑 Motif 范畴的纤维函子的(在这个 fpqc topos 上的)gerbe。这个同伦型在(刚才描述的)基底 topos 上的相对上同调,是拟凝聚的(甚至是“凝聚的”),并且可以等同于 $X$ 在 $K$ 上的 Motif 上同调。若用 $X$ 的一个复点(特征零的 $K$),取 Betti 纤维函子,则相应的纤维同伦型应当典范地同构于经超越途径与 $X\otimes_K\mathbb{C}$ 相关联的 $\mathbb{Q}$-同伦型(不计挠现象……)——至少当 $X\otimes_K\mathbb{C}$ 是 1-连通的时候。},它
的“上同调”应当不是别的,正是 $X$ 的 Motif 上同调;这属于“Motif”工地上一个有意思的概念性
工作,朝着一条分明不同的方向\dmark{}\origpage{1215}(而且很大程度上多半独立)——不同于那项中心任务:
Motif 范畴的实际构造,以及为它们建立六种运算的形式主义。

\textbf{工地六:标准猜想。}如我在前面一条脚注(打字稿第1202页脚注(*))里解释过的,这些猜
想可以在两种不同的意义上理解。首先是字面的意义,如我1967年在孟买讨论班上陈述它们的那样
\fnote{1008}{(*)}{Algebraic~Geometry,~Bombay~1968,~Oxford~University~Press~(1969)。}。照这个样子,在我看来,它们概括了代数闭链理论眼下最关键的问题——至少从这些闭链的所谓“同
调”等价的观点看是这样。

陈述这些猜想的时候,我的主要动机,其实并非朝向闭链本身,而是朝向它们(或许……)提供的
手段:在一域上建立一套半单 Motif 的理论,满足那些本该在十五年、二十年前就“人所共知”(至今
却仍讳莫如深……)的要求。我将在《数学反思》第3卷里指出这些猜想的种种减弱的变体——它们足
以建成这样一套理论(其中最弱的那个,几乎就是为了这个目的而必要且充分的)。如我在别处强
调过的:即便初始形态的猜想被证明在某个给定的域 $k$ 上成立(比如 $k$ 有限,甚至一切 $k$),这也
绝不意味着,应当称作“Motif 的”上同调类\fnote{1009}{(*)}{我想我可以提出一个合理的定义,界定射影光滑代数簇上的 Motif 上同调类——至少当基域特征零时如此。对一般情形,关键的(前文已经谈过的)是有限基域的情形。知道了后者之中 Motif 类的描述(且听下文分解),我想就能给出“那一个”Motif 类的恰切定义。可参照打字稿第1202页脚注(*)处的评注。}(人们可以指望这些类使 Hodge 猜想、Tate 猜想之类种种猜想
为真)就必然是代数的。倘若有一天发现存在非代数的 Motif 上同调类,那多半意味着:代数闭链在
Motif 理论中——也就是在代数簇上同调的算术-几何研究中——的分量,要比理论初创时我有理
由相信的更轻一些。话虽如此,我如今隐约望见的 Motif 理论的实际构造,先天地独立于关于
代数闭链的时下诸猜想(Hodge、Tate 或“标准”型)。

\pmark\origpage{1216}但这并不妨碍:一方面,标准猜想及其种种变体;另一方面,Hodge、Tate 猜想及其众多的变体
——这些猜想,尤其包含代数闭链的\textbf{存在性}命题(即上同调类的代数性),或者(在修改版
里)所谓“Motif 的”上同调类的存在性命题——彼此紧密相联,并且同主要“系数类型”的描述、乃
至同 Motif 范畴本身的描述,紧紧相连\fnote{1010}{(*)}{这并不与我刚才的断言相抵触——我说过,我隐约望见的 Motif 范畴(比方说一个域上的)构造,与所论的种种猜想“独立”(即“技术上”或“逻辑上”的独立)。我说的那些“紧密的联系”(比如,我在 Hodge、Tate 型猜想中看到的十二种变体,恰好提示了同样多的不同类型的上同调“系数”),属于启发而非技术——正如(所谓“猜想性的”)Lefschetz--Verdier公式与 Frobenius 对应的迹公式之间的联系一样。在后一情形,这个根本的启发式联系(并非逻辑依赖的联系),已在注“操演”的两条子注“\textbf{真正的}数学……”、“……以及‘废话’”(n\textsuperscript{o}s~169$_5$、169$_6$)里得到郑重的指明。}。

这里,又是一件近二十年来早该做的澄清、归置、通报的活儿,没有被做(尤其,没有公之于众),
做的人却宁愿到今天还在埋葬丰饶的观念(当它们还不曾发表),或者拆毁它们(当它们发表了),
好独占(眼前的)利益与(日后的)功劳,而不肯把这些迷人的问题通报出来、交到大家手里——
这些问题,对我们理解几何、拓扑与算术之间的纽带,至关紧要。我看出,这里欠缺的,绝不是能
力,甚至不是光辉的天分,而是一种朴素的诚实,以及在对一个派发声望与权力的“科学共同体”的
关系里,某种起码的\textbf{廉耻}——在那些并不因此觉得自己欠下丝毫义务、欠下丝毫“服务”式“回
报”的人身上。所以,虽然我同这个题目断了往来已十五年以上,虽然我已经说不上“局中”二字,
费力重新入水的却还是我:至少要在《数学反思》第3卷里,把我那些更年轻、更有天分的人漏下的,
尽力修补;到底,把他们没有慷慨去做的事做了。

到这里,我相信我已经把这些“工地”巡视了一遍——从建造那门“算术几何”的眼光看(这门几何
的基石,是我整个六十年代铺下的),它们如今(而且自我告别数学舞台那一刻起)正是“\dmark{}\origpage{1217}最灼热”
的。我绝无意说,我已经把那些实质问题——那些也许只有我看得见、而又放在我心
上的问题——粗粗走遍了。据我所知,它们仍停留在我离去时的原地,其中许多甚至还没来得及在
文献里得到显式的陈述。其中我要特别指出\textbf{离散Riemann--Roch猜想}(概形框架下的)\fnote{1011}{(*)}{这一猜想似乎第一次被显式陈述,是在那道名字意味深长的注“屠杀”的子注 n\textsuperscript{o}~87$_1$ 里——因为这条猜想,本就是 SGA~5 遭屠杀的诸事项之一,在Illusie的编订版里消失得连一个\textbf{名字}都没留下。}。同样,还有把\textbf{局部与几何整体类域论}推广为一条本质
上“几何”性质的\textbf{对偶}命题(经典“算术”命题作为推论给出)的工作。1976年致Larry~Breen的几封信里谈到了它,那几封信将作为附录收入《追寻栈》第一章(因此会出在
《数学反思》第4卷)。在这些命题里,主要的工作将在于对人们所处理的“系数”范畴的细心描述。其中
一个重要角色,由Serre发现的一种自对偶性扮演\fnote{1012}{(*)}{除了Serre这个漂亮的念头之外,影响我的,还有 Lang 在几何整体类域论中、Serre在局部类域论中引入的“几何”观点。}:特征 $p>0$ 的域 $k$ 上、幂幺代数群范畴里(差根同源不计)的一种自对偶(在我看来,除了当年听我讲过的屈指可数的几个人,至今无人知晓)。把这类命题推广到更高维的问题,(至少对我)是个彻底的谜(不过Milne在代数曲面的情形,或许有些眉目……)。

这些对偶问题,我记得可以上溯到五十年代末。那时我还顺手开创了一个与之相联的构造:与一个
域上有限型概形(先从这里开始……)相关联的、pro-代数群的“\textbf{雅可比}”(链)\textbf{复形}——
借与那些局部环相关联的合宜的“局部雅可比簇”,类比于我几年前在凝聚对偶里构造过的“残差”
或“对偶化”复形。所有这些对偶问题,在六十年代,都被推到了第二线,尤其是为了
\dmark{}\origpage{1218}发展平展与 $\ell$-进上同调的“废话”和 topos 的语言。我纲领里相对局部与整体雅可比的那一小段,约在1977年由Contou-Carrère完成(没有提我这卑不足道的人);而一见到Deligne与Raynaud对它摆出的脸色,他就赶忙收了起来\fnote{1013}{(*)}{见注“棺材3——或有点太相对的Jacobi簇”(n\textsuperscript{o}~95),及其子注 n\textsuperscript{o}~95$^{1}$,其中谈到Contou-Carrère在数学大世界里的几段倒霉遭遇。}。如今,要重新拾起、发展那些打着我太明显印记的观念(哪怕当事人力图把这印记
藏得最好),是需要一点勇气的。唯一死磕到底的,是Zoghman~Mebkhout;而他遭到的命运——在反常
层会议上登峰造极的那些“绝技”——足够清楚地显示了这里头的风险。

倘若我想要开列一张1955到1970年间我发现的漂亮问题的清单(它们当中的一些,当年我曾在周
围这里那里谈论过),我还得再花上几天;若想稍稍说得明白些,把来龙去脉都交代了,恐怕要几
个星期。这里不是做这件事的地方,我也怀疑我永远会去做。更何况,倘若我盼着有一天(谁知道
呢!)哪位年轻数学家肯在其中一个问题上扑进去,练练手,再让自己为人所知,那还是让他自己
重新发现的好——犯不上冒险,让人给贴上某种标签。

当心将来等着我们的那些反常层会议……

% ================================================================
%  18.7 Les fruits du soir (原书 p.908–916)
%  源文: _work/chunks/chunk14.md
% ================================================================
\section{(7) 暮熟的果实}

\subsection{(1) 尊重}   % 18.7.1

\notehead{179}{}

(4月2日)已经五个星期了(自2月26日那道“沉默”之注起——它揭开了集结在“四种运算”名目
之下的那组笔记),我一直在重新巡视有关\emph{葬礼}的种种“物质性”的、或者(多少)“技术性”
的事实。在“四种运算”里,我把自己限定在严格意义下的“诈骗”这一面——也就是越过了那道
“门槛”的一面;那道门槛,同名的注(n\textsuperscript{o}~172)里谈过,它分开\textbf{恶劣的心性}(表现为
“自动排斥”的反射,往往不顾一个数学家最起码的直觉)与公然的\textbf{蓄意刁难}、与铁板钉钉的剽
窃。在刚写完的“荒废的工地”里,我面对的主要是\emph{葬礼}的“第一层次”,门槛以下的那一面:
一场广阔视野\dmark{}\origpage{1219}与一些强有力执念观念的埋葬——当然,没有谁非得接着做它们不可,人人都有权
无视或遗忘它们——只是这样一来,也就“把自己埋葬了”,让自己的工作(至少让那部分直接蒙受
被拒视野的工作)或多或少归于不育。

在那里,我总算有了走完一遍的印象!就“工地巡视”(撂荒的那些)而言,它让我对我著作的\emph{葬
礼}有了一份原原本本的把握,同时又让我同十五年来失了音讯的一些主题,多少重新接上了头。这
尤其让我看清了:在接下来几卷《数学反思》里我打算白纸黑字写下的东西,孰急孰缓。我的宗旨,当
然不再是给一门临盆中的科学打细致的地基——这种事我做得够多了;倘若再没有旁人肯像我当年
那样,把自己交给这样一项任务,那就各人自便,大家都自便!我的宗旨毋宁是:提炼某些执念观
念——它们服务于一个1955到1970年间诞生的总体视野;而今天我重新找到它们(多亏我某些昔日学生的努力,又蒙全体默认)时,它们或被遗忘,或沦为笑柄,或被肆无忌惮地据为己有、
残害得失去大半元气。今天重拾它们,我终于松开了缰绳,放开心中那股求知的冲动——六十年代
里,为了没完没了的“服务”任务,我常常只给它一份可怜的口粮。那个时代一去不返了——然而我
知道,在我数学激情的这个新阶段里,服务的冲动,丝毫不减当年。我依旧要为那个由渴望知道的
头脑结成的理想“共同体”服务\fnote{1014}{(*)}{关于“数学共同体”,我第一次表明看法,是在《收获与播种》第一部分“‘数学共同体’:虚构与现实”一节(n\textsuperscript{o}~10)里。在这里我援引一个“由渴望知道的头脑结成的理想共同体”,看上去似乎我又退回到某种东西上去了——它的虚构性质,在所引那一节里曾显示得清清楚楚。但在《愚妄与更新》第八部分里,我有生以来第一次(迟做总比不做好……)确认了我自己的“求知历险”在数学层面上的一个集体维度。(关于这一点,见“孤独的历险”与“一段过去的重量”两节,n\textsuperscript{o}s~47、50,尤其打字稿第134、135页。)同样清楚的是:经历这场集体历险的那个“共同体”(或“集体”),与任何社会学意义上的实体——它体现于特定“时代”的特定“圈子”,带着特定的“心态”,或者(今天)带着特定的权力结构与阶级利益——完全是两回事。我所援引的这个“理想共同体”,“在空间与时间上都无边界”,对我而言丝毫不比社会学实体更不“真实”。它甚至更其根本,因为正是它(如我在同一句话的后文里所写)“继续给我的数学投入以一种比个人消遣、比自我扩张手段更深的意义”。它并不比我自己更“虚构”——我自觉是它的一分子,而且比当年更清醒。“虚构”不在于觉察到这样一个“共同体”的存在,而在于把它同一个我曾自认同属的圈子混为一谈。},它继
续给我的数学投入,以一种比个人消遣、比自我扩张的手段更深的意义。

\pmark\origpage{1220}在这些投入里,当然,“老板”也并没有比当年更缺席。面对那些在我眼里曾是数学世界里
“我的亲近者”的人的恶意与嘲弄,面对那些我爱过、也曾无条件信任的人对我最起码体面感的屡
屡伤害,我心中涌起一股按捺不住的冲动:在那些丢失了敬意的人面前,\textbf{以对这些活生生的、
茁壮而美丽的东西的敬意,来见证我对自己的敬意}——这些东西,是我亲手带到光天化日之下
的。也许,见证这份敬意的最好方式,就是在这段还归于我的宝贵岁月里,做这些东西的仆
人。因此,未来几年里,我打算在《数学反思》的后续各卷中发展的数学思考,将既是重新拾起一个\textbf{孩
童的游戏},又是一件\textbf{服务的礼物},又是一个\textbf{敬意的举动}。

在给\emph{葬礼}画上句号之前,我还想在“物质事实”之外,就这场思考教给我的东西,作一小结。
我先看它教给我的关于他人的事,最后再谈它教给我的关于我自己的事。

在整个思考过程中大白于天下的种种事实里,至今最触目的,仍然是七十年代、八十年代数学世界
里\textbf{风气与人心的败坏}。这种败坏,借成百上千桩“小节”表现出来,如同过去八九年间,一阵阵
回到我耳边的那些“小节”——它们却足以让人错愕,足以催生《收获与播种》第一部分的思考及
其根本追问:事情是怎么(又是何时)走到这一步的?而在这场我今天目睹的阴险而无情的败坏之
中,我的角色是什么,我的位置又在哪里?

\pmark\origpage{1221}这场败坏,在“SGA~4½--SGA~5”那样的行动里,或者在(更难以置信的)反常层会议那场
行动里,达到了顶点——其犬儒与轻蔑,远远超出了我所能想象的一切;哪怕就在我身不由己地发现
它们的前一天,也想象不到。

在这里,不适合回头去谈那些“小节”(其中不少,已在思考途中顺带指出),也不适合谈那些大行
动(它们自有小伎俩伺候)。显露在两者——小节与大骗局——之中的,是同一种精神。在可容忍与
下作之间有时不妨画一道“门槛”,这道门槛本身就相当脆弱、相当人为,是一根护栏;反正(看样
子)已经没有人在乎它了。经由这场把我本人牵连至深的\emph{葬礼},我得以也许比以往任何时候都
更近地看清这种精神——它当然既不是这场葬礼(为我这卑不足道的人发起)的专利,也不是数学
家世界的专利。对此我能说的只是:我不知道,在我们这个时代之外的任何时代,这种精神曾统治
过那个世界,或其他任何一门科学。这大概是众多征兆之一,表明一种文明——以及那无论如何还
在给它赋予意义的东西——分解到了终局。

\pmark 近来,我的思绪不止一次地流连于一个奇怪的巧合:十五年前我告别数学舞台,是受了科学
界某种腐败的冲击——那种腐败,我曾长期选择闭目不见(自以为离它远着呢)。而我是在那所我
本打算终老于此的机构里,猝然同它打了个照面\fnote{1015}{(*)}{关于这一点,见注“有益的挣脱”(n\textsuperscript{o}~42),以及“兄弟与夫妻——或双重署名”及其子注(n\textsuperscript{o}s~134、134$_{1}$)。}。在那里,事关科学家们与军事机器之间那几近普遍的、有利可图的
串通。军方对整个科学世界的这种阴险把持,同样是个新现象,只是(至少就我们今日所知的规模而
言)自第二次世界大战以来才出现。当然,那场“冲击”之所以搅乱了我预定的轨迹(我自己、人人
都预定的轨迹),以至于触发我无可挽回地离开那个我一直自认同属的世界(除了一项默而不宣的
保留……),是因为我心中有一种紧迫而急切的更新之需——这一点,我是事后回望才意识到的
\origpage{1222}。后来,我曾倾向于把触发那次不寻常离开的特殊诱因说得轻些。然而我也知道,那些惯性之力
有多么巨大(同时又是多么无形):它们要把我们永远拴在同一条“轨迹”上,敌视内在的更新——
这也让我掂量出,把我从一条划得那样牢固的轨迹上拔出来,需要何等猛烈的内心震撼。

我想说的是:触发我离开的那桩“特殊诱因”,肯定不是没有\textbf{意义}。至少在我离开后的头几个
月,多半甚至头一年里,这个意义强烈地在场。后来,在新印象的涌入之中、在这场头一波汹涌更新
的动势里,这个意义自然地退居幕后,终于从我眼前消失。可是,即使我不再感知我过去或当下的行
为及其果实的某个“意义”,这个意义并不因此消失。而我重返数学活动——它不可避免地带来同我
离弃的那个世界的更多接触——又冷不丁把我带回那段被遗忘的过去。因为这次“回归”(同当年那
次离开一样出人意料……)最初的收获之一,就是在我那个曾经属于自己的世界里,发现了另
一场腐败——一种据我所知它从未有过的腐败。若要给这件新东西起个名字,涌到我嘴边的是:\textbf{敬
意的丧失}。过去几年里,我不止一次痛切地感到它:当我看到“我昔日所爱的某人,不动声色地踩
压我如今所爱的另一个人——在那人身上,他本该认出我”。在对\emph{葬礼}的思考里,我又一次次重
新遇到它,而且语气更凶,这一回是冲着我亲手生下的某些东西,或者冲着某个胆敢从中汲取灵感的
继承者。在这些时刻,我真正领教了这种丢了敬意的精神的“气息”与“气味”。可我也明白,这种精
神“并不只在我家四邻吹拂”——尽管,恰恰是它吹到我身上、吹到我所爱的人身上,我才真正“认
识”它——正如人们只有亲口尝过,才知道苦果的味道。今天,这种精神已经成了时代精神……

我看得很清楚:这两场腐败——触发我离开的那场,与我“归来”时等着我的那场——并非没有关
联。若要用言词勾勒这层关联的隐隐感受,我要说:\dmark{}\origpage{1223}在科学家们面对军方的金钱(只提这一面)及
其便利时的那种顺手就沾的态度里,我察觉到一种对自己的不敬——个人层面与集体层面皆然
\fnote{1016}{(*)}{在这里,我冒着得罪某些昔日朋友的风险——他们把这种“顺手就沾的态度”当成自己的,却当然不觉得自己亏欠了对自己
的敬意!再说,也绝非确定:别的时代的科学家们,倘若集体地面对同量级的“诱惑”,反应就会不同。常言道,机会造就窃贼!}。而正是在对自己的敬意的丧失里,我认出了对他人的敬意、对出自
他人之手或造物之手的作品的敬意,所以丧失的根由。

我并不声称“理解”了这两场腐败中的任何一场。一方面,有“时代精神”——它的独特动势,在我
看来几乎完全脱出个人行动所能及。这股集体动势,对我始终是个彻底的谜,我从未想过要去探测
它。另一方面,是每一个具体的存在——秉着他的知觉与创造之能,驮着他那全部特殊条件作用的重
负——回应这时代精神的方式;他(自觉或不自觉地)把这种回应,做成了他个人历险的关键要素之
一。

在整个思考过程中,我花了很久,试图去把握\emph{葬礼}的两大主角各自的选择,以及这些选择背
后的力:死者,与葬礼主祭\fnote{1017}{(**)}{(6月22日)到了“最后一刻”,在“家庭相册”一注(n\textsuperscript{o}~173)的 c.(“众人之上的那一位——或首肯”)、d. 与 e. 三部分里,第三个“主角”终于向我显形。}。确定无疑的是,一路走来我学到了东
西;但绝不能说我完成了任务。我甚至可以说,就我的主角而言,我肯定\textbf{没有}完全成功。我收集
拼图的碎片,把它们拼了起来,我甚至相信碎片是对的、拼法大体不错——但对\textbf{整体}的认识,
始终缺着。它至今还是一堆碎片,而这些碎片,对我而言仍是\textbf{陌生的}——陌生于我这个人、我
的经历,因此也就是不可理解的。已做的功课,无疑会有助于我在别的场合多少认路,在该留神的地
方留神(年纪越大,我越发觉要留神的地方真多……)。但这一切,离真正的理解还差得远。于是我
心里冒出一个问题:往这个方向下的功夫,最终是不是一场障眼戏——或者至少,那个\textbf{目标}(在
某种冲突情境里“理解他人”)是不是\dmark{}\origpage{1224}一场障眼戏(虽然走过的\textbf{路}倒是教益丰富……)。我对自己说:真正理解\textbf{那个人}身上的冲突(或者任何一个曾与我亲密相系、而我看
到同样矛盾炸开的人身上的冲突),大概同时也就是\textbf{理解冲突本身}。而且我深知,这样的理解,
不可能来自对他人的冥想(他人永远逃脱我直接的认知),只能来自对我自身的冥想。“阴阳之钥”
这场长考若要显出它的丰饶,靠的不是对他人其人的偶然瞥视,而恰恰是回头看我自己的生命、我
自己的经历,以及我当年对它们的理解。

\subsection{馈赠}   % 18.7.2. Le don

\notehead{180}{}

(4月3日)说到底,我并不觉得有兴致,用几行或几页文字,为\emph{葬礼}中我的主角向我显形的一
切,作一番回顾。照现在的情形,这恐怕不过是一场文体的练习,而对那最残缺不全的理解的更新,
并无多少助益。眼下,我倒是急于抵达这场关于\emph{葬礼}的思考的句号!

当然,我也知道,这个句号绝不等于\emph{葬礼}本身的终结。接下来的几个月,随着这些孤独的果
实般的笔记招来形形色色的回声,想必会有孤独的思考给不了的惊喜与教益。也说不定,来到我这里
的惊喜并不都是苦味的;也许不远的将来还给我备着某种欢乐——正因为它稀罕,才愈发珍贵;正
如仅仅去年(丰收之年!),我就得到过欢乐:收到旧日同事与朋友中的三位——我素来特别敬重或
亲近的——写来的满纸热忱的信\fnote{1018}{(*)}{这里说的是 D.~Mumford、I.M.~Gelfand与 J.~Murre的来信。}。

至于《收获与播种》对数学世界的“时代精神”会有什么全局效应——无论多么微末——几乎不必
说,我对这一点没有半点幻想。至多,这套笔记的出版,或许能叫某桩空前的横逆收场,叫某桩太过刺
眼的不正常得以纠正——而且,我这么说也许还是乐观了。也可能,死者本人的不期复现——他被当
作死了多少年、都馊了——会给那团裹着他双手之作的、压低嗓门的嘲弄合唱,添一个\dmark{}\origpage{1225}终止,或至少
添一个更知趣的弱音器。而倘若这次复现不同时终结那场对一种视野、对一些强而丰
饶的观念的时髦抵制,那它或许至少能打动某一位比别人更大度的年轻数学家,让他毫无保留地从中
汲取灵感(不惜惹人嫌弃),怀着敬意把它们化为己有。

然而,我写《收获与播种》,不是为了这些事里的任何一件——它们当中有的兴许会附带着来,谁
知道呢!我写它,是“为我自己”的,当然,如同我写的一切——作为一场摸索着寻找自身的理解的
手段。但与此同时,对他人的念想——对那些我爱过、而我的历险把我带往\textbf{别处}的那天撇下的
人的念想——贯穿了《收获与播种》写作的始终\fnote{1019}{(*)}{这一念想,在《愚妄与更新》(《收获与播种》的第一部分)里不止一次得到表达。它兴许在后文各部分里不那么显眼,却同样在场。}。这些笔记,既是一场思考,有时也是一场冥
想,对我而言,曾经并且始终是一份\textbf{馈赠}——在我自己之外,送给那些我向他说话的人。当然
我也知道,这份馈赠,也许除了我自己,无人会接。但我并不因此后悔写了它。再说,今天若没有被
某个它属意的人接住,也许明天会被接住。这份既自发、又经长久酝酿的见证,其中每一页、每一个
字都来在其时、来在其所——它明天不会比今天更不真。而无论今天还是明天,若有出乎意料的、
被欣喜迎接的事,那将是得悉:我的馈赠被接住了,哪怕只被一个人——他在我之中认出了他自己……

\subsection{(3) 信使(2)}   % 18.7.3

\notehead{181}{}

比起\emph{葬礼}图景的“前景”,我更不觉得有兴致,就另两个“平面”——一个由“我那班生意兴
隆的学生组成的、扛着铁锹绳索的队伍”,一个由“教团全体”组成——对我在他们身上的所悟与所
惑,作一番原原本本的回顾。关于后者,以及它在\emph{葬礼}里的角色,我在注“掘墓人——或教团全
体”(n\textsuperscript{o}~97)里已经说得相当周详\fnote{1020}{(**)}{(6月22日)我对教团尚且模糊的知觉,最近在已引的注“家庭相册”(n\textsuperscript{o}~173)的 c.、d.、e. 三部分里,出人意料地变得具体了。}。至于我对我那些亲爱的前学生的角色与动机的困惑,它们最清楚地显
现在注“沉默”(n\textsuperscript{o}~84)里,却\dmark{}\origpage{1226}没有在思考的任何后续时刻被认真重审。所以,恰恰
是在这个层次——\emph{葬礼}图景的“第二平面”——我的功课最欠火候!\fnote{1021}{(*)}{(6月22日)不过,关于图景“第二平面”的思考,确有一份(微薄的)续篇:6月19日的注“一场屠杀中的五篇论文——或孝道”(n\textsuperscript{o}~176$_{7}$)。}。在那里,没有一桩功课可与我在
“掘墓人……”那条注里做的相比。图景的这一部分,后来在两道更晚的注里、在阴阳动势的光照下
得到了深化:“天赐的机缘——或封神”与“拒认(一)——或召回”(n\textsuperscript{o}s~151、152)。

“掘墓人——或教团全体”这条注,是思考“第一口气”里写下的最后一道,想必也是它的顶点。可
是,时隔将近一年,我已经不再相信:在我这卑不足道之人的\emph{葬礼}背后(它被看作“对异端的
报复”),那股显得相当明显的集体动机,在集体意志的层面上,真正触到了\emph{葬礼}的\textbf{神
经}。使我生疑的是:在我每一个学生的情形里,这股动机在我看来完全缺席,或者,同其他在场的
力量相比,无关痛痒\fnote{1022}{(**)}{这一事实,出现在注“天鹅绒的脚掌——或微笑”(n\textsuperscript{o}~137)的思考里,打字稿第644--645页。}。可是,\emph{葬礼}里最触目的事实之一,恰恰是它的三个前后相续的“平面”之间
存在的那种“全体一致”:它们的作为与不作为互相衔接、互相补足(仿佛由一个“天衣无缝”的共
同意志指挥着),同一场名副其实的葬礼上一般无二!在如此显著的全体一致里,在心境与作为的如
此整齐划一里,人们当能猜到一个共同的动机,一根同一的“神经”,在驱动着这一个与那一个。

我无意暗示:那股我这里那里察觉到的“弥漫的怨怼”——它起于我的“异端”,(浅浅地)被感受
为开小差,(更深地)被感受为不可容忍的挑衅——全然子虚乌有、全然不起作用。但我如今怀疑,
它的作用是否是决定性的;怀疑它是否就是那根共同的“神经”——难道它人人皆有,却偏偏不
在那些\textbf{在葬礼里戏份最重的人}身上!也就是说,不在我那些学生身上——他们恰恰是那
份遗产的第一批受托人……

这桩“缘由”(我那场“异端”——表面看还相当讲理)在我看来,同那股暴戾之气完全不成比
例——那股气,\dmark{}\origpage{1227}我是在“屠杀一个‘光辉讨论班’”那样的行动里、在教团默许的目光下感到的;同反常层会议上那场当众鼓掌的、同样凶暴的不义,也不成比例。这也不是说,我曾是一
个可憎的同事或老板,怕我怕到积攒的敌意不敢发作、非得等我死了葬了才一齐倾泻到我头上、倾泻
到人们多少“认出我”的那些人头上。从我这里那里收到的种种回声看,既没有\textbf{畏惧}的痕迹
——仿佛我这人曾叫人害怕,如今报了迟到的一箭之仇\fnote{1023}{(*)}{的确,在“愚妄与更新”里,我长篇谈过那种“知名人士”身边——从一个我也说不清何时开始的时刻——笼罩着的\textbf{畏惧},我自己也时或在我周围察觉它的迹象。但那是附着于“知名度”本身的弥散的畏惧,而不是对我这个人本身的——只要稍能建立一点个人接触,它就消失了。我印象里,在个人接触的层面上,人们与其说怕我,不如说把我看作“好脾气”。我确信,在那位“失手——或二十年之后”一节(n\textsuperscript{o}~27)里谈到的学生那里,情形并无不同:每次见面,某种“发怵”持续了很久。这种发怵,今天我看作一种弥漫性的内在不安(“Unsicherheit”)的征兆,它后来在对支配与轻蔑的姿态里找到了补偿与出口。在他众多的学生当中,我有机会认识的三个,个个都结结实实吃过他那副貌似“无缘无故”的恶意姿态的苦头。显然,那种安营扎寨、统治着数学圈子各处的精神,助长了这类悖谬的行径;而这些行径又回过头来塑造这种精神,给它打上那种叫人错愕的、闷声的野蛮印记……};也没有什么稍微具体的行迹可以
归\textbf{罪}于我、足以滋养某种敌意或暴戾(而那种敌意,却从来不报姓名)。

这正是我称作“无缘无故”或“无因”的暴力的典型情境。这种暴力,最终来到“阴阳之钥”这场长
考(它本身构成《收获与播种》的心脏)关注的中心,想来绝非偶然。这种暴力,我认识它远非昨
日始;我头一次同它打照面,并不在我的数学生涯里。即便我有时忘了它在人间世界的存在,那也
从来长不了:它自己会很快前来提醒我,好意不忘。至于当下——借着一个奇怪的“巧合”(我
承认,它常常不合时宜,或至少不受欢迎……),我不记得我此生曾在什么时候,同这样一种暴力的
\dmark{}\origpage{1228}熟悉的征兆打过这么执拗、这么重复、这么磨人的交道——自从我“重返数学”以来,尤其自写《收
获与播种》以来;而最近的这几个月、这几个星期,愈发如此。

\pmark 当然,这里有一个执拗的讯息,它一遍遍回来,而且想必会一直回来,直到被听见。在阴阳长考
的最后几个星期里,我开始侧耳倾听它——同时知道,它要对我说的,我还没有听完。此后两个月
里,一种地下的功课想必在无声中进行。我觉得,那\textbf{要紧而隐秘}\fnote{1024}{(*)}{写下这一行时,我心里清楚,“隐秘”一词在这里是个权宜之计,是对“共识”的让步。我常常有机会发现:当我得知一件我一生都不知道的事时,那件事绝非“隐秘”,恰恰相反,它明摆着、显而易见,有时甚至刺眼——只是我不肯看它而已。发现新事物的情形,大抵如此,无论那是数学的工作,还是认识自己的工作。这种失明的根由、这种常识或直觉官能的阻滞,绝不在这些官能的缺陷,而在于精神的一种近乎不可克服的惰性:它不肯离开共识已然铺好的车辙——无论那是全社会接纳的共识,还是自己所处的那个较小的圈子里的共识,甚至,它们只在我们内心缔结并加盖了印章,如同某位“老板”为他一己之便、同他自己缔结的条约条款……}的东西,开始从表面可见的(或
至少不那么难于承认的)杂事里沉淀出来。“侏儒与巨人”的形象(我朋友皮埃尔供认的)继续萦绕
着我。在这形象背后,我相信辨出一股力量可观的原型:它像童年所受压抑投下的影子,或者
影子之一。它的角色,该是为创造之力所受的压抑提供一个出口与补偿——那压抑,早已内化为那句
“说不出口的无能之信”……在这个预感到的原型里,我仿佛触到一股强大的引擎,驱动着
那些无端施暴的行径,打击那个被当作“巨人”、被当作完好无损之力的携带者的人——这些行径的
爆发,不需要别的“缘由”,只需要一个\textbf{合宜的时机},当下手的风险显得为零、或微不足道。

也许我已经说得太多,而这几行字,才刚刚拂过一道纤细而执拗的直觉——它向我示意:有一桩功
课要做,它还横在我的前头。对这桩功课而言,\emph{葬礼}只是材料之一,同许多别的材料一道,来自
我那所谓“私人”的生活\dmark{}\origpage{1229}。这里不是继续它、甚至不是接近它的地方。它的位置,不在这些
注定要发表的笔记里。

\subsection{(4) 失乐园(2)}   % 18.7.4

\notehead{182}{}

(4月4日)在这场已许诺的回顾——我的思考教给我关于他人的事——里,我的思绪不由自主地
执拗地回到我自己的身上。这对我是个好兆头——它标志着我心中那股回到要紧之事的强烈需要。
对我的认识,额外地带来对他人的理解,而不是相反。自我开始冥想以来,不止一次,“理解他人”
的心思,成了那桩要紧功课——认识我自己——的消遣与岔路。

在有意地回到我自己之前(逆着我那股要抵达那个有名“句号”的不耐!),我还想收入一份最近到
我手里的、关于我朋友皮埃尔的见证。这是我离开数学舞台以来,所知的此类见证中的孤例。它给我
朋友的照明,与我通过别的途径所知的,大异其趣。它也十分适时地再次提醒我:现实,永远比我步
履蹒跚地试图给它画的像更复杂、更丰富\fnote{1025}{(*)}{我绝无意借此暗示:人们(包括我自己不停地)为现实画一幅尽可能“忠实”的像、并随各种“信息”的到来不断调整这幅像的努力,是徒劳而无效的。恰恰相反,这里有一种效力极大的辩证法,让我们同现实接触、去“认识”现实。只有当这幅像(它按事物的本性,自带一份惰性)完全僵死、板结之时,它才成为把握现实的障碍,或者更确切地说,成为让我们的把握官能\textbf{失效}、把我们确实拥有的对现实的知识“蒸发”掉的\textbf{有效手段}。}。

这份见证不是直接的。它关乎一位外国数学家同Deligne的一次(或多或少偶然的)会面留下的印象;
那位同事把这次会面对我的通信人讲了(我猜还是热乎的),通信人又在一封信里把经过转述给我。
经我的通信人与那位讲述会面的同事(下文称他“Z”)同意,我在这里翻译信中涉及这次会面的部分。
我的通信人推测,这一幕应当发生在1981年。(注意:这也正是反常层会议之年——不过通信人与我
之间,倒不曾谈起过那场讨论会。)

\begin{quote}
\pmark\origpage{1230}“……有一天,Z.~去 Bures 参加一个会,发现自己置身一间上茶的屋子[看
来是法国高等科学研究所的‘茶室’],屋里有很多数学家。只见门开了,Deligne走进来。Z.~先生把
这一幕讲得相当生动:他看上去蔫蔫的,双臂低垂,人觉出他周身有一种孤立。别的人都定定地看
着他,有点像看一只珍禽,又没人知道能对他说点什么。Z.~靠窗独自坐着,Deligne有些迟疑,在他
身旁坐了下来。Z.~不知道跟他说什么好。这时他起了个念头,就简单地说:他觉得围绕‘平展拓扑’
等等的那一整套想法,还有您带来的那些新思想,非同寻常。[“您”,这里及下文,指的都是我,Grothendieck——我的通信人是转述给我的。]Deligne的眼睛顿时亮了,他对 Z.~说:是啊,那是数学里
最好的东西之一;听您的报告\fnote{1026}{(*)}{同前,这里的“您”指的就是我,Grothendieck。},那才叫美……他接着说:您想想这个,再想想那个……数出一大
串事情来——Z.~自己(据他向我说的)一窍不通,但他看见了谈话对方身上骤然涌起的热情。Deligne又说:多可惜,您\textsuperscript{(*)}退了!他确信,晶体上同调和许多别的东西,不至于落到现在这副半
生不熟的样子;要是您\textsuperscript{(*)}真还接着干,它们如今早就是站得稳稳的构造了,就像平展上同调
那样……”
\end{quote}

这个故事里有两件事击中了我。其一是那种孤立的印象——它似乎深深击中了 Z.~先生。这一印象究
竟出自Deligne生涯里某个特别的时刻,还是这样的孤立终于浸透了他同所有同侪的关系,我很难说。
我没有收到过任何指向后者的别的见证。

另一件同样触目的事,在回到我这里的种种回声之中同样独一无二,就是提到我的名字与某段过去
时,那份热情的骤然涌现。那段过去,他早已决意宣告无效作废。还有他在这段过去里的根。而在那段
过去里,也还有一种童年的清新——那种被他逐出自己的“成人”生活、\dmark{}\origpage{1231}逐出这个重要而受人仰慕的
男子的生活的清新。在他周围,识相的人想必都不去提那段过去,不去提他还只是众多学生
中的一个、还怀着一腔美好激情的那些年月——正如在阔人宅邸里,满屋老样式的家具,谁也不去
提寒素甚至窘迫的起家……

而偏偏是这个陌生人,鬼使神差地坐在他身旁,忽然热切地谈起这件从来无人谈起的事(至少不当着
他的面谈……)!当然,这就像那团精选的、拘谨的空气骤然消散了;一个陌生人的热情,唤醒了他
身上同样的热情,并且——在一瞬之间——把他重新接通到那眼遥远的泉源,那眼他被认为早已忘
却、早已失去的泉源……

% ================================================================
%  18.8 Découverte d'un passé (原书 p.916–926)
%  源文: _work/chunks/chunk15.md
% ================================================================
\section{发现一段过去}

\subsection{(1) 第一口气息——或确认}   % 18.8.1

\notehead{183}{}

我终于来到这场一个多月前开始的回顾与盘点里最切身私己的部分。还剩要快速巡视的,是这场思
考就\textbf{我自己}教给我的东西。

思考让我发现的第一件事,是一段\textbf{过去}——我的数学家过去;此前我从不曾费心在上面停留
过哪怕一瞬。在那光鲜无事、看似平淡的表面之下,我又一次看见了深处的种种展露出来:它们通常被忽略、被抹去(仿佛一把扫帚悄悄扫过),从我们习惯为自己、为周遭描画的那幅舒舒服服的自觉图像里扫了出去。在那些从未检视过(至少在我的数学生涯里如此)的“涂痕”(或扫出门外的尘屑……)之中,有一种阴险、有时咄咄逼
人的东西:我在同某些朋友的关系里的自负。打一开始,这种自负就取了一种数学精英主义的形状;
它秘而不宣,我毫无察觉,因为我的态度看上去太天经地义了。这种精英主义(或如Chevalley与盖杰所
称的“绩优主义”),想必逐年变硬。它结晶为那种“竞技”的姿态——我是在思考“第一口气”接近
尾声时才终于发现的。在一副规规矩矩的门面底下,这种姿态认可的,是一种占有式的、嫉妒的心性,
指向那些被当作我的“禁猎场”的东西,以及那些凭其出色的品质、我乐于放进禁猎场来的人。

\pmark\origpage{1232}所幸,这些非常“老板”的心性,并没有穷尽1948到1970年间,我同数学世界里的朋友、同事
与学生,同数学本身之间关系的内容——差得远呢。尽管如此,它们构成一条阴险的低音线,去年以
前我从未费神去注意它;直到《收获与播种》的第一部分(“第一口气”)。这场渐进的发现,在“竞技
数学”一节(n\textsuperscript{o}~40)里达到顶点。这一节,在我看来标志着思考中一次质变的时刻。当时我立刻
把它感受为一次\textbf{翻越垭口}——它蓦地为我打开一片新的全景……

再退后一年回望,如今我把1948到1970年间、我在数学家中度过的这段漫长的数学生涯的第一
期,看作一场\textbf{交易}:以我本来所属(人人本属)的“长子权”——去完完整整地(倘若这是我
的选择)经历一场独特而独一无二的历险——换来的,是一碗“扁豆汤”:与一个田园诗般、虚构的
“数学共同体”认同(我希望这认同毫无保留,却从未完全做到……),而那个共同体同时又发放
种种舒服的好处\fnote{1027}{(*)}{这就是前面一条脚注(打字稿第1219页注(*))里谈到的那层暧昧。}。借这幅图像,我并不声称已把那段岁月说尽——它太丰厚,套不进一句现成
的公式。但我觉得,它抓住了第一次思考阶段里首次显形的一个重要侧面。这个侧面,又出现在《收
获与播种》这一部分(事后)取得的名字“愚妄与更新”里。

这第一阶段最切身、最深沉的部分,由最后三个“章”\fnote{1028}{(**)}{当然(如我在《收获与播种》引论里所说明的),这些“章”——把由共同主题或特殊亲缘联系起来的相邻各节归拢在一处——是事后设立的,在写完(那将只是)《收获与播种》第一部分之后。在《愚妄与更新》里,我偶尔把它们当作《收获与播种》的“部”来援引(不要把它们同那五个“部”——“愚妄与更新”等等——混为一谈;1984年2月至今的整个思考,归拢为那五个部)。}:第六至第八:“收获”、“孩子自娱”与“孤独的历险”。在“收获”里,我先同我生命中的某些
时刻(不只是我的数学生涯,这一回不是)重新接上头——一些满蓄更新之力的时刻。仿佛受一股
未知之力、一道隐秘而急迫的声音驱使,我寻找着重拾那些\textbf{天真}的心境\dmark{}\origpage{1233},好跨过那道我
隐约还感到横在前面的门槛。当然,当时我无从预料:接下来还等着我的,是发现我对数学本身
的一种占有式的态度。我继续爬坡,不急不缓,仿佛双脚踩着一条只有它们自己才“看得见”的无形
之路。不必对自己说,我也知道它引我去该去的地方;而脚下的雾,一步一步,渐渐散了。

就这样,我抵达了旅途中这道新的门槛,或者说这道\textbf{垭口}:

\begin{quote}
“……一到那个地方,我就有了那种来到瞭望台的人的感觉:从那里,他望得见自己一路走来的景
色——行走之时,他每一步只能看见其中一角。而此刻,有了这种对广度与空间的知觉,这是一
种解放……”
\end{quote}

刚一越过思考的这个敏感点,它便深化为一场对我自身的冥想。第二天,我已感到需要引入“老板”
与“工匠”的形象——别名“孩子”;这形象,两三年来对我已相当熟悉。但我远没有料到,它在接下
来的思考里将证明何等有用——那时,近两个月来我一直自以为马上就要收尾,好马上回到我的数
学笔记,回到《追寻栈》去!

在构成“孩子自娱”这一“章”的四节里,我重新接上了我同数学关系的某些侧面与波折。这些,近
三年前(1981年7月至12月)我已细细探测过,可此后有的是时间,大半又忘了。这一回,我的宗
旨主要是让自己进入状态,去探测我那次不期然而重返长期数学投入的意义,并试着在如今支配我
生命的两种激情之间——表面看互不相容——“给自己”定位:数学,与冥想。

顺便说,这两股激情的“互斥”,如今在我看来已不像两年前那样严酷了。在《追寻栈》里,
数学思考有时让位于、甚至化作一场多少个人化的思考;在那里,我这个有感受、有情感、有(不只
是数学的)好奇、有命运的存在,不再完全缺席。反过来,在\dmark{}\origpage{1234}《收获与播种》这场对我自身的思考里,思考本身又把我同旧日的数学之爱重新接通,并且这里那里,成了数学思考的引信
\fnote{1029}{(*)}{(5月10日)这些“思考的引信”,至今已结了果:我对若干搁置了十五年的主题的理解,得到了更新。}。

这些共存的可能——乃至共生的可能——这两股我心中求知冲动的不同表达,按事物的本性,或许
注定要相当有限。但无论如何,去年的思考(甚至更早,三年前那场长考)已让我看清:这两股
激情绝不互相敌对,甚至本质并无不同。在思考的最后一部分“孤独的历险”里,我竭力贴身把
握:这两股激情究竟在何处不同,它们各自为我打开的“历险”又如何不同。正是在这个追问
中,我发现了一个明摆着、而我一辈子佯装不知的事实:数学是“\textbf{一场集体的历险}”;我个人的
数学历险,只有系于它所属的那场更宏大的集体历险,才获得意义。

说实话,在“孤独的历险”一节里,起初我只是顺带掠过这个事实;当时的宗旨,倒是要用言词把
另一件事勾勒出来——那件事我素来深知,却始终难以全盘接受:冥想,才是一场\textbf{孤独的}\textbf{历
险}。给一件“已知”的事作表述,这番功夫绝没有白费,远非如此!它深化了这份认识,并且顺带让
我发现(至少对我而言)一个明摆着的新事实:把我系于\textbf{另一场}历险——数学的历险,它倒是
集体的——的那条纽带(当时,我、或者说当时的我之中的某人或某物,本想同它保持距离……)。

地基由此打好了:第二天,在“分裂的确认”一节里,我径直进入我困惑的中心。首先是确认:“老板”
的下注,哪怕他想自欺(那本更合他的天性……),也只能是那场集体历险——唯一能给他“回报”
的历险。“孩子天性孤独”;\dmark{}\origpage{1235}唯有孩子,才会被一场世上无人想要的历险、一种切实可触且常常明摆着
的知识所吸引,而这知识,他无法同任何人分享。而此刻,在我这桩“事业”里,“小
家伙的偏好”,偏偏不合“老板”的意,正落在这一边。

这个确认,引出对我之内一场\textbf{分裂}的确认:\textbf{老板与孩子的分裂}。第一次,我在如此专注、
如此严格的心境里,作出这样一个确认。这不是一道\textbf{判词}——依照这样那样的“看法”或哲
学或随便什么拟出、自命多少具有普遍效力的判词。它只是一个简单的\textbf{确认},出自对一个极特
殊的个案的细心检视:我这卑不足道的人,在我发展的某个阶段。也许有一天,这道分裂会消失,
而老板照旧去做必要的事,任工匠孩子随意做工。这不是我今天操心的事,也不必是。一天的难处一
天当就够了……

(4月5日)确实,这道分裂,差不多九年前,已在一个梦里,借一场以震撼人心的力量搬演出来的
寓言,向我揭示过。那是在发现冥想——这种长期被我无视、随时待命的力量——的两天之后;正
是掘到这个梦的意义的深处,我才重新找到我之内那个不分裂的东西,\textbf{另一个}我——沉默而
不可见了那么久的,“一位至亲之人,被人当作死了一辈子……”(见《愚妄与更新》)。当时出现的
新东西、要紧的东西,并\textbf{不是}那道我熟得不能再熟的分裂,也不是梦以那样大的力量向我揭示
的、这道分裂的本性——它化身成两个亲爱而熟悉的存在,谁也没有名字,而他们本是\textbf{同一
个}——而是那场\textbf{重逢}:在四个钟头高强度冥想(如同分娩的剧烈阵痛)之后来到的重逢。

我当时就知道,此后几天、几星期里也一直知道:重逢并不等于分裂的终结。但托重逢的福,我用
新的眼睛看这道分裂了——它固然要紧,但在另一个更根本的真实面前,终究是“次要的”:那是我
重新找到的我之内那个\textbf{不分裂}的、不可摧毁的统一,我后来认出它就是“孩子”。这一双重的
\dmark{}\origpage{1236}认识,当时鲜活而敏锐地在场。随后几年里,它钝化了——那份对这道“次要”却真真切
切的分裂的认识,渐渐被收藏起来。而当“老板”由着自己一头扎进冥想(那匹出名的“三脚马”……),
他很想暗示(却始终没有胆量、或没有笨到明说……):有了冥想和这一切,分裂从此已成过往,几乎
可以说没有分裂了——顶多这里一小道涂痕、那里一道,好吧,咱们不否认,但总归几乎跟没有一样;
只消看看小家伙工匠多么兴高采烈地埋头干活,看看一团和气的老板怎样踮着脚尖走路,生怕惊扰
了他——真可谓琴瑟和鸣!我不免疑惑:去年那份思考——翻越垭口(“竞技数学”)之前的那份,
尤其当我出人意料地回顾“我的激情”时(同名的节,n\textsuperscript{o}~35)——是不是恰恰还带着这样
的调子:光照逼得人对玫瑰色起一丝疑心……

不管怎样,这道“分裂的确认”当时恰到好处地让我重新接上了一个我多年渐渐看不见的真实。同
时,它让我以新的照明、新的眼睛,重新看见八年前我曾知觉得清清楚楚的那道分裂。我可以毫无保
留、毫无怀疑地这样说,因为我记得清清楚楚:作那个“确认”的时刻,同重逢的插曲、同那插曲就
某道分裂及其本性教给我的东西,没有任何联想!这一联想,直到稍后——我重新拾起前一天的笔
记时——才姗姗而来。这再清楚不过地表明:那次插曲中显现的认识的“次要”(而且不受欢迎!)
的内容,被收藏得多么干净。而且当年做来想必格外容易:重逢那道关键转折之后,没有人就这份
内容作过思考;而最能表达它的那幅形象——“老板”与“工匠孩子”——要再过好几年才出现。

如今我似乎看清:这道重新作出的“分裂的确认”,才是《收获与播种》第一部分里我对自己学
到的最要紧的东西。这个确认,不过几行字,是思考这一部分里最短的几节之一。\dmark{}\origpage{1237}人们也许
会想:既然只为了走到这一步,或许犯不着再用一百五十页,去追索自负在我数学生涯中的种种曲
折显形。就流行的“常识”而论,这千真万确。可同样真确的是:这种用修枝刀裁出来的“常识”,根
本无力把握一场发现工作的幽深曲折——无论是认识自己,还是更粗朴的数学发现的工作\fnote{1030}{(*)}{科学发现的工作,在我看来之所以比认识自己的工作“更粗朴”,(我想)有两个缘故。其一,它几乎只调用我们的智性官能,也就是我们存在中极小的一部分。(而且,科学工作倾向于让这部分官能增生肥大,以其余官能、以人格的整体平衡为代价,及至把人变成一种电子计算机怪物……)其二,科学工作所调动的、对抗揭示真实的内在抵抗,通常与对抗自我认识的抵抗,不成比例。这也是为什么,“科学历险”很少——今日几乎从不——是一场“求真的历险”:一场调用我们谦卑的能力、调用我们承担不期而至之真实的勇气的历险——首先对我们自己,然后对外部世界。}。我有一个切身的信念:在这场《收获与播种》的长考
里,每件事都来在其所、来在其时,由先来的一切所准备、所酝酿。

\subsection{(2) 第二口气息——或调查}   % 18.8.2

\notehead{184}{}

(4月6日)随着这道关于分裂的短短确认——去年3月末(一年多前),我起初以为自己已经写完
了《收获与播种》。我远没有料到,它还要再来五倍之多(按页数算)!随后几天,我做些杂事,思绪
开始回到数学主题。可是,思考中还有一个悬着的“小尾巴”,继续在我脑子里打转。越过一层看似
纯然细节的困惑,我隐约感到:对于“老板”转向一项长期数学投入的背后那股力量,我还没有真正
走完一遍。或者说,即便我把它的主要机关都翻了出来,我的理解仍旧苍白而游移——只因为在这
件事上“摆”得还不够久,没能让它渗得更深。这个“最后的尾巴”,将成为我回过头去看那团模糊印
象的门径。思考的重新接续,发生在当时(并且此后三个星期里)被当作《收获与播种》终篇的那一
节\dmark{}\origpage{1238},它旋即得到名字“一段过去的重量”。这个名字,说尽了那场意外的发现:我的数学家过去的
这份\textbf{重量},以及把我继续系于那场集体历险的纽带之力。而且,那天我所瞥见的,还只是
冰山露出的那一点小小峰尖;那沉在水下的庞大部分,要在随后的几个月、整整一年里,才渐渐
显形……

为思考的第一口气收尾的这一节,同时像第二口气的引信与召唤。这份“过去的重量”,显然扎根于
我对一份著作的眷恋——与其说眷恋那部完成了的、做到终点的著作,不如说眷恋一些执念观念与
一些洞见;它们的丰饶与力量,我感觉得到、我“熟知”入骨,而我多少含混地、多年以来就意识
到:它们却在贫瘠而干旱的土地上萎顿着,那土地暗地里、阴险地怀着敌意……于是,“一段过去的
重量”这场思考——它唤回我的好记忆:那部著作,以及我同它的纽带——成为一道长注的机缘;自
“离去”以来,我第一次就这部著作、就它遭到的命运表明态度。十年来、十五年来含混感到的东西,
终于成形,化作言词——言词有时迟迟不肯来到;而一旦白纸黑字写下,它们清清楚楚地向我道出一
个讯息,直到那时我一直在回避去知悉的讯息。后来,鉴于这条一气写成的长注的篇幅,我把它分成
两道,取名“我的孤儿们”与“拒受遗产——或矛盾的代价”(n\textsuperscript{o}s~46、47)。

\enlargethispage{2pt}
可以认为,这双注,就是\emph{葬礼}思考的发令枪\fnote{1031}{(*)}{可惜,这一情节在“葬礼(I)(或中国皇帝的新衣)”的目录里看不出:那道双注构成送葬行列之二(“孤儿们”),而非行列之一(“身后的学生”)。缘由在于“愚妄与更新”最后一节“一段过去的重量”(n\textsuperscript{o}~50)——那几道注本应是对它的评说——之中,诸“注”(n\textsuperscript{o}s~44至47)引用的先后次序。}。三星期后,4月19日,思考继续推进:那本“难忘文集”LN~900激起的
震动,为 Motif 的掘墓复生、由“新爸爸”Deligne主持,作了证。思考的这“第二口气”,一直高密度
地持续到5月底、6月中——当时我再一次自以为马上要画上句号,货真价实的句号!——以一场病\dmark{}\origpage{1239}
中插曲告终\fnote{1032}{(*)}{关于这场病中插曲,见“事件——或身体与精神”与“陷阱——或便利与衰竭”两条注(n\textsuperscript{o}s~98、99)。}。

这第二口气,严格说不是一场对我自己、对我过去的思考,而毋宁是一场对刚发现的\emph{葬礼}的
“调查”,同时是一场努力:把那些明摆着却(大概鉴于我那拔不掉的天真!)令人瞠目、令人难以置
信的事实,边走边尽量“消化”。不过,它毕竟就我自己教给了我一些东西:主要是让我切身掂出,我
对我的过去、对我的著作的眷恋之力。我被触到了痛处:眼看着著作被一块块撕走——有的进了垃
圾堆,有的供人取笑,有的被肆无忌惮地据为己有,当作谁都可以拿的小玩意……

那时我知道,我并没有像我一度相信的那样(在翻过某个“垭口”、眼前展开广阔全景之后
的那阵欢腾里\fnote{1033}{(**)}{这阵欢腾,表达于“旋转木马收摊了!”一节(n\textsuperscript{o}~41);五六星期后,在注“一脚踏进旋转木马”(n\textsuperscript{o}~72)里被压低了调门。})“离开旋转木马”。换句话说,我得以掂量这份过去的全部\textbf{重量},以及把我系
于其上的那些唯我机制的全体之力。这是个大意外!

然而,在思考的第二阶段,我关于自己还有另一项发现,补足第一阶段的所获。第一阶段,我翻出的
主要是我身上自负态度的某个“反面”:以\textbf{排斥}的姿态,对待某些同事甚至朋友——因为这样
那样的缘故,我不把他们算进我自觉身属(当然,是默而不宣的!)的“精英”的世界。同一枚奖牌
的\textbf{正面},则是我在同更年轻的数学家(尤其是我的学生)的关系里的一种\textbf{纵容}与暧昧:我
可以说把他们“吸纳”进了“我的世界”;或者因为他们手段出色,或者仅仅因为我收他们做了学
生,他们从此便被我知觉为处于我的“庇护”之下。我是在5月10日的注“高升”(n\textsuperscript{o}~63$'$)里开始
点到这种态度的,继而是5月27日的注“格格不入者”(n\textsuperscript{o}~67$'$),两道都写我同我年轻而\dmark{}\origpage{1240}光辉的
朋友皮埃尔的关系。这场思考在6月1日的注“暧昧”(n\textsuperscript{o}~63$''$)里深化,转向我同
我的学生们的一般关系。正是在那里,我终于察觉某种暧昧:它从未被我察觉、检视过,一直跟着我
到近几年。而且,就在最近,我又在稍有不同的语境里同这种暧昧打了个照面:在子注“驱逐(2)”
(n\textsuperscript{o}~169$^{1}$)(其第二部分,日期为3月16日)里。我不得不在那里确认:把我本人从 SGA 讨论班
中排挤出去——那个讨论班,是我十年生命投入的总和\fnote{1034}{(*)}{关于这场排挤最新的一桩,见注“殡仪馆——‘为科学服务’”(n\textsuperscript{o}~175)。}——这一排挤,主要由我旧日学生中最亲近的某几
位操办,不过是我在他们面前乐于经营的那种暧昧态度的自然收成:关于他们在那部规模恢弘的
SGA 里、与我在其中的各自位置——他们中的某一位,曾在其中投入一两年光阴。

\subsection{(3) 第三口气息——或发现暴力}   % 18.8.3

\notehead{185}{}

还剩要巡视的,是思考的“第三口气”就我自己教给我的东西。第三口气始于去年9月22日(病中
插曲之后),眼下正接近尾声\fnote{1035}{(**)}{《收获与播种》的第五部分,我在这里暂且撇开;它最初本是葬礼仪式内部——甚至就在“阴阳之钥”内部——的一段“枝蔓”。写这篇回顾时,这一部分尚未完成,故未收入。(6月22日)随后几星期里事情显明:“阴阳之钥”之后的“四种运算”部分(即葬礼(3)),构成《收获与播种》的“第四口气”,同样不在这篇最终回顾之列。}。这里首先指的,是在“阴阳之钥”里继续的那场思考;在我看来,
它是《收获与播种》最切身、最深沉的部分。没有任何刻意的安排,我的这个人、我与世界的关系,
在那里最常处于注意的中心。当注意力似乎偶尔岔开,转向一些表面上更一般的主题,或在我朋友皮埃尔其人身上流连,它却总是回到中心、回到那位演员兼观察者、那个感受着、知觉着、追问着、探
测着的人身上,仿佛被一股无形之力吸回。所以,先于一切、尽管并不以之为目标,它是一场\textbf{对我
的生命、对我自己的冥想}——从一个意外的入口进入:那场\emph{葬礼}。

这也是思考中在我看来最丰饶的部分,是我从中获益最多的部分。许多“已知”的东西彼此各安其
位\dmark{}\origpage{1241};一些仅仅瞥见或预感到的东西,一些“知道”却被忽视、被淹没在日常杂沓的昏暗光晕里的东西,
开始从暗处浮出,既显出它们的重量,也显出它们的轮廓。像是一次新的敞开,一次向未知再次出发
的邀请——而在此刻,那场有名的“长期数学投入”眼看要让认识自己的工作一放就是好几年……

在这里,不适合原原本本地重述这场长考的各个阶段,也不适合为它教给我的一切开一张“清
单”。我宁可几句话说说:对认识我自己而言——作为一场仍在日复一日、月复一月、年复一年继续
的成熟的材料——什么在我看来最要紧。

这场思考开始时,我本意是开一个“括号”(顶多一两道注的工夫……),好让读者(顺便也让我自
己)“入水”,浸进一种阴阳(或“阴柔-阳刚”)的事物辩证洞见。开这个括号的缘由,是需要用阴
阳的直觉,去定位一个触目的印象:它来自对某一篇“悼词”的检视\fnote{1036}{(*)}{关于这篇“悼词”(其恭维之语,调配得法,施用有道……),见同名的两条注(n\textsuperscript{o}s~104、105),以及部分概括了它的注“宝石们”(n\textsuperscript{o}s~170(iii))。}:那是一种刻意的“角色对调”,发
生在一场原本的阴阳关系里。括号于10月2日打开。直到11月10日——在关于阴阳的把戏在我
的生命里、在一般的存在里、最后在数学发现的把戏里的整整一百页密不透风的思考之后——时
机才算成熟,可以至少把这个六个月前冒出的联想\textbf{说出来}\fnote{1037}{(**)}{在注“对调(3)——或阴葬阳”(n\textsuperscript{o}~137)里。};至于在知根知底的状态下去探测它,还要再过十四
天\fnote{1038}{(***)}{在注“天鹅绒的脚掌——或微笑”(n\textsuperscript{o}~137)的开头。}。(而著名的阴阳括号,要到将近两个月后、1月14日,才终于合拢——甚至过了一阵,我才察觉它早已合上了……)

\pmark\origpage{1242}很快,无须寻找也无须计划,“冲突”——在人的生命里、在人身上的冲突——占据了注意的
中心。\emph{葬礼}的发现骤然而强力地调动起来的唯我能量,在那里充当了一支援兵,让我再一次、
就在当下,直面多年以来一直向我发问的那个“冲突之谜”\fnote{1039}{(*)}{这一“发问”,主要开始被察觉,是在我1979年8月至1980年3月间持续的那场对我父母的漫长冥想之中。}。早在此前那些年里,这个谜已渐渐浮到我
想探测、想理解之事的前排,能走多远走多远——却从来没有“跨出那一步”、把自己整个儿投进去
……

渐渐地,思考之中显形出来的,是我生命里仿佛那枚“硬核”的东西,这个谜那可怕的震中,仿佛
“恶之谜”的心脏本身:一种可以称作“无缘无故”或“无因”的暴力——一种似乎单以伤人、害人、
毁人为快的暴力,一种从不报出姓名的暴力,常常闷声不响,顶着一副天真无辜、和蔼可亲的面孔,
因而在触及与摧残上愈发有效——“天鹅绒里的利爪”,纤细、敏捷、毫不留情……在“天鹅绒里的
利爪”那组注(n\textsuperscript{o}s~137--140)的思考中,注意最终落到的,正是这种暴力;直到《阴阳之钥》
终卷,它也一直处于注意的中心。它构成那把最高音:最后一道注里,那条“业”的“没有尽头的锁链”
——父母传给子女,子女传给孙辈,一代又一代,自太初以来——被唤起。

这是我一生中第一次,直面这个“无恨亦无情”的暴力之谜——一种深植于人类生活的暴力,自我
少年时代起,就在我生命里打下了洗不掉的印记。这也是我第一次,在我自己的存在里,确认这道印
记。这同时也是对一个简单事实的确认:这种暴力的\textbf{存在},它那可怕的无所不在——在我的生
命里,也在每一个人的生命里\fnote{1040}{(**)}{这一确认,构成注“无恨亦无容赦”(n\textsuperscript{o}~157)那场思考的高峰时刻。}。而这一个简简单单的确认,本身就在胚胎里含着对这个可怕事实的一
种\textbf{接受}。或许正是在这个确认里,藏着整部《收获与播种》的思考中我所\dmark{}\origpage{1243}学到的最重要的东西
(或者说,至少\textbf{开始}学到的东西)。

这并不是一个终局,不是一场思考的顶点。毋宁说,这仍是第一步,把我带过一道通向未知的门
槛。对我的行路与成熟而言,这个谦卑的步子,在我看来,分量比我所瞥见的那些“回答”的胚胎更
大——在随后几天里,我瞥见了“无因的暴力”之“因”这一问题的答案的胚胎\fnote{1041}{(*)}{见同名之注(n\textsuperscript{o}~159)。}。这个问题本身,只有
在充分看见并承认为其发问的那个事实的存在与分量之后,才获得它的全部意义——那分量,远比
一个“心理机制”的问题沉重。

有人会说,我这是跑题了:对一个(或我自以为)关乎人人的普遍心理事实的确认,属于客观知识的
领地,那是各门科学的专属(心理学、精神病学、社会学,或随便什么),不属于那出了名的“自我认
识”的领地(一种被感觉为含混而不可捉摸、倘若不算全然想入非非的领地)。可我看得很清楚——不
是含混不可捉摸地看,而是像看一个熟悉的、明摆着的数学事实一样清楚……:脱开认识自己,这样
的确认就失去了它活的意义——失去那使它不至于沦为一场哲学-心理学的文体练习、沦为一条
“论纲”的展开(当然十分有意思,如此等等……)的东西。这个确认,本身是一场\textbf{发现},一场贴
身的发现:世上没有任何人能替我作这个发现,我也无法替世上任何别人作它。这场发现,是那次发
现自己的旅程中,最近的一步——或者几乎是最新的。它把我相对于一件要紧而可怕的事定位:这件
事打过我烙印,而我直到那时一直执意忽视它,仿佛我是由于某种特殊的倒霉(也许系于我这卑不
足道的人的某些特点),才一生暴露于它;仿佛别人暴露于它、或施它于人——只要我肯费心睁眼看
看周围——也无非同样的倒霉。

而且这肯定不是偶然:对暴力的思考刚一起头,思考自身的内在逻辑,就引着我——也是平生第一
次——回头去看我记得的那几桩个案:我自己也曾\dmark{}\origpage{1244}\textbf{施}这种“超出理解的暴力”于人,当然,不
曾多想\fnote{1042}{(*)}{见《阴阳之钥》里“天鹅绒里的利爪”部分之后的那道注“正义者的暴力”(n\textsuperscript{o}~141)。}。这次回望的意义,不在于给我机会自责(还是在众目睽睽之下!)——这种事,我反倒完
全没有去做。而在于它为我打开了一扇门,通向对暴力更深的理解——一扇门,往后的日子里,何时
跨过它,只由我定。

\subsection{(4) 忠诚——或女性气质的数学}   % 18.8.4

\notehead{186}{}

这就是,以发现自己的旅程的眼光看,在我看来最要紧的东西。阴阳思考的最后阶段——以暴力为
中心——贯穿最后四部分:“天鹅绒里的利爪”、“暴力——或游戏与刺棒”、“另一个自己”、“冲突与
发现——或恶之谜”,从12月7日到1月14日(它们略多于《阴阳之钥》的三分之一)。

退后看去,我似乎觉得:《阴阳之钥》前八部分的主要角色,是最终把我带到这场关键的思考。我在
这准备性的部分里展开的许多东西,是多年以来我熟悉的事;却仍须“重提”,好让一位初来的读者
能够跟上,好给思考一种内在的融贯——否则它就有缺失之虞,或只有我一个人看得见。不时地,文
体显出这些内在心境的痕迹:写的人急着赶紧把重提打发完,好终于抵达“正题”——虽然,那些所谓
的重提,常常分量重得多,值得我在上面多停留片刻,比我那么急着要去的那个“正题”更值得(而
无论急不急,我到那“正题”,还是晚了一个多月……)。这些心境,我看在三节相连的部分里最为
可触:“夫妻”、“我们的母亲死神”、“拒与受”。即使在那里,真的,重新接上那些本该“已知”的事
情时,我仍旧不由自主同时重新认出了它们,而且有时是在新的光线下——哪怕是为了一些初看如
此非个人化的事情,比如清点那些由直接的亲缘联起来的阴阳“夫妻对”(或“锁孔”)——那些通往
世界的门。

但正是在接下来的三部分里(它们也正好先于以暴力为中心的最后四部分),我\dmark{}\origpage{1245}重新驶入此前未曾
勘探的岸域:“阴阳数学”、“阴阳的对调”、“主人们与仆人们”。

正是在这三部分的第一部分里,发生了那场“大意外”:它将为\emph{葬礼}的意义——或至少某种意
义——投下新的光。事情是这样的:在我对待数学的进路里,更一般地,在我走向发现世界的自发举
步里,我的存在的基调是\textbf{阴,“阴柔的”}。换个说法:尽管我的自我的条件化结构、我这桩事业
的“老板”是阳(不说“彻头彻尾的男子汉”吧),我本来的天性、我之中的“孩子”(他同时是那个
把孩子在玩耍中发现的东西做出形状的工匠……)却以“阴柔”为主。当然,单是这个特点,并不足
以把我个人对待数学的“风格”同任何别人区分开来。我甚至觉得,即使在数学家当中,这个基调(或
“主导”)本是阴的,也并不那么罕见。而我的情形里反常的(在我看来)倒是:在我的发现步履里、
尤其在我的数学工作里,我一生都对这本来天性完全忠诚,从不曾起过一丁点儿要修改它、纠正它的
念头——无论是出于内心检查员的意愿(反正他从来什么也没看见:在一门“男人之间”的事情里,
谁会怀疑一种“阴柔”的感受力与创造性的进路呢!),还是出于迁就外部世界、尤其是科学世界里流
行的好品味的法度。我毫不怀疑:主要正是靠着对我自己天性的这份忠诚——至少在我生命中这个
有限的领域里\fnote{1043}{(*)}{如我在《收获与播种》各处有机会一再说起的,自我八岁起(直到1976年,我四十八岁)支配我生命的两股最强的唯我之力之一,是对我身上“阴柔”特质的压抑,以利于那些被感受为“阳刚”的特质。只是到了“阴阳之钥”的思考里,我才意识到:这种压抑,并没有施行在我的数学工作上(后来,也没有施行在冥想、认识自己的工作上)。我存在中本来“阴柔”的主导,得以在一项通常被感受(也难怪)为最“阳刚”的活动里,尽情尽兴!(关于这一点,见注“最‘男子汉’的艺术”,n\textsuperscript{o}~119。)},我的数学创造力才
得以充分地、无阻碍地铺展,像一棵茁壮的树,稳稳扎根在大地里,依着昼夜、风与季节的节律自由
舒展。事情是这样成全的,尽管我的“天分”颇为平平,尽管开端根本算不上\dmark{}\origpage{1246}\textbf{吉兆}\fnote{1044}{(*)}{我说“天分平平”,绝非故作谦虚。这件事我一遍遍有机会确认:无论是同那些光辉的数学家打交道——他们把握要领、接受并消化新思想,比我敏锐得多——还是同这样一些没有名姓、没有像样数学底子的学生打交道,他们的好奇心与数学发明力一时被唤起。
我的“开端”(至少是我同数学家世界接触的开端,1948年),我在“受欢迎的异乡人”一节(n\textsuperscript{o}~9)里谈了一点。不过,还要早三年,1945年,我的“数学生涯”就开始了:我的精力的大半,投给了数学研究工作。直到1949、1950年前后,作为在法国的一个外国人,我以数学谋生的前景,还成问题得很。倘若这样的机会没有出现,我打算去学木匠——那是一门合我口味的营生。}。

当我在注“上涨的海……”(n\textsuperscript{o}~122)里,对我的数学进路作出这项出人意料的确认时
\fnote{1045}{(**)}{另见后文注“箭与浪”(n\textsuperscript{o}~130)。},它来得有点像一桩意外的奇事,有点“在边上”——因为在我的生活里,同他人的
关系,全都打着我阳的、超级阳的抉择的印记。只是在这场以冲突的动势为中心的思考的继续里,
借着对\emph{葬礼}的一次回望,我才意识到:我的数学家同侪同我的这个人、尤其同我的著作的关
系,被打上了这个不寻常的特点何等深的烙印——在他们心里唤起矜持的反射(若不是排斥),面对
一种被暗暗感受为“不合适”(不说“失体”)的进路风格。这类反应,早在我初入数学世界时就很有
些;只是在那宽厚的年月,它们被当时盛行的敬重他人的氛围调和着——至少在我落脚的那些数学
圈子里。后来,鉴于“Grothendieck结果之力”(借Borel致Mebkhout信中的话,那封信里提到了这些
“矜持”),它们也只好被压着。而自我告别数学舞台以来,它们成了常规,有时还在某副得体的低声
调门(那种低声调门,至今仍是规矩)背后舒舒服服地铺展:昔日的敬意早已风化、消失多时,而当事
人(据推定已死且已葬)也不在场,不能接话了……\emph{葬礼}的这个出人意料的面向——以我这卑
不足道的人为象征,埋葬“数学中的阴柔”——在“天赐的机缘——或封神”与“拒认——或召回”两
道注(n\textsuperscript{o}s~151、152)里得到探测;那是12月23日与24日,正在关于暴力的冥想的正当中。

\pmark\origpage{1247}还剩我这个人最后的一个面向想谈一谈——它是在写《阴阳之钥》时、在刚引的最后那一部
分“主人们与仆人们”里(那一部分紧挨着由“天鹅绒里的利爪”开启的思考转折之前)显形的。这
就是“服务的冲动”,以及它在我对数学投入的抉择里扮演的头等角色,以及它作为一股活的力,在
那些浩大而没有尽头的奠基任务里运转——我之后,还没有别人找到勇气(或谦卑……)去接手并把
它们继续下去。这个以非凡之力在我身上在场的面向,雄辩地见证了我本来天性中“阴柔”的主导:
它保存了下来(乃至避难……)在数学活动里(谁也不会想到去那里找它……)。

此刻我心中闪过一个念头:说不定,正是这股冲动,也以它非唯我的那一份,参与了那次“转向”——
转向高强度的数学活动,而把冥想的功课撂到了后台,一放就是不知多久。冥想,按其本性,是一项
孤独的功课;一项(在我看来)除非自欺、便无法纳入“为众人服务”、为某个“渴望知道的存在的理
想共同体”服务的投资眼光的功课。这么看来,似乎有一股深层的冲动,不同于唯我的求证之心、求
许之心,一股表达着个人与其所属的种属之间的深层纽带的冲动——它注定要在一场旷日持久的冥
想功课里受挫,按我对冥想的理解。而这也许是一条附加的缘由——除了那些来自唯我结构(“老
板”的心性)的缘由(它们本身已经足够强大)——使得这样的功课显得如此稀罕,以至于我拿不
准,可曾在别人身上见过它的痕迹。

% ================================================================
%  18.9 De Profundis (原书 p.926–929)
%  源文: _work/chunks/chunk16.md
% ================================================================
\section{自深处}

\subsection{(1) 感恩}

\notehead{187}{}

(4月7日)这场回顾与盘点——盘点全部《收获与播种》的思考教给了我什么——我相信已经走完
了一遍。唯有一处我从这份回顾里剔除出去,即《收获与播种》的第五部分\fnote{1046}{(*)}{(6月22日)还有第四部分(我眼下正在写的)也一样!见打字稿第1240页脚注(**)。},它至今没有写完。它最初本是“阴阳之钥”里的一段“枝
蔓”,后来延宕了整整一个月,落实为关于 C.G.~Jung自传的一百来页“读书笔记”。由于
这条枝蔓的收尾仍\dmark{}\origpage{1248}遥遥无期,我便把它搁置了。我最盼望的,是把\emph{葬礼}善始善终:写完、打
字、印出、四处寄发——然后,再没有人提它!

何况我预感到,这第五部分将就这同一个\emph{葬礼}给我带来意想不到的光亮——可不是嘛!——
通过我打算在其中对Jung与Sigmund~Freud之间关系的考察:有多年之久,弗洛伊德曾是尚
在寻找道路的年轻Jung的师父。初读自传里论及这层关系的那一章,我什么也没看出来——随后,
若干蹊跷之处引起了我的注意,我折回去看了几处,又把这一章重读了一遍。显然,这层关系绷
满了暧昧——弗洛伊德本人似乎强烈地感到了它,而Jung却津津有味地全然无视它(活脱脱像个
初来的讨论班新生……),把弗洛伊德的不安一概记在他那“神经症”的账上(他还饶有兴致地
把这条“神经症”描绘得有声有色,也许甚至绘声绘色得过了头,反而不大像真的……)。话虽如
此,种种联想随之而来,系于我的朋友、(同时也是)非学生Deligne与我的关系之上;这些联想我
打算追下去,兴许还要稍加开掘。我隐隐觉得,就所牵动的心理机制而言,\emph{葬礼}上发生的事,
绝非一桩绝无仅有、不合典型的巧合——恰恰相反!而我预感到,Jung与弗洛伊德的那层关系,
在这方面没准能提供更多的光亮。

不过对我而言,至少眼下,这第五部分(它也许会叫“Jung——或一场冒险的陷入泥沼”\fnote{1047}{(*)}{本想写“陷入泥沼”(enlisement),却看见自己写出了“葬礼”(enterrement)。这个口误提示的新名字——“Jung——或一场冒险的葬礼”——未必不贴切,甚至没准比原先拟的那个更切中要害。}),已经不再是
那场\emph{葬礼},哪怕它从中生出——我甚至要说:它已不再是《收获与播种》!它是“\textbf{此后}”
(l'Après)——同种种回声一样,那些回声,随着“愚妄与更新”、“葬礼(I)——或中国皇帝的新衣”
和“葬礼(III)——或四种运算”这三部分的寄出,势必要纷纷回到我这里,其中肯定少不了一些
尖酸刻薄的话\fnote{1048}{(**)}{不要把《收获与播种》的第四部分(副题为“四种运算”),与同在该部分之内、集结在这一名目之下
的那组笔记(n\textsuperscript{o}s~167$'$--176$_{7}$)混为一谈。}。到这里,等这部分打完清稿,就已经是一千来页,甚至更多——已经够多的
了!一天的难处一天当就够了……

\pmark\origpage{1249}这种赶紧收场、“把它寄出去”的急切,无疑首先是一匹嗅到火药味的战马的急切,急不可
耐地要冲进战阵\fnote{1049}{(*)}{类似的心绪,在“愚妄与更新”最后一节“一段过去的重量”(n\textsuperscript{o}~50)里已经谈过,只是视角稍有不
同:在那里,冲进战阵的不是战马,而是一头公牛,追赶着一块在它“鼻子跟前晃来晃去”的红布……}。但也许还有更深一层的,是一种愿望:巴望某一段过去从此脱开我而去。这一千来页,以惊
人的方式把这段过去的全部\textbf{重量}凝成了实体——把这项工作做到终了,直到最后一桩庶务
(最后一桩,想必就是把《收获与播种》寄给我寄发清单上预先开列的一百三十位收件人……\fnote{1050}{(**)}{到那时,那著名的“重量”将变得更加“惊人”:一上来就是二十万页(200×1000),而不是一千页!}),在我
看来,几乎出于本能,也正是我把它\textbf{抛掉}的时刻。幻觉?唯有将来知道……

就这样,我来到了那些“终止和弦”面前——在那个有名的“句号”之前;那个“句号”,一年多
来我一直以为它就在眼前,却一天天、一星期一星期、一个月一个月地被推远,总有料想不到的
东西涌来,要占一席之地。

还剩什么要说,在这些终止和弦里?有一份感激,它要借“感谢”说出。这场思考是孤独的果实,
然而我得到过许多方式的帮助。

最显而易见的帮助,来自Zoghman~Mebkhout,同样是以许多方式:凭他的耐心,他领我“入了门”,
走进“好上帝-Mebkhout定理”周身的哲学;凭他对我的信任,不管不顾地把他在我昔日学生们那里
遭遇到的困难与挫折向我和盘托出;凭他帮我厘清一片枝蔓丛生的数学文献——我早已同它断了
往来;最后,凭他自打得知这项工作那一刻起,就对它投以无私而友善的关注——他在其中看见
我倾注了身心,而首先(我相信),他认出并接纳了其中的\textbf{见证}。

我同样感激Pierre~Deligne,他曾专程前来看我,并(去年10月)读阅了当时已写出的\emph{葬礼}
部分,向我提出意见\dmark{}\origpage{1250}\fnote{1051}{(*)}{关于这次造访以及Deligne带给我的种种澄清,见“葬礼(III)”中构成“最后的义务(或造访)”一节的那两条注(n\textsuperscript{o}s~163、164)。}。这次造访对我多有助益。

最后,帮助过我的,还有朗格多克科学技术大学那些为我打字的秘书们的好意与友善氛围:Boulet
小姐(Mlle~Boulet)、Boucher夫人、Brun小姐、Cellier夫人、
Lacan小姐、Mori夫人。她们中有两位,不惜动用个人时间,为按时交
件承担了一部分打字工作,还不肯为此收取报酬——这份情义令我深深感动。此外,整部《收获
与播种》笔记的后一半,是Lacan小姐一个人打出来的,其用心与效率堪称典范。对她们每一位,
我都乐于在此表达我的感激。

我还想着所有那些人,所有那些女人和男人——在我工作的许多时刻,他们曾以常常是不受欢迎
的方式\fnote{1052}{(**)}{这几个月的笔记里,这里那里曾提到这类“打扰”。关于这一点,特别见注“信使(2)”(n\textsuperscript{o}~181)。}打断我的工作与安宁。当然,这些“打扰”本身——它们在某些时刻考验过我,有一些至今还在
我心里留着悲哀的残余——在我的工作里同样有它们的角色要扮演,同样带来一个讯息,听不听、
化不化,全在我自己。当悲哀或怨怼化为感激,我便会知道,这个讯息已被接纳……

\subsection{(2) 朋友}

\notehead{188}{}

\emph{葬礼}的这些终极和弦,近一年来早已名字现成:\emph{自深处}!(De~Profundis)在《引论》
(I~7,“葬仪的安排”)里,我甚至走得更远,竟(也许失之轻率……)预告说,死者那“完全的满
足”将构成“那场难忘葬礼的终曲与最后的和弦”。写下这句预言时(仿佛事已办讫),我是可以
原谅的——写那几行的时候(去年5月),这看上去确实是一个为期极近的预言,因为我自以为正
要抵达“自深处”的那些终止和弦。

诚然,比去年更切身地(当时,思考的“第二口气”行将告一段落),我意识到\dmark{}\origpage{1251}自己离真正“走
完”这葬礼还有多远——若把纯然物质性的事实除外的话(那些事实,我自觉得已经“尽在掌握”,
颇有几分自负\fnote{1053}{(*)}{(5月10日)然而,这几行写下之后,又有一个多月过去,耗在把新冒出来的事实勉力“塞进去”,临
时添了足足二十来条子注!})。倘若确实像我不时感到的那样,理解\emph{葬礼},也就意味着“理解冲
突”,那么很可能,我余生的岁月将不足以完成这番“走完”——至少走不到深处。

如此说来,我可以说,我此刻写下这条终极之注时的心境,已全然不同于写《\emph{葬礼}引论》时
的心境。这是否意味着,我是在“完全的满足”之情不在场的情况下,为这场思考收尾的?

我不以为然。愿景一旦深化,某项曾催生这愿景、为这深化作过铺垫的工作——它本可以被看
作“已臻完成”——便会在一个“彼岸”显现之时,显出自身的\textbf{未完成}。然而,工作的\textbf{意
义},以及它让我们尝到的满足或不满足,并不在它的完成之中,也不取决于这项工作注定完
成与否。工作的意义就在工作本身,就在\textbf{当下此刻}——在我们做事时的心境里,在我们倾
注其中的爱里(或爱的缺席里……)——而不在一个我们无从企及的假想的将来。

去年3月,甚至在我发现\emph{葬礼}之前,我在《引论》(I~1,“梦与实现”,打字稿第~iv~页)里写道:

\begin{quote}
“……我离开这项工作,带着一个知道已把工作做到终点的人的完全满足。没有任何一件事,
无论它多么‘微小’,被我避而不谈;没有任何一件事,是我念念于心要说、却终究没有说的,
以至于在此刻,会在我心里留下一星不满、一星遗憾的残余,无论它们多么‘微小’。”
\end{quote}

如今我知道,这份我原以为“做到终点”的工作,今天仍未到终点,也许永远到不了。但我也知
道,这毕竟是一件次要的事。那份“完全的满足”——在我写下这些力图贴它最近的话语的那一刻,
我曾强烈地感到过它——它一路伴随我写完了\dmark{}\origpage{1252}《收获与播种》。它是我的一位老朋友,早在
我数学生涯的全部岁月里,它就陪伴着我,低声告诉我:路走对了。后来,在冥想的功课里,我
又找到了它——正是同一位。

当我听不见它,工作便失去意义。正因如此,它的声音于我弥足珍贵,也正因如此,我在工作时
小心在意,从不让自己远离它。靠了这个,工作在我一生中始终是喜悦的源泉——那个全然以身
相许者的“完全的满足”。

在即将结束的这项工作里,情形并无二致——这项工作,即是“\textbf{收获}”,同时又是
“\textbf{播种}”。


\end{document}
